\documentclass[12pt,a4paper]{article}
\usepackage[utf8]{inputenc}
\usepackage[T1]{fontenc}
\usepackage[italian]{babel}
\usepackage{amsmath,amssymb,amsfonts,amsthm}
\usepackage{mathtools}
\usepackage{geometry}
\usepackage{setspace}
\onehalfspacing
\usepackage{booktabs}
\usepackage{enumitem}
\usepackage{hyperref}
\usepackage{xcolor}
\usepackage{tcolorbox}
\usepackage{fancyhdr}
\usepackage{caption}
\usepackage{graphicx}
\usepackage{float}
\usepackage{tikz}
\usepackage{pgfplots}
\pgfplotsset{compat=1.18}
\usetikzlibrary{positioning, arrows.meta, shapes.geometric, calc, decorations.pathreplacing, patterns, fit, backgrounds}
\usepgfplotslibrary{fillbetween}
\geometry{margin=2.5cm, headheight=14.5pt}
\hypersetup{colorlinks=true, linkcolor=blue!60!black, urlcolor=blue!60!black, bookmarks=true, bookmarksnumbered=true, bookmarksopen=true, bookmarksopenlevel=2}
\pagestyle{fancy}
\fancyhf{}
\fancyhead[L]{\small Deep Learning -- Preparazione Esame}
\fancyhead[R]{\small Paper: Perceiver \& Perceiver IO}
\fancyfoot[C]{\thepage}
\tcbuselibrary{skins,breakable}
\newtcolorbox{formulabox}[1][]{colback=blue!2!white, colframe=blue!50!black, title=#1, fonttitle=\bfseries, breakable}
\newtcolorbox{notabox}[1][]{colback=yellow!2!white, colframe=orange!50!black, title=#1, fonttitle=\bfseries, breakable}
\newtcolorbox{proofbox}[1][]{colback=black!2!white, colframe=black!30!white, title=#1, fonttitle=\itshape, breakable}
\newtcolorbox{examplebox}[1][]{colback=green!2!white, colframe=green!50!black, title=#1, fonttitle=\bfseries, breakable}
\newtheorem{theorem}{Teorema}[section]
\newtheorem{proposition}[theorem]{Proposizione}
\newtheorem{lemma}[theorem]{Lemma}
\newtheorem{definition}[theorem]{Definizione}
\newtheorem{corollary}[theorem]{Corollario}

% Pipeline bar for forward/backward pass steps
\newtcolorbox{pipelinebox}[1][]{
    colback=gray!5!white,
    colframe=gray!60!black,
    title={\small Pipeline --- Passo corrente evidenziato},
    fonttitle=\bfseries,
    breakable,
    left=4pt, right=4pt, top=2pt, bottom=2pt,
    #1
}

\title{\Huge\textbf{Perceiver e Perceiver IO}\\\vspace{0.5cm}\LARGE Jaegle et al.\ (ICML 2021 \& ICLR 2022)\\\vspace{0.3cm}\large Dispensa completa per la preparazione dell'esame}
\author{Dispensa preparazione esame}
\date{}

\begin{document}
\maketitle

\section*{Riferimenti bibliografici}
\begin{itemize}[leftmargin=*]
    \item \textbf{Jaegle et al.}\ (2021): \textit{Perceiver: General Perception with Iterative Attention}, ICML 2021
    \item \textbf{Vaswani et al.}\ (2017): \textit{Attention Is All You Need}
    \item \textbf{Dosovitskiy et al.}\ (2021): \textit{An Image is Worth 16x16 Words} (ViT)
    \item \textbf{Tancik et al.}\ (2020): \textit{Fourier Features Let Networks Learn High Frequency Functions in Low Dimensional Domains}
    \item \textbf{He et al.}\ (2016): \textit{Deep Residual Learning for Image Recognition} (ResNet)
    \item \textbf{Bahdanau, Cho \& Bengio} (2015): \textit{Neural Machine Translation by Jointly Learning to Align and Translate}
    \item \textbf{Hendrycks \& Gimpel} (2016): \textit{Gaussian Error Linear Units (GELUs)}
    \item \textbf{You et al.}\ (2020): \textit{Large Batch Optimization for Deep Learning: Training BERT in 76 Minutes} (LAMB)
    \item \textbf{Kingma \& Ba} (2015): \textit{Adam: A Method for Stochastic Optimization}
    \item \textbf{Jaegle et al.}\ (2022): \textit{Perceiver IO: A General Architecture for Structured Inputs and Outputs}, ICLR 2022
    \item \textbf{Devlin et al.}\ (2019): \textit{BERT: Pre-training of Deep Bidirectional Transformers for Language Understanding}
    \item \textbf{Teed \& Deng}\ (2020): \textit{RAFT: Recurrent All-Pairs Field Transforms for Optical Flow}
    \item \textbf{Vinyals et al.}\ (2019): \textit{Grandmaster Level in StarCraft II Using Multi-Agent Reinforcement Learning} (AlphaStar)
    \item \textbf{Brock et al.}\ (2021): \textit{High-Performance Large-Scale Image Recognition Without Normalization} (NFNet)
    \item \textbf{BB24} (Bishop \& Bishop, \textit{Deep Learning -- Foundations and Concepts}, 2024): Sezioni 12.1--12.2
\end{itemize}

\tableofcontents
\newpage


%% ============================================================
%%  SEZIONE 1 --- INTRODUZIONE E MOTIVAZIONE
%% ============================================================
\section{Introduzione e Motivazione} \label{sec:intro}

\subsection{Il problema delle architetture specializzate} \label{sec:specializzate}

Immaginiamo di voler costruire un sistema di intelligenza artificiale capace di ``percepire'' il mondo: deve saper riconoscere immagini, capire l'audio, analizzare video e persino interpretare nuvole di punti 3D. La domanda che ci si pone \`e: \textit{servono tanti cervelli diversi, uno per ogni senso, o ne basta uno solo?}

Fino al 2021, la risposta della comunit\`a di ricerca era chiara: servono cervelli diversi. Questa situazione si chiama \textbf{frammentazione architetturale}: per ogni modalit\`a percettiva --- immagini, audio, video, nuvole di punti 3D, testo --- si \`e sviluppata un'architettura altamente specializzata.

Le \textbf{reti convoluzionali} (CNN, si veda l'\hyperref[app:cnn]{Appendice~F}) dominano la visione artificiale sfruttando la struttura a griglia bidimensionale delle immagini: applicano filtri locali (kernel) che scorrono sull'immagine, condividendo i pesi e catturando pattern locali come bordi e texture. Le \textbf{reti ricorrenti} (RNN, LSTM, GRU --- si vedano le \hyperref[app:rnn]{Appendici~C}, \hyperref[app:lstm]{D} e \hyperref[app:gru]{E}) sono state progettate per le sequenze temporali, processando un elemento alla volta e mantenendo uno stato nascosto che accumula informazione nel tempo. Architetture come \textbf{PointNet/PointNet++} sono specifiche per le nuvole di punti 3D, sfruttando la geometria locale dei punti nello spazio tridimensionale. I \textbf{Transformer} (Vaswani et al., 2017 --- si veda l'\hyperref[app:transformer]{Appendice~H}) hanno mostrato una notevole versatilit\`a, ma rimangono legati ad assunzioni strutturali, in particolare l'uso di \emph{tokenizzazione} specifica per dominio (patch per le immagini, token per il testo).

Perch\'e tutto questo \`e un problema? Perch\'e questa specializzazione presenta diversi svantaggi:
\begin{enumerate}[label=(\roman*)]
    \item \textbf{Costo ingegneristico}: ogni nuova modalit\`a richiede la progettazione, l'implementazione e l'ottimizzazione di una nuova architettura. In pratica, quello che succede \`e che ogni volta che emerge un nuovo tipo di dato, la comunit\`a deve ``reinventare la ruota''.
    \item \textbf{Mancanza di trasferibilit\`a}: le intuizioni architetturali sviluppate per una modalit\`a non si trasferiscono facilmente ad altre.
    \item \textbf{Difficolt\`a nella fusione multi-modale}: combinare informazioni da sorgenti diverse (es.\ audio + video) richiede moduli di fusione \emph{ad hoc} e non banali.
    \item \textbf{Assunzioni implicite}: le CNN assumono localit\`a spaziale e invarianza traslazionale; le RNN assumono un ordinamento sequenziale; tali \emph{inductive bias} possono essere vantaggiosi ma anche limitanti.
\end{enumerate}

\begin{notabox}[Analogia con la percezione biologica]
Il sistema percettivo umano elabora stimoli visivi, uditivi, tattili e olfattivi convergendo in rappresentazioni integrate nella corteccia associativa. Non esiste un ``cervello separato'' per ogni senso. Il Perceiver si ispira a questa idea: un'unica architettura capace di gestire qualsiasi tipo di input sensoriale.
\end{notabox}

\begin{notabox}[In parole semplici: il problema delle architetture specializzate]
Detto in parole povere: \`e come se per leggere un libro in inglese servisse un paio di occhiali diverso da quello per leggere un libro in francese. Ogni lingua ha il suo paio di occhiali. Se arriva una nuova lingua, devi fabbricare un nuovo paio. Il Perceiver vuole creare un unico paio di occhiali universale che funziona con qualsiasi lingua --- anzi, con qualsiasi tipo di informazione.
\end{notabox}

\subsection{Il limite computazionale dei Transformer} \label{sec:limite_transformer}

Facciamo un passo indietro e chiediamoci: perch\'e non usare semplicemente un Transformer per tutto? In fondo, il Transformer (Vaswani et al., 2017) \`e gi\`a molto versatile. Il problema \`e un limite fondamentale: la complessit\`a computazionale. Vediamo cosa succede concretamente.

Il Transformer standard utilizza la \textbf{self-attention} su una sequenza di $M$ elementi. La complessit\`a computazionale e di memoria della self-attention \`e:

\begin{formulabox}[Complessit\`a della self-attention standard]
\[
    \text{Self-Attention:} \quad \mathcal{O}(M^2 \cdot d)
\]
dove $M$ \`e il numero di token (lunghezza della sequenza) e $d$ \`e la dimensione del modello. Lo spazio di memoria per la matrice di attenzione \`e $\mathcal{O}(M^2)$.
\end{formulabox}

In pratica, quello che succede \`e che raddoppiando la lunghezza dell'input, il costo della self-attention \textbf{quadruplica}. Questa crescita quadratica pone un limite fondamentale all'applicabilit\`a dei Transformer a input di grandi dimensioni (per l'analisi numerica dettagliata si veda la \hyperref[sec:complessita]{Sezione~\ref*{sec:complessita}}). Per capire quanto sia grave, vediamo i numeri concreti:

\begin{examplebox}[Dimensioni tipiche degli input percettivi]
\begin{itemize}
    \item \textbf{Immagine ImageNet} ($224 \times 224$): $M = 50{,}176$ pixel. La matrice di attenzione avrebbe $50{,}176^2 \approx 2.5 \times 10^9$ elementi --- \emph{impossibile} da gestire in memoria anche su GPU moderne.
    \item \textbf{Audio} (1 secondo a 48kHz): $M = 48{,}000$ campioni.
    \item \textbf{Video} (1 secondo, $224 \times 224$ a 30fps): $M \approx 1{,}505{,}280$ pixel.
    \item \textbf{Point cloud} (ModelNet40): $M \approx 2{,}048$ punti (relativamente piccolo).
\end{itemize}
\end{examplebox}

Per ovviare a questo problema, il \textbf{Vision Transformer (ViT)} (Dosovitskiy et al., 2021 --- si veda l'\hyperref[app:vit]{Appendice~I}) suddivide l'immagine in \emph{patch} di $16 \times 16$ pixel, riducendo la sequenza da $50{,}176$ a $196$ token. Tuttavia, questa tokenizzazione:
\begin{itemize}
    \item Introduce un \emph{inductive bias} specifico per le immagini (griglia regolare di patch).
    \item Perde informazione a livello sub-patch.
    \item Non \`e immediatamente generalizzabile ad audio, video o point cloud.
\end{itemize}

% TikZ: Complexity comparison plot (log scale)
\begin{figure}[H]
\centering
\begin{tikzpicture}
\begin{axis}[
    width=13cm, height=7cm,
    xlabel={Numero di elementi input $M$},
    ylabel={Operazioni (scala log)},
    xmin=100, xmax=60000,
    ymin=1e4, ymax=1e11,
    xmode=log, ymode=log,
    legend pos=north west,
    grid=major,
    grid style={gray!30},
    thick
]
% Transformer O(M^2) with d=512
\addplot[red, very thick, domain=100:60000, samples=50] {x*x};
\addlegendentry{Transformer $\mathcal{O}(M^2)$}
% Perceiver O(MN) with N=512
\addplot[blue, very thick, domain=100:60000, samples=50] {x*512};
\addlegendentry{Perceiver $\mathcal{O}(MN)$, $N{=}512$}
% Linear O(M)
\addplot[green!60!black, very thick, dashed, domain=100:60000, samples=50] {x*9};
\addlegendentry{CNN locale $\mathcal{O}(Mk^2)$, $k{=}3$}
% Vertical line at ImageNet
\draw[gray, dashed, thick] (axis cs:50176,1e4) -- (axis cs:50176,1e11);
\node[font=\scriptsize, rotate=90, gray] at (axis cs:42000,1e8) {ImageNet $M{=}50{,}176$};
\end{axis}
\end{tikzpicture}
\caption{Confronto della complessit\`a computazionale in scala logaritmica. Il Transformer standard (rosso) cresce quadraticamente, il Perceiver (blu) linearmente in $M$, la CNN locale (verde) ancora pi\`u economicamente ma senza dipendenze globali.}
\label{fig:complexity_comparison}
\end{figure}


\subsection{L'idea chiave: il latent bottleneck} \label{sec:latent_bottleneck}

Ora che abbiamo capito il problema (la self-attention \`e troppo costosa per input grandi), vediamo la soluzione proposta dal Perceiver. L'idea \`e semplice ma potente.

Immaginiamo di avere 50.000 libri da leggere (i pixel). Invece di far parlare ogni libro con ogni altro libro ($50.000^2$ conversazioni!), nominiamo 512 ``rappresentanti'' che leggono i libri e poi discutono tra di loro. Le conversazioni diventano $50.000 \times 512$ (leggere) + $512^2$ (discutere) --- enormemente meno di $50.000^2$.

In pratica, il Perceiver propone una soluzione elegante al problema della complessit\`a quadratica tramite l'introduzione di un \textbf{array latente} di dimensione controllata. L'intuizione formale \`e la seguente:

\begin{formulabox}[Idea centrale del Perceiver]
Invece di applicare la self-attention direttamente sull'input di $M$ elementi, si introduce un array latente $\mathbf{z} \in \mathbb{R}^{N \times D}$ con $N \ll M$. L'informazione dall'input viene estratta tramite \textbf{cross-attention} (costo $\mathcal{O}(MN)$) e poi elaborata tramite \textbf{self-attention} solo tra i latenti (costo $\mathcal{O}(N^2)$).
\[
    \underbrace{\mathcal{O}(M^2)}_{\text{Transformer standard}} \quad \longrightarrow \quad \underbrace{\mathcal{O}(MN) + \mathcal{O}(N^2)}_{\text{Perceiver, con } N \ll M}
\]
\end{formulabox}

Questo \emph{latent bottleneck} (collo di bottiglia latente) disaccoppia la dimensionalit\`a dell'elaborazione dalla dimensionalit\`a dell'input, permettendo di processare input arbitrariamente grandi a costo computazionale controllato.

\subsection{Obiettivi del paper} \label{sec:obiettivi}

Gli autori si propongono di dimostrare che:
\begin{enumerate}[label=\arabic*.]
    \item Un'unica architettura pu\`o raggiungere risultati competitivi su \textbf{immagini, audio, video e point cloud} senza modifiche strutturali specifiche per ogni dominio.
    \item L'architettura \`e \textbf{scalabile} rispetto alla dimensione dell'input, con complessit\`a lineare in $M$.
    \item L'architettura non fa \textbf{assunzioni sulla struttura spaziale} dell'input (provato dall'esperimento su ImageNet permutato, \hyperref[sec:imagenet_perm]{Sezione~10.2}).
    \item Il \textbf{weight sharing} (condivisione dei pesi) tra iterazioni agisce come regolarizzatore efficace (\hyperref[sec:iterative]{Sezione~3.4}).
\end{enumerate}

\bigskip

% TikZ: Specialized architectures vs Perceiver
\begin{figure}[H]
\centering
\begin{tikzpicture}[
    modbox/.style={draw, rounded corners, minimum width=2.3cm, minimum height=0.9cm, font=\small\bfseries, text=white},
    archbox/.style={draw, rounded corners, minimum width=2.3cm, minimum height=0.9cm, font=\small, fill=gray!10},
    arrowstyle/.style={-{Stealth[length=3mm]}, thick},
    >=Stealth
]

% Left side: Specialized
\node[font=\large\bfseries] at (-4.5, 4.2) {Approccio tradizionale};

\node[modbox, fill=red!70!black] (img) at (-6, 3) {Immagini};
\node[modbox, fill=blue!70!black] (aud) at (-6, 1.8) {Audio};
\node[modbox, fill=green!60!black] (vid) at (-6, 0.6) {Video};
\node[modbox, fill=purple!70!black] (pc) at (-6, -0.6) {Point Cloud};

\node[archbox] (cnn) at (-2.5, 3) {CNN};
\node[archbox] (rnn) at (-2.5, 1.8) {RNN / CNN-1D};
\node[archbox] (vidarch) at (-2.5, 0.6) {3D CNN};
\node[archbox] (pn) at (-2.5, -0.6) {PointNet++};

\draw[arrowstyle] (img) -- (cnn);
\draw[arrowstyle] (aud) -- (rnn);
\draw[arrowstyle] (vid) -- (vidarch);
\draw[arrowstyle] (pc) -- (pn);

% Separator
\draw[dashed, thick, gray] (0.2, 4.5) -- (0.2, -1.5);

% Right side: Perceiver
\node[font=\large\bfseries] at (4.5, 4.2) {Approccio Perceiver};

\node[modbox, fill=red!70!black] (img2) at (2.5, 3) {Immagini};
\node[modbox, fill=blue!70!black] (aud2) at (2.5, 1.8) {Audio};
\node[modbox, fill=green!60!black] (vid2) at (2.5, 0.6) {Video};
\node[modbox, fill=purple!70!black] (pc2) at (2.5, -0.6) {Point Cloud};

\node[archbox, fill=orange!20, draw=orange!70!black, minimum width=3cm, minimum height=2.5cm, font=\large\bfseries] (perc) at (6, 1.2) {Perceiver};

\draw[arrowstyle, red!70!black] (img2) -- (perc.west |- img2);
\draw[arrowstyle, blue!70!black] (aud2) -- (perc.west |- aud2);
\draw[arrowstyle, green!60!black] (vid2) -- (perc.west |- vid2);
\draw[arrowstyle, purple!70!black] (pc2) -- (perc.west |- pc2);

\end{tikzpicture}
\caption{Confronto tra l'approccio tradizionale (un'architettura specializzata per ogni modalit\`a) e il Perceiver (un'unica architettura per tutte le modalit\`a).}
\label{fig:specialized_vs_perceiver}
\end{figure}


% TikZ: CNN vs Transformer vs Perceiver comparison
\begin{figure}[H]
\centering
\begin{tikzpicture}[
    colhead/.style={font=\small\bfseries, minimum height=0.8cm},
    imgbox/.style={draw, thick, minimum width=2cm, minimum height=2cm},
    procbox/.style={draw, rounded corners=3pt, minimum width=2.5cm, minimum height=0.6cm, font=\scriptsize},
    arrowstyle/.style={-{Stealth[length=2mm]}, thick},
    >=Stealth,
    scale=0.82, transform shape
]

% Column 1: CNN
\node[colhead] at (0, 5.5) {CNN (ResNet)};
\node[imgbox, fill=red!10] (img1) at (0, 4) {};
\node[font=\tiny] at (0, 4) {Immagine};
% Small kernel overlay
\draw[blue, very thick] (-0.4, 3.6) rectangle (0.0, 4.0);
\node[font=\tiny, blue] at (-0.2, 3.3) {kernel $3\times 3$};
\draw[arrowstyle] (0, 2.8) -- (0, 2.2);
\node[procbox, fill=red!15] at (0, 1.8) {Filtri locali $\to$ Pooling};
\draw[arrowstyle] (0, 1.4) -- (0, 0.8);
\node[procbox, fill=red!15] at (0, 0.4) {Feature gerarchiche};
\draw[arrowstyle] (0, 0.0) -- (0, -0.6);
\node[procbox, fill=green!15] at (0, -1.0) {Classificazione};
\node[font=\tiny, gray, text width=2.5cm, align=center] at (0, -1.8) {Locale $\to$ Globale\\Solo griglie 2D};

% Column 2: Transformer (ViT)
\node[colhead] at (5, 5.5) {Transformer (ViT)};
\node[imgbox, fill=orange!10] (img2) at (5, 4) {};
\node[font=\tiny] at (5, 4) {Immagine};
% Patch grid
\draw[orange, thick] (4.3, 3.3) -- (4.3, 4.7);
\draw[orange, thick] (5.0, 3.3) -- (5.0, 4.7);
\draw[orange, thick] (5.7, 3.3) -- (5.7, 4.7);
\draw[orange, thick] (3.9, 3.7) -- (6.1, 3.7);
\draw[orange, thick] (3.9, 4.3) -- (6.1, 4.3);
\node[font=\tiny, orange] at (5, 3.0) {patch $16\times 16$};
\draw[arrowstyle] (5, 2.8) -- (5, 2.2);
\node[procbox, fill=orange!15] at (5, 1.8) {Self-Att. tra patch};
\node[font=\tiny, gray] at (5, 1.2) {$196 \times 196$};
\draw[arrowstyle] (5, 0.8) -- (5, 0.2);
\node[procbox, fill=orange!15] at (5, -0.2) {Token CLS};
\draw[arrowstyle] (5, -0.6) -- (5, -0.8);
\node[procbox, fill=green!15] at (5, -1.0) {Classificazione};
\node[font=\tiny, gray, text width=2.5cm, align=center] at (5, -1.8) {Globale da subito\\Ma perde dettaglio sub-patch};

% Column 3: Perceiver
\node[colhead] at (10, 5.5) {Perceiver};
\node[imgbox, fill=blue!10] (img3) at (10, 4) {};
\node[font=\tiny] at (10, 4) {Immagine};
% Individual pixels
\foreach \x in {9.2, 9.5, 9.8, 10.1, 10.4, 10.7} {
    \fill[blue!40] (\x, 4.3) circle (1.5pt);
    \fill[blue!40] (\x, 4.0) circle (1.5pt);
    \fill[blue!40] (\x, 3.7) circle (1.5pt);
}
\node[font=\tiny, blue] at (10, 3.0) {pixel singoli + Fourier};
\draw[arrowstyle] (10, 2.8) -- (10, 2.2);
\node[procbox, fill=blue!15] at (10, 1.8) {Cross-Att. $\to$ Latenti};
\node[font=\tiny, gray] at (10, 1.2) {$512 \times 50176$};
\draw[arrowstyle] (10, 0.8) -- (10, 0.2);
\node[procbox, fill=blue!15] at (10, -0.2) {Self-Att. latenti + Avg Pool};
\draw[arrowstyle] (10, -0.6) -- (10, -0.8);
\node[procbox, fill=green!15] at (10, -1.0) {Classificazione};
\node[font=\tiny, gray, text width=2.8cm, align=center] at (10, -1.8) {Globale da subito\\Opera su TUTTI i pixel\\Qualsiasi modalit\`a};

\end{tikzpicture}
\caption{Confronto tra CNN, Transformer (ViT) e Perceiver nel processamento di un'immagine. La CNN usa filtri locali su una griglia; il ViT divide in patch e applica self-attention; il Perceiver opera direttamente su tutti i pixel tramite cross-attention con un array latente compatto.}
\label{fig:cnn_vs_transformer_vs_perceiver}
\end{figure}

\begin{figure}[H]
\centering
\includegraphics[width=0.95\textwidth]{paper_figures/fig2_modalities.png}
\caption{Il Perceiver addestrato su diverse modalit\`a (dal paper originale, Figure 2): immagini da ImageNet (sinistra), video e audio da AudioSet (centro), e point cloud 3D da ModelNet40 (destra). La stessa architettura viene usata per tutti i domini, senza alcuna modifica strutturale.}
\label{fig:paper_modalities}
\end{figure}



\newpage
%% ============================================================
%%  SEZIONE 2 --- PREREQUISITI FONDAMENTALI
%% ============================================================
\section{Prerequisiti Fondamentali} \label{sec:prerequisiti}

Prima di tuffarci nell'architettura del Perceiver, dobbiamo assicurarci di avere tutte le basi necessarie. Questa sezione spiega in modo autocontenuto tutti i concetti che servono: dalla softmax all'attenzione, dalla normalizzazione alle connessioni residue. Se avete gi\`a familiarit\`a con questi concetti, potete sfogliarla rapidamente; se no, leggetela con attenzione, perch\'e ogni concetto verr\`a usato ampiamente nelle sezioni successive.

Per una trattazione pi\`u approfondita di ciascun argomento, si rimanda alle appendici: \hyperref[app:perceptron]{A (Perceptrone)}, \hyperref[app:feedforward]{B (MLP)}, \hyperref[app:rnn]{C (RNN)}, \hyperref[app:lstm]{D (LSTM)}, \hyperref[app:gru]{E (GRU)}, \hyperref[app:cnn]{F (CNN)}, \hyperref[app:resnet]{G (ResNet)}, \hyperref[app:transformer]{H (Transformer)}, \hyperref[app:vit]{I (ViT)}.

\subsection{La funzione Softmax} \label{sec:softmax}

Iniziamo dal concetto pi\`u basilare ma fondamentale. Immaginiamo di avere un modello che, guardando un'immagine, produce dei ``punteggi'' per ogni classe: 2.0 per ``gatto'', 1.0 per ``cane'', 0.1 per ``uccello''. Questi numeri non hanno un significato immediato come probabilit\`a --- possono essere negativi, possono essere molto grandi, e non sommano a 1. Come li trasformiamo in qualcosa di interpretabile?

La risposta \`e la \textbf{funzione softmax}: una funzione che trasforma un vettore di numeri reali arbitrari in una distribuzione di probabilit\`a (valori positivi che sommano a 1).

\begin{formulabox}[Softmax: definizione e propriet\`a]
Dato un vettore $\mathbf{s} = (s_1, s_2, \ldots, s_n) \in \mathbb{R}^n$:
\[
    \text{softmax}(s_i) = \frac{e^{s_i}}{\sum_{j=1}^{n} e^{s_j}} \qquad \text{per } i = 1, \ldots, n
\]

\textbf{Propriet\`a}:
\begin{itemize}
    \item Tutti gli output sono \textbf{positivi}: $\text{softmax}(s_i) > 0$ per ogni $i$.
    \item \textbf{Somma a 1}: $\sum_{i=1}^n \text{softmax}(s_i) = 1$.
    \item \`E una \textbf{generalizzazione della funzione logistica} al caso multi-classe.
    \item Se un $s_i$ \`e molto pi\`u grande degli altri, $\text{softmax}(s_i) \approx 1$ (``winner takes all'').
    \item \`E \textbf{invariante per traslazione}: $\text{softmax}(\mathbf{s} + c) = \text{softmax}(\mathbf{s})$ per ogni costante $c$.
\end{itemize}
\end{formulabox}

\begin{examplebox}[Esempio numerico della softmax]
Dato $\mathbf{s} = (2.0, 1.0, 0.1)$:
\begin{align*}
    e^{2.0} &= 7.389, \quad e^{1.0} = 2.718, \quad e^{0.1} = 1.105 \\
    \text{somma} &= 7.389 + 2.718 + 1.105 = 11.212 \\
    \text{softmax}(\mathbf{s}) &= \left(\frac{7.389}{11.212},\; \frac{2.718}{11.212},\; \frac{1.105}{11.212}\right) = (0.659, 0.242, 0.099)
\end{align*}
La prima componente, essendo la pi\`u grande, riceve il peso maggiore (65.9\%).
\end{examplebox}

% TikZ: Softmax visualization - bar chart
\begin{figure}[H]
\centering
\begin{tikzpicture}[scale=0.85, transform shape]

% Left side: raw scores
\node[font=\small\bfseries] at (1.5, 4.2) {Score grezzi (logit)};
\fill[red!60] (0, 0) rectangle (1, 2.0);
\fill[orange!60] (1.3, 0) rectangle (2.3, 1.0);
\fill[yellow!60] (2.6, 0) rectangle (3.6, 0.1);
\node[font=\scriptsize] at (0.5, -0.3) {2.0};
\node[font=\scriptsize] at (1.8, -0.3) {1.0};
\node[font=\scriptsize] at (3.1, -0.3) {0.1};
\node[font=\tiny] at (0.5, 2.3) {gatto};
\node[font=\tiny] at (1.8, 1.3) {cane};
\node[font=\tiny] at (3.1, 0.4) {uccello};
\draw[thick] (0, 0) -- (3.6, 0);

% Arrow
\draw[-{Stealth[length=3mm]}, very thick, blue!60!black] (4.2, 1.5) -- (5.8, 1.5);
\node[font=\small\bfseries, blue!60!black] at (5.0, 2.0) {softmax};

% Right side: probabilities
\node[font=\small\bfseries] at (8.5, 4.2) {Probabilit\`a (sommano a 1)};
\fill[red!60] (7, 0) rectangle (8, 3.3);
\fill[orange!60] (8.3, 0) rectangle (9.3, 1.21);
\fill[yellow!60] (9.6, 0) rectangle (10.6, 0.495);
\node[font=\scriptsize\bfseries] at (7.5, 3.5) {65.9\%};
\node[font=\scriptsize] at (8.8, 1.5) {24.2\%};
\node[font=\scriptsize] at (10.1, 0.8) {9.9\%};
\node[font=\tiny] at (7.5, -0.3) {gatto};
\node[font=\tiny] at (8.8, -0.3) {cane};
\node[font=\tiny] at (10.1, -0.3) {uccello};
\draw[thick] (7, 0) -- (10.6, 0);

% Sum annotation
\draw[decorate, decoration={brace, amplitude=4pt}, thick, green!60!black]
    (10.8, 0) -- (10.8, 3.3) node[midway, right=5pt, font=\scriptsize, green!60!black] {$= 100\%$};

\end{tikzpicture}
\caption{Visualizzazione della softmax: i punteggi grezzi (logit) vengono trasformati in probabilit\`a che sommano a 100\%. Il valore pi\`u alto (2.0) riceve la fetta maggiore (65.9\%), ma anche gli altri contribuiscono.}
\label{fig:softmax_visualization}
\end{figure}

\begin{notabox}[In parole semplici: la softmax]
Detto in parole povere, la softmax \`e come trasformare i voti di un esame in percentuali che sommano a 100\%. Se uno studente prende 8, un altro 6 e un altro 5, la softmax direbbe: il primo ha il 65\% del ``merito'', il secondo il 24\%, il terzo il 10\%. La softmax ``amplifica'' le differenze: anche se i voti erano relativamente vicini, il migliore prende una fetta molto pi\`u grande.
\end{notabox}

\textbf{Derivata della softmax} (necessaria per la \hyperref[sec:backward]{backpropagation}):
\[
    \frac{\partial \text{softmax}(s_i)}{\partial s_j} = \text{softmax}(s_i)\left(\delta_{ij} - \text{softmax}(s_j)\right)
\]
dove $\delta_{ij}$ \`e il delta di Kronecker ($\delta_{ij} = 1$ se $i = j$, 0 altrimenti).

\begin{proofbox}[Dimostrazione della derivata della softmax]
Sia $p_i = \frac{e^{s_i}}{\sum_k e^{s_k}}$. Calcoliamo $\frac{\partial p_i}{\partial s_j}$.

\textbf{Caso $i = j$} (derivata rispetto al proprio score):
\begin{align*}
    \frac{\partial p_i}{\partial s_i} &= \frac{e^{s_i} \cdot \sum_k e^{s_k} - e^{s_i} \cdot e^{s_i}}{(\sum_k e^{s_k})^2} = \frac{e^{s_i}}{\sum_k e^{s_k}} - \frac{e^{2s_i}}{(\sum_k e^{s_k})^2} \\
    &= p_i - p_i^2 = p_i(1 - p_i)
\end{align*}

\textbf{Caso $i \neq j$} (derivata rispetto allo score di un altro elemento):
\begin{align*}
    \frac{\partial p_i}{\partial s_j} &= \frac{0 - e^{s_i} \cdot e^{s_j}}{(\sum_k e^{s_k})^2} = -\frac{e^{s_i}}{\sum_k e^{s_k}} \cdot \frac{e^{s_j}}{\sum_k e^{s_k}} = -p_i \cdot p_j
\end{align*}

\textbf{Forma compatta}: combinando i due casi con il delta di Kronecker:
\[
    \frac{\partial p_i}{\partial s_j} = p_i(\delta_{ij} - p_j) \qquad \blacksquare
\]
\end{proofbox}

\begin{notabox}[Forma matriciale dello Jacobiano della softmax]
Lo Jacobiano completo della softmax \`e la matrice $\mathbf{J} \in \mathbb{R}^{n \times n}$ con elementi $J_{ij} = p_i(\delta_{ij} - p_j)$:
\[
    \mathbf{J} = \text{diag}(\mathbf{p}) - \mathbf{p}\mathbf{p}^\top
\]
dove $\text{diag}(\mathbf{p})$ \`e la matrice diagonale con $p_i$ sulla diagonale. Questa formula compatta \`e usata nell'implementazione efficiente del backward della self-attention (\hyperref[sec:grad_sa]{Sezione~7.5}).
\end{notabox}

\subsubsection{Propriet\`a di stabilit\`a numerica della softmax}

Un dettaglio tecnico importante per l'implementazione: in pratica, la formula ``ingenua'' della softmax pu\`o causare problemi numerici. Se uno degli score \`e molto grande (es. $s_i = 1000$), allora $e^{1000}$ \`e un numero talmente enorme che il computer non riesce a rappresentarlo (overflow!). La soluzione \`e il cosiddetto ``trucco del massimo'':

In pratica, si calcola la softmax con il ``trucco del massimo'' per evitare overflow numerico:
\[
    \text{softmax}(s_i) = \frac{e^{s_i - \max_j s_j}}{\sum_{k} e^{s_k - \max_j s_j}}
\]
Poich\'e la softmax \`e invariante per traslazione ($\text{softmax}(\mathbf{s} + c) = \text{softmax}(\mathbf{s})$), sottrarre il massimo non cambia il risultato ma evita che $e^{s_i}$ diventi $+\infty$ (float overflow).


\subsection{Il meccanismo di attenzione} \label{sec:attention}

Ora che abbiamo capito la softmax, possiamo affrontare il concetto pi\`u importante del deep learning moderno: l'\textbf{attenzione}. Senza comprendere l'attenzione, il Perceiver resta incomprensibile --- perch\'e l'attenzione \`e letteralmente l'unica operazione che il Perceiver usa per elaborare l'informazione.

L'attenzione \`e stata introdotta nel deep learning da Bahdanau, Cho e Bengio (2015) per risolvere il \textbf{collo di bottiglia} dei modelli encoder-decoder: l'intera informazione di una sequenza sorgente veniva compressa in un unico vettore di contesto di dimensione fissa, causando perdita di informazione per sequenze lunghe.

L'idea fondamentale \`e tanto semplice quanto potente: permettere al modello di ``rivisitare'' l'intera sequenza di input ad ogni passo di decodifica, concentrandosi di volta in volta sulle parti pi\`u rilevanti. In parole povere: invece di leggere tutto il libro e ricordare solo un riassunto, il modello pu\`o ``risfogliare'' il libro ogni volta che ne ha bisogno, guardando le pagine pi\`u utili.

\subsubsection{Scaled Dot-Product Attention (QKV)}

Ora che abbiamo capito la softmax, possiamo descrivere il meccanismo completo di attenzione. Il meccanismo usato nei Transformer e nel Perceiver si chiama \textbf{QKV attention} (Query-Key-Value), e funziona in tre passi: (1) calcola la ``compatibilit\`a'' tra query e chiavi, (2) normalizza con softmax, (3) usa i pesi risultanti per combinare i valori.

\begin{formulabox}[Scaled Dot-Product Attention]
Date tre matrici:
\begin{itemize}
    \item $\mathbf{Q} \in \mathbb{R}^{n_q \times d_k}$ --- le \textbf{Query} (``domande''): cosa sto cercando?
    \item $\mathbf{K} \in \mathbb{R}^{n_{kv} \times d_k}$ --- le \textbf{Key} (``chiavi''): cosa contiene ogni elemento?
    \item $\mathbf{V} \in \mathbb{R}^{n_{kv} \times d_v}$ --- i \textbf{Value} (``valori''): l'informazione effettiva di ogni elemento.
\end{itemize}
L'attenzione si calcola come:
\[
    \text{Attention}(\mathbf{Q}, \mathbf{K}, \mathbf{V}) = \text{softmax}\!\left(\frac{\mathbf{Q}\,\mathbf{K}^\top}{\sqrt{d_k}}\right)\,\mathbf{V}
\]
\end{formulabox}

\begin{notabox}[title={In parole semplici: Q, K, V}]
Pensiamo all'attenzione come a una ricerca in biblioteca:
\begin{itemize}
    \item La \textbf{Query} (Q) \`e la domanda che fai al bibliotecario: ``Cerco un libro sul Rinascimento''.
    \item La \textbf{Key} (K) \`e l'etichetta su ogni libro: ``Storia dell'arte'', ``Fisica quantistica'', ``Cucina italiana''...
    \item Il \textbf{Value} (V) \`e il contenuto effettivo del libro.
\end{itemize}
Il meccanismo di attenzione confronta la tua domanda (Q) con tutte le etichette (K), trova i libri pi\`u pertinenti (tramite softmax), e ti restituisce una ``media pesata'' dei loro contenuti (V). I libri pi\`u rilevanti pesano di pi\`u nella risposta.
\end{notabox}

\textbf{Perch\'e lo scaling $\sqrt{d_k}$?} Il prodotto scalare $\mathbf{q}_i \cdot \mathbf{k}_j$ cresce in magnitudine con $d_k$: se $q$ e $k$ hanno componenti con media 0 e varianza 1, allora $\mathbf{q} \cdot \mathbf{k}$ ha media 0 e varianza $d_k$. Per $d_k$ grande, i valori diventano molto grandi in valore assoluto, causando una softmax ``appuntita'' (quasi one-hot) con gradienti quasi nulli. Lo scaling per $\sqrt{d_k}$ riporta la varianza a 1, stabilizzando l'addestramento.

\begin{examplebox}[Effetto dello scaling: esempio con $d_k = 128$]
Supponiamo $\mathbf{q}, \mathbf{k} \in \mathbb{R}^{128}$ con componenti $\sim \mathcal{N}(0, 1)$.

\textbf{Senza scaling}: $\mathbf{q} \cdot \mathbf{k}$ ha deviazione standard $\sqrt{128} \approx 11.3$. Valori tipici dello score: $\pm 20$. Softmax di $[20, -20, 5]$: $\approx [1.000, 0.000, 0.000]$ --- quasi one-hot, gradiente $\approx 0$ per tutti tranne un elemento.

\textbf{Con scaling $\sqrt{d_k} = \sqrt{128} \approx 11.3$}: score diviso per 11.3. Valori tipici: $\pm 1.8$. Softmax di $[1.8, -1.8, 0.4]$: $\approx [0.632, 0.017, 0.156]$ --- distribuzione ``morbida'', gradienti informativi per l'ottimizzazione.
\end{examplebox}

\begin{examplebox}[Esempio numerico: effetto dello scaling sulla softmax]
Consideriamo 3 score grezzi prima dello scaling: $\mathbf{s} = [4.0, 2.0, 1.0]$. Vediamo cosa succede con e senza la divisione per $\sqrt{d_k}$.

\textbf{Caso 1 --- Senza scaling} ($d_k$ qualsiasi, nessuna divisione):
\begin{align*}
    e^{4.0} &= 54.598, \quad e^{2.0} = 7.389, \quad e^{1.0} = 2.718 \\
    \text{somma} &= 54.598 + 7.389 + 2.718 = 64.705 \\
    \text{softmax} &= \left[\frac{54.598}{64.705},\; \frac{7.389}{64.705},\; \frac{2.718}{64.705}\right] = [0.844,\; 0.114,\; 0.042]
\end{align*}
La distribuzione \`e molto \textbf{appuntita}: il primo elemento domina con peso $0.844$. I gradienti del secondo e terzo elemento sono quasi nulli.

\textbf{Caso 2 --- Con scaling} per $\sqrt{d_k} = \sqrt{64} = 8$:
\begin{align*}
    \mathbf{s}_{\text{scaled}} &= [4.0/8,\; 2.0/8,\; 1.0/8] = [0.500,\; 0.250,\; 0.125] \\
    e^{0.500} &= 1.649, \quad e^{0.250} = 1.284, \quad e^{0.125} = 1.133 \\
    \text{somma} &= 1.649 + 1.284 + 1.133 = 4.066 \\
    \text{softmax} &= \left[\frac{1.649}{4.066},\; \frac{1.284}{4.066},\; \frac{1.133}{4.066}\right] = [0.406,\; 0.316,\; 0.279]
\end{align*}
La distribuzione \`e molto pi\`u \textbf{uniforme}: tutti e tre gli elementi contribuiscono significativamente. I gradienti sono informativi per \emph{tutti} gli elementi.

\textbf{Confronto riassuntivo}:
\begin{center}
\begin{tabular}{@{}lcccc@{}}
\toprule
& Peso 1 & Peso 2 & Peso 3 & Entropia \\
\midrule
Senza scaling & 0.844 & 0.114 & 0.042 & 0.604 (bassa) \\
Con $\sqrt{d_k}{=}8$ & 0.406 & 0.316 & 0.279 & 1.077 (alta) \\
\bottomrule
\end{tabular}
\end{center}

\textbf{Conclusione}: senza scaling, con $d_k$ grande gli score crescono in magnitudine e la softmax diventa quasi one-hot (``winner takes all''). Lo scaling per $\sqrt{d_k}$ previene questo fenomeno, permettendo al modello di distribuire l'attenzione su pi\`u elementi e mantenendo gradienti informativi durante il training.
\end{examplebox}

\subsubsection{Self-attention vs Cross-attention} \label{sec:self_vs_cross}

La domanda che ci si pone \`e: da dove vengono Q, K, V? La risposta distingue due tipi fondamentali di attenzione:

\begin{itemize}[leftmargin=*]
    \item \textbf{Self-attention}: Q, K, V provengono \emph{dalla stessa sequenza}. Matrice di attenzione: $\mathbb{R}^{n \times n}$ (quadratica).

    \item \textbf{Cross-attention}: le query provengono da una sequenza, mentre le chiavi e i valori provengono da un'altra. Matrice di attenzione: $\mathbb{R}^{n_q \times n_{kv}}$ (rettangolare). Questo \`e il meccanismo chiave del Perceiver (si veda la \hyperref[sec:cross_attention]{Sezione~3.2}).
\end{itemize}

% TikZ: Self-attention vs Cross-attention side by side
\begin{figure}[H]
\centering
\begin{tikzpicture}[
    seqbox/.style={draw, rounded corners=2pt, minimum width=1.8cm, minimum height=0.5cm, font=\scriptsize},
    projbox/.style={draw, rounded corners=2pt, minimum width=1cm, minimum height=0.4cm, font=\tiny, fill=gray!10},
    arrowstyle/.style={-{Stealth[length=2mm]}, thick},
    >=Stealth,
    scale=0.85, transform shape
]

% LEFT: Self-attention
\node[font=\small\bfseries] at (2, 5) {Self-Attention};
\node[seqbox, fill=blue!20, minimum width=3cm] (seq1) at (2, 4) {Stessa sequenza $\mathbf{z}$};
\node[projbox, fill=red!15] (q1) at (0, 2.5) {Q};
\node[projbox, fill=green!15] (k1) at (2, 2.5) {K};
\node[projbox, fill=orange!15] (v1) at (4, 2.5) {V};

\draw[arrowstyle, red!60!black] (seq1.south) -- (q1.north) node[midway, left, font=\tiny] {$\mathbf{W}_Q$};
\draw[arrowstyle, green!60!black] (seq1.south) -- (k1.north) node[midway, right, font=\tiny] {$\mathbf{W}_K$};
\draw[arrowstyle, orange!60!black] (seq1.south) -- (v1.north) node[midway, right, font=\tiny] {$\mathbf{W}_V$};

\node[seqbox, fill=yellow!15] (att1) at (2, 1) {Attenzione $n \times n$};
\draw[arrowstyle] (q1.south) -- (att1.north);
\draw[arrowstyle] (k1.south) -- (att1.north);
\draw[arrowstyle] (v1.south) -- (att1.north);

\node[font=\tiny, gray, text width=3cm, align=center] at (2, 0.0) {Q, K, V dalla \textbf{stessa} sorgente\\Matrice \textbf{quadrata}};

% Separator
\draw[dashed, thick, gray] (5.5, 5.5) -- (5.5, -0.5);

% RIGHT: Cross-attention
\node[font=\small\bfseries] at (9, 5) {Cross-Attention};
\node[seqbox, fill=blue!20, minimum width=2.2cm] (lat) at (7, 4) {Latenti $\mathbf{z}$};
\node[seqbox, fill=red!20, minimum width=2.2cm] (inp) at (11, 4) {Input $\mathbf{x}$};

\node[projbox, fill=red!15] (q2) at (7, 2.5) {Q};
\node[projbox, fill=green!15] (k2) at (9.5, 2.5) {K};
\node[projbox, fill=orange!15] (v2) at (11, 2.5) {V};

\draw[arrowstyle, red!60!black] (lat.south) -- (q2.north) node[midway, left, font=\tiny] {$\mathbf{W}_Q$};
\draw[arrowstyle, green!60!black] (inp.south) -- (k2.north) node[midway, left, font=\tiny] {$\mathbf{W}_K$};
\draw[arrowstyle, orange!60!black] (inp.south) -- (v2.north) node[midway, right, font=\tiny] {$\mathbf{W}_V$};

\node[seqbox, fill=yellow!15, minimum width=2.5cm] (att2) at (9, 1) {Attenzione $N \times M$};
\draw[arrowstyle] (q2.south) -- (att2.north);
\draw[arrowstyle] (k2.south) -- (att2.north);
\draw[arrowstyle] (v2.south) -- (att2.north);

\node[font=\tiny, gray, text width=3.5cm, align=center] at (9, 0.0) {Q dai \textbf{latenti}, K/V dall'\textbf{input}\\Matrice \textbf{rettangolare} ($N \ll M$)};

\end{tikzpicture}
\caption{Confronto diretto tra self-attention e cross-attention. Nella self-attention, Q, K, V derivano dalla stessa sequenza e la matrice \`e quadrata ($n \times n$). Nella cross-attention del Perceiver, le query vengono dai latenti (pochi) e le chiavi/valori dall'input (molti), producendo una matrice rettangolare ($N \times M$) molto pi\`u efficiente.}
\label{fig:self_vs_cross_attention}
\end{figure}

\begin{examplebox}[Analogia: self-attention vs cross-attention]
\textbf{Self-attention = un gruppo di studenti che discute tra di loro}. Ogni studente (token) pu\`o porre domande a qualsiasi altro studente. Ogni studente \`e sia ``chi fa domande'' (Q) sia ``chi risponde'' (K, V). La matrice di attenzione \`e $n \times n$: ognuno presta attenzione a tutti gli altri.

\textbf{Cross-attention = uno studente che consulta un'enciclopedia}. Lo studente (latente, Q) fa domande, e l'enciclopedia (input, K/V) fornisce risposte. Lo studente non modifica l'enciclopedia, e l'enciclopedia non fa domande. La matrice di attenzione \`e $n_q \times n_{kv}$: poche domande (pochi latenti), molte possibili risposte (molti pixel/token).
\end{examplebox}

\begin{notabox}[Perch\'e la cross-attention \`e computazionalmente efficiente]
La chiave dell'efficienza del Perceiver \`e che le \textbf{query sono poche} ($N = 512$) anche se le \textbf{chiavi/valori sono molti} ($M = 50{,}176$). La matrice di attenzione ha dimensione $N \times M = 512 \times 50{,}176 \approx 25.7$M elementi, invece di $M \times M = 50{,}176^2 \approx 2.5 \times 10^9$ per la self-attention su tutti i pixel. Questo \`e possibile perch\'e il Perceiver ``comprime'' l'informazione attraverso un numero limitato di query.
\end{notabox}

\begin{notabox}[In parole semplici: self-attention vs cross-attention]
\textbf{Self-attention} = un gruppo di studenti che discute tra di loro. Ogni studente (token) pu\`o porre domande a qualsiasi altro studente. Se ci sono 50.000 studenti, servono $50.000^2 = 2.5$ miliardi di ``conversazioni''. \textbf{Impossibile} gestire tanti studenti!

\textbf{Cross-attention} = 512 studenti (latenti) che consultano un'enciclopedia con 50.000 voci (pixel). Ogni studente sfoglia l'enciclopedia cercando le voci rilevanti per la sua domanda. Servono solo $512 \times 50.000 = 25.6$ milioni di ``consultazioni'' --- $100\times$ meno della self-attention!
\end{notabox}


\begin{examplebox}[Esempio numerico completo: Multi-Head Attention con 2 teste]
Consideriamo un esempio ridotto con $H = 2$ teste, $d = 8$ (dimensione del modello) e $d_k = d/H = 4$ (dimensione per testa). L'input \`e composto da 3 token, ciascuno di 8 dimensioni.

\textbf{Input}: $\mathbf{Z} \in \mathbb{R}^{3 \times 8}$ (3 token, 8 dimensioni ciascuno):
\[
\mathbf{Z} = \begin{pmatrix} 1.0 & 0.5 & -0.3 & 0.8 & 0.2 & -0.1 & 0.6 & 0.4 \\ -0.2 & 0.7 & 0.1 & -0.5 & 0.9 & 0.3 & -0.4 & 0.8 \\ 0.3 & -0.6 & 0.4 & 0.2 & -0.1 & 0.5 & 0.7 & -0.3 \end{pmatrix}
\]

\textbf{Passo 1 --- Proiezione Q, K, V (con pesi semplificati):}

Per semplicit\`a, le matrici di proiezione sono $\mathbf{W}_Q, \mathbf{W}_K, \mathbf{W}_V \in \mathbb{R}^{8 \times 8}$. Dopo la proiezione, otteniamo $\mathbf{Q}, \mathbf{K}, \mathbf{V} \in \mathbb{R}^{3 \times 8}$.

\textbf{Passo 2 --- Divisione in 2 teste:}

Ogni vettore di 8 dimensioni viene \textbf{suddiviso} in 2 blocchi da 4:
\begin{align*}
\text{Testa 1}: \quad &\mathbf{Q}_1, \mathbf{K}_1, \mathbf{V}_1 \in \mathbb{R}^{3 \times 4} \quad \text{(primi 4 elementi)} \\
\text{Testa 2}: \quad &\mathbf{Q}_2, \mathbf{K}_2, \mathbf{V}_2 \in \mathbb{R}^{3 \times 4} \quad \text{(ultimi 4 elementi)}
\end{align*}

Supponiamo che dopo la proiezione e la divisione:

\textit{Testa 1}:
\[
\mathbf{Q}_1 = \begin{pmatrix} 0.8 & 0.3 & -0.2 & 0.5 \\ -0.1 & 0.6 & 0.4 & -0.3 \\ 0.2 & -0.4 & 0.3 & 0.1 \end{pmatrix}, \quad \mathbf{K}_1 = \begin{pmatrix} 0.5 & 0.2 & -0.1 & 0.4 \\ 0.3 & -0.3 & 0.6 & 0.1 \\ -0.2 & 0.5 & 0.2 & -0.4 \end{pmatrix}
\]

\textbf{Passo 3 --- Attenzione per la Testa 1:}

$\mathbf{S}_1 = \frac{\mathbf{Q}_1 \mathbf{K}_1^\top}{\sqrt{d_k}} = \frac{\mathbf{Q}_1 \mathbf{K}_1^\top}{\sqrt{4}} = \frac{\mathbf{Q}_1 \mathbf{K}_1^\top}{2}$

Calcoliamo $\mathbf{Q}_1 \mathbf{K}_1^\top$: il prodotto scalare $\mathbf{q}_1 \cdot \mathbf{k}_1 = 0.8 \times 0.5 + 0.3 \times 0.2 + (-0.2) \times (-0.1) + 0.5 \times 0.4 = 0.40 + 0.06 + 0.02 + 0.20 = 0.68$

Dopo aver calcolato tutti i prodotti scalari e diviso per 2, supponiamo:
\[
\mathbf{S}_1 = \begin{pmatrix} 0.34 & 0.15 & 0.08 \\ 0.12 & 0.28 & 0.19 \\ 0.06 & 0.22 & 0.31 \end{pmatrix}
\]

Applichiamo la softmax riga per riga:
\[
\mathbf{A}_1 = \text{softmax}(\mathbf{S}_1) \approx \begin{pmatrix} 0.392 & 0.324 & 0.284 \\ 0.310 & 0.363 & 0.327 \\ 0.288 & 0.339 & 0.373 \end{pmatrix}
\]

$\text{head}_1 = \mathbf{A}_1 \mathbf{V}_1 \in \mathbb{R}^{3 \times 4}$ (media pesata dei valori).

\textbf{Passo 4 --- Stessa procedura per la Testa 2:}

$\text{head}_2 = \text{softmax}\!\left(\frac{\mathbf{Q}_2 \mathbf{K}_2^\top}{2}\right) \mathbf{V}_2 \in \mathbb{R}^{3 \times 4}$

Ogni testa ha il suo ``punto di vista'': la testa 1 potrebbe catturare relazioni sintattiche, la testa 2 relazioni semantiche.

\textbf{Passo 5 --- Concatenazione:}

$\text{MultiHead} = [\text{head}_1 \;|\; \text{head}_2] \in \mathbb{R}^{3 \times 8}$

I due output da 4 dimensioni vengono affiancati per ricostruire un vettore di 8 dimensioni.

\textbf{Passo 6 --- Proiezione con $\mathbf{W}_O \in \mathbb{R}^{8 \times 8}$:}

$\text{output} = \text{MultiHead} \cdot \mathbf{W}_O \in \mathbb{R}^{3 \times 8}$

La matrice $\mathbf{W}_O$ ``mescola'' le informazioni delle 2 teste, permettendo al modello di combinare i diversi punti di vista in una rappresentazione unificata. L'output ha la \textbf{stessa dimensione} dell'input ($3 \times 8$), permettendo la connessione residua $\mathbf{y} = \mathbf{Z} + \text{output}$.
\end{examplebox}

% TikZ: Multi-head attention detailed diagram
\begin{figure}[H]
\centering
\begin{tikzpicture}[
    tensor/.style={draw, thick, rounded corners=2pt, minimum height=0.55cm, font=\scriptsize},
    op/.style={draw, circle, fill=gray!15, minimum size=0.5cm, font=\scriptsize},
    arrowstyle/.style={-{Stealth[length=2mm]}, thick},
    >=Stealth,
    scale=0.82, transform shape
]

% Input
\node[tensor, fill=blue!15, minimum width=2.5cm] (input) at (0, 0) {Input $\mathbf{Z} \in \mathbb{R}^{3 \times 8}$};

% Split into Q, K, V
\node[tensor, fill=red!10, minimum width=1.5cm] (Q) at (-3, -1.5) {$\mathbf{Q}$};
\node[tensor, fill=green!10, minimum width=1.5cm] (K) at (0, -1.5) {$\mathbf{K}$};
\node[tensor, fill=orange!10, minimum width=1.5cm] (V) at (3, -1.5) {$\mathbf{V}$};

\draw[arrowstyle] (input) -- (Q) node[midway, left, font=\tiny] {$\mathbf{W}_Q$};
\draw[arrowstyle] (input) -- (K) node[midway, left, font=\tiny] {$\mathbf{W}_K$};
\draw[arrowstyle] (input) -- (V) node[midway, right, font=\tiny] {$\mathbf{W}_V$};

% Head 1
\node[tensor, fill=yellow!20, minimum width=2.5cm] (h1) at (-2.5, -3.5) {Testa 1 ($d_k{=}4$)};
\node[font=\tiny, gray] at (-2.5, -4.1) {$\mathbf{Q}_1, \mathbf{K}_1, \mathbf{V}_1 \in \mathbb{R}^{3 \times 4}$};

% Head 2
\node[tensor, fill=yellow!20, minimum width=2.5cm] (h2) at (2.5, -3.5) {Testa 2 ($d_k{=}4$)};
\node[font=\tiny, gray] at (2.5, -4.1) {$\mathbf{Q}_2, \mathbf{K}_2, \mathbf{V}_2 \in \mathbb{R}^{3 \times 4}$};

% Arrows to heads
\draw[arrowstyle, dashed] (Q.south) -- (h1.north);
\draw[arrowstyle, dashed] (K.south) -- (h1.north);
\draw[arrowstyle, dashed] (V.south) -- (h1.north);
\draw[arrowstyle, dashed] (Q.south) -- (h2.north);
\draw[arrowstyle, dashed] (K.south) -- (h2.north);
\draw[arrowstyle, dashed] (V.south) -- (h2.north);

% Attention outputs
\node[tensor, fill=cyan!15, minimum width=2cm] (out1) at (-2.5, -5.5) {head$_1 \in \mathbb{R}^{3 \times 4}$};
\node[tensor, fill=cyan!15, minimum width=2cm] (out2) at (2.5, -5.5) {head$_2 \in \mathbb{R}^{3 \times 4}$};

\draw[arrowstyle] (h1) -- (out1) node[midway, right, font=\tiny] {Attn};
\draw[arrowstyle] (h2) -- (out2) node[midway, right, font=\tiny] {Attn};

% Concatenation
\node[tensor, fill=purple!15, minimum width=3cm] (concat) at (0, -7) {Concat $\in \mathbb{R}^{3 \times 8}$};
\draw[arrowstyle] (out1) -- (concat);
\draw[arrowstyle] (out2) -- (concat);

% W_O projection
\node[tensor, fill=green!20, minimum width=3cm] (final) at (0, -8.5) {Output $\in \mathbb{R}^{3 \times 8}$};
\draw[arrowstyle] (concat) -- (final) node[midway, right, font=\tiny] {$\mathbf{W}_O$};

% Dimension annotations
\node[font=\tiny, gray, text width=3cm, align=center] at (-5.5, -3.5) {Split: prime 4 dim\\per ogni Q, K, V};
\node[font=\tiny, gray, text width=3cm, align=center] at (5.5, -3.5) {Split: ultime 4 dim\\per ogni Q, K, V};

\end{tikzpicture}
\caption{Meccanismo di Multi-Head Attention con $H = 2$ teste. L'input viene proiettato in Q, K, V, poi suddiviso in teste di dimensione $d_k = d/H$. Ogni testa esegue l'attenzione indipendentemente. I risultati vengono concatenati e proiettati con $\mathbf{W}_O$ per produrre l'output finale.}
\label{fig:multi_head_attention_detail}
\end{figure}


\subsection{Layer Normalization} \label{sec:layernorm}

Ora che abbiamo capito l'attenzione, parliamo di un componente che sembra tecnico ma \`e assolutamente essenziale: la \textbf{Layer Normalization}. Senza di essa, le reti profonde come il Perceiver diventano instabili e impossibili da addestrare. Vediamo cosa fa e perch\'e.

La \textbf{Layer Normalization} (Ba et al., 2016) normalizza le attivazioni lungo la dimensione delle feature:

\begin{formulabox}[Layer Normalization]
Per un vettore $\mathbf{h} \in \mathbb{R}^D$:
\[
    \text{LN}(\mathbf{h}) = \boldsymbol{\gamma} \odot \frac{\mathbf{h} - \mu_{\mathbf{h}}}{\sqrt{\sigma^2_{\mathbf{h}} + \epsilon}} + \boldsymbol{\beta}
\]
dove:
\begin{align*}
    \mu_{\mathbf{h}} &= \frac{1}{D} \sum_{j=1}^{D} h_j \qquad \text{(media sulle feature)} \\
    \sigma^2_{\mathbf{h}} &= \frac{1}{D} \sum_{j=1}^{D} (h_j - \mu_{\mathbf{h}})^2 \qquad \text{(varianza sulle feature)} \\
    \boldsymbol{\gamma}, \boldsymbol{\beta} &\in \mathbb{R}^D \qquad \text{(parametri apprendibili: scala e shift)}
\end{align*}
$\epsilon \approx 10^{-5}$ \`e una costante per la stabilit\`a numerica.
\end{formulabox}

La LayerNorm prende ogni vettore latente e: (1) calcola media e deviazione standard sulle $D$ componenti; (2) normalizza ogni valore sottraendo la media e dividendo per la deviazione standard; (3) applica parametri appresi $\boldsymbol{\gamma}$ (scala) e $\boldsymbol{\beta}$ (shift) che permettono al modello di ``de-normalizzare'' dove serve.

\textbf{Pre-norm vs Post-norm}: nel Transformer originale, la LN \`e applicata \emph{dopo} l'attenzione/MLP (post-norm). Nel Perceiver e in architetture pi\`u recenti (GPT-2, ecc.), la LN \`e applicata \emph{prima} (pre-norm). Lo stile pre-norm \`e pi\`u stabile per reti molto profonde perch\'e ogni sotto-blocco riceve input normalizzati.

\begin{notabox}[In parole semplici: la Layer Normalization]
Detto in parole povere, la LayerNorm \`e come un ``regolatore di volume''. Se un vettore ha valori troppo grandi (volume troppo alto), la LayerNorm li abbassa; se troppo piccoli (volume troppo basso), li alza. Porta tutto a una scala ``comoda'' (media 0, varianza 1) prima di passare al blocco successivo. Senza questo regolatore, dopo 56 blocchi i valori sarebbero o esplosi a numeri enormi o svaniti a zero.
\end{notabox}

\begin{examplebox}[Esempio numerico della Layer Normalization]
Consideriamo un vettore $\mathbf{h} = [2.0, 4.0, 6.0, 8.0]$ ($D = 4$), con $\boldsymbol{\gamma} = [1, 1, 1, 1]$ e $\boldsymbol{\beta} = [0, 0, 0, 0]$ (inizializzazione standard):

\textbf{Passo 1 --- Media}: $\mu = \frac{2 + 4 + 6 + 8}{4} = 5.0$

\textbf{Passo 2 --- Varianza}: $\sigma^2 = \frac{(2-5)^2 + (4-5)^2 + (6-5)^2 + (8-5)^2}{4} = \frac{9 + 1 + 1 + 9}{4} = 5.0$

\textbf{Passo 3 --- Normalizzazione} (con $\epsilon = 10^{-5}$):
\begin{align*}
    \text{LN}(h_1) &= \frac{2.0 - 5.0}{\sqrt{5.0}} = \frac{-3.0}{2.236} = -1.342 \\
    \text{LN}(h_2) &= \frac{4.0 - 5.0}{\sqrt{5.0}} = \frac{-1.0}{2.236} = -0.447 \\
    \text{LN}(h_3) &= \frac{6.0 - 5.0}{\sqrt{5.0}} = \frac{1.0}{2.236} = 0.447 \\
    \text{LN}(h_4) &= \frac{8.0 - 5.0}{\sqrt{5.0}} = \frac{3.0}{2.236} = 1.342
\end{align*}

\textbf{Risultato}: $\text{LN}(\mathbf{h}) = [-1.342, -0.447, 0.447, 1.342]$ con media $\approx 0$ e varianza $\approx 1$.

I parametri apprendibili $\boldsymbol{\gamma}$ e $\boldsymbol{\beta}$ possono poi ``de-normalizzare'' dove serve: dopo il training, il modello potrebbe apprendere $\gamma_1 = 2.0$, $\beta_1 = 0.5$, ottenendo $\text{output}_1 = 2.0 \times (-1.342) + 0.5 = -2.184$.
\end{examplebox}

\begin{examplebox}[Esempio numerico dettagliato: Layer Normalization con gamma e beta]
Consideriamo un vettore latente $\mathbf{h} = [2.0, -1.0, 0.5, 1.5]$ ($D = 4$), con $\epsilon = 10^{-5}$.

\textbf{Passo 1 --- Media}:
\[
\mu = \frac{2.0 + (-1.0) + 0.5 + 1.5}{4} = \frac{3.0}{4} = 0.75
\]

\textbf{Passo 2 --- Varianza}:
\begin{align*}
\sigma^2 &= \frac{(2.0 - 0.75)^2 + (-1.0 - 0.75)^2 + (0.5 - 0.75)^2 + (1.5 - 0.75)^2}{4} \\
&= \frac{1.5625 + 3.0625 + 0.0625 + 0.5625}{4} = \frac{5.25}{4} = 1.3125
\end{align*}
$\sqrt{\sigma^2 + \epsilon} = \sqrt{1.3125} \approx 1.1456$

\textbf{Passo 3 --- Normalizzazione} (caso identit\`a: $\boldsymbol{\gamma} = [1,1,1,1]$, $\boldsymbol{\beta} = [0,0,0,0]$):
\begin{align*}
\hat{h}_1 &= \frac{2.0 - 0.75}{1.1456} = \frac{1.25}{1.1456} \approx 1.091 \\
\hat{h}_2 &= \frac{-1.0 - 0.75}{1.1456} = \frac{-1.75}{1.1456} \approx -1.527 \\
\hat{h}_3 &= \frac{0.5 - 0.75}{1.1456} = \frac{-0.25}{1.1456} \approx -0.218 \\
\hat{h}_4 &= \frac{1.5 - 0.75}{1.1456} = \frac{0.75}{1.1456} \approx 0.654
\end{align*}

\textbf{Risultato} (caso identit\`a): $\text{LN}(\mathbf{h}) = [1.091, -1.527, -0.218, 0.654]$

Verifica: media $\approx 0$, varianza $\approx 1$ (come atteso).

\textbf{Passo 4 --- Con $\boldsymbol{\gamma} = [2, 2, 2, 2]$ e $\boldsymbol{\beta} = [1, 1, 1, 1]$} (caso non banale):
\begin{align*}
\text{out}_1 &= 2 \times 1.091 + 1 = 3.182 \\
\text{out}_2 &= 2 \times (-1.527) + 1 = -2.055 \\
\text{out}_3 &= 2 \times (-0.218) + 1 = 0.563 \\
\text{out}_4 &= 2 \times 0.654 + 1 = 2.309
\end{align*}

\textbf{Risultato} (con $\gamma{=}2$, $\beta{=}1$): $\text{LN}(\mathbf{h}) = [3.182, -2.055, 0.563, 2.309]$

Osservazione: con $\boldsymbol{\gamma} = [2,2,2,2]$ e $\boldsymbol{\beta} = [1,1,1,1]$, la media diventa $1$ e la deviazione standard diventa $2$. I parametri $\boldsymbol{\gamma}$ e $\boldsymbol{\beta}$ permettono al modello di ``ri-scalare'' l'output normalizzato alla scala pi\`u utile per il blocco successivo. Durante il training, il modello apprende automaticamente i valori ottimali di $\boldsymbol{\gamma}$ e $\boldsymbol{\beta}$ per ogni layer.
\end{examplebox}


\subsection{Residual Connections (Connessioni Residue)} \label{sec:residual}

Ora che abbiamo capito come funzionano l'attenzione e la normalizzazione, dobbiamo affrontare un ultimo ingrediente fondamentale: le \textbf{residual connections}. Senza di esse, il Perceiver (che ha fino a 56 blocchi di trasformazione!) non potrebbe essere addestrato. Vediamo perch\'e.

Introdotte da He et al.\ (2016) con ResNet (si veda l'\hyperref[app:resnet]{Appendice~G}), le \textbf{residual connections} aggiungono l'input di un blocco direttamente al suo output:

\begin{formulabox}[Residual Connection]
\[
    \mathbf{y} = \mathbf{x} + f(\mathbf{x})
\]
dove $f$ \`e la trasformazione del blocco (attenzione, MLP, ecc.) e $\mathbf{x}$ \`e l'input originale.
\end{formulabox}

\textbf{Perch\'e sono essenziali?}
\begin{enumerate}
    \item \textbf{Flusso del gradiente}: durante la backpropagation, il gradiente di $\mathbf{y}$ rispetto a $\mathbf{x}$ \`e $\frac{\partial \mathbf{y}}{\partial \mathbf{x}} = \mathbf{I} + \frac{\partial f}{\partial \mathbf{x}}$. Il termine $\mathbf{I}$ (identit\`a) garantisce che il gradiente possa fluire direttamente attraverso lo skip connection, mitigando il vanishing gradient.
    \item \textbf{Apprendimento incrementale}: il blocco deve apprendere solo il ``residuo'' $f(\mathbf{x})$, cio\`e la \emph{correzione} da applicare all'input.
    \item \textbf{Reti molto profonde}: senza residual connections, reti con pi\`u di 20--30 layer diventano impossibili da addestrare. Con le residual connections, il Perceiver usa fino a 56 blocchi di trasformazione ($T = 8$ iterazioni $\times$ $(L+1) = 7$ blocchi per iterazione).
\end{enumerate}

\begin{examplebox}[Esempio: flusso del gradiente con e senza residual connection]
Consideriamo una catena di 10 blocchi, ciascuno con funzione $f(\mathbf{x}) = 0.5 \cdot \mathbf{x}$ (attenuazione).

\textbf{Senza residual connection}: il gradiente attraverso 10 blocchi \`e:
\[
    \frac{\partial \mathbf{y}}{\partial \mathbf{x}} = (0.5)^{10} = 0.000977 \approx 10^{-3}
\]
Il gradiente si \`e attenuato di $1000\times$: \textbf{vanishing gradient}.

\textbf{Con residual connection}: ogni blocco calcola $\mathbf{y} = \mathbf{x} + 0.5 \cdot \mathbf{x} = 1.5 \cdot \mathbf{x}$. Il gradiente locale \`e $1 + 0.5 = 1.5$. Dopo 10 blocchi:
\[
    \frac{\partial \mathbf{y}}{\partial \mathbf{x}} = (1.5)^{10} = 57.7
\]
Il gradiente \textbf{non svanisce}. In realt\`a, nella pratica i gradienti non si moltiplicano cos\`i direttamente (il percorso diretto $\mathbf{I}$ fornisce sempre un contributo costante), ma l'esempio illustra il principio.
\end{examplebox}

\begin{notabox}[Residual connections nel Perceiver: conteggio]
In una singola iterazione del Perceiver:
\begin{itemize}
    \item 1 residual connection dopo la cross-attention
    \item 1 residual connection dopo l'MLP post-cross-attention
    \item $L = 6$ residual connections dopo la self-attention nel Latent Transformer
    \item $L = 6$ residual connections dopo l'MLP nel Latent Transformer
\end{itemize}
Totale per iterazione: $2 + 12 = 14$ residual connections.
Con $T = 8$ iterazioni: $8 \times 14 = \textbf{112 residual connections}$ totali nella rete!
Senza di esse, il training sarebbe impossibile.
\end{notabox}

\begin{notabox}[In parole semplici: le residual connections]
Detto in parole povere, la residual connection \`e come tenere una \textbf{copia di backup} prima di fare modifiche a un documento. Se le modifiche migliorano il documento, bene --- si tiene la versione modificata. Se le modifiche peggiorano il documento, non \`e un problema --- la copia di backup \`e ancora l\`i. In pratica, ogni blocco del Perceiver dice: ``prendo l'input, ci applico una trasformazione, e poi sommo il risultato all'input originale''. Il blocco deve solo imparare la \emph{correzione} da applicare, non ricostruire tutto da zero.
\end{notabox}

\begin{examplebox}[Esempio numerico: residual connection e analisi del gradiente]
Consideriamo un esempio concreto con vettori a 3 dimensioni.

\textbf{Forward pass}:
\begin{itemize}
\item Input: $\mathbf{x} = [1.0, 2.0, 3.0]$
\item Il blocco (ad es. un MLP) produce: $f(\mathbf{x}) = [0.1, -0.3, 0.5]$
\item Output con residual: $\mathbf{y} = \mathbf{x} + f(\mathbf{x}) = [1.0 + 0.1,\; 2.0 + (-0.3),\; 3.0 + 0.5] = [1.1, 1.7, 3.5]$
\end{itemize}

\textbf{Backward pass --- analisi del gradiente}:

Il gradiente di $\mathbf{y}$ rispetto a $\mathbf{x}$ \`e:
\[
\frac{\partial \mathbf{y}}{\partial \mathbf{x}} = \mathbf{I} + \frac{\partial f}{\partial \mathbf{x}} = \underbrace{\begin{pmatrix} 1 & 0 & 0 \\ 0 & 1 & 0 \\ 0 & 0 & 1 \end{pmatrix}}_{\text{identit\`a (skip)}} + \underbrace{\mathbf{J}_f(\mathbf{x})}_{\text{Jacobiano del blocco}}
\]

Supponiamo che il Jacobiano di $f$ sia:
\[
\mathbf{J}_f(\mathbf{x}) = \begin{pmatrix} 0.05 & 0.01 & -0.02 \\ -0.03 & 0.08 & 0.01 \\ 0.02 & -0.01 & 0.06 \end{pmatrix}
\]

Allora il gradiente totale \`e:
\[
\frac{\partial \mathbf{y}}{\partial \mathbf{x}} = \begin{pmatrix} 1.05 & 0.01 & -0.02 \\ -0.03 & 1.08 & 0.01 \\ 0.02 & -0.01 & 1.06 \end{pmatrix}
\]

I valori sulla diagonale sono tutti $\approx 1$: il gradiente non svanisce mai!

\textbf{Senza residual connection}: $\mathbf{y} = f(\mathbf{x})$, il gradiente sarebbe solo $\mathbf{J}_f(\mathbf{x})$ con valori diagonali $\approx 0.05\text{--}0.08$. Dopo 10 blocchi in cascata: $(0.06)^{10} \approx 6 \times 10^{-13}$ --- il gradiente \`e praticamente \textbf{svanito}.

\textbf{Con residual connection}: dopo 10 blocchi, il gradiente diagonale \`e $\approx (1.06)^{10} \approx 1.79$ --- perfettamente utilizzabile per l'ottimizzazione.
\end{examplebox}


\subsection{Inizializzazione dei pesi: Xavier e TruncatedNormal} \label{sec:inizializzazione}

Prima di parlare dell'architettura, dobbiamo affrontare un ultimo prerequisito: come inizializziamo i pesi della rete? Questa domanda potrebbe sembrare banale, ma in realt\`a \`e cruciale. Immaginiamo di costruire una catena di 56 blocchi (come nel Perceiver): se i pesi iniziali sono troppo grandi, le attivazioni esplodono come una reazione a catena; se troppo piccoli, svaniscono come un eco che si perde. Il Perceiver usa due schemi di inizializzazione attentamente calibrati:

\begin{formulabox}[title={Xavier Uniform (Glorot \& Bengio, 2010)}]
Per un layer con $n_{\text{in}}$ input e $n_{\text{out}}$ output:
\[
    W_{ij} \sim \mathcal{U}\!\left(-\sqrt{\frac{6}{n_{\text{in}} + n_{\text{out}}}},\; \sqrt{\frac{6}{n_{\text{in}} + n_{\text{out}}}}\right)
\]
\textbf{Obiettivo}: mantenere la varianza delle attivazioni costante attraverso i layer. Se $\text{Var}(\text{input}) = 1$, allora $\text{Var}(\text{output}) \approx 1$ (n\'e esplosione n\'e scomparsa).

Usata per: $\mathbf{W}_Q, \mathbf{W}_K, \mathbf{W}_V, \mathbf{W}_O, \mathbf{W}_1, \mathbf{W}_2, \mathbf{W}_{\text{embed}}, \mathbf{W}_{\text{class}}$.
\end{formulabox}

\begin{formulabox}[TruncatedNormal]
\[
    W_{ij} \sim \text{TruncatedNormal}(\mu, \sigma, a, b)
\]
dove i campioni da $\mathcal{N}(\mu, \sigma^2)$ vengono rigettati se cadono fuori da $[a, b]$.

Usata per: l'array latente $\mathbf{z}_0$ con $\mu = 0$, $\sigma = 0.02$, $a = -0.04$, $b = 0.04$.
\end{formulabox}

\begin{examplebox}[Xavier init per la proiezione Q del Perceiver]
$\mathbf{W}_Q \in \mathbb{R}^{1024 \times 1024}$: $n_{\text{in}} = 1024$, $n_{\text{out}} = 1024$.
\[
    W_{ij} \sim \mathcal{U}\!\left(-\sqrt{\frac{6}{2048}},\; \sqrt{\frac{6}{2048}}\right) = \mathcal{U}(-0.054, 0.054)
\]
I pesi iniziali sono numeri molto piccoli nell'intervallo $[-0.054, 0.054]$. La matrice $\mathbf{W}_Q$ iniziale \`e essenzialmente ``rumore strutturato'' che verr\`a raffinato durante il training.
\end{examplebox}

\begin{notabox}[In parole semplici: l'inizializzazione dei pesi]
Detto in parole povere, inizializzare i pesi \`e come accordare uno strumento musicale prima di suonare. Se le corde sono troppo tese (pesi troppo grandi), il suono \`e stridulo e distorto (attivazioni esplose). Se troppo lente (pesi troppo piccoli), il suono \`e impercettibile (attivazioni svanite). Xavier init trova la tensione giusta per ogni corda, in modo che il segnale si propaghi senza distorsioni attraverso tutti i 56 blocchi della rete.
\end{notabox}

% TikZ: Weight initialization distributions
\begin{figure}[H]
\centering
\begin{tikzpicture}
\begin{axis}[
    width=14cm, height=7cm,
    xlabel={Valore del peso $w$},
    ylabel={Densit\`a di probabilit\`a},
    xmin=-0.12, xmax=0.12,
    ymin=0, ymax=25,
    legend pos=north east,
    grid=major,
    grid style={gray!20},
    thick,
    samples=200
]
% Xavier Normal with sigma = 0.031 (for n_in=n_out=1024)
\addplot[blue, very thick, domain=-0.12:0.12] {1/(0.031*sqrt(2*pi)) * exp(-x^2/(2*0.031^2))};
\addlegendentry{Xavier Normal ($\sigma \approx 0.031$)}

% TruncatedNormal(0, 0.02, -0.04, 0.04)
\addplot[red, very thick, domain=-0.04:0.04] {1/(0.02*sqrt(2*pi)) * exp(-x^2/(2*0.02^2)) / 0.9545};
\addlegendentry{TruncatedNormal($\sigma{=}0.02$, $|w|{<}0.04$)}

% Zero outside truncation
\addplot[red, very thick, domain=-0.12:-0.041] {0};
\addplot[red, very thick, domain=0.041:0.12] {0};

% Uniform[-0.054, 0.054] for Xavier Uniform
\addplot[green!60!black, very thick, domain=-0.054:0.054] {1/(2*0.054)};
\addlegendentry{Xavier Uniform ($[-0.054, 0.054]$)}

% Zero outside uniform
\addplot[green!60!black, very thick, domain=-0.12:-0.055] {0};
\addplot[green!60!black, very thick, domain=0.055:0.12] {0};

% Truncation lines
\draw[red, dashed, thin] (axis cs:-0.04,0) -- (axis cs:-0.04,22);
\draw[red, dashed, thin] (axis cs:0.04,0) -- (axis cs:0.04,22);

\end{axis}
\end{tikzpicture}
\caption{Confronto delle distribuzioni di inizializzazione dei pesi. La Xavier Normal (blu) \`e una gaussiana centrata in 0. La TruncatedNormal (rosso, usata per $\mathbf{z}_0$) \`e simile ma tronca i valori oltre $\pm 0.04$ (linee tratteggiate rosse). La Xavier Uniform (verde) ha densit\`a costante nell'intervallo $[-0.054, 0.054]$. Tutti e tre gli schemi producono pesi piccoli, evitando esplosione o scomparsa delle attivazioni.}
\label{fig:weight_init_distributions}
\end{figure}

% Immagine Xavier Uniform dagli appunti
\begin{figure}[H]
\centering
\includegraphics[width=0.7\textwidth]{paper_figures/docx_xavier_uniform.png}
\caption{Inizializzazione Xavier Uniform per le matrici di peso $\mathbf{W}_K$ e $\mathbf{W}_V$. Con $C = D = 512$, il range di inizializzazione \`e $[-0.076, +0.076]$. Questa scelta evita che i segnali esplodano o collassino durante le prime iterazioni di training.}
\label{fig:docx_xavier_uniform}
\end{figure}



\newpage
%% ============================================================
%%  SEZIONE 3 --- ARCHITETTURA DEL PERCEIVER
%% ============================================================
\section{Architettura del Perceiver} \label{sec:architettura}

Vediamo cosa succede quando un'immagine entra nel Perceiver.

\subsection{Panoramica generale} \label{sec:panoramica}

L'architettura del Perceiver \`e costruita su due componenti fondamentali (Jaegle et al., 2021, Sezione~3.1):
\begin{itemize}[leftmargin=*]
    \item un \textbf{modulo di cross-attention} che mappa un byte array (es.\ un array di pixel) e un latent array in un latent array;
    \item una \textbf{torre Transformer} (Latent Transformer) che mappa un latent array in un latent array.
\end{itemize}
La dimensione del byte array \`e determinata dall'input ed \`e generalmente grande (es.\ le immagini ImageNet a risoluzione $224$ hanno $M = 50.176$ pixel), mentre la dimensione del latent array \`e un iperparametro tipicamente molto pi\`u piccolo ($N = 512$). Il modello applica cross-attention e Transformer in alternanza, proiettando l'array ad alta dimensionalit\`a attraverso un \textbf{collo di bottiglia} nello spazio latente prima di elaborarlo con un Transformer profondo, per poi usare la rappresentazione risultante per interrogare nuovamente l'input. Poich\'e i pesi vengono opzionalmente condivisi tra le iterazioni (dalla seconda in poi), il modello pu\`o essere interpretato come una \textbf{RNN srotolata in profondit\`a} usando lo stesso input, anzich\'e nel tempo.

\begin{figure}[H]
\centering
\includegraphics[width=0.85\textwidth]{paper_figures/docx_architecture_full.png}
\label{fig:docx_architecture_full}
\end{figure}

I valori di $N$, $D$, $T$, $L$, $H$ sono \textbf{identici} per tutti i dataset --- l'unica cosa che cambia \`e la dimensione dell'input ($M$ e $C$). Questo \`e il punto centrale del paper: la stessa architettura funziona su modalit\`a completamente diverse.

\begin{table}[htbp]
\centering
\caption{Valori degli iperparametri del Perceiver per i diversi dataset.}
\begin{tabular}{@{}lccccccc@{}}
\toprule
\textbf{Dataset} & $M$ & $C$ & $N$ & $D$ & $T$ & $L$ & $H$ \\
\midrule
ImageNet & 50{,}176 & 261 & 512 & 1024 & 8 & 6 & 8 \\
AudioSet (raw) & 480 & 128 & 512 & 1024 & 8 & 6 & 8 \\
AudioSet (mel) & 4{,}800 & --- & 512 & 1024 & 8 & 6 & 8 \\
ModelNet40 & 2{,}048 & 390 & 512 & 1024 & 8 & 6 & 8 \\
\bottomrule
\end{tabular}
\end{table}

% ====================================================================
%  PIPELINE DEL FORWARD PASS
% ====================================================================

\newpage
\subsection{Pipeline del Forward Pass} \label{sec:pipeline_overview}

Il flusso completo del Perceiver, dall'input alla predizione, segue questa pipeline:

\begin{center}
\fbox{\parbox{0.92\textwidth}{\centering
\textbf{\colorbox{yellow!40}{Input}} $\;\to\;$ Fourier Features $\;\to\;$ Proiezione Lineare $\;\to\;$ Latent Array $\;\to\;$ Cross-Attention $\;\to\;$ Residual + MLP $\;\to\;$ Latent Transformer $\;\to\;$ Residual + MLP $\;\to\;$ \textit{(ripetizione $T$ volte)} $\;\to\;$ Pooling $\;\to\;$ Classificatore $\;\to\;$ Loss%
}}
\end{center}

Nelle sezioni seguenti analizziamo ogni componente nel dettaglio.

% ====================================================================
%  DIAGRAMMI COMPONENTE PER COMPONENTE DELL'ARCHITETTURA PERCEIVER
% ====================================================================

\subsubsection{Input Processing} \label{sec:input_processing}

Il Perceiver riceve in ingresso due strutture dati: il \textbf{byte array} (l'input) e il \textbf{latent array} (la memoria del modello). Le figure seguenti mostrano come queste due strutture alimentano il modulo di cross-attention tramite le proiezioni Q, K, V.

\begin{figure}[H]
\centering
\begin{minipage}[t]{0.48\textwidth}
\centering
\includegraphics[width=\textwidth]{paper_figures/docx_byte_array_crop.png}
\end{minipage}
\hfill
\begin{minipage}[t]{0.48\textwidth}
\centering
\includegraphics[width=\textwidth]{paper_figures/docx_latent_array_crop.png}
\end{minipage}
\label{fig:byte_and_latent_arrays}
\end{figure}

\paragraph{Byte Array $\mathbf{x} \in \mathbb{R}^{M \times C}$ (in verde).}
Qualsiasi input --- immagine, audio, point cloud --- viene appiattito in una matrice dove ogni riga \`e un elemento (un pixel, un campione, un punto) e ogni colonna un canale. Per ImageNet: $M = 224 \times 224 = 50.176$ righe e $C = 3$ colonne (R, G, B). Il byte array fornisce \textbf{Key} e \textbf{Value} alla cross-attention. Questa operazione di \textit{flatten} elimina qualsiasi assunzione sulla struttura spaziale: il Perceiver vede i pixel come 50.176 vettori indipendenti, senza sapere che formano una griglia 2D.

\paragraph{Latent Array $\mathbf{z}_0 \in \mathbb{R}^{N \times D}$ (in viola).}
Una matrice $512 \times 1024$ di parametri \textbf{appresi}, inizializzata con \texttt{TruncatedNormal(0, 0.02)} e identica per ogni input all'inizio di ogni forward pass. Il latent array fornisce le \textbf{Query} alla cross-attention. Pensatelo come 512 ``slot di memoria'' vuoti che verranno riempiti con le informazioni pi\`u rilevanti estratte dall'input. Con $N = 512 \ll M = 50.176$, i latenti forzano una compressione di circa $100\times$: questo \`e il collo di bottiglia che rende il Perceiver efficiente.

\bigskip

Ma se il byte array perde l'informazione spaziale con il flatten, come fa il modello a sapere dove si trova ogni pixel? La risposta viene dal prossimo componente: le \textbf{Fourier features}.


% ---- DIAGRAMMA 2: FOURIER FEATURES ----

\begin{figure}[H]
\centering
\begin{tikzpicture}[scale=0.85, transform shape]

% --- Pixel singolo ---
\fill[red!15, draw=red!60!black, thick, rounded corners=3pt] (0, 0) rectangle (2.6, 1.5);
\node[font=\small\bfseries] at (1.3, 1.05) {1 pixel};
\node[font=\scriptsize] at (1.3, 0.45) {$[R, G, B]$};
\node[below, font=\scriptsize, text=gray] at (1.3, -0.15) {3 valori};

% --- Freccia verso Fourier ---
\draw[-{Stealth[length=3mm]}, very thick, blue!60!black] (3.0, 0.75) -- (4.8, 0.75);

% --- Box Fourier Transform ---
\fill[yellow!12, draw=orange!60!black, thick, rounded corners=5pt] (5.0, -0.2) rectangle (9.5, 1.7);
\node[font=\small\bfseries, orange!70!black] at (7.25, 1.2) {Fourier Features};
\node[font=\scriptsize, align=center] at (7.25, 0.4) {$\sin(f_k \pi d),\; \cos(f_k \pi d)$\\per $k=1,\ldots,64$};

% --- Freccia verso output ---
\draw[-{Stealth[length=3mm]}, very thick, blue!60!black] (9.8, 0.75) -- (11.3, 0.75);

% --- Output arricchito ---
\fill[green!8, draw=green!50!black, thick, rounded corners=3pt] (11.5, -0.5) rectangle (16.5, 2.0);
\node[font=\small\bfseries, green!40!black] at (14.0, 1.55) {Pixel arricchito};

% Tre blocchi separati per chiarezza
\fill[red!10] (11.8, 0.5) rectangle (12.8, 1.1);
\draw[red!40!black] (11.8, 0.5) rectangle (12.8, 1.1);
\node[font=\tiny] at (12.3, 0.8) {RGB (3)};

\fill[blue!10] (13.0, 0.5) rectangle (14.3, 1.1);
\draw[blue!40!black] (13.0, 0.5) rectangle (14.3, 1.1);
\node[font=\tiny] at (13.65, 0.8) {FF asse $x$ (129)};

\fill[orange!10] (14.5, 0.5) rectangle (15.8, 1.1);
\draw[orange!40!black] (14.5, 0.5) rectangle (15.8, 1.1);
\node[font=\tiny] at (15.15, 0.8) {FF asse $y$ (129)};

\node[font=\scriptsize, text=gray] at (14.0, -0.1) {Totale: $3 + 129 + 129 = 261$};

% --- Annotazione dimensionale ---
\draw[decorate, decoration={brace, amplitude=5pt}, thick, purple!60!black]
    (16.7, 2.0) -- (16.7, -0.5)
    node[midway, right=6pt, font=\footnotesize, purple!60!black] {$C = 261$};

\end{tikzpicture}
\caption{Componente 2 --- Fourier Features. Ogni pixel passa da 3 valori (solo colore) a 261 valori che includono la codifica posizionale tramite funzioni sinusoidali a 64 frequenze per ciascun asse.}
\label{fig:fourier_diagram}
\end{figure}

Perch\'e arricchire ogni pixel? Perch\'e un pixel con valori $[R, G, B]$ contiene solo informazione sul colore, ma non sa \textit{dove} si trova nell'immagine. Due pixel identicamente rossi in posizioni diverse dell'immagine sarebbero indistinguibili. Le \textbf{Fourier features} risolvono questo problema aggiungendo a ogni pixel una codifica della sua posizione spaziale tramite funzioni sinusoidali a frequenze diverse.

In concreto, per ogni asse spaziale (asse $x$ e asse $y$), si calcolano $K_{\max}$ coppie di $(\sin, \cos)$ a frequenze crescenti $f_1, f_2, \ldots, f_K$. Questo produce $2 \times K_{\max}$ valori per asse, pi\`u un valore per la posizione normalizzata stessa. Con $K_{\max} = 64$ bande di frequenza e 2 assi spaziali, si ottengono $2 \times (2 \times 64 + 1) = 258$ valori posizionali, che sommati ai 3 canali di colore danno $C = 261$.

Dopo questa fase, il byte array ha dimensione $M \times C = 50.176 \times 261$: ogni pixel ``sa'' sia il proprio colore sia la propria posizione. Si noti che le frequenze basse catturano la posizione ``a grandi linee'' (met\`a sinistra vs.\ met\`a destra), mentre le frequenze alte codificano posizioni precise a livello sub-pixel. Questo approccio \`e ispirato dal lavoro di Tancik et al.\ (2020) --- si veda la \hyperref[sec:fourier]{Sezione~4} per i dettagli matematici completi.


% Il Latent Array è già descritto nella sezione Input Processing sopra.


% ---- DIAGRAMMA 4: IL MODULO DI CROSS-ATTENTION ----

\begin{figure}[H]
\centering
\begin{tikzpicture}[scale=0.85, transform shape,
    tensor/.style={draw, thick, rounded corners=2pt, minimum height=0.65cm, font=\footnotesize},
    proj/.style={draw, thick, fill=orange!12, rounded corners=2pt, minimum height=0.55cm, font=\scriptsize},
    op/.style={draw, thick, rounded corners=2pt, fill=gray!12, minimum height=0.55cm, font=\scriptsize},
    arrowstyle/.style={-{Stealth[length=2.5mm]}, thick}
]

% === SORGENTI ===
% Input (byte array)
\node[tensor, fill=red!12, minimum width=2.8cm] (x) at (0, 0) {$\mathbf{x} \in \mathbb{R}^{M \times C}$};
\node[font=\scriptsize, below=0.08cm of x, red!60!black] {Input ($50.176 \times 261$)};

% Latent array
\node[tensor, fill=blue!12, minimum width=2.8cm] (z) at (0, 4) {$\mathbf{z} \in \mathbb{R}^{N \times D}$};
\node[font=\scriptsize, above=0.08cm of z, blue!60!black] {Latenti ($512 \times 1024$)};

% === PROIEZIONI ===
% W_K
\node[proj] (wk) at (4.5, 0.8) {$\times\, \mathbf{W}_K$};
\node[font=\tiny, above=0.05cm of wk] {$C \times d_k$};
\draw[arrowstyle, red!60!black] (x.east) -- ++(0.8,0) |- (wk.west);

% W_V
\node[proj] (wv) at (4.5, -0.8) {$\times\, \mathbf{W}_V$};
\node[font=\tiny, below=0.05cm of wv] {$C \times d_v$};
\draw[arrowstyle, red!60!black] (x.east) -- ++(0.8,0) |- (wv.west);

% W_Q
\node[proj] (wq) at (4.5, 4) {$\times\, \mathbf{W}_Q$};
\node[font=\tiny, above=0.05cm of wq] {$D \times d_k$};
\draw[arrowstyle, blue!60!black] (z.east) -- (wq.west);

% === RISULTATI PROIEZIONI ===
\node[tensor, fill=red!8, minimum width=1.8cm] (K) at (7.5, 0.8) {$\mathbf{K}$};
\node[font=\tiny, above=0.05cm of K] {$M \times d_k$};
\draw[arrowstyle] (wk.east) -- (K.west);

\node[tensor, fill=red!8, minimum width=1.8cm] (V) at (7.5, -0.8) {$\mathbf{V}$};
\node[font=\tiny, below=0.05cm of V] {$M \times d_v$};
\draw[arrowstyle] (wv.east) -- (V.west);

\node[tensor, fill=blue!8, minimum width=1.8cm] (Q) at (7.5, 4) {$\mathbf{Q}$};
\node[font=\tiny, above=0.05cm of Q] {$N \times d_k$};
\draw[arrowstyle] (wq.east) -- (Q.west);

% === SCORE: Q * K^T ===
\node[op, minimum width=2.0cm] (score) at (11, 2.4) {$\mathbf{Q} \cdot \mathbf{K}^\top$};
\draw[arrowstyle, blue!50!black] (Q.east) -- ++(0.8,0) |- (score.west);
\draw[arrowstyle, red!50!black] (K.east) -- ++(0.8,0) |- (score.west);

\node[tensor, fill=purple!10, minimum width=1.6cm] (S) at (14, 2.4) {$\mathbf{S}$};
\node[font=\tiny, above=0.05cm of S] {$N \times M$};
\draw[arrowstyle] (score.east) -- (S.west);

% === SCALE + SOFTMAX ===
\node[op, minimum width=1.8cm] (scale) at (14, 1.2) {$\div \sqrt{d_k}$};
\draw[arrowstyle] (S.south) -- (scale.north);

\node[op, minimum width=1.8cm] (soft) at (14, 0) {softmax};
\draw[arrowstyle] (scale.south) -- (soft.north);

\node[tensor, fill=purple!15, minimum width=1.6cm] (A) at (14, -1.2) {$\mathbf{A}$};
\node[font=\tiny, below=0.05cm of A] {$N \times M$};
\draw[arrowstyle] (soft.south) -- (A.north);

% === A * V ===
\node[op, minimum width=1.6cm] (av) at (14, -2.6) {$\mathbf{A} \cdot \mathbf{V}$};
\draw[arrowstyle] (A.south) -- (av.north);
\draw[arrowstyle, red!50!black] (V.east) -- ++(3.5,0) |- (av.east);

% === OUTPUT ===
\node[tensor, fill=green!15, minimum width=2.2cm] (out) at (14, -4.0) {Output $\in \mathbb{R}^{N \times d_v}$};
\node[font=\scriptsize, below=0.08cm of out, green!40!black] {$512 \times d_v$ --- Compresso!};
\draw[arrowstyle, green!50!black] (av.south) -- (out.north);

% === EVIDENZIAZIONE COMPRESSIONE ===
\node[draw=green!60!black, fill=green!5, rounded corners=4pt, font=\scriptsize, align=center, text width=3.8cm]
    (comp) at (7, -4.0) {\textbf{Compressione!}\\$M = 50.176$ righe input\\$\downarrow$\\$N = 512$ righe output};
\draw[-{Stealth[length=2mm]}, thick, green!50!black, dashed] (comp.east) -- (out.west);

% === LEGENDA COLORI ===
\node[font=\scriptsize, red!60!black] at (0, -3.0) {{\color{red!60!black}\rule{0.5cm}{2pt}} percorso input};
\node[font=\scriptsize, blue!60!black] at (0, -3.5) {{\color{blue!60!black}\rule{0.5cm}{2pt}} percorso latenti};
\node[font=\scriptsize, green!50!black] at (0, -4.0) {{\color{green!50!black}\rule{0.5cm}{2pt}} output};

\end{tikzpicture}
\caption{Componente 4 --- Il modulo di Cross-Attention (dettaglio completo). Le Query provengono dai latenti (blu), le Key e Value dall'input (rosso). La matrice di attenzione $\mathbf{A}$ \`e rettangolare ($N \times M$), e l'output ha $N$ righe: l'informazione viene compressa da $M$ a $N$ elementi.}
\label{fig:cross_attention_detail}
\end{figure}

% Immagine dei tre passi della cross-attention dagli appunti
\begin{figure}[H]
\centering
\includegraphics[width=0.7\textwidth]{paper_figures/docx_cross_attention_steps.png}
\caption{I tre passi fondamentali della cross-attention nel Perceiver: (1) Calcolo di Q, K, V tramite proiezioni lineari (1D convolution); (2) Calcolo della mappa di attenzione $\mathbf{A} = \text{softmax}(\mathbf{Q}\mathbf{K}^\top)$ --- la matrice passa da quadrata $M \times M$ a rettangolare $N \times M$; (3) Calcolo delle feature di output $\mathbf{F} = \mathbf{A}\mathbf{V}$ con dimensione $N \times D$.}
\label{fig:docx_cross_attention_steps}
\end{figure}

La cross-attention \`e il momento in cui i latenti ``leggono'' l'input. Vediamo i tre passi chiave con le dimensioni concrete per ImageNet:

\begin{enumerate}[label=\textbf{Passo \arabic*:}, leftmargin=*]
    \item \textbf{Proiezione delle Query} ($\mathbf{z} \to \mathbf{Q}$): i latenti $\mathbf{z} \in \mathbb{R}^{512 \times 1024}$ vengono proiettati tramite $\mathbf{W}_Q \in \mathbb{R}^{1024 \times 64}$ per ottenere $\mathbf{Q} \in \mathbb{R}^{512 \times 64}$. Ogni latente ``chiede'' qualcosa all'input.
    \item \textbf{Proiezione delle Key} ($\mathbf{x} \to \mathbf{K}$): l'input $\mathbf{x} \in \mathbb{R}^{50.176 \times 261}$ viene proiettato tramite $\mathbf{W}_K \in \mathbb{R}^{261 \times 64}$ per ottenere $\mathbf{K} \in \mathbb{R}^{50.176 \times 64}$. Ogni pixel ``etichetta'' s\'e stesso.
    \item \textbf{Proiezione dei Value} ($\mathbf{x} \to \mathbf{V}$): analogamente, $\mathbf{W}_V \in \mathbb{R}^{261 \times 64}$ produce $\mathbf{V} \in \mathbb{R}^{50.176 \times 64}$. Ogni pixel ``prepara'' il proprio contenuto informativo.
    \item \textbf{Score} ($\mathbf{Q} \cdot \mathbf{K}^\top$): si ottiene una matrice $\mathbf{S} \in \mathbb{R}^{512 \times 50.176}$. Ogni elemento $S_{ij}$ misura quanto il latente $i$ \`e ``interessato'' al pixel $j$. Questa matrice \`e \textbf{rettangolare}, non quadrata.
    \item \textbf{Scaling e Softmax}: si divide per $\sqrt{d_k} = \sqrt{64} = 8$ e si applica softmax per riga, ottenendo i pesi di attenzione $\mathbf{A} \in \mathbb{R}^{512 \times 50.176}$ con righe che sommano a 1.
    \item \textbf{Output} ($\mathbf{A} \cdot \mathbf{V}$): il risultato \`e $\in \mathbb{R}^{512 \times 64}$. Ogni riga \`e una media pesata di tutti i 50.176 pixel, dove i pesi sono stati appresi dal modello.
\end{enumerate}

Il punto cruciale \`e che l'output ha $N = 512$ righe, non $M = 50.176$. L'informazione \`e stata \textbf{compressa}: i 50.176 pixel sono stati ``riassunti'' in 512 vettori latenti. Ogni latente ha selezionato e combinato le informazioni pi\`u rilevanti dall'intero input.


% ---- DIAGRAMMA 5: IL LATENT TRANSFORMER ----

\begin{figure}[H]
\centering
\begin{tikzpicture}[scale=0.9, transform shape,
    block/.style={draw, thick, rounded corners=3pt, minimum width=3.2cm, minimum height=0.7cm, font=\footnotesize},
    addcirc/.style={draw, thick, circle, minimum size=0.5cm, font=\footnotesize, fill=yellow!15},
    arrowstyle/.style={-{Stealth[length=2.5mm]}, thick}
]

% === Input ===
\node[block, fill=blue!12] (inp) at (0, 0) {$\mathbf{z} \in \mathbb{R}^{N \times D}$};
\node[font=\scriptsize, left=0.3cm of inp, blue!60!black] {Input};

% === LayerNorm 1 ===
\node[block, fill=gray!10] (ln1) at (0, 1.3) {LayerNorm};
\draw[arrowstyle] (inp.north) -- (ln1.south);

% === QKV split ===
\node[block, fill=orange!12, minimum width=4.5cm] (qkv) at (0, 2.6) {$\mathbf{W}_Q, \mathbf{W}_K, \mathbf{W}_V$ (tutti da $\mathbf{z}$)};
\draw[arrowstyle] (ln1.north) -- (qkv.south);

% === Multi-Head Self-Attention ===
\node[block, fill=purple!12, minimum width=4.2cm] (mhsa) at (0, 3.9) {Multi-Head Self-Attention};
\draw[arrowstyle] (qkv.north) -- (mhsa.south);

% Annotazione attenzione quadrata
\node[draw=purple!50!black, fill=purple!5, rounded corners=3pt, font=\scriptsize, align=center, text width=3.5cm]
    (sqnote) at (5.5, 3.9) {Matrice di attenzione:\\$N \times N = 512 \times 512$\\\textbf{quadrata!}};
\draw[-{Stealth[length=2mm]}, thick, purple!40!black, dashed] (sqnote.west) -- (mhsa.east);

% === Residual 1 ===
\node[addcirc] (add1) at (0, 5.2) {$+$};
\draw[arrowstyle] (mhsa.north) -- (add1.south);
% Skip connection 1
\draw[thick, blue!40!black] (inp.east) -- ++(2.8,0) |- (add1.east);
\node[font=\tiny, blue!40!black] at (3.5, 2.6) {residual};

% === LayerNorm 2 ===
\node[block, fill=gray!10] (ln2) at (0, 6.3) {LayerNorm};
\draw[arrowstyle] (add1.north) -- (ln2.south);

% === FFN ===
\node[block, fill=cyan!12, minimum width=4.0cm] (ffn) at (0, 7.6) {FFN: $D \to 4D \to D$ (GELU)};
\draw[arrowstyle] (ln2.north) -- (ffn.south);

% Annotazione FFN
\node[draw=cyan!50!black, fill=cyan!5, rounded corners=3pt, font=\scriptsize, align=center, text width=3.2cm]
    (ffnnote) at (5.5, 7.6) {$1024 \to 4096 \to 1024$\\con attivazione GELU};
\draw[-{Stealth[length=2mm]}, thick, cyan!40!black, dashed] (ffnnote.west) -- (ffn.east);

% === Residual 2 ===
\node[addcirc] (add2) at (0, 8.9) {$+$};
\draw[arrowstyle] (ffn.north) -- (add2.south);
% Skip connection 2
\draw[thick, blue!40!black] (add1.west) -- ++(-2.5,0) |- (add2.west);
\node[font=\tiny, blue!40!black] at (-3.2, 7.0) {residual};

% === Output ===
\node[block, fill=blue!12] (outp) at (0, 10.0) {$\mathbf{z}' \in \mathbb{R}^{N \times D}$};
\draw[arrowstyle] (add2.north) -- (outp.south);
\node[font=\scriptsize, left=0.3cm of outp, blue!60!black] {Output};

% === Brace per L blocchi ===
\draw[decorate, decoration={brace, amplitude=8pt, mirror}, very thick, red!60!black]
    (-3.8, 0.8) -- (-3.8, 9.3)
    node[midway, left=10pt, font=\small\bfseries, red!60!black, align=center] {$\times\, L$\\blocchi};

% === Differenza chiave ===
\node[draw=red!60!black, fill=red!5, rounded corners=4pt, font=\scriptsize, align=center, text width=4.2cm]
    (diff) at (5.5, 1.3) {\textbf{Differenza dalla cross-att.:}\\Q, K, V provengono \textbf{tutti}\\dai latenti (non dall'input!)};
\draw[-{Stealth[length=2mm]}, thick, red!50!black, dashed] (diff.west) -- (ln1.east);

\end{tikzpicture}
\caption{Componente 5 --- Il Latent Transformer. Un blocco Transformer standard (LayerNorm $\to$ Multi-Head Self-Attention $\to$ residual $\to$ LayerNorm $\to$ FFN $\to$ residual) applicato $L$ volte. La differenza chiave rispetto alla cross-attention \`e che Q, K, V provengono tutti dai latenti, e la matrice di attenzione \`e quadrata ($N \times N$).}
\label{fig:latent_transformer_detail}
\end{figure}

Dopo la cross-attention, i latenti hanno ``assorbito'' informazione dall'input. Ma questa informazione \`e ancora ``grezza'' --- ogni latente ha raccolto informazioni indipendentemente dagli altri. Il Latent Transformer permette ai latenti di comunicare tra di loro tramite \textbf{self-attention}.

Strutturalmente, il Latent Transformer \`e un blocco Transformer standard con architettura pre-norm:
\begin{enumerate}[label=(\roman*), leftmargin=*]
    \item \textbf{LayerNorm} normalizza le attivazioni (si veda la \hyperref[sec:layernorm]{Sezione~2.3}).
    \item \textbf{Multi-Head Self-Attention} permette ai 512 latenti di ``parlarsi''. La matrice di attenzione \`e $512 \times 512$ --- perfettamente gestibile in memoria.
    \item \textbf{Connessione residua} preserva le informazioni gi\`a acquisite (si veda la \hyperref[sec:residual]{Sezione~2.4}).
    \item \textbf{Feed-Forward Network (FFN)} con espansione $D \to 4D \to D$ e attivazione GELU elabora ogni latente indipendentemente, arricchendo la rappresentazione.
    \item \textbf{Seconda connessione residua} chiude il blocco.
\end{enumerate}

Questo blocco viene ripetuto $L$ volte (tipicamente $L = 6$ o $L = 8$). La profondit\`a del Latent Transformer \`e cruciale: essendo applicato solo ai $N = 512$ latenti (e non ai $M = 50.176$ pixel), \`e \textbf{computazionalmente economico} e pu\`o essere reso molto profondo senza problemi di memoria.


% ---- DIAGRAMMA 6: PIPELINE COMPLETA CON WEIGHT SHARING ----

\begin{figure}[H]
\centering
\begin{tikzpicture}[scale=0.55, transform shape,
    block/.style={draw, thick, rounded corners=3pt, minimum height=0.9cm, font=\footnotesize, minimum width=1.3cm},
    arrowstyle/.style={-{Stealth[length=2.5mm]}, thick}
]

% === z_0 ===
\node[block, fill=blue!12] (z0) at (0, 2) {$\mathbf{z}_0$};
\node[font=\tiny, above=0.05cm of z0] {$512 \times 1024$};

% === Iterazione 1 (pesi unici) ===
\node[block, fill=orange!15, minimum width=1.5cm] (ca1) at (2.5, 2) {CA$_1$};
\node[block, fill=cyan!15, minimum width=1.5cm] (lt1) at (4.5, 2) {LT$_1$};
\draw[arrowstyle] (z0.east) -- (ca1.west);
\draw[arrowstyle] (ca1.east) -- (lt1.west);

\node[block, fill=blue!12] (z1) at (6.5, 2) {$\mathbf{z}_1$};
\draw[arrowstyle] (lt1.east) -- (z1.west);

% === Iterazione 2 (pesi condivisi) ===
\node[block, fill=orange!15, minimum width=1.5cm] (ca2) at (9.0, 2) {CA$_2$};
\node[block, fill=cyan!15, minimum width=1.5cm] (lt2) at (11.0, 2) {LT$_2$};
\draw[arrowstyle] (z1.east) -- (ca2.west);
\draw[arrowstyle] (ca2.east) -- (lt2.west);

\node[block, fill=blue!12] (z2) at (13.0, 2) {$\mathbf{z}_2$};
\draw[arrowstyle] (lt2.east) -- (z2.west);

% === Puntini ===
\node[font=\large] at (14.3, 2) {$\cdots$};

% === Iterazione T ===
\node[block, fill=orange!15, minimum width=1.5cm] (caT) at (16.0, 2) {CA$_T$};
\node[block, fill=cyan!15, minimum width=1.5cm] (ltT) at (18.0, 2) {LT$_T$};
\draw[arrowstyle] (14.7, 2) -- (caT.west);
\draw[arrowstyle] (caT.east) -- (ltT.west);

\node[block, fill=blue!12] (zT) at (20.0, 2) {$\mathbf{z}_T$};
\draw[arrowstyle] (ltT.east) -- (zT.west);

% === Input x alimenta ogni CA ===
\node[block, fill=red!12, minimum width=2.8cm] (x) at (9.0, -1.0) {$\mathbf{x} \in \mathbb{R}^{M \times C}$};
\node[font=\scriptsize, below=0.05cm of x, red!50!black] {Input (stesso per ogni iterazione)};

\draw[arrowstyle, red!40!black] (x.north) -- ++(0, 0.5) -| (ca1.south);
\draw[arrowstyle, red!40!black] (x.north) -- ++(0, 0.5) -| (ca2.south);
\draw[arrowstyle, red!40!black] (x.north) -- ++(0, 0.5) -| (caT.south);

% === Brace pesi unici ===
\draw[decorate, decoration={brace, amplitude=5pt}, thick, green!60!black]
    (1.7, 3.3) -- (5.3, 3.3)
    node[midway, above=7pt, font=\scriptsize\bfseries, green!60!black] {Pesi unici};

% === Brace pesi condivisi ===
\draw[decorate, decoration={brace, amplitude=5pt}, thick, blue!60!black]
    (8.2, 3.3) -- (18.8, 3.3)
    node[midway, above=7pt, font=\scriptsize\bfseries, blue!60!black] {Pesi condivisi (weight sharing)};

% === Classificazione ===
\draw[arrowstyle] (zT.east) -- ++(0.6, 0);

\node[block, fill=green!10, minimum width=1.6cm] (pool) at (22.0, 2) {Avg Pool};
\draw[arrowstyle] (20.6, 2) -- (pool.west);

\node[font=\scriptsize, below=0.05cm of pool] {$\mathbf{h} \in \mathbb{R}^D$};

\node[block, fill=green!15, minimum width=1.6cm] (cls) at (24.3, 2) {$\mathbf{W}_{\text{cls}}$};
\draw[arrowstyle] (pool.east) -- (cls.west);

\node[block, fill=green!20, minimum width=1.3cm] (soft) at (26.3, 2) {softmax};
\draw[arrowstyle] (cls.east) -- (soft.west);

\node[font=\scriptsize, above=0.05cm of cls] {$D \to 1000$};
\node[font=\scriptsize, right=0.1cm of soft] {$\mathbf{p} \in \mathbb{R}^{1000}$};

% === Analogia RNN ===
\node[draw=gray, fill=gray!5, rounded corners=4pt, font=\scriptsize, align=center, text width=6cm]
    (rnn) at (9.0, -3.2) {Analogia: il weight sharing rende il Perceiver\\simile a una \textbf{RNN srotolata}, dove lo ``stato nascosto''\\sono i latenti e l'``input a ogni passo'' \`e sempre $\mathbf{x}$};

\end{tikzpicture}
\caption{Componente 6 --- Pipeline completa con weight sharing. L'array latente $\mathbf{z}_0$ viene raffinato iterativamente: ogni iterazione consiste in un modulo di cross-attention (CA) seguito da un Latent Transformer (LT). L'input $\mathbf{x}$ alimenta ogni CA. Dalla seconda iterazione, i pesi sono condivisi. L'output finale passa per average pooling e una testa di classificazione lineare.}
\label{fig:pipeline_complete}
\end{figure}

Osserviamo ora il flusso completo dell'architettura. Il Perceiver applica $T$ iterazioni di cross-attention + Latent Transformer, dove ogni iterazione raffina progressivamente la rappresentazione latente:

\begin{enumerate}[label=\textbf{Iterazione \arabic*:}, leftmargin=*]
    \item $\mathbf{z}_0 \xrightarrow{\text{CA}_1(\mathbf{z}_0, \mathbf{x})} \xrightarrow{\text{LT}_1} \mathbf{z}_1$ --- prima lettura dell'input, con \textbf{pesi unici}.
    \item $\mathbf{z}_1 \xrightarrow{\text{CA}_2(\mathbf{z}_1, \mathbf{x})} \xrightarrow{\text{LT}_2} \mathbf{z}_2$ --- seconda lettura, con \textbf{pesi condivisi} con le iterazioni successive.
    \item[] $\vdots$
    \item[$T.$] $\mathbf{z}_{T-1} \xrightarrow{\text{CA}_T(\mathbf{z}_{T-1}, \mathbf{x})} \xrightarrow{\text{LT}_T} \mathbf{z}_T$ --- ultima lettura.
\end{enumerate}

Il \textbf{weight sharing} (condivisione dei pesi) a partire dalla seconda iterazione ha due vantaggi fondamentali:
\begin{itemize}[leftmargin=*]
    \item \textbf{Riduzione parametri}: invece di avere $T$ copie distinte di cross-attention e Latent Transformer, se ne hanno solo 2 (una per la prima iterazione, una condivisa per le successive). Questo riduce drasticamente il numero di parametri del modello.
    \item \textbf{Regolarizzazione}: il weight sharing agisce come regolarizzatore, impedendo l'overfitting e migliorando la generalizzazione. Sperimentalmente, gli autori mostrano che il weight sharing \textbf{migliora} le performance (non le peggiora!), il che suggerisce che il modello senza sharing era sovra-parametrizzato.
\end{itemize}

L'analogia con le \textbf{RNN} \`e illuminante: il Perceiver con weight sharing \`e equivalente a una RNN srotolata nel tempo, dove lo ``stato nascosto'' sono i latenti $\mathbf{z}_t$ e l'``input a ogni passo temporale'' \`e sempre lo stesso byte array $\mathbf{x}$. A differenza delle RNN classiche, per\`o, il Perceiver utilizza l'attenzione (non operazioni ricorrenti punto a punto), il che evita il problema del \textit{vanishing gradient}.

Dopo l'ultima iterazione, il vettore $\mathbf{z}_T \in \mathbb{R}^{512 \times 1024}$ viene compresso tramite \textbf{average pooling} lungo la dimensione $N$, ottenendo $\mathbf{h} \in \mathbb{R}^{1024}$, e infine proiettato tramite un layer lineare $\mathbf{W}_{\text{cls}} \in \mathbb{R}^{1024 \times 1000}$ seguito da softmax per produrre la distribuzione di probabilit\`a sulle 1000 classi di ImageNet.


\subsection{Il modulo di Cross-Attention} \label{sec:cross_attention}

Eccoci al componente pi\`u importante del Perceiver. Il modulo di cross-attention \`e il meccanismo che permette all'array latente di ``leggere'' l'input. La differenza con la self-attention (dove Q, K, V provengono dalla stessa sequenza, come spiegato nella \hyperref[sec:attention]{Sezione~2.2}) \`e cruciale: nella cross-attention le \textbf{query} provengono dai latenti e le \textbf{chiavi/valori} dall'input.

Perch\'e questa asimmetria \`e cos\`i importante? Perch\'e ci sono \textbf{poche query} ($N = 512$) ma \textbf{molte chiavi/valori} ($M = 50.176$). Questo significa che la matrice di attenzione \`e rettangolare ($512 \times 50.176$) anzich\'e quadrata ($50.176 \times 50.176$), risparmiando un fattore $M/N \approx 100$ in memoria e computazione.

\begin{formulabox}[Cross-Attention del Perceiver]
\begin{align}
    \mathbf{Q} &= \mathbf{z}\, \mathbf{W}_Q + \mathbf{b}_Q &&\in \mathbb{R}^{N \times d_k} \label{eq:ca_Q}\\[4pt]
    \mathbf{K} &= \mathbf{x}\, \mathbf{W}_K + \mathbf{b}_K &&\in \mathbb{R}^{M \times d_k} \label{eq:ca_K}\\[4pt]
    \mathbf{V} &= \mathbf{x}\, \mathbf{W}_V + \mathbf{b}_V &&\in \mathbb{R}^{M \times d_v} \label{eq:ca_V}\\[4pt]
    \mathbf{A} &= \text{softmax}\!\left(\frac{\mathbf{Q}\,\mathbf{K}^\top}{\sqrt{d_k}}\right) &&\in \mathbb{R}^{N \times M} \label{eq:ca_A}\\[4pt]
    \text{CrossAttn}(\mathbf{z}, \mathbf{x}) &= \mathbf{A}\,\mathbf{V} &&\in \mathbb{R}^{N \times d_v} \label{eq:ca_out}
\end{align}
dove $\mathbf{W}_Q \in \mathbb{R}^{D \times d_k}$, $\mathbf{W}_K \in \mathbb{R}^{C \times d_k}$, $\mathbf{W}_V \in \mathbb{R}^{C \times d_v}$.
\end{formulabox}

\begin{notabox}[Osservazione cruciale sulle dimensioni]
La matrice di attenzione $\mathbf{A} \in \mathbb{R}^{N \times M}$ ha dimensione $N \times M$ (non $M \times M$ come nella self-attention standard). Poich\'e $N \ll M$, questa matrice \`e molto pi\`u piccola. L'output ha dimensione $N \times d_v$: l'informazione viene ``compressa'' dall'input ($M$ elementi) ai latenti ($N$ elementi).

\textbf{Direzione della compressione}: l'input ha $M$ righe e l'output ha $N$ righe, con $N \ll M$. La cross-attention ``seleziona e comprime'' le $M$ righe dell'input in $N$ righe di output. Ogni riga di output \`e una media pesata di tutte le $M$ righe dell'input, dove i pesi dipendono dalla compatibilit\`a query-chiave.
\end{notabox}

% TikZ: Cross-attention data flow diagram
\begin{figure}[H]
\centering
\begin{tikzpicture}[
    tensor/.style={draw, thick, rounded corners=2pt, minimum height=0.7cm, font=\footnotesize},
    op/.style={draw, circle, fill=gray!15, minimum size=0.7cm, font=\footnotesize},
    arrowstyle/.style={-{Stealth[length=2.5mm]}, thick},
    >=Stealth,
    scale=0.9, transform shape
]
% Latent array z
\node[tensor, fill=blue!15, minimum width=2.2cm] (z) at (0, 4) {$\mathbf{z} \in \mathbb{R}^{N \times D}$};
\node[font=\scriptsize, above=0.1cm of z] {Latenti (N=512, D=1024)};

% Input x
\node[tensor, fill=red!15, minimum width=3cm] (x) at (0, 0) {$\mathbf{x} \in \mathbb{R}^{M \times C}$};
\node[font=\scriptsize, below=0.1cm of x] {Byte array (M=50176, C=261)};

% W_Q projection
\node[tensor, fill=orange!15, minimum width=1.8cm] (wq) at (4, 4) {$\mathbf{W}_Q$};
\node[font=\scriptsize, above=0.1cm of wq] {$D \times d_k$};

% W_K projection
\node[tensor, fill=orange!15, minimum width=1.8cm] (wk) at (4, 1.5) {$\mathbf{W}_K$};
\node[font=\scriptsize, above=0.1cm of wk] {$C \times d_k$};

% W_V projection
\node[tensor, fill=orange!15, minimum width=1.8cm] (wv) at (4, -1) {$\mathbf{W}_V$};
\node[font=\scriptsize, below=0.1cm of wv] {$C \times d_v$};

% Q, K, V results
\node[tensor, fill=blue!10, minimum width=2cm] (Q) at (8, 4) {$\mathbf{Q} \in \mathbb{R}^{N \times d_k}$};
\node[tensor, fill=red!10, minimum width=2cm] (K) at (8, 1.5) {$\mathbf{K} \in \mathbb{R}^{M \times d_k}$};
\node[tensor, fill=red!10, minimum width=2cm] (V) at (8, -1) {$\mathbf{V} \in \mathbb{R}^{M \times d_v}$};

% Dot product
\node[op] (dot) at (11, 2.75) {$\cdot$};
\node[font=\scriptsize, right=0.1cm of dot] {$/\sqrt{d_k}$};

% Softmax
\node[tensor, fill=yellow!20, minimum width=1.5cm] (sm) at (13.5, 2.75) {softmax};

% Attention matrix
\node[tensor, fill=green!15, minimum width=2.2cm] (A) at (13.5, 1) {$\mathbf{A} \in \mathbb{R}^{N \times M}$};

% Final matmul
\node[op] (matmul) at (11, -1) {$\times$};

% Output
\node[tensor, fill=green!20, minimum width=2.5cm] (out) at (13.5, -1) {Out $\in \mathbb{R}^{N \times d_v}$};

% Arrows
\draw[arrowstyle] (z) -- (wq);
\draw[arrowstyle] (x) -- (wk);
\draw[arrowstyle] (x) -- (wv);
\draw[arrowstyle] (wq) -- (Q);
\draw[arrowstyle] (wk) -- (K);
\draw[arrowstyle] (wv) -- (V);
\draw[arrowstyle] (Q.east) -| (dot);
\draw[arrowstyle] (K.east) -| (dot);
\draw[arrowstyle] (dot) -- (sm);
\draw[arrowstyle] (sm) -- (A);
\draw[arrowstyle] (A.south) -| (matmul);
\draw[arrowstyle] (V) -- (matmul);
\draw[arrowstyle] (matmul) -- (out);

% Key annotation
\node[font=\scriptsize, text=blue!70!black] at (2, 5) {Q dai latenti};
\node[font=\scriptsize, text=red!70!black] at (2, -1.8) {K, V dall'input};

\end{tikzpicture}
\caption{Flusso di dati nella cross-attention del Perceiver. Le query (Q) provengono dai latenti, le chiavi (K) e i valori (V) dall'input. La matrice di attenzione risultante ha dimensione $N \times M$ (non $M \times M$), garantendo complessit\`a lineare in $M$.}
\label{fig:cross_attention_flow}
\end{figure}


\begin{notabox}[In parole semplici: la cross-attention del Perceiver]
Immaginiamo 512 studenti (i latenti) in una biblioteca con 50.176 libri (i pixel). Ogni studente ha una domanda diversa (Q): ``dove sono i bordi?'', ``che colori ci sono in alto?'', ``c'\`e un volto?''. Ogni libro ha un'etichetta (K) che descrive il suo contenuto, e il contenuto vero (V). Ogni studente scorre tutte le etichette, trova i libri pi\`u rilevanti per la sua domanda, e prende appunti (media pesata dei V). Alla fine, i 512 studenti hanno ciascuno un riassunto personalizzato della biblioteca.
\end{notabox}

\subsection{Il Latent Transformer} \label{sec:latent_transformer}

Ora che i latenti hanno estratto informazione dall'input tramite la cross-attention, dobbiamo chiederci: cosa succede dopo? I latenti hanno estratto informazione dall'input, ma questa informazione \`e ancora ``grezza''. Serve un passo ulteriore: i latenti devono ``parlare tra loro'' per integrare le informazioni.

Dopo ogni modulo di cross-attention, l'array latente aggiornato viene elaborato da una \textbf{torre di self-attention} (Latent Transformer). Questo \`e essenzialmente un Transformer standard (si veda l'\hyperref[app:transformer]{Appendice~H}) che opera esclusivamente sull'array latente $\mathbf{z} \in \mathbb{R}^{N \times D}$.

\begin{formulabox}[Blocco del Latent Transformer]
Per ogni layer $\ell = 1, \ldots, L$ del Latent Transformer:
\begin{align}
    \mathbf{z}' &= \mathbf{z} + \text{MultiHeadSelfAttention}\!\left(\text{LayerNorm}(\mathbf{z})\right) \label{eq:lt_sa} \\[4pt]
    \mathbf{z}'' &= \mathbf{z}' + \text{FFN}\!\left(\text{LayerNorm}(\mathbf{z}')\right) \label{eq:lt_ffn}
\end{align}
dove il Feed-Forward Network (FFN) ha la struttura standard:
\[
    \text{FFN}(\mathbf{h}) = \text{GELU}\!\left(\mathbf{h}\,\mathbf{W}_1 + \mathbf{b}_1\right)\,\mathbf{W}_2 + \mathbf{b}_2
\]
con $\mathbf{W}_1 \in \mathbb{R}^{D \times D_{\text{ff}}}$, $\mathbf{W}_2 \in \mathbb{R}^{D_{\text{ff}} \times D}$ e tipicamente $D_{\text{ff}} = 4D$.
\end{formulabox}

% TikZ: Latent Transformer block detail
\begin{figure}[H]
\centering
\begin{tikzpicture}[
    block/.style={draw, rounded corners=3pt, fill=#1, minimum width=3.5cm, minimum height=0.6cm, font=\small},
    block/.default={orange!15},
    arrowstyle/.style={-{Stealth[length=2.5mm]}, thick},
    >=Stealth,
    scale=0.85, transform shape
]

% Input
\node[font=\small\bfseries] (zin) at (0, 0) {$\mathbf{z}$ input};

% LN1
\node[block=gray!15] (ln1) at (0, 1.2) {LayerNorm};
\draw[arrowstyle] (zin) -- (ln1);

% Self-Attention
\node[block=orange!20] (sa) at (0, 2.6) {Multi-Head Self-Attention};
\draw[arrowstyle] (ln1) -- (sa);

% Add1
\node[circle, draw, fill=green!15, minimum size=0.5cm, font=\small] (add1) at (0, 4.0) {$+$};
\draw[arrowstyle] (sa) -- (add1);

% Skip connection 1
\draw[arrowstyle, blue!60!black] (-2.5, 0) -- (-2.5, 4.0) -- (add1.west);
\node[font=\scriptsize, blue!60!black, rotate=90] at (-2.8, 2) {skip connection};

% LN2
\node[block=gray!15] (ln2) at (0, 5.2) {LayerNorm};
\draw[arrowstyle] (add1) -- (ln2);

% MLP expansion
\node[block=yellow!15] (mlp1) at (0, 6.4) {Linear $D \to 4D$ + GELU};
\draw[arrowstyle] (ln2) -- (mlp1);

% MLP compression
\node[block=yellow!15] (mlp2) at (0, 7.6) {Linear $4D \to D$};
\draw[arrowstyle] (mlp1) -- (mlp2);

% Add2
\node[circle, draw, fill=green!15, minimum size=0.5cm, font=\small] (add2) at (0, 8.8) {$+$};
\draw[arrowstyle] (mlp2) -- (add2);

% Skip connection 2
\draw[arrowstyle, blue!60!black] (2.8, 4.0) -- (2.8, 8.8) -- (add2.east);
\node[font=\scriptsize, blue!60!black, rotate=90] at (3.1, 6.4) {skip connection};

% Line connecting add1 output to skip2 start
\draw[blue!60!black, thick] (add1.east) -- (2.8, 4.0);

% Output
\node[font=\small\bfseries] (zout) at (0, 9.8) {$\mathbf{z}$ output};
\draw[arrowstyle] (add2) -- (zout);

% Dimension annotations
\node[font=\scriptsize, gray] at (4.5, 0) {$\mathbb{R}^{512 \times 1024}$};
\node[font=\scriptsize, gray] at (4.5, 2.6) {$\mathbb{R}^{512 \times 1024}$};
\node[font=\scriptsize, gray] at (4.5, 6.4) {$\mathbb{R}^{512 \times 4096}$};
\node[font=\scriptsize, gray] at (4.5, 7.6) {$\mathbb{R}^{512 \times 1024}$};
\node[font=\scriptsize, gray] at (4.5, 9.8) {$\mathbb{R}^{512 \times 1024}$};

% Brace for MLP
\draw[decorate, decoration={brace, amplitude=5pt, mirror}, thick, orange!70!black]
    (5.5, 6.0) -- (5.5, 8.0) node[midway, right=6pt, font=\scriptsize, orange!70!black] {FFN (MLP)};

\end{tikzpicture}
\caption{Dettaglio di un singolo blocco del Latent Transformer. Schema pre-norm: LayerNorm $\to$ Self-Attention $\to$ residual $\to$ LayerNorm $\to$ MLP (espansione $D \to 4D$ + GELU + compressione $4D \to D$) $\to$ residual. Le dimensioni restano $\mathbb{R}^{512 \times 1024}$ in ingresso e uscita.}
\label{fig:latent_transformer_block}
\end{figure}


\begin{notabox}[In parole semplici: il Latent Transformer]
Continuando l'analogia della biblioteca: dopo che i 512 studenti hanno preso i loro appunti dai libri (cross-attention), si riuniscono per una \textbf{riunione di gruppo} dove condividono le informazioni. Lo studente che ha trovato ``zampe'' parla con quello che ha trovato ``orecchie appuntite'', e insieme capiscono ``\`e un gatto!''. Questo scambio di informazioni avviene 6 volte ($L=6$), permettendo relazioni sempre pi\`u complesse.
\end{notabox}

\subsection{Cross-Attention iterativa e Weight Sharing} \label{sec:iterative}

L'architettura completa consiste in $T$ ripetizioni della sequenza Cross-Attention $\to$ Latent Transformer. Ma perch\'e ripetere? Perch\'e una sola ``lettura'' dell'immagine non basta. \`E come rileggere un testo difficile: ad ogni rilettura si colgono dettagli nuovi, perch\'e ci si avvicina al testo con una comprensione pi\`u profonda.

I pesi della prima cross-attention sono unici; le iterazioni $2, \ldots, T$ condividono gli stessi pesi.

\begin{formulabox}[Weight sharing formale]
\begin{align}
    \theta_{\text{CA}}^{(1)} &\neq \theta_{\text{CA}}^{(t)} \quad \text{per } t \geq 2 \quad \text{(prima cross-attention unica)} \\
    \theta_{\text{CA}}^{(t)} &= \theta_{\text{CA}}^{(2)} \quad \text{per ogni } t \geq 2 \quad \text{(pesi condivisi)} \\
    \theta_{\text{LT}}^{(t)} &= \theta_{\text{LT}}^{(1)} \quad \text{per ogni } t \geq 1 \quad \text{(pesi condivisi)}
\end{align}
\end{formulabox}

Il weight sharing riduce i parametri da $\sim$326M a $\sim$45M e \textbf{migliora} le prestazioni di validazione (da 72.9\% senza WS a 78.0\% con WS su ImageNet), agendo come una potente regolarizzazione.

\begin{notabox}[In parole semplici: il weight sharing]
Il weight sharing \`e come \textbf{rileggere lo stesso libro pi\`u volte con lo stesso cervello}. Non compri un cervello nuovo ad ogni rilettura --- usi sempre lo stesso. Questo ``vincolo'' \`e in realt\`a un vantaggio: il cervello deve imparare a fare una buona lettura \emph{in generale}, non una lettura specializzata per la prima volta e un'altra per la seconda. Il risultato \`e che il modello generalizza meglio, con 7 volte meno parametri.
\end{notabox}

% TikZ: Unrolled architecture with weight sharing annotations
\begin{figure}[H]
\centering
\begin{tikzpicture}[
    cabox/.style={draw, rounded corners=3pt, fill=orange!20, minimum width=1.6cm, minimum height=0.6cm, font=\scriptsize},
    ltbox/.style={draw, rounded corners=3pt, fill=cyan!20, minimum width=1.6cm, minimum height=0.6cm, font=\scriptsize},
    arrowstyle/.style={-{Stealth[length=2mm]}, thick},
    >=Stealth,
    scale=0.85, transform shape
]

% Input
\node[draw, fill=red!15, rounded corners, minimum width=1.5cm, minimum height=0.5cm, font=\scriptsize] (input) at (0, -1.5) {$\mathbf{x} \in \mathbb{R}^{M \times C}$};

% Iteration 1
\node[cabox] (ca1) at (0, 0) {CA$_1$};
\node[ltbox] (lt1) at (0, 1.2) {LT $\times L$};
\node[font=\tiny, red!70!black, right=0.2cm of ca1] {$\theta_{\text{CA}}^{(1)}$};
\node[font=\tiny, blue!70!black, right=0.2cm of lt1] {$\theta_{\text{LT}}$};

% Iteration 2
\node[cabox] (ca2) at (2.5, 0) {CA$_2$};
\node[ltbox] (lt2) at (2.5, 1.2) {LT $\times L$};
\node[font=\tiny, purple!70!black, right=0.2cm of ca2] {$\theta_{\text{CA}}^{(s)}$};
\node[font=\tiny, blue!70!black, right=0.2cm of lt2] {$\theta_{\text{LT}}$};

% Iteration 3
\node[cabox] (ca3) at (5, 0) {CA$_3$};
\node[ltbox] (lt3) at (5, 1.2) {LT $\times L$};
\node[font=\tiny, purple!70!black, right=0.2cm of ca3] {$\theta_{\text{CA}}^{(s)}$};
\node[font=\tiny, blue!70!black, right=0.2cm of lt3] {$\theta_{\text{LT}}$};

% Dots
\node[font=\large] at (7, 0.6) {$\cdots$};

% Iteration T
\node[cabox] (caT) at (9, 0) {CA$_T$};
\node[ltbox] (ltT) at (9, 1.2) {LT $\times L$};
\node[font=\tiny, purple!70!black, right=0.2cm of caT] {$\theta_{\text{CA}}^{(s)}$};
\node[font=\tiny, blue!70!black, right=0.2cm of ltT] {$\theta_{\text{LT}}$};

% Output
\node[draw, fill=green!15, rounded corners, minimum width=1.2cm, minimum height=0.5cm, font=\scriptsize] (out) at (11.5, 0.6) {Pool+Class};

% Arrows
\draw[arrowstyle] (input) -- (ca1);
\draw[arrowstyle] (ca1) -- (lt1);
\draw[arrowstyle] (lt1.east) -| ($(ca2.north west)+(-0.2,0.3)$) -- (ca2.north);
\draw[arrowstyle] (ca2) -- (lt2);
\draw[arrowstyle] (lt2.east) -| ($(ca3.north west)+(-0.2,0.3)$) -- (ca3.north);
\draw[arrowstyle] (ca3) -- (lt3);
\draw[arrowstyle] (caT) -- (ltT);
\draw[arrowstyle] (ltT) -- (out);

% Input feeding arrows
\draw[arrowstyle, gray, dashed] (input.east) -| (ca2.south);
\draw[arrowstyle, gray, dashed] (input.east) -| (ca3.south);
\draw[arrowstyle, gray, dashed] (input.east) -| (caT.south);

% Weight sharing braces
\draw[decorate, decoration={brace, amplitude=5pt}, thick, red!60!black]
    (-0.8, -0.5) -- (-0.8, 0.5) node[midway, left=6pt, font=\tiny, red!60!black, align=center] {Pesi\\unici};

\draw[decorate, decoration={brace, amplitude=5pt}, thick, purple!60!black]
    (1.7, -0.5) -- (10, -0.5) node[midway, below=6pt, font=\scriptsize, purple!60!black] {$\theta_{\text{CA}}^{(s)}$ condivisi tra iterazioni 2..T};

\draw[decorate, decoration={brace, amplitude=5pt, mirror}, thick, blue!60!black]
    (-0.8, 2.0) -- (10, 2.0) node[midway, above=6pt, font=\scriptsize, blue!60!black] {$\theta_{\text{LT}}$ condivisi tra TUTTE le iterazioni 1..T};

\end{tikzpicture}
\caption{Architettura srotolata del Perceiver con $T=8$ iterazioni (mostrate solo 4 per chiarezza). La prima cross-attention (CA$_1$) ha pesi unici (rosso); le successive (CA$_2$\ldots CA$_T$) condividono gli stessi pesi (viola). Il Latent Transformer (LT) condivide i pesi in tutte le iterazioni (blu). L'input $\mathbf{x}$ viene ri-letto ad ogni iterazione (frecce grigie tratteggiate).}
\label{fig:unrolled_architecture}
\end{figure}

% TikZ: Weight sharing parameter count diagram
\begin{figure}[H]
\centering
\begin{tikzpicture}[
    box/.style={draw, rounded corners=3pt, minimum width=2.8cm, minimum height=1.0cm, font=\small, text centered},
    arrowstyle/.style={-{Stealth[length=2.5mm]}, thick},
    >=Stealth
]

% Without WS
\node[box, fill=red!15] (nows) at (0, 0) {\begin{tabular}{c}\textbf{Senza WS}\\$\sim$326M param.\end{tabular}};

% With WS
\node[box, fill=green!15] (ws) at (6, 0) {\begin{tabular}{c}\textbf{Con WS}\\$\sim$45M param.\end{tabular}};

% Arrow
\draw[arrowstyle, very thick] (nows) -- (ws) node[midway, above, font=\small\bfseries] {$\div\; 7.3\times$};

% Accuracy annotations
\node[font=\small, red!70!black] at (0, -1.5) {Val accuracy: 72.9\%};
\node[font=\small, green!60!black] at (6, -1.5) {Val accuracy: 78.0\%};

\end{tikzpicture}
\caption{Impatto del weight sharing sui parametri e sulle prestazioni. Controintuitivamente, meno parametri portano a migliori prestazioni, confermando che il WS agisce come una forma di regolarizzazione.}
\label{fig:ws_parameter_count}
\end{figure}

\begin{examplebox}[Calcolo dettagliato dei parametri: con e senza weight sharing]
Vediamo il conteggio esatto dei parametri per ogni componente del Perceiver su ImageNet ($N = 512$, $D = 1024$, $D_{\text{ff}} = 4096$, $C_{\text{embed}} = 512$, $H = 8$, $L = 6$, $T = 8$).

\textbf{Componente 1 --- Cross-Attention (1 modulo):}
\begin{itemize}
\item $\mathbf{W}_Q \in \mathbb{R}^{1024 \times 1024}$: $1{,}048{,}576$ pesi + $1{,}024$ bias
\item $\mathbf{W}_K \in \mathbb{R}^{512 \times 1024}$: $524{,}288$ pesi + $1{,}024$ bias
\item $\mathbf{W}_V \in \mathbb{R}^{512 \times 1024}$: $524{,}288$ pesi + $1{,}024$ bias
\item $\mathbf{W}_O \in \mathbb{R}^{1024 \times 1024}$: $1{,}048{,}576$ pesi + $1{,}024$ bias
\item LayerNorm latenti: $2 \times 1024 = 2{,}048$ ($\gamma + \beta$)
\item LayerNorm input: $2 \times 512 = 1{,}024$
\item MLP post-CA: $1024 \times 4096 + 4096 + 4096 \times 1024 + 1024 = 8{,}394{,}752$
\item LayerNorm pre-MLP: $2 \times 1024 = 2{,}048$
\item \textbf{Subtotale CA: $\approx 11.6$M parametri}
\end{itemize}

\textbf{Componente 2 --- Latent Transformer (1 blocco SA):}
\begin{itemize}
\item $\mathbf{W}_Q, \mathbf{W}_K, \mathbf{W}_V, \mathbf{W}_O$: $4 \times (1024 \times 1024 + 1024) = 4{,}198{,}400$
\item MLP: $1024 \times 4096 + 4096 + 4096 \times 1024 + 1024 = 8{,}394{,}752$
\item LayerNorm ($\times 2$): $2 \times 2 \times 1024 = 4{,}096$
\item \textbf{Subtotale 1 blocco: $\approx 12.6$M}
\item \textbf{Subtotale torre $L{=}6$: $6 \times 12.6 \approx 75.5$M}
\end{itemize}

\textbf{Componente 3 --- Array latente $\mathbf{z}_0$}: $512 \times 1024 = 524{,}288 \approx 0.5$M

\textbf{Componente 4 --- Embedding}: $261 \times 512 + 512 = 134{,}144 \approx 0.1$M

\textbf{Componente 5 --- Classificatore}: $1024 \times 1000 + 1000 = 1{,}025{,}000 \approx 1.0$M

\medskip

\textbf{Totale SENZA weight sharing} ($T = 8$ iterazioni indipendenti):
\begin{align*}
&\underbrace{11.6\text{M}}_{\text{CA}_1} + \underbrace{7 \times 11.6\text{M}}_{\text{CA}_{2\ldots 8}} + \underbrace{8 \times 75.5\text{M}}_{\text{8 torri LT}} + \underbrace{0.5\text{M}}_{\mathbf{z}_0} + \underbrace{0.1\text{M}}_{\text{embed}} + \underbrace{1.0\text{M}}_{\text{class}} \\
&= 11.6 + 81.2 + 604.0 + 1.6 \approx \textbf{698M parametri}
\end{align*}

\textbf{Totale CON weight sharing}:
\begin{align*}
&\underbrace{11.6\text{M}}_{\text{CA}_1\text{ (unica)}} + \underbrace{11.6\text{M}}_{\text{CA}_{\text{shared}}\text{ (1 sola copia)}} + \underbrace{75.5\text{M}}_{\text{LT (1 sola torre)}} + \underbrace{0.5\text{M}}_{\mathbf{z}_0} + \underbrace{0.1\text{M}}_{\text{embed}} + \underbrace{1.0\text{M}}_{\text{class}} \\
&= 11.6 + 11.6 + 75.5 + 1.6 \approx \textbf{100M parametri}
\end{align*}

\textbf{Fattore di riduzione}: $698 / 100 \approx 7\times$ meno parametri con weight sharing.

\textbf{Nota}: la configurazione reale del paper usa $D_{\text{latent}} < D$ in alcune varianti, ottenendo $\sim$45M con WS. I numeri sopra si riferiscono alla configurazione con $d_k = D = 1024$.
\end{examplebox}

% TikZ: Comparison of attention matrix sizes
\begin{figure}[H]
\centering
\begin{tikzpicture}[
    scale=0.75, transform shape
]

% Transformer full self-attention
\node[font=\small\bfseries] at (0, 4.5) {Transformer};
\node[font=\small\bfseries] at (0, 4.0) {$50{,}176 \times 50{,}176$};
\fill[red!30] (-2.5, -2.5) rectangle (2.5, 2.5);
\draw[red!70!black, thick] (-2.5, -2.5) rectangle (2.5, 2.5);
\node[font=\scriptsize, red!70!black] at (0, 0) {$\sim$2.5 miliardi};
\node[font=\scriptsize, red!70!black] at (0, -0.5) {$\approx 10$ GB};
\node[font=\scriptsize, red!70!black, fill=red!10, draw=red!50, rounded corners] at (0, -3.2) {IMPOSSIBILE};

% ViT self-attention
\node[font=\small\bfseries] at (6.5, 4.5) {ViT (patch $16{\times}16$)};
\node[font=\small\bfseries] at (6.5, 4.0) {$196 \times 196$};
\fill[orange!30] (6.0, 2.0) rectangle (7.0, 3.0);
\draw[orange!70!black, thick] (6.0, 2.0) rectangle (7.0, 3.0);
\node[font=\scriptsize, orange!70!black] at (6.5, 2.5) {38K};
\node[font=\scriptsize, orange!70!black] at (6.5, 1.5) {$\approx 0.15$ MB};
\node[font=\scriptsize, orange!70!black, fill=orange!10, draw=orange!50, rounded corners] at (6.5, -3.2) {Fattibile (perde dettaglio)};

% Perceiver cross-attention
\node[font=\small\bfseries] at (12, 4.5) {Perceiver};
\node[font=\small\bfseries] at (12, 4.0) {$512 \times 50{,}176$};
\fill[green!30] (10.2, 1.5) rectangle (13.8, 2.0);
\draw[green!60!black, thick] (10.2, 1.5) rectangle (13.8, 2.0);
\node[font=\scriptsize, green!60!black] at (12, 1.75) {25.7M};
\node[font=\scriptsize, green!60!black] at (12, 1.0) {$\approx 103$ MB};
\node[font=\scriptsize, green!60!black, fill=green!10, draw=green!50, rounded corners] at (12, -3.2) {Efficiente (pixel-level)};

% Scale reference
\node[font=\tiny, gray] at (6.5, -4.0) {Le dimensioni dei rettangoli sono indicative della dimensione della matrice di attenzione};

\end{tikzpicture}
\caption{Confronto delle dimensioni della matrice di attenzione: il Transformer standard ($50{,}176 \times 50{,}176$) richiederebbe $\sim$10 GB (impossibile); ViT ($196 \times 196$) \`e fattibile ma perde dettaglio sub-patch; il Perceiver ($512 \times 50{,}176$) \`e efficiente e lavora a livello di singolo pixel.}
\label{fig:attention_matrix_sizes_comparison}
\end{figure}


\newpage
%% ============================================================
%%  SEZIONE 4 --- POSITION ENCODING CON FOURIER FEATURES
%% ============================================================
\section{Position Encoding con Fourier Features} \label{sec:fourier}

Ora che abbiamo descritto l'architettura del Perceiver (cross-attention, Latent Transformer e weight sharing), dobbiamo affrontare un problema fondamentale che non abbiamo ancora risolto: come fa il Perceiver a ``sapere'' dove si trovano i pixel nell'immagine? Senza questa informazione, il modello non potrebbe distinguere un volto da una versione ``mescolata'' dello stesso volto!

\subsection{Il problema della permutation invariance} \label{sec:perm_invariance}

Facciamo un passo indietro e chiediamoci: il Perceiver ``sa'' dove si trovano i pixel nell'immagine? La risposta, sorprendentemente, \`e no --- almeno non senza un aiuto esplicito. Vediamo perch\'e.

L'operazione di attenzione (\hyperref[sec:attention]{Sezione~2.2}) \`e \textbf{invariante rispetto alle permutazioni}: se permutiamo le righe di $\mathbf{x}$, l'output della cross-attention rimane invariato. Questo significa che il Perceiver, di per s\'e, non ha alcuna nozione di posizione spaziale o temporale. Per molti compiti percettivi, la posizione \`e un'informazione cruciale.

\begin{proofbox}[Dimostrazione: la cross-attention \`e permutation-invariant]
Sia $\pi$ una permutazione degli indici $\{1, \ldots, M\}$ e sia $\mathbf{P}_\pi$ la matrice di permutazione corrispondente. Se permutiamo le righe dell'input: $\mathbf{x}' = \mathbf{P}_\pi \mathbf{x}$.

Le chiavi e i valori diventano:
\[
    \mathbf{K}' = \mathbf{x}' \mathbf{W}_K = \mathbf{P}_\pi \mathbf{x} \mathbf{W}_K = \mathbf{P}_\pi \mathbf{K}, \qquad \mathbf{V}' = \mathbf{P}_\pi \mathbf{V}
\]

Lo score: $\mathbf{Q} {\mathbf{K}'}^\top = \mathbf{Q} (\mathbf{P}_\pi \mathbf{K})^\top = \mathbf{Q} \mathbf{K}^\top \mathbf{P}_\pi^\top$.

La softmax \`e applicata per riga, e la permutazione delle colonne non cambia il risultato della combinazione lineare finale:
\[
    \text{softmax}\!\left(\frac{\mathbf{Q} \mathbf{K}^\top \mathbf{P}_\pi^\top}{\sqrt{d_k}}\right) \mathbf{P}_\pi \mathbf{V} = \text{softmax}\!\left(\frac{\mathbf{Q} \mathbf{K}^\top}{\sqrt{d_k}}\right) \mathbf{P}_\pi^\top \mathbf{P}_\pi \mathbf{V} = \text{softmax}\!\left(\frac{\mathbf{Q} \mathbf{K}^\top}{\sqrt{d_k}}\right) \mathbf{V}
\]
perch\'e $\mathbf{P}_\pi^\top \mathbf{P}_\pi = \mathbf{I}$ (le matrici di permutazione sono ortogonali) e la softmax per riga \`e invariante alla permutazione delle colonne dello score (la somma nella softmax \`e su tutti gli elementi, indipendentemente dall'ordine). $\blacksquare$
\end{proofbox}

\textbf{Conseguenza}: senza encoding posizionale, il Perceiver non potrebbe distinguere tra un'immagine di un volto e la stessa immagine con i pixel mescolati casualmente. Per compiti di classificazione, la posizione relativa dei pixel \`e fondamentale.

\begin{notabox}[In parole semplici: la permutation invariance]
Immaginiamo di avere le tessere di un mosaico. Senza posizione, il modello vede solo ``ci sono 100 tessere rosse, 200 tessere blu, 50 tessere verdi'' --- ma non sa se formano un tramonto o un paesaggio astratto. \`E la \emph{disposizione} delle tessere che conta, non solo il loro colore. L'encoding posizionale ``numera'' ogni tessera, permettendo al modello di ricostruire il mosaico.
\end{notabox}

\begin{examplebox}[Esempio: perch\'e la permutation invariance \`e un problema]
Consideriamo due immagini $3 \times 3$ (per semplicit\`a):

\textbf{Immagine A} (una ``L''): pixel bianchi in alto a destra e in basso; pixel neri nella colonna sinistra.

\textbf{Immagine B} (una ``T''): pixel bianchi in basso e nella colonna centrale; pixel neri nelle colonne laterali.

Senza informazione posizionale, le due immagini hanno \emph{gli stessi pixel} ma in posizioni diverse. Il Perceiver senza PE calcolerebbe la stessa distribuzione di attenzione per entrambe, producendo la stessa rappresentazione e quindi la stessa classificazione. Con le Fourier features, ogni pixel ``sa dove si trova'' e le due immagini producono rappresentazioni diverse.
\end{examplebox}

% TikZ: Why Fourier features - without vs with position
\begin{figure}[H]
\centering
\begin{tikzpicture}[
    pixbox/.style={draw, thick, minimum width=1.2cm, minimum height=1.2cm, rounded corners=2pt},
    arrowstyle/.style={-{Stealth[length=2.5mm]}, thick},
    >=Stealth,
    scale=0.85, transform shape
]

% LEFT: Without position
\node[font=\small\bfseries] at (2, 4.5) {Senza posizione};
\node[pixbox, fill=red!40] (p1) at (0.5, 3) {};
\node[font=\tiny, text width=1.2cm, align=center] at (0.5, 3) {R=200 G=50 B=50};
\node[font=\tiny, gray] at (0.5, 2.0) {pixel (10, 20)};

\node[pixbox, fill=red!40] (p2) at (3.5, 3) {};
\node[font=\tiny, text width=1.2cm, align=center] at (3.5, 3) {R=200 G=50 B=50};
\node[font=\tiny, gray] at (3.5, 2.0) {pixel (180, 100)};

\node[font=\scriptsize, text=red!70!black, text width=3.5cm, align=center] at (2, 1.0) {Il modello vede \textbf{identici}! Non sa chi \`e in alto e chi \`e in basso};
\draw[very thick, red!60!black] (0.5, 1.6) -- (3.5, 1.6);
\node[font=\large, red!60!black] at (2, 1.6) {$=$};

% Separator
\draw[dashed, thick, gray] (5.5, 5) -- (5.5, 0.5);

% RIGHT: With Fourier features
\node[font=\small\bfseries] at (9, 4.5) {Con Fourier features};
\node[pixbox, fill=red!40] (p3) at (7, 3) {};
\node[font=\tiny, text width=1.2cm, align=center] at (7, 3) {R=200 G=50 B=50};
\node[font=\tiny, blue!70!black, text width=2.5cm, align=center] at (7, 1.8) {+ pos: $[-0.91, \sin(0.3\pi), \ldots]$};
\node[font=\tiny, gray] at (7, 1.2) {258 valori};

\node[pixbox, fill=red!40] (p4) at (11, 3) {};
\node[font=\tiny, text width=1.2cm, align=center] at (11, 3) {R=200 G=50 B=50};
\node[font=\tiny, blue!70!black, text width=2.5cm, align=center] at (11, 1.8) {+ pos: $[0.61, \sin(2.1\pi), \ldots]$};
\node[font=\tiny, gray] at (11, 1.2) {258 valori};

\node[font=\scriptsize, text=green!50!black, text width=4cm, align=center] at (9, 0.6) {Il modello li vede \textbf{diversi}!\\Ogni pixel ha un ``codice postale''\\unico basato sulla sua posizione};
\draw[very thick, green!60!black] (7, 2.5) -- (11, 2.5);
\node[font=\large, green!60!black] at (9, 2.5) {$\neq$};

\end{tikzpicture}
\caption{Perch\'e servono le Fourier features. A sinistra: senza informazione posizionale, due pixel rossi in posizioni diverse sono indistinguibili per il modello. A destra: con le Fourier features concatenate, ogni pixel porta un ``codice postale'' unico, permettendo al modello di distinguerli.}
\label{fig:why_fourier_features}
\end{figure}

\subsection{Fourier feature positional encoding} \label{sec:fourier_encoding}

\begin{formulabox}[Fourier feature encoding]
Per un punto con coordinate $(x_1, x_2, \ldots, x_P)$ in $P$ dimensioni spaziali, con $K$ bande di frequenza:
\[
    \text{PE}(x_1, \ldots, x_P) = \Big[\underbrace{x_1, \ldots, x_P}_{\text{raw}},\;
    \underbrace{\sin(f_1 \pi x_1), \cos(f_1 \pi x_1), \ldots}_{\text{FF per } x_1},\; \ldots\Big]
\]
La dimensione totale dell'encoding \`e:
\[
    d_{\text{PE}} = P \cdot (2K + 1)
\]
dove il $+1$ \`e dovuto alla coordinata raw $x_d$ inclusa per ogni dimensione.
\end{formulabox}

L'encoding posizionale viene \textbf{concatenato} ai canali dell'input (non sommato, come nel Transformer originale):
\[
    \tilde{\mathbf{x}} = [\mathbf{x} \;;\; \text{PE}] \in \mathbb{R}^{M \times (C + d_{\text{PE}})}
\]

\begin{examplebox}[Esempio: Fourier features per ImageNet]
Per ImageNet ($224 \times 224$ pixel, $P = 2$ dimensioni, $K = 64$ bande):
\begin{align*}
    d_{\text{PE}} &= 2 \times (2 \times 64 + 1) = 2 \times 129 = 258 \\
    C_{\text{RGB}} &= 3 \\
    C_{\text{totale}} &= 3 + 258 = 261
\end{align*}
Quindi ogni pixel \`e rappresentato da un vettore di 261 dimensioni: 3 per i canali RGB e 258 per l'encoding posizionale. L'input al Perceiver \`e $\tilde{\mathbf{x}} \in \mathbb{R}^{50{,}176 \times 261}$.
\end{examplebox}

\subsection{Esempio numerico completo: Fourier features per il pixel (112, 56)} \label{sec:fourier_numerico}

Per consolidare la comprensione, seguiamo passo per passo il calcolo delle Fourier features per un pixel specifico dell'immagine ImageNet. Questo esempio \`e importante per l'esame: mostra che si padroneggiano i dettagli dell'implementazione, non solo i concetti generali.

\begin{examplebox}[Calcolo dettagliato delle Fourier features per il pixel (112, 56)]
\textbf{Passo 1: Coordinate normalizzate.} Mappiamo le coordinate $(i, j)$ nell'intervallo $[-1, 1]$:
\begin{align*}
    x_1 &= \frac{2 \times 112}{223} - 1 = \frac{224}{223} - 1 = 1.00448 - 1 = 0.00448 \\
    x_2 &= \frac{2 \times 56}{223} - 1 = \frac{112}{223} - 1 = 0.50224 - 1 = -0.4978
\end{align*}

\textbf{Passo 2: Bande di frequenza.} Le $K = 64$ frequenze sono distribuite logaritmicamente tra $f_1 = 1.0$ e $f_{64} = \mu / 2 = 112.0$:
\[
    f_k = \exp\!\left(\frac{k-1}{63} \cdot \ln(112)\right) \quad \text{per } k = 1, \ldots, 64
\]
Valori selezionati: $f_1 = 1.0$, $f_{16} \approx 4.47$, $f_{32} \approx 10.58$, $f_{48} \approx 25.03$, $f_{64} = 112.0$.

\textbf{Passo 3: Calcolo sin/cos per la coordinata $x_1 = 0.00448$.}

Per $f_1 = 1.0$:
\begin{align*}
    \sin(f_1 \pi x_1) &= \sin(1.0 \times \pi \times 0.00448) = \sin(0.01408) \approx 0.01408 \\
    \cos(f_1 \pi x_1) &= \cos(0.01408) \approx 0.99990
\end{align*}

Per $f_{32} = 10.58$:
\begin{align*}
    \sin(10.58 \times \pi \times 0.00448) &= \sin(0.14886) \approx 0.14831 \\
    \cos(10.58 \times \pi \times 0.00448) &= \cos(0.14886) \approx 0.98894
\end{align*}

Per $f_{64} = 112.0$:
\begin{align*}
    \sin(112.0 \times \pi \times 0.00448) &= \sin(1.5762) \approx 0.99997 \\
    \cos(112.0 \times \pi \times 0.00448) &= \cos(1.5762) \approx 0.00765
\end{align*}

\textbf{Passo 4: Calcolo sin/cos per la coordinata $x_2 = -0.4978$.} Procedimento analogo per tutte le 64 frequenze.

\textbf{Passo 5: Assemblaggio del vettore finale.} Per ogni coordinata, abbiamo: 1 valore raw + $64 \times 2 = 128$ valori sin/cos $= 129$ valori. Con 2 coordinate: $2 \times 129 = 258$ valori posizionali.

Il vettore finale del pixel (112, 56) con valori RGB $(0.45, 0.32, 0.78)$:
\[
    \tilde{\mathbf{x}}_{(112,56)} = [\underbrace{0.45, 0.32, 0.78}_{3\text{ RGB}},\;
    \underbrace{0.004, 0.014, 0.999, \ldots}_{129\text{ per }x_1},\;
    \underbrace{-0.498, -0.999, \ldots}_{129\text{ per }x_2}] \in \mathbb{R}^{261}
\]
\end{examplebox}

% TikZ: Pixel enrichment with Fourier features
\begin{figure}[H]
\centering
\begin{tikzpicture}[
    box/.style={draw, rounded corners=3pt, minimum height=0.7cm, font=\small},
    arrowstyle/.style={-{Stealth[length=2.5mm]}, thick},
    >=Stealth,
    scale=0.9, transform shape
]

% Original pixel
\node[box, fill=red!20, minimum width=1.5cm] (pixel) at (0, 0) {Pixel};
\node[font=\scriptsize, below=0.1cm of pixel] {$[\text{R}, \text{G}, \text{B}]$};
\node[font=\scriptsize, below=0.5cm of pixel] {$\mathbb{R}^{3}$};

% Arrow
\draw[arrowstyle, very thick] (1.5, 0) -- (3.5, 0);
\node[font=\scriptsize, above=0.1cm] at (2.5, 0) {+ Fourier PE};

% Enriched pixel
\node[box, fill=green!20, minimum width=5cm] (enriched) at (7, 0) {Pixel arricchito};
\node[font=\scriptsize, below=0.1cm of enriched] {$[\text{R}, \text{G}, \text{B}, x_1, \sin(f_1\pi x_1), \cos(f_1\pi x_1), \ldots, x_2, \sin(f_1\pi x_2), \ldots]$};
\node[font=\scriptsize, below=0.5cm of enriched] {$\mathbb{R}^{261}$};

% Breakdown
\draw[decorate, decoration={brace, amplitude=5pt, mirror}, thick]
    (4.2, -1.5) -- (5.0, -1.5) node[midway, below=6pt, font=\scriptsize] {3 RGB};
\draw[decorate, decoration={brace, amplitude=5pt, mirror}, thick]
    (5.2, -1.5) -- (7.4, -1.5) node[midway, below=6pt, font=\scriptsize] {129 per $x_1$};
\draw[decorate, decoration={brace, amplitude=5pt, mirror}, thick]
    (7.6, -1.5) -- (9.8, -1.5) node[midway, below=6pt, font=\scriptsize] {129 per $x_2$};

% Multiplication factor
\node[draw, circle, fill=yellow!20, minimum size=0.8cm, font=\small\bfseries] at (2.5, -2.5) {$\times 87$};
\node[font=\scriptsize] at (2.5, -3.2) {Da 3 a 261 dimensioni};

\end{tikzpicture}
\caption{Arricchimento di un singolo pixel con Fourier features. Il pixel passa da $\mathbb{R}^3$ (solo colore) a $\mathbb{R}^{261}$ (colore + posizione codificata con funzioni sinusoidali a 64 frequenze per ciascuna delle 2 coordinate spaziali). La dimensionalit\`a aumenta di un fattore $87\times$.}
\label{fig:pixel_enrichment}
\end{figure}

% TikZ: Fourier feature generation flow
\begin{figure}[H]
\centering
\begin{tikzpicture}[
    box/.style={draw, rounded corners=3pt, minimum height=0.7cm, font=\small},
    arrowstyle/.style={-{Stealth[length=2.5mm]}, thick},
    >=Stealth,
    scale=0.85, transform shape
]

% Coordinate
\node[box, fill=blue!15, minimum width=2cm] (coord) at (0, 0) {$x_d \in [-1,1]$};

% Frequency bands
\node[box, fill=orange!15, minimum width=2.5cm] (freq) at (4, 1.5) {$f_1 = 1.0$};
\node[box, fill=orange!15, minimum width=2.5cm] (freq2) at (4, 0.5) {$f_2 \approx 1.08$};
\node[font=\small] at (4, -0.3) {$\vdots$};
\node[box, fill=orange!15, minimum width=2.5cm] (freqK) at (4, -1.2) {$f_{64} = 112.0$};

% Multiply
\node[font=\small] at (2, 0.8) {$\times \pi$};
\draw[arrowstyle, gray] (coord.east) -- (freq.west);
\draw[arrowstyle, gray] (coord.east) -- (freq2.west);
\draw[arrowstyle, gray] (coord.east) -- (freqK.west);

% Sin/Cos
\node[box, fill=green!15, minimum width=1.5cm] (sin1) at (8, 2) {$\sin(\cdot)$};
\node[box, fill=green!15, minimum width=1.5cm] (cos1) at (8, 1) {$\cos(\cdot)$};
\node[font=\small] at (8, -0.3) {$\vdots$};
\node[box, fill=green!15, minimum width=1.5cm] (sinK) at (8, -0.8) {$\sin(\cdot)$};
\node[box, fill=green!15, minimum width=1.5cm] (cosK) at (8, -1.7) {$\cos(\cdot)$};

\draw[arrowstyle] (freq.east) -- (sin1.west);
\draw[arrowstyle] (freq.east) -- (cos1.west);
\draw[arrowstyle] (freqK.east) -- (sinK.west);
\draw[arrowstyle] (freqK.east) -- (cosK.west);

% Concatenation
\node[box, fill=yellow!20, minimum width=2.5cm] (concat) at (12, 0) {\begin{tabular}{c}Concatena\\$[x_d, \sin, \cos, \ldots]$\end{tabular}};
\draw[arrowstyle] (sin1.east) -- (concat.north west);
\draw[arrowstyle] (cos1.east) -- (concat.west);
\draw[arrowstyle] (sinK.east) -- (concat.west);
\draw[arrowstyle] (cosK.east) -- (concat.south west);
\draw[arrowstyle] (coord.east) .. controls (1.5,-2.5) and (10,-2.5) .. (concat.south);

% Output
\node[box, fill=purple!15, minimum width=2cm] (out) at (12, -2.5) {$\mathbb{R}^{129}$};
\draw[arrowstyle] (concat) -- (out);
\node[font=\scriptsize] at (12, -3.3) {1 raw + 128 sin/cos};

\end{tikzpicture}
\caption{Schema di generazione delle Fourier features per una singola coordinata $x_d$. La coordinata viene moltiplicata per $K=64$ frequenze, ciascuna produce un valore $\sin$ e un valore $\cos$. Con la coordinata raw, si ottengono $1 + 2 \times 64 = 129$ valori. Per 2 coordinate ($P=2$): $2 \times 129 = 258$ valori.}
\label{fig:fourier_generation}
\end{figure}

\begin{notabox}[title={Perch\'e le Fourier features e non semplici coordinate $(x, y)$?}]
Due pixel adiacenti hanno coordinate quasi identiche: $(0.004, -0.498)$ e $(0.004, -0.489)$. La differenza \`e solo $0.009$, troppo piccola perch\'e la rete la ``noti'' facilmente. Con le Fourier features ad alta frequenza ($f_{64} = 112$), le stesse coordinate producono valori di $\sin$ e $\cos$ molto diversi, rendendo i pixel facilmente distinguibili. Questa \`e l'intuizione del \textbf{criterio di Nyquist}: per distinguere $R$ posizioni equispaziate su una griglia, servono frequenze fino a $R/2$.
\end{notabox}

\begin{examplebox}[Perch\'e le bande di frequenza multiple sono necessarie]
\textbf{Con una sola frequenza bassa} ($f = 1$): $\sin(\pi \times 0.004) = 0.0126$ e $\sin(\pi \times 0.013) = 0.0408$. La differenza \`e 0.028 --- poco distinguibile.

\textbf{Con una frequenza alta} ($f = 112$): $\sin(112\pi \times 0.004) = 0.9999$ e $\sin(112\pi \times 0.013) = -0.077$. La differenza \`e 1.077 --- \textbf{enormemente distinguibile}!

\textbf{Ma solo frequenze alte non bastano}: due pixel lontani (es.\ $(0, 0)$ e $(0.5, 0)$) possono avere lo stesso valore ad alta frequenza (perch\'e $\sin$ \`e periodico). Le frequenze basse distinguono posizioni lontane; le frequenze alte distinguono posizioni vicine. L'insieme completo di 64 bande crea un ``codice'' univoco per ogni posizione.

Questa idea \`e analoga alla \textbf{trasformata di Fourier}: una funzione arbitraria pu\`o essere rappresentata come somma di sinusoidi a frequenze diverse. Le Fourier features ``campionano'' questa rappresentazione a 64 frequenze discrete, sufficiente per codificare univocamente ogni pixel nella griglia $224 \times 224$.
\end{examplebox}

\begin{notabox}[In parole semplici: perch\'e servono tante frequenze]
Detto in parole povere, le diverse frequenze sono come una ``scala di zoom''. La frequenza bassa ($f_1 = 1$) \`e come una mappa del mondo: distingue continenti ma non citt\`a. La frequenza media ($f_{32} \approx 56$) \`e come una mappa regionale: distingue citt\`a ma non quartieri. La frequenza alta ($f_{64} = 112$) \`e come una mappa stradale: distingue quartieri e singoli edifici. Combinando tutte le 64 frequenze, ogni posizione ha un ``indirizzo'' unico --- come avere contemporaneamente la mappa del mondo, regionale e stradale.
\end{notabox}

\subsubsection{Confronto: Fourier features vs learned positional embedding} \label{sec:ff_vs_learned}

\begin{table}[htbp]
\centering
\caption{Confronto tra Fourier features e positional embedding appreso.}
\begin{tabular}{@{}p{3.5cm}p{5cm}p{5cm}@{}}
\toprule
\textbf{Propriet\`a} & \textbf{Fourier features} & \textbf{Learned PE} \\
\midrule
Parametri aggiuntivi & 0 (calcolate analiticamente) & $M \times d_{\text{PE}}$ (apprendibili) \\
ImageNet Top-1 & 78.0\% & 70.9\% ($-7.1$) \\
Generalizzazione a $M$ diversi & S\`i (funzionano per qualsiasi $M$) & No (fissati per $M$ specifico) \\
Invarianza alla permutazione & S\`i (concatenate) & S\`i (concatenate) \\
Interpretabilit\`a & Alta (frequenze note) & Bassa (vettori opachi) \\
Multi-dimensionale & Naturale ($P$ coordinate) & Richiede design specifico \\
\bottomrule
\end{tabular}
\end{table}

Il gap di 7.1 punti percentuali \`e enorme --- per avere un termine di paragone, corrisponde alla differenza tra un modello ``buono'' e un modello ``mediocre'' su ImageNet. Questo dimostra che le Fourier features sono molto superiori alle posizioni apprese per il Perceiver. Il motivo principale \`e che con $M = 50{,}176$ posizioni, apprendere un embedding univoco per ciascuna richiede $50{,}176 \times 258 \approx 12.9$M parametri aggiuntivi che devono essere ottimizzati --- un compito arduo quando ogni posizione appare in media solo $\sim$24 volte per epoca ($1.2$M immagini / $50{,}176$ posizioni).


\newpage
%% ============================================================
%%  SEZIONE 5 --- FORWARD PASS COMPLETO
%% ============================================================
\section{Forward Pass Completo: Passo per Passo} \label{sec:forward}

Questa \`e probabilmente la sezione pi\`u importante della dispensa. Qui seguiremo, passo dopo passo, il viaggio di un'immagine attraverso l'intero Perceiver --- dall'ingresso come griglia di pixel all'uscita come distribuzione di probabilit\`a su 1000 classi. Per ogni passo, mostriamo la \textbf{pipeline completa} con il passo corrente evidenziato, le \textbf{dimensioni esatte} dei tensori, e una \textbf{spiegazione discorsiva} del perch\'e quel passo serve.

Il forward pass trasforma un'immagine RGB di $224 \times 224$ pixel in una distribuzione di probabilit\`a su 1000 classi. Il percorso coinvolge:
\begin{enumerate}[label=\arabic*.]
    \item \textbf{Preparazione dell'input} (A1--A3): l'immagine viene appiattita, arricchita con Fourier features e proiettata in uno spazio di embedding.
    \item \textbf{Inizializzazione dei latenti} (A4): un array appreso di 512 vettori viene caricato come ``memoria di lavoro''.
    \item \textbf{Cross-attention + MLP} (A5--A6): i latenti ``leggono'' l'immagine e elaborano l'informazione estratta.
    \item \textbf{Latent Transformer} (A7): i latenti ``parlano tra loro'' per integrare l'informazione.
    \item \textbf{Iterazione} (A8): i passi 3--4 vengono ripetuti $T = 8$ volte per raffinare la rappresentazione.
    \item \textbf{Output} (A9--A12): la rappresentazione viene mediata, classificata e confrontata con la label vera.
\end{enumerate}

\textbf{Pipeline completa}:
\begin{center}
\small
\fbox{\parbox{0.95\textwidth}{\centering
$\text{Input} \to \text{Fourier} \to \text{Embed} \to \text{Latent Array} \to \text{Cross-Att} \to \text{Res{+}MLP}$\\[2pt]
$\to \text{Latent Transf.} \to \text{Res{+}MLP} \to (\times T) \to \text{Pool} \to \text{Classifier} \to \text{Softmax} \to \text{Loss}$
}}
\end{center}

% TikZ: Complete vertical forward pass pipeline with dimensions
\begin{figure}[H]
\centering
\begin{tikzpicture}[
    step/.style={draw, rounded corners=3pt, fill=#1, minimum width=5.5cm, minimum height=0.6cm, font=\small},
    step/.default={gray!10},
    arrowstyle/.style={-{Stealth[length=2.5mm]}, thick},
    dim/.style={font=\scriptsize, gray},
    >=Stealth,
    scale=0.88, transform shape
]

\node[step=red!15] (s1) at (0, 12) {A1. Immagine RGB};
\node[dim, right=0.3cm of s1] {$224 \times 224 \times 3$};

\node[step=red!10] (s2) at (0, 10.8) {A1. Flatten (byte array)};
\node[dim, right=0.3cm of s2] {$50{,}176 \times 3$};

\node[step=yellow!20] (s3) at (0, 9.6) {A2. Fourier Features};
\node[dim, right=0.3cm of s3] {$50{,}176 \times 261$};

\node[step=orange!15] (s4) at (0, 8.4) {A3. Proiezione Lineare};
\node[dim, right=0.3cm of s4] {$50{,}176 \times 512$};

\node[step=blue!15] (s5) at (0, 7.2) {A4. Latent Array (appreso)};
\node[dim, right=0.3cm of s5] {$512 \times 1024$};

\node[step=orange!20] (s6) at (0, 6.0) {A5. Cross-Attention};
\node[dim, right=0.3cm of s6] {$512 \times 1024$};

\node[step=green!15] (s7) at (0, 4.8) {A6. Residual + MLP (FFN)};
\node[dim, right=0.3cm of s7] {$512 \times 1024$};

\node[step=cyan!15] (s8) at (0, 3.6) {A7. Latent Transformer ($L{=}6$ blocchi SA)};
\node[dim, right=0.3cm of s8] {$512 \times 1024$};

\node[step=purple!15] (s9) at (0, 2.4) {A8. Ripetizione ($\times T{=}8$)};
\node[dim, right=0.3cm of s9] {$512 \times 1024$};

\node[step=green!20] (s10) at (0, 1.2) {A9. Global Average Pooling};
\node[dim, right=0.3cm of s10] {$1024$};

\node[step=blue!20] (s11) at (0, 0) {A10. Classificatore Lineare};
\node[dim, right=0.3cm of s11] {$1000$};

\node[step=red!20] (s12) at (0, -1.2) {A11--A12. Softmax + Cross-Entropy Loss};
\node[dim, right=0.3cm of s12] {scalare $\mathcal{L}$};

% Arrows
\draw[arrowstyle] (s1) -- (s2);
\draw[arrowstyle] (s2) -- (s3);
\draw[arrowstyle] (s3) -- (s4);
\draw[arrowstyle] (s4) -- (s6);
\draw[arrowstyle] (s5) -- (s6);
\draw[arrowstyle] (s6) -- (s7);
\draw[arrowstyle] (s7) -- (s8);
\draw[arrowstyle] (s8) -- (s9);
\draw[arrowstyle] (s9) -- (s10);
\draw[arrowstyle] (s10) -- (s11);
\draw[arrowstyle] (s11) -- (s12);

% Iteration loop
\draw[arrowstyle, dashed, blue!60!black] (3.8, 2.4) -- (3.8, 6.0) -- (s6.east);
\node[font=\scriptsize, blue!60!black, rotate=90] at (4.1, 4.2) {iterazione $t=2,\ldots,T$};

% Input feeding annotation
\draw[arrowstyle, dashed, gray] (-3.8, 8.4) -- (-3.8, 6.0) -- (s6.west);
\node[font=\scriptsize, gray, rotate=90] at (-4.1, 7.2) {input $\mathbf{x}$ ri-letto};

\end{tikzpicture}
\caption{Pipeline verticale completa del forward pass del Perceiver per ImageNet. Le dimensioni di ogni tensore sono annotate a destra. Il ciclo tratteggiato blu indica la ripetizione iterativa: ad ogni iterazione, i latenti vengono ri-aggiornati tramite cross-attention sull'input.}
\label{fig:forward_pipeline_vertical}
\end{figure}


%% --- A1: Byte Array ---
\subsection{A1. Dall'immagine al Byte Array} \label{sec:byte_array}

\begin{pipelinebox}
\textbf{Input} $\to$ Fourier $\to$ Embedding $\to$ Latent Array $\to$ Cross-Att $\to$ Res+MLP $\to$ Latent Transf.\ $\to$ Res+MLP $\to$ ($\times T$) $\to$ Pool $\to$ Classifier $\to$ Softmax $\to$ Loss
\end{pipelinebox}

Come descritto nella \hyperref[sec:panoramica]{Sezione~\ref*{sec:panoramica}} (e illustrato nella Figura~\ref{fig:byte_array_diagram}), l'immagine $224 \times 224 \times 3$ viene appiattita in un byte array $\mathbf{x}_{\text{raw}} \in \mathbb{R}^{50{,}176 \times 3}$ tramite flatten. Ogni riga corrisponde a un pixel, ogni colonna a un canale RGB. Il pixel alla posizione $(i, j)$ diventa la riga $i \times 224 + j$.

% Figure di flatten gi\`a presente nella Sezione 3 (Figura~\ref{fig:byte_array_diagram})

\begin{table}[htbp]
\centering
\caption{Il byte array per diverse modalit\`a: stessa astrazione, dati diversi.}
\begin{tabular}{@{}llccl@{}}
\toprule
\textbf{Modalit\`a} & \textbf{Sorgente} & $M$ & $C$ & \textbf{Significato di una riga} \\
\midrule
Immagine & ImageNet $224^2$ & 50{,}176 & 3 & Un pixel: $[R, G, B]$ \\
Audio (raw) & WAV 48kHz & 480 & 128 & 128 campioni audio contigui \\
Audio (mel) & Mel spectrogram & 4{,}800 & var & Un frame spettrale \\
Point cloud & ModelNet40 & 2{,}048 & 3 & Un punto: $[x, y, z]$ \\
Video & $224^2 \times 30$fps & $\sim$1.5M & 3 & Un pixel spazio-temporale \\
Testo & Sequenza di token & $\sim$512 & $d_{\text{token}}$ & Un token embedding \\
\bottomrule
\end{tabular}
\end{table}

\textbf{L'universalit\`a del byte array}: questo \`e forse il concetto pi\`u potente dell'intero paper. Qualsiasi dato che pu\`o essere rappresentato come una matrice $M \times C$ pu\`o essere processato dal Perceiver. Questo include non solo le modalit\`a standard (immagini, audio, testo), ma anche dati non convenzionali come grafi (ogni nodo come riga), misure di sensori, dati tabulari, ecc. La sola cosa che cambia da una modalit\`a all'altra \`e la scelta delle Fourier features (dimensionalit\`a e frequenze).

\begin{examplebox}[Come convertire diverse modalit\`a in byte array]

\textbf{Immagine a colori} ($H \times W \times 3$): flatten $\to$ $M = H \times W$ righe, $C = 3$ colonne (RGB). Ogni riga = un pixel.

\textbf{Audio mono a 48kHz}: segmentiamo in finestre di $w$ campioni $\to$ $M = \lceil S / w \rceil$ righe (dove $S$ = campioni totali), $C = w$ colonne. Ogni riga = una finestra temporale.

\textbf{Nuvola di punti 3D}: $M$ = numero di punti, $C = 3$ (coordinate $x, y, z$). Opzionalmente, $C = 6$ se includiamo anche le normali $n_x, n_y, n_z$.

\textbf{Testo}: dopo tokenizzazione, $M$ = numero di token, $C = d_{\text{embed}}$ (dimensione dell'embedding del token). Le Fourier features codificano la posizione sequenziale.

\textbf{Video}: appiattire tutti i frame in una sequenza di pixel spazio-temporali: $M = T_{\text{frame}} \times H \times W$, $C = 3$. Le Fourier features codificano sia la posizione spaziale $(x, y)$ sia la posizione temporale $t$ ($P = 3$ dimensioni).

\textbf{Dati tabulari}: $M$ = numero di feature (colonne della tabella), $C$ = dimensione dell'embedding di ogni feature. Le Fourier features non sono necessarie (nessuna struttura spaziale).
\end{examplebox}


\begin{notabox}[In parole semplici: il byte array]
Detto in parole povere, il byte array \`e come prendere una foto e trasformarla in un foglio Excel. Ogni riga del foglio \`e un pixel, e le colonne sono i valori di rosso, verde e blu. Il Perceiver non ``vede'' una foto --- vede un foglio Excel con 50.176 righe e 3 colonne. Questo \`e sia il suo limite (non ``sa'' che i pixel erano in una griglia) sia la sua forza (pu\`o leggere qualsiasi foglio Excel, indipendentemente da cosa contenga).
\end{notabox}


%% --- A2: Fourier Features ---
\subsection{A2. Fourier Features --- Arricchimento posizionale} \label{sec:fourier_forward}

\begin{pipelinebox}
Input $\to$ \textbf{Fourier} $\to$ Embedding $\to$ Latent Array $\to$ Cross-Att $\to$ Res+MLP $\to$ Latent Transf.\ $\to$ Res+MLP $\to$ ($\times T$) $\to$ Pool $\to$ Classifier $\to$ Softmax $\to$ Loss
\end{pipelinebox}

Ogni pixel del byte array contiene solo informazione di colore (RGB), ma non di posizione. Per risolvere questo problema, il Perceiver usa le \textbf{Fourier features} (teoria completa, motivazione e confronto con learned embeddings nella \hyperref[sec:fourier]{Sezione~\ref*{sec:fourier}}). Qui applichiamo le formule a un esempio numerico concreto.

\begin{examplebox}[Esempio numerico dettagliato: Fourier features per un pixel]
Consideriamo il pixel alla posizione $(112, 56)$:

\textbf{Passo 1: Coordinate normalizzate} (ogni coordinata viene mappata in $[-1, 1]$):
\[
    x = \frac{2 \times 112}{223} - 1 \approx 0.00448, \qquad y = \frac{2 \times 56}{223} - 1 \approx -0.4978
\]

\textbf{Passo 2: Calcolo delle Fourier features}. Con $K = 64$ bande di frequenza $f_1, f_2, \ldots, f_{64}$ distribuite linearmente tra $1$ e $\mu/2 = 112$, calcoliamo $\sin(f_k \pi x)$ e $\cos(f_k \pi x)$ per ogni $k$ e ogni coordinata.

Le frequenze crescono esponenzialmente: $f_1 = 1$ (bassa frequenza, distingue ``sinistra vs destra''), $f_{32} \approx 56$ (media frequenza), $f_{64} = 112$ (alta frequenza, distingue pixel adiacenti --- criterio di Nyquist).

\textbf{Passo 3}: L'encoding posizionale totale ha $2 \times 128 = 256$ valori sinusoidali + 2 coordinate raw = \textbf{258 valori}.

\textbf{Passo 4: Concatenazione con RGB}. Il pixel con valori RGB $(0.45, 0.32, 0.78)$ diventa:
\[
    \tilde{\mathbf{x}}_i = [\underbrace{0.45,\; 0.32,\; 0.78}_{\text{RGB (3 valori)}},\; \underbrace{\text{258 Fourier features}}_{\text{posizione}}] \in \mathbb{R}^{261}
\]
\end{examplebox}

\begin{formulabox}[Dimensioni dopo le Fourier features]
\[
    \mathbf{x}_{\text{raw}} \in \mathbb{R}^{50{,}176 \times 3} \quad \xrightarrow{\text{+Fourier}} \quad \tilde{\mathbf{x}} \in \mathbb{R}^{50{,}176 \times 261}
\]
Il canale $C$ passa da \textbf{3} (solo RGB) a \textbf{261} (RGB + posizioni Fourier). Ora ogni pixel ``sa dove si trova'' nell'immagine.
\end{formulabox}

\begin{examplebox}[Viaggio completo di UN pixel attraverso l'intero Perceiver]
Seguiamo il pixel alla posizione $(100, 50)$ con colore RGB $= (0.8, 0.2, 0.1)$ (rosso intenso) attraverso l'intera pipeline.

\textbf{Passo 1 --- Flatten}: il pixel (100, 50) in un'immagine $224 \times 224$ diventa la riga numero $100 \times 224 + 50 = 22{,}450$ del byte array. Il suo vettore: $\mathbf{x}_{22450} = [0.8, 0.2, 0.1] \in \mathbb{R}^3$.

\textbf{Passo 2 --- Coordinate normalizzate}:
\[
x_{\text{norm}} = \frac{2 \times 100}{223} - 1 \approx -0.103, \qquad y_{\text{norm}} = \frac{2 \times 50}{223} - 1 \approx -0.552
\]

\textbf{Passo 3 --- Fourier features (prime 3 bande di frequenza su $K{=}64$)}:
\begin{align*}
f_1 {=} 1: &\quad \sin(\pi \!\cdot\! 1 \!\cdot\! (-0.103)) = -0.321,\;\; \cos(\ldots) = 0.947 \\
f_2 {\approx} 2.75: &\quad \sin(\pi \!\cdot\! 2.75 \!\cdot\! (-0.103)) = -0.803,\;\; \cos(\ldots) = 0.596 \\
f_3 {\approx} 4.49: &\quad \sin(\pi \!\cdot\! 4.49 \!\cdot\! (-0.103)) = -0.990,\;\; \cos(\ldots) = -0.139
\end{align*}
Questo si ripete per tutte le 64 bande e per entrambe le coordinate $(x, y)$, producendo $2 \times 2 \times 64 + 2 = 258$ valori di posizione.

\textbf{Passo 4 --- Concatenazione}:
\[
\tilde{\mathbf{x}}_{22450} = [\underbrace{0.8, 0.2, 0.1}_{\text{RGB}},\; \underbrace{-0.103, -0.552}_{\text{coord.}},\; \underbrace{-0.321, 0.947, \ldots}_{\text{256 FF}}] \in \mathbb{R}^{261}
\]

\textbf{Passo 5 --- Proiezione lineare}: $\mathbf{x}_{22450}^{\text{embed}} = \tilde{\mathbf{x}}_{22450} \cdot \mathbf{W}_{\text{embed}} + \mathbf{b}_{\text{embed}} \in \mathbb{R}^{512}$

Il nostro pixel \`e ora un vettore di 512 dimensioni. Ha ``perso'' la sua identit\`a come pixel rosso ed \`e diventato una rappresentazione astratta che codifica sia il colore sia la posizione.

\textbf{Passo 6 --- Cross-attention}: questo pixel diventa una riga di $\mathbf{K}$ e una riga di $\mathbf{V}$:
\begin{align*}
\mathbf{k}_{22450} &= \mathbf{x}_{22450}^{\text{embed}} \cdot \mathbf{W}_K \in \mathbb{R}^{1024} \quad \text{(la ``descrizione'' del pixel)} \\
\mathbf{v}_{22450} &= \mathbf{x}_{22450}^{\text{embed}} \cdot \mathbf{W}_V \in \mathbb{R}^{1024} \quad \text{(il ``contenuto'' del pixel)}
\end{align*}

Ogni latente $i$ calcola un peso di attenzione per questo pixel:
\[
\alpha_{i, 22450} = \frac{\exp(\mathbf{q}_i \cdot \mathbf{k}_{22450} / \sqrt{1024})}{\sum_{j=1}^{50176} \exp(\mathbf{q}_i \cdot \mathbf{k}_j / \sqrt{1024})}
\]

Se il latente $i$ ``cerca rosso'' e il nostro pixel \`e rosso, $\alpha_{i, 22450}$ sar\`a alto. Il contenuto del pixel ($\mathbf{v}_{22450}$) verr\`a ``assorbito'' dal latente con peso proporzionale ad $\alpha$.

\textbf{Risultato}: il pixel non esiste pi\`u come entit\`a separata. La sua informazione \`e stata ``distribuita'' tra i 512 latenti, ciascuno con il proprio peso di attenzione. Il pixel rosso ha contribuito maggiormente ai latenti che ``cercavano rosso'' o ``cercavano la regione in alto a sinistra''.
\end{examplebox}

% TikZ: One pixel journey through Perceiver
\begin{figure}[H]
\centering
\begin{tikzpicture}[
    step/.style={draw, rounded corners=3pt, fill=#1, minimum width=5cm, minimum height=0.55cm, font=\small},
    step/.default={gray!10},
    arrowstyle/.style={-{Stealth[length=2mm]}, thick},
    dim/.style={font=\scriptsize, gray},
    >=Stealth,
    scale=0.85, transform shape
]

\node[step=red!20] (s1) at (0, 10) {Pixel (100, 50): RGB = (0.8, 0.2, 0.1)};
\node[dim, right=0.2cm of s1] {$\mathbb{R}^{3}$};

\node[step=red!10] (s2) at (0, 8.5) {Flatten: riga 22{,}450 del byte array};
\node[dim, right=0.2cm of s2] {$\mathbb{R}^{3}$};

\node[step=yellow!20] (s3) at (0, 7) {Fourier: coord. + sin/cos};
\node[dim, right=0.2cm of s3] {$\mathbb{R}^{261}$};

\node[step=orange!15] (s4) at (0, 5.5) {Proiezione $\mathbf{W}_{\text{embed}}$};
\node[dim, right=0.2cm of s4] {$\mathbb{R}^{512}$};

\node[step=green!15] (s5) at (0, 4) {Diventa chiave $\mathbf{k}_{22450}$ e valore $\mathbf{v}_{22450}$};
\node[dim, right=0.2cm of s5] {$\mathbb{R}^{1024}$ ciascuno};

\node[step=cyan!15] (s6) at (0, 2.5) {Cross-attention: peso $\alpha_{i, 22450}$ per latente $i$};
\node[dim, right=0.2cm of s6] {scalare per latente};

\node[step=blue!15] (s7) at (0, 1) {Assorbito nei 512 latenti ($\alpha \cdot \mathbf{v}$)};
\node[dim, right=0.2cm of s7] {contributo a $\mathbb{R}^{512 \times 1024}$};

\draw[arrowstyle] (s1) -- (s2);
\draw[arrowstyle] (s2) -- (s3);
\draw[arrowstyle] (s3) -- (s4);
\draw[arrowstyle] (s4) -- (s5);
\draw[arrowstyle] (s5) -- (s6);
\draw[arrowstyle] (s6) -- (s7);

% Left annotation
\node[font=\scriptsize, text=red!70!black, rotate=90] at (-4, 5.5) {Il pixel ``esiste'' ancora come entit\`a};
\draw[thick, red!50, dashed] (-3.5, 4.5) -- (-3.5, 10.5);

\node[font=\scriptsize, text=blue!70!black, rotate=90] at (-4, 1.75) {Dissolto nei latenti};
\draw[thick, blue!50, dashed] (-3.5, 0.5) -- (-3.5, 3.5);

\end{tikzpicture}
\caption{Il viaggio di un singolo pixel attraverso il Perceiver: dal pixel RGB alla dissoluzione nei latenti tramite cross-attention. Le dimensioni crescono (da 3 a 261 a 512 a 1024) man mano che il pixel viene arricchito con informazione posizionale e proiettato in spazi di dimensione crescente.}
\label{fig:one_pixel_journey}
\end{figure}

% Immagine del byte array che produce K e V dagli appunti
\begin{figure}[H]
\centering
\includegraphics[width=0.7\textwidth]{paper_figures/docx_byte_array_kv.png}
\caption{Il Byte Array ($M \times C$, in verde) produce le matrici Key (K) e Value (V) tramite proiezioni lineari. Ogni riga del byte array (un pixel, un campione audio, un punto 3D) viene proiettata nello spazio delle chiavi e dei valori.}
\label{fig:docx_byte_array_kv}
\end{figure}

%% --- A3: Proiezione Lineare ---
\subsection{A3. Proiezione Lineare dell'Input} \label{sec:proiezione}

\begin{pipelinebox}
Input $\to$ Fourier $\to$ \textbf{Embedding} $\to$ Latent Array $\to$ Cross-Att $\to$ Res+MLP $\to$ Latent Transf.\ $\to$ Res+MLP $\to$ ($\times T$) $\to$ Pool $\to$ Classifier $\to$ Softmax $\to$ Loss
\end{pipelinebox}

Dopo l'arricchimento con Fourier features, viene applicata una \textbf{proiezione lineare} che porta la dimensione del canale da 261 a 512:

\begin{formulabox}[Proiezione lineare dell'input]
\[
    \mathbf{x}_{\text{embed}} = \tilde{\mathbf{x}} \cdot \mathbf{W}_{\text{embed}} + \mathbf{b}_{\text{embed}}
\]
dove $\mathbf{W}_{\text{embed}} \in \mathbb{R}^{261 \times 512}$, $\mathbf{b}_{\text{embed}} \in \mathbb{R}^{512}$. Risultato: $\mathbf{x}_{\text{embed}} \in \mathbb{R}^{50{,}176 \times 512}$.

\textbf{Parametri di questa proiezione}: $261 \times 512 + 512 = 134{,}144 \approx 0.13$M.

\textbf{Inizializzazione}: Xavier Uniform, che inizializza i pesi nell'intervallo:
\[
    w \sim \mathcal{U}\!\left(-\sqrt{\frac{6}{261 + 512}},\; \sqrt{\frac{6}{261 + 512}}\right) = \mathcal{U}(-0.088, 0.088)
\]
\end{formulabox}

\begin{notabox}[In parole semplici: la proiezione lineare]
Detto in parole povere, la proiezione lineare \`e come tradurre una lingua in un'altra. I pixel ``parlano'' in un linguaggio a 261 dimensioni (3 di colore + 258 di posizione), ma i moduli interni del Perceiver ``capiscono'' solo il linguaggio a 512 dimensioni. La proiezione \`e il ``traduttore'' che converte ogni pixel nel formato giusto, mescolando informazioni di colore e posizione in una rappresentazione integrata.
\end{notabox}

\textbf{Perch\'e passare da uno spazio 261 $\to$ 512?}
\begin{enumerate}
    \item \textbf{Uniformit\`a del modello}: tutti i blocchi successivi (Q, K, V, latenti) lavorano con embedding fissi $D = 512$. Avere la stessa dimensione semplifica tutte le moltiplicazioni matriciali e garantisce compatibilit\`a tra i moduli.
    \item \textbf{Maggiore capacit\`a rappresentativa}: 512 dimensioni offrono pi\`u ``spazio'' per codificare pattern complessi. Passare da 261 a 512 \`e come dare al modello pi\`u ``carta'' su cui scrivere.
    \item \textbf{Compatibilit\`a con i latenti}: i latenti hanno dimensione $D = 512$. Se l'input avesse dimensione diversa, la cross-attention richiederebbe matrici di proiezione con forme diverse per Q e K/V, complicando il design.
    \item \textbf{Stabilit\`a del training}: la proiezione lineare (inizializzata con Xavier init) distribuisce meglio l'informazione nello spazio 512-dimensionale, evitando che i valori siano troppo concentrati o sparsi.
    \item \textbf{Combinazione colore + posizione}: la proiezione lineare \textbf{mescola} i 3 valori RGB e i 258 valori posizionali in una rappresentazione integrata. Il modello impara automaticamente quanto ``pesare'' il colore rispetto alla posizione.
\end{enumerate}

\begin{examplebox}[Esempio: cosa fa concretamente la proiezione lineare]
Un singolo pixel con vettore $\tilde{\mathbf{x}}_i \in \mathbb{R}^{261}$ viene trasformato:
\[
    [\mathbf{x}_{\text{embed}}]_i = \tilde{\mathbf{x}}_i \cdot \mathbf{W}_{\text{embed}} + \mathbf{b}_{\text{embed}} \in \mathbb{R}^{512}
\]
Ogni componente $j$ dell'output \`e una \emph{combinazione lineare} di tutti i 261 valori dell'input:
\[
    [\mathbf{x}_{\text{embed}}]_{i,j} = \sum_{k=1}^{261} \tilde{x}_{i,k} \cdot W_{k,j} + b_j = \underbrace{\sum_{k=1}^{3} \tilde{x}_{i,k} W_{k,j}}_{\text{contributo RGB}} + \underbrace{\sum_{k=4}^{261} \tilde{x}_{i,k} W_{k,j}}_{\text{contributo posizionale}} + b_j
\]
Dopo il training, il modello potrebbe apprendere che certe dimensioni dell'embedding sono dominate dal colore (pesi grandi per $k = 1, 2, 3$) e altre dalla posizione (pesi grandi per $k > 3$).
\end{examplebox}

\begin{notabox}[Nota: la proiezione \`e opzionale]
In alcune configurazioni del Perceiver, la proiezione lineare dell'input \`e omessa: le Fourier features vengono fornite direttamente alla cross-attention, che ha le proprie matrici di proiezione $\mathbf{W}_K$ e $\mathbf{W}_V$. Queste matrici svolgono implicitamente lo stesso ruolo. La proiezione esplicita \`e preferita quando $C + d_{\text{PE}}$ \`e molto diverso da $D$, per garantire che le dimensioni siano compatibili e l'informazione sia distribuita uniformemente.
\end{notabox}


%% --- A4: Latent Array ---
\subsection{A4. Inizializzazione dell'Array Latente} \label{sec:latent_init}

\begin{pipelinebox}
Input $\to$ Fourier $\to$ Embedding $\to$ \textbf{Latent Array} $\to$ Cross-Att $\to$ Res+MLP $\to$ Latent Transf.\ $\to$ Res+MLP $\to$ ($\times T$) $\to$ Pool $\to$ Classifier $\to$ Softmax $\to$ Loss
\end{pipelinebox}

Questo \`e un concetto chiave che merita una spiegazione attenta. L'array latente $\mathbf{z}_0 \in \mathbb{R}^{N \times D}$ \`e una matrice di vettori latenti interamente \textbf{appresa} (parametri del modello). In pratica, quello che succede \`e che il modello ha una ``memoria di lavoro'' predefinita che usa come punto di partenza per elaborare qualsiasi input.
\begin{itemize}
    \item $N = 512$: numero di latenti (slot / ``unit\`a di memoria'').
    \item $D = 1024$: dimensione del canale/embedding dei latenti.
\end{itemize}

Questi vettori \textbf{non derivano direttamente dall'input} e non hanno un legame spaziale con token/pixel specifici: sono una base compatta che il modello usa per riassumere l'input. Sono gli stessi per ogni campione del dataset. L'adattamento all'input specifico avviene esclusivamente attraverso la \hyperref[sec:cross_attention]{cross-attention}.

\begin{formulabox}[Inizializzazione dei latenti]
\[
    [\mathbf{z}_0]_{ij} \sim \text{TruncatedNormal}(\mu = 0,\; \sigma = 0.02,\; \text{bounds} = [-0.04, 0.04])
\]
per $i = 1, \ldots, N$, $j = 1, \ldots, D$. Questi sono parametri del modello ottimizzati durante l'addestramento tramite backpropagation.
\end{formulabox}

\textbf{Perch\'e si usa TruncatedNormal$(0, 0.02)$?} La scelta non \`e casuale --- ci sono 4 ragioni fondamentali:
\begin{enumerate}
    \item \textbf{Simmetria iniziale}: media zero garantisce che i latenti non abbiano bias iniziale verso nessuna direzione. All'inizio del training, tutti i latenti ``partono uguali'' e si differenzieranno durante l'addestramento.
    \item \textbf{Scala appropriata}: $\sigma = 0.02$ \`e piccola ma non troppo. Valori troppo grandi causerebbero instabilit\`a nei primi passi di training; valori troppo piccoli renderebbero i latenti indistinguibili.
    \item \textbf{Troncamento}: eliminare i valori fuori da $[-0.04, 0.04]$ evita outlier estremi che potrebbero destabilizzare la prima iterazione di cross-attention.
    \item \textbf{Compatibilit\`a con LayerNorm}: i valori piccoli e centrati su zero sono ben gestiti dalla \hyperref[sec:layernorm]{Layer Normalization} applicata prima della cross-attention.
\end{enumerate}

\begin{examplebox}[Dimensioni e memoria dell'array latente]
L'array latente $\mathbf{z}_0 \in \mathbb{R}^{512 \times 1024}$:
\begin{itemize}
    \item \textbf{Numero di parametri}: $512 \times 1024 = 524{,}288 \approx 0.5$M.
    \item \textbf{Memoria}: $524{,}288 \times 4$ byte (float32) $= 2.1$ MB.
    \item \textbf{Confronto con l'input}: l'input arricchito $\mathbb{R}^{50{,}176 \times 261}$ occupa $50{,}176 \times 261 \times 4 \approx 52.4$ MB. I latenti sono $\sim 25\times$ pi\`u compatti.
\end{itemize}

\textbf{Analogia}: l'array latente \`e come la ``memoria di lavoro'' di un essere umano. Non possiamo tenere a mente tutti i $50{,}176$ pixel dell'immagine simultaneamente, ma possiamo tenere a mente $\sim$512 ``concetti'' (``c'\`e rosso in alto a sinistra'', ``c'\`e un bordo diagonale'', ecc.) che riassumono l'informazione essenziale.
\end{examplebox}

% TikZ: Latent array as memory slots
\begin{figure}[H]
\centering
\begin{tikzpicture}[
    slot/.style={draw, minimum width=0.45cm, minimum height=1.0cm, font=\tiny},
    arrowstyle/.style={-{Stealth[length=2mm]}, thick},
    >=Stealth,
    scale=0.85, transform shape
]

% Before cross-attention: empty slots
\node[font=\small\bfseries] at (3, 3.5) {Prima della cross-attention};
\foreach \i in {0,...,11} {
    \node[slot, fill=gray!10] at (\i*0.55, 2) {};
}
\node[font=\tiny] at (6.2, 2) {$\cdots$};
\foreach \i in {0,...,3} {
    \node[slot, fill=gray!10] at (7.0+\i*0.55, 2) {};
}
\node[font=\scriptsize, gray] at (4.5, 1.0) {512 slot ``vuoti'' (inizializzazione casuale $\sim 0.02$)};

% Arrow
\draw[arrowstyle, very thick, blue!60!black] (4.5, 0.6) -- (4.5, -0.2);
\node[font=\scriptsize, blue!60!black] at (6.5, 0.2) {Cross-Attention (lettura dall'immagine)};

% After cross-attention: filled slots
\node[font=\small\bfseries] at (3, -0.8) {Dopo la cross-attention};
\node[slot, fill=red!30] at (0, -1.8) {};
\node[slot, fill=blue!30] at (0.55, -1.8) {};
\node[slot, fill=green!30] at (1.1, -1.8) {};
\node[slot, fill=orange!30] at (1.65, -1.8) {};
\node[slot, fill=purple!20] at (2.2, -1.8) {};
\node[slot, fill=yellow!30] at (2.75, -1.8) {};
\node[slot, fill=red!20] at (3.3, -1.8) {};
\node[slot, fill=cyan!30] at (3.85, -1.8) {};
\node[slot, fill=pink!30] at (4.4, -1.8) {};
\node[slot, fill=brown!20] at (4.95, -1.8) {};
\node[slot, fill=olive!20] at (5.5, -1.8) {};
\node[slot, fill=teal!20] at (6.05, -1.8) {};
\node[font=\tiny] at (6.6, -1.8) {$\cdots$};
\node[slot, fill=violet!20] at (7.0, -1.8) {};
\node[slot, fill=lime!20] at (7.55, -1.8) {};
\node[slot, fill=magenta!20] at (8.1, -1.8) {};
\node[slot, fill=gray!30] at (8.65, -1.8) {};
\node[font=\scriptsize, gray] at (4.5, -2.8) {512 slot ``riempiti'' con informazione dall'immagine};

% Labels for some slots
\node[font=\tiny, text=red!70!black] at (0, -3.2) {bordi};
\node[font=\tiny, text=blue!70!black] at (1.1, -3.2) {cielo};
\node[font=\tiny, text=green!70!black] at (2.2, -3.2) {erba};
\node[font=\tiny, text=orange!70!black] at (3.3, -3.2) {oggetto};

\end{tikzpicture}
\caption{L'array latente come ``slot di memoria''. Prima della cross-attention, i 512 slot contengono valori quasi casuali (inizializzazione TruncatedNormal). Dopo la cross-attention, ogni slot si \`e ``riempito'' con informazione specifica estratta dall'immagine: un latente potrebbe codificare i bordi, un altro il cielo, un altro l'oggetto principale.}
\label{fig:latent_memory_slots}
\end{figure}

% Immagine della matrice latente con valori reali dagli appunti
\begin{figure}[H]
\centering
\includegraphics[width=0.4\textwidth]{paper_figures/docx_latent_matrix.png}
\caption{Esempio della matrice latente $\mathbf{z}_0 \in \mathbb{R}^{512 \times 512}$ con valori reali. I valori sono molto piccoli (tipicamente tra $-0.03$ e $+0.03$) perch\'e inizializzati con TruncatedNormal($\mu=0, \sigma=0.02$). Ogni riga rappresenta un latente, ogni colonna una dimensione del vettore latente.}
\label{fig:docx_latent_matrix}
\end{figure}

\begin{notabox}[I latenti sono condivisi tra tutte le immagini del dataset]
Un punto fondamentale: $\mathbf{z}_0$ \`e lo \textbf{stesso} per ogni immagine del dataset. Non ci sono latenti ``specializzati per gatti'' e latenti ``specializzati per aerei'': gli stessi 512 vettori vengono usati come punto di partenza per \emph{qualsiasi} immagine. L'adattamento all'immagine specifica avviene interamente attraverso la cross-attention: i latenti ``leggono'' l'immagine corrente e si aggiornano di conseguenza. Dopo la prima iterazione, $\mathbf{z}_1$ \`e gi\`a diverso per immagini diverse; dopo 8 iterazioni, $\mathbf{z}_8$ codifica informazione altamente specifica per l'immagine corrente.
\end{notabox}


%% --- A5: Cross-Attention ---
\subsection{A5. Blocco di Cross-Attention --- Zoom completo} \label{sec:ca_detail}

\begin{pipelinebox}
Input $\to$ Fourier $\to$ Embedding $\to$ Latent Array $\to$ \textbf{Cross-Attention} $\to$ Res+MLP $\to$ Latent Transf.\ $\to$ Res+MLP $\to$ ($\times T$) $\to$ Pool $\to$ Classifier $\to$ Softmax $\to$ Loss
\end{pipelinebox}

Siamo arrivati al momento cruciale: la cross-attention, il meccanismo che permette ai latenti di ``leggere'' l'immagine (si vedano le formule generali nella \hyperref[sec:cross_attention]{Sezione~\ref*{sec:cross_attention}} e il confronto con la self-attention nella \hyperref[sec:self_vs_cross]{Sezione~\ref*{sec:self_vs_cross}}). Qui vediamo cosa succede concretamente, passo per passo, con i numeri del Perceiver su ImageNet.

\subsubsection{Passo 1: Pre-LayerNorm}

Il primo passo \`e normalizzare sia i latenti che l'input. Perch\'e? Perch\'e i latenti e l'input possono avere scale molto diverse: i latenti iniziano con valori $\sim 0.02$ (TruncatedNormal), mentre i pixel RGB variano tra 0 e 1, e le Fourier features variano tra $-1$ e $1$. Senza normalizzazione, i prodotti scalari nell'attenzione sarebbero dominati dalle componenti con valori pi\`u grandi, distorcendo la distribuzione dell'attenzione.

Si applica la \hyperref[sec:layernorm]{Layer Normalization} sia ai latenti che all'input:
\[
    \mathbf{z}_{\text{norm}} = \text{LN}(\mathbf{z}_0) \in \mathbb{R}^{512 \times 1024}, \qquad \mathbf{x}_{\text{norm}} = \text{LN}(\mathbf{x}_{\text{embed}}) \in \mathbb{R}^{50{,}176 \times 512}
\]

\subsubsection{Passo 2--4: Proiezioni Q, K, V} \label{sec:qkv_proiezioni}

Ora arriviamo al cuore della cross-attention: le proiezioni. Questo \`e il momento in cui i latenti e l'input vengono ``tradotti'' nello stesso linguaggio (query, chiavi e valori) per potersi ``confrontare''.

Le matrici di peso apprese $\mathbf{W}_K$ e $\mathbf{W}_V$ \textbf{non contengono conoscenza della rete} all'inizio: sono inizializzate con piccoli numeri casuali (es.\ Xavier Uniform). Quando creiamo le matrici di peso in PyTorch (es.\ \texttt{nn.Linear}), vengono inizializzate automaticamente con Xavier Uniform. \`E durante il training che questi pesi apprendono a estrarre informazione rilevante.

\begin{formulabox}[title={Proiezioni Q, K, V della cross-attention con dimensioni}]
\begin{align*}
    \mathbf{Q} &= \underbrace{\mathbf{z}_{\text{norm}}}_{512 \times 1024} \cdot \underbrace{\mathbf{W}_Q}_{1024 \times 1024} &&= \mathbb{R}^{512 \times 1024} & \text{(Q dai \textbf{latenti})} \\
    \mathbf{K} &= \underbrace{\mathbf{x}_{\text{norm}}}_{50176 \times 512} \cdot \underbrace{\mathbf{W}_K}_{512 \times 1024} &&= \mathbb{R}^{50{,}176 \times 1024} & \text{(K dal \textbf{byte array})} \\
    \mathbf{V} &= \underbrace{\mathbf{x}_{\text{norm}}}_{50176 \times 512} \cdot \underbrace{\mathbf{W}_V}_{512 \times 1024} &&= \mathbb{R}^{50{,}176 \times 1024} & \text{(V dal \textbf{byte array})}
\end{align*}

\textbf{Nota sulle dimensioni}: $\mathbf{W}_Q \in \mathbb{R}^{D \times d_k}$ dove $D = 1024$ (dimensione dei latenti) e $d_k = 1024$. $\mathbf{W}_K, \mathbf{W}_V \in \mathbb{R}^{C_{\text{embed}} \times d_k}$ dove $C_{\text{embed}} = 512$ (dimensione dell'input dopo la proiezione). Le dimensioni diverse di $D$ e $C_{\text{embed}}$ non sono un problema perch\'e Q e K vengono portati alla stessa dimensione $d_k$ dalle rispettive proiezioni.
\end{formulabox}

% Immagine zoom sulle proiezioni QKV dagli appunti
\begin{figure}[H]
\centering
\includegraphics[width=0.7\textwidth]{paper_figures/docx_qkv_projections.png}
\caption{Zoom sulle proiezioni Q, K, V. Il Latent Array ($N \times D$, in viola) genera le Query tramite $f_Q$. Il Byte Array ($M \times C$, in verde) genera Key e Value tramite $f_K$ e $f_V$. Si noti che $N \ll M$: questa asimmetria \`e la chiave dell'efficienza.}
\label{fig:docx_qkv_projections}
\end{figure}

\textbf{Parametri delle proiezioni Q, K, V}:
\begin{itemize}
    \item $\mathbf{W}_Q$: $1024 \times 1024 = 1{,}048{,}576$ parametri
    \item $\mathbf{W}_K$: $512 \times 1024 = 524{,}288$ parametri
    \item $\mathbf{W}_V$: $512 \times 1024 = 524{,}288$ parametri
    \item Totale (per un modulo di cross-attention): $\sim 2.1$M parametri (solo per le proiezioni)
\end{itemize}

\begin{notabox}[title={Intuizione: Query, Key, Value}]
Ogni latente \`e una \textbf{Query} ($\mathbf{Q}$): rappresenta una ``domanda'' del tipo: ``Quale parte dell'immagine mi serve?''. Ogni token d'input \`e una \textbf{Key} ($\mathbf{K}$): rappresenta una ``descrizione'' del contenuto del token (es: ``sono un pixel rosso in alto a sinistra''). Il \textbf{Value} ($\mathbf{V}$) \`e l'informazione effettiva che il pixel ``porterebbe'' se venisse selezionato.

Le matrici di peso $\mathbf{W}_Q$, $\mathbf{W}_K$, $\mathbf{W}_V$ sono gli ``occhiali'' attraverso cui il modello guarda i latenti e l'input: trasformano i dati grezzi in rappresentazioni utili per il confronto.
\end{notabox}

\subsubsection{Passo 5: Scaled Dot-Product Attention} \label{sec:scaled_dot}

\begin{formulabox}[Score di attenzione]
\[
    \mathbf{S} = \frac{\mathbf{Q} \cdot \mathbf{K}^\top}{\sqrt{d_k}} \in \mathbb{R}^{512 \times 50{,}176}
\]
\end{formulabox}

Vediamo cosa succede concretamente in questo passo cruciale.

\textbf{Perch\'e facciamo la scaled-dot product attention?} Il prodotto scalare $\mathbf{S}_{ij} = \mathbf{q}_i \cdot \mathbf{k}_j$ misura quanto la ``domanda'' del latente $i$-esimo \`e compatibile con la ``descrizione'' dell'input $j$-esimo. Se il latente $i$ sta ``cercando rosso'' e il pixel $j$ \`e ``rosso'', il prodotto scalare sar\`a alto. Se il pixel $j$ \`e ``verde'', il prodotto scalare sar\`a basso.

La divisione per $\sqrt{d_k} = \sqrt{1024} = 32$ serve a evitare che gli score diventino troppo grandi: senza scaling, la softmax produrrebbe pesi quasi one-hot (un solo pixel con peso $\approx 1$, tutti gli altri $\approx 0$), rendendo i gradienti quasi nulli.

\begin{proofbox}[title={Perch\'e $\sqrt{d_k}$ e non un altro valore?}]
Supponiamo che le componenti di $\mathbf{q}$ e $\mathbf{k}$ siano i.i.d.\ con media 0 e varianza 1. Il prodotto scalare \`e:
\[
    \mathbf{q} \cdot \mathbf{k} = \sum_{i=1}^{d_k} q_i k_i
\]
Poich\'e $q_i$ e $k_i$ sono indipendenti: $\text{Var}(q_i k_i) = \text{Var}(q_i) \cdot \text{Var}(k_i) = 1 \cdot 1 = 1$ (dato che le medie sono 0).

Per la somma di $d_k$ variabili indipendenti: $\text{Var}(\mathbf{q} \cdot \mathbf{k}) = d_k$.

Quindi $\mathbf{q} \cdot \mathbf{k}$ ha deviazione standard $\sqrt{d_k}$. Dividendo per $\sqrt{d_k}$, riportiamo la deviazione standard a 1, indipendentemente dalla dimensione $d_k$.

Per $d_k = 1024$, senza scaling gli score avrebbero deviazione standard $\sqrt{1024} = 32$. Con valori di score $\sim 32$, la softmax sarebbe quasi one-hot, con gradienti $\sim 0$ per quasi tutti gli elementi. Lo scaling riporta la deviazione standard a 1, producendo una distribuzione softmax pi\`u ``morbida'' e gradienti informativi.
\end{proofbox}

\subsubsection{Passo 6: Softmax per riga} \label{sec:softmax_ca}

Ora abbiamo una matrice di score $512 \times 50.176$ con numeri qualsiasi (anche negativi). La domanda che ci si pone \`e: come trasformiamo questi numeri in ``pesi di importanza'' utilizzabili? Servono due propriet\`a: (1) essere positivi, (2) sommare a 1 per ogni latente (cos\`i che possiamo interpretarli come ``quanto ogni pixel \`e importante per questo latente'').

La \hyperref[sec:softmax]{softmax} fa esattamente questo --- ed \`e la stessa funzione che abbiamo gi\`a incontrato nella \hyperref[sec:softmax]{Sezione~2.1}:

\[
    A_{ij} = \frac{\exp(S_{ij})}{\sum_{j'=1}^{50{,}176} \exp(S_{ij'})} \qquad \text{Risultato: } \mathbf{A} \in \mathbb{R}^{512 \times 50{,}176}, \quad \sum_j A_{ij} = 1 \;\;\forall\, i
\]

Ogni riga $i$ della matrice $\mathbf{A}$ \`e una distribuzione di probabilit\`a sui $50{,}176$ pixel: indica ``quanto ogni pixel \`e rilevante per il latente $i$''.

\begin{examplebox}[Visualizzazione della matrice di attenzione]
La matrice $\mathbf{A} \in \mathbb{R}^{512 \times 50{,}176}$ pu\`o essere ``reshapata'' in $\mathbb{R}^{512 \times 224 \times 224}$ per visualizzarla come 512 mappe di calore sull'immagine. Ogni mappa mostra ``dove guarda'' un latente:

\begin{itemize}
    \item \textbf{Latente 1}: mappa concentrata nell'angolo in alto a sinistra $\to$ questo latente si \`e ``specializzato'' per quella regione.
    \item \textbf{Latente 47}: mappa diffusa su tutto l'oggetto principale $\to$ cattura informazione globale sull'oggetto.
    \item \textbf{Latente 200}: mappa concentrata sui bordi $\to$ si \`e specializzato per le texture di contorno.
    \item \textbf{Latente 400}: mappa diffusa e quasi uniforme $\to$ cattura informazione di sfondo.
\end{itemize}

Le mappe di attenzione del paper (Figure 3) confermano questa specializzazione emergente: i latenti si differenziano automaticamente durante il training, senza supervisione esplicita.
\end{examplebox}

\begin{notabox}[title={La matrice $\mathbf{A}$ come mappa di importanza}]
La matrice di attenzione codifica una forma di \textbf{soft selection}: ogni latente ``seleziona'' softly (con pesi continui) i pixel rilevanti. Questo \`e fondamentalmente diverso dalla hard selection (come il cropping o il pooling), perch\'e:
\begin{enumerate}
    \item \`E \textbf{differenziabile}: i pesi $A_{ij}$ possono essere ottimizzati via gradient descent.
    \item \`E \textbf{adattiva}: la selezione cambia per ogni immagine (i pesi dipendono dall'input tramite Q e K).
    \item \`E \textbf{morbida}: nessun pixel viene completamente ignorato (tutti hanno peso $> 0$).
\end{enumerate}
\end{notabox}

% Immagine zoom sulla mappa di attenzione dagli appunti
\begin{figure}[H]
\centering
\includegraphics[width=0.4\textwidth]{paper_figures/docx_attention_map.png}
\caption{Zoom sul calcolo della mappa di attenzione. La matrice $\mathbf{A} = \text{softmax}(\mathbf{Q}\mathbf{K}^\top)$ ha dimensione $N \times M$ (rettangolare, non quadrata). La complessit\`a passa da $\mathcal{O}(M^2)$ del Transformer standard a $\mathcal{O}(MN)$ nel Perceiver.}
\label{fig:docx_attention_map}
\end{figure}

Il punto cruciale \`e che la matrice $\mathbf{A}$ \`e \textbf{rettangolare} ($N \times M$) e non quadrata ($M \times M$): con $N = 512$ e $M = 50{,}176$, la matrice contiene $\sim$25.7 milioni di elementi anzich\'e $\sim$2.5 miliardi.

\subsubsection{Passo 7: Combinazione pesata}

Questo \`e il momento magico: i pesi di attenzione vengono usati per ``pesare'' i valori di tutti i pixel, producendo l'output della cross-attention.

\[
    \underbrace{\text{Out}}_{512 \times 1024} = \underbrace{\mathbf{A}}_{512 \times 50{,}176} \times \underbrace{\mathbf{V}}_{50{,}176 \times 1024}
\]

Ogni latente $i$ produce il suo output come \textbf{media pesata} dei valori di tutti i pixel: $\text{out}_i = \sum_{j=1}^{50{,}176} A_{ij} \mathbf{v}_j$. I pixel pi\`u ``rilevanti'' (peso $A_{ij}$ alto) contribuiscono di pi\`u; i pixel irrilevanti ($A_{ij} \approx 0$) sono quasi ignorati.

% Immagine zoom sul calcolo delle feature F=AV dagli appunti
\begin{figure}[H]
\centering
\includegraphics[width=0.4\textwidth]{paper_figures/docx_compute_features.png}
\caption{Zoom sul calcolo delle feature ($\mathbf{F} = \mathbf{A}\mathbf{V}$). L'output ha dimensione $N \times D$: l'informazione di $M$ elementi dell'input \`e stata compressa in $N$ latenti. La complessit\`a di questo passo \`e $\mathcal{O}(MN)$ invece di $\mathcal{O}(M^2)$.}
\label{fig:docx_compute_features}
\end{figure}

La Figura~\ref{fig:docx_compute_features} visualizza il passo finale della cross-attention: la moltiplicazione $\mathbf{F} = \mathbf{A}\mathbf{V}$. La matrice di attenzione $\mathbf{A}$ ($N \times M$) moltiplica la matrice dei valori $\mathbf{V}$ ($M \times D$), producendo un output di dimensione $N \times D$. Questo \`e il momento in cui avviene la compressione: l'informazione di tutti gli $M$ pixel dell'input viene ``riassunta'' in $N$ vettori latenti, ciascuno dei quali ha catturato una media pesata delle feature pi\`u rilevanti dall'input.

\subsubsection{Passo 8--9: Proiezione di output e residual connection}

Siamo quasi alla fine del primo blocco di cross-attention. L'ultimo passo \`e combinare l'output della cross-attention con l'input originale tramite la residual connection.

\[
    \mathbf{z}_{\text{after\_ca}} = \underbrace{\mathbf{z}_0}_{\text{originale}} + \underbrace{\text{Out} \cdot \mathbf{W}_O}_{\text{output cross-attention}}
\]

La \hyperref[sec:residual]{connessione residua} somma l'input originale all'output della cross-attention. Questo step serve a:
\begin{enumerate}
    \item \textbf{Non perdere l'informazione originaria dei latenti}: anche se la cross-attention ``sbaglia'', il latente mantiene la sua rappresentazione precedente.
    \item \textbf{Prevenire la degenerazione}: se l'attenzione produce valori poco informativi, la residual connection preserva il segnale.
    \item \textbf{Stabilizzare il flusso del gradiente}: durante la backpropagation, il gradiente pu\`o fluire direttamente attraverso lo skip connection (si veda la \hyperref[sec:residual]{Sezione~2.4}).
\end{enumerate}

\begin{examplebox}[Intuizione: come funziona la cross-attention]
Con 3 pixel e 1 latente che ``cerca il rosso'': Pixel A (rosso, score 4.2), Pixel B (verde, score $-1.5$), Pixel C (arancione, score 2.1). Dopo softmax: $\alpha_A \approx 0.89$, $\alpha_B \approx 0.003$, $\alpha_C \approx 0.107$. Il latente ``assorbe'' soprattutto l'informazione dal pixel rosso.
\end{examplebox}

\begin{examplebox}[Esempio numerico completo: cross-attention con 3 pixel e 1 latente]
Consideriamo un caso semplificato con 1 latente (``cerca rosso'') e 3 pixel, con $d_k = 3$:

\textbf{Query} (dal latente, dopo proiezione $\mathbf{W}_Q$):
$\mathbf{q}_1 = [1.5,\; 0.2,\; -0.3]$ --- codifica ``cerco rosso''

\textbf{Key} (dai pixel, dopo proiezione $\mathbf{W}_K$):
\begin{align*}
    \mathbf{k}_A &= [1.8,\; 0.1,\; -0.5] \quad \text{(pixel A: rosso vivo)} \\
    \mathbf{k}_B &= [-0.3,\; 1.7,\; 0.4] \quad \text{(pixel B: verde)} \\
    \mathbf{k}_C &= [0.9,\; 0.5,\; -0.1] \quad \text{(pixel C: arancione)}
\end{align*}

\textbf{Passo 1 --- Score (prodotto scalare $\mathbf{q} \cdot \mathbf{k}$)}:
\begin{align*}
    s_A &= 1.5 \times 1.8 + 0.2 \times 0.1 + (-0.3) \times (-0.5) = 2.70 + 0.02 + 0.15 = 2.87 \\
    s_B &= 1.5 \times (-0.3) + 0.2 \times 1.7 + (-0.3) \times 0.4 = -0.45 + 0.34 - 0.12 = -0.23 \\
    s_C &= 1.5 \times 0.9 + 0.2 \times 0.5 + (-0.3) \times (-0.1) = 1.35 + 0.10 + 0.03 = 1.48
\end{align*}

\textbf{Passo 2 --- Scaling}: $\sqrt{d_k} = \sqrt{3} \approx 1.732$
\[
    \hat{s}_A = 2.87 / 1.732 = 1.657, \quad \hat{s}_B = -0.23 / 1.732 = -0.133, \quad \hat{s}_C = 1.48 / 1.732 = 0.855
\]

\textbf{Passo 3 --- Softmax}:
\begin{align*}
    e^{1.657} &= 5.243, \quad e^{-0.133} = 0.875, \quad e^{0.855} = 2.351 \\
    \text{somma} &= 5.243 + 0.875 + 2.351 = 8.469 \\
    \alpha_A &= 5.243 / 8.469 = \mathbf{0.619} \\
    \alpha_B &= 0.875 / 8.469 = \mathbf{0.103}, \quad \alpha_C = 2.351 / 8.469 = \mathbf{0.278}
\end{align*}

\textbf{Passo 4 --- Media pesata dei Value}. Se $\mathbf{v}_A = [1.0, 0.0]$, $\mathbf{v}_B = [0.0, 1.0]$, $\mathbf{v}_C = [0.6, 0.3]$:
\[
    \text{out} = 0.619 \!\times\! [1.0, 0.0] + 0.103 \!\times\! [0.0, 1.0] + 0.278 \!\times\! [0.6, 0.3]
\]
\[
    = [0.619 + 0 + 0.167,\; 0 + 0.103 + 0.083] = [0.786,\; 0.186]
\]

Il latente ha estratto un vettore che \`e prevalentemente ``rosso'' ($0.786$), con un contributo minore dalle altre componenti.
\end{examplebox}

% TikZ: Intuitive cross-attention diagram
\begin{figure}[H]
\centering
\begin{tikzpicture}[
    pixel/.style={draw, rounded corners=2pt, minimum width=1.5cm, minimum height=0.7cm, font=\small},
    latent/.style={draw, rounded corners=2pt, fill=blue!15, minimum width=2cm, minimum height=0.7cm, font=\small},
    arrowstyle/.style={-{Stealth[length=2mm]}, thick},
    >=Stealth
]

% Latent
\node[latent] (lat) at (5, 4) {Latente: ``cerca rosso''};

% Pixels
\node[pixel, fill=red!30] (pA) at (0, 0) {Pixel A (rosso)};
\node[pixel, fill=green!30] (pB) at (5, 0) {Pixel B (verde)};
\node[pixel, fill=orange!30] (pC) at (10, 0) {Pixel C (arancione)};

% Attention weights
\draw[arrowstyle, red!70!black, line width=2.5pt] (lat.south) -- (pA.north)
    node[midway, left, font=\small\bfseries, red!70!black] {$\alpha = 0.619$};
\draw[arrowstyle, green!50!black, line width=0.5pt] (lat.south) -- (pB.north)
    node[midway, right, font=\small, green!50!black] {$\alpha = 0.103$};
\draw[arrowstyle, orange!70!black, line width=1.5pt] (lat.south) -- (pC.north)
    node[midway, right, font=\small, orange!70!black] {$\alpha = 0.278$};

% Annotation
\node[font=\small\itshape, gray] at (5, -1.2) {Lo spessore della freccia \`e proporzionale al peso di attenzione};
\node[font=\small\itshape, gray] at (5, -1.7) {Il latente ``assorbe'' soprattutto l'informazione dal pixel rosso};

\end{tikzpicture}
\caption{Esempio intuitivo di cross-attention con 1 latente e 3 pixel. Il latente ``cerca rosso'' assegna il peso maggiore ($0.619$) al pixel rosso, un peso intermedio ($0.278$) al pixel arancione (contiene del rosso), e un peso minimo ($0.103$) al pixel verde.}
\label{fig:intuitive_cross_attention}
\end{figure}

\begin{examplebox}[Esempio numerico completo: cross-attention con dimensioni reali]
Consideriamo un esempio con $M = 4$ pixel, $N = 2$ latenti, $d = 3$ (dimensioni ridotte per chiarezza). Senza multi-head per semplicit\`a.

\textbf{Input --- 4 pixel con colori semantici:}
\[
\mathbf{x} = \begin{pmatrix} 1 & 0 & 0 \\ 0 & 1 & 0 \\ 0 & 0 & 1 \\ 1 & 1 & 0 \end{pmatrix} \quad \begin{array}{l} p_1 = \text{rosso} \\ p_2 = \text{verde} \\ p_3 = \text{blu} \\ p_4 = \text{giallo} \end{array}
\]

\textbf{Latenti iniziali ($N = 2$):}
\[
\mathbf{z} = \begin{pmatrix} 0.1 & 0.2 & 0.1 \\ 0.2 & 0.1 & 0.3 \end{pmatrix} \quad \begin{array}{l} z_1: \text{``cerca colori caldi?''} \\ z_2: \text{``cerca colori freddi?''} \end{array}
\]

\textbf{Matrici di peso} (semplificate, $3 \times 3$):
\[
\mathbf{W}_Q = \begin{pmatrix} 1 & 0 & 0 \\ 0 & 1 & 0 \\ 0 & 0 & 1 \end{pmatrix}, \quad
\mathbf{W}_K = \begin{pmatrix} 1 & 0 & 0.5 \\ 0 & 1 & 0 \\ 0.5 & 0 & 1 \end{pmatrix}, \quad
\mathbf{W}_V = \begin{pmatrix} 1 & 0 & 0 \\ 0 & 1 & 0 \\ 0 & 0 & 1 \end{pmatrix}
\]

\textbf{Passo 1 --- Calcolo di Q ($2 \times 3$):}
\[
\mathbf{Q} = \mathbf{z} \cdot \mathbf{W}_Q = \begin{pmatrix} 0.1 & 0.2 & 0.1 \\ 0.2 & 0.1 & 0.3 \end{pmatrix}
\]

\textbf{Passo 2 --- Calcolo di K ($4 \times 3$):}
\[
\mathbf{K} = \mathbf{x} \cdot \mathbf{W}_K = \begin{pmatrix} 1 & 0 & 0.5 \\ 0 & 1 & 0 \\ 0.5 & 0 & 1 \\ 1 & 1 & 0.5 \end{pmatrix}
\]

\textbf{Passo 3 --- Calcolo di V ($4 \times 3$):}
$\mathbf{V} = \mathbf{x} \cdot \mathbf{W}_V = \mathbf{x}$ (perch\'e $\mathbf{W}_V = \mathbf{I}$).

\textbf{Passo 4 --- Score $\mathbf{Q}\mathbf{K}^\top$ ($2 \times 4$):}
\begin{align*}
\mathbf{S}_{11} &= \mathbf{q}_1 \cdot \mathbf{k}_1 = 0.1 \times 1 + 0.2 \times 0 + 0.1 \times 0.5 = 0.15 \\
\mathbf{S}_{12} &= \mathbf{q}_1 \cdot \mathbf{k}_2 = 0.1 \times 0 + 0.2 \times 1 + 0.1 \times 0 = 0.20 \\
\mathbf{S}_{13} &= \mathbf{q}_1 \cdot \mathbf{k}_3 = 0.1 \times 0.5 + 0.2 \times 0 + 0.1 \times 1 = 0.15 \\
\mathbf{S}_{14} &= \mathbf{q}_1 \cdot \mathbf{k}_4 = 0.1 \times 1 + 0.2 \times 1 + 0.1 \times 0.5 = 0.35
\end{align*}

Analogamente per la riga 2:
\[
\mathbf{S} = \begin{pmatrix} 0.15 & 0.20 & 0.15 & 0.35 \\ 0.35 & 0.10 & 0.40 & 0.35 \end{pmatrix}
\]

\textbf{Passo 5 --- Scaling per $\sqrt{d_k} = \sqrt{3} \approx 1.732$:}
\[
\mathbf{S}_{\text{scaled}} = \begin{pmatrix} 0.087 & 0.116 & 0.087 & 0.202 \\ 0.202 & 0.058 & 0.231 & 0.202 \end{pmatrix}
\]

\textbf{Passo 6 --- Softmax riga per riga:}

\textit{Riga 1}: $e^{0.087} = 1.091$, $e^{0.116} = 1.123$, $e^{0.087} = 1.091$, $e^{0.202} = 1.224$. Somma $= 4.529$.
\[
\mathbf{A}_1 = \tfrac{1}{4.529}[1.091,\; 1.123,\; 1.091,\; 1.224] = [0.241,\; 0.248,\; 0.241,\; 0.270]
\]

\textit{Riga 2}: $e^{0.202} = 1.224$, $e^{0.058} = 1.060$, $e^{0.231} = 1.260$, $e^{0.202} = 1.224$. Somma $= 4.768$.
\[
\mathbf{A}_2 = \tfrac{1}{4.768}[1.224,\; 1.060,\; 1.260,\; 1.224] = [0.257,\; 0.222,\; 0.264,\; 0.257]
\]

\textbf{Matrice di attenzione completa} ($2 \times 4$):
\[
\mathbf{A} = \begin{pmatrix} 0.241 & 0.248 & 0.241 & 0.270 \\ 0.257 & 0.222 & 0.264 & 0.257 \end{pmatrix}
\]

\textbf{Passo 7 --- Output $\mathbf{A} \cdot \mathbf{V}$ ($2 \times 3$):}
\begin{align*}
\text{out}_1 &= 0.241 \times [1,0,0] + 0.248 \times [0,1,0] + 0.241 \times [0,0,1] + 0.270 \times [1,1,0] \\
&= [0.241 + 0.270,\; 0.248 + 0.270,\; 0.241] = [0.511,\; 0.518,\; 0.241]
\end{align*}

\textbf{Interpretazione}:
\begin{itemize}
\item Il \textbf{latente 1} ha prestato attenzione soprattutto al pixel giallo ($\alpha = 0.270$), poi al verde ($0.248$). Ha estratto una miscela di colori caldi.
\item Il \textbf{latente 2} ha prestato attenzione soprattutto al pixel blu ($\alpha = 0.264$), poi al rosso e al giallo ($0.257$). Ha estratto una miscela leggermente sbilanciata verso il freddo.
\item In questo esempio le differenze sono piccole perch\'e i latenti iniziali sono simili e i pesi sono semplici. Con il training, le query diventerebbero molto pi\`u selettive!
\end{itemize}
\end{examplebox}

% TikZ: Attention matrix heatmap
\begin{figure}[H]
\centering
\begin{tikzpicture}[scale=0.9, transform shape]

% Title
\node[font=\small\bfseries] at (3.5, 4) {Matrice di attenzione $\mathbf{A} \in \mathbb{R}^{2 \times 4}$ (latenti $\times$ pixel)};

% Column labels (pixels)
\node[font=\scriptsize\bfseries] at (1.5, 3.2) {$p_1$ (rosso)};
\node[font=\scriptsize\bfseries] at (3, 3.2) {$p_2$ (verde)};
\node[font=\scriptsize\bfseries] at (4.5, 3.2) {$p_3$ (blu)};
\node[font=\scriptsize\bfseries] at (6, 3.2) {$p_4$ (giallo)};

% Row labels (latents)
\node[font=\scriptsize\bfseries, anchor=east] at (0.5, 2.25) {$z_1$};
\node[font=\scriptsize\bfseries, anchor=east] at (0.5, 1.25) {$z_2$};

% Heatmap cells - Row 1 (z1)
\fill[green!24] (0.75, 1.75) rectangle (2.25, 2.75);
\fill[green!25] (2.25, 1.75) rectangle (3.75, 2.75);
\fill[green!24] (3.75, 1.75) rectangle (5.25, 2.75);
\fill[green!27] (5.25, 1.75) rectangle (6.75, 2.75);

% Heatmap cells - Row 2 (z2)
\fill[blue!26] (0.75, 0.75) rectangle (2.25, 1.75);
\fill[blue!22] (2.25, 0.75) rectangle (3.75, 1.75);
\fill[blue!26] (3.75, 0.75) rectangle (5.25, 1.75);
\fill[blue!26] (5.25, 0.75) rectangle (6.75, 1.75);

% Grid
\draw[thick] (0.75, 0.75) grid[step=1.5] (6.75, 2.75);
\draw[very thick] (0.75, 0.75) rectangle (6.75, 2.75);

% Values in cells - Row 1
\node[font=\small] at (1.5, 2.25) {0.241};
\node[font=\small] at (3.0, 2.25) {0.248};
\node[font=\small] at (4.5, 2.25) {0.241};
\node[font=\small\bfseries] at (6.0, 2.25) {0.270};

% Values in cells - Row 2
\node[font=\small] at (1.5, 1.25) {0.257};
\node[font=\small] at (3.0, 1.25) {0.222};
\node[font=\small\bfseries] at (4.5, 1.25) {0.264};
\node[font=\small] at (6.0, 1.25) {0.257};

% Annotations
\node[font=\scriptsize, gray] at (3.75, 0.2) {Ogni riga somma a 1 (softmax). I valori in grassetto indicano il peso massimo.};

\end{tikzpicture}
\caption{Heatmap della matrice di attenzione $\mathbf{A}$ dall'esempio numerico. Il latente $z_1$ presta la massima attenzione al pixel giallo ($p_4$, peso $0.270$). Il latente $z_2$ presta la massima attenzione al pixel blu ($p_3$, peso $0.264$). La distribuzione \`e quasi uniforme perch\'e i latenti iniziali non sono ancora addestrati.}
\label{fig:attention_heatmap_example}
\end{figure}


%% --- A6: MLP ---
\subsection{A6. Feed-Forward MLP dopo la Cross-Attention} \label{sec:mlp_forward}

\begin{pipelinebox}
Input $\to$ Fourier $\to$ Embedding $\to$ Latent Array $\to$ Cross-Att $\to$ \textbf{Residual + MLP} $\to$ Latent Transf.\ $\to$ Res+MLP $\to$ ($\times T$) $\to$ Pool $\to$ Classifier $\to$ Softmax $\to$ Loss
\end{pipelinebox}

\subsubsection{Pre-Norm con LayerNorm}

Prima di entrare nell'MLP, si applica di nuovo la \hyperref[sec:layernorm]{LayerNorm}:
\[
    \mathbf{z}_{\text{norm}} = \text{LN}(\mathbf{z}_{\text{after\_ca}}) \in \mathbb{R}^{512 \times 1024}
\]
La LayerNorm prende ogni vettore latente e: (1) calcola media e deviazione standard sulle 1024 componenti; (2) normalizza ogni valore; (3) applica parametri appresi $\boldsymbol{\gamma}$ e $\boldsymbol{\beta}$.

\subsubsection{Perch\'e serve l'MLP?}

La domanda che ci si pone \`e: se l'attenzione gi\`a mescola le informazioni, perch\'e serve anche l'MLP?

L'attenzione serve a \textbf{mescolare informazione} tra token/latenti. Per\`o matematicamente \`e una \textbf{combinazione lineare} dei valori. Se ci fermassimo l\`i, il potere espressivo sarebbe limitato: una sequenza di operazioni di attenzione senza non-linearit\`a sarebbe equivalente a una singola combinazione lineare. Con $\mathbf{W}_1$, $\mathbf{W}_2$ e la GELU (attivazione non lineare), il modello diventa capace di rappresentare \textbf{funzioni non lineari complesse}.

% TikZ: MLP expansion/compression hourglass
\begin{figure}[H]
\centering
\begin{tikzpicture}[
    arrowstyle/.style={-{Stealth[length=2.5mm]}, thick},
    >=Stealth,
    scale=0.85, transform shape
]

% Input bar (512 = D)
\node[font=\small\bfseries] at (0, 6) {Input: $D = 1024$};
\fill[blue!25] (-1.024, 5.0) rectangle (1.024, 5.4);
\draw[thick] (-1.024, 5.0) rectangle (1.024, 5.4);
\node[font=\scriptsize] at (0, 4.7) {1024 dimensioni};

% Arrow down to expansion
\draw[arrowstyle] (0, 4.5) -- (0, 3.8);
\node[font=\scriptsize, right] at (0.3, 4.15) {$\mathbf{W}_1$ (espansione)};

% Expanded bar (4D = 4096)
\node[font=\small\bfseries, red!60!black] at (0, 3.5) {Espansione: $4D = 4096$};
\fill[red!25] (-4.096, 2.4) rectangle (4.096, 2.8);
\draw[thick] (-4.096, 2.4) rectangle (4.096, 2.8);
\node[font=\scriptsize] at (0, 2.1) {4096 dimensioni --- pi\`u ``spazio di lavoro''};

% GELU
\draw[arrowstyle] (0, 1.9) -- (0, 1.4);
\node[font=\scriptsize, fill=yellow!20, draw, rounded corners=2pt] at (0, 1.0) {GELU (non-linearit\`a)};

% Arrow down to compression
\draw[arrowstyle] (0, 0.6) -- (0, -0.1);
\node[font=\scriptsize, right] at (0.3, 0.25) {$\mathbf{W}_2$ (compressione)};

% Output bar (D = 1024)
\node[font=\small\bfseries] at (0, -0.5) {Output: $D = 1024$};
\fill[green!25] (-1.024, -1.5) rectangle (1.024, -1.1);
\draw[thick] (-1.024, -1.5) rectangle (1.024, -1.1);
\node[font=\scriptsize] at (0, -1.8) {1024 dimensioni (come l'input)};

% Hourglass shape overlay
\draw[thick, blue!40, dashed] (-1.024, 5.0) -- (-4.096, 2.8);
\draw[thick, blue!40, dashed] (1.024, 5.0) -- (4.096, 2.8);
\draw[thick, blue!40, dashed] (-4.096, 2.4) -- (-1.024, -1.1);
\draw[thick, blue!40, dashed] (4.096, 2.4) -- (1.024, -1.1);

% Annotation
\node[font=\scriptsize, gray, text width=4cm, align=center] at (5.5, 2.6) {``Clessidra inversa'':\\si espande, si trasforma,\\e si ricomprime};

\end{tikzpicture}
\caption{L'MLP come ``clessidra inversa'': espande da 1024 a 4096 dimensioni (creando spazio per trasformazioni ricche), applica la non-linearit\`a GELU, e ricomprime a 1024 dimensioni (distillando l'informazione trasformata). Le dimensioni in input e output sono identiche, permettendo la connessione residua.}
\label{fig:mlp_hourglass}
\end{figure}

\begin{formulabox}[MLP --- dettaglio completo]
\textbf{Espansione} ($D \to D_{\text{ff}} = 4D = 4096$):
\[
\underbrace{\mathbf{h}}_{512 \times 4096} = \underbrace{\mathbf{z}_{\text{norm}}}_{512 \times 1024} \times \underbrace{\mathbf{W}_1}_{1024 \times 4096} + \mathbf{b}_1
\]

\textbf{Attivazione GELU}: $\mathbf{h}_{\text{act}} = \text{GELU}(\mathbf{h}) \in \mathbb{R}^{512 \times 4096}$

L'espansione a 4096 dimensioni crea ``spazio di lavoro'' per trasformazioni ricche. La GELU \`e una funzione liscia che permette piccoli valori negativi (nessun dying neuron).

\textbf{Compressione} ($D_{\text{ff}} \to D = 1024$):
\[
\underbrace{\mathbf{z}_{\text{ffn}}}_{512 \times 1024} = \underbrace{\mathbf{h}_{\text{act}}}_{512 \times 4096} \times \underbrace{\mathbf{W}_2}_{4096 \times 1024} + \mathbf{b}_2
\]

La compressione riporta alla dimensione standard, ``distillando'' le trasformazioni non lineari.
\end{formulabox}

\begin{examplebox}[Esempio numerico semplificato: un singolo latente attraverso l'MLP]
Consideriamo un singolo latente con $D = 4$ e $D_{\text{ff}} = 8$ (dimensioni ridotte per chiarezza):

\textbf{Input} (dopo LayerNorm): $\mathbf{z}_{\text{norm}} = [0.5, -0.3, 0.8, -0.1]$

\textbf{Espansione}: $\mathbf{h} = \mathbf{z}_{\text{norm}} \cdot \mathbf{W}_1 + \mathbf{b}_1 \in \mathbb{R}^{8}$

Supponiamo $\mathbf{h} = [1.2, -0.5, 0.3, -1.8, 0.7, -0.1, 2.0, -0.9]$

\textbf{GELU}: applichiamo la GELU (\hyperref[sec:gelu]{Sezione~6.2}) elemento per elemento:
\begin{align*}
    \text{GELU}(1.2) &\approx 1.095, \quad \text{GELU}(-0.5) \approx -0.154, \quad \text{GELU}(0.3) \approx 0.185 \\
    \text{GELU}(-1.8) &\approx -0.038, \quad \text{GELU}(0.7) \approx 0.520, \quad \text{GELU}(-0.1) \approx -0.046 \\
    \text{GELU}(2.0) &\approx 1.955, \quad \text{GELU}(-0.9) \approx -0.121
\end{align*}
$\mathbf{h}_{\text{act}} = [1.095, -0.154, 0.185, -0.038, 0.520, -0.046, 1.955, -0.121]$

\textbf{Compressione}: $\mathbf{z}_{\text{ffn}} = \mathbf{h}_{\text{act}} \cdot \mathbf{W}_2 + \mathbf{b}_2 \in \mathbb{R}^{4}$

Supponiamo $\mathbf{z}_{\text{ffn}} = [0.3, -0.1, 0.5, 0.2]$.

\textbf{Osservazione}: l'MLP ha trasformato il vettore 4-dimensionale in uno spazio 8-dimensionale (pi\`u ``espressivo''), applicato una non-linearit\`a, e poi compresso di nuovo a 4 dimensioni. Questa sequenza expand-activate-compress \`e la ``ricetta'' standard per introdurre non-linearit\`a nel Transformer.
\end{examplebox}

\begin{notabox}[In parole semplici: l'MLP (expand-activate-compress)]
Detto in parole povere, l'MLP funziona come un ``laboratorio di trasformazione'' in tre fasi:
\begin{enumerate}
    \item \textbf{Espansione} ($1024 \to 4096$): il vettore viene ``srotolato'' in uno spazio pi\`u grande, come stendere un foglio accartocciato per vedere tutti i dettagli.
    \item \textbf{GELU} (non-linearit\`a): in questo spazio espanso, il modello pu\`o fare trasformazioni non lineari complesse --- combinare informazioni, sopprimere rumori, enfatizzare pattern.
    \item \textbf{Compressione} ($4096 \to 1024$): il vettore viene ``ricompresso'' alla dimensione originale, tenendo solo le trasformazioni utili.
\end{enumerate}
Senza l'MLP, il Perceiver sarebbe limitato a combinazioni lineari --- come cercare di disegnare curve usando solo il righello.
\end{notabox}

\subsubsection{Connessione residua dopo l'MLP}

\[
    \mathbf{z}_{\text{after\_ffn}} = \mathbf{z}_{\text{after\_ca}} + \mathbf{z}_{\text{ffn}}
\]

La connessione residua assicura che l'informazione del passo precedente (dopo la cross-attention) non venga persa. Anche se l'MLP produce un output poco utile, la residual connection preserva il segnale originale.

Continuando l'esempio numerico di sopra:
\[
    \mathbf{z}_{\text{after\_ffn}} = \underbrace{[0.5, -0.3, 0.8, -0.1]}_{\text{skip (originale)}} + \underbrace{[0.3, -0.1, 0.5, 0.2]}_{\text{MLP output}} = [0.8, -0.4, 1.3, 0.1]
\]
Il residuo \`e una ``correzione'' al vettore originale.


%% --- A7: Latent Transformer ---
\subsection{A7. Blocco Latent Transformer --- Dettaglio} \label{sec:latent_detail}

\begin{pipelinebox}
Input $\to$ Fourier $\to$ Embedding $\to$ Latent Array $\to$ Cross-Att $\to$ Res+MLP $\to$ \textbf{Latent Transformer} $\to$ Res+MLP $\to$ ($\times T$) $\to$ Pool $\to$ Classifier $\to$ Softmax $\to$ Loss
\end{pipelinebox}

Il Latent Transformer \`e una torre di $L = 6$ blocchi di self-attention che operano esclusivamente sui latenti (formule e schema del singolo blocco nella \hyperref[sec:latent_transformer]{Sezione~\ref*{sec:latent_transformer}}). La differenza cruciale con la cross-attention \`e che qui Q, K, V vengono \textbf{tutti dai latenti}: i latenti ``parlano tra di loro'' per integrare le informazioni estratte dall'input. La matrice di attenzione \`e quadrata $512 \times 512$, molto pi\`u piccola della $512 \times 50{,}176$ della cross-attention. La normalizzazione viene rieseguita perch\'e \textbf{prima di ogni blocco} si normalizza: il blocco di latent ``non si fida'' dell'output del cross-attention.

\subsubsection{Multi-Head Attention nel dettaglio} \label{sec:multi_head_detail}

Facciamo un passo indietro e chiediamoci: perch\'e usare 8 ``teste'' di attenzione e non una sola? La risposta \`e che una singola testa pu\`o catturare solo un tipo di relazione alla volta. La multi-head attention permette al modello di catturare contemporaneamente \textbf{relazioni diverse} tra i latenti, come se avesse 8 paia di occhiali con lenti diverse:

\begin{formulabox}[Multi-Head Attention --- meccanismo completo]
Con $H = 8$ teste e $D = 1024$, ogni testa opera su un sottospazio di $d_k = d_v = D / H = 128$ dimensioni:

\textbf{Per ogni testa $h = 1, \ldots, 8$}:
\begin{align*}
    \mathbf{Q}_h &= \mathbf{z}_{\text{norm}} \cdot \mathbf{W}_Q^{(h)} \in \mathbb{R}^{512 \times 128} \\
    \mathbf{K}_h &= \mathbf{z}_{\text{norm}} \cdot \mathbf{W}_K^{(h)} \in \mathbb{R}^{512 \times 128} \\
    \mathbf{V}_h &= \mathbf{z}_{\text{norm}} \cdot \mathbf{W}_V^{(h)} \in \mathbb{R}^{512 \times 128}
\end{align*}

\textbf{Attenzione per testa}: $\text{head}_h = \text{softmax}\!\left(\frac{\mathbf{Q}_h \mathbf{K}_h^\top}{\sqrt{128}}\right) \mathbf{V}_h \in \mathbb{R}^{512 \times 128}$

\textbf{Concatenazione}: $\text{MultiHead} = [\text{head}_1; \text{head}_2; \ldots; \text{head}_8] \in \mathbb{R}^{512 \times 1024}$

\textbf{Proiezione finale}: $\text{output} = \text{MultiHead} \cdot \mathbf{W}_O \in \mathbb{R}^{512 \times 1024}$

dove $\mathbf{W}_O \in \mathbb{R}^{1024 \times 1024}$ mescola le informazioni provenienti dalle diverse teste.
\end{formulabox}

\begin{notabox}[Perch\'e pi\`u teste sono meglio di una sola?]
Ogni testa pu\`o specializzarsi per catturare un tipo diverso di relazione tra i latenti:
\begin{itemize}
    \item \textbf{Testa 1}: potrebbe catturare relazioni tra latenti che codificano lo stesso colore.
    \item \textbf{Testa 2}: potrebbe catturare relazioni tra latenti che codificano regioni spaziali vicine.
    \item \textbf{Testa 3}: potrebbe catturare relazioni semantiche (es.\ latenti che codificano ``bordo'' e ``texture'').
    \item \textbf{Testa 4--8}: altre relazioni emergenti durante il training.
\end{itemize}
Con una singola testa di dimensione 1024, tutte queste relazioni dovrebbero essere catturate nello stesso spazio, limitando il potere espressivo. Con 8 teste da 128, ogni testa ha il suo ``spazio privato'' per catturare un tipo di relazione.
\end{notabox}

\subsubsection{Perch\'e $L = 6$ blocchi di self-attention?}

La domanda che ci si pone \`e: perch\'e non basta un solo blocco di self-attention? Con un solo blocco ($L = 1$), i latenti possono interagire solo a coppie. Con $L = 6$ blocchi, l'informazione pu\`o propagarsi attraverso catene di interazioni pi\`u complesse, come in una conversazione che si approfondisce turno dopo turno:
\begin{itemize}
    \item \textbf{Blocco 1}: il latente ``gatto'' interagisce con il latente ``zampe'' (relazione diretta).
    \item \textbf{Blocco 2}: il latente ``zampe'' (che ha gi\`a informazione dal ``gatto'') interagisce con il latente ``erba'' (relazione transitiva: gatto $\to$ zampe $\to$ erba).
    \item \textbf{Blocchi 3--6}: relazioni ancora pi\`u complesse e astratte si formano, permettendo rappresentazioni semanticamente ricche.
\end{itemize}
Il numero $L = 6$ \`e una scelta empirica: sufficiente per catturare relazioni complesse, ma non cos\`i grande da causare vanishing gradient (mitigato dalle residual connections).

\begin{examplebox}[Esempio numerico: self-attention nel Latent Transformer]
Consideriamo 3 latenti con $d = 4$ dimensioni (semplificato). A differenza della cross-attention, qui Q, K, V provengono \textbf{tutti dalla stessa sorgente} (i latenti).

\textbf{Input} (3 latenti dopo la cross-attention):
\[
\mathbf{z} = \begin{pmatrix} 0.5 & -0.3 & 0.8 & 0.1 \\ -0.2 & 0.7 & 0.1 & -0.5 \\ 0.3 & -0.1 & 0.4 & 0.6 \end{pmatrix}
\]

Supponiamo $\mathbf{W}_Q = \mathbf{W}_K = \mathbf{W}_V = \mathbf{I}$ (identit\`a, per semplicit\`a), cosicch\'e $\mathbf{Q} = \mathbf{K} = \mathbf{V} = \mathbf{z}$.

\textbf{Score $\mathbf{Q}\mathbf{K}^\top$ ($3 \times 3$, matrice QUADRATA)}:
\begin{align*}
\mathbf{S}_{11} &= \mathbf{z}_1 \cdot \mathbf{z}_1 = 0.25 + 0.09 + 0.64 + 0.01 = 0.99 \quad \text{(auto-attenzione)} \\
\mathbf{S}_{12} &= \mathbf{z}_1 \cdot \mathbf{z}_2 = -0.10 - 0.21 + 0.08 - 0.05 = -0.28 \\
\mathbf{S}_{13} &= \mathbf{z}_1 \cdot \mathbf{z}_3 = 0.15 + 0.03 + 0.32 + 0.06 = 0.56
\end{align*}

Matrice completa (scaling per $\sqrt{d_k} = \sqrt{4} = 2$):
\[
\mathbf{S}_{\text{scaled}} = \frac{1}{2}\begin{pmatrix} 0.99 & -0.28 & 0.56 \\ -0.28 & 0.79 & -0.05 \\ 0.56 & -0.05 & 0.62 \end{pmatrix} = \begin{pmatrix} 0.495 & -0.140 & 0.280 \\ -0.140 & 0.395 & -0.025 \\ 0.280 & -0.025 & 0.310 \end{pmatrix}
\]

\textbf{Softmax (riga per riga)}:

Riga 1: $e^{0.495} = 1.641$, $e^{-0.140} = 0.869$, $e^{0.280} = 1.323$. Somma $= 3.833$.
\[
\mathbf{A}_1 = [0.428,\; 0.227,\; 0.345]
\]

\textbf{Matrice di attenzione completa ($3 \times 3$, quadrata)}:
\[
\mathbf{A}_{\text{self}} = \begin{pmatrix} 0.428 & 0.227 & 0.345 \\ 0.227 & 0.430 & 0.343 \\ 0.345 & 0.243 & 0.412 \end{pmatrix}
\]

\textbf{Osservazione fondamentale}: ogni latente presta la massima attenzione \textbf{a s\'e stesso} (valori diagonali pi\`u alti: $0.428$, $0.430$, $0.412$). Questo \`e tipico della self-attention: i vettori sono pi\`u simili a s\'e stessi che agli altri.

\textbf{Confronto con la cross-attention}:
\begin{center}
\begin{tabular}{@{}lcc@{}}
\toprule
& Self-Attention & Cross-Attention \\
\midrule
Matrice $\mathbf{A}$ & $3 \times 3$ (quadrata) & $2 \times 4$ (rettangolare) \\
Q viene da & latenti & latenti \\
K, V vengono da & latenti (stessa sorgente!) & input (sorgente diversa!) \\
Scopo & ``parlare tra loro'' & ``leggere l'input'' \\
\bottomrule
\end{tabular}
\end{center}
\end{examplebox}


%% --- A8: Iterazione ---
\subsection{A8. Ripetizione Iterativa e Weight Sharing} \label{sec:ripetizione}

\begin{pipelinebox}
Input $\to$ Fourier $\to$ Embedding $\to$ Latent Array $\to$ Cross-Att $\to$ Res+MLP $\to$ Latent Transf.\ $\to$ Res+MLP $\to$ \textbf{(ripetizione $T$ volte)} $\to$ Pool $\to$ Classifier $\to$ Softmax $\to$ Loss
\end{pipelinebox}

L'intero ciclo (cross-attention + latent transformer) viene ripetuto $T = 8$ volte per ImageNet con weight sharing (si vedano le formule e lo schema nella \hyperref[sec:iterative]{Sezione~\ref*{sec:iterative}}). Ad ogni iterazione, i latenti ``rileggono'' l'input e raffinano la loro rappresentazione.

\begin{examplebox}[Cosa succede nelle 8 iterazioni --- intuizione]
\textbf{Iterazione 1} (pesi unici): i latenti ``leggono'' l'immagine per la prima volta. Ogni latente estrae informazione grezza: colori dominanti, posizioni grossolane, texture generali. La rappresentazione \`e ancora ``sfocata'' e poco discriminativa.

\textbf{Iterazione 2}: i latenti, che ora contengono una prima bozza dell'immagine, ``rileggono'' i pixel. Questa volta, grazie all'informazione gi\`a estratta, possono fare domande pi\`u specifiche: ``C'\`e un bordo nella regione dove ho trovato il rosso?''. La rappresentazione diventa pi\`u raffinata.

\textbf{Iterazioni 3--4}: i latenti catturano relazioni tra parti diverse dell'immagine. Un latente che codifica ``zampe'' e un latente che codifica ``orecchie'' iniziano a ``parlare'' durante il Latent Transformer, formando la rappresentazione ``gatto''.

\textbf{Iterazioni 5--6}: la rappresentazione diventa semanticamente coerente. I latenti si specializzano: alcuni codificano lo sfondo, altri l'oggetto principale, altri i dettagli.

\textbf{Iterazioni 7--8}: raffinamento finale. Le correzioni sono piccole ma importanti per la classificazione accurata. Il modello ``rilegge'' i dettagli sottili che servono a distinguere razze simili di cani, specie di uccelli, ecc.

Questo processo \`e simile a come un essere umano ``guarda'' un'immagine: prima una visione d'insieme (``\`e un animale''), poi i dettagli (``ha le orecchie appuntite''), poi il riconoscimento (``\`e un gatto siamese'').
\end{examplebox}

% TikZ: Iterative refinement visualization
\begin{figure}[H]
\centering
\begin{tikzpicture}[
    scale=0.85, transform shape
]

% Iteration 1
\node[font=\small\bfseries] at (0, 3.5) {Iterazione 1};
\fill[gray!30] (-1.5, -0.5) rectangle (1.5, 2.5);
\draw[thick] (-1.5, -0.5) rectangle (1.5, 2.5);
% Blurry content - simple shapes
\fill[red!20, opacity=0.5] (-0.5, 0.5) circle (0.8);
\fill[green!20, opacity=0.5] (0.5, 1.5) circle (0.6);
\node[font=\tiny, gray] at (0, -1.0) {Visione ``sfocata''};
\node[font=\tiny, gray] at (0, -1.5) {Colori dominanti};

% Iteration 2
\node[font=\small\bfseries] at (4.5, 3.5) {Iterazione 2};
\fill[gray!20] (3, -0.5) rectangle (6, 2.5);
\draw[thick] (3, -0.5) rectangle (6, 2.5);
% Slightly sharper
\fill[red!35] (3.8, 0.3) rectangle (5.2, 1.7);
\fill[green!30] (4.5, 1.5) circle (0.4);
\draw[thick, gray!50] (3.8, 0.3) -- (5.2, 0.3) -- (5.2, 1.7);
\node[font=\tiny, gray] at (4.5, -1.0) {Bordi grossolani};
\node[font=\tiny, gray] at (4.5, -1.5) {Regioni separate};

% Iteration 4
\node[font=\small\bfseries] at (9, 3.5) {Iterazione 4};
\fill[gray!10] (7.5, -0.5) rectangle (10.5, 2.5);
\draw[thick] (7.5, -0.5) rectangle (10.5, 2.5);
% Sharper shapes
\fill[red!50] (8.2, 0.2) rectangle (9.8, 1.3);
\fill[brown!40] (8.5, 1.3) -- (9.0, 2.0) -- (9.5, 1.3) -- cycle;
\draw[thick, gray!70] (8.2, 0.2) rectangle (9.8, 1.3);
\fill[green!40] (8.3, -0.2) rectangle (10.2, 0.2);
\node[font=\tiny, gray] at (9, -1.0) {Strutture riconoscibili};
\node[font=\tiny, gray] at (9, -1.5) {``Una casa?''};

% Iteration 8
\node[font=\small\bfseries] at (13.5, 3.5) {Iterazione 8};
\fill[white] (12, -0.5) rectangle (15, 2.5);
\draw[thick] (12, -0.5) rectangle (15, 2.5);
% Sharp content
\fill[red!60] (12.7, 0.2) rectangle (14.3, 1.3);
\fill[brown!50] (13.0, 1.3) -- (13.5, 2.2) -- (14.0, 1.3) -- cycle;
\draw[thick] (12.7, 0.2) rectangle (14.3, 1.3);
\draw[thick] (13.0, 1.3) -- (13.5, 2.2) -- (14.0, 1.3);
\fill[blue!30] (13.1, 0.6) rectangle (13.5, 1.0);
\fill[blue!30] (13.7, 0.6) rectangle (14.1, 1.0);
\fill[green!50] (12.3, -0.3) rectangle (14.7, 0.2);
\node[font=\tiny, gray] at (13.5, -1.0) {Dettagli fini};
\node[font=\tiny, gray] at (13.5, -1.5) {``Una casa rossa!''};

% Arrows between iterations
\draw[-{Stealth}, very thick, gray!60] (1.7, 1) -- (2.8, 1);
\draw[-{Stealth}, very thick, gray!60] (6.2, 1) -- (7.3, 1);
\draw[-{Stealth}, very thick, gray!60] (10.7, 1) -- (11.8, 1);

% Quality bar at bottom
\fill[red!30] (-1.5, -2.5) rectangle (1.5, -2.2);
\fill[orange!30] (3, -2.5) rectangle (6, -2.2);
\fill[yellow!40] (7.5, -2.5) rectangle (10.5, -2.2);
\fill[green!40] (12, -2.5) rectangle (15, -2.2);
\node[font=\tiny] at (0, -2.8) {Bassa qualit\`a};
\node[font=\tiny] at (13.5, -2.8) {Alta qualit\`a};
\draw[-{Stealth}, thin, gray] (-1.5, -3.1) -- (15, -3.1) node[right, font=\tiny] {Qualit\`a};

\end{tikzpicture}
\caption{Visualizzazione concettuale del raffinamento iterativo. Ad ogni iterazione, la rappresentazione latente diventa pi\`u precisa: dall'iterazione 1 (visione sfocata, solo colori dominanti) all'iterazione 8 (dettagli fini, riconoscimento completo). I pesi condivisi costringono ogni iterazione a migliorare incrementalmente la rappresentazione.}
\label{fig:iterative_refinement}
\end{figure}

\textbf{Formalizzazione}: lo schema iterativo formale e le formule del weight sharing ($\theta_{\text{CA}}^{(1)} \neq \theta_{\text{CA}}^{(t\geq 2)}$, $\theta_{\text{LT}}$ condivisi) sono nella \hyperref[sec:iterative]{Sezione~\ref*{sec:iterative}}. L'effetto chiave: il WS riduce i parametri da $\sim$326M a $\sim$45M e migliora le prestazioni da 72.9\% a 78.0\% su ImageNet, agendo come regolarizzazione.


%% --- A9: Pooling ---
\subsection{A9. Average Pooling} \label{sec:pooling}

\begin{pipelinebox}
Input $\to$ Fourier $\to$ Embedding $\to$ Latent Array $\to$ Cross-Att $\to$ Res+MLP $\to$ Latent Transf.\ $\to$ Res+MLP $\to$ ($\times T$) $\to$ \textbf{Pooling} $\to$ Classifier $\to$ Softmax $\to$ Loss
\end{pipelinebox}

Siamo quasi alla fine del forward pass! Dopo $T = 8$ iterazioni, l'array latente $\mathbf{z}_T \in \mathbb{R}^{512 \times 1024}$ contiene una rappresentazione ricca dell'intera immagine distribuita su 512 vettori latenti. Ma per classificare l'immagine, il modello deve produrre un \textbf{singolo vettore} che riassuma tutta questa informazione. La domanda che ci si pone \`e: come ``comprimiamo'' 512 vettori in uno solo?

\begin{formulabox}[Global Average Pooling sui latenti]
\[
    \mathbf{h} = \frac{1}{N} \sum_{i=1}^{N} \mathbf{z}_{T,i} = \frac{1}{512} \sum_{i=1}^{512} \mathbf{z}_{T,i} \in \mathbb{R}^{1024}
\]
Il pooling \`e la scelta pi\`u semplice: nessun parametro aggiuntivo, contributo uniforme di tutti i latenti. Ogni componente $j$ del vettore $\mathbf{h}$ \`e la media della $j$-esima componente dei 512 latenti:
\[
    h_j = \frac{1}{512} \sum_{i=1}^{512} [\mathbf{z}_{T}]_{i,j} \qquad \text{per } j = 1, \ldots, 1024
\]
\end{formulabox}

% TikZ: Average pooling visualization
\begin{figure}[H]
\centering
\begin{tikzpicture}[
    vec/.style={draw, minimum width=0.35cm, minimum height=1.2cm, font=\tiny},
    arrowstyle/.style={-{Stealth[length=2mm]}, thick},
    >=Stealth,
    scale=0.85, transform shape
]

% 512 latent vectors (show a few)
\node[font=\small\bfseries] at (3, 3.0) {512 latenti $\in \mathbb{R}^{1024}$};
\node[vec, fill=red!20] (v1) at (0, 1.5) {};
\node[vec, fill=blue!20] (v2) at (0.5, 1.5) {};
\node[vec, fill=green!20] (v3) at (1.0, 1.5) {};
\node[vec, fill=orange!20] (v4) at (1.5, 1.5) {};
\node[vec, fill=purple!20] (v5) at (2.0, 1.5) {};
\node[vec, fill=cyan!20] (v6) at (2.5, 1.5) {};
\node[font=\small] at (3.2, 1.5) {$\cdots$};
\node[vec, fill=yellow!20] (v7) at (3.8, 1.5) {};
\node[vec, fill=pink!20] (v8) at (4.3, 1.5) {};
\node[vec, fill=teal!20] (v9) at (4.8, 1.5) {};
\node[vec, fill=lime!20] (v10) at (5.3, 1.5) {};

% Arrow down
\draw[arrowstyle, very thick] (2.5, 0.6) -- (2.5, -0.2);
\node[font=\scriptsize] at (4.5, 0.2) {$\mathbf{h} = \frac{1}{512} \sum_i \mathbf{z}_i$};

% Single output vector
\node[font=\small\bfseries] at (3, -0.7) {1 vettore $\in \mathbb{R}^{1024}$};
\fill[green!30] (1.0, -1.8) rectangle (4.0, -1.2);
\draw[thick] (1.0, -1.8) rectangle (4.0, -1.2);
\node[font=\scriptsize] at (2.5, -1.5) {vettore medio $\mathbf{h}$};

% Annotation
\node[font=\scriptsize, gray, text width=5cm, align=center] at (2.5, -2.5) {Ogni componente $j$ del vettore $\mathbf{h}$ \`e\\la media della $j$-esima componente\\di tutti i 512 latenti};

\end{tikzpicture}
\caption{Global average pooling: i 512 vettori latenti (ciascuno di 1024 dimensioni) vengono mediati componente per componente, producendo un unico vettore di 1024 dimensioni che riassume l'intera rappresentazione dell'immagine.}
\label{fig:average_pooling}
\end{figure}

\begin{notabox}[In parole semplici: l'average pooling]
Detto in parole povere, il pooling \`e come chiedere a 512 persone che hanno osservato un'immagine di dare il loro ``riassunto'' e poi fare la \textbf{media} di tutti i riassunti. Se 300 persone dicono ``c'\`e un gatto'' e 212 dicono altre cose, la media sar\`a fortemente influenzata da ``gatto''. Il risultato \`e un unico vettore che cattura il ``consenso'' di tutti i latenti.
\end{notabox}

\begin{notabox}[Alternative al global average pooling]
\textbf{Perch\'e la media e non il max pooling?} Il max pooling seleziona il valore massimo, perdendo informazione sugli altri latenti. La media \`e pi\`u ``democratica'': ogni latente contribuisce ugualmente. Nella pratica, il global average pooling funziona bene per la classificazione.

\textbf{Perch\'e non un token CLS?} Il ViT (\hyperref[app:vit]{Appendice~I}) usa un token CLS speciale che accumula l'informazione globale. Il Perceiver potrebbe fare lo stesso (aggiungendo un latente ``speciale''), ma la media semplice funziona altrettanto bene e non richiede parametri aggiuntivi.

\textbf{Limite}: il pooling perde la struttura spaziale. Dopo la media, non sappiamo pi\`u ``quale latente codifica quale parte dell'immagine''. Questo \`e accettabile per la classificazione (serve solo la classe), ma limita l'uso per compiti come la segmentazione (risolto dal Perceiver IO con decoder cross-attention).
\end{notabox}


%% --- A10: Classificatore ---
\subsection{A10. Testa di Classificazione} \label{sec:classificatore}

\begin{pipelinebox}
Input $\to$ Fourier $\to$ Embedding $\to$ Latent Array $\to$ Cross-Att $\to$ Res+MLP $\to$ Latent Transf.\ $\to$ Res+MLP $\to$ ($\times T$) $\to$ Pool $\to$ \textbf{Classifier} $\to$ Softmax $\to$ Loss
\end{pipelinebox}

\begin{formulabox}[Classificatore lineare]
\[
    \text{logits} = \mathbf{h} \cdot \mathbf{W}_{\text{class}} + \mathbf{b}_{\text{class}} \in \mathbb{R}^{1000}
\]
dove $\mathbf{W}_{\text{class}} \in \mathbb{R}^{1024 \times 1000}$, $\mathbf{b}_{\text{class}} \in \mathbb{R}^{1000}$.

\textbf{Parametri}: $1024 \times 1000 + 1000 = 1{,}025{,}000 \approx 1$M.
\end{formulabox}

\textbf{Perch\'e un classificatore lineare e non un MLP?} La domanda \`e legittima: se abbiamo usato MLP non lineari in tutti i blocchi precedenti, perch\'e l'ultimo strato \`e lineare?

Il vettore $\mathbf{h} \in \mathbb{R}^{1024}$, dopo 8 iterazioni di cross-attention e 48 blocchi di self-attention, contiene gi\`a una rappresentazione altamente elaborata e non lineare dell'immagine. Un ulteriore MLP non lineare aggiungerebbe parametri e rischio di overfitting senza beneficio significativo. La scelta di un classificatore lineare \`e standard nei Transformer (ViT usa la stessa strategia) e riflette l'idea che la rete ``profonda'' ha gi\`a estratto feature sufficientemente discriminative.

\begin{notabox}[In parole semplici: il classificatore lineare]
Detto in parole povere, il classificatore \`e come l'ultimo esaminatore di una lunga catena di ispezioni. Dopo che 56 ispettori (i blocchi del Perceiver) hanno gi\`a analizzato, trasformato e raffinato la rappresentazione dell'immagine, l'ultimo esaminatore deve solo ``mettere un voto'' --- non ha bisogno di fare analisi complesse, perch\'e tutto il lavoro \`e gi\`a stato fatto. Un semplice prodotto matriciale $\mathbf{h} \cdot \mathbf{W}_{\text{class}}$ basta per mappare il vettore di 1024 dimensioni nelle 1000 classi di ImageNet.
\end{notabox}

\begin{notabox}[Inizializzazione del classificatore]
Il classificatore viene inizializzato con pesi piccoli (Xavier init) e bias a zero. All'inizio del training, i logit sono tutti vicini a zero, quindi la softmax produce una distribuzione quasi uniforme ($p_k \approx 1/1000 = 0.001$ per ogni classe). Questo significa che la loss iniziale \`e circa $-\log(0.001) = 6.9$, che \`e il valore massimo atteso. Man mano che il training procede, i logit divergono e la loss diminuisce.
\end{notabox}


%% --- A11: Softmax ---
\subsection{A11. Softmax e Predizione} \label{sec:softmax_forward}

\begin{pipelinebox}
Input $\to$ Fourier $\to$ Embedding $\to$ Latent Array $\to$ Cross-Att $\to$ Res+MLP $\to$ Latent Transf.\ $\to$ Res+MLP $\to$ ($\times T$) $\to$ Pool $\to$ Classifier $\to$ \textbf{Softmax} $\to$ Loss
\end{pipelinebox}

I 1000 logit prodotti dal classificatore vengono trasformati in probabilit\`a tramite la funzione softmax (si veda la \hyperref[sec:softmax]{Sezione~\ref*{sec:softmax}} per definizione, propriet\`a e esempio numerico):
\[
    p_k = \text{softmax}(\text{logit}_k) = \frac{\exp(\text{logit}_k)}{\sum_{k'=1}^{1000} \exp(\text{logit}_{k'})} \qquad \hat{y} = \arg\max_k p_k
\]
La classe predetta \`e quella con probabilit\`a massima.


%% --- A12: Loss ---
\subsection{A12. Funzione di Loss} \label{sec:loss}

\begin{pipelinebox}
Input $\to$ Fourier $\to$ Embedding $\to$ Latent Array $\to$ Cross-Att $\to$ Res+MLP $\to$ Latent Transf.\ $\to$ Res+MLP $\to$ ($\times T$) $\to$ Pool $\to$ Classifier $\to$ Softmax $\to$ \textbf{Loss}
\end{pipelinebox}

Per la classificazione su ImageNet si usa la \textbf{cross-entropy loss}:
\[
    \mathcal{L} = -\sum_{k=1}^{1000} y_{\text{true},k} \cdot \log(p_k) = -\log(p_c)
\]
dove $c$ \`e l'indice della classe corretta. La loss \`e bassa quando il modello assegna alta probabilit\`a alla classe giusta ($p_c \to 1 \implies \mathcal{L} \to 0$) e diverge quando il modello \`e sicuro ma sbagliato ($p_c \to 0 \implies \mathcal{L} \to +\infty$).

\medskip
\noindent Per la trattazione completa con esempi numerici passo-passo, gradienti, e le altre funzioni di loss usate dal Perceiver IO (MSE per ricostruzione, EPE per optical flow, loss combinata multimodale), si veda la \hyperref[sec:cross_entropy_loss]{\textbf{Sezione~\ref*{sec:cross_entropy_loss} --- Approfondimento sulle funzioni di Loss}}.


\newpage
%% ============================================================
%%  SEZIONE 6 --- FUNZIONI DI ATTIVAZIONE
%% ============================================================
\section{Funzioni di Attivazione: ReLU vs GELU vs GEGLU} \label{sec:attivazioni}

Senza non-linearit\`a, una rete neurale con 100 strati lineari sarebbe equivalente a un singolo strato lineare --- tutta quella profondit\`a sarebbe inutile.

Il blocco MLP del Perceiver (\hyperref[sec:mlp_forward]{Sezione~5.6}) utilizza la funzione di attivazione \textbf{GELU}. Questa sezione analizza in dettaglio le tre principali funzioni di attivazione usate nei Transformer moderni, spiegando perch\'e la GELU \`e stata scelta per il Perceiver.

\subsection{ReLU (Rectified Linear Unit)} \label{sec:relu}

Iniziamo dalla funzione di attivazione pi\`u semplice e storicamente pi\`u importante, per poi capire perch\'e il Perceiver ha scelto qualcosa di diverso.

\begin{formulabox}[ReLU: formula e derivata]
\[
    \text{ReLU}(x) = \max(0, x) = \begin{cases} x & \text{se } x > 0 \\ 0 & \text{se } x \leq 0 \end{cases}
\]
Derivata: $\text{ReLU}'(x) = \begin{cases} 1 & x > 0 \\ 0 & x < 0 \\ \text{indefinita} & x = 0 \end{cases}$

In pratica, si usa la convenzione $\text{ReLU}'(0) = 0$.
\end{formulabox}

\textbf{Vantaggi della ReLU}:
\begin{enumerate}
    \item \textbf{Semplicit\`a computazionale}: un solo confronto e una moltiplicazione.
    \item \textbf{Gradiente costante}: per $x > 0$, il gradiente \`e esattamente 1, evitando il vanishing gradient.
    \item \textbf{Sparsit\`a}: molte attivazioni sono esattamente 0, rendendo la rappresentazione sparsa.
\end{enumerate}

\textbf{Problema critico: i ``neuroni morti'' (dying ReLU)}: se un neurone riceve sempre input negativi (es.\ a causa di un learning rate troppo alto che spinge il bias in zona negativa), la sua uscita \`e sempre 0 e il suo gradiente \`e sempre 0. Il neurone non si aggiorner\`a mai pi\`u --- \`e ``morto'' permanentemente. Questo pu\`o accadere al 10--40\% dei neuroni in reti profonde.

\textbf{Derivata discontinua in $x = 0$}: la ReLU non \`e differenziabile in 0. Anche se in pratica questo non causa problemi gravi (la probabilit\`a che un input sia esattamente 0 \`e trascurabile), la discontinuit\`a del gradiente pu\`o rendere l'ottimizzazione meno stabile rispetto a funzioni lisce.

\subsection{GELU (Gaussian Error Linear Unit)} \label{sec:gelu}

Ora vediamo l'attivazione che il Perceiver usa effettivamente. Proposta da Hendrycks \& Gimpel (2016), la GELU \`e diventata l'attivazione standard nei Transformer moderni (GPT, BERT, Perceiver). La domanda che ci si pone \`e: perch\'e non usare semplicemente la ReLU, che \`e pi\`u semplice?

\begin{formulabox}[GELU: definizione esatta]
\textbf{Definizione esatta}: $\text{GELU}(x) = x \cdot \Phi(x)$

dove $\Phi(x) = \frac{1}{2}\left[1 + \text{erf}\!\left(\frac{x}{\sqrt{2}}\right)\right]$ \`e la CDF della distribuzione normale standard $\mathcal{N}(0,1)$.

\textbf{Approssimazione} (usata in pratica per efficienza):
\[
\text{GELU}(x) \approx 0.5\, x \left[1 + \tanh\!\left(\sqrt{\frac{2}{\pi}}\left(x + 0.044715\, x^3\right)\right)\right]
\]

\textbf{Derivata}: $\text{GELU}'(x) = \Phi(x) + x \cdot \phi(x)$

dove $\phi(x) = \frac{1}{\sqrt{2\pi}} e^{-x^2/2}$ \`e la PDF della normale standard.
\end{formulabox}

\textbf{Perch\'e la GELU \`e migliore della ReLU --- 4 ragioni fondamentali}:
\begin{enumerate}
    \item \textbf{Liscia e differenziabile ovunque}: nessuna discontinuit\`a nel gradiente. L'ottimizzatore pu\`o calcolare gradienti esatti in ogni punto, rendendo il training pi\`u stabile.
    \item \textbf{Nessun dying neuron}: la GELU permette piccoli valori negativi per input negativi. Per esempio, $\text{GELU}(-1) \approx -0.159 \neq 0$. Questo garantisce che il gradiente non sia mai esattamente zero, permettendo al neurone di ``riprendersi''.
    \item \textbf{Self-gating}: la formula $x \cdot \Phi(x)$ pu\`o essere interpretata come un meccanismo di \emph{gating}: il valore $x$ viene modulato dalla probabilit\`a $\Phi(x)$ che un campione dalla normale standard sia $\leq x$. Per $x$ molto positivo, $\Phi(x) \approx 1$ e l'output \`e $\approx x$; per $x$ molto negativo, $\Phi(x) \approx 0$ e l'output \`e $\approx 0$.
    \item \textbf{Interpretazione come dropout deterministico}: la GELU pu\`o essere vista come una versione deterministica del dropout stocastico: ogni neurone viene ``mantenuto'' con probabilit\`a proporzionale al suo valore. Valori grandi e positivi vengono sempre mantenuti; valori piccoli o negativi vengono parzialmente soppressi.
\end{enumerate}

\subsection{GEGLU (Gated GELU)} \label{sec:geglu}

Per completezza, menzioniamo anche la GEGLU, un'estensione della GELU che il Perceiver \emph{non} usa, ma che \`e importante conoscere perch\'e \`e usata in architetture pi\`u recenti.

\begin{formulabox}[GEGLU: formula e struttura]
\[
\text{GEGLU}(\mathbf{x}, \mathbf{W}_1, \mathbf{W}_2) = \text{GELU}(\mathbf{x}\, \mathbf{W}_1) \odot (\mathbf{x}\, \mathbf{W}_2)
\]
Due rami paralleli:
\begin{itemize}
    \item Ramo ``gate'': $\text{GELU}(\mathbf{x}\, \mathbf{W}_1)$ --- decide \emph{quanto} lasciare passare.
    \item Ramo ``contenuto'': $\mathbf{x}\, \mathbf{W}_2$ --- l'informazione effettiva.
    \item Il prodotto elemento per elemento $\odot$ combina i due rami.
\end{itemize}
Il GEGLU richiede \textbf{due} matrici di peso $\mathbf{W}_1, \mathbf{W}_2 \in \mathbb{R}^{D \times D_{\text{ff}}}$ nel primo strato dell'MLP, raddoppiando i parametri rispetto alla GELU.
\end{formulabox}

Il Perceiver usa GELU (non GEGLU) perch\'e GEGLU richiede il doppio dei parametri nel primo layer dell'MLP. Con il weight sharing, ogni parametro aggiuntivo \`e particolarmente costoso.

\begin{notabox}[In parole semplici: ReLU vs GELU]
Immaginiamo un interruttore della luce. La ReLU \`e un interruttore \textbf{on/off}: se il segnale \`e positivo, passa tutto; se \`e negativo, viene bloccato completamente (il neurone ``muore''). La GELU \`e pi\`u simile a un \textbf{dimmer} (regolatore di intensit\`a): i segnali positivi passano quasi completamente, i segnali negativi vengono fortemente attenuati ma non azzerati del tutto. Questo piccolo ``spiraglio'' per i valori negativi permette ai neuroni di ``riprendersi'' durante il training, evitando che muoiano permanentemente.
\end{notabox}

% TikZ: GELU vs ReLU vs LeakyReLU comparison plot
\begin{figure}[H]
\centering
\begin{tikzpicture}
\begin{axis}[
    width=13cm, height=8cm,
    xlabel={$x$},
    ylabel={$f(x)$},
    xmin=-4, xmax=4,
    ymin=-1.5, ymax=4,
    legend pos=north west,
    grid=major,
    grid style={gray!20},
    thick,
    samples=200,
    axis lines=middle,
    every axis x label/.style={at={(ticklabel* cs:1.0)}, anchor=west},
    every axis y label/.style={at={(ticklabel* cs:1.0)}, anchor=south}
]

% ReLU
\addplot[red, very thick, domain=-4:0] {0};
\addplot[red, very thick, domain=0:4] {x};
\addlegendentry{ReLU: $\max(0, x)$}

% LeakyReLU (alpha=0.1)
\addplot[orange, very thick, dashed, domain=-4:0] {0.1*x};
\addplot[orange, very thick, dashed, domain=0:4] {x};
\addlegendentry{LeakyReLU ($\alpha{=}0.1$)}

% GELU (approximation)
\addplot[blue, very thick, domain=-4:4] {0.5*x*(1 + tanh(deg(sqrt(2/pi)*(x + 0.044715*x^3))/57.2958))};
\addlegendentry{GELU: $x \cdot \Phi(x)$}

% Highlight the critical region
\fill[yellow!15, opacity=0.4] (axis cs:-2,-1.5) rectangle (axis cs:0,1);

% Annotation for dead zone
\draw[-{Stealth}, thick, red!70!black] (axis cs:-1.5, -1.0) -- (axis cs:-1.5, 0);
\node[font=\scriptsize, red!70!black, text width=2cm, align=center] at (axis cs:-1.5, -1.3) {ReLU: zona morta\\(gradiente = 0)};

% Annotation for GELU survival
\draw[-{Stealth}, thick, blue!70!black] (axis cs:-0.8, 0.8) -- (axis cs:-0.8, -0.12);
\node[font=\scriptsize, blue!70!black, text width=2cm, align=center] at (axis cs:-0.8, 1.1) {GELU: valore\\$\neq 0$};

% Region label
\node[font=\scriptsize, fill=yellow!20, draw=yellow!50, rounded corners] at (axis cs:-1, 0.5) {Regione critica};

\end{axis}
\end{tikzpicture}
\caption{Confronto tra ReLU, LeakyReLU e GELU. La regione evidenziata in giallo ($x \in [-2, 0]$) mostra dove le tre funzioni differiscono maggiormente. La ReLU (rosso) \`e esattamente zero per $x < 0$ (zona morta), la LeakyReLU (arancione) ha una pendenza ridotta $0.1x$, e la GELU (blu) ha una curva liscia che permette piccoli valori negativi. Il Perceiver usa la GELU per evitare neuroni morti.}
\label{fig:gelu_relu_leakyrelu_comparison}
\end{figure}

\subsection{Confronto numerico: ReLU vs GELU su 7 valori} \label{sec:confronto_attivazioni}

Per consolidare l'intuizione, vediamo un confronto numerico dettagliato. Questa tabella \`e molto utile per l'esame: mostra concretamente la differenza tra le due funzioni.

\begin{examplebox}[Confronto ReLU vs GELU per 7 valori di input]
\begin{center}
\begin{tabular}{@{}rrrrrr@{}}
\toprule
$x$ & ReLU($x$) & ReLU$'(x)$ & GELU($x$) & GELU$'(x)$ & Differenza \\
\midrule
$-2.0$ & $0.000$ & $0$ & $-0.045$ & $0.085$ & GELU $\neq 0$! \\
$-1.0$ & $0.000$ & $0$ & $-0.159$ & $0.133$ & GELU $\neq 0$! \\
$-0.5$ & $0.000$ & $0$ & $-0.154$ & $0.243$ & GELU $\neq 0$! \\
$0.0$ & $0.000$ & --- & $0.000$ & $0.500$ & Uguali \\
$0.5$ & $0.500$ & $1$ & $0.346$ & $0.757$ & GELU pi\`u cauto \\
$1.0$ & $1.000$ & $1$ & $0.841$ & $1.083$ & Quasi uguali \\
$2.0$ & $2.000$ & $1$ & $1.955$ & $1.085$ & Quasi uguali \\
\bottomrule
\end{tabular}
\end{center}

\textbf{Osservazione chiave}: per $x < 0$, la ReLU produce esattamente 0 (neurone ``morto''), mentre la GELU produce un piccolo valore negativo, mantenendo un gradiente non nullo. Per $x > 1$, le due funzioni convergono. La transizione \`e liscia per la GELU, a gradino per la ReLU.
\end{examplebox}

% TikZ/pgfplots: ReLU vs GELU comparison plot
\begin{figure}[H]
\centering
\begin{tikzpicture}
\begin{axis}[
    width=12cm, height=7cm,
    xlabel={$x$},
    ylabel={$f(x)$},
    legend pos=north west,
    grid=major,
    grid style={gray!30},
    domain=-3:3,
    samples=200,
    xmin=-3, xmax=3,
    ymin=-1, ymax=3,
    thick
]
% ReLU
\addplot[red, very thick] {max(0, x)};
\addlegendentry{ReLU: $\max(0, x)$}
% GELU approximation
\addplot[blue, very thick] {0.5*x*(1 + tanh(deg(sqrt(2/pi)*(x + 0.044715*x*x*x))/57.2958))};
\addlegendentry{GELU: $x \cdot \Phi(x)$}
% Identity reference
\addplot[gray, thin, dashed] {x};
\addlegendentry{Identit\`a: $y = x$}
% Zero line
\draw[gray, thin] (-3, 0) -- (3, 0);
\end{axis}
\end{tikzpicture}
\caption{Confronto tra ReLU (rosso) e GELU (blu). La ReLU ha una discontinuit\`a netta in $x=0$ e produce esattamente 0 per $x < 0$. La GELU \`e liscia, con una transizione graduale che permette piccoli valori negativi (visibili per $x \in [-2, 0]$). Per $x > 1$, le due funzioni convergono.}
\label{fig:relu_vs_gelu}
\end{figure}

\begin{figure}[H]
\centering
\begin{tikzpicture}
\begin{axis}[
    width=12cm, height=6cm,
    xlabel={$x$},
    ylabel={$f'(x)$},
    legend pos=south east,
    grid=major,
    grid style={gray!30},
    domain=-3:3,
    samples=200,
    xmin=-3, xmax=3,
    ymin=-0.3, ymax=1.5,
    thick
]
% ReLU derivative
\addplot[red, very thick] coordinates {(-3, 0) (-0.001, 0) (0.001, 1) (3, 1)};
\addlegendentry{ReLU$'(x)$: gradino}
% GELU derivative (approximation: Phi(x) + x * phi(x))
\addplot[blue, very thick] {0.5*(1 + tanh(deg(sqrt(2/pi)*(x + 0.044715*x*x*x))/57.2958)) + x * 0.3989 * exp(-x*x/2)};
\addlegendentry{GELU$'(x)$: liscia}
\end{axis}
\end{tikzpicture}
\caption{Confronto delle derivate. La derivata della ReLU \`e un gradino: 0 per $x<0$ (gradiente nullo = neurone morto!) e 1 per $x>0$. La derivata della GELU \`e continua, con un picco sopra 1 intorno a $x \approx 0.5$ e valori positivi anche per $x < 0$ (il gradiente fluisce sempre).}
\label{fig:relu_gelu_derivative}
\end{figure}

\begin{table}[htbp]
\centering
\caption{Confronto riassuntivo delle funzioni di attivazione.}
\begin{tabular}{@{}lccc@{}}
\toprule
\textbf{Propriet\`a} & \textbf{ReLU} & \textbf{GELU} & \textbf{GEGLU} \\
\midrule
Formula & $\max(0, x)$ & $x \cdot \Phi(x)$ & $\text{GELU}(x W_1) \odot (x W_2)$ \\
Liscia? & No ($x=0$) & S\`i & S\`i \\
Dying neuron? & S\`i & No & No \\
Parametri extra MLP & 0 & 0 & $\times 2$ nel primo layer \\
Usata in & CNN classiche & Perceiver, BERT, GPT & PaLM, T5-v1.1 \\
\bottomrule
\end{tabular}
\end{table}


\newpage
%% ============================================================
%%  SEZIONE 7 --- BACKWARD PASS
%% ============================================================
\section{Backward Pass: Derivazione dei Gradienti} \label{sec:backward}

Fino ad ora abbiamo visto come il Perceiver trasforma un'immagine in una predizione (forward pass). Ma come fa il modello a \emph{imparare}? La risposta \`e il backward pass --- il processo che calcola ``quanto e in che direzione'' aggiornare ciascuno dei 45 milioni di parametri.

Questa sezione descrive il flusso dei gradienti durante la backpropagation, seguendo il percorso inverso rispetto al \hyperref[sec:forward]{forward pass}. In pratica, quello che succede \`e che il segnale di errore ``risale'' dalla loss verso l'input, indicando a ogni parametro ``sei stato troppo grande o troppo piccolo''.

\textbf{Obiettivo del backward pass}: calcolare $\frac{\partial \mathcal{L}}{\partial \theta}$ per ogni parametro $\theta$ del modello (pesi delle proiezioni Q/K/V, pesi dell'MLP, parametri della LayerNorm, pesi del classificatore, array latente iniziale). Una volta calcolati tutti i gradienti, l'ottimizzatore LAMB (\hyperref[sec:ottimizzatori]{Sezione~8.2}) li usa per aggiornare i parametri.

\textbf{Principio della chain rule}: il backward pass si basa su un unico principio matematico: la regola della catena (chain rule) del calcolo differenziale. \`E lo strumento che permette di ``scomporre'' la derivata di una funzione composta in prodotti di derivate pi\`u semplici. Se $\mathcal{L} = f(g(h(\theta)))$, allora:
\[
    \frac{\partial \mathcal{L}}{\partial \theta} = \frac{\partial f}{\partial g} \cdot \frac{\partial g}{\partial h} \cdot \frac{\partial h}{\partial \theta}
\]
Ogni ``blocco'' del backward riceve un gradiente dall'alto ($\frac{\partial \mathcal{L}}{\partial \text{output}}$) e produce due output:
\begin{enumerate}
    \item Il \textbf{gradiente rispetto ai propri parametri} (per l'aggiornamento tramite l'ottimizzatore).
    \item Il \textbf{gradiente rispetto al proprio input} (per la propagazione al blocco precedente).
\end{enumerate}
Questo schema si ripete identico per ogni blocco, dal classificatore fino all'embedding dell'input.

\begin{notabox}[In parole semplici: il backward pass]
Immaginiamo di aver cucinato una torta che non \`e venuta buona (loss alta). Il backward pass \`e come ripercorrere tutti i passi della ricetta all'indietro: ``Forse ho messo troppo sale? Forse il forno era troppo caldo? Forse la farina era sbagliata?''. Per ogni ingrediente (parametro), il backward calcola ``di quanto cambierebbe il gusto (la loss) se cambiassi questo ingrediente di un pochino''. Questo ``di quanto'' \`e il gradiente, e l'ottimizzatore lo usa per aggiustare la ricetta al prossimo tentativo.
\end{notabox}

\textbf{Ordine del backward}: inverso rispetto al forward. Si parte dalla loss e si risale verso l'input. Per le $T$ iterazioni con weight sharing, il backward procede da $t = T$ a $t = 1$, accumulando i gradienti dei pesi condivisi.

\textbf{Nota sulla notazione}: usiamo la notazione $\delta_{\mathbf{x}} = \frac{\partial \mathcal{L}}{\partial \mathbf{x}}$ per indicare il gradiente della loss rispetto a una quantit\`a $\mathbf{x}$. Il simbolo $\delta$ indica sempre un gradiente che ``scende'' dall'alto (dalla loss) verso il basso (verso i parametri). Quando scriviamo $\delta_{\text{ffn}}$, intendiamo ``il gradiente della loss rispetto all'output dell'FFN''.

\textbf{Due tipi di gradienti}: per ogni blocco, il backward calcola due tipi di gradiente:
\begin{enumerate}
    \item \textbf{Gradiente rispetto ai parametri} ($\partial \mathcal{L} / \partial \mathbf{W}$): indica come aggiornare i pesi del blocco. Questo viene ``consumato'' dall'ottimizzatore e non si propaga ulteriormente.
    \item \textbf{Gradiente rispetto all'input} ($\partial \mathcal{L} / \partial \mathbf{x}$): indica l'errore che deve essere propagato al blocco precedente. Questo \`e la ``staffetta'' che passa al blocco successivo nel backward.
\end{enumerate}

\textbf{Parametri totali da aggiornare}: con weight sharing, il Perceiver ha $\sim$45M parametri. Ad ogni passo di training, il backward calcola il gradiente di \emph{tutti} i 45M parametri, e l'ottimizzatore LAMB li aggiorna. Il costo computazionale del backward \`e circa $2\times$ quello del forward (perch\'e deve calcolare sia i gradienti dei pesi sia i gradienti degli input per la propagazione).

Questo significa che un singolo passo di training (forward + backward + aggiornamento) richiede circa $3\times$ le operazioni del solo forward. Per 2.4 milioni di passi totali su ImageNet, il costo computazionale \`e enorme, ma gestibile su acceleratori moderni.

\textbf{Memoria necessaria}: durante il backward, il framework deve ``ricordare'' tutte le attivazioni intermedie calcolate nel forward (per calcolare i gradienti). Per il Perceiver su ImageNet, questo richiede circa 2--4 GB di memoria GPU per un singolo campione (a seconda del batch size e dell'implementazione). La matrice di cross-attention $\mathbb{R}^{512 \times 50{,}176}$ \`e la struttura pi\`u grande da memorizzare.

\begin{table}[htbp]
\centering
\caption{Strutture da memorizzare durante il forward per il backward pass.}
\small
\begin{tabular}{@{}p{3.2cm}p{3cm}rp{4.5cm}@{}}
\toprule
\textbf{Struttura} & \textbf{Dimensione} & \textbf{Mem.} & \textbf{Note} \\
\midrule
Input $\tilde{\mathbf{x}}$ & $50\text{K} \times 261$ & 52 MB & Per tutte le $T$ iterazioni \\
Attenzione CA & $512 \times 50\text{K}$ & 103 MB & Una per iterazione \\
Q, K, V (CA) & varie & 200 MB & Per gradiente proiezioni \\
Latenti interm. & $512 \times 1024$ & 2.1 MB & Per ogni blocco residual \\
Attenzione SA & $512 \times 512$ & 1.0 MB & Per testa, per blocco \\
Attivazioni MLP & $512 \times 4096$ & 8.4 MB & Pre-GELU (backward) \\
\bottomrule
\end{tabular}
\end{table}

La componente dominante \`e la matrice di cross-attention ($103$ MB per iterazione, $\times 8 = 824$ MB totali) e le proiezioni K, V dell'input ($\sim$200 MB per iterazione). Con checkpoint di attivazione (gradient checkpointing), si pu\`o ridurre la memoria ricalcolando alcune attivazioni durante il backward, al costo di un $\sim$30\% in pi\`u di tempo di computazione.

\begin{notabox}[In parole semplici: la memoria del backward pass]
Detto in parole povere, il backward pass ha bisogno di ``ricordare'' tutto ci\`o che \`e successo durante il forward pass. \`E come dover risolvere un'equazione a ritroso: per calcolare ``se cambio questo peso, come cambia la loss?'', devo sapere quali erano i valori intermedi calcolati durante il forward. Questo ``ricordo'' occupa memoria GPU --- la parte pi\`u ingombrante \`e la matrice di cross-attention ($512 \times 50.176$), che pesa $\sim$100 MB per ciascuna delle 8 iterazioni.
\end{notabox}

\textbf{Pipeline del backward} (direzione inversa):
\begin{center}
\small
\fbox{\parbox{0.95\textwidth}{\centering
$\text{Loss} \leftarrow \text{Softmax} \leftarrow \text{Classifier} \leftarrow \text{Pool} \leftarrow (\times T)$\\[2pt]
$\leftarrow \text{Res{+}MLP} \leftarrow \text{Latent Transf.} \leftarrow \text{Res{+}MLP} \leftarrow \text{Cross-Att}$
}}
\end{center}

% TikZ: Complete backward pass pipeline
\begin{figure}[H]
\centering
\begin{tikzpicture}[
    step/.style={draw, rounded corners=3pt, fill=#1, minimum width=5.5cm, minimum height=0.6cm, font=\small},
    step/.default={gray!10},
    arrowstyle/.style={-{Stealth[length=2.5mm]}, thick, red!60!black},
    dim/.style={font=\scriptsize, gray},
    >=Stealth,
    scale=0.88, transform shape
]

\node[step=red!20] (s1) at (0, 0) {B1. $\partial\mathcal{L}/\partial\text{logits}$: Loss + Softmax};
\node[dim, right=0.3cm of s1] {$\delta \in \mathbb{R}^{1000}$};

\node[step=blue!20] (s2) at (0, 1.2) {B2. $\partial\mathcal{L}/\partial\mathbf{W}_{\text{class}}$, $\partial\mathcal{L}/\partial\mathbf{h}$};
\node[dim, right=0.3cm of s2] {$\delta_h \in \mathbb{R}^{1024}$};

\node[step=green!20] (s3) at (0, 2.4) {B3. $\partial\mathcal{L}/\partial\mathbf{z}_{T,i}$: Pooling};
\node[dim, right=0.3cm of s3] {$\delta_{z} \in \mathbb{R}^{512 \times 1024}$};

\node[step=cyan!15] (s4) at (0, 3.6) {B4. Latent Transformer backward};
\node[dim, right=0.3cm of s4] {Residual + MLP + LN};

\node[step=orange!15] (s5) at (0, 4.8) {B5. Self-Attention backward};
\node[dim, right=0.3cm of s5] {$\partial A$, softmax Jac., $\partial Q/K/V$};

\node[step=orange!20] (s6) at (0, 6.0) {B6. Cross-Attention backward};
\node[dim, right=0.3cm of s6] {$\partial\mathbf{z}|_Q$, $\partial\mathbf{x}|_{K,V}$};

\node[step=yellow!20] (s7) at (0, 7.2) {B7. Embedding backward};
\node[dim, right=0.3cm of s7] {$\partial\mathbf{W}_{\text{embed}}$};

\node[step=purple!15] (s8) at (0, 8.4) {B8. Weight sharing: accumulo gradienti};
\node[dim, right=0.3cm of s8] {$\sum_{t=2}^{T} \partial\theta^{(t)}$};

% Arrows (upward = backward direction)
\draw[arrowstyle] (s1) -- (s2);
\draw[arrowstyle] (s2) -- (s3);
\draw[arrowstyle] (s3) -- (s4);
\draw[arrowstyle] (s4) -- (s5);
\draw[arrowstyle] (s5) -- (s6);
\draw[arrowstyle] (s6) -- (s7);
\draw[arrowstyle] (s7) -- (s8);

% Iteration loop
\draw[arrowstyle, dashed, blue!60!black] (3.8, 3.6) -- (3.8, 6.0) -- (s6.east);
\node[font=\scriptsize, blue!60!black, rotate=90] at (4.1, 4.8) {$t = T, T{-}1, \ldots, 1$};

\end{tikzpicture}
\caption{Pipeline del backward pass completo. I gradienti fluiscono dalla loss verso l'input, in ordine inverso rispetto al forward pass. Le frecce rosse indicano la direzione del flusso dei gradienti. Il ciclo tratteggiato blu indica la ripetizione all'indietro per le $T$ iterazioni.}
\label{fig:backward_pipeline}
\end{figure}


\subsection{B1. Gradiente della Loss} \label{sec:grad_loss}

Il backward pass inizia dalla fine: dalla loss. Il primo gradiente da calcolare \`e forse il pi\`u elegante di tutto il deep learning. Vediamo perch\'e.

\begin{pipelinebox}[title={\small Backward --- Passo corrente evidenziato}]
\textbf{Loss} $\leftarrow$ Softmax $\leftarrow$ Classifier $\leftarrow$ Pool $\leftarrow$ ($\times T$) $\leftarrow$ Res+MLP $\leftarrow$ Latent Transf.\ $\leftarrow$ Res+MLP $\leftarrow$ Cross-Att
\end{pipelinebox}

\begin{formulabox}[Gradiente combinato softmax + cross-entropy]
Il gradiente della cross-entropy loss rispetto ai logit ha una forma elegante:
\[
    \frac{\partial \mathcal{L}}{\partial \text{logit}_k} = p_k - y_{\text{true},k}
\]
\end{formulabox}

\begin{proofbox}[Dimostrazione completa: gradiente combinato softmax + cross-entropy]
\textbf{Obiettivo}: dimostrare che $\frac{\partial \mathcal{L}}{\partial \text{logit}_k} = p_k - y_k$ dove $\mathcal{L} = -\sum_k y_k \log p_k$ e $p_k = \text{softmax}(\text{logit}_k)$.

\textbf{Passo 1}: ricordiamo la derivata della softmax (\hyperref[sec:softmax]{Sezione~2.1}):
\[
    \frac{\partial p_i}{\partial \text{logit}_j} = p_i(\delta_{ij} - p_j)
\]

\textbf{Passo 2}: applichiamo la chain rule alla cross-entropy:
\begin{align*}
    \frac{\partial \mathcal{L}}{\partial \text{logit}_j} &= -\sum_{i=1}^{K} y_i \cdot \frac{1}{p_i} \cdot \frac{\partial p_i}{\partial \text{logit}_j} \\
    &= -\sum_{i=1}^{K} y_i \cdot \frac{1}{p_i} \cdot p_i(\delta_{ij} - p_j) \\
    &= -\sum_{i=1}^{K} y_i \cdot (\delta_{ij} - p_j) \\
    &= -\sum_{i=1}^{K} y_i \delta_{ij} + p_j \sum_{i=1}^{K} y_i
\end{align*}

\textbf{Passo 3}: poich\'e $y$ \`e one-hot, $\sum_i y_i = 1$ e $\sum_i y_i \delta_{ij} = y_j$:
\[
    \frac{\partial \mathcal{L}}{\partial \text{logit}_j} = -y_j + p_j = p_j - y_j \qquad \blacksquare
\]
\end{proofbox}

\begin{notabox}[In parole semplici: il gradiente della loss]
Detto in parole povere, il gradiente $p_k - y_k$ ha un significato intuitivo bellissimo. Per la classe corretta (dove $y_k = 1$): il gradiente \`e $p_k - 1$, che \`e sempre negativo --- il modello deve \emph{aumentare} la probabilit\`a di questa classe. Pi\`u $p_k$ \`e basso (predizione sbagliata), pi\`u il gradiente \`e grande (correzione forte). Per le classi sbagliate (dove $y_k = 0$): il gradiente \`e $p_k - 0 = p_k$, che \`e sempre positivo --- il modello deve \emph{diminuire} la probabilit\`a di queste classi.
\end{notabox}

\textbf{Perch\'e ci serve questo gradiente?} La backpropagation parte dalla loss e procede all'indietro. Il primo gradiente da calcolare \`e ``quanto cambia la loss se cambio ciascun logit?''. La risposta \`e elegantemente semplice: $p_k - y_{\text{true},k}$. Se il modello assegna $p_c = 0.9$ alla classe corretta $c$, il gradiente per quella classe \`e $0.9 - 1 = -0.1$ (piccolo, il modello \`e gi\`a buono). Se $p_c = 0.01$, il gradiente \`e $0.01 - 1 = -0.99$ (grande, il modello deve correggersi molto).

\begin{examplebox}[Esempio numerico: gradiente della loss per 3 classi]
Supponiamo 3 classi, classe corretta = 1, logits $= [2.0, 1.0, 0.1]$.

\textbf{Softmax}: $\mathbf{p} = [0.659, 0.242, 0.099]$ (calcolato nella \hyperref[sec:softmax]{Sezione~2.1}).

\textbf{Label one-hot}: $\mathbf{y} = [1, 0, 0]$.

\textbf{Gradiente}: $\delta = \mathbf{p} - \mathbf{y} = [0.659 - 1,\; 0.242 - 0,\; 0.099 - 0] = [-0.341,\; 0.242,\; 0.099]$.

\textbf{Interpretazione}: la classe corretta (classe 1) ha gradiente negativo ($-0.341$): il modello deve \emph{aumentare} il logit corrispondente. Le classi errate hanno gradiente positivo: il modello deve \emph{diminuire} i loro logit.
\end{examplebox}


\subsection{B2. Gradiente della Testa di Classificazione} \label{sec:grad_class}

Ora che abbiamo il gradiente rispetto ai logit, dobbiamo propagarlo all'indietro attraverso il classificatore lineare. Questo passo produce tre gradienti: uno per i pesi, uno per il bias, e uno per l'input (che verr\`a propagato ulteriormente al pooling).

\begin{pipelinebox}[title={\small Backward --- Passo corrente evidenziato}]
Loss $\leftarrow$ Softmax $\leftarrow$ \textbf{Classifier} $\leftarrow$ Pool $\leftarrow$ ($\times T$) $\leftarrow$ Res+MLP $\leftarrow$ Latent Transf.\ $\leftarrow$ Res+MLP $\leftarrow$ Cross-Att
\end{pipelinebox}

\begin{formulabox}[Gradienti del classificatore lineare]
Con $\delta_{\text{logits}} = \mathbf{p} - \mathbf{y}_{\text{true}} \in \mathbb{R}^{1000}$:
\begin{align*}
    \frac{\partial \mathcal{L}}{\partial \mathbf{W}_{\text{class}}} &= \mathbf{h}^\top \cdot \delta_{\text{logits}} \in \mathbb{R}^{1024 \times 1000} \\
    \frac{\partial \mathcal{L}}{\partial \mathbf{b}_{\text{class}}} &= \delta_{\text{logits}} \in \mathbb{R}^{1000} \\
    \frac{\partial \mathcal{L}}{\partial \mathbf{h}} &= \delta_{\text{logits}} \cdot \mathbf{W}_{\text{class}}^\top \in \mathbb{R}^{1024}
\end{align*}
\end{formulabox}

\textbf{Perch\'e ci servono questi tre gradienti?} Questa \`e una domanda fondamentale per capire il backward pass:
\begin{itemize}
    \item Il primo ($\partial \mathcal{L} / \partial \mathbf{W}_{\text{class}}$) ci dice come aggiornare i pesi del classificatore. Viene ``consumato'' dall'ottimizzatore LAMB.
    \item Il secondo ($\partial \mathcal{L} / \partial \mathbf{b}_{\text{class}}$) ci dice come aggiornare il bias del classificatore. Anche questo viene consumato dall'ottimizzatore.
    \item Il terzo ($\partial \mathcal{L} / \partial \mathbf{h}$) \`e il segnale di errore che deve propagarsi ``all'indietro'' verso il pooling e poi verso i latenti. Questo \`e il gradiente ``di transito'' che continua il viaggio del backward pass.
\end{itemize}

\begin{proofbox}[Derivazione dei gradienti del classificatore]
Dato $\text{logits} = \mathbf{h} \cdot \mathbf{W}_{\text{class}} + \mathbf{b}_{\text{class}}$ e $\delta_{\text{logits}} = \mathbf{p} - \mathbf{y} \in \mathbb{R}^{1000}$:

\textbf{Gradiente rispetto a $\mathbf{W}_{\text{class}}$}:
\[
    \text{logit}_k = \sum_{j=1}^{1024} h_j \cdot W_{j,k} + b_k \implies \frac{\partial \text{logit}_k}{\partial W_{j,k}} = h_j
\]
Quindi: $\frac{\partial \mathcal{L}}{\partial W_{j,k}} = h_j \cdot \delta_{\text{logits},k}$, in forma matriciale: $\frac{\partial \mathcal{L}}{\partial \mathbf{W}_{\text{class}}} = \mathbf{h}^\top \cdot \delta_{\text{logits}} \in \mathbb{R}^{1024 \times 1000}$.

\textbf{Gradiente rispetto a $\mathbf{b}_{\text{class}}$}: $\frac{\partial \text{logit}_k}{\partial b_k} = 1$, quindi $\frac{\partial \mathcal{L}}{\partial \mathbf{b}_{\text{class}}} = \delta_{\text{logits}} \in \mathbb{R}^{1000}$.

\textbf{Gradiente rispetto a $\mathbf{h}$} (per la propagazione ai layer precedenti):
\[
    \frac{\partial \text{logit}_k}{\partial h_j} = W_{j,k} \implies \frac{\partial \mathcal{L}}{\partial h_j} = \sum_{k=1}^{1000} \delta_{\text{logits},k} \cdot W_{j,k}
\]
In forma matriciale: $\frac{\partial \mathcal{L}}{\partial \mathbf{h}} = \delta_{\text{logits}} \cdot \mathbf{W}_{\text{class}}^\top \in \mathbb{R}^{1024}$.
\end{proofbox}

\begin{examplebox}[Esempio numerico: backward del classificatore]
Supponiamo (semplificato) $D = 3$, 4 classi, classe corretta = 2.

$\mathbf{h} = [0.5, -0.3, 0.8]$ (output del pooling).

$\mathbf{W}_{\text{class}} = \begin{pmatrix} 0.2 & -0.1 & 0.4 & 0.1 \\ -0.3 & 0.5 & 0.2 & -0.2 \\ 0.1 & 0.3 & -0.1 & 0.4 \end{pmatrix}$ ($3 \times 4$).

\textbf{Forward}: $\text{logits} = [0.5 \times 0.2 + (-0.3) \times (-0.3) + 0.8 \times 0.1, \ldots] = [0.27, 0.09, 0.06, 0.43]$.

\textbf{Softmax}: $\mathbf{p} \approx [0.266, 0.222, 0.215, 0.297]$.

\textbf{Gradiente}: $\delta = \mathbf{p} - \mathbf{y} = [0.266, 0.222 - 1, 0.215, 0.297] = [0.266, -0.778, 0.215, 0.297]$.

\textbf{Backward a $\mathbf{h}$}: $\frac{\partial \mathcal{L}}{\partial \mathbf{h}} = \delta \cdot \mathbf{W}_{\text{class}}^\top \in \mathbb{R}^{3}$ --- questo vettore indica ``come aggiornare il pooling output''.
\end{examplebox}


\subsection{B3. Gradiente del Pooling} \label{sec:grad_pool}

Il gradiente attraverso il pooling \`e concettualmente semplice ma importante: deve ``distribuire'' il segnale di errore dal singolo vettore $\mathbf{h}$ a tutti i 512 latenti.

\begin{pipelinebox}[title={\small Backward --- Passo corrente evidenziato}]
Loss $\leftarrow$ Softmax $\leftarrow$ Classifier $\leftarrow$ \textbf{Pooling} $\leftarrow$ ($\times T$) $\leftarrow$ Res+MLP $\leftarrow$ Latent Transf.\ $\leftarrow$ Res+MLP $\leftarrow$ Cross-Att
\end{pipelinebox}

\begin{formulabox}[Gradiente dell'average pooling]
\[
    \frac{\partial \mathcal{L}}{\partial \mathbf{z}_{T,i}} = \frac{1}{N} \cdot \frac{\partial \mathcal{L}}{\partial \mathbf{h}} \qquad \forall\, i = 1, \ldots, N
\]
Il gradiente \`e uniforme: ogni latente riceve la stessa porzione ($1/N = 1/512$).
\end{formulabox}

\textbf{Perch\'e ci serve il backward del pooling?} Il backward serve a propagare l'errore non solo a $\mathbf{h}$, ma a \textbf{ciascun latente}. Se ci fermassimo a $\mathbf{h}$, sapremmo solo quanto l'errore dipende dal vettore medio, ma non sapremmo come aggiornare i \textbf{singoli latenti} $\mathbf{z}_{T,1}, \ldots, \mathbf{z}_{T,512}$.

Lo Jacobiano del pooling \`e semplice: poich\'e $\mathbf{h} = \frac{1}{N}\sum_i \mathbf{z}_{T,i}$, la derivata di $\mathbf{h}$ rispetto a qualsiasi $\mathbf{z}_{T,i}$ \`e $\frac{1}{N} \mathbf{I}$ (matrice identit\`a scalata). Questo significa che il gradiente si ``distribuisce uniformemente'' a tutti i latenti: ogni latente contribuisce in egual misura alla loss.

\begin{proofbox}[Derivazione dello Jacobiano del pooling]
Il pooling \`e $h_j = \frac{1}{N} \sum_{i=1}^{N} z_{i,j}$ per ogni componente $j$.

Lo Jacobiano rispetto al latente $i$-esimo:
\[
    \frac{\partial h_j}{\partial z_{i,k}} = \frac{1}{N} \cdot \delta_{jk}
\]
dove $\delta_{jk}$ \`e il delta di Kronecker. In forma matriciale:
\[
    \frac{\partial \mathbf{h}}{\partial \mathbf{z}_{T,i}} = \frac{1}{N} \mathbf{I}_{D \times D}
\]
Applicando la chain rule:
\[
    \frac{\partial \mathcal{L}}{\partial \mathbf{z}_{T,i}} = \frac{\partial \mathcal{L}}{\partial \mathbf{h}} \cdot \frac{\partial \mathbf{h}}{\partial \mathbf{z}_{T,i}} = \frac{1}{N} \cdot \frac{\partial \mathcal{L}}{\partial \mathbf{h}}
\]
Il gradiente \`e lo stesso per tutti i latenti $i = 1, \ldots, N$: $\frac{1}{512} \cdot \delta_{\mathbf{h}}$.
\end{proofbox}

\begin{examplebox}[Esempio numerico: backward del pooling con 3 latenti in $\mathbb{R}^2$]
Latenti: $\mathbf{z}_1 = [2, 4]$, $\mathbf{z}_2 = [6, 0]$, $\mathbf{z}_3 = [1, 5]$.

\textbf{Forward}: $\mathbf{h} = \frac{1}{3}([2,4] + [6,0] + [1,5]) = \frac{1}{3}[9, 9] = [3, 3]$.

\textbf{Backward}: supponiamo $\frac{\partial \mathcal{L}}{\partial \mathbf{h}} = [-0.6, 0.3]$.

Allora \textbf{ogni} latente riceve lo stesso gradiente:
\[
    \frac{\partial \mathcal{L}}{\partial \mathbf{z}_1} = \frac{\partial \mathcal{L}}{\partial \mathbf{z}_2} = \frac{\partial \mathcal{L}}{\partial \mathbf{z}_3} = \frac{1}{3} \cdot [-0.6, 0.3] = [-0.2, 0.1]
\]
Tutti e tre i latenti verranno aggiornati nella stessa direzione, ma le loro differenziazioni verranno mantenute dai gradienti che fluiscono attraverso le iterazioni precedenti (cross-attention e self-attention), dove ogni latente ha interazioni diverse.
\end{examplebox}


\subsection{B4. Gradiente del Latent Transformer} \label{sec:grad_latent}

Ora il gradiente deve attraversare il Latent Transformer. Questo \`e il punto pi\`u delicato del backward, perch\'e ci sono 6 blocchi di self-attention, ciascuno con residual connections, LayerNorm e MLP. Procediamo un passo alla volta.

\begin{pipelinebox}[title={\small Backward --- Passo corrente evidenziato}]
Loss $\leftarrow$ Softmax $\leftarrow$ Classifier $\leftarrow$ Pool $\leftarrow$ ($\times T$) $\leftarrow$ Res+MLP $\leftarrow$ \textbf{Latent Transf.} $\leftarrow$ Res+MLP $\leftarrow$ Cross-Att
\end{pipelinebox}

\subsubsection{Residual post-MLP}

Per la connessione residua $\mathbf{z}'' = \mathbf{z}' + \text{FFN}(\text{LN}(\mathbf{z}'))$, il gradiente si \textbf{copia uguale nei due rami} perch\'e la residual connection \`e una semplice somma:
\[
    \frac{\partial \mathcal{L}}{\partial \mathbf{z}'}\bigg|_{\text{skip}} = \frac{\partial \mathcal{L}}{\partial \mathbf{z}''}, \qquad \frac{\partial \mathcal{L}}{\partial \mathbf{z}'}\bigg|_{\text{MLP}} = \frac{\partial \mathcal{L}}{\partial \mathbf{z}''} \cdot \frac{\partial \text{FFN}}{\partial \mathbf{z}'}
\]
Il gradiente totale \`e la somma dei due rami:
\[
    \frac{\partial \mathcal{L}}{\partial \mathbf{z}'} = \underbrace{\frac{\partial \mathcal{L}}{\partial \mathbf{z}'}\bigg|_{\text{skip}}}_{\text{copia diretta del gradiente}} + \underbrace{\frac{\partial \mathcal{L}}{\partial \mathbf{z}'}\bigg|_{\text{MLP}}}_{\text{gradiente trasformato dal MLP}}
\]
Questa \`e la propriet\`a fondamentale delle residual connections: il gradiente non si attenua mai completamente, perch\'e una copia passa sempre ``direttamente'' attraverso lo skip.

\begin{notabox}[In parole semplici: il backward della residual connection]
Detto in parole povere: quando il gradiente arriva a una residual connection, si ``duplica'' in due copie. Una copia passa direttamente al blocco precedente (come una ``autostrada'' per il gradiente), l'altra attraversa il blocco (MLP o attenzione) e viene trasformata. Le due copie vengono poi sommate. Anche se la copia che attraversa il blocco si attenua molto (gradiente piccolo), la copia diretta arriva sempre ``fresca'' e intatta. \`E per questo che il Perceiver pu\`o avere 56 blocchi senza vanishing gradient.
\end{notabox}

\begin{notabox}[Perch\'e la somma dei due rami funziona]
Consideriamo la funzione $f(\mathbf{x}) = \mathbf{x} + g(\mathbf{x})$ (residual connection). Per la chain rule:
\[
    \frac{\partial \mathcal{L}}{\partial \mathbf{x}} = \frac{\partial \mathcal{L}}{\partial f} \cdot \frac{\partial f}{\partial \mathbf{x}} = \frac{\partial \mathcal{L}}{\partial f} \cdot \left(\mathbf{I} + \frac{\partial g}{\partial \mathbf{x}}\right) = \frac{\partial \mathcal{L}}{\partial f} + \frac{\partial \mathcal{L}}{\partial f} \cdot \frac{\partial g}{\partial \mathbf{x}}
\]
Il primo termine \`e il gradiente che ``salta'' il blocco; il secondo \`e il gradiente che ``attraversa'' il blocco. Anche se $\frac{\partial g}{\partial \mathbf{x}} \approx 0$ (blocco poco informativo), il primo termine garantisce che il gradiente fluisca. Questo \`e il motivo per cui il Perceiver pu\`o avere 56+ blocchi senza vanishing gradient.
\end{notabox}

\subsubsection{MLP backward --- passo per passo} \label{sec:mlp_backward}

Vediamo ora il backward dell'MLP, che \`e pi\`u meccanico ma molto istruttivo. Il principio \`e semplice: si procede in ordine inverso rispetto al forward. Se il forward fa espansione $\to$ GELU $\to$ compressione, il backward fa compressione$^{-1}$ $\to$ GELU$'$ $\to$ espansione$^{-1}$.

Il backward dell'MLP procede in ordine inverso rispetto al forward: prima la compressione, poi l'attivazione, poi l'espansione. Per ogni passo, calcoliamo sia il gradiente rispetto all'\textbf{input} (per la propagazione all'indietro) sia il gradiente rispetto ai \textbf{pesi} (per l'aggiornamento dei parametri).

\begin{formulabox}[Gradienti del blocco MLP --- dettaglio completo]
Sia $\delta_{\text{ffn}} = \frac{\partial \mathcal{L}}{\partial \mathbf{z}_{\text{ffn}}} \in \mathbb{R}^{512 \times 1024}$ il gradiente in ingresso (dal residual).

\textbf{1. Compressione backward} (secondo strato lineare: $\mathbf{z}_{\text{ffn}} = \mathbf{h}_{\text{act}} \cdot \mathbf{W}_2 + \mathbf{b}_2$):
\begin{align*}
    \delta_{\text{h\_act}} &= \delta_{\text{ffn}} \cdot \mathbf{W}_2^\top \in \mathbb{R}^{512 \times 4096} \quad \text{(gradiente per la propagazione)} \\
    \frac{\partial \mathcal{L}}{\partial \mathbf{W}_2} &= \mathbf{h}_{\text{act}}^\top \cdot \delta_{\text{ffn}} \in \mathbb{R}^{4096 \times 1024} \quad \text{(gradiente dei pesi)} \\
    \frac{\partial \mathcal{L}}{\partial \mathbf{b}_2} &= \sum_{i=1}^{512} [\delta_{\text{ffn}}]_i \in \mathbb{R}^{1024} \quad \text{(gradiente del bias)}
\end{align*}

\textbf{2. Attivazione GELU backward} ($\mathbf{h}_{\text{act}} = \text{GELU}(\mathbf{h})$):
\[
    \delta_{\mathbf{h}} = \delta_{\text{h\_act}} \odot \text{GELU}'(\mathbf{h}) \in \mathbb{R}^{512 \times 4096}
\]
dove $\text{GELU}'(x) = \Phi(x) + x \cdot \phi(x)$ (\hyperref[sec:gelu]{Sezione~6.2}). Il prodotto \`e \textbf{elemento per elemento} ($\odot$): ogni componente del gradiente viene moltiplicata per la derivata locale della GELU.

\textbf{3. Espansione backward} (primo strato lineare: $\mathbf{h} = \mathbf{z}_{\text{norm}} \cdot \mathbf{W}_1 + \mathbf{b}_1$):
\begin{align*}
    \delta_{\text{z\_norm}} &= \delta_{\mathbf{h}} \cdot \mathbf{W}_1^\top \in \mathbb{R}^{512 \times 1024} \quad \text{(gradiente per la propagazione)} \\
    \frac{\partial \mathcal{L}}{\partial \mathbf{W}_1} &= \mathbf{z}_{\text{norm}}^\top \cdot \delta_{\mathbf{h}} \in \mathbb{R}^{1024 \times 4096} \quad \text{(gradiente dei pesi)} \\
    \frac{\partial \mathcal{L}}{\partial \mathbf{b}_1} &= \sum_{i=1}^{512} [\delta_{\mathbf{h}}]_i \in \mathbb{R}^{4096} \quad \text{(gradiente del bias)}
\end{align*}
\end{formulabox}

\begin{notabox}[Dimensioni dei gradienti dell'MLP]
\begin{center}
\begin{tabular}{@{}llll@{}}
\toprule
\textbf{Gradiente} & \textbf{Dimensione} & \textbf{Tipo} & \textbf{Uso} \\
\midrule
$\delta_{\text{ffn}}$ & $512 \times 1024$ & Input & Dal residual \\
$\partial\mathbf{W}_2$ & $4096 \times 1024$ & Peso & Aggiornamento $\mathbf{W}_2$ \\
$\delta_{\text{h\_act}}$ & $512 \times 4096$ & Intermedio & Propagazione \\
$\delta_{\mathbf{h}}$ & $512 \times 4096$ & Intermedio & Dopo GELU' \\
$\partial\mathbf{W}_1$ & $1024 \times 4096$ & Peso & Aggiornamento $\mathbf{W}_1$ \\
$\delta_{\text{z\_norm}}$ & $512 \times 1024$ & Output & Verso la LayerNorm \\
\bottomrule
\end{tabular}
\end{center}
\end{notabox}

\begin{examplebox}[Esempio numerico: backward pass completo attraverso un MLP]
Consideriamo un MLP semplificato con $D = 2$ e $D_{\text{ff}} = 4$. Seguiamo il forward e poi il backward con numeri concreti.

\textbf{Pesi del MLP}:
\[
\mathbf{W}_1 = \begin{pmatrix} 0.5 & -0.3 & 0.2 & 0.4 \\ -0.1 & 0.6 & -0.4 & 0.3 \end{pmatrix}, \quad
\mathbf{W}_2 = \begin{pmatrix} 0.3 & -0.2 \\ 0.5 & 0.1 \\ -0.1 & 0.4 \\ 0.2 & -0.3 \end{pmatrix}
\]
$\mathbf{b}_1 = [0, 0, 0, 0]$, $\mathbf{b}_2 = [0, 0]$ (bias nulli per semplicit\`a).

\textbf{Input} (dopo LayerNorm): $\mathbf{z}_{\text{norm}} = [1.0, -0.5]$

\textbf{--- FORWARD ---}

\textbf{1. Espansione}: $\mathbf{h} = \mathbf{z}_{\text{norm}} \cdot \mathbf{W}_1$
\begin{align*}
h_1 &= 1.0 \times 0.5 + (-0.5) \times (-0.1) = 0.55 \\
h_2 &= 1.0 \times (-0.3) + (-0.5) \times 0.6 = -0.60 \\
h_3 &= 1.0 \times 0.2 + (-0.5) \times (-0.4) = 0.40 \\
h_4 &= 1.0 \times 0.4 + (-0.5) \times 0.3 = 0.25
\end{align*}
$\mathbf{h} = [0.55, -0.60, 0.40, 0.25]$

\textbf{2. GELU}: applico la GELU elemento per elemento:
\begin{align*}
\text{GELU}(0.55) &\approx 0.55 \times \Phi(0.55) \approx 0.55 \times 0.709 = 0.390 \\
\text{GELU}(-0.60) &\approx -0.60 \times \Phi(-0.60) \approx -0.60 \times 0.274 = -0.165 \\
\text{GELU}(0.40) &\approx 0.40 \times \Phi(0.40) \approx 0.40 \times 0.655 = 0.262 \\
\text{GELU}(0.25) &\approx 0.25 \times \Phi(0.25) \approx 0.25 \times 0.599 = 0.150
\end{align*}
$\mathbf{h}_{\text{act}} = [0.390, -0.165, 0.262, 0.150]$

\textbf{3. Compressione}: $\mathbf{z}_{\text{ffn}} = \mathbf{h}_{\text{act}} \cdot \mathbf{W}_2$
\begin{align*}
z_{\text{ffn},1} &= 0.390 \times 0.3 + (-0.165) \times 0.5 + 0.262 \times (-0.1) + 0.150 \times 0.2 \\
&= 0.117 - 0.083 - 0.026 + 0.030 = 0.038 \\
z_{\text{ffn},2} &= 0.390 \times (-0.2) + (-0.165) \times 0.1 + 0.262 \times 0.4 + 0.150 \times (-0.3) \\
&= -0.078 - 0.017 + 0.105 - 0.045 = -0.035
\end{align*}
$\mathbf{z}_{\text{ffn}} = [0.038, -0.035]$

\textbf{--- BACKWARD ---}

Supponiamo che il gradiente dal residual sia $\delta_{\text{ffn}} = [0.8, -0.4]$.

\textbf{4. Compressione backward}: $\delta_{\text{h\_act}} = \delta_{\text{ffn}} \cdot \mathbf{W}_2^\top$
\begin{align*}
\delta_{\text{h\_act}} &= [0.8, -0.4] \cdot \begin{pmatrix} 0.3 & 0.5 & -0.1 & 0.2 \\ -0.2 & 0.1 & 0.4 & -0.3 \end{pmatrix} \\
&= [0.8 \times 0.3 + (-0.4) \times (-0.2), \ldots] = [0.32, 0.36, -0.24, 0.28]
\end{align*}

Gradiente dei pesi: $\frac{\partial \mathcal{L}}{\partial \mathbf{W}_2} = \mathbf{h}_{\text{act}}^\top \cdot \delta_{\text{ffn}} \in \mathbb{R}^{4 \times 2}$

\textbf{5. GELU backward}: $\delta_{\mathbf{h}} = \delta_{\text{h\_act}} \odot \text{GELU}'(\mathbf{h})$

Calcoliamo $\text{GELU}'(h_i) = \Phi(h_i) + h_i \cdot \phi(h_i)$:
\begin{align*}
\text{GELU}'(0.55) &= 0.709 + 0.55 \times 0.332 = 0.892 \\
\text{GELU}'(-0.60) &= 0.274 + (-0.60) \times 0.333 = 0.074 \\
\text{GELU}'(0.40) &= 0.655 + 0.40 \times 0.368 = 0.802 \\
\text{GELU}'(0.25) &= 0.599 + 0.25 \times 0.387 = 0.696
\end{align*}

$\delta_{\mathbf{h}} = [0.32 \times 0.892,\; 0.36 \times 0.074,\; -0.24 \times 0.802,\; 0.28 \times 0.696] = [0.285, 0.027, -0.192, 0.195]$

\textbf{6. Espansione backward}: $\delta_{\text{z\_norm}} = \delta_{\mathbf{h}} \cdot \mathbf{W}_1^\top \in \mathbb{R}^{2}$
\begin{align*}
\delta_{\text{z\_norm}} &= [0.285, 0.027, -0.192, 0.195] \cdot \mathbf{W}_1^\top \\
&= [0.285 \times 0.5 + 0.027 \times (-0.3) + (-0.192) \times 0.2 + 0.195 \times 0.4, \ldots] \\
&= [0.173, -0.042]
\end{align*}

Gradiente dei pesi: $\frac{\partial \mathcal{L}}{\partial \mathbf{W}_1} = \mathbf{z}_{\text{norm}}^\top \cdot \delta_{\mathbf{h}} \in \mathbb{R}^{2 \times 4}$

\textbf{Riepilogo}: dal gradiente $\delta_{\text{ffn}} = [0.8, -0.4]$ abbiamo ottenuto:
\begin{itemize}
\item $\frac{\partial \mathcal{L}}{\partial \mathbf{W}_2}$ e $\frac{\partial \mathcal{L}}{\partial \mathbf{W}_1}$: per aggiornare i pesi dell'MLP
\item $\delta_{\text{z\_norm}} = [0.173, -0.042]$: per propagare il gradiente alla LayerNorm sottostante
\end{itemize}
\end{examplebox}

% TikZ: Gradient flow diagram with residual connections
\begin{figure}[H]
\centering
\begin{tikzpicture}[
    block/.style={draw, rounded corners=3pt, fill=#1, minimum width=2.5cm, minimum height=0.6cm, font=\small},
    block/.default={gray!10},
    arrowstyle/.style={-{Stealth[length=2.5mm]}, thick},
    gradarrow/.style={-{Stealth[length=2.5mm]}, thick, red!60!black},
    >=Stealth,
    scale=0.85, transform shape
]

% Forward blocks
\node[block=blue!15] (x) at (0, 0) {$\mathbf{z}$ input};
\node[block=orange!15] (attn) at (0, 2) {Self-Attention};
\node[circle, draw, fill=green!15, minimum size=0.5cm, font=\small] (add1) at (0, 3.5) {$+$};
\node[block=yellow!15] (mlp) at (0, 5) {MLP (FFN)};
\node[circle, draw, fill=green!15, minimum size=0.5cm, font=\small] (add2) at (0, 6.5) {$+$};
\node[block=blue!15] (out) at (0, 8) {$\mathbf{z}$ output};

% Forward arrows
\draw[arrowstyle] (x) -- (attn);
\draw[arrowstyle] (attn) -- (add1);
\draw[arrowstyle] (add1) -- (mlp);
\draw[arrowstyle] (mlp) -- (add2);
\draw[arrowstyle] (add2) -- (out);

% Skip connections (forward)
\draw[arrowstyle, blue!40, dashed] (-2.5, 0) -- (-2.5, 3.5) -- (add1.west);
\draw[arrowstyle, blue!40, dashed] (2.5, 3.5) -- (2.5, 6.5) -- (add2.east);

% Backward gradient arrows (right side)
\draw[gradarrow] (6, 8) -- (6, 6.5) node[midway, right, font=\scriptsize, red!60!black] {$\delta_{\text{out}}$};
\draw[gradarrow] (5.5, 6.5) -- (5.5, 5.5) node[midway, right, font=\scriptsize, red!60!black] {$\delta_{\text{out}}$};
\draw[gradarrow] (6.5, 6.5) -- (6.5, 5) -- (6.5, 3.7) node[midway, right, font=\scriptsize, red!60!black] {$\delta_{\text{out}}$ (skip!)};
\draw[gradarrow] (5.5, 4.5) -- (5.5, 3.5);

\draw[gradarrow] (5, 3.5) -- (5, 2.5) node[midway, right, font=\scriptsize, red!60!black] {$\delta_{\text{sum}}$};
\draw[gradarrow] (7, 3.5) -- (7, 1.5) -- (7, 0.5) node[midway, right, font=\scriptsize, red!60!black] {$\delta_{\text{sum}}$ (skip!)};

% Labels
\node[font=\scriptsize, blue!60!black] at (-3.3, 1.8) {Skip 1};
\node[font=\scriptsize, blue!60!black] at (3.3, 5) {Skip 2};

% Gradient split annotations
\node[font=\scriptsize, red!70!black, fill=red!5, draw=red!30, rounded corners, text width=3cm, align=center] at (8.5, 6.5) {Il gradiente si \textbf{duplica} al residual: una copia salta, una attraversa};

\node[font=\scriptsize, green!50!black, fill=green!5, draw=green!30, rounded corners, text width=2.5cm, align=center] at (8.5, 3.5) {I due rami si \textbf{sommano}: gradiente mai zero};

\end{tikzpicture}
\caption{Flusso del gradiente attraverso le residual connections nel backward pass. Ad ogni connessione residua il gradiente si ``duplica'': una copia attraversa il blocco (Attention o MLP), l'altra salta direttamente. I due rami si sommano al punto di biforcazione. Questo garantisce che il gradiente non svanisca mai, anche dopo 56+ blocchi.}
\label{fig:gradient_flow_residual}
\end{figure}

\subsubsection{LayerNorm backward} \label{sec:ln_backward}

Il backward della LayerNorm \`e forse il pi\`u ``antipatico'' matematicamente, ma \`e importante capirne l'intuizione. Il backward della Layer Normalization (\hyperref[sec:layernorm]{Sezione~2.3}) \`e pi\`u complesso di quello del MLP perch\'e la normalizzazione introduce dipendenze tra le componenti del vettore: cambiare $h_j$ modifica la media $\mu$ e la varianza $\sigma^2$, che influenzano \emph{tutte} le componenti normalizzate.

\begin{formulabox}[Gradiente della Layer Normalization]
Siano $\hat{\mathbf{h}} = \frac{\mathbf{h} - \mu}{\sqrt{\sigma^2 + \epsilon}}$ i valori normalizzati e $\delta_{\hat{\mathbf{h}}} = \frac{\partial \mathcal{L}}{\partial \hat{\mathbf{h}}}$ il gradiente in ingresso.

\textbf{Gradienti dei parametri apprendibili}:
\begin{align*}
    \frac{\partial \mathcal{L}}{\partial \boldsymbol{\gamma}} &= \sum_{i} \delta_{\hat{\mathbf{h}},i} \odot \frac{\mathbf{h}_i - \mu_i}{\sqrt{\sigma_i^2 + \epsilon}} \qquad \text{(gradiente per la scala)} \\
    \frac{\partial \mathcal{L}}{\partial \boldsymbol{\beta}} &= \sum_i \delta_{\hat{\mathbf{h}},i} \qquad \text{(gradiente per lo shift)}
\end{align*}

\textbf{Gradiente rispetto all'input} (per la propagazione all'indietro):
\[
    \frac{\partial \mathcal{L}}{\partial \mathbf{h}_j} = \frac{\boldsymbol{\gamma}}{\sqrt{\sigma^2 + \epsilon}} \left(\delta_{\hat{\mathbf{h}},j} - \frac{1}{D}\sum_{k} \delta_{\hat{\mathbf{h}},k} - \frac{\hat{h}_j}{D}\sum_{k} \delta_{\hat{\mathbf{h}},k} \cdot \hat{h}_k\right)
\]
\end{formulabox}

\begin{notabox}[Interpretazione del gradiente della LayerNorm]
La formula del gradiente della LN ha tre termini:
\begin{enumerate}
    \item $\delta_{\hat{\mathbf{h}},j}$: il gradiente ``diretto'' che passerebbe se non ci fosse la normalizzazione.
    \item $-\frac{1}{D}\sum_k \delta_{\hat{\mathbf{h}},k}$: sottrae la media dei gradienti (centra il gradiente a media zero, come la LN centra le attivazioni).
    \item $-\frac{\hat{h}_j}{D}\sum_k \delta_{\hat{\mathbf{h}},k} \hat{h}_k$: termine di ``decorrelazione'' che tiene conto dell'effetto della varianza.
\end{enumerate}
In pratica, la LN backward ``normalizza'' anche il gradiente, stabilizzando il flusso del gradiente attraverso reti molto profonde. Questo \`e uno dei motivi per cui la pre-norm rende il Perceiver (con fino a 56 blocchi) addestrabile.
\end{notabox}


\subsection{B5. Gradiente della Self-Attention} \label{sec:grad_sa}

Ora affrontiamo la parte pi\`u complessa del backward pass: il gradiente attraverso la self-attention. Questa sezione \`e tecnicamente impegnativa, ma la buona notizia \`e che il principio \`e sempre lo stesso: applicare la chain rule passo per passo, in ordine inverso rispetto al forward.

\begin{pipelinebox}[title={\small Backward --- Passo corrente evidenziato}]
Loss $\leftarrow$ Softmax $\leftarrow$ Classifier $\leftarrow$ Pool $\leftarrow$ ($\times T$) $\leftarrow$ Res+MLP $\leftarrow$ \textbf{Self-Attention (Latent Transf.)} $\leftarrow$ Res+MLP $\leftarrow$ Cross-Att
\end{pipelinebox}

Il backward della self-attention \`e la parte pi\`u complessa del backward pass. La buona notizia \`e che si tratta di applicare la chain rule meccanicamente, passo per passo. La cattiva notizia \`e che ci sono molti passi! Procediamo in ordine inverso rispetto al forward, partendo dall'output e risalendo verso Q, K, V.

\textbf{Struttura del backward della self-attention} (6 sotto-passi):
\begin{enumerate}
    \item Proiezione di output ($\mathbf{W}_O$) $\to$ gradienti per $\mathbf{W}_O$ e per le teste concatenate
    \item Prodotto $\mathbf{A} \cdot \mathbf{V}$ $\to$ gradienti per $\mathbf{A}$ e per $\mathbf{V}$
    \item Softmax $\to$ gradiente per gli score (usando lo Jacobiano della softmax)
    \item Scaling $\to$ gradiente per gli score non scalati
    \item Score $\mathbf{Q}\mathbf{K}^\top$ $\to$ gradienti per $\mathbf{Q}$ e $\mathbf{K}$
    \item Proiezioni Q, K, V $\to$ gradienti per $\mathbf{W}_Q$, $\mathbf{W}_K$, $\mathbf{W}_V$ e per $\mathbf{z}_{\text{norm}}$
\end{enumerate}

\begin{formulabox}[B5.1 --- Proiezione di output]
L'output della multi-head attention passa attraverso una proiezione $\mathbf{W}_O$:
\[
    \text{attn\_out} = \text{heads\_concat} \cdot \mathbf{W}_O
\]
Il backward \`e:
\[
    \frac{\partial \mathcal{L}}{\partial \text{heads\_concat}} = \delta_{\text{attn\_out}} \cdot \mathbf{W}_O^\top, \qquad \frac{\partial \mathcal{L}}{\partial \mathbf{W}_O} = \text{heads\_concat}^\top \cdot \delta_{\text{attn\_out}}
\]
\end{formulabox}

\begin{formulabox}[title={B5.2 --- Prodotto $\mathbf{A} \cdot \mathbf{V}$}]
Partendo da $\delta_{\text{attn}} \in \mathbb{R}^{N \times d_v}$ (gradiente rispetto all'output di una testa):
\begin{align*}
    \frac{\partial \mathcal{L}}{\partial \mathbf{A}} &= \delta_{\text{attn}} \cdot \mathbf{V}^\top \in \mathbb{R}^{N \times N} \\[4pt]
    \frac{\partial \mathcal{L}}{\partial \mathbf{V}} &= \mathbf{A}^\top \cdot \delta_{\text{attn}} \in \mathbb{R}^{N \times d_v}
\end{align*}
\textbf{Intuizione}: il gradiente rispetto ad $\mathbf{A}$ ci dice ``quanto cambiare i pesi di attenzione''; il gradiente rispetto a $\mathbf{V}$ ci dice ``quanto cambiare i valori''.
\end{formulabox}

\begin{formulabox}[B5.3 --- Softmax backward (Jacobiano)]
Per ogni riga $i$ della matrice di attenzione, la softmax ha un Jacobiano complesso. Il gradiente si calcola efficientemente:
\[
    \delta_{\mathbf{s},i} = \mathbf{a}_i \odot \left(\delta_{\mathbf{A},i} - \langle \delta_{\mathbf{A},i}, \mathbf{a}_i \rangle \cdot \mathbf{1}\right)
\]
dove $\mathbf{a}_i$ \`e la $i$-esima riga di $\mathbf{A}$ (pesi di attenzione), $\odot$ \`e il prodotto elemento per elemento, e $\langle \cdot, \cdot \rangle$ \`e il prodotto scalare.

\textbf{Dimostrazione sintetica}: partendo dalla derivata della softmax $\frac{\partial a_k}{\partial s_j} = a_k(\delta_{kj} - a_j)$, e sommando su $k$:
\begin{align*}
    \delta_{s_j} &= \sum_k \frac{\partial \mathcal{L}}{\partial a_k} \cdot \frac{\partial a_k}{\partial s_j} = \sum_k \delta_{A,k} \cdot a_k(\delta_{kj} - a_j) \\
    &= \delta_{A,j} \cdot a_j - a_j \sum_k \delta_{A,k} \cdot a_k = a_j \left(\delta_{A,j} - \langle \delta_{\mathbf{A}}, \mathbf{a} \rangle\right)
\end{align*}
\end{formulabox}

\begin{formulabox}[B5.4 --- Score scaled]
Il gradiente rispetto allo scaling \`e diretto:
\[
    \delta_{\mathbf{S}} = \frac{1}{\sqrt{d_k}} \cdot \delta_{\mathbf{s}} \in \mathbb{R}^{N \times N}
\]
\end{formulabox}

\begin{formulabox}[B5.5 --- Gradienti di Q e K]
Dato lo score $\mathbf{S} = \mathbf{Q} \cdot \mathbf{K}^\top$:
\begin{align*}
    \frac{\partial \mathcal{L}}{\partial \mathbf{Q}} &= \frac{1}{\sqrt{d_k}} \delta_{\mathbf{S}} \cdot \mathbf{K} \in \mathbb{R}^{N \times d_k} \\[4pt]
    \frac{\partial \mathcal{L}}{\partial \mathbf{K}} &= \frac{1}{\sqrt{d_k}} \delta_{\mathbf{S}}^\top \cdot \mathbf{Q} \in \mathbb{R}^{N \times d_k}
\end{align*}
\end{formulabox}

\begin{formulabox}[title={B5.6 --- Proiezioni Q, K, V e combinazione}]
I gradienti delle matrici di proiezione:
\begin{align*}
    \frac{\partial \mathcal{L}}{\partial \mathbf{W}_Q} &= \mathbf{z}_{\text{norm}}^\top \cdot \frac{\partial \mathcal{L}}{\partial \mathbf{Q}} \in \mathbb{R}^{D \times d_k} \\[4pt]
    \frac{\partial \mathcal{L}}{\partial \mathbf{W}_K} &= \mathbf{z}_{\text{norm}}^\top \cdot \frac{\partial \mathcal{L}}{\partial \mathbf{K}} \in \mathbb{R}^{D \times d_k} \\[4pt]
    \frac{\partial \mathcal{L}}{\partial \mathbf{W}_V} &= \mathbf{z}_{\text{norm}}^\top \cdot \frac{\partial \mathcal{L}}{\partial \mathbf{V}} \in \mathbb{R}^{D \times d_v}
\end{align*}

Nella self-attention, Q, K, V vengono tutti dalla stessa sorgente $\mathbf{z}_{\text{norm}}$, quindi il gradiente totale \`e la somma dei tre percorsi:
\[
\frac{\partial \mathcal{L}}{\partial \mathbf{z}_{\text{norm}}} = \frac{\partial \mathcal{L}}{\partial \mathbf{Q}} \cdot \mathbf{W}_Q^\top + \frac{\partial \mathcal{L}}{\partial \mathbf{K}} \cdot \mathbf{W}_K^\top + \frac{\partial \mathcal{L}}{\partial \mathbf{V}} \cdot \mathbf{W}_V^\top
\]
I tre contributi si sommano perch\'e $\mathbf{z}_{\text{norm}}$ influenza l'output attraverso tutti e tre i percorsi (Q, K, V). Questo \`e un'applicazione diretta della \textbf{regola della somma} per la chain rule multivariata.
\end{formulabox}

\begin{notabox}[Riepilogo dimensionale del backward della self-attention]
\begin{center}
\begin{tabular}{@{}lll@{}}
\toprule
\textbf{Gradiente} & \textbf{Dimensione} & \textbf{Nota} \\
\midrule
$\delta_{\text{attn\_out}}$ & $512 \times 1024$ & Dal residual \\
$\delta_{\text{heads\_concat}}$ & $512 \times 1024$ & Via $\mathbf{W}_O^\top$ \\
$\partial\mathbf{A}$ & $512 \times 512$ (per testa) & $\delta \cdot \mathbf{V}^\top$ \\
$\delta_{\mathbf{s}}$ & $512 \times 512$ (per testa) & Softmax Jacobiano \\
$\partial\mathbf{Q}, \partial\mathbf{K}$ & $512 \times 128$ (per testa) & Score backward \\
$\partial\mathbf{V}$ & $512 \times 128$ (per testa) & $\mathbf{A}^\top \cdot \delta$ \\
$\partial\mathbf{z}_{\text{norm}}$ & $512 \times 1024$ & Somma Q+K+V \\
\bottomrule
\end{tabular}
\end{center}
\end{notabox}


\subsection{B6. Gradiente della Cross-Attention} \label{sec:grad_ca}

Il backward della cross-attention \`e simile a quello della self-attention, ma con una differenza cruciale: Q proviene da una sorgente diversa rispetto a K e V. Vediamo cosa cambia.

\begin{pipelinebox}[title={\small Backward --- Passo corrente evidenziato}]
Loss $\leftarrow$ Softmax $\leftarrow$ Classifier $\leftarrow$ Pool $\leftarrow$ ($\times T$) $\leftarrow$ Res+MLP $\leftarrow$ Latent Transf.\ $\leftarrow$ Res+MLP $\leftarrow$ \textbf{Cross-Att}
\end{pipelinebox}

\begin{formulabox}[Differenza chiave: gradiente della cross-attention]
Nella \textbf{self-attention}, Q, K, V provengono tutti dalla stessa sorgente $\mathbf{z}$, quindi il gradiente rispetto a $\mathbf{z}_{\text{norm}}$ \`e la somma dei tre contributi (da Q, K e V).

Nella \textbf{cross-attention}, Q proviene dai latenti $\mathbf{z}$, K e V dall'input $\mathbf{x}$. Quindi i gradienti si \textbf{separano}:
\begin{align*}
    \frac{\partial \mathcal{L}}{\partial \mathbf{z}} &= \frac{\partial \mathcal{L}}{\partial \mathbf{z}}\bigg|_Q \quad \text{(solo dal percorso delle query)} \\[4pt]
    \frac{\partial \mathcal{L}}{\partial \mathbf{x}_{\text{embed}}} &= \frac{\partial \mathcal{L}}{\partial \mathbf{x}}\bigg|_K + \frac{\partial \mathcal{L}}{\partial \mathbf{x}}\bigg|_V \quad \text{(da chiavi e valori)}
\end{align*}
I latenti ricevono gradiente solo attraverso Q; l'input riceve gradiente solo attraverso K e V.
\end{formulabox}

\begin{formulabox}[Gradienti espliciti della cross-attention]
I gradienti delle proiezioni nella cross-attention sono:

\textbf{Gradiente rispetto a $\mathbf{W}_Q$} (solo dai latenti):
\[
    \frac{\partial \mathcal{L}}{\partial \mathbf{W}_Q} = \mathbf{z}_{\text{norm}}^\top \cdot \frac{\partial \mathcal{L}}{\partial \mathbf{Q}} \in \mathbb{R}^{D \times d_k}
\]

\textbf{Gradiente rispetto a $\mathbf{W}_K$} (dall'input):
\[
    \frac{\partial \mathcal{L}}{\partial \mathbf{W}_K} = \mathbf{x}_{\text{norm}}^\top \cdot \frac{\partial \mathcal{L}}{\partial \mathbf{K}} \in \mathbb{R}^{C \times d_k}
\]

\textbf{Gradiente rispetto a $\mathbf{W}_V$} (dall'input):
\[
    \frac{\partial \mathcal{L}}{\partial \mathbf{W}_V} = \mathbf{x}_{\text{norm}}^\top \cdot \frac{\partial \mathcal{L}}{\partial \mathbf{V}} \in \mathbb{R}^{C \times d_v}
\]
\end{formulabox}

\begin{notabox}[Perch\'e la separazione Q vs K/V \`e importante]
Nella cross-attention, i latenti ``pongono domande'' (Q) e l'input ``risponde'' (K, V). Durante il backward, il gradiente riflette questa asimmetria: i latenti vengono aggiornati solo in base a ``quanto bene hanno posto le domande'' (gradiente via Q), mentre i pesi di proiezione dell'input vengono aggiornati in base a ``quanto bene hanno descritto e fornito informazione'' (gradiente via K e V). Questa separazione \`e ci\`o che rende il Perceiver un \textbf{lettore asimmetrico}: i latenti sono ``attivi'' (fanno domande) e l'input \`e ``passivo'' (fornisce risposte).
\end{notabox}


\begin{notabox}[In parole semplici: il backward della cross-attention]
Detto in parole povere: durante il backward, il Perceiver ``riflette'' su come ha letto l'immagine. I latenti ricevono un feedback su ``hai fatto le domande giuste?'' (gradiente via Q), e l'input riceve un feedback su ``hai dato le descrizioni giuste?'' e ``hai fornito le informazioni giuste?'' (gradienti via K e V). Questa separazione riflette la natura asimmetrica della cross-attention: i latenti sono ``attivi'' (chiedono) e l'input \`e ``passivo'' (risponde).
\end{notabox}

\subsection{B7. Gradiente della proiezione dell'input (Embedding)} \label{sec:grad_embed}

\begin{pipelinebox}[title={\small Backward --- Passo corrente evidenziato}]
Loss $\leftarrow$ \ldots\ $\leftarrow$ Cross-Att $\leftarrow$ \textbf{Embedding} $\leftarrow$ Fourier $\leftarrow$ Input
\end{pipelinebox}

Il gradiente della proiezione lineare dell'input (\hyperref[sec:proiezione]{Sezione~5.3}) \`e:
\begin{formulabox}[Gradiente della proiezione dell'input]
\begin{align*}
    \frac{\partial \mathcal{L}}{\partial \mathbf{W}_{\text{embed}}} &= \tilde{\mathbf{x}}^\top \cdot \frac{\partial \mathcal{L}}{\partial \mathbf{x}_{\text{embed}}} \in \mathbb{R}^{261 \times 512} \\[4pt]
    \frac{\partial \mathcal{L}}{\partial \mathbf{b}_{\text{embed}}} &= \sum_i \frac{\partial \mathcal{L}}{\partial [\mathbf{x}_{\text{embed}}]_i} \in \mathbb{R}^{512}
\end{align*}
Nota: il gradiente rispetto all'input grezzo $\tilde{\mathbf{x}}$ \textbf{non serve} per aggiornare parametri (le Fourier features non hanno parametri apprendibili), ma in framework come PyTorch viene comunque calcolato per completezza del grafo computazionale.
\end{formulabox}


\subsection{B8. Weight Sharing e Accumulo dei Gradienti} \label{sec:grad_ws}

L'ultimo aspetto del backward pass riguarda il weight sharing. Quando gli stessi pesi sono usati in pi\`u iterazioni, i gradienti da tutte le iterazioni devono essere \textbf{accumulati} prima di aggiornare i parametri. Vediamo come funziona.

\begin{formulabox}[Gradiente con weight sharing]
Quando i pesi sono condivisi tra le iterazioni $t = 2, \ldots, T$:
\[
    \frac{\partial \mathcal{L}}{\partial \theta_{\text{shared}}} = \sum_{t=2}^{T} \frac{\partial \mathcal{L}}{\partial \theta^{(t)}}
\]
L'aggiornamento avviene su un \textbf{singolo set} di parametri condivisi, ma il gradiente viene accumulato da tutte le $T-1$ iterazioni. Questo \`e esattamente ci\`o che accade nella backpropagation through time (BPTT) delle RNN (\hyperref[app:rnn]{Appendice~C}).
\end{formulabox}

\begin{notabox}[Il weight sharing come BPTT]
L'analogia con le RNN \`e diretta:
\begin{itemize}
    \item \textbf{RNN}: $\mathbf{h}_t = f(\mathbf{h}_{t-1}, \mathbf{x}_t; \theta)$ dove $\mathbf{x}_t$ cambia ad ogni passo.
    \item \textbf{Perceiver}: $\mathbf{z}_t = g(\mathbf{z}_{t-1}, \mathbf{x}; \theta_{\text{shared}})$ dove $\mathbf{x}$ \`e \textbf{costante} (lo stesso input viene riletto).
\end{itemize}
In entrambi i casi, il gradiente dei pesi condivisi \`e la somma dei gradienti di tutte le iterazioni. La differenza chiave \`e che nel Perceiver l'input non cambia, quindi le iterazioni successive raffinano progressivamente la rappresentazione dello \emph{stesso} input, anzich\'e elaborare input diversi come nella RNN.

Il Latent Transformer ha pesi condivisi tra \emph{tutte} le iterazioni ($t = 1, \ldots, T$), quindi accumula $T$ contributi:
\[
    \frac{\partial \mathcal{L}}{\partial \theta_{\text{LT}}} = \sum_{t=1}^{T} \frac{\partial \mathcal{L}}{\partial \theta_{\text{LT}}^{(t)}}
\]
\end{notabox}

\begin{notabox}[In parole semplici: accumulo dei gradienti nel weight sharing]
Detto in parole povere: quando gli stessi pesi sono usati in 8 iterazioni, durante il backward ciascuna iterazione ``suggerisce'' una correzione diversa per gli stessi pesi. Il gradiente finale \`e la \textbf{somma} di tutti questi suggerimenti: ``l'iterazione 2 dice di alzare questo peso, l'iterazione 5 dice di abbassarlo, l'iterazione 7 dice di alzarlo poco...'' Il risultato \`e un compromesso che funziona ragionevolmente bene per tutte le iterazioni, ed \`e proprio questo compromesso che agisce come regolarizzazione.
\end{notabox}

\begin{examplebox}[Conteggio parametri: senza vs con weight sharing]
\textbf{Senza WS} ($T = 8$ iterazioni, $L = 6$ layer SA per iterazione):
\begin{itemize}
    \item Cross-attention: 8 moduli unici $\times$ parametri CA $\approx 8 \times 12\text{M} = 96\text{M}$
    \item Latent Transformer: 8 torri uniche $\times$ 6 blocchi $\times$ parametri blocco $\approx 8 \times 6 \times 4.7\text{M} = 226\text{M}$
    \item Classificatore + embedding + latenti: $\approx 4\text{M}$
    \item \textbf{Totale: $\sim$326M parametri}
\end{itemize}

\textbf{Con WS}:
\begin{itemize}
    \item Cross-attention: 2 moduli (1 unico + 1 condiviso) $\times 12\text{M} = 24\text{M}$
    \item Latent Transformer: 1 torre condivisa $\times$ 6 blocchi $\times 4.7\text{M} \approx 17\text{M}$
    \item Classificatore + embedding + latenti: $\approx 4\text{M}$
    \item \textbf{Totale: $\sim$45M parametri (riduzione $7.3\times$)}
\end{itemize}
\end{examplebox}


\newpage
%% ============================================================
%%  SEZIONE 8 --- PROCESSO DI ADDESTRAMENTO
%% ============================================================
\section{Processo di Addestramento} \label{sec:training}

Come si passa da un modello con pesi casuali (che classifica a caso) a un modello che raggiunge il 78\% di accuratezza su ImageNet?

In pratica, quello che succede durante l'addestramento \`e un ciclo continuo: il modello fa una predizione (forward pass), misura l'errore (loss), calcola come correggere i pesi (backward pass), e aggiorna i pesi (ottimizzatore). Questo ciclo si ripete milioni di volte.

% TikZ: Training loop circular diagram
\begin{figure}[H]
\centering
\begin{tikzpicture}[
    stepbox/.style={draw, rounded corners=5pt, fill=#1, minimum width=2.8cm, minimum height=0.8cm, font=\small\bfseries, text=white},
    arrowstyle/.style={-{Stealth[length=3mm]}, very thick, #1},
    scale=0.85, transform shape
]

% Four steps in a circle
\node[stepbox=blue!70!black] (fw) at (0, 2.5) {Forward Pass};
\node[stepbox=red!70!black] (loss) at (3.5, 0) {Calcolo Loss};
\node[stepbox=orange!70!black] (bw) at (0, -2.5) {Backward Pass};
\node[stepbox=green!60!black] (upd) at (-3.5, 0) {Aggiornamento};

% Arrows
\draw[arrowstyle=blue!60!black] (fw.east) .. controls (3, 2) and (3.5, 1.2) .. (loss.north);
\node[font=\scriptsize, blue!60!black] at (3.2, 1.8) {predizione};

\draw[arrowstyle=red!60!black] (loss.south) .. controls (3, -1.2) and (2, -2.5) .. (bw.east);
\node[font=\scriptsize, red!60!black] at (3.2, -1.8) {errore};

\draw[arrowstyle=orange!60!black] (bw.west) .. controls (-2, -2.5) and (-3.5, -1.2) .. (upd.south);
\node[font=\scriptsize, orange!60!black] at (-3.2, -1.8) {gradienti};

\draw[arrowstyle=green!60!black] (upd.north) .. controls (-3.5, 1.2) and (-2, 2.5) .. (fw.west);
\node[font=\scriptsize, green!60!black] at (-3.2, 1.8) {pesi migliori};

% Center annotation
\node[font=\scriptsize, gray, text width=2cm, align=center] at (0, 0) {$\sim$2.4M\\iterazioni};

\end{tikzpicture}
\caption{Il ciclo di addestramento del Perceiver. Ad ogni iterazione: (1) il forward pass produce una predizione, (2) la loss misura l'errore, (3) il backward pass calcola i gradienti, (4) l'ottimizzatore LAMB aggiorna i pesi. Questo ciclo si ripete circa 2.4 milioni di volte per 120 epoche su ImageNet.}
\label{fig:training_loop}
\end{figure}

\subsection{Dataset e configurazione} \label{sec:dataset}

\begin{formulabox}[Dettagli dell'addestramento su ImageNet]
\begin{itemize}
    \item \textbf{Dataset}: ImageNet ILSVRC-2012, $\sim$1.2 milioni di immagini, 1000 classi.
    \item \textbf{Batch size}: 64 immagini per batch.
    \item \textbf{Epoche}: 120 epoche. Iterazioni totali: $\sim$2.25M.
    \item \textbf{Ottimizzatore}: LAMB con learning rate $= 0.004$.
    \item \textbf{Schedule}: decadimento per fattore 10 alle epoche 84, 102 e 114.
    \item \textbf{Data augmentation}: RandAugment, random crop $224 \times 224$, flip orizzontale.
\end{itemize}
\end{formulabox}

\subsubsection{Learning rate schedule} \label{sec:lr_schedule}

Un aspetto cruciale che spesso viene trascurato \`e che il learning rate non \`e costante durante il training. Variare il learning rate durante il training \`e una delle tecniche pi\`u importanti per ottenere buone prestazioni. Il Perceiver usa un \textbf{step decay schedule}, che \`e il pi\`u semplice tra i possibili schemi:

\begin{formulabox}[Learning rate schedule del Perceiver]
\[
    \eta(t) = \begin{cases}
        0.004 & \text{epoche } 1\text{--}83 \\
        0.0004 & \text{epoche } 84\text{--}101 \quad (\div 10) \\
        0.00004 & \text{epoche } 102\text{--}113 \quad (\div 100) \\
        0.000004 & \text{epoche } 114\text{--}120 \quad (\div 1000)
    \end{cases}
\]
\end{formulabox}

\textbf{Perch\'e decrescere il learning rate?} All'inizio del training, il modello \`e lontano dall'ottimo: serve un LR grande per esplorare rapidamente lo spazio dei parametri. Man mano che ci si avvicina all'ottimo, un LR grande causerebbe oscillazioni attorno al minimo. Riducendo il LR, il modello ``rallenta'' e converge verso un minimo pi\`u preciso.

\begin{notabox}[In parole semplici: il learning rate schedule]
Immaginiamo di cercare un tesoro in un campo enorme. All'inizio fai passi grandi per esplorare rapidamente zone diverse (LR alto). Quando sei vicino al tesoro, fai passi piccoli e cauti per non superarlo (LR basso). Il Perceiver fa esattamente questo: passi grandi per le prime 83 epoche, poi riduce di 10 volte, poi di 100 volte, poi di 1000 volte.
\end{notabox}

\subsubsection{Data Augmentation} \label{sec:data_augment}

La domanda che ci si pone \`e: 1.2 milioni di immagini bastano per addestrare 45 milioni di parametri? La risposta \`e: a malapena. Per evitare che il modello ``memorizzi'' le immagini del training set invece di imparare pattern generali (overfitting), si usa la \textbf{data augmentation}: trasformazioni casuali che ``moltiplicano'' artificialmente la variet\`a del dataset.

\begin{itemize}
    \item \textbf{Random crop $224 \times 224$}: l'immagine viene ridimensionata a una scala casuale e poi ritagliata a $224 \times 224$. Questo espone il modello a scale diverse dell'oggetto.
    \item \textbf{Flip orizzontale} (probabilit\`a 50\%): l'immagine viene specchiata orizzontalmente. Un gatto che guarda a sinistra diventa un gatto che guarda a destra.
    \item \textbf{RandAugment}: applica trasformazioni casuali (rotazione, traslazione, contrasto, saturazione, ecc.) con intensit\`a e combinazioni casuali.
\end{itemize}

\begin{notabox}[Data augmentation e il Perceiver]
La data augmentation \`e particolarmente importante per il Perceiver perch\'e l'architettura non ha inductive bias spaziali (no localit\`a, no invarianza traslazionale). Mentre una CNN ``sa'' automaticamente che un gatto spostato di 5 pixel \`e lo stesso gatto (grazie ai filtri condivisi), il Perceiver deve \emph{apprendere} questa invarianza dai dati. La data augmentation fornisce gli esempi necessari per questo apprendimento.
\end{notabox}

\subsubsection{Calcolo del numero di iterazioni} \label{sec:num_iterazioni}

Per avere un'idea concreta della scala del training, calcoliamo quante iterazioni servono e quanto tempo richiedono.

\begin{examplebox}[Calcolo delle iterazioni di training]
\begin{align*}
    \text{Immagini totali} &= 1{,}281{,}167 \approx 1.28\text{M} \\
    \text{Batch size} &= 64 \\
    \text{Iterazioni per epoca} &= \lceil 1{,}281{,}167 / 64 \rceil = 20{,}019 \\
    \text{Epoche totali} &= 120 \\
    \text{Iterazioni totali} &= 120 \times 20{,}019 \approx 2.4\text{M}
\end{align*}
Ad ogni iterazione, il modello processa 64 immagini, ciascuna con $50{,}176$ pixel $\times$ 261 feature = $\sim$13M valori per immagine, $\sim$839M valori per batch. Il forward + backward pass di ogni iterazione richiede $\sim$10--20 secondi su una TPU v3 (stima).

\textbf{Tempo totale stimato}: $2.4$M iter $\times$ 15 sec/iter $\approx$ 10{,}000 ore GPU $\approx$ 416 giorni su 1 GPU. In pratica, il training viene distribuito su pi\`u acceleratori.
\end{examplebox}


\subsection{Ottimizzatori: SGD vs Adam vs LAMB} \label{sec:ottimizzatori}

Ora che abbiamo capito cosa calcola il backward pass (i gradienti), dobbiamo chiederci: come li usiamo per aggiornare i pesi? La risposta \`e l'ottimizzatore. Questa sottosezione descrive in dettaglio i tre principali ottimizzatori rilevanti per il training del Perceiver, partendo dal pi\`u semplice (SGD) fino al pi\`u sofisticato (LAMB).

\subsubsection{SGD con Momentum}

Partiamo dall'ottimizzatore pi\`u semplice. Anche se il Perceiver non lo usa, capire SGD \`e essenziale per apprezzare i miglioramenti introdotti da Adam e LAMB.

\begin{formulabox}[SGD con Momentum --- algoritmo completo]
\textbf{Inizializzazione}: $\mathbf{v}_0 = \mathbf{0}$

\textbf{Ad ogni passo $t$}:
\begin{enumerate}
    \item Calcola il gradiente: $\mathbf{g}_t = \nabla_\theta \mathcal{L}(\theta_t)$
    \item Aggiorna la velocit\`a: $\mathbf{v}_t = \beta \cdot \mathbf{v}_{t-1} + \mathbf{g}_t$
    \item Aggiorna i parametri: $\theta_{t+1} = \theta_t - \eta \cdot \mathbf{v}_t$
\end{enumerate}
Iperparametri: $\eta$ (learning rate), $\beta = 0.9$ (momentum).
\end{formulabox}

\textbf{Intuizione: la palla che rotola gi\`u dalla collina.} SGD senza momentum \`e come una palla che si muove solo in base alla pendenza locale: oscilla nelle valli strette e avanza lentamente nelle zone piatte. Il momentum aggiunge ``inerzia'': la palla accumula velocit\`a nella direzione prevalente, smorzando le oscillazioni trasversali e accelerando nella direzione giusta.

\textbf{Limite di SGD}: tutti i parametri condividono lo stesso learning rate globale $\eta$. Se un parametro ha gradienti grandi (es.\ il classificatore) e un altro ha gradienti piccoli (es.\ i latenti iniziali), lo stesso $\eta$ non pu\`o essere ottimale per entrambi.

\subsubsection{Adam (Adaptive Moment Estimation)}

Ora passiamo all'ottimizzatore che ha rivoluzionato il deep learning. Adam (Kingma \& Ba, 2015) risolve il problema principale di SGD: adatta il learning rate \textbf{per ogni parametro singolarmente}, anzich\'e usare un unico learning rate globale.

\begin{formulabox}[Adam --- algoritmo completo in 4 passi]
\textbf{Inizializzazione}: $\mathbf{m}_0 = \mathbf{0}$, $\mathbf{v}_0 = \mathbf{0}$, $t = 0$

\textbf{Ad ogni passo $t$}:

\textbf{Passo 1 --- Primo momento (media mobile dei gradienti)}:
\[
    \mathbf{m}_t = \beta_1 \mathbf{m}_{t-1} + (1-\beta_1) \mathbf{g}_t
\]
Intuizione: $\mathbf{m}_t$ \`e una stima della \emph{direzione media} del gradiente. Stabilizza l'aggiornamento mediando le oscillazioni.

\textbf{Passo 2 --- Secondo momento (media mobile dei gradienti al quadrato)}:
\[
    \mathbf{v}_t = \beta_2 \mathbf{v}_{t-1} + (1-\beta_2) \mathbf{g}_t^2
\]
Intuizione: $\mathbf{v}_t$ stima la \emph{varianza} del gradiente. Parametri con gradienti variabili riceveranno un learning rate pi\`u piccolo.

\textbf{Passo 3 --- Bias correction} (necessaria perch\'e $\mathbf{m}_0 = \mathbf{v}_0 = \mathbf{0}$):
\[
    \hat{\mathbf{m}}_t = \frac{\mathbf{m}_t}{1 - \beta_1^t}, \qquad \hat{\mathbf{v}}_t = \frac{\mathbf{v}_t}{1 - \beta_2^t}
\]

\textbf{Passo 4 --- Aggiornamento dei parametri}:
\[
    \theta_{t+1} = \theta_t - \eta \cdot \frac{\hat{\mathbf{m}}_t}{\sqrt{\hat{\mathbf{v}}_t} + \epsilon}
\]
Iperparametri: $\beta_1 = 0.9$, $\beta_2 = 0.999$, $\epsilon = 10^{-8}$, $\eta = 0.001$.
\end{formulabox}

\textbf{Perch\'e Adam funziona}: il denominatore $\sqrt{\hat{\mathbf{v}}_t} + \epsilon$ normalizza l'aggiornamento per la ``scala'' tipica dei gradienti di quel parametro. Un parametro con gradienti tipicamente grandi ($\hat{\mathbf{v}}_t$ grande) riceve un learning rate effettivo piccolo; un parametro con gradienti tipicamente piccoli ($\hat{\mathbf{v}}_t$ piccolo) riceve un learning rate effettivo grande. Questo ``adattamento'' rende Adam molto pi\`u robusto alla scelta dell'iperparametro $\eta$.

\textbf{Problema con batch grandi}: Adam tende a ``sotto-performare'' con batch molto grandi ($>$ 8192). Il motivo \`e che il trust ratio (rapporto tra norma dei pesi e norma dell'aggiornamento) diventa molto sbilanciato tra layer diversi. Questo ha motivato la creazione di LAMB.

\begin{notabox}[In parole semplici: Adam]
Detto in parole povere, Adam \`e come un insegnante che d\`a voti personalizzati a ogni studente (parametro). Se uno studente ha prestazioni molto variabili (gradienti oscillanti), l'insegnante \`e pi\`u cauto nel giudicarlo (learning rate pi\`u piccolo). Se uno studente ha prestazioni stabili, l'insegnante si fida di pi\`u del suo trend (learning rate pi\`u grande). Questo ``adattamento personalizzato'' rende Adam molto pi\`u efficace del semplice SGD.
\end{notabox}

\begin{examplebox}[Esempio numerico: un passo di Adam]
Consideriamo un singolo parametro $\theta_0 = 1.5$, con gradiente $g_1 = 0.4$.

\textbf{Inizializzazione}: $m_0 = 0$, $v_0 = 0$, $\eta = 0.001$, $\beta_1 = 0.9$, $\beta_2 = 0.999$.

\textbf{Passo 1 ($t = 1$)}:
\begin{align*}
    m_1 &= 0.9 \times 0 + 0.1 \times 0.4 = 0.04 \\
    v_1 &= 0.999 \times 0 + 0.001 \times 0.4^2 = 0.00016 \\
    \hat{m}_1 &= 0.04 / (1 - 0.9^1) = 0.04 / 0.1 = 0.4 \\
    \hat{v}_1 &= 0.00016 / (1 - 0.999^1) = 0.00016 / 0.001 = 0.16 \\
    \theta_1 &= 1.5 - 0.001 \times \frac{0.4}{\sqrt{0.16} + 10^{-8}} = 1.5 - 0.001 \times 1.0 = 1.499
\end{align*}

\textbf{Osservazione}: la bias correction \`e fondamentale al primo passo! Senza di essa, $\hat{m}_1 = 0.04$ (10$\times$ pi\`u piccolo del gradiente reale $0.4$), e l'aggiornamento sarebbe 10$\times$ troppo piccolo.

\textbf{Passo 2} con $g_2 = -0.2$ (gradiente di segno opposto):
\begin{align*}
    m_2 &= 0.9 \times 0.04 + 0.1 \times (-0.2) = 0.036 - 0.02 = 0.016 \\
    v_2 &= 0.999 \times 0.00016 + 0.001 \times 0.04 = 0.000200 \\
    \hat{m}_2 &= 0.016 / (1 - 0.81) = 0.016 / 0.19 = 0.0842 \\
    \hat{v}_2 &= 0.000200 / (1 - 0.998) = 0.000200 / 0.002 = 0.100 \\
    \theta_2 &= 1.499 - 0.001 \times 0.0842 / (\sqrt{0.100} + 10^{-8}) = 1.499 - 0.000266 = 1.4987
\end{align*}

L'aggiornamento al passo 2 \`e pi\`u piccolo perch\'e il primo momento $\hat{m}_2 \approx 0.084$ media i due gradienti di segno opposto ($0.4$ e $-0.2$), smorzando le oscillazioni.
\end{examplebox}

\subsubsection{LAMB (Layer-wise Adaptive Moments for Batch training)}

Infine, arriviamo all'ottimizzatore scelto per il Perceiver. Perch\'e non basta Adam? Il problema emerge con batch di dimensione grande: Adam non tiene conto della \emph{scala} dei pesi di ogni layer, causando aggiornamenti sproporzionati.

LAMB (You et al., 2020) estende Adam con una normalizzazione \textbf{layer-wise} che permette il training con batch enormi:

\begin{formulabox}[LAMB --- algoritmo completo in 6 passi]
\textbf{Inizializzazione}: come Adam ($\mathbf{m}_0 = \mathbf{0}$, $\mathbf{v}_0 = \mathbf{0}$)

\textbf{Passi 1--3}: identici ad Adam (primo momento, secondo momento, bias correction).

\textbf{Passo 4 --- Aggiornamento tipo Adam} (senza applicarlo ancora):
\[
    \mathbf{r}_t = \frac{\hat{\mathbf{m}}_t}{\sqrt{\hat{\mathbf{v}}_t} + \epsilon} + \lambda \, \theta_t
\]
dove $\lambda$ \`e il weight decay. L'aggiornamento $\mathbf{r}_t$ \`e calcolato ma \textbf{non ancora applicato}.

\textbf{Passo 5 --- Trust ratio} (la novit\`a di LAMB):
\[
    \phi = \frac{\|\theta_t\|_2}{\|\mathbf{r}_t\|_2}
\]
Il trust ratio $\phi$ \`e calcolato \textbf{per ogni layer} separatamente. Se $\|\theta_t\| = 0$, si pone $\phi = 1$.

\textbf{Passo 6 --- Aggiornamento scalato}:
\[
    \theta_{t+1} = \theta_t - \eta \cdot \phi \cdot \mathbf{r}_t
\]
\end{formulabox}

\textbf{Perch\'e il trust ratio \`e fondamentale per il Perceiver}: nel Perceiver coesistono moduli con scale molto diverse:
\begin{itemize}
    \item \textbf{Latenti iniziali}: $\|\theta\| \approx \sqrt{N \cdot D} \times 0.02 \approx 14.5$ (valori piccoli, TruncatedNormal con $\sigma = 0.02$).
    \item \textbf{Pesi di proiezione} ($\mathbf{W}_Q, \mathbf{W}_K, \mathbf{W}_V$): $\|\theta\| \approx \sqrt{D \cdot d_k} \times 0.03 \approx 30$ (Xavier init).
    \item \textbf{Classificatore}: $\|\theta\| \approx \sqrt{D \cdot 1000} \times 0.03 \approx 30$.
\end{itemize}
Senza il trust ratio, un learning rate ottimale per un layer sarebbe troppo grande (o piccolo) per un altro. LAMB riscala automaticamente l'aggiornamento in modo che sia proporzionale alla norma dei pesi di quel layer.

\begin{notabox}[In parole semplici: LAMB e il trust ratio]
Detto in parole povere, LAMB \`e come Adam ma con un ``guardiano'' per ogni layer. Adam dice ``aggiorna questo parametro di tanto'', ma il guardiano (trust ratio) controlla: ``aspetta, questo layer ha pesi molto piccoli, un aggiornamento cos\`i grande lo destabilizzerebbe --- ridimensiono''. Per un altro layer con pesi grandi, il guardiano dice ``questo aggiornamento \`e proporzionato, pu\`o passare''. Questo bilanciamento automatico \`e fondamentale nel Perceiver, dove coesistono i latenti (pesi $\sim 0.02$) e le proiezioni (pesi $\sim 0.05$).
\end{notabox}

\begin{examplebox}[Esempio numerico: un passo di LAMB]
Consideriamo un layer con $\theta_t = [1.0, -0.5, 0.3]$ e gradiente $\mathbf{g}_t = [0.1, -0.2, 0.05]$.

\textbf{Passo 1}: $\mathbf{m}_1 = 0.1 \times [0.1, -0.2, 0.05] = [0.01, -0.02, 0.005]$ (con $\beta_1 = 0.9$, $\mathbf{m}_0 = 0$).

\textbf{Passo 2}: $\mathbf{v}_1 = 0.001 \times [0.01, 0.04, 0.0025] = [10^{-5}, 4 \times 10^{-5}, 2.5 \times 10^{-6}]$.

\textbf{Passo 3}: $\hat{\mathbf{m}}_1 = \mathbf{m}_1 / (1 - 0.9) = [0.1, -0.2, 0.05]$.
$\hat{\mathbf{v}}_1 = \mathbf{v}_1 / (1 - 0.999) = [0.01, 0.04, 0.0025]$.

\textbf{Passo 4}: $\mathbf{r}_1 = [0.1/0.1, -0.2/0.2, 0.05/0.05] = [1.0, -1.0, 1.0]$ (semplificato, senza weight decay).

\textbf{Passo 5}: $\|\theta_t\|_2 = \sqrt{1.0 + 0.25 + 0.09} = 1.158$, $\|\mathbf{r}_1\|_2 = \sqrt{3} = 1.732$.
$\phi = 1.158 / 1.732 = 0.669$.

\textbf{Passo 6}: $\theta_1 = [1.0, -0.5, 0.3] - 0.004 \times 0.669 \times [1.0, -1.0, 1.0] = [0.997, -0.497, 0.297]$.

L'aggiornamento \`e piccolo e proporzionato alla scala dei pesi del layer.
\end{examplebox}

\subsection{Confronto tabulare degli ottimizzatori} \label{sec:confronto_ottimizzatori}

Raccogliamo ora i tre ottimizzatori in una tabella comparativa. Questa tabella \`e utile per l'esame: mostra in modo compatto le differenze chiave tra SGD, Adam e LAMB, e spiega perch\'e il Perceiver ha scelto LAMB.

\begin{table}[htbp]
\centering
\caption{Confronto tra SGD con Momentum, Adam e LAMB.}
\small
\begin{tabular}{@{}p{3cm}p{3cm}p{3cm}p{3.5cm}@{}}
\toprule
\textbf{Propriet\`a} & \textbf{SGD+Mom.} & \textbf{Adam} & \textbf{LAMB} \\
\midrule
LR adattivo param. & No & S\`i & S\`i \\
LR adattivo layer & No & No & S\`i (trust ratio) \\
Batch size & 32--256 & 64--512 & Fino a 32K \\
Memoria extra & $1\!\times\!|\theta|$ & $2\!\times\!|\theta|$ & $2\!\times\!|\theta|$ \\
Iperparametri & $\eta, \beta$ & $\eta, \beta_1, \beta_2$ & $\eta, \beta_1, \beta_2, \lambda$ \\
Nel Perceiver & No & No & \textbf{S\`i} \\
\bottomrule
\end{tabular}
\end{table}


\newpage
%% ============================================================
%%  SEZIONE 9 --- ANALISI DELLA COMPLESSITA
%% ============================================================
\section{Analisi della Complessit\`a Computazionale} \label{sec:complessita}

Quanto costa computazionalmente il Perceiver rispetto alle alternative?

\subsection{Complessit\`a del Transformer standard} \label{sec:compl_transf}

Costo del Transformer standard:

\begin{formulabox}[Complessit\`a del Transformer standard]
\[
    \mathcal{O}_{\text{Transformer}} = L_{\text{tot}} \cdot \left[\underbrace{M^2 \cdot d}_{\text{self-attention}} + \underbrace{M \cdot d^2}_{\text{FFN}}\right]
\]
Per $M \gg d$, il termine dominante \`e $\mathcal{O}(L_{\text{tot}} \cdot M^2 \cdot d)$.
\end{formulabox}

\subsection{Complessit\`a del Perceiver} \label{sec:compl_perceiver}

Vediamo ora cosa succede con il Perceiver. La differenza chiave \`e che il termine quadratico $M^2$ viene sostituito dal termine $MN$, dove $N = 512 \ll M = 50.176$.

\begin{formulabox}[Complessit\`a del Perceiver]
\[
    \mathcal{O}_{\text{Perceiver}} = T \cdot \bigg[\underbrace{M \cdot N \cdot d_k}_{\text{cross-attention}} + \underbrace{L \cdot N^2 \cdot D}_{\text{latent SA}} + \underbrace{L \cdot N \cdot D^2}_{\text{latent FFN}}\bigg]
\]
Poich\'e $N$, $D$, $L$, $T$ e $d_k$ sono tutti indipendenti da $M$, la complessit\`a rispetto a $M$ \`e \textbf{lineare}: $\mathcal{O}(M)$.
\end{formulabox}

\subsection{Confronto tabulare} \label{sec:confronto_compl}

Confronto numerico concreto:

\begin{table}[htbp]
\centering
\caption{Confronto della complessit\`a computazionale.}
\begin{tabular}{@{}lccc@{}}
\toprule
\textbf{Architettura} & \textbf{Complessit\`a in $M$} & \textbf{ImageNet ($M$)} & \textbf{Praticabile?} \\
\midrule
Transformer (pixel) & $\mathcal{O}(M^2)$ & $50{,}176^2 \approx 2.5 \times 10^9$ & No \\
ViT-B/16 (patch) & $\mathcal{O}(M_p^2)$ & $196^2 = 38{,}416$ & S\`i \\
Perceiver & $\mathcal{O}(MN)$ & $50{,}176 \times 512 \approx 2.6 \times 10^7$ & S\`i \\
CNN (ResNet) & $\mathcal{O}(M \cdot k^2)$ & $50{,}176 \times 9 \approx 4.5 \times 10^5$ & S\`i \\
\bottomrule
\end{tabular}
\end{table}

\begin{notabox}[Interpretazione della tabella di complessit\`a]
\textbf{Transformer su pixel}: con $M = 50{,}176$, la matrice di attenzione richiede $2.5 \times 10^9$ float32 $= 10$ GB solo per un singolo layer. Con 12 layer e un batch di 64 immagini, servirebbero circa 7.7 TB di memoria GPU. \textbf{Completamente impraticabile} con la tecnologia attuale.

\textbf{ViT-B/16}: riducendo l'immagine a 196 patch, la matrice di attenzione ha solo $38{,}416$ elementi ($\sim$0.15 MB). Molto praticabile, ma richiede la suddivisione in patch (inductive bias).

\textbf{Perceiver}: la matrice di cross-attention ha $\sim$25.7M elementi ($\sim$103 MB). Pi\`u grande del ViT, ma comunque gestibile. E soprattutto, opera direttamente sui pixel senza patch.

\textbf{CNN (ResNet)}: la complessit\`a per pixel \`e solo $k^2 = 9$ (kernel $3 \times 3$), ma non cattura dipendenze globali. Per collegare pixel distanti, servono molti layer (50--152) con pooling intermedi.
\end{notabox}

\subsection{Scalabilit\`a rispetto alla dimensione dell'input} \label{sec:scalabilita}

Un vantaggio cruciale del Perceiver \`e che la complessit\`a cresce \textbf{linearmente} con $M$. In pratica, quello che succede \`e che raddoppiare la risoluzione dell'immagine raddoppia (circa) il costo computazionale, anzich\'e quadruplicarlo come nel Transformer standard. Questa differenza \`e enorme per applicazioni come il video ad alta risoluzione.

\begin{notabox}[In parole semplici: perch\'e la scalabilit\`a lineare \`e cos\`i importante]
Immaginiamo di gestire una mailing list. Con il Transformer, ogni persona deve parlare con ogni altra persona (complessit\`a $M^2$): se la lista raddoppia da 100 a 200 persone, le conversazioni passano da 10.000 a 40.000. Con il Perceiver, ogni persona parla solo con 512 ``delegati'' (complessit\`a $MN$): se la lista raddoppia, le conversazioni passano da 51.200 a 102.400 --- esattamente il doppio, non il quadruplo. Per un video con 1.5 milioni di pixel-frame, questa differenza \`e tra ``impossibile'' e ``fattibile''.
\end{notabox}

\begin{examplebox}[Scalabilit\`a: cosa succede con input pi\`u grandi?]
\begin{center}
\begin{tabular}{@{}lccc@{}}
\toprule
\textbf{Input} & $M$ & \textbf{Transf. $\mathcal{O}(M^2)$} & \textbf{Perceiver $\mathcal{O}(MN)$} \\
\midrule
Immagine $64 \times 64$ & 4{,}096 & $1.7 \times 10^7$ & $2.1 \times 10^6$ \\
Immagine $224 \times 224$ & 50{,}176 & $2.5 \times 10^9$ & $2.6 \times 10^7$ \\
Immagine $512 \times 512$ & 262{,}144 & $6.9 \times 10^{10}$ & $1.3 \times 10^8$ \\
Video 1s $224 \times 224$ & 1{,}505{,}280 & $2.3 \times 10^{12}$ & $7.7 \times 10^8$ \\
\bottomrule
\end{tabular}
\end{center}

Passando da $64 \times 64$ a $224 \times 224$ ($12\times$ pi\`u pixel), il Transformer diventa $150\times$ pi\`u costoso, mentre il Perceiver diventa solo $12\times$ pi\`u costoso. Per il video ($1.5$M pixel), il Transformer \`e completamente proibitivo ($2.3 \times 10^{12}$), mentre il Perceiver resta trattabile ($7.7 \times 10^8$).
\end{examplebox}

Il ViT ``bara'' riducendo la risoluzione a livello di patch, mentre il Perceiver opera \textbf{direttamente sui pixel} pur mantenendo una complessit\`a trattabile.

\begin{notabox}[In parole semplici: il confronto di complessit\`a]
Facciamo un'analogia con le telefonate. In una conferenza telefonica con il Transformer, tutti i $50.176$ partecipanti devono parlare con tutti gli altri --- $50.176^2 = 2.5$ miliardi di chiamate! Con il Perceiver, si nominano 512 ``moderatori'' (i latenti): ogni partecipante chiama solo i 512 moderatori ($50.176 \times 512 = 25.7$ milioni di chiamate), e i moderatori parlano tra loro ($512^2 = 262.144$ chiamate). Il costo totale \`e $\sim$100 volte minore, e soprattutto scala \textbf{linearmente} col numero di partecipanti.
\end{notabox}

\begin{examplebox}[Esempio numerico dettagliato: complessit\`a per ImageNet]
\textbf{Transformer standard (se operasse su pixel)}:
\begin{align*}
    \text{Self-attention} &= M^2 \cdot d = 50{,}176^2 \times 512 = 1.29 \times 10^{12} \text{ operazioni} \\
    \text{Memoria matrice att.} &= M^2 = 50{,}176^2 \approx 2.52 \times 10^{9} \text{ float32}\\
    &= \textbf{10.1 GB per singolo layer!}
\end{align*}
Con 12 layer e batch size 64, servirebbero $\sim$7.7 TB solo per le matrici di attenzione. \textbf{Impossibile.}

\textbf{Perceiver} ($N = 512$, $D = 1024$, $d_k = 1024$):
\begin{align*}
    \text{Cross-attention} &= M \cdot N \cdot d_k = 50{,}176 \times 512 \times 1024 = 2.63 \times 10^{10} \\
    \text{Latent SA (1 layer)} &= N^2 \cdot D = 512^2 \times 1024 = 2.68 \times 10^{8} \\
    \text{Latent SA ($L = 6$)} &= 6 \times 2.68 \times 10^{8} = 1.61 \times 10^{9} \\
    \text{Totale 1 iterazione} &\approx 2.63 \times 10^{10} + 1.61 \times 10^{9} \approx 2.79 \times 10^{10} \\
    \text{Totale $T = 8$} &\approx 8 \times 2.79 \times 10^{10} = 2.23 \times 10^{11}
\end{align*}

\textbf{Rapporto}: Transformer / Perceiver $\approx 1.29 \times 10^{12} / 2.23 \times 10^{11} \approx 5.8\times$ pi\`u costoso, e questo solo per \textbf{un layer} del Transformer vs.\ \textbf{tutte le 8 iterazioni} del Perceiver!

\textbf{Memoria matrice attenzione}: CA = $N \times M = 512 \times 50{,}176 \approx 25.7$M float32 $= 103$ MB. Totalmente gestibile su una GPU moderna.
\end{examplebox}

\subsection{Analisi del numero di parametri} \label{sec:num_parametri}

Vediamo ora un conteggio dettagliato di tutti i parametri del Perceiver. Questo tipo di analisi \`e molto utile sia per l'esame (mostra comprensione profonda dell'architettura) sia nella pratica (aiuta a stimare i costi di training e deploy).

\begin{examplebox}[Parametri del Perceiver su ImageNet --- conteggio dettagliato]

\textbf{Con weight sharing} ($T = 8$, $L = 6$, $N = 512$, $D = 1024$, $D_{\text{ff}} = 4 \times 1024 = 4096$):

\textbf{1. Cross-attention (prima, pesi unici)}:
\begin{itemize}
    \item $\mathbf{W}_Q \in \mathbb{R}^{1024 \times 1024}$: $1{,}048{,}576$ parametri
    \item $\mathbf{W}_K \in \mathbb{R}^{512 \times 1024}$: $524{,}288$
    \item $\mathbf{W}_V \in \mathbb{R}^{512 \times 1024}$: $524{,}288$
    \item $\mathbf{W}_O \in \mathbb{R}^{1024 \times 1024}$: $1{,}048{,}576$
    \item MLP post-CA: $1024 \times 4096 + 4096 \times 1024 = 8{,}388{,}608$
    \item LayerNorm ($\times 3$): $3 \times 2 \times 1024 = 6{,}144$
    \item Bias vari: $\sim 10{,}000$
    \item \textbf{Subtotale CA$_1$: $\sim$11.6M}
\end{itemize}

\textbf{2. Cross-attention (condivisa, stesso conteggio)}: $\sim$11.6M

\textbf{3. Latent Transformer (1 torre di $L = 6$ blocchi, condivisa)}:
\begin{itemize}
    \item Per blocco: $\mathbf{W}_Q, \mathbf{W}_K, \mathbf{W}_V, \mathbf{W}_O$ ($\times 1024^2$ ciascuno) + MLP ($2 \times 1024 \times 4096$) + LN ($\times 2$)
    \item Per blocco: $\approx 4 \times 1{,}048{,}576 + 8{,}388{,}608 + 4{,}096 \approx 12.6$M
    \item Per torre ($L = 6$): $6 \times 12.6$M $\approx$ 75.5M ... ma con weight sharing solo 1 torre!
    \item \textbf{Subtotale LT: $\sim$17M} (nota: il paper usa $D_{\text{latent}} < D$ in alcune configurazioni)
\end{itemize}

\textbf{4. Embedding input}: $261 \times 512 = 133{,}632$ $\approx 0.1$M

\textbf{5. Latent array}: $512 \times 1024 = 524{,}288$ $\approx 0.5$M

\textbf{6. Classificatore}: $1024 \times 1000 + 1000 = 1{,}025{,}000$ $\approx 1.0$M

\textbf{Totale con WS}: $\sim$45M parametri.
\textbf{Senza WS}: ogni iterazione ha CA e LT separati $\to$ $\sim$326M parametri.

Per confronto: ResNet-50 $\sim$25M, ViT-B/16 $\sim$86M.
\end{examplebox}


\newpage
%% ============================================================
%%  SEZIONE 10 --- ESPERIMENTI
%% ============================================================
\section{Esperimenti} \label{sec:esperimenti}

Ora arriva il momento della verit\`a: funziona davvero? Dopo tutta la teoria, vediamo come il Perceiver se la cava nella pratica, confrontandolo con le architetture specializzate su diversi benchmark.

\subsection{ImageNet} \label{sec:imagenet}

\begin{table}[htbp]
\centering
\caption{Risultati su ImageNet (classificazione top-1).}
\begin{tabular}{@{}lccc@{}}
\toprule
\textbf{Modello} & \textbf{Input} & \textbf{Top-1 (\%)} & \textbf{Parametri} \\
\midrule
ResNet-50 & $224 \times 224$ pixel (conv) & 77.6 & $\sim$25M \\
ViT-B/16 & $14 \times 14$ patch & 77.9 & $\sim$86M \\
ResNet-50 (Fourier features) & pixel + FF & 73.5 & --- \\
ViT-B/16 (Fourier features) & pixel + FF & 76.7 & --- \\
Transformer (64$\times$64, FF) & pixel + FF & 57.0 & --- \\
\textbf{Perceiver (Fourier)} & $224 \times 224$ pixel & \textbf{78.0} & $\sim$45M \\
Perceiver (learned PE) & $224 \times 224$ pixel & 70.9 & $\sim$45M \\
\bottomrule
\end{tabular}
\end{table}

Il Perceiver raggiunge 78.0\% operando \textbf{direttamente sui pixel}, senza patch o convoluzioni, con la stessa architettura usata su audio, video e point cloud. Questo risultato \`e notevole: il Perceiver supera sia ResNet-50 che ViT-B/16, pur essendo un'architettura ``generalista'' senza alcun inductive bias specifico per le immagini.

\begin{notabox}[In parole semplici: i risultati su ImageNet]
Detto in parole povere: il Perceiver \`e come uno studente che non ha mai studiato specificatamente per l'esame di matematica, n\'e per quello di fisica, n\'e per quello di chimica --- ma che ha sviluppato un metodo di studio cos\`i efficace da prendere un buon voto in tutti e tre. Non il voto pi\`u alto in assoluto (uno studente che ha studiato solo matematica potrebbe fare leggermente meglio in matematica), ma un voto impressionante considerando che usa lo stesso identico metodo per tutto.
\end{notabox}

\begin{notabox}[Analisi dei risultati ImageNet]
\textbf{ResNet-50 (77.6\%)}: architettura specializzata per immagini, sfrutta convoluzioni locali, invarianza traslazionale e gerarchia di feature. Ha solo $\sim$25M parametri.

\textbf{ViT-B/16 (77.9\%)}: opera su patch $16 \times 16$, riducendo la sequenza a 196 token. Raggiunge prestazioni simili a ResNet con pi\`u parametri ($\sim$86M).

\textbf{ResNet-50 con FF (73.5\%)}: ResNet modificata per operare su pixel con Fourier features. Perde 4 punti perch\'e i filtri $3 \times 3$ lavorano su pixel adiacenti che, dopo l'appiattimento, non sono pi\`u spazialmente coerenti nel modo atteso dalla CNN.

\textbf{Transformer 64$\times$64 FF (57.0\%)}: Transformer diretto su pixel, ma solo su immagini ridotte a $64 \times 64$ (per evitare $\mathcal{O}(M^2)$ proibitivo). La risoluzione ridotta limita le prestazioni.

\textbf{Perceiver FF (78.0\%)}: la migliore architettura generalista. Opera direttamente su tutti i $50{,}176$ pixel a piena risoluzione. Le Fourier features forniscono l'informazione posizionale.

\textbf{Perceiver learned PE (70.9\%)}: la versione con positional encoding appreso perde 7 punti. Questo dimostra l'importanza cruciale delle Fourier features: apprendere $50{,}176$ embedding posizionali distinti \`e difficile.
\end{notabox}

\begin{figure}[H]
\centering
\includegraphics[width=0.7\textwidth]{paper_figures/table1_imagenet.png}
\caption{Table 1 dal paper originale: accuratezza top-1 su ImageNet. Il Perceiver con Fourier features (78.0\%) \`e competitivo con ResNet-50 (77.6\%) e ViT-B-16 (77.9\%), pur operando direttamente sui pixel senza convoluzioni o patch.}
\label{fig:paper_table1}
\end{figure}

La Figura~\ref{fig:paper_table1} riporta la tabella originale dal paper, che include anche le configurazioni con diversi valori di cross-attend e con augmentation. Si noti come il Perceiver con 8 cross-attend e Fourier features raggiunga il risultato migliore (78.0\%), superando modelli che sfruttano inductive bias specifici per le immagini. Questo conferma che il latent bottleneck, combinato con le Fourier features, \`e sufficiente per catturare informazione visiva a livello competitivo.


\subsection{ImageNet Permutato: l'esperimento chiave} \label{sec:imagenet_perm}

Questo \`e forse l'esperimento pi\`u elegante e convincente del paper. La domanda che ci si pone \`e: cosa succede se mescoliamo casualmente tutti i pixel di ogni immagine? Se un'architettura si basa sulla struttura spaziale (come le CNN), le sue prestazioni dovrebbero crollare. Se \`e veramente ``agnostica alla struttura'', dovrebbe funzionare esattamente come prima.

Si applica la \emph{stessa permutazione casuale} a tutti i pixel di tutte le immagini del dataset.

\begin{table}[htbp]
\centering
\caption{Risultati su ImageNet permutato.}
\begin{tabular}{@{}lccc@{}}
\toprule
\textbf{Modello} & \textbf{Top-1 originale (\%)} & \textbf{Top-1 permutato (\%)} & \textbf{$\Delta$} \\
\midrule
ResNet-50 (FF) & 73.5 & 39.4 & $-34.1$ \\
ViT-B/16 (FF) & 76.7 & 61.7 & $-15.0$ \\
\textbf{Perceiver (Fourier)} & \textbf{78.0} & \textbf{78.0} & $\mathbf{0.0}$ \\
Perceiver (learned PE) & 70.9 & 70.9 & $0.0$ \\
\bottomrule
\end{tabular}
\end{table}

\begin{proofbox}[Perch\'e il Perceiver \`e invariante alla permutazione]
\begin{itemize}
    \item \textbf{ResNet-50 ($-34$ punti)}: le convoluzioni con kernel $3 \times 3$ si basano sulla \emph{localit\`a spaziale}. Il kernel $3 \times 3$ calcola il prodotto scalare con 9 pixel adiacenti. Dopo la permutazione, quei 9 pixel ``adiacenti'' nella griglia non sono pi\`u spazialmente coerenti: il kernel trova rumore invece di bordi e texture. Le feature map dei primi layer diventano essenzialmente casuali, e la CNN non riesce pi\`u ad estrarre informazione utile.
    \item \textbf{ViT ($-15$ punti)}: il ViT suddivide l'immagine in patch $16 \times 16$. Dopo la permutazione, ogni patch contiene 256 pixel provenienti da posizioni casuali dell'immagine originale: la patch non rappresenta pi\`u una regione spaziale coerente. Il positional embedding (che \`e \emph{sommato}, non concatenato) fornisce ancora informazione sulla posizione della patch, ma la proiezione lineare della patch produce un embedding quasi privo di informazione visiva.
    \item \textbf{Perceiver (0 punti)}: opera su \textbf{pixel individuali}, non su regioni. Le Fourier features sono \emph{concatenate} (non sommate) a ogni pixel: ogni pixel ``porta con s\'e'' la propria coordinata spaziale come parte del suo vettore di feature. La cross-attention \`e permutation-invariant:
    \[
        \sum_{j=1}^{M} A_{ij} \mathbf{v}_j = \sum_{j=1}^{M} A_{i,\pi(j)} \mathbf{v}_{\pi(j)}
    \]
    perch\'e \`e una somma su tutti gli elementi, indipendentemente dal loro ordine.
\end{itemize}
\end{proofbox}

% TikZ: Permuted ImageNet visualization
\begin{figure}[H]
\centering
\begin{tikzpicture}[
    arrowstyle/.style={-{Stealth[length=2.5mm]}, thick},
    >=Stealth,
    scale=0.85, transform shape
]

% Original image (stylized grid)
\node[font=\small\bfseries] at (0, 3) {Immagine originale};
\draw[thick] (-1.2, -0.5) rectangle (1.2, 2.5);
% Colored grid cells to suggest structure
\fill[red!30] (-1.2, 1.5) rectangle (-0.4, 2.5);
\fill[red!20] (-0.4, 1.5) rectangle (0.4, 2.5);
\fill[blue!30] (-1.2, 0.5) rectangle (-0.4, 1.5);
\fill[green!30] (0.4, 0.5) rectangle (1.2, 1.5);
\fill[yellow!30] (-1.2, -0.5) rectangle (-0.4, 0.5);
\fill[orange!30] (0.4, -0.5) rectangle (1.2, 0.5);
\draw[gray, thin] (-0.4, -0.5) -- (-0.4, 2.5);
\draw[gray, thin] (0.4, -0.5) -- (0.4, 2.5);
\draw[gray, thin] (-1.2, 0.5) -- (1.2, 0.5);
\draw[gray, thin] (-1.2, 1.5) -- (1.2, 1.5);

% Arrow
\draw[arrowstyle, very thick] (1.8, 1) -- (3.2, 1) node[midway, above, font=\scriptsize] {Permutazione};
\node[font=\scriptsize] at (2.5, 0.5) {$\pi$ fissa};

% Permuted image (shuffled colors)
\node[font=\small\bfseries] at (5, 3) {Immagine permutata};
\draw[thick] (3.8, -0.5) rectangle (6.2, 2.5);
\fill[green!30] (3.8, 1.5) rectangle (4.6, 2.5);
\fill[yellow!30] (4.6, 1.5) rectangle (5.4, 2.5);
\fill[orange!30] (5.4, 1.5) rectangle (6.2, 2.5);
\fill[red!30] (3.8, 0.5) rectangle (4.6, 1.5);
\fill[blue!30] (5.4, 0.5) rectangle (6.2, 1.5);
\fill[red!20] (3.8, -0.5) rectangle (4.6, 0.5);
\draw[gray, thin] (4.6, -0.5) -- (4.6, 2.5);
\draw[gray, thin] (5.4, -0.5) -- (5.4, 2.5);
\draw[gray, thin] (3.8, 0.5) -- (6.2, 0.5);
\draw[gray, thin] (3.8, 1.5) -- (6.2, 1.5);

% Results bar chart
\node[font=\small\bfseries] at (10.5, 3) {$\Delta$ accuracy};

% ResNet bar
\fill[red!50] (8.5, -0.5) rectangle (9.2, 2.2);
\node[font=\scriptsize, rotate=90] at (8.85, 0.8) {ResNet $-34$};

% ViT bar
\fill[orange!50] (9.5, -0.5) rectangle (10.2, 1.3);
\node[font=\scriptsize, rotate=90] at (9.85, 0.3) {ViT $-15$};

% Perceiver bar
\fill[green!50] (10.5, -0.5) rectangle (11.2, -0.5);
\node[font=\scriptsize] at (10.85, -0.8) {Perceiver};
\node[font=\scriptsize\bfseries, green!60!black] at (10.85, -0.2) {$\Delta = 0$};

% Zero line
\draw[dashed, gray] (8.2, -0.5) -- (11.5, -0.5);
\node[font=\tiny, gray] at (11.8, -0.5) {0};

\end{tikzpicture}
\caption{Esperimento di ImageNet permutato. La stessa permutazione casuale dei pixel viene applicata a tutte le immagini. ResNet perde 34 punti, ViT perde 15 punti, il Perceiver perde 0 punti (completamente invariante alla permutazione).}
\label{fig:permuted_imagenet}
\end{figure}

\begin{figure}[H]
\centering
\includegraphics[width=0.7\textwidth]{paper_figures/table2_permuted.png}
\caption{Table 2 dal paper originale: confronto su ImageNet standard e permutato. Il Perceiver con Fourier features mantiene esattamente la stessa accuratezza (78.0\%) su entrambe le versioni, dimostrando la totale invarianza alla permutazione. ResNet-50 crolla da 73.5\% a 39.4\%, e ViT da 76.7\% a 61.7\%.}
\label{fig:paper_table2}
\end{figure}

La Figura~\ref{fig:paper_table2} mostra la tabella originale del paper che confronta direttamente le prestazioni su ImageNet standard e permutato. I numeri parlano da soli: il Perceiver con Fourier features mantiene \textbf{esattamente} la stessa accuratezza ($78.0\%$) in entrambe le condizioni, mentre le architetture che si basano sulla struttura spaziale subiscono cali drammatici. Questo risultato \`e forse la dimostrazione pi\`u elegante dell'assenza di inductive bias strutturali nel Perceiver: il modello non ``sa'' che i pixel formano una griglia 2D, e quindi non viene influenzato quando questa griglia viene distrutta dalla permutazione.


\subsection{AudioSet: classificazione audio} \label{sec:audioset}

Ora viene la parte pi\`u impressionante: il Perceiver usa la \textbf{stessa identica architettura} anche per l'audio. Non cambia nulla tranne il formato dell'input. Vediamo come.

\begin{formulabox}[AudioSet: configurazione dell'input]
\textbf{Raw waveform}: 48kHz $\times$ 1.28s = 61{,}440 campioni, segmentati in 480 vettori di 128. Input: $M = 480$, $C = 128$.

\textbf{Mel spectrogram}: $M = 4{,}800$ input.

L'architettura \`e \textbf{identica} a ImageNet: cambia solo $M$ e $C$.
\end{formulabox}

\begin{table}[htbp]
\centering
\caption{Perceiver su AudioSet (mAP).}
\begin{tabular}{@{}lc@{}}
\toprule
\textbf{Configurazione} & \textbf{mAP} \\
\midrule
Perceiver (raw audio) & 38.4 \\
Perceiver (mel spectrogram) & 38.4 \\
Perceiver (video only) & 25.8 \\
Perceiver (audio + video) & 43.5 \\
\bottomrule
\end{tabular}
\end{table}

La fusione multi-modale si realizza semplicemente \textbf{concatenando} i byte array delle diverse modalit\`a: $\mathbf{x}_{\text{multi}} = [\mathbf{x}_{\text{audio}}; \mathbf{x}_{\text{video}}]$.

\begin{notabox}[In parole semplici: la fusione multimodale nel Perceiver]
Detto in parole povere, la fusione multimodale nel Perceiver \`e come mescolare due mazzi di carte. Il mazzo ``audio'' ha 480 carte e il mazzo ``video'' ha 3136 carte. Li mettiamo insieme in un unico mazzo di 3616 carte e li diamo al Perceiver. I latenti ``pescano'' carte dal mazzo unificato, scegliendo liberamente se pescare carte audio o video in base a ci\`o che gli serve. Un latente potrebbe specializzarsi per ``applauso + gesto delle mani'' (correlazione cross-modale che n\'e l'audio n\'e il video da soli contengono). Nessun modulo di fusione speciale --- solo concatenazione e attenzione!
\end{notabox}

\begin{examplebox}[Walkthrough: da un file WAV al Perceiver]
\textbf{Passo 1}: si carica un file audio WAV a 48kHz, durata 1.28 secondi. Campioni totali: $48{,}000 \times 1.28 = 61{,}440$.

\textbf{Passo 2}: si segmenta in finestre di 128 campioni ciascuna (senza sovrapposizione): $61{,}440 / 128 = 480$ segmenti.

\textbf{Passo 3}: l'input al Perceiver \`e $\mathbf{x}_{\text{audio}} \in \mathbb{R}^{480 \times 128}$, dove $M = 480$ e $C = 128$.

\textbf{Passo 4}: si aggiungono le Fourier features 1D ($P = 1$, $K = 64$): $d_{\text{PE}} = 1 \times 129 = 129$, $C_{\text{tot}} = 128 + 129 = 257$.

\textbf{Passo 5}: il Perceiver processa $\mathbb{R}^{480 \times 257}$ con la stessa architettura di ImageNet (cambia solo $M$ e $C$).
\end{examplebox}

\begin{notabox}[Fusione audio-video multimodale]
La fusione multimodale nel Perceiver \`e elegantemente semplice: si concatenano i byte array delle diverse modalit\`a lungo la dimensione $M$:
\[
    \mathbf{x}_{\text{multi}} = \begin{pmatrix} \mathbf{x}_{\text{audio}} \\ \mathbf{x}_{\text{video}} \end{pmatrix} \in \mathbb{R}^{(M_{\text{audio}} + M_{\text{video}}) \times C_{\text{max}}}
\]
Le Fourier features distinguono la modalit\`a di provenienza tramite frequenze diverse. La cross-attention permette ai latenti di ``leggere'' simultaneamente da audio e video, apprendendo automaticamente le correlazioni cross-modali. Risultato: 43.5 mAP (vs.\ 38.4 solo audio e 25.8 solo video), dimostrando che la fusione cattura informazione complementare.

\textbf{Dettaglio implementativo}: quando le due modalit\`a hanno dimensioni di canale diverse ($C_{\text{audio}} \neq C_{\text{video}}$), si esegue il padding a zero della modalit\`a con meno canali, oppure si proiettano entrambe le modalit\`a alla stessa dimensione prima della concatenazione. Le Fourier features includono una dimensione aggiuntiva che identifica la modalit\`a (es.\ una coordinata binaria ``audio=0, video=1''), permettendo al modello di distinguere gli elementi delle due sorgenti.
\end{notabox}

\begin{examplebox}[Esempio: fusione audio-video in pratica]
\textbf{Audio}: 480 segmenti $\times$ 257 feature (128 campioni + 129 Fourier) = $\mathbb{R}^{480 \times 257}$.

\textbf{Video}: supponiamo 16 frame $\times$ 196 patch = 3136 token $\times$ 257 feature = $\mathbb{R}^{3136 \times 257}$.

\textbf{Concatenazione}: $\mathbb{R}^{(480 + 3136) \times 257} = \mathbb{R}^{3616 \times 257}$.

I latenti ($512 \times 1024$) ``leggono'' simultaneamente audio e video tramite cross-attention:
\[
    \mathbf{A} \in \mathbb{R}^{512 \times 3616}
\]
Ogni latente pu\`o prestare attenzione sia a token audio che a token video. Un latente potrebbe specializzarsi per ``suono di applauso + gesto delle mani'' (correlazione cross-modale), catturando informazione che n\'e l'audio n\'e il video da soli contengono.
\end{examplebox}

\begin{examplebox}[Esempio numerico: preparazione dell'input AudioSet]
Vediamo concretamente come un segnale audio viene preparato per il Perceiver.

\textbf{Passo 1 --- Acquisizione}: file WAV mono a 48 kHz, durata 1.28 secondi.

Campioni totali: $48{,}000 \times 1.28 = 61{,}440$ valori float (ampiezza del segnale tra $-1$ e $+1$).

\textbf{Passo 2 --- Segmentazione in blocchi}:

Dividiamo i $61{,}440$ campioni in blocchi contigui di 128 campioni ciascuno (senza sovrapposizione):
\[
\frac{61{,}440}{128} = 480 \text{ blocchi}
\]

Ogni blocco contiene 128 valori float consecutivi, ad esempio il blocco 0:
\[
\text{blocco}_0 = [0.012, -0.034, 0.051, 0.008, -0.022, \ldots, 0.041] \in \mathbb{R}^{128}
\]

Ogni blocco rappresenta $128 / 48{,}000 \approx 2.67$ millisecondi di audio.

\textbf{Passo 3 --- Input grezzo}: $\mathbf{x}_{\text{audio}} \in \mathbb{R}^{480 \times 128}$, dove $M = 480$ (token temporali) e $C = 128$ (campioni per token).

\textbf{Passo 4 --- Fourier features 1D} (posizione temporale):
\begin{itemize}
\item $P = 1$ dimensione (il tempo)
\item $K = 64$ bande di frequenza
\item Per ogni blocco $i$, la posizione normalizzata \`e $t_i = \frac{2i}{479} - 1 \in [-1, 1]$
\item Fourier features: $2 \times 64 = 128$ valori sin/cos + 1 coordinata raw = 129 valori
\end{itemize}

\textbf{Passo 5 --- Concatenazione}:
\[
\tilde{\mathbf{x}}_{\text{audio}} \in \mathbb{R}^{480 \times (128 + 129)} = \mathbb{R}^{480 \times 257}
\]

\textbf{Confronto con ImageNet}:
\begin{center}
\begin{tabular}{@{}lccc@{}}
\toprule
& $M$ (token) & $C$ (canali) & $C_{\text{tot}}$ (con Fourier) \\
\midrule
ImageNet & 50{,}176 & 3 (RGB) & 261 \\
AudioSet (raw) & 480 & 128 (campioni) & 257 \\
\bottomrule
\end{tabular}
\end{center}

L'architettura del Perceiver \`e \textbf{identica} in entrambi i casi: cambia solo $M$ e $C$. I 480 token audio vengono trattati esattamente come i 50{,}176 token immagine: diventano chiavi e valori nella cross-attention, e i 512 latenti li ``leggono'' alla stessa maniera.
\end{examplebox}

\begin{figure}[H]
\centering
\includegraphics[width=0.85\textwidth]{paper_figures/table3_audioset_crop.png}
\caption{Table 3 dal paper: prestazioni del Perceiver su AudioSet (mAP). Il Perceiver con raw audio raggiunge 38.4 mAP, competitivo con modelli specializzati. La fusione audio+video raggiunge 43.5 mAP.}
\label{fig:paper_table3_audioset}
\end{figure}

Come si vede dalla Tabella~\ref{fig:paper_table3_audioset}, il Perceiver raggiunge prestazioni competitive su AudioSet anche senza alcuna modifica architetturale rispetto alla configurazione usata per ImageNet. Il risultato pi\`u notevole \`e che la fusione multimodale audio+video (43.5 mAP) supera significativamente entrambe le singole modalit\`a, confermando che la semplice concatenazione dei byte array permette ai latenti di catturare correlazioni cross-modali utili.


\subsection{ModelNet40: classificazione di point cloud 3D} \label{sec:modelnet}

E per le nuvole di punti 3D? Stesso discorso: il Perceiver usa la stessa architettura, cambiando solo l'input.

\begin{formulabox}[ModelNet40: configurazione]
$M \approx 2{,}048$ punti, $C = 3$ coordinate ($x, y, z$), 40 classi (sedia, tavolo, aereo, ecc.). Fourier features 3D con $P = 3$ dimensioni, $K = 64$ bande e max resolution $\mu = 1120$. $d_{\text{PE}} = 3 \times 129 = 387$, $C_{\text{tot}} = 3 + 387 = 390$.
\end{formulabox}

\begin{examplebox}[Walkthrough: da una point cloud 3D al Perceiver]
\textbf{Passo 1}: si carica la nuvola di punti 3D di un oggetto (es.\ una sedia): 2048 punti, ciascuno con coordinate $(x, y, z)$. Input grezzo: $\mathbb{R}^{2048 \times 3}$.

\textbf{Passo 2}: si normalizzano le coordinate in $[-1, 1]$ rispetto al bounding box dell'oggetto.

\textbf{Passo 3}: si calcolano le Fourier features 3D. Per ogni punto, 3 coordinate $\times$ 129 valori ciascuna = 387 valori. Concatenati con le 3 coordinate raw: 390 valori per punto.

\textbf{Passo 4}: l'input al Perceiver \`e $\mathbb{R}^{2048 \times 390}$. Architettura identica!

\textbf{Passo 5}: il classificatore ha 40 output (invece di 1000 per ImageNet).
\end{examplebox}

\begin{table}[htbp]
\centering
\caption{Risultati su ModelNet40.}
\begin{tabular}{@{}lcc@{}}
\toprule
\textbf{Modello} & \textbf{Accuratezza (\%)} & \textbf{Specializzato per 3D?} \\
\midrule
PointNet++ & 91.9 & S\`i \\
\textbf{Perceiver} & 85.7 & \textbf{No} \\
ViT-B/2 (FF) & 78.9 & No \\
ResNet-50 (FF) & 66.3 & No \\
\bottomrule
\end{tabular}
\end{table}

Il gap di $\sim$6 punti riflette il trade-off generalit\`a vs.\ prestazioni: il Perceiver non sfrutta la geometria locale, ma \`e la stessa architettura usata per immagini e audio.

\begin{notabox}[In parole semplici: il trade-off del Perceiver sulle point cloud]
Detto in parole povere: PointNet++ \`e come un esperto di scultura che ``tocca'' l'oggetto 3D, sentendo le curve e gli spigoli locali (raggruppamento gerarchico). Il Perceiver \`e come qualcuno che guarda l'oggetto da lontano --- vede la forma generale ma perde i dettagli fini. Il Perceiver ``compensa'' con pi\`u parametri e pi\`u iterazioni, ma l'esperto di scultura, con la sua conoscenza specializzata, resta pi\`u preciso. Il punto \`e che il Perceiver non ha bisogno di essere ri-progettato per passare dalle point cloud alle immagini o all'audio --- usa sempre lo stesso ``sguardo da lontano''.
\end{notabox}

\begin{examplebox}[Confronto diretto: Perceiver vs PointNet++ su ModelNet40]
\begin{center}
\begin{tabular}{@{}p{4cm}p{5cm}p{4.5cm}@{}}
\toprule
\textbf{Aspetto} & \textbf{PointNet++} & \textbf{Perceiver} \\
\midrule
Input & 2048 punti $(x,y,z)$ & 2048 punti $(x,y,z)$ + FF \\
Preprocessing & Farthest point sampling + KNN grouping & Flatten + Fourier features \\
Inductive bias & Localit\`a 3D, gerarchia & Nessuno \\
Architettura & Set Abstraction layers (multi-scala) & Cross-att + Latent Transformer \\
Campo recettivo & Locale $\to$ globale (gerarchico) & Globale fin dalla prima cross-att \\
Parametri & $\sim$1.5M & $\sim$45M \\
Accuratezza & 91.9\% & 85.7\% \\
Specializzato? & S\`i (solo point cloud) & No (stessa architettura di ImageNet) \\
\bottomrule
\end{tabular}
\end{center}
\end{examplebox}

\begin{notabox}[Perch\'e il Perceiver perde contro PointNet++]
PointNet++ sfrutta tre \emph{inductive bias} specifici per le point cloud 3D che il Perceiver non ha:
\begin{enumerate}
    \item \textbf{Raggruppamento gerarchico}: PointNet++ costruisce una gerarchia di regioni locali (set abstraction layers), catturando strutture a scale diverse.
    \item \textbf{Campionamento adattivo}: usa farthest point sampling per selezionare punti rappresentativi.
    \item \textbf{Feature locali}: estrae feature da vicinati locali prima di aggregarle globalmente.
\end{enumerate}
Il Perceiver tratta ogni punto come un elemento indipendente: non ``sa'' che i punti vicini nello spazio 3D sono correlati. Questa \`e una scelta deliberata: rinunciare agli inductive bias specifici per ottenere un'architettura universale.
\end{notabox}

\begin{figure}[H]
\centering
\includegraphics[width=0.5\textwidth]{paper_figures/table4_modelnet40_crop.png}
\caption{Table 4 dal paper: accuratezza su ModelNet40. Il Perceiver raggiunge 85.7\%, inferiore a PointNet++ (91.9\%) ma superiore a tutti i modelli generici (ViT, ResNet con Fourier features).}
\label{fig:paper_table4_modelnet}
\end{figure}

La Tabella~\ref{fig:paper_table4_modelnet} mostra il confronto diretto dal paper originale. Il Perceiver si colloca nettamente al di sopra degli altri modelli generici (ViT-B/2 con Fourier features al 78.9\%, ResNet-50 al 66.3\%), ma rimane a $-6.2$ punti percentuali da PointNet++, che sfrutta inductive bias specifici per le point cloud come il raggruppamento gerarchico e il campionamento adattivo. Questo gap quantifica il ``costo della generalit\`a'': rinunciare alla conoscenza specifica del dominio ha un prezzo misurabile, ma il vantaggio \`e che la stessa architettura funziona su qualsiasi modalit\`a.


\subsection{Studi di ablazione dal paper} \label{sec:ablation}

Gli studi di ablazione sono fondamentali per capire \emph{perch\'e} il Perceiver funziona: rimuovendo o modificando un componente alla volta, possiamo misurare il suo contributo. Vediamo cosa succede concretamente.

\begin{table}[htbp]
\centering
\caption{Ablazione: effetto del numero di cross-attend (Table 5 del paper).}
\begin{tabular}{@{}lcc@{}}
\toprule
\textbf{Numero cross-attend $T$} & \textbf{Top-1 (\%)} & \textbf{Nota} \\
\midrule
$T = 1$ & 72.7 & Una sola lettura dell'input \\
$T = 2$ & 76.0 & Significativo miglioramento \\
$T = 4$ & 77.5 & Rendimenti decrescenti \\
$T = 8$ & 78.0 & Configurazione standard \\
$T = 16$ & 77.8 & Leggero calo (overfitting?) \\
\bottomrule
\end{tabular}
\end{table}

\begin{table}[htbp]
\centering
\caption{Ablazione: effetto del weight sharing (Table 7 del paper).}
\begin{tabular}{@{}lccc@{}}
\toprule
\textbf{Configurazione} & \textbf{Top-1 (\%)} & \textbf{Parametri} & \textbf{Note} \\
\midrule
Senza WS & 72.9 & $\sim$326M & Overfitting severo \\
WS solo LT & 75.4 & $\sim$172M & Miglioramento parziale \\
WS completo & 78.0 & $\sim$45M & Migliore! \\
\bottomrule
\end{tabular}
\end{table}

\begin{notabox}[Interpretazione degli studi di ablazione]
\textbf{Numero di cross-attend}: il salto maggiore \`e tra $T=1$ e $T=2$ (+3.3 punti), suggerendo che ``rileggere'' l'input almeno una volta \`e fondamentale. Oltre $T=8$, i rendimenti sono decrescenti o negativi, probabilmente perch\'e le iterazioni aggiuntive non estraggono informazione nuova ma possono causare overfitting.

\textbf{Weight sharing}: controintuitivamente, meno parametri portano a migliori prestazioni. Questo \`e un chiaro segnale di \textbf{regolarizzazione}: senza WS, il modello ha $7\times$ pi\`u parametri e overfitta il training set. Con WS, i parametri condivisi sono ``costretti'' a generalizzare, funzionando bene su tutte le iterazioni.

\textbf{Mappe di attenzione} (Figure 3 del paper): le mappe di attenzione della prima cross-attention mostrano che diversi latenti si specializzano per regioni semantiche diverse dell'immagine. Alcuni latenti ``guardano'' lo sfondo, altri l'oggetto principale, altri ancora i bordi. Nelle iterazioni successive, le mappe diventano pi\`u raffinate e semanticamente coerenti.
\end{notabox}

\begin{figure}[H]
\centering
\includegraphics[width=0.5\textwidth]{paper_figures/table5_cross_attends.png}
\caption{Table 5 dal paper: effetto del numero di cross-attend senza latent transformer. Con solo cross-attention (senza self-attention sui latenti), il modello raggiunge solo 45.3\% con 8 stadi. Questo dimostra che il latent transformer \`e essenziale.}
\label{fig:paper_table5}
\end{figure}

La Tabella~\ref{fig:paper_table5} evidenzia un risultato fondamentale: la cross-attention da sola non \`e sufficiente. Senza il latent transformer, anche aumentando il numero di cross-attend a 8, le prestazioni rimangono molto basse (45.3\%). Questo dimostra che il ruolo del latent transformer --- elaborare e raffinare le rappresentazioni latenti tramite self-attention --- \`e essenziale per le prestazioni del Perceiver.

\begin{figure}[H]
\centering
\includegraphics[width=0.55\textwidth]{paper_figures/table6_arrangement.png}
\caption{Table 6 dal paper: effetto del numero e della disposizione dei cross-attend. ``Interleaved'' (distribuiti nella rete) vs ``at start'' (tutti all'inizio). 8 cross-attend interleaved raggiungono 78.0\%, il miglior risultato.}
\label{fig:paper_table6}
\end{figure}

La Tabella~\ref{fig:paper_table6} confronta due strategie di posizionamento dei moduli di cross-attention: ``interleaved'' (alternati con blocchi di self-attention latente lungo tutta la rete) e ``at start'' (tutti concentrati all'inizio della rete). La disposizione interleaved risulta sistematicamente superiore, perch\'e permette al modello di ``rileggere'' l'input dopo ogni fase di elaborazione latente, aggiornando progressivamente la propria rappresentazione.

\begin{figure}[H]
\centering
\includegraphics[width=0.9\textwidth]{paper_figures/fig5_hyperparameter_sweeps.png}
\caption{Figure 5 dal paper: effetto degli iperparametri su ImageNet (modello ridotto). Da sinistra: dimensione latente, dimensione canale latente, numero self-attend per blocco, numero cross-attend. Aumentare la dimensione e il numero di blocchi generalmente migliora le prestazioni.}
\label{fig:paper_fig5_sweeps}
\end{figure}

La Figura~\ref{fig:paper_fig5_sweeps} mostra gli sweep degli iperparametri condotti dagli autori su una versione ridotta del modello. Si osserva che aumentare la dimensione latente $N$ e la dimensione dei canali $D$ porta a miglioramenti consistenti, ma con rendimenti decrescenti. Il numero di self-attend per blocco ($L$) ha un effetto positivo fino a circa 6 layer, oltre i quali le prestazioni si stabilizzano. Questi risultati guidano la scelta degli iperparametri della configurazione finale del Perceiver.

\begin{figure}[H]
\centering
\includegraphics[width=0.5\textwidth]{paper_figures/table7_weight_sharing.png}
\caption{Table 7 dal paper: il weight sharing migliora la validazione (da 72.9\% a 78.0\%) riducendo l'overfitting (train da 87.7\% a 79.5\%) e i parametri da 326.2M a 44.9M. Questo risultato contrasta con i nostri esperimenti su CIFAR-10, dove il WS peggiora le prestazioni.}
\label{fig:paper_table7}
\end{figure}

La Tabella~\ref{fig:paper_table7} \`e particolarmente istruttiva: mostra che senza weight sharing il modello overfitta in modo severo (87.7\% training ma solo 72.9\% validazione, un gap di quasi 15 punti), mentre con weight sharing completo il gap si riduce a soli 1.5 punti (79.5\% training, 78.0\% validazione). Questo risultato dimostra che il weight sharing agisce come un potente regolarizzatore sul modello full-scale, anche se nei nostri esperimenti su CIFAR-10 (Sezione~\ref{sec:exp_ws}) abbiamo osservato l'effetto opposto, a causa della scala molto diversa del modello e del dataset.

\begin{figure}[H]
\centering
\includegraphics[width=0.95\textwidth]{paper_figures/fig3_attention_maps.png}
\caption{Attention maps dal paper originale (Figure 3) per un modello con 8 cross-attention su ImageNet. Riga 1: immagine originale e close-up. Righe 2-4: mappe di attenzione della prima (blu), seconda (arancione) e ottava (verde) cross-attention. I primi layer catturano strutture a bassa frequenza (il soggetto \`e visibile); i layer successivi mostrano pattern a ``tartan'' ad alta frequenza.}
\label{fig:paper_attention_maps}
\end{figure}

La Figura~\ref{fig:paper_attention_maps} rende visibile un fenomeno affascinante: il Perceiver sviluppa spontaneamente una gerarchia di rappresentazioni, dai pattern a bassa frequenza (strutture globali, sagoma del soggetto) nei primi layer di cross-attention, fino ai pattern ad alta frequenza (dettagli fini, texture a ``tartan'') nei layer successivi. Questo comportamento ricorda la gerarchia delle CNN, dove i primi layer catturano bordi e i layer profondi catturano feature semantiche --- ma nel Perceiver questa gerarchia \textbf{emerge dal training}, senza essere imposta dall'architettura. Nella Sezione~\ref{sec:exp_attention} confronteremo queste attention maps con quelle ottenute dai nostri esperimenti su CIFAR-10.


\newpage
%% ============================================================
%%  SEZIONE 11 --- I NOSTRI ESPERIMENTI
%% ============================================================
\section{I Nostri Esperimenti: Implementazione e Risultati} \label{sec:nostri_esperimenti}

In questa sezione documentiamo l'intero percorso sperimentale che abbiamo condotto per verificare, riprodurre e approfondire i risultati del paper originale. Abbiamo implementato da zero sia il Perceiver che il Perceiver IO in PyTorch, addestrando i modelli su quattro dataset diversi (CIFAR-10, ModelNet40, WikiText-103 e GLUE), conducendo un totale di 17 esperimenti distinti. Lo scopo non \`e stato semplicemente ``far girare il codice'', ma interrogare le scelte progettuali degli autori: il positional encoding \`e davvero cruciale? Il weight sharing aiuta sempre? Le Fourier features sono robuste alla permutazione?

I risultati che presentiamo contengono alcune conferme importanti dei risultati del paper, ma anche scoperte sorprendenti che sfidano alcune intuizioni consolidate.

%% ------------------------------------------------------------
\subsection{Panoramica del progetto} \label{sec:panoramica_progetto}

\subsubsection{Cosa abbiamo implementato}

Il progetto consiste in un'implementazione completa \emph{from scratch} in PyTorch di due architetture:
\begin{itemize}[leftmargin=*]
    \item \textbf{Perceiver} (Jaegle et al., 2021): il modello descritto in dettaglio nelle sezioni precedenti, con cross-attention iterativa, latent transformer e weight sharing.
    \item \textbf{Perceiver IO} (Jaegle et al., 2022): l'estensione del Perceiver che aggiunge un \emph{decoder} con cross-attention inversa, permettendo output strutturati di dimensione arbitraria (e non solo classificazione con pooling).
\end{itemize}

\subsubsection{Struttura del codice}

\begin{formulabox}[Organizzazione del codice sorgente]
\begin{verbatim}
src/
  perceiver/
    perceiver.py     # Perceiver per classificazione
    encoder.py       # PerceiverEncoder: cross-att + transformer
    attention.py     # MultiHeadAttention, Cross/SelfAttention
  perceiver_io/
    perceiver_io.py  # PerceiverIO con decoder
  data/
    cifar10.py       # DataModule CIFAR-10 con Fourier PE
    modelnet40.py    # DataModule ModelNet40 (point cloud 3D)
    glue_tasks.py    # DataModule GLUE (NLP)
train.py             # Loop di addestramento principale
reproduce.py         # Script per riprodurre tutti gli esperimenti
\end{verbatim}
\end{formulabox}

\subsubsection{Hardware e software}

Tutti gli esperimenti sono stati condotti su una singola GPU consumer (NVIDIA RTX) con PyTorch 2.x. I tempi di addestramento variano da circa 3 ore (esperimenti CIFAR-10 semplici) a oltre 7 ore (esperimenti con permutazione). L'intero ciclo sperimentale ha richiesto circa 40 ore di GPU.

\subsubsection{Dataset utilizzati}

\begin{table}[htbp]
\centering
\caption{Dataset utilizzati nei nostri esperimenti.}
\label{tab:dataset_nostri}
\begin{tabular}{@{}llccc@{}}
\toprule
\textbf{Dataset} & \textbf{Dominio} & \textbf{Train} & \textbf{Test/Val} & \textbf{Classi/Task} \\
\midrule
CIFAR-10 & Immagini $32 \times 32$ & 50\,000 & 10\,000 & 10 classi \\
ModelNet40 & Point cloud 3D & 9\,843 & 2\,468 & 40 classi \\
WikiText-103 & Testo (MLM) & 103M token & --- & Vocabolario \\
WikiText-2 & Testo (MLM) & 2M token & --- & Vocabolario \\
GLUE & NLP (8 task) & Variabile & Variabile & Classificazione/Regressione \\
\bottomrule
\end{tabular}
\end{table}


%% ------------------------------------------------------------
\subsection{Architettura implementata: confronto con il paper} \label{sec:arch_confronto}

Poich\'e CIFAR-10 \`e un dataset molto pi\`u piccolo di ImageNet ($32 \times 32$ pixel, 50\,000 immagini vs $224 \times 224$ pixel, 1.3M immagini), abbiamo ridotto significativamente le dimensioni del modello per evitare overfitting. La tabella seguente confronta i nostri iperparametri con quelli del paper originale.

\begin{table}[htbp]
\centering
\caption{Confronto degli iperparametri: il nostro modello CIFAR-10 vs il Perceiver del paper su ImageNet.}
\label{tab:iperparametri_confronto}
\begin{tabular}{@{}lcc@{}}
\toprule
\textbf{Iperparametro} & \textbf{Nostro (CIFAR-10)} & \textbf{Paper (ImageNet)} \\
\midrule
Numero latenti $N$ & 96 & 512 \\
Dimensione latente $D$ & 384 & 1024 \\
Stadi cross-attention $T$ & 4 & 8 \\
Blocchi transformer per stadio $L$ & 4 & 6 \\
Numero teste $H$ & 3 & 8 \\
Dropout & 0.2 & 0.0 \\
Ottimizzatore & LAMB & LAMB \\
Learning rate & 0.004 & 0.004 \\
Scheduler & MultiStep & Cosine \\
Epoche & 120 & 210 \\
Batch size & 64 & 1024 \\
Fourier bands $K$ & 64 & 64 \\
Max frequency & 32.0 & 110.0 \\
\midrule
\textbf{Parametri totali (con WS)} & \textbf{$\sim$3.35M} & \textbf{$\sim$45M} \\
\bottomrule
\end{tabular}
\end{table}

\begin{notabox}[Perch\'e dimensioni cos\`i diverse?]
La scelta di ridurre il modello \`e guidata da tre considerazioni pratiche:
\begin{enumerate}[label=(\roman*)]
    \item \textbf{Dimensione del dataset}: CIFAR-10 ha solo 50\,000 immagini di training. Un modello da 45M di parametri andrebbe in overfitting severo.
    \item \textbf{Dimensione dell'input}: un'immagine CIFAR-10 ha $32 \times 32 = 1{,}024$ pixel, contro i $224 \times 224 = 50{,}176$ di ImageNet. Non servono 512 latenti per comprimere 1\,024 pixel.
    \item \textbf{Risorse computazionali}: con una singola GPU consumer, batch size di 1024 non \`e praticabile; 64 \`e un compromesso ragionevole.
\end{enumerate}
\end{notabox}

\subsubsection{Dimensione dell'input: da dove viene 76?}

La dimensione dell'input per ogni pixel \`e $d_{\text{input}} = 76$, calcolata come:
\begin{itemize}
    \item \textbf{RGB flattened}: $3 \times 4 = 12$ (3 canali colore $\times$ 4 dimensioni per il flatten, dato che l'immagine $32 \times 32$ viene appiattita come sequenza di pixel, ognuno descritto da 3 valori RGB moltiplicati per la struttura dei patch)
    \item \textbf{Fourier features}: $64$ bande $= 2 \times 32 = 64$ componenti (seno e coseno per ogni frequenza sulle 2 coordinate spaziali)
    \item \textbf{Totale}: $12 + 64 = 76$
\end{itemize}

\begin{center}
\begin{tikzpicture}[
    box/.style={rectangle, draw, rounded corners=3pt, minimum height=0.8cm, minimum width=2cm, font=\small},
    arrow/.style={-Stealth, thick}
]
\node[box, fill=red!15] (rgb) at (0,0) {RGB: 12 dim};
\node[box, fill=blue!15] (fourier) at (4.5,0) {Fourier PE: 64 dim};
\node[box, fill=green!15, minimum width=4cm] (concat) at (2.25,-1.5) {Concatenazione: 76 dim};
\node[box, fill=orange!15, minimum width=4cm] (cross) at (2.25,-3) {Input alla cross-attention};
\draw[arrow] (rgb.south) -- (concat.north west);
\draw[arrow] (fourier.south) -- (concat.north east);
\draw[arrow] (concat.south) -- (cross.north);
\end{tikzpicture}
\end{center}


%% ------------------------------------------------------------
\subsection{Esperimento 1: L'importanza del Positional Encoding} \label{sec:exp_pe_importance}

La prima domanda che ci siamo posti \`e forse la pi\`u fondamentale: \textit{quanto conta il positional encoding per il Perceiver?} Il paper originale sottolinea l'importanza delle Fourier features, ma non fornisce un'ablazione diretta che confronti ``con PE'' vs ``senza PE'' sullo stesso setup. Noi l'abbiamo fatto.

\subsubsection{Setup sperimentale}

Abbiamo confrontato due esperimenti identici in tutto tranne che per l'input:
\begin{itemize}
    \item \textbf{exp3A\_fourier\_control}: input con Fourier PE concatenato, dimensione 76
    \item \textbf{exp3B\_rgb\_only}: input con solo i valori RGB, dimensione 12, \textbf{nessun positional encoding}
\end{itemize}
Tutti gli altri iperparametri (architettura, ottimizzatore, scheduler, epoche) sono identici.

\subsubsection{Risultati}

\begin{center}
\begin{tikzpicture}
\begin{axis}[
    ybar,
    bar width=1.8cm,
    width=10cm,
    height=7cm,
    ymin=50, ymax=80,
    ylabel={Accuracy (\%)},
    symbolic x coords={Fourier PE, Senza PE},
    xtick=data,
    nodes near coords,
    nodes near coords align={vertical},
    every node near coord/.append style={font=\bfseries\small},
    title={Effetto del Positional Encoding su CIFAR-10},
    enlarge x limits=0.4,
    ymajorgrids=true,
    grid style={dashed, gray!40},
]
\addplot[fill=blue!60, draw=blue!80] coordinates {(Fourier PE, 72.02) (Senza PE, 61.34)};
\end{axis}
\draw[<->, very thick, red] (6.8,1.0) -- node[right, font=\small\bfseries, text=red] {$-10.68$ pp} (6.8,3.2);
\end{tikzpicture}
\end{center}

Il risultato \`e netto: senza positional encoding, l'accuratezza crolla di \textbf{10.68 punti percentuali} (una riduzione relativa del 14.8\%). Questo \`e un dato molto significativo.

\begin{examplebox}[Risultato chiave: il PE \`e cruciale]
\textbf{Con Fourier PE}: 72.02\% \quad vs \quad \textbf{Senza PE}: 61.34\% \quad $\Rightarrow$ \quad $\Delta = -10.68$ pp

Senza positional encoding, il Perceiver non sa \textit{dove} si trovano i pixel nell'immagine. Vede un insieme non ordinato di 1\,024 triplette RGB, senza alcuna informazione sulla disposizione spaziale. \`E come se qualcuno vi mostrasse tutti i tasselli di un mosaico sparsi su un tavolo, senza dirvi come vanno assemblati: potete vedere i colori, ma non potete ricostruire l'immagine.

Il fatto che il modello raggiunga comunque il 61.34\% senza PE suggerisce che anche le sole statistiche dei colori (istogrammi RGB, correlazioni tra canali) contengono informazione discriminante --- ma non sufficiente per una classificazione accurata.
\end{examplebox}

\subsubsection{Confronto con il paper}

Il paper originale riporta un confronto tra Fourier PE (78.0\%) e Learned PE (70.9\%) su ImageNet (Table~2), mostrando un gap di 7.1 punti percentuali. Il nostro esperimento affronta una domanda ancora pi\`u radicale: cosa succede se si \textbf{rimuove completamente} il positional encoding? Il calo che osserviamo (72.02\% $\to$ 61.34\%, ovvero $-10.68$ pp) \`e ancora pi\`u drammatico del gap tra Fourier e Learned riportato nel paper. Questo \`e coerente: nel paper, il Learned PE fornisce comunque \emph{qualche} informazione posizionale --- semplicemente la codifica in modo meno efficiente. Nel nostro caso, il modello non ha \emph{nessuna} informazione su dove si trovano i pixel, il che \`e una condizione molto pi\`u estrema.

In pratica, questo risultato conferma e \textbf{rafforza} il messaggio del paper: il positional encoding non \`e un accessorio opzionale, ma un componente strutturale senza il quale il Perceiver non pu\`o funzionare correttamente.

\subsubsection{Cosa impara il modello senza PE?}

Una domanda legittima \`e: come fa il modello a raggiungere comunque il 61.34\% senza alcuna informazione posizionale? La risposta \`e che le sole statistiche dei colori contengono una quantit\`a sorprendente di informazione discriminante. Senza PE, il Perceiver vede l'input come un \textbf{insieme non ordinato} di 1\,024 triplette RGB --- in pratica, un istogramma dei colori con informazioni sulle correlazioni tra canali.

Detto in parole semplici, il modello pu\`o distinguere ``prevalentemente blu'' (cielo, oceano) da ``prevalentemente verde'' (foresta, prato) e da ``prevalentemente grigio'' (automobili, camion). Ma non pu\`o distinguere un gatto che guarda a sinistra da un gatto che guarda a destra, perch\'e questa distinzione richiede sapere \textit{dove} si trovano i pixel chiari e scuri nell'immagine. La posizione \`e la ``colla'' che trasforma un sacchetto di pixel in un'immagine strutturata.

\begin{notabox}[Analogia: puzzle senza posizioni]
Immaginate di avere tutti i tasselli di un puzzle da 1\,024 pezzi sparsi casualmente su un tavolo. Senza PE, il Perceiver \`e come un giocatore bendato che pu\`o solo toccare ogni tassello e sentire il colore (liscio = chiaro, ruvido = scuro), ma non sa in quale riga e colonna va posizionato. Con il Fourier PE, ogni tassello ha scritto sul retro il suo numero di riga e di colonna: il giocatore pu\`o ricostruire l'immagine completa. La differenza di 10.68 punti percentuali misura esattamente quanto vale questa informazione posizionale.
\end{notabox}


%% ------------------------------------------------------------
\subsection{Esperimento 2: Fourier vs Learned PE su input permutato} \label{sec:exp_fourier_vs_learned}

Il paper originale confronta Fourier features e learned positional encoding su ImageNet, trovando che le Fourier features sono nettamente superiori ($78.0\%$ vs $70.9\%$). Ma cosa succede quando l'input \`e \textbf{permutato}? Ovvero, quando i pixel vengono mescolati con una permutazione casuale fissa (la stessa per tutte le immagini)?

Questo esperimento testa una propriet\`a fondamentale: se il PE codifica correttamente la posizione ``originale'' di ogni pixel, il modello dovrebbe funzionare \textbf{indipendentemente dall'ordine} in cui i pixel vengono presentati.

\subsubsection{Setup sperimentale}

\begin{itemize}
    \item \textbf{exp6\_fourier\_permuted}: Fourier PE + input permutato
    \item \textbf{exp2\_learned\_pe\_permuted}: Learned PE + input permutato
\end{itemize}
In entrambi i casi, si applica la stessa permutazione casuale fissa a tutti i pixel di tutte le immagini.

\subsubsection{Risultati}

\begin{center}
\begin{tikzpicture}
\begin{axis}[
    ybar,
    bar width=1.8cm,
    width=10cm,
    height=7cm,
    ymin=70, ymax=82,
    ylabel={Accuracy (\%)},
    symbolic x coords={Fourier PE, Learned PE},
    xtick=data,
    nodes near coords,
    nodes near coords align={vertical},
    every node near coord/.append style={font=\bfseries\small},
    title={Fourier vs Learned PE su input permutato (CIFAR-10)},
    enlarge x limits=0.4,
    ymajorgrids=true,
    grid style={dashed, gray!40},
]
\addplot[fill=teal!60, draw=teal!80] coordinates {(Fourier PE, 78.12) (Learned PE, 77.60)};
\end{axis}
\end{tikzpicture}
\end{center}

\begin{examplebox}[Risultato chiave: entrambi i PE funzionano con permutazione]
\textbf{Fourier PE permutato}: 78.12\% \quad vs \quad \textbf{Learned PE permutato}: 77.60\% \quad $\Rightarrow$ \quad $\Delta = 0.52$ pp

La differenza \`e minima (0.52 punti percentuali), il che suggerisce che \textbf{entrambi i tipi di PE sono capaci di codificare la posizione indipendentemente dall'ordine di presentazione}. Questo conferma la tesi del paper: il Perceiver con PE adeguato \`e genuinamente invariante alla permutazione.

Nota: su ImageNet, il gap Fourier vs Learned \`e di 7.1 punti (78.0\% vs 70.9\%). Su CIFAR-10 con permutazione, il gap \`e solo di 0.52 punti. Questo potrebbe essere dovuto al fatto che CIFAR-10 ha solo 1\,024 posizioni da codificare (vs 50\,176 di ImageNet), rendendo il compito del learned PE molto pi\`u semplice.
\end{examplebox}

\subsubsection{Confronto con il paper}

Su ImageNet, il paper riporta Fourier PE 78.0\% vs Learned PE 70.9\%, con un gap di 7.1 punti percentuali (Table~2). Sul nostro CIFAR-10 permutato, il gap \`e di soli 0.52 punti (78.12\% vs 77.60\%). Perch\'e una differenza cos\`i marcata?

La spiegazione \`e legata alla \textbf{dimensione dello spazio posizionale}. CIFAR-10 ha solo $32 \times 32 = 1{,}024$ posizioni distinte da codificare: il Learned PE deve imparare 1\,024 vettori, un compito perfettamente alla portata di qualsiasi ottimizzatore. ImageNet ha invece $224 \times 224 = 50{,}176$ posizioni --- quasi 50 volte di pi\`u. Memorizzare 50\,176 embedding distinti \`e enormemente pi\`u difficile, e il Learned PE non riesce a convergere a una soluzione altrettanto buona.

In pratica, il Learned PE ``funziona'' fintanto che il numero di posizioni \`e piccolo. Quando lo spazio posizionale cresce, le Fourier features mantengono il loro vantaggio perch\'e codificano la posizione tramite una \textbf{formula analitica} che non ha parametri da imparare e scala naturalmente a qualsiasi risoluzione.

\subsubsection{Implicazione pratica}

Per chi deve scegliere tra Fourier e Learned PE in un progetto reale, la regola empirica che emerge dai nostri risultati \`e chiara:
\begin{itemize}
    \item \textbf{Input piccoli} (poche migliaia di posizioni): il Learned PE \`e una scelta valida e raggiunge prestazioni comparabili.
    \item \textbf{Input grandi} (immagini ad alta risoluzione, audio lungo, video): le Fourier features sono essenziali perch\'e scalano naturalmente senza richiedere parametri aggiuntivi e, soprattutto, senza necessit\`a di ri-addestramento quando la risoluzione cambia.
\end{itemize}

\begin{formulabox}[Costo parametrico del Learned PE]
Il Learned PE deve imparare una matrice $\mathbf{E} \in \mathbb{R}^{N_{\text{pos}} \times D}$, dove $N_{\text{pos}}$ \`e il numero di posizioni e $D$ la dimensione dell'embedding:
\begin{itemize}
    \item \textbf{CIFAR-10}: $N_{\text{pos}} = 1{,}024$, $D = 384$ $\Rightarrow$ $1{,}024 \times 384 = 393{,}216$ parametri ($\sim$393K)
    \item \textbf{ImageNet}: $N_{\text{pos}} = 50{,}176$, $D = 1{,}024$ $\Rightarrow$ $50{,}176 \times 1{,}024 = 51{,}380{,}224$ parametri ($\sim$51.4M)
\end{itemize}
Nel caso di ImageNet, il Learned PE da solo ha pi\`u parametri dell'intero modello Perceiver con weight sharing ($\sim$45M)! Questo spiega sia la difficolt\`a di ottimizzazione sia il rischio di overfitting del Learned PE su grandi input.

Le Fourier features, al contrario, hanno \textbf{zero parametri apprendibili}: le frequenze sono fissate analiticamente, e la codifica posizionale \`e deterministica.
\end{formulabox}


%% ------------------------------------------------------------
\subsection{Esperimento 3: Effetto del Weight Sharing} \label{sec:exp_ws}

Il weight sharing \`e una delle scelte architetturali pi\`u importanti del Perceiver. Nel paper originale, gli autori riportano che il weight sharing \textbf{migliora} le prestazioni su ImageNet (da $72.9\%$ senza WS a $78.0\%$ con WS), fungendo da regolarizzatore. Ma questo vale anche su un dataset pi\`u piccolo come CIFAR-10?

\subsubsection{Setup sperimentale}

\begin{itemize}
    \item \textbf{exp4A\_ws\_control}: con weight sharing, 3.35M parametri
    \item \textbf{exp4B\_no\_ws}: senza weight sharing, 8.67M parametri (2.59$\times$ pi\`u parametri)
\end{itemize}

\subsubsection{Risultati}

\begin{center}
\begin{tikzpicture}
\begin{axis}[
    ybar,
    bar width=1.5cm,
    width=10cm,
    height=7cm,
    ymin=60, ymax=80,
    ylabel={Accuracy (\%)},
    symbolic x coords={Con WS (3.35M), Senza WS (8.67M)},
    xtick=data,
    nodes near coords,
    nodes near coords align={vertical},
    every node near coord/.append style={font=\bfseries\small},
    title={Effetto del Weight Sharing su CIFAR-10},
    enlarge x limits=0.4,
    ymajorgrids=true,
    grid style={dashed, gray!40},
]
\addplot[fill=violet!60, draw=violet!80] coordinates {({Con WS (3.35M)}, 68.49) ({Senza WS (8.67M)}, 73.85)};
\end{axis}
\draw[<->, very thick, red] (6.8,1.6) -- node[right, font=\small\bfseries, text=red] {$+5.36$ pp} (6.8,3.1);
\end{tikzpicture}
\end{center}

\begin{examplebox}[Risultato chiave: il WS non sempre aiuta]
\textbf{Con weight sharing}: 68.49\% (3.35M param) \quad vs \quad \textbf{Senza weight sharing}: 73.85\% (8.67M param)

Senza weight sharing, l'accuratezza \textbf{aumenta} di 5.36 punti percentuali, ma al costo di 2.59$\times$ pi\`u parametri.

Questo risultato \textbf{contraddice parzialmente} il paper, dove il weight sharing migliora le prestazioni su ImageNet. La spiegazione pi\`u probabile \`e legata alla dimensione del dataset:
\begin{itemize}
    \item Su \textbf{ImageNet} (1.3M immagini, 45M parametri): il weight sharing funge da regolarizzatore efficace, prevenendo l'overfitting su un modello molto grande.
    \item Su \textbf{CIFAR-10} (50K immagini, 3.35M--8.67M parametri): il modello \`e gi\`a piccolo e i dati sono relativamente semplici. La regolarizzazione del WS diventa eccessiva, limitando la capacit\`a espressiva del modello.
\end{itemize}
\end{examplebox}

\begin{notabox}[Lezione sul weight sharing]
Il weight sharing \`e un compromesso classico bias-varianza: aggiunge un forte inductive bias (``tutti gli stadi devono fare la stessa trasformazione''), che pu\`o essere benefico quando il modello \`e grande e i dati sono tanti, ma penalizzante quando il modello \`e piccolo e ha bisogno di tutta la sua capacit\`a espressiva. In altre parole: il weight sharing ``regolarizza'' --- e come ogni regolarizzazione, funziona bene solo quando il modello \`e a rischio di overfitting, non quando \`e gi\`a sotto-parametrizzato.
\end{notabox}

\subsubsection{Confronto dettagliato con il paper}

Il paper riporta risultati molto dettagliati sull'effetto del weight sharing nella Table~7:
\begin{itemize}
    \item \textbf{Senza WS}: 72.9\% validation, 87.7\% training $\Rightarrow$ gap train-val di 14.8 pp, chiaro segnale di \textbf{overfitting severo}.
    \item \textbf{Con WS}: 78.0\% validation, 79.5\% training $\Rightarrow$ gap train-val di soli 1.5 pp, modello \textbf{ben regolarizzato}.
\end{itemize}

La differenza chiave con il nostro setup \`e la \textbf{scala}: il modello del paper senza weight sharing ha circa 326M di parametri (ogni stadio ha i propri pesi indipendenti), addestrato su 1.3M immagini. Il nostro modello senza WS ha solo 8.67M di parametri su 50K immagini. A 8.67M di parametri, il nostro modello non soffre di overfitting cos\`i severo da necessitare una regolarizzazione aggressiva come il weight sharing. Al contrario, il weight sharing riduce i parametri a 3.35M, limitando eccessivamente la capacit\`a espressiva.

\begin{examplebox}[Calcolo del rapporto parametri/dati]
Il \textbf{rapporto parametri/immagine} ci d\`a una misura intuitiva di quanto il modello sia ``sovra-parametrizzato'' rispetto ai dati:
\begin{itemize}
    \item \textbf{Paper, senza WS}: $326\text{M} / 1.3\text{M} \approx 250$ parametri per immagine $\Rightarrow$ forte rischio di overfitting
    \item \textbf{Paper, con WS}: $45\text{M} / 1.3\text{M} \approx 35$ parametri per immagine $\Rightarrow$ ben bilanciato
    \item \textbf{Nostro, senza WS}: $8.67\text{M} / 50\text{K} \approx 173$ parametri per immagine $\Rightarrow$ rischio moderato
    \item \textbf{Nostro, con WS}: $3.35\text{M} / 50\text{K} \approx 67$ parametri per immagine $\Rightarrow$ possibilmente sotto-parametrizzato
\end{itemize}
Il modello del paper senza WS ha un rapporto di 250 params/immagine, il che spiega l'overfitting severo (87.7\% train vs 72.9\% val). Con WS, il rapporto scende a 35 e l'overfitting scompare. Nel nostro caso, anche senza WS il rapporto \`e solo 173 --- non abbastanza alto da causare overfitting grave, quindi la regolarizzazione del WS risulta eccessiva e controproducente.
\end{examplebox}

\begin{notabox}[Regola pratica: quando usare il weight sharing?]
Dai nostri esperimenti e dal confronto con il paper, emerge una regola empirica:
\begin{itemize}
    \item \textbf{Usare il WS} quando il rapporto parametri/dati \`e alto ($> 200$) e si osserva un significativo gap tra accuratezza di training e di validazione (segnale di overfitting).
    \item \textbf{Non usare il WS} quando il modello \`e gi\`a relativamente piccolo rispetto ai dati e non mostra segni di overfitting. In questo caso, ogni stadio della cross-attention iterativa beneficia dall'avere pesi indipendenti, potendo specializzarsi su livelli diversi di astrazione (feature grossolane nei primi stadi, feature fini negli ultimi).
\end{itemize}
In pratica, il weight sharing \`e uno strumento di \textbf{regolarizzazione}, non un principio architetturale universale. Come ogni regolarizzatore (dropout, weight decay, data augmentation), il suo effetto dipende dal regime operativo.
\end{notabox}


%% ------------------------------------------------------------
\subsection{Esperimento 4: Robustezza alla permutazione --- un risultato sorprendente} \label{sec:exp_permutation}

Questo \`e forse il risultato pi\`u inatteso dell'intero progetto. Confrontiamo il Perceiver con Fourier PE su input standard (pixel nell'ordine naturale, riga per riga) e su input permutato (pixel in ordine casuale fisso).

\subsubsection{Setup sperimentale}

\begin{itemize}
    \item \textbf{exp1\_baseline\_fourier}: input nell'ordine standard, Fourier PE
    \item \textbf{exp6\_fourier\_permuted}: input permutato, Fourier PE
\end{itemize}

\subsubsection{Risultati}

\begin{center}
\begin{tikzpicture}
\begin{axis}[
    ybar,
    bar width=1.8cm,
    width=10cm,
    height=7cm,
    ymin=60, ymax=85,
    ylabel={Accuracy (\%)},
    symbolic x coords={Standard, Permutato},
    xtick=data,
    nodes near coords,
    nodes near coords align={vertical},
    every node near coord/.append style={font=\bfseries\small},
    title={Ordine standard vs permutato con Fourier PE (CIFAR-10)},
    enlarge x limits=0.4,
    ymajorgrids=true,
    grid style={dashed, gray!40},
]
\addplot[fill=orange!60, draw=orange!80] coordinates {(Standard, 69.69) (Permutato, 78.12)};
\end{axis}
\draw[<->, very thick, red] (6.8,1.2) -- node[right, font=\small\bfseries, text=red] {$+8.43$ pp!} (6.8,3.6);
\end{tikzpicture}
\end{center}

\begin{examplebox}[Risultato sorprendente: la permutazione MIGLIORA l'accuratezza!]
\textbf{Input standard}: 69.69\% \quad vs \quad \textbf{Input permutato}: 78.12\% \quad $\Rightarrow$ \quad $\Delta = +8.43$ pp

Questo risultato \`e \textbf{controintuitivo}: permutare casualmente i pixel dovrebbe, nel migliore dei casi, non cambiare nulla (se il modello \`e veramente invariante alla permutazione) e nel peggiore dei casi peggiorare le cose. Invece, la permutazione \textbf{migliora} l'accuratezza di quasi 8.5 punti!

\textbf{La nostra ipotesi}: Le immagini CIFAR-10 sono molto piccole ($32 \times 32$), e le Fourier features con max\_freq$=32.0$ generano frequenze che ``risuonano'' con la griglia regolare dei pixel. In pratica, le frequenze creano pattern di aliasing sulla griglia regolare che confondono il modello. La permutazione rompe questa regolarit\`a, eliminando l'aliasing e permettendo alle Fourier features di funzionare in modo pi\`u pulito.

Un'analogia: immaginate di suonare una nota su un pianoforte che ha esattamente la stessa frequenza di risonanza del tavolo su cui \`e appoggiato. Il tavolo vibra e crea un suono sgradevole (interferenza). Se spostate il pianoforte su una superficie irregolare, la risonanza scompare e il suono \`e pi\`u pulito.
\end{examplebox}

\begin{notabox}[Confronto con il paper originale]
Su ImageNet, il paper riporta che il Perceiver con Fourier PE mantiene \textbf{esattamente} la stessa accuracy (78.0\%) su input standard e permutato. Quindi il modello \`e ``perfettamente invariante'' alla permutazione.

Nel nostro caso su CIFAR-10, non solo il modello \`e invariante, ma la permutazione \textbf{aiuta}. Questo effetto potrebbe essere specifico di:
\begin{itemize}
    \item La piccola dimensione delle immagini CIFAR-10 ($32 \times 32$)
    \item Il rapporto max\_freq/risoluzione (32.0/32 = 1.0, che crea un fenomeno di aliasing)
    \item La dimensione ridotta del modello
\end{itemize}
Sarebbe interessante ripetere l'esperimento con max\_freq diverse per confermare questa ipotesi.
\end{notabox}

\subsubsection{Confronto con il paper}

Il paper riporta nella Table~2 che il Perceiver con Fourier features ottiene \textbf{esattamente la stessa accuratezza} (78.0\%) su ImageNet standard e ImageNet permutato. Questa \`e la dimostrazione che il modello \`e genuinamente invariante alla permutazione: l'ordine in cui i pixel vengono presentati non conta, perch\'e tutta l'informazione posizionale \`e codificata nelle Fourier features.

Il nostro risultato su CIFAR-10 va \textbf{oltre}: non solo il modello \`e invariante alla permutazione, ma la permutazione \emph{migliora} l'accuratezza di 8.43 punti (69.69\% $\to$ 78.12\%). Questo \`e un risultato genuinamente nuovo che non \`e presente nel paper originale e merita un'analisi approfondita.

\subsubsection{Ipotesi 1: Aliasing delle Fourier features}

La nostra ipotesi principale riguarda l'\textbf{aliasing}. Su una griglia regolare $32 \times 32$ con \texttt{max\_freq}$= 32$, la frequenza pi\`u alta ha un periodo di esattamente 1 pixel. Questo significa che la componente sinusoidale $\sin(2\pi \cdot 32 \cdot x)$ valutata su una griglia con passo $\Delta x = 1/32$ genera pattern perfettamente periodici che ``risuonano'' con la griglia stessa. In pratica, si creano artefatti di aliasing sistematici (effetto Moir\'e) che confondono il modello.

La permutazione rompe la regolarit\`a della griglia: i pixel non sono pi\`u equispaziati, e le componenti ad alta frequenza delle Fourier features non creano pi\`u pattern di interferenza sistematici. Su ImageNet ($224 \times 224$ con \texttt{max\_freq}$= 112$), le dimensioni molto maggiori della griglia attenuano naturalmente questo effetto.

\subsubsection{Ipotesi 2: Data augmentation implicita}

La permutazione agisce come una forma molto forte di \textbf{data augmentation}. Poich\'e le Fourier features preservano la posizione ``vera'' di ogni pixel (il pixel che si trovava in posizione $(i, j)$ mantiene la sua codifica posizionale Fourier per $(i, j)$ anche quando viene spostato in una posizione diversa della sequenza), il modello dispone ancora di tutta l'informazione necessaria per ricostruire la disposizione spaziale.

Tuttavia, l'ordine ``grezzo'' della sequenza di input \`e randomizzato. Questo forza la cross-attention a fare \textbf{pi\`u affidamento} sul positional encoding (le Fourier features) e meno sull'ordine sequenziale dei pixel. In un certo senso, il modello ``impara meglio'' a usare le Fourier features proprio perch\'e non ha altra scelta.

\subsubsection{Ipotesi 3: Regolarizzazione}

Senza permutazione, il modello potrebbe imparare \textit{scorciatoie} basate sull'ordine sequenziale dei pixel. Per esempio: ``i primi 32 pixel sono sempre la prima riga dell'immagine, e le immagini di aeroplani tendono ad avere cielo (blu) nella prima riga''. Queste scorciatoie funzionano sul training set ma non generalizzano bene.

La permutazione \textbf{elimina} queste scorciatoie, forzando il modello a imparare feature pi\`u generali e robuste. Questo \`e concettualmente simile al dropout (che impedisce al modello di fare affidamento su singoli neuroni) e alla data augmentation (che impedisce di memorizzare dettagli irrilevanti degli esempi di training).

\begin{proofbox}[Verifica formale dell'invarianza alla permutazione]
Verifichiamo formalmente che la cross-attention con Fourier features \`e invariante alla permutazione. Sia $\widetilde{\mathbf{X}} \in \mathbb{R}^{M \times (C + D_{\text{PE}})}$ l'input con Fourier features concatenate, dove ogni riga $\widetilde{\mathbf{x}}_i = [\text{rgb}_i \,\|\, \text{fourier}(i)]$ contiene sia il valore RGB del pixel $i$ sia la sua codifica posizionale Fourier.

Sia $\mathbf{P}_\pi$ una matrice di permutazione che riordina le righe di $\widetilde{\mathbf{X}}$. La cross-attention calcola:
\[
\text{CrossAttn}(\mathbf{Z}, \mathbf{P}_\pi \widetilde{\mathbf{X}}) = \text{softmax}\!\left(\frac{\mathbf{Z} \mathbf{W}_Q (\mathbf{P}_\pi \widetilde{\mathbf{X}} \mathbf{W}_K)^\top}{\sqrt{d_k}}\right) \mathbf{P}_\pi \widetilde{\mathbf{X}} \mathbf{W}_V
\]

Poich\'e $\mathbf{P}_\pi$ permuta le \textbf{righe} (i pixel), e poich\'e ogni riga $\widetilde{\mathbf{x}}_i$ contiene gi\`a la propria codifica posizionale Fourier (che non dipende dalla posizione nella sequenza ma dalla posizione \emph{originale} nell'immagine), l'output della softmax \`e lo stesso a meno di una ri-permutazione delle colonne. Poich\'e la softmax normalizza su tutte le posizioni e il risultato finale \`e una somma pesata, l'output \`e \textbf{identico}:
\[
\text{CrossAttn}(\mathbf{Z}, \mathbf{P}_\pi \widetilde{\mathbf{X}}) = \text{CrossAttn}(\mathbf{Z}, \widetilde{\mathbf{X}})
\]

Questa identit\`a vale perch\'e (1) le Fourier features sono ``ancorate'' alla posizione originale del pixel, non alla posizione nella sequenza, e (2) l'attenzione calcola una media pesata che \`e invariante all'ordine dei termini.
\end{proofbox}


%% ------------------------------------------------------------
\subsection{Riepilogo CIFAR-10: tutti gli esperimenti} \label{sec:riepilogo_cifar}

Prima di passare agli altri dataset, raccogliamo tutti i risultati CIFAR-10 in un'unica tabella e in un grafico riassuntivo.

\begin{table}[htbp]
\centering
\caption{Riepilogo completo degli esperimenti CIFAR-10. Tutti i modelli hanno la stessa architettura base ($N=96$, $D=384$, $T=4$, $L=4$, $H=3$), 120 epoche di addestramento.}
\label{tab:riepilogo_cifar}
\small
\begin{tabular}{@{}llccccr@{}}
\toprule
\textbf{Esperimento} & \textbf{PE} & \textbf{Perm.} & \textbf{WS} & \textbf{Dim Input} & \textbf{Parametri} & \textbf{Acc (\%)} \\
\midrule
exp3A\_fourier\_control & Fourier & No & S\`i & 76 & 3.35M & 72.02 \\
exp1\_baseline\_fourier & Fourier & No & S\`i & 76 & 3.35M & 69.69 \\
exp4A\_ws\_control & Fourier & No & S\`i & 76 & 3.35M & 68.49 \\
exp3B\_rgb\_only & Nessuno & No & S\`i & 12 & 3.30M & 61.34 \\
exp4B\_no\_ws & Fourier & No & No & 76 & 8.67M & 73.85 \\
exp6\_fourier\_permuted & Fourier & S\`i & S\`i & 76 & 3.35M & \textbf{78.12} \\
exp2\_learned\_pe\_permuted & Learned & S\`i & S\`i & 76 & 3.35M & 77.60 \\
\midrule
Perceiver IO & Fourier & No & S\`i & 76 & $\sim$3.35M & 78.20 \\
\bottomrule
\end{tabular}
\end{table}

\begin{center}
\begin{tikzpicture}
\begin{axis}[
    xbar,
    bar width=0.5cm,
    width=13cm,
    height=10cm,
    xmin=55, xmax=85,
    xlabel={Accuracy (\%)},
    ytick={1,2,3,4,5,6,7,8},
    yticklabels={
        {Perceiver IO},
        {exp6 Fourier+Perm},
        {exp2 Learned+Perm},
        {exp4B No WS},
        {exp3A Fourier ctrl},
        {exp1 Baseline},
        {exp4A WS ctrl},
        {exp3B RGB only}
    },
    y dir=reverse,
    nodes near coords,
    nodes near coords align={horizontal},
    every node near coord/.append style={font=\small\bfseries, anchor=west},
    title={CIFAR-10: Accuratezza di tutti gli esperimenti},
    xmajorgrids=true,
    grid style={dashed, gray!40},
    enlarge y limits=0.08,
]
\addplot[fill=green!60, draw=green!80] coordinates {(78.20, 1)};
\addplot[fill=orange!60, draw=orange!80] coordinates {(78.12, 2)};
\addplot[fill=orange!40, draw=orange!60] coordinates {(77.60, 3)};
\addplot[fill=violet!50, draw=violet!70] coordinates {(73.85, 4)};
\addplot[fill=blue!50, draw=blue!70] coordinates {(72.02, 5)};
\addplot[fill=blue!30, draw=blue!50] coordinates {(69.69, 6)};
\addplot[fill=blue!20, draw=blue!40] coordinates {(68.49, 7)};
\addplot[fill=red!40, draw=red!60] coordinates {(61.34, 8)};
\end{axis}
\end{tikzpicture}
\end{center}

\begin{notabox}[Osservazioni chiave dal riepilogo CIFAR-10]
Dal grafico emergono chiaramente tre ``cluster'' di prestazioni:
\begin{enumerate}[label=(\roman*)]
    \item \textbf{Top tier} ($\geq 77\%$): gli esperimenti con permutazione (exp6, exp2) e il Perceiver IO. La permutazione, come discusso, rompe l'aliasing delle Fourier features.
    \item \textbf{Mid tier} ($68$--$74\%$): gli esperimenti con Fourier PE senza permutazione e l'esperimento senza WS. C'\`e una variabilit\`a di circa 4 punti tra le diverse run con gli stessi iperparametri (exp1 vs exp3A vs exp4A), il che suggerisce una certa sensibilit\`a all'inizializzazione casuale.
    \item \textbf{Low tier} ($\leq 62\%$): l'esperimento senza PE (exp3B). La mancanza di informazione posizionale \`e chiaramente il fattore limitante.
\end{enumerate}
Notiamo anche che exp1, exp3A e exp4A hanno gli stessi iperparametri ma accuratezze diverse (69.69\%, 72.02\%, 68.49\%). Questa variabilit\`a di circa 3.5 punti percentuali \`e coerente con la letteratura sul training di modelli basati su attention con dataset piccoli: la stocasticit\`a dell'inizializzazione e dell'ordine dei mini-batch ha un impatto significativo.
\end{notabox}

\begin{figure}[H]
\centering
\includegraphics[width=0.9\textwidth]{paper_figures/our_cifar10_chart.png}
\caption{Risultati della nostra ablation study su CIFAR-10. Il grafico mostra l'accuratezza di validazione per tutti e 7 gli esperimenti Perceiver pi\`u il Perceiver IO. L'esperimento con permutazione e Fourier PE (exp6) raggiunge il risultato migliore.}
\label{fig:our_cifar10}
\end{figure}

La Figura~\ref{fig:our_cifar10} offre una visione d'insieme dei risultati della nostra ablation study. Si distinguono chiaramente i tre cluster di prestazioni menzionati sopra: gli esperimenti con permutazione dominano la classifica, seguiti dagli esperimenti standard con Fourier PE, e infine l'esperimento senza positional encoding. Questo grafico conferma visivamente che la scelta del positional encoding \`e il fattore pi\`u determinante per le prestazioni del Perceiver.

\begin{figure}[H]
\centering
\includegraphics[width=0.9\textwidth]{paper_figures/our_convergence_chart.png}
\caption{Curve di convergenza dei nostri esperimenti CIFAR-10. Si nota come gli esperimenti con permutazione (exp2, exp6) convergano pi\`u lentamente ma raggiungano accuratezze finali pi\`u alte.}
\label{fig:our_convergence}
\end{figure}

Le curve di convergenza (Figura~\ref{fig:our_convergence}) rivelano un aspetto importante della dinamica di addestramento: gli esperimenti con permutazione partono pi\`u lentamente nelle prime epoche, ma superano gli altri esperimenti nella fase finale del training. Questo suggerisce che la permutazione rende il task inizialmente pi\`u difficile (il modello deve ``riordinare'' i pixel prima di poterli interpretare), ma una volta appresa questa riorganizzazione, il modello converge a una soluzione pi\`u generalizzata.


%% ------------------------------------------------------------
\subsection{ModelNet40: Point Cloud 3D} \label{sec:exp_modelnet}

Per verificare la generalit\`a del Perceiver su dati non-immagine, lo abbiamo addestrato su ModelNet40, un dataset di nuvole di punti 3D. Ogni oggetto \`e rappresentato come un insieme di 1\,024 punti nello spazio tridimensionale $(x, y, z)$, con il compito di classificarlo in una delle 40 categorie (es.\ sedia, tavolo, aereo, ecc.).

Questo esperimento \`e particolarmente significativo perch\'e le nuvole di punti sono \textbf{intrinsecamente non ordinate}: a differenza dei pixel di un'immagine, non c'\`e un ordine ``naturale'' per i punti 3D. Il Perceiver, che tratta l'input come un insieme non ordinato (con PE per codificare la posizione), \`e concettualmente ideale per questo dominio.

\subsubsection{Risultati}

\begin{table}[htbp]
\centering
\caption{Risultati su ModelNet40 (classificazione 3D, 40 classi). Il modello ha 5.93M parametri, 200 epoche.}
\label{tab:modelnet_risultati}
\begin{tabular}{@{}lccl@{}}
\toprule
\textbf{Esperimento} & \textbf{Best Acc (\%)} & \textbf{Epoca migliore} & \textbf{Note} \\
\midrule
Baseline (senza augmentation) & 84.16 (84.24) & 74 & --- \\
Con rotazione casuale & 83.06 (83.14) & 45 & $-1.10$ pp \\
Con traslazione casuale & 82.90 (83.67) & 62 & $-1.26$ pp \\
\midrule
\textbf{Paper originale} & \textbf{85.7} & --- & PointNet++-like \\
\bottomrule
\end{tabular}
\end{table}

\begin{center}
\begin{tikzpicture}
\begin{axis}[
    ybar,
    bar width=1.4cm,
    width=12cm,
    height=7cm,
    ymin=78, ymax=88,
    ylabel={Accuracy (\%)},
    symbolic x coords={Paper, Baseline, Rotazione, Traslazione},
    xtick=data,
    nodes near coords,
    nodes near coords align={vertical},
    every node near coord/.append style={font=\bfseries\small},
    title={ModelNet40: confronto con il paper e effetto delle augmentation},
    enlarge x limits=0.25,
    ymajorgrids=true,
    grid style={dashed, gray!40},
]
\addplot[fill=gray!50, draw=gray!70] coordinates {(Paper, 85.70)};
\addplot[fill=blue!60, draw=blue!80] coordinates {(Baseline, 84.16)};
\addplot[fill=orange!50, draw=orange!70] coordinates {(Rotazione, 83.06)};
\addplot[fill=red!50, draw=red!70] coordinates {(Traslazione, 82.90)};
\end{axis}
\end{tikzpicture}
\end{center}

\begin{examplebox}[Risultato chiave: vicinanza al paper su ModelNet40]
Il nostro baseline raggiunge \textbf{84.16\%}, solo 1.54 punti sotto il risultato del paper (\textbf{85.7\%}). Considerando che non abbiamo replicato esattamente il setup del paper (diverse augmentation, possibilmente diverso numero di punti campionati), questo risultato \`e molto soddisfacente.

L'aspetto sorprendente \`e che le \textbf{augmentation peggiorano} leggermente le prestazioni ($-1.10$ pp per la rotazione, $-1.26$ pp per la traslazione). Questo suggerisce che il modello ha gi\`a appreso una certa invarianza geometrica dai dati, e l'augmentation aggiuntiva introduce rumore piuttosto che regolarizzazione utile. Alternativamente, le augmentation potrebbero richiedere pi\`u epoche per essere sfruttate efficacemente.
\end{examplebox}

\subsubsection{Confronto con il paper}

Il paper originale riporta un'accuratezza di \textbf{85.7\%} su ModelNet40. Il nostro modello baseline raggiunge \textbf{84.16\%}, con un gap di soli 1.54 punti percentuali. Questo risultato \`e notevole considerando che abbiamo utilizzato un modello con dimensioni ridotte rispetto alla configurazione del paper. La vicinanza al risultato del paper dimostra che il Perceiver \`e efficace anche su dati 3D e che le Fourier features sono una codifica posizionale adatta anche per coordinate spaziali tridimensionali.

\subsubsection{Effetto dell'augmentation: perch\'e peggiora?}

Il fatto che la data augmentation (rotazione $-1.10$ pp, traslazione $-1.26$ pp) \textbf{peggiori} le prestazioni sembra controintuitivo, dato che l'augmentation \`e generalmente considerata benefica. La spiegazione pi\`u probabile \`e legata all'interazione tra la dimensione del dataset e il numero di epoche:

ModelNet40 ha solo \textbf{9\,843 campioni di training} --- un dataset molto piccolo. Con sole 200 epoche, il modello ha un budget computazionale limitato. L'augmentation rende il task di apprendimento pi\`u difficile (ogni esempio appare in pose leggermente diverse ad ogni epoca), e il modello non ha abbastanza iterazioni per sfruttare questa diversit\`a aggiuntiva.

In pratica, l'augmentation ``diluisce'' il segnale di apprendimento: invece di vedere lo stesso oggetto 200 volte nella stessa posizione (consolidando la rappresentazione), il modello vede 200 versioni leggermente diverse dello stesso oggetto, senza avere il tempo di generalizzare correttamente da questa variabilit\`a. Con un numero maggiore di epoche (ad esempio 500--1000), ci aspetteremmo che l'augmentation inizi a portare benefici.

\subsubsection{Confronto con PointNet++}

Per contestualizzare i risultati, \`e utile confrontare il Perceiver con PointNet++ (Qi et al., 2017), il modello specializzato pi\`u noto per le point cloud. PointNet++ raggiunge \textbf{91.9\%} su ModelNet40, con un gap di 7.7 punti percentuali rispetto al nostro baseline.

Questo gap \`e il \textbf{prezzo della generalit\`a}. PointNet++ sfrutta la struttura geometrica locale delle point cloud attraverso meccanismi specializzati: raggruppamento gerarchico dei punti vicini (\emph{set abstraction}), costruzione di feature locali a scale multiple, e aggregazione progressiva dell'informazione dal locale al globale. Questi inductive bias sono potentissimi per dati 3D, ma non trasferibili ad altri domini.

Il Perceiver, al contrario, tratta ogni punto in modo completamente indipendente, affidandosi unicamente alle Fourier features per l'informazione spaziale. Non ha nozione di ``vicinanza'' tra punti, non costruisce gerarchie locali, e non sfrutta la simmetria rotazionale. Eppure raggiunge l'84.16\% --- a dimostrazione che un'architettura \textbf{generale} pu\`o avvicinarsi a modelli \textbf{specializzati}, pur senza superarli.

\begin{figure}[H]
\centering
\includegraphics[width=0.7\textwidth]{paper_figures/our_modelnet40_chart.png}
\caption{Risultati dei nostri esperimenti su ModelNet40. Il baseline (84.16\%) \`e molto vicino al risultato del paper (85.7\%). L'aggiunta di augmentation (rotazione, traslazione) peggiora leggermente le prestazioni, probabilmente a causa del dataset piccolo e del numero limitato di epoche.}
\label{fig:our_modelnet40}
\end{figure}

La Figura~\ref{fig:our_modelnet40} visualizza i risultati dei nostri tre esperimenti su ModelNet40 rispetto al benchmark del paper. Il gap di soli 1.54 punti percentuali tra il nostro baseline e il risultato pubblicato \`e particolarmente incoraggiante, considerando le differenze nel setup sperimentale. Il grafico rende anche evidente l'effetto negativo della data augmentation nel nostro regime di training: sia la rotazione casuale che la traslazione peggiorano le prestazioni, un risultato controintuitivo che abbiamo discusso in dettaglio nella Sezione~\ref{sec:exp_modelnet}.


%% ------------------------------------------------------------
\subsection{Perceiver IO: Estensione multi-task} \label{sec:exp_perceiver_io}

Il Perceiver IO (Jaegle et al., 2022) estende il Perceiver aggiungendo un \textbf{decoder} basato su cross-attention. Mentre il Perceiver originale produce un singolo vettore di output (tramite average pooling), il Perceiver IO pu\`o produrre output strutturati di dimensione arbitraria, rendendolo adatto a task come il masked language modeling (MLM) e la classificazione di sequenze testuali.

\subsubsection{Architettura del decoder}

Il decoder del Perceiver IO funziona ``al contrario'' rispetto all'encoder:
\begin{itemize}
    \item \textbf{Query}: provengono dall'output desiderato (task-specific query)
    \item \textbf{Key e Value}: provengono dai latenti $\mathbf{z}_T$ prodotti dall'encoder
\end{itemize}
Questo permette di ``estrarre'' dai latenti esattamente l'informazione necessaria per ogni task specifico.

\subsubsection{Cosa cambia nel Perceiver IO rispetto al Perceiver originale}

La differenza fondamentale tra Perceiver e Perceiver IO \`e nel modo in cui l'output viene prodotto. Nel Perceiver originale, i latenti $\mathbf{z}_T$ vengono semplicemente mediati (average pooling) e passati a un classificatore lineare. Questo funziona bene per la classificazione (dove serve un singolo vettore), ma \`e inadeguato per task che richiedono output strutturati (ad esempio, predire un token per ogni posizione dell'input in un task di MLM).

Il Perceiver IO risolve questo problema con un \textbf{decoder cross-attention}: un array di query $\mathbf{O} \in \mathbb{R}^{O \times D_O}$ (specifico per il task) viene usato come query, mentre i latenti $\mathbf{z}_T$ servono come key e value:
\[
\text{output} = \text{CrossAttn}(\mathbf{O}, \mathbf{z}_T) = \text{softmax}\!\left(\frac{\mathbf{O} \mathbf{W}_Q^{\text{dec}} (\mathbf{z}_T \mathbf{W}_K^{\text{dec}})^\top}{\sqrt{d_k}}\right) \mathbf{z}_T \mathbf{W}_V^{\text{dec}}
\]

Detto in parole semplici: il decoder ``interroga'' i latenti ponendo domande specifiche per il task. Per la classificazione, l'output array ha una sola riga (``a quale classe appartiene questa immagine?''). Per il MLM, ha tante righe quanti sono i token mascherati (``quale carattere va in posizione 15? E in posizione 42?''). Questa flessibilit\`a \`e ci\`o che rende il Perceiver IO un modello genuinamente multi-task.

\subsubsection{CIFAR-10 con Perceiver IO}

Sul task di classificazione CIFAR-10, il Perceiver IO raggiunge \textbf{78.20\%}, leggermente superiore al miglior risultato del Perceiver standard (78.12\% con permutazione). La differenza \`e minima e potrebbe essere dovuta alla stocasticit\`a del training, ma conferma che il decoder del Perceiver IO non introduce overhead significativo per task di classificazione semplice.

\subsubsection{Masked Language Modeling su WikiText-103}

Abbiamo pre-addestrato il Perceiver IO su WikiText-103 con l'obiettivo di Masked Language Modeling: il modello deve predire i token mascherati dato il contesto circostante.

\begin{formulabox}[Risultati MLM]
\begin{itemize}
    \item \textbf{WikiText-103}: 82.20\% di accuratezza nella predizione dei token mascherati (49 epoche)
    \item \textbf{WikiText-2}: 19.54\% (solo 3 epoche --- addestramento insufficiente)
\end{itemize}
L'accuratezza dell'82.20\% su WikiText-103 indica che il Perceiver IO riesce a catturare pattern linguistici significativi, anche se non raggiunge le prestazioni di modelli specializzati come BERT. Va sottolineato che il nostro modello ha molti meno parametri e non utilizza un tokenizer ottimizzato.
\end{formulabox}

\subsubsection{Analisi del risultato MLM}

Il risultato dell'82.20\% di accuratezza a livello di byte merita un'analisi pi\`u attenta. \`E importante notare che questa \textbf{non} \`e un'accuratezza a livello di parola: il Perceiver IO processa il testo come una sequenza di byte (256 possibili valori), non come una sequenza di token WordPiece o BPE come BERT.

A livello di byte, molte predizioni sono ``facili'': gli spazi tra le parole sono altamente prevedibili (dopo una parola lunga, quasi certamente segue uno spazio), le lettere comuni come `e', `t', `a' hanno distribuzioni contestuali molto concentrate, e la struttura sintattica dell'inglese rende molti caratteri quasi deterministici dato il contesto. Questo ``gonfia'' la metrica rispetto a un'accuratezza a livello di parola.

Tuttavia, il risultato \`e comunque significativo: mostra che il Perceiver IO, con la stessa architettura usata per classificare immagini e point cloud, riesce a catturare pattern linguistici non banali --- dalla struttura delle parole inglesi alle dipendenze sintattiche locali.

\subsubsection{Fine-tuning su GLUE benchmark}

Dopo il pre-training su WikiText-103, abbiamo effettuato il fine-tuning del Perceiver IO sugli 8 task del benchmark GLUE (General Language Understanding Evaluation), che coprono diversi aspetti della comprensione del linguaggio naturale.

\begin{table}[htbp]
\centering
\caption{Risultati del Perceiver IO sui task GLUE dopo fine-tuning dal pre-training WikiText-103.}
\label{tab:glue_risultati}
\begin{tabular}{@{}llccl@{}}
\toprule
\textbf{Task} & \textbf{Tipo} & \textbf{Metrica} & \textbf{Risultato} & \textbf{Epoche} \\
\midrule
QQP & Parafrasi (coppie) & Accuracy & 75.65\% & 24 \\
CoLA & Accettabilit\`a linguistica & Accuracy & 69.13\% & 1 \\
MRPC & Parafrasi (coppie) & Accuracy & 68.38\% & 1 \\
SST-2 & Sentiment analysis & Accuracy & 61.24\% & 13 \\
QNLI & NLI (domanda-contesto) & Accuracy & 59.93\% & 17 \\
RTE & NLI (testual entailment) & Accuracy & 52.71\% & 4 \\
MNLI & NLI (multi-genre) & Accuracy & 46.47\% & 22 \\
STS-B & Similarit\`a semantica & MSE & 2.28 & 4 \\
\bottomrule
\end{tabular}
\end{table}

\begin{center}
\begin{tikzpicture}
\begin{axis}[
    ybar,
    bar width=0.7cm,
    width=14cm,
    height=7.5cm,
    ymin=30, ymax=85,
    ylabel={Accuracy (\%)},
    symbolic x coords={QQP, CoLA, MRPC, SST-2, QNLI, RTE, MNLI},
    xtick=data,
    x tick label style={rotate=30, anchor=east, font=\small},
    nodes near coords,
    nodes near coords align={vertical},
    every node near coord/.append style={font=\scriptsize\bfseries},
    title={Perceiver IO: Risultati sui task GLUE},
    enlarge x limits=0.1,
    ymajorgrids=true,
    grid style={dashed, gray!40},
]
\addplot[fill=cyan!60, draw=cyan!80] coordinates {
    (QQP, 75.65)
    (CoLA, 69.13)
    (MRPC, 68.38)
    (SST-2, 61.24)
    (QNLI, 59.93)
    (RTE, 52.71)
    (MNLI, 46.47)
};
\end{axis}
\end{tikzpicture}
\end{center}

\begin{notabox}[Analisi dei risultati GLUE]
I risultati sui task GLUE sono modesti rispetto ai modelli stato dell'arte (BERT-base raggiunge $\sim$80\% in media sui task GLUE), ma vanno contestualizzati:
\begin{enumerate}[label=(\roman*)]
    \item Il nostro Perceiver IO ha \textbf{molti meno parametri} di BERT-base (110M).
    \item Il pre-training \`e stato effettuato su WikiText-103 (103M token), non su BooksCorpus + Wikipedia ($\sim$3.3B token come BERT).
    \item Non utilizziamo un tokenizer ottimizzato (WordPiece/BPE) ma un encoding a livello di carattere/byte.
    \item Il task QQP (75.65\%) e CoLA (69.13\%) mostrano risultati ragionevoli, suggerendo che il Perceiver IO cattura almeno alcune strutture linguistiche.
    \item MNLI (46.47\%) \`e il task pi\`u difficile, coerentemente con la sua natura multi-genre e il bisogno di ragionamento complesso.
    \item Il fatto che CoLA e MRPC raggiungano buoni risultati in solo 1 epoca suggerisce che il pre-training ha fornito una buona inizializzazione per questi task.
\end{enumerate}
\end{notabox}

\subsubsection{Analisi per task GLUE}

Analizziamo pi\`u nel dettaglio i risultati, raggruppando i task per difficolt\`a:

\textbf{Task ``facili'' (accuracy $> 65\%$):} QQP (75.65\%), CoLA (69.13\%), MRPC (68.38\%). Questi task hanno in comune la natura binaria della classificazione e la relativa semplicit\`a del pattern da riconoscere. QQP (Quora Question Pairs) chiede se due domande sono parafrasi l'una dell'altra --- un compito che si pu\`o in parte risolvere con overlap lessicale. CoLA (Corpus of Linguistic Acceptability) chiede se una frase \`e grammaticalmente corretta --- un pattern spesso catturabile a livello di carattere/byte. MRPC (Microsoft Research Paraphrase Corpus) \`e simile a QQP.

\textbf{Task ``difficili'' (accuracy $< 65\%$):} SST-2 (61.24\%), QNLI (59.93\%), RTE (52.71\%), MNLI (46.47\%). Questi task richiedono una comprensione semantica pi\`u profonda. MNLI (Multi-Genre Natural Language Inference) \`e particolarmente difficile perch\'e richiede ragionamento su relazioni logiche (entailment, contradiction, neutral) tra coppie di frasi provenienti da generi diversi, con 3 classi invece di 2. Il Perceiver IO fatica su questi task perch\'e il processing a livello di byte perde le \textbf{unit\`a semantiche} (parole, morfemi) che modelli come BERT catturano nativamente attraverso il tokenizer WordPiece.

Detto in parole semplici: per capire se ``Il gatto insegue il topo'' implica ``Un animale si muove'', bisogna prima riconoscere le parole ``gatto'' e ``insegue'' come unit\`a semantiche, poi collegare ``gatto'' $\to$ ``animale'' e ``insegue'' $\to$ ``si muove''. Il Perceiver IO, lavorando a livello di byte, deve prima ricostruire le parole dai singoli caratteri e \emph{poi} fare il ragionamento semantico --- un compito enormemente pi\`u arduo.


%% ------------------------------------------------------------
\subsection{Analisi delle Attention Maps} \label{sec:exp_attention}

Oltre ai risultati quantitativi, abbiamo condotto un'analisi qualitativa delle mappe di attenzione del Perceiver durante il forward pass su CIFAR-10. Le mappe di attenzione della cross-attention mostrano \textit{quali pixel dell'input ogni latente sta ``guardando''} ad ogni stadio di iterazione.

\subsubsection{Cosa mostrano le attention maps}

Per ogni stadio di cross-attention $t = 1, \ldots, T$, i pesi di attenzione formano una matrice $\mathbf{A}^{(t)} \in \mathbb{R}^{N \times M}$, dove $A^{(t)}_{ij}$ indica quanto il latente $i$ ``attende'' al pixel $j$. Riorganizzando la riga $i$-esima sulla griglia $32 \times 32$, otteniamo una mappa spaziale che mostra i pattern di attenzione di ogni latente.

\subsubsection{Osservazioni principali}

Dall'analisi delle 24 visualizzazioni generate (una per ogni combinazione di stadio $\times$ esperimento), emergono pattern coerenti:

\begin{enumerate}[label=\arabic*.]
    \item \textbf{Primi stadi} ($t=1$): i pattern di attenzione sono \textbf{ampi e diffusi}, coprendo regioni estese dell'immagine. I latenti catturano informazione globale (colore dominante, luminosit\`a complessiva, contorni grossolani). Questo corrisponde a feature a \textbf{bassa frequenza}.

    \item \textbf{Stadi intermedi} ($t=2,3$): i pattern diventano pi\`u \textbf{focalizzati}, con i latenti che si ``specializzano'' su regioni specifiche dell'immagine. Si osservano pattern che assomigliano a bordi, angoli e texture locali. Questo corrisponde a feature a \textbf{media frequenza}.

    \item \textbf{Ultimi stadi} ($t=4$): i pattern sono molto \textbf{selettivi e localizzati}, con pochi pixel fortemente attesi e il resto dell'immagine quasi ignorato. I latenti catturano dettagli fini e feature discriminanti specifiche per la classe. Questo corrisponde a feature ad \textbf{alta frequenza}.
\end{enumerate}

\begin{notabox}[Analogia con le CNN]
Questa progressione da feature globali a feature locali ricorda la gerarchia delle CNN, dove i primi layer catturano bordi e texture semplici, e i layer pi\`u profondi catturano pattern complessi e specifici per la classe. Tuttavia, c'\`e una differenza fondamentale: nelle CNN questa gerarchia \`e \textbf{imposta dall'architettura} (il campo recettivo cresce layer dopo layer), mentre nel Perceiver emerge \textbf{spontaneamente} dal meccanismo di attenzione iterativa. Il modello ``impara'' da solo che conviene prima guardare il quadro generale e poi concentrarsi sui dettagli.
\end{notabox}

\begin{figure}[H]
\centering
\includegraphics[width=0.8\textwidth]{paper_figures/our_exp1_perceiver_style.png}
\caption{Visualizzazione delle attention maps in stile Perceiver per il nostro esperimento exp1 (baseline Fourier) su CIFAR-10. Ogni riga mostra le mappe di attenzione di uno dei 96 latenti su un'immagine di input.}
\label{fig:our_exp1_attention}
\end{figure}

\begin{figure}[H]
\centering
\includegraphics[width=0.8\textwidth]{paper_figures/our_exp6_perceiver_style.png}
\caption{Visualizzazione delle attention maps per exp6 (Fourier + permutazione). Nonostante i pixel siano permutati, i pattern di attenzione mostrano struttura coerente, confermando che il modello ricostruisce le relazioni spaziali tramite le Fourier features.}
\label{fig:our_exp6_attention}
\end{figure}

Le nostre attention maps (Figure~\ref{fig:our_exp1_attention} e~\ref{fig:our_exp6_attention}) mostrano pattern qualitativamente simili a quelli del paper originale (Figura~\ref{fig:paper_attention_maps}), anche se operano su immagini molto pi\`u piccole ($32 \times 32$ vs $224 \times 224$). Nel caso del baseline Fourier (exp1), si osservano latenti che si specializzano per regioni diverse dell'immagine, con alcuni che catturano lo sfondo e altri che si concentrano sull'oggetto principale. Il confronto con exp6 (Fourier + permutazione) \`e particolarmente istruttivo: nonostante i pixel siano stati mescolati casualmente, le mappe di attenzione rivelano che il modello \textbf{ricostruisce implicitamente} le relazioni spaziali grazie alle Fourier features concatenate a ogni pixel. Questo \`e una conferma visiva diretta dell'invarianza alla permutazione del Perceiver.

\begin{figure}[H]
\centering
\includegraphics[width=0.9\textwidth]{paper_figures/our_attention_comparative.png}
\caption{Analisi comparativa dell'evoluzione dell'attenzione attraverso i nostri esperimenti. Si confrontano i pattern di attenzione tra diversi tipi di PE e configurazioni.}
\label{fig:our_attention_comparative}
\end{figure}

La Figura~\ref{fig:our_attention_comparative} offre un confronto sistematico dell'evoluzione dei pattern di attenzione attraverso i diversi esperimenti. Questa analisi comparativa permette di visualizzare come il tipo di positional encoding e la presenza o assenza di permutazione influenzino la struttura dell'attenzione. Si nota come i pattern siano pi\`u definiti e spazialmente coerenti negli esperimenti con Fourier features, mentre risultano pi\`u diffusi e meno strutturati negli esperimenti senza PE.


%% ------------------------------------------------------------
\subsection{Lezioni apprese e conclusioni sperimentali} \label{sec:lezioni_apprese}

Dopo 17 esperimenti e circa 40 ore di GPU, raccogliamo le lezioni principali del nostro percorso sperimentale.

\subsubsection{Tabella riassuntiva globale}

\begin{table}[htbp]
\centering
\caption{Riepilogo globale di tutti gli esperimenti condotti nel progetto.}
\label{tab:riepilogo_globale}
\small
\begin{tabular}{@{}llcl@{}}
\toprule
\textbf{Dataset} & \textbf{Modello} & \textbf{Miglior risultato} & \textbf{Confronto paper} \\
\midrule
CIFAR-10 (7 exp) & Perceiver & 78.12\% & n/a (paper usa ImageNet) \\
CIFAR-10 & Perceiver IO & 78.20\% & --- \\
ModelNet40 & Perceiver & 84.16\% & 85.7\% ($-1.54$ pp) \\
WikiText-103 (MLM) & Perceiver IO & 82.20\% & --- \\
GLUE (8 task) & Perceiver IO & 75.65\% (best, QQP) & --- \\
\bottomrule
\end{tabular}
\end{table}

\subsubsection{Cosa abbiamo confermato dal paper}

\begin{enumerate}[label=\arabic*.]
    \item \textbf{Il PE \`e cruciale}: senza positional encoding, l'accuratezza crolla di $\sim$10 punti. Il Perceiver \`e costruito per essere agnostico alla struttura, ma ha \emph{bisogno} di sapere dove si trovano gli elementi dell'input.

    \item \textbf{Fourier features $\approx$ Learned PE su input permutato}: la differenza di 0.52 punti \`e trascurabile. Le Fourier features sono preferibili perch\'e non richiedono parametri aggiuntivi e generalizzano a risoluzioni diverse.

    \item \textbf{Invarianza alla permutazione}: il Perceiver con PE adeguato mantiene (o addirittura migliora) le prestazioni quando i pixel vengono permutati, confermando l'assenza di inductive bias strutturali.

    \item \textbf{Generalit\`a multi-dominio}: la stessa architettura funziona su immagini 2D, point cloud 3D e testo, come promesso dal paper.

    \item \textbf{Vicinanza ai risultati del paper}: su ModelNet40, il nostro risultato (84.16\%) \`e molto vicino al paper (85.7\%).
\end{enumerate}

\subsubsection{Cosa abbiamo scoperto di diverso}

\begin{enumerate}[label=\arabic*.]
    \item \textbf{Weight sharing non sempre aiuta}: su CIFAR-10, rimuovere il WS migliora l'accuratezza di 5.36 punti, contrariamente al paper dove il WS migliora le prestazioni su ImageNet. L'effetto del WS dipende dal rapporto tra dimensione del modello e complessit\`a del dataset.

    \item \textbf{La permutazione migliora su CIFAR-10}: un risultato inatteso che probabilmente \`e legato all'interazione tra Fourier features e la piccola dimensione della griglia ($32 \times 32$). Su ImageNet ($224 \times 224$), questo effetto potrebbe non manifestarsi.

    \item \textbf{Variabilit\`a tra run}: esperimenti con iperparametri identici mostrano una variabilit\`a di $\sim$3.5 punti, suggerendo che su dataset piccoli l'inizializzazione casuale ha un ruolo significativo.
\end{enumerate}

\subsubsection{Tabella di confronto sistematico con il paper}

La tabella seguente riassume in modo sistematico il confronto tra i claim del paper originale e i nostri risultati sperimentali.

\begin{table}[htbp]
\centering
\caption{Confronto sistematico: risultati del paper vs nostri risultati.}
\label{tab:confronto_paper_nostro}
\small
\begin{tabular}{@{}p{3cm}p{3.5cm}p{3.5cm}l@{}}
\toprule
\textbf{Finding} & \textbf{Claim del paper} & \textbf{Nostro risultato} & \textbf{Confermato?} \\
\midrule
Importanza del PE & Cruciale per le prestazioni & Confermato: $-10.68$ pp senza PE & S\`i \\
\addlinespace
Fourier $>$ Learned & Gap di 7.1 pp su ImageNet & Gap di soli 0.52 pp su CIFAR-10 & Parzialmente \\
\addlinespace
WS migliora accuracy & +5.1 pp su ImageNet & $-5.36$ pp su CIFAR-10 & \textbf{No} \\
\addlinespace
Invarianza permutazione & Stessa accuracy (78.0\%) & Permutato \textbf{migliore} (+8.43 pp) & Superato \\
\addlinespace
Generalit\`a multi-modale & Funziona su pi\`u domini & CIFAR-10, ModelNet40, WikiText, GLUE & S\`i \\
\addlinespace
ModelNet40 & 85.7\% & 84.16\% ($-1.54$ pp) & S\`i \\
\bottomrule
\end{tabular}
\end{table}

Questa tabella mette in evidenza un punto importante: i risultati del paper non si trasferiscono automaticamente a setup diversi. Il weight sharing, che \`e presentato come un componente chiave del Perceiver nel paper, risulta controproducente nel nostro setup pi\`u piccolo. Allo stesso modo, il gap tra Fourier e Learned PE, drammatico su ImageNet, \`e trascurabile su CIFAR-10. Questo sottolinea l'importanza di \textbf{riprodurre e verificare} i risultati pubblicati, piuttosto che accettarli acriticamente.

\begin{examplebox}[Messaggio finale: cosa portare all'esame]
I nostri esperimenti dimostrano che il Perceiver \`e un'architettura \textbf{genuinamente generale}: la stessa struttura funziona su immagini, point cloud e testo. Le Fourier features sono il componente pi\`u critico, fornendo l'informazione posizionale senza inductive bias strutturali. Il weight sharing \`e una scelta di regolarizzazione il cui beneficio dipende dal regime di training (modello grande/dati tanti vs modello piccolo/dati pochi).

Il risultato pi\`u interessante dal punto di vista scientifico \`e che la permutazione pu\`o \textbf{migliorare} le prestazioni su input di piccola dimensione, probabilmente a causa di fenomeni di aliasing tra le frequenze di Fourier e la griglia regolare dei pixel. Questo suggerisce che la relazione tra positional encoding e struttura dell'input \`e pi\`u complessa di quanto il paper originale lasci intendere.
\end{examplebox}


\newpage
%% ============================================================
%%  SEZIONE 12 --- DOMANDE D'ESAME
%% ============================================================
\section{Domande per la Preparazione dell'Esame} \label{sec:domande}

Dopo aver analizzato sia i risultati del paper originale che i nostri esperimenti, \`e il momento di mettere alla prova la comprensione di tutti i concetti trattati. Questa sezione contiene domande tipiche che potrebbero essere poste all'esame, con tracce di risposta dettagliate. Il consiglio \`e di provare a rispondere a ogni domanda \emph{senza guardare la traccia}, e poi confrontare. Per ogni domanda, cerchiamo di identificare i ``punti chiave'' che il professore si aspetta.

\subsection{Domande di comprensione} \label{sec:domande_comp}

Queste domande coprono i concetti fondamentali del Perceiver. Per ogni domanda, \`e importante non solo \emph{cosa} rispondere, ma anche \emph{come} strutturare la risposta: partire dal problema, spiegare la soluzione, e fornire i dettagli tecnici (formule, dimensioni, numeri).

\begin{enumerate}[label=\textbf{D\arabic*.}]

\item \textbf{Qual \`e la complessit\`a della self-attention standard e perch\'e \`e un problema per input ad alta dimensionalit\`a? Come risolve il Perceiver questo problema?}

\begin{proofbox}[Traccia di risposta]
La self-attention ha complessit\`a $\mathcal{O}(M^2 d)$, dove $M$ \`e il numero di token. Per un'immagine $224 \times 224$, $M = 50{,}176$ e la matrice di attenzione avrebbe $\sim$2.5 miliardi di elementi. Il Perceiver introduce un array latente con $N \ll M$ elementi, usando cross-attention ($\mathcal{O}(MN)$) per comprimere l'input nei latenti, e self-attention ($\mathcal{O}(N^2)$) solo sui latenti. Con $N = 512$, la matrice di cross-attention ha solo $\sim$25.7M elementi. (Si veda la \hyperref[sec:limite_transformer]{Sezione~1.2} e la \hyperref[sec:compl_perceiver]{Sezione~9.2}.)
\end{proofbox}

\item \textbf{Spiega la differenza tra self-attention e cross-attention nel Perceiver, specificando da dove provengono Q, K, V.}

\begin{proofbox}[Traccia di risposta]
Nella self-attention del Latent Transformer (\hyperref[sec:latent_detail]{Sezione~5.7}): $\mathbf{Q} = \mathbf{z}\mathbf{W}_Q$, $\mathbf{K} = \mathbf{z}\mathbf{W}_K$, $\mathbf{V} = \mathbf{z}\mathbf{W}_V$ (tutti dai latenti). Matrice di attenzione: $N \times N$.

Nella cross-attention (\hyperref[sec:ca_detail]{Sezione~5.5}): $\mathbf{Q} = \mathbf{z}\mathbf{W}_Q$ (dai latenti), $\mathbf{K} = \mathbf{x}\mathbf{W}_K$, $\mathbf{V} = \mathbf{x}\mathbf{W}_V$ (dall'input). Matrice di attenzione: $N \times M$.
\end{proofbox}

\item \textbf{Perch\'e le Fourier features sono preferibili alle posizioni apprese nel Perceiver? Che ruolo gioca il criterio di Nyquist?}

\begin{proofbox}[Traccia di risposta]
Le Fourier features (\hyperref[sec:fourier]{Sezione~4}) codificano la posizione con funzioni $\sin$ e $\cos$ a frequenze multiple, fornendo una rappresentazione univoca e ad alta risoluzione di ogni posizione. Le posizioni apprese (70.9\%) perdono 7 punti rispetto alle Fourier (78.0\%) perch\'e non riescono a codificare univocamente $50{,}176$ posizioni. Il criterio di Nyquist richiede frequenza $\geq R/2$ per distinguere posizioni adiacenti su una griglia di $R$ punti.
\end{proofbox}

\item \textbf{Spiega il weight sharing nel Perceiver e la sua analogia con le RNN. Qual \`e il suo effetto su parametri e prestazioni?}

\begin{proofbox}[Traccia di risposta]
La prima cross-attention ha pesi unici; le iterazioni 2--$T$ condividono gli stessi pesi (\hyperref[sec:iterative]{Sezione~3.4}). Analogia RNN (\hyperref[app:rnn]{Appendice~C}): $\mathbf{z}_t = g(\mathbf{z}_{t-1}, \mathbf{x}; \theta_{\text{shared}})$, dove per\`o l'input $\mathbf{x}$ \`e costante. Il WS riduce i parametri da $\sim$326M a $\sim$45M e \textbf{migliora} la validazione da 72.9\% a 78.0\%, agendo come regolarizzazione.
\end{proofbox}

\item \textbf{Perch\'e il Perceiver \`e invariante alla permutazione dell'input mentre ResNet e ViT no?}

\begin{proofbox}[Traccia di risposta]
Il Perceiver opera su pixel individuali con Fourier features \emph{concatenate} (\hyperref[sec:fourier_forward]{Sezione~5.2}): ogni pixel porta con s\'e la propria etichetta posizionale. La cross-attention \`e permutation-invariant (\hyperref[sec:perm_invariance]{Sezione~4.1}): $\sum_j A_{ij} \mathbf{v}_j$ non dipende dall'ordine degli addendi. ResNet perde $-34$ punti perch\'e le convoluzioni con kernel $3 \times 3$ assumono localit\`a spaziale; ViT perde $-15$ punti perch\'e le patch $16 \times 16$ contengono pixel casuali dopo la permutazione (\hyperref[sec:imagenet_perm]{Sezione~10.2}).
\end{proofbox}

\item \textbf{Descrivi il ruolo della residual connection nel Perceiver e spiega perch\'e la rete non potrebbe funzionare senza.}

\begin{proofbox}[Traccia di risposta]
Il Perceiver ha fino a 112 residual connections totali (\hyperref[sec:residual]{Sezione~2.4}): 14 per iterazione $\times$ 8 iterazioni. Senza di esse, il gradiente dovrebbe attraversare $\sim$56 blocchi di trasformazione, attenuandosi esponenzialmente (vanishing gradient). La residual connection garantisce che il gradiente possa fluire direttamente attraverso lo skip ($\frac{\partial \mathbf{y}}{\partial \mathbf{x}} = \mathbf{I} + \frac{\partial f}{\partial \mathbf{x}}$), e permette l'apprendimento incrementale: ogni blocco deve apprendere solo il ``residuo'' (la correzione).
\end{proofbox}

\end{enumerate}

\subsection{Domande aperte e avanzate} \label{sec:domande_aperte}

Queste domande richiedono riflessione critica e capacit\`a di collegare concetti diversi. Sono il tipo di domande che distingue un voto ``buono'' da un voto ``eccellente'' all'esame. Per rispondere bene, \`e fondamentale mostrare di aver compreso non solo \emph{come} funziona il Perceiver, ma anche \emph{perch\'e} sono state fatte certe scelte e \emph{quali} sono le conseguenze.

\begin{enumerate}[resume, label=\textbf{D\arabic*.}]

\item \textbf{Discuti il trade-off tra generalit\`a e prestazioni nel Perceiver. In quali scenari il Perceiver supererebbe le architetture specializzate?}

\begin{proofbox}[Traccia di risposta]
Il Perceiver rinuncia agli inductive bias specifici in favore della generalit\`a. Scenari favorevoli: grandi dataset (dove il modello pu\`o apprendere le strutture dai dati), compiti multi-modali (fusione nativa), modalit\`a non standard (point cloud, grafi), transfer cross-modale. Limite: il bottleneck latente ($N$ fisso) pu\`o perdere informazione fine. Su ModelNet40, il gap con PointNet++ (\hyperref[sec:modelnet]{Sezione~10.4}) \`e di 6.2 punti, mostrando che per compiti con forte struttura geometrica locale, gli inductive bias specializzati restano vantaggiosi.
\end{proofbox}

\item \textbf{Come potresti estendere il Perceiver per produrre output strutturati (es.\ segmentazione)?}

\begin{proofbox}[Traccia di risposta]
Il global average pooling (\hyperref[sec:pooling]{Sezione~5.9}) perde la struttura spaziale. Soluzione (anticipando il Perceiver IO): introdurre un decoder cross-attention dove le query provengono dall'output desiderato ($M$ query, una per pixel) e K,V dai latenti. Questo ``decomprime'' l'informazione da $N$ latenti a $M$ output. La matrice di attenzione del decoder sarebbe $\mathbb{R}^{M \times N}$, costo $\mathcal{O}(MN)$, mantenendo la complessit\`a lineare.
\end{proofbox}

\item \textbf{Spiega il backward pass completo della cross-entropy loss combinata con softmax. Perch\'e il risultato \`e cos\`i elegante?}

\begin{proofbox}[Traccia di risposta]
La dimostrazione completa \`e nella \hyperref[sec:grad_loss]{Sezione~7.1}. Si parte da $\mathcal{L} = -\sum_k y_k \log p_k$ e $p_k = \text{softmax}(\text{logit}_k)$. Usando la chain rule e la derivata della softmax $\frac{\partial p_i}{\partial \text{logit}_j} = p_i(\delta_{ij} - p_j)$, si ottiene $\frac{\partial \mathcal{L}}{\partial \text{logit}_k} = p_k - y_k$. L'eleganza deriva dalla ``cancellazione'' tra numeratore e denominatore della softmax. Questo risultato semplifica enormemente l'implementazione del backward pass.
\end{proofbox}

\item \textbf{Calcola il numero di operazioni nella cross-attention del Perceiver per ImageNet e confrontalo con la self-attention di un Transformer.}

\begin{proofbox}[Traccia di risposta]
\textbf{Cross-attention}: $\mathbf{Q} \cdot \mathbf{K}^\top$ ha dimensione $N \times M$, costo $\mathcal{O}(N \cdot M \cdot d_k) = 512 \times 50{,}176 \times 1024 \approx 2.63 \times 10^{10}$. Poi $\mathbf{A} \cdot \mathbf{V}$ ha costo $\mathcal{O}(N \cdot M \cdot d_v)$, simile. Totale: $\sim 5.3 \times 10^{10}$ operazioni.

\textbf{Self-attention su pixel}: $\mathbf{Q} \cdot \mathbf{K}^\top$ avrebbe dimensione $M \times M$, costo $\mathcal{O}(M^2 \cdot d_k) = 50{,}176^2 \times 1024 \approx 2.58 \times 10^{12}$. Il Perceiver \`e $\sim 50\times$ pi\`u efficiente.

\textbf{Self-attention nel Latent Transformer}: $\mathcal{O}(N^2 \cdot D) = 512^2 \times 1024 \approx 2.68 \times 10^{8}$. Molto pi\`u economica della cross-attention, perch\'e $N \ll M$.
\end{proofbox}

\item \textbf{Perch\'e il Perceiver usa GELU anzich\'e ReLU? Quali sono le conseguenze per il backward pass?}

\begin{proofbox}[Traccia di risposta]
La GELU (\hyperref[sec:gelu]{Sezione~6.2}) \`e preferita perch\'e: (1) \`e liscia (derivata continua ovunque), facilitando l'ottimizzazione; (2) non soffre di dying neurons; (3) ha un meccanismo di self-gating ($x \cdot \Phi(x)$). Per il backward pass, la derivata della GELU \`e $\text{GELU}'(x) = \Phi(x) + x \cdot \phi(x)$: non \`e mai esattamente zero, quindi il gradiente fluisce sempre. Con la ReLU, il 10--40\% dei neuroni con input negativo avrebbe gradiente nullo, bloccando l'apprendimento.
\end{proofbox}

\item \textbf{Spiega la differenza tra il positional encoding del Transformer e le Fourier features del Perceiver. Perch\'e il Perceiver \emph{concatena} anzich\'e \emph{sommare}?}

\begin{proofbox}[Traccia di risposta]
Il Transformer originale \emph{somma} il PE ai token embedding: $\mathbf{z}_i = \mathbf{x}_i + \text{PE}(i)$. Questo \`e possibile perch\'e entrambi hanno la stessa dimensione $d$. Il Perceiver \emph{concatena}: $\tilde{\mathbf{x}}_i = [\mathbf{x}_i; \text{PE}(i)]$, aumentando la dimensionalit\`a da $C$ a $C + d_{\text{PE}}$. La concatenazione \`e fondamentale per l'invarianza alla permutazione: con la somma, permutare i token e poi sommare PE darebbe PE diversi (il PE dipende dall'indice); con la concatenazione, ogni pixel ``porta con s\'e'' le proprie coordinate, indipendentemente dall'ordine in cui viene processato.
\end{proofbox}

\item \textbf{Descrivi il ruolo del trust ratio in LAMB e spiega perch\'e \`e particolarmente importante per il Perceiver.}

\begin{proofbox}[Traccia di risposta]
Il trust ratio $\phi = \|\theta_t\|_2 / \|\mathbf{r}_t\|_2$ (\hyperref[sec:ottimizzatori]{Sezione~8.2}) normalizza l'aggiornamento layer per layer. Nel Perceiver, i latenti iniziali hanno $\|\theta\| \approx 14.5$ (piccoli, $\sigma = 0.02$), mentre i pesi di proiezione hanno $\|\theta\| \approx 30$ (Xavier init). Senza trust ratio, un learning rate globale sarebbe o troppo aggressivo per i latenti (li destabilizza) o troppo cauto per le proiezioni (convergenza lenta). LAMB riscala automaticamente, garantendo che l'aggiornamento sia proporzionale alla scala di ogni layer.
\end{proofbox}

\end{enumerate}


\newpage
%% ============================================================
%%  SEZIONE 13 --- CONCLUSIONI
%% ============================================================
\section{Conclusioni} \label{sec:conclusioni}

Dopo aver analizzato il paper originale, i risultati degli autori e i nostri esperimenti indipendenti, facciamo un passo indietro e guardiamo il quadro complessivo. Il Perceiver rappresenta un passo significativo verso architetture di deep learning \textbf{generali e agnostiche rispetto alla modalit\`a}. La scommessa fondamentale del paper \`e: ``possiamo rinunciare a tutti gli inductive bias specifici per ogni dominio e costruire un'unica architettura che funziona su tutto?''. La risposta, come abbiamo visto, \`e un s\`i qualificato. I contributi principali sono:

Per un'analisi completa dell'evoluzione di questa architettura, si rimanda alla \textbf{Parte~II} di questa dispensa, che tratta in dettaglio il \textbf{Perceiver IO} (Jaegle et al., 2022) e il suo meccanismo di output basato su query, che supera i limiti del Perceiver originale discussi nella Sezione~\ref{sec:limiti}.


\begin{enumerate}[label=\arabic*.]
    \item \textbf{Latent bottleneck} (\hyperref[sec:latent_bottleneck]{Sezione~1.3}): l'introduzione di un array latente con $N \ll M$ riduce la complessit\`a da $\mathcal{O}(M^2)$ a $\mathcal{O}(MN)$.

    \item \textbf{Cross-attention asimmetrica} (\hyperref[sec:cross_attention]{Sezione~3.2}): le query provengono dai latenti ($N$ righe) e le chiavi/valori dall'input ($M$ righe), comprimendo l'informazione.

    \item \textbf{Iterazione e weight sharing} (\hyperref[sec:iterative]{Sezione~3.4}): l'alternanza ripetuta con pesi condivisi riduce i parametri di un ordine di grandezza e funge da regolarizzazione.

    \item \textbf{Fourier features} (\hyperref[sec:fourier]{Sezione~4}): encoding posizionale multi-dimensionale senza inductive bias strutturali.

    \item \textbf{Generalit\`a verificata} (\hyperref[sec:esperimenti]{Sezione~10}): stessa architettura competitiva su immagini (78.0\% su ImageNet), audio (38.4 mAP su AudioSet), e point cloud 3D (85.7\% su ModelNet40), con prestazioni identiche su input permutati.
\end{enumerate}

\subsection{Riepilogo del forward e backward pass} \label{sec:riepilogo_fb}

Per facilitare il ripasso, raccogliamo in due tabelle tutte le dimensioni del forward pass e tutti i parametri del modello. Queste tabelle sono il ``cheat sheet'' definitivo per l'esame --- raccomandiamo di stamparle e tenerle a portata di mano durante lo studio.

\begin{table}[htbp]
\centering
\caption{Riepilogo dimensionale completo del forward pass del Perceiver per ImageNet.}
\begin{tabular}{@{}clll@{}}
\toprule
\textbf{Passo} & \textbf{Operazione} & \textbf{Input} & \textbf{Output} \\
\midrule
A1 & Flatten & $224 \times 224 \times 3$ & $50{,}176 \times 3$ \\
A2 & Fourier features & $50{,}176 \times 3$ & $50{,}176 \times 261$ \\
A3 & Proiezione lineare & $50{,}176 \times 261$ & $50{,}176 \times 512$ \\
A4 & Latent array (appreso) & --- & $512 \times 1024$ \\
A5 & Pre-LN (latenti) & $512 \times 1024$ & $512 \times 1024$ \\
A5 & Pre-LN (input) & $50{,}176 \times 512$ & $50{,}176 \times 512$ \\
A5 & Q = z LN $\cdot$ W$_Q$ & $512 \times 1024 \cdot 1024 \times 1024$ & $512 \times 1024$ \\
A5 & K = x LN $\cdot$ W$_K$ & $50{,}176 \times 512 \cdot 512 \times 1024$ & $50{,}176 \times 1024$ \\
A5 & V = x LN $\cdot$ W$_V$ & $50{,}176 \times 512 \cdot 512 \times 1024$ & $50{,}176 \times 1024$ \\
A5 & Score = QK$^\top/\sqrt{d_k}$ & $512 \times 1024 \cdot 1024 \times 50{,}176$ & $512 \times 50{,}176$ \\
A5 & A = softmax(Score) & $512 \times 50{,}176$ & $512 \times 50{,}176$ \\
A5 & Out = A $\cdot$ V & $512 \times 50{,}176 \cdot 50{,}176 \times 1024$ & $512 \times 1024$ \\
A5 & Residual & $512 \times 1024$ & $512 \times 1024$ \\
A6 & MLP (expand) & $512 \times 1024$ & $512 \times 4096$ \\
A6 & GELU & $512 \times 4096$ & $512 \times 4096$ \\
A6 & MLP (compress) & $512 \times 4096$ & $512 \times 1024$ \\
A6 & Residual & $512 \times 1024$ & $512 \times 1024$ \\
A7 & Latent SA $\times 6$ & $512 \times 1024$ & $512 \times 1024$ \\
A8 & $\times T = 8$ iter. & --- & --- \\
A9 & Avg pool & $512 \times 1024$ & $1024$ \\
A10 & Classificatore & $1024$ & $1000$ \\
A11 & Softmax & $1000$ & $1000$ \\
A12 & Cross-entropy & $1000$ & scalare $\mathcal{L}$ \\
\bottomrule
\end{tabular}
\end{table}

\begin{table}[htbp]
\centering
\caption{Riepilogo dei parametri apprendibili del Perceiver (con weight sharing).}
\begin{tabular}{@{}llr@{}}
\toprule
\textbf{Componente} & \textbf{Dimensioni principali} & \textbf{Parametri} \\
\midrule
Latent array $\mathbf{z}_0$ & $512 \times 1024$ & 0.5M \\
Input embedding & $261 \times 512 + 512$ & 0.13M \\
Cross-attention (unica) & W$_Q$, W$_K$, W$_V$, W$_O$, MLP, LN & $\sim$11.6M \\
Cross-attention (condivisa) & Stesso set di parametri & $\sim$11.6M \\
Latent Transformer ($L=6$) & 6 $\times$ (SA + MLP + LN) & $\sim$17M \\
Classificatore & $1024 \times 1000 + 1000$ & $\sim$1.0M \\
\midrule
\textbf{Totale} & & $\sim$\textbf{45M} \\
\bottomrule
\end{tabular}
\end{table}


\subsection{Limiti del Perceiver} \label{sec:limiti}

Nonostante i suoi contributi, il Perceiver presenta alcuni limiti importanti:

\begin{enumerate}[label=(\alph*)]
    \item \textbf{Bottleneck informativo}: l'array latente con $N = 512$ vettori deve codificare \emph{tutta} l'informazione di un input con $M = 50{,}176$ elementi. Questo comporta necessariamente una perdita di informazione fine, limitando le prestazioni su compiti che richiedono dettaglio pixel-level (es.\ segmentazione).

    \item \textbf{Solo output globale}: il global average pooling (\hyperref[sec:pooling]{Sezione~5.9}) produce un singolo vettore, impedendo output strutturati (segmentazione, detection, generazione). Questo limite fondamentale \`e analizzato in dettaglio e risolto dal \textbf{Perceiver IO} nella \hyperref[sec:perceiver_io_preview]{Sezione~\ref*{sec:perceiver_io_preview}}, che introduce un decoder cross-attention con output query specifiche per task.

    \item \textbf{Costo della cross-attention}: nonostante la complessit\`a sia lineare in $M$, la cross-attention su $50{,}176$ pixel viene ripetuta $T = 8$ volte. Il costo totale ($\sim 2.2 \times 10^{11}$ operazioni) \`e significativo, anche se inferiore al Transformer diretto.

    \item \textbf{Gap con architetture specializzate}: su ModelNet40, il Perceiver perde 6.2 punti contro PointNet++ (\hyperref[sec:modelnet]{Sezione~10.4}). Su ImageNet, raggiunge ResNet-50 ma senza superarlo significativamente. Le architetture specializzate restano superiori quando gli inductive bias sono appropriati.

    \item \textbf{Scelta degli iperparametri}: $N$ (numero di latenti), $T$ (iterazioni), $L$ (layer SA) sono iperparametri cruciali. Non esiste una regola automatica per sceglierli: dipendono dalla complessit\`a del task e dalla dimensione dell'input.
\end{enumerate}

\subsection{Posizione del Perceiver nel panorama delle architetture} \label{sec:panorama_architetture}

Per capire davvero il contributo del Perceiver, \`e utile posizionarlo nel contesto pi\`u ampio delle architetture di deep learning. Dove si colloca rispetto alle CNN, agli RNN, ai Transformer e al ViT?

\begin{table}[htbp]
\centering
\caption{Confronto riassuntivo delle principali architetture di deep learning.}
\begin{tabular}{@{}p{2.2cm}p{2.5cm}p{2cm}p{1.5cm}p{2.5cm}p{2cm}@{}}
\toprule
\textbf{Architettura} & \textbf{Inductive bias} & \textbf{Complessit\`a} & \textbf{Input} & \textbf{Dipendenze globali} & \textbf{Multimodale} \\
\midrule
CNN & Localit\`a, traslaz. & $\mathcal{O}(Mk^2)$ & Griglia 2D & Dopo molti layer & No \\
RNN/LSTM & Sequenzialit\`a & $\mathcal{O}(Md^2)$ & Sequenze & Degradate (VG) & No \\
Transformer & Nessuno (+ PE) & $\mathcal{O}(M^2d)$ & Token & S\`i (layer 1) & Con design \\
ViT & Patch 2D & $\mathcal{O}(P^2d)$ & Patch immagine & S\`i (layer 1) & No \\
PointNet++ & Geometria locale & $\mathcal{O}(M\log M)$ & Punti 3D & Gerarchica & No \\
\textbf{Perceiver} & \textbf{Nessuno} & $\mathcal{O}(MNd)$ & \textbf{Qualsiasi} & \textbf{S\`i (via latenti)} & \textbf{Nativo} \\
\bottomrule
\end{tabular}
\end{table}

Il Perceiver \`e l'architettura con il minor numero di inductive bias e la maggiore generalit\`a, al costo di un gap prestazionale rispetto alle architetture specializzate (es.\ $-6$ punti vs PointNet++ su ModelNet40). La scommessa del Perceiver \`e che, con dataset sufficientemente grandi, il modello pu\`o \emph{apprendere} dai dati le stesse strutture che le architetture specializzate ``incorporano'' nel design.

\begin{notabox}[Il trade-off fondamentale: inductive bias vs generalit\`a]
Le architetture nel deep learning si collocano su uno spettro:
\begin{itemize}
    \item \textbf{Massimo inductive bias} (CNN, PointNet++): prestazioni eccellenti sul proprio dominio, ma non trasferibili ad altri. Richiedono meno dati (il design codifica conoscenza a priori).
    \item \textbf{Inductive bias moderato} (ViT): patch 2D come unica assunzione. Prestazioni eccellenti con grandi dataset; sotto-performano con pochi dati rispetto alle CNN.
    \item \textbf{Nessun inductive bias} (Perceiver): massima generalit\`a. Richiedono pi\`u dati per raggiungere le stesse prestazioni. Brillano su compiti multi-modali e su input non standard.
\end{itemize}

La tendenza del campo \`e verso architetture con meno inductive bias ma pi\`u dati e compute: GPT-4, DALL-E 3, Gemini sono tutti modelli ``generalisti'' che apprendono le strutture dai dati anzich\'e dal design. Il Perceiver \`e un precursore di questa tendenza nel dominio della percezione.
\end{notabox}


\subsection{Evoluzione: dal Perceiver al Perceiver IO} \label{sec:perceiver_io_preview}

Prima di concludere, guardiamo brevemente al futuro. Il Perceiver ha un limite importante: pu\`o solo produrre un output ``globale'' (un vettore, una classe). Ma cosa succede se vogliamo un output strutturato, come una segmentazione pixel per pixel? Il Perceiver IO (Jaegle et al., 2022) risolve questo problema estendendo il Perceiver in due direzioni:
\begin{itemize}
    \item \textbf{Output decoder}: aggiunge una cross-attention decoder dove le query sono un ``output array'' $\mathbf{o} \in \mathbb{R}^{O \times D_o}$ e K, V provengono dai latenti. Questo permette output strutturati con $O$ elementi (es.\ $O = M$ per segmentazione, $O = 1$ per classificazione).
    \item \textbf{Proiezioni flessibili}: le dimensioni delle proiezioni Q, K, V possono variare tra encoder e decoder, permettendo di disaccoppiare completamente l'architettura dalla dimensionalit\`a di input e output.
\end{itemize}

\begin{formulabox}[Decoder cross-attention del Perceiver IO]
\[
    \text{output} = \text{softmax}\!\left(\frac{\mathbf{o}\mathbf{W}_Q^{\text{dec}} (\mathbf{z}_T \mathbf{W}_K^{\text{dec}})^\top}{\sqrt{d_k}}\right) \mathbf{z}_T \mathbf{W}_V^{\text{dec}} \in \mathbb{R}^{O \times D_o}
\]
dove $\mathbf{o} \in \mathbb{R}^{O \times D_o}$ \`e l'output query array (appreso o derivato dal task).

\textbf{Complessit\`a}: $\mathcal{O}(O \cdot N \cdot D)$. Per $O = M$ (segmentazione): $\mathcal{O}(MN)$, stessa complessit\`a dell'encoder.
\end{formulabox}

Il Perceiver IO dimostra che l'idea del latent bottleneck \`e generalizzabile: non solo per la \emph{lettura} dell'input (encoder cross-attention), ma anche per la \emph{scrittura} dell'output (decoder cross-attention). Questo schema encoder-latent-decoder \`e diventato un paradigma influente nell'architettura dei modelli multi-modali.

\begin{notabox}[In parole semplici: dal Perceiver al Perceiver IO]
Detto in parole povere: il Perceiver originale \`e come un lettore che pu\`o leggere qualsiasi libro e rispondere ``di che genere \`e?'' (classificazione). Il Perceiver IO \`e come un lettore che pu\`o anche \emph{riassumere} il libro capitolo per capitolo, \emph{tradurlo} in un'altra lingua, o \emph{rispondere a domande specifiche} su qualsiasi pagina. La differenza \`e il decoder: invece di comprimere tutto in un unico vettore, il Perceiver IO pu\`o ``scrivere'' output strutturati di qualsiasi dimensione, usando i latenti come fonte di informazione.
\end{notabox}

\begin{notabox}[Dal Perceiver IO ai modelli attuali]
L'idea del Perceiver (latent bottleneck + cross-attention) ha influenzato molti sviluppi successivi:
\begin{itemize}
    \item \textbf{Flamingo} (DeepMind, 2022): usa cross-attention perceiver-like per fondere immagini e testo in un modello multimodale.
    \item \textbf{DALL-E 2 / Stable Diffusion}: usano encoder simili al Perceiver per condizionare la generazione di immagini.
    \item \textbf{Gato} (DeepMind, 2022): agente generalista che usa un'unica architettura per diverse modalit\`a e compiti, ispirandosi al principio del Perceiver.
\end{itemize}
\end{notabox}

\subsection{Riepilogo schematico} \label{sec:riepilogo}

Per concludere, raccogliamo in due tabelle le scelte architetturali del Perceiver e le formule chiave. Queste tabelle sono un ottimo punto di partenza per il ripasso prima dell'esame.

\begin{table}[htbp]
\centering
\caption{Riepilogo delle scelte architetturali del Perceiver e le loro motivazioni.}
\begin{tabular}{@{}p{3.5cm}p{4cm}p{5cm}@{}}
\toprule
\textbf{Scelta} & \textbf{Alternativa} & \textbf{Motivazione} \\
\midrule
Cross-attention & Self-attention su tutto & $\mathcal{O}(MN)$ vs $\mathcal{O}(M^2)$ \\
Fourier features concat. & PE sommato (Transformer) & Invarianza alla permutazione \\
Weight sharing & Pesi separati per iterazione & Regolarizzazione (72.9\% $\to$ 78.0\%) \\
Pre-norm & Post-norm & Stabilit\`a per reti profonde (56 blocchi) \\
GELU & ReLU & No dying neurons, gradiente liscio \\
LAMB & Adam / SGD & Stabilit\`a con scale diverse tra layer \\
Global avg pool & Token CLS & Nessun parametro aggiuntivo \\
Latenti apprendibili & Proiezione dell'input & Indipendenza dall'input specifico \\
\bottomrule
\end{tabular}
\end{table}

\begin{notabox}[Messaggio chiave per l'esame]
Il Perceiver dimostra che la complessit\`a quadratica della self-attention non \`e un limite intrinseco: tramite il disaccoppiamento tra input ($M$ elementi) e rappresentazione ($N$ latenti), si ottiene complessit\`a lineare in $M$ mantenendo la capacit\`a di catturare dipendenze globali. L'architettura rinuncia consapevolmente a tutti gli inductive bias specifici per dominio, affidandosi interamente all'attenzione e all'encoding posizionale.

Per l'esame, \`e essenziale saper descrivere: (1) l'intero \hyperref[sec:forward]{forward pass} con le dimensioni di ogni tensore, (2) il meccanismo di \hyperref[sec:ca_detail]{cross-attention} e perch\'e \`e lineare in $M$, (3) il ruolo delle \hyperref[sec:fourier]{Fourier features} e il confronto con il PE dei Transformer, (4) il \hyperref[sec:iterative]{weight sharing} e la sua analogia con le RNN, (5) l'esperimento di \hyperref[sec:imagenet_perm]{ImageNet permutato}, (6) il \hyperref[sec:backward]{backward pass} almeno per la loss, il classificatore e il pooling, (7) le \hyperref[sec:attivazioni]{funzioni di attivazione} GELU vs ReLU, (8) gli \hyperref[sec:ottimizzatori]{ottimizzatori} (SGD vs Adam vs LAMB).
\end{notabox}

\subsection{Formule chiave da ricordare} \label{sec:formule_chiave}

Chiudiamo con le 10 formule pi\`u importanti da ricordare per l'esame. Raccomandiamo di saperle scrivere a memoria e di saper spiegare il significato di ogni simbolo.

\begin{formulabox}[Le 10 formule fondamentali del Perceiver]

\textbf{1. Complessit\`a} (\hyperref[sec:compl_perceiver]{Sez.~9.2}): $\mathcal{O}(MN + N^2) \ll \mathcal{O}(M^2)$

\textbf{2. Cross-attention} (\hyperref[sec:cross_attention]{Sez.~3.2}): $\text{CA}(\mathbf{z}, \mathbf{x}) = \text{softmax}\!\left(\frac{\mathbf{z}\mathbf{W}_Q (\mathbf{x}\mathbf{W}_K)^\top}{\sqrt{d_k}}\right) \mathbf{x}\mathbf{W}_V$

\textbf{3. Fourier features} (\hyperref[sec:fourier_encoding]{Sez.~4.2}): $\text{PE}(x_d) = [x_d, \sin(f_1\pi x_d), \cos(f_1\pi x_d), \ldots]$, dimensione $= P(2K+1)$

\textbf{4. LayerNorm} (\hyperref[sec:layernorm]{Sez.~2.3}): $\text{LN}(\mathbf{h}) = \boldsymbol{\gamma} \odot \frac{\mathbf{h} - \mu}{\sqrt{\sigma^2 + \epsilon}} + \boldsymbol{\beta}$

\textbf{5. Residual} (\hyperref[sec:residual]{Sez.~2.4}): $\mathbf{y} = \mathbf{x} + f(\mathbf{x})$, gradiente: $\frac{\partial \mathbf{y}}{\partial \mathbf{x}} = \mathbf{I} + \frac{\partial f}{\partial \mathbf{x}}$

\textbf{6. GELU} (\hyperref[sec:gelu]{Sez.~6.2}): $\text{GELU}(x) = x \cdot \Phi(x)$, derivata: $\Phi(x) + x\phi(x)$

\textbf{7. Cross-entropy loss} (\hyperref[sec:loss]{Sez.~5.12}): $\mathcal{L} = -\log(p_c)$

\textbf{8. Gradiente softmax + CE} (\hyperref[sec:grad_loss]{Sez.~7.1}): $\frac{\partial \mathcal{L}}{\partial \text{logit}_k} = p_k - y_k$

\textbf{9. Weight sharing} (\hyperref[sec:grad_ws]{Sez.~7.8}): $\frac{\partial \mathcal{L}}{\partial \theta_{\text{shared}}} = \sum_{t=2}^{T} \frac{\partial \mathcal{L}}{\partial \theta^{(t)}}$

\textbf{10. Average pooling} (\hyperref[sec:pooling]{Sez.~5.9}): $\mathbf{h} = \frac{1}{N}\sum_{i=1}^{N} \mathbf{z}_{T,i}$, gradiente: $\frac{\partial \mathcal{L}}{\partial \mathbf{z}_{T,i}} = \frac{1}{N} \frac{\partial \mathcal{L}}{\partial \mathbf{h}}$
\end{formulabox}


\newpage
%% ============================================================
%%  APPENDICI --- FONDAMENTI DI DEEP LEARNING
%% ============================================================
\newpage
\part*{Parte II --- Perceiver IO (Jaegle et al., 2022)}
\addcontentsline{toc}{part}{Parte II --- Perceiver IO}



%% ============================================================
%%  SEZIONE 1 --- MOTIVAZIONE
%% ============================================================
\section{Motivazione: dal Perceiver al Perceiver IO}

\begin{figure}[H]
\centering
\includegraphics[width=0.85\textwidth]{paper_figures/perceiver_io_fig1_domains.png}
\caption{Il Perceiver IO pu\`o essere applicato a domini con grande variet\`a di input e output (Jaegle et al., 2022, Figura 1): comprensione del linguaggio multi-task, visione densa (optical flow), autoencoding multimodale (video+audio+label), e output simbolici come StarCraft II. La stessa architettura gestisce tutti questi task.}
\label{fig:perceiver_io_domains}
\end{figure}

\subsection{Limiti del Perceiver originale}

Come abbiamo visto nella Parte I (\hyperref[sec:limiti]{Sezione~\ref*{sec:limiti}} e \hyperref[sec:pooling]{Sezione~\ref*{sec:pooling}}), il Perceiver originale comprime l'input in un array latente e poi estrae un singolo vettore tramite average pooling. Questo funziona per la classificazione, ma \`e inadeguato per output strutturati. In questa sezione approfondiamo perch\'e questo limite \`e cos\`i importante e come il Perceiver IO lo risolve.

Il problema centrale \`e che molti task reali richiedono output \textbf{strutturati}, non un singolo vettore:

\begin{itemize}
    \item \textbf{Optical flow}: l'output deve essere un vettore di spostamento 2D \emph{per ogni pixel} dell'immagine
    \item \textbf{Modelli di linguaggio}: l'output \`e una sequenza di token, uno per ogni posizione
    \item \textbf{Segmentazione semantica}: l'output \`e una mappa di classi per-pixel
    \item \textbf{Ricostruzione multimodale}: l'output pu\`o includere video, audio e label simultaneamente
    \item \textbf{Object detection}: l'output \`e un insieme di bounding box con classi associate
    \item \textbf{Giochi strategici}: l'output comprende azioni per centinaia di entit\`a diverse
\end{itemize}

\begin{notabox}[Il problema centrale]
Il Perceiver originale comprime \emph{tutto} l'input in un array latente $\mathbf{z} \in \mathbb{R}^{N \times D}$ e poi estrae un \emph{singolo vettore} tramite average pooling. Ma se il task richiede un output con $O$ elementi (es.\ $O = 224 \times 224 = 50{,}176$ per la segmentazione), come si fa a ``decomprimere'' l'informazione dal latente a uno spazio cos\`i grande?
\end{notabox}

\subsection{Output strutturati: la sfida}

In un Transformer standard encoder-decoder (come quello di Vaswani et al., 2017 per la traduzione), il decoder produce output strutturati attraverso un meccanismo \emph{autoregressivo}: genera un token alla volta, usando attenzione causale. Questo approccio ha due problemi:

\begin{enumerate}
    \item \textbf{Sequenzialit\`a}: la generazione autoregressiva \`e intrinsecamente sequenziale, non parallelizzabile
    \item \textbf{Specificit\`a}: il decoder \`e progettato per un tipo specifico di output (sequenze di token)
\end{enumerate}

Il Perceiver IO propone una soluzione elegante e completamente diversa: un meccanismo di output basato su \textbf{query}, in cui ogni punto dell'output ``interroga'' indipendentemente l'array latente per estrarre l'informazione di cui ha bisogno.

\subsection{Il paradigma Read-Process-Write}

L'intuizione principale del Perceiver IO pu\`o essere riassunta con il paradigma \textbf{Read-Process-Write}:

\begin{enumerate}
    \item \textbf{Read (Encode)}: leggi l'input e comprimi l'informazione nell'array latente tramite cross-attention
    \item \textbf{Process}: processa l'informazione latente con self-attention (il ``ragionamento'' avviene qui)
    \item \textbf{Write (Decode)}: scrivi l'output usando query che interrogano l'array latente tramite cross-attention
\end{enumerate}

\begin{center}
\begin{tikzpicture}[
    block/.style={draw, rounded corners, minimum width=3.8cm, minimum height=1.4cm, align=center, font=\small\bfseries},
    arr/.style={-{Stealth[length=3mm]}, thick},
    label/.style={font=\footnotesize\itshape, text=gray!70!black}
]
    % Input
    \node[block, fill=red!10, draw=red!60!black] (input) at (0,0) {Input Array\\$\mathbf{x} \in \mathbb{R}^{M \times C}$};
    % Encoder
    \node[block, fill=orange!15, draw=orange!60!black] (enc) at (5.5,0) {Encoder\\(Cross-Attention)};
    % Latent
    \node[block, fill=blue!15, draw=blue!60!black] (lat) at (5.5,-3) {Processor\\(Self-Attention $\times L$)};
    % Decoder
    \node[block, fill=green!15, draw=green!60!black] (dec) at (5.5,-6) {Decoder\\(Cross-Attention)};
    % Output queries
    \node[block, fill=purple!10, draw=purple!60!black] (queries) at (0,-6) {Output Queries\\$\mathbf{q} \in \mathbb{R}^{O \times E}$};
    % Output
    \node[block, fill=teal!10, draw=teal!60!black] (output) at (11,-6) {Output Array\\$\mathbf{y} \in \mathbb{R}^{O \times F}$};

    % Arrows
    \draw[arr] (input) -- (enc) node[midway, above, label] {Read};
    \draw[arr] (enc) -- (lat) node[midway, right, label] {$\mathbf{z}^{(0)} \in \mathbb{R}^{N \times D}$};
    \draw[arr] (lat) -- (dec) node[midway, right, label] {$\mathbf{z}^{(L)} \in \mathbb{R}^{N \times D}$};
    \draw[arr] (queries) -- (dec) node[midway, above, label] {Write};
    \draw[arr] (dec) -- (output) node[midway, above, label] {Output};

    % Annotations
    \node[font=\scriptsize, text=red!60!black, below=0.1cm of input] {Dipende dal task};
    \node[font=\scriptsize, text=purple!60!black, below=0.1cm of queries] {Dipende dal task};
    \node[font=\scriptsize, text=blue!60!black, right=0.1cm of lat] {\begin{tabular}{l}Iperparametro\\$N \ll M, O$\end{tabular}};
\end{tikzpicture}
\end{center}

\begin{notabox}[In parole semplici: il paradigma Read-Process-Write]
Immagina di essere un \textbf{giornalista} che deve scrivere un articolo:
\begin{enumerate}
    \item \textbf{Read (Leggere)}: leggi tutte le fonti (interviste, documenti, dati) e prendi appunti su un quaderno. Il quaderno ha solo $N$ pagine (i latenti), molto meno delle fonti originali. Devi decidere cosa \`e importante e cosa omettere.
    \item \textbf{Process (Elaborare)}: rileggi i tuoi appunti pi\`u volte ($L$ volte), collegando le informazioni tra le pagine. ``Ah, questo dato della pagina 3 si collega a quell'intervista della pagina 7!'' Questo \`e il ragionamento.
    \item \textbf{Write (Scrivere)}: quando qualcuno ti fa una domanda specifica (output query), consulti i tuoi appunti e rispondi. Se ti chiedono ``Di cosa parlava il ministro?'', cerchi nelle pagine rilevanti. Se ti chiedono ``Qual era il dato economico?'', cerchi in pagine diverse. \textbf{Domande diverse $\to$ risposte diverse}, anche se gli appunti sono gli stessi.
\end{enumerate}
\end{notabox}

\subsection{Confronto visivo: Perceiver vs.\ Perceiver IO}

\begin{center}
\begin{tikzpicture}[
    block/.style={draw, rounded corners, minimum width=2.5cm, minimum height=1cm, align=center, font=\scriptsize\bfseries},
    arr/.style={-{Stealth[length=2.5mm]}, thick},
    title/.style={font=\normalsize\bfseries, text=black}
]
    % --- Perceiver (left) ---
    \node[title] at (2.5,4.5) {Perceiver (2021)};

    \node[block, fill=red!10] (p_in) at (2.5,3.5) {Input\\$\mathbf{x} \in \mathbb{R}^{M \times C}$};
    \node[block, fill=orange!15] (p_enc) at (2.5,2) {Cross-Att.\\(Encoder)};
    \node[block, fill=blue!15] (p_lat) at (2.5,0.5) {Self-Att. $\times L$\\(Processor)};
    \node[block, fill=gray!15] (p_pool) at (2.5,-1) {Average Pool};
    \node[block, fill=gray!15] (p_mlp) at (2.5,-2.5) {MLP};
    \node[block, fill=teal!10] (p_out) at (2.5,-4) {Classe\\$\mathbf{y} \in \mathbb{R}^{K}$};

    \draw[arr] (p_in) -- (p_enc);
    \draw[arr] (p_enc) -- (p_lat);
    \draw[arr] (p_lat) -- (p_pool);
    \draw[arr] (p_pool) -- (p_mlp);
    \draw[arr] (p_mlp) -- (p_out);

    % Limitation
    \node[font=\scriptsize, text=red!70!black, align=center] at (2.5,-5) {Solo classificazione!\\Output: singolo vettore};

    % --- Perceiver IO (right) ---
    \node[title] at (9,4.5) {Perceiver IO (2022)};

    \node[block, fill=red!10] (io_in) at (9,3.5) {Input\\$\mathbf{x} \in \mathbb{R}^{M \times C}$};
    \node[block, fill=orange!15] (io_enc) at (9,2) {Cross-Att.\\(Encoder)};
    \node[block, fill=blue!15] (io_lat) at (9,0.5) {Self-Att. $\times L$\\(Processor)};
    \node[block, fill=green!15] (io_dec) at (9,-1.5) {Cross-Att.\\(Decoder)};
    \node[block, fill=teal!10] (io_out) at (9,-3.5) {Output\\$\mathbf{y} \in \mathbb{R}^{O \times F}$};
    \node[block, fill=purple!10] (io_q) at (13,-1.5) {Query\\$\mathbf{q} \in \mathbb{R}^{O \times E}$};

    \draw[arr] (io_in) -- (io_enc);
    \draw[arr] (io_enc) -- (io_lat);
    \draw[arr] (io_lat) -- (io_dec);
    \draw[arr] (io_dec) -- (io_out);
    \draw[arr] (io_q) -- (io_dec);

    % Advantage
    \node[font=\scriptsize, text=green!50!black, align=center] at (9,-5) {Output arbitrari!\\Classificazione, flusso ottico,\\linguaggio, multimodale \ldots};

    % Separator
    \draw[dashed, gray] (5.8,5) -- (5.8,-5.5);
\end{tikzpicture}
\end{center}

La differenza chiave \`e nel \textbf{meccanismo di decodifica}: il Perceiver usa un semplice average pooling (che produce un singolo vettore), mentre il Perceiver IO usa una \textbf{cross-attention con query di output} che pu\`o produrre un numero arbitrario di vettori di output.


%% ============================================================
%%  SEZIONE 2 --- ARCHITETTURA PERCEIVER IO
%% ============================================================
\section{Architettura Perceiver IO}

L'architettura Perceiver IO segue un flusso Encode-Process-Decode.

\begin{figure}[H]
\centering
\includegraphics[width=0.85\textwidth]{paper_figures/perceiver_io_fig2_architecture.png}
\caption{Architettura del Perceiver IO (Jaegle et al., 2022, Figura 2). L'input array ($M \times C$) viene codificato nello spazio latente ($N \times D$) tramite cross-attention (Encode), processato con $L$ blocchi di self-attention (Process), e infine decodificato tramite cross-attention con le output query ($O \times E$) per produrre l'output array ($O \times F$) (Decode). L'innovazione chiave rispetto al Perceiver originale \`e il modulo Decode, che permette output strutturati di dimensione arbitraria.}
\label{fig:perceiver_io_arch}
\end{figure}

\subsection{Overview: Encode-Process-Decode}

L'architettura Perceiver IO opera su quattro array principali:

\begin{formulabox}[Array fondamentali del Perceiver IO]
\begin{align}
    \textbf{Input array:} \quad & \mathbf{x} \in \mathbb{R}^{M \times C} \quad \text{($M$ elementi, $C$ canali per elemento)} \\
    \textbf{Latent array:} \quad & \mathbf{z} \in \mathbb{R}^{N \times D} \quad \text{($N$ latenti, $D$ canali per latente)} \\
    \textbf{Output query array:} \quad & \mathbf{q} \in \mathbb{R}^{O \times E} \quad \text{($O$ punti di output, $E$ feature per query)} \\
    \textbf{Output array:} \quad & \mathbf{y} \in \mathbb{R}^{O \times F} \quad \text{($O$ output, $F$ canali per output)}
\end{align}
\end{formulabox}

\begin{notabox}[Parametri del task vs.\ iperparametri]
\begin{itemize}
    \item $M$, $C$: determinati dall'\textbf{input} (es.\ $M = 50{,}176$ pixel, $C = 3$ canali RGB). \textbf{Propriet\`a del task.}
    \item $O$, $E$, $F$: determinati dall'\textbf{output} desiderato (es.\ $O = 50{,}176$ spostamenti, $F = 2$). \textbf{Propriet\`a del task.}
    \item $N$, $D$: \textbf{iperparametri} scelti dal progettista. $N$ \`e tipicamente \emph{molto minore} di $M$ e $O$ (es.\ $N = 256$ o $N = 512$). Questo \`e ci\`o che rende il modello trattabile.
\end{itemize}
\end{notabox}

L'architettura completa \`e composta da tre moduli in cascata:

\begin{center}
\begin{tikzpicture}[
    module/.style={draw, thick, rounded corners=5pt, minimum width=4cm, minimum height=2.5cm, align=center},
    arr/.style={-{Stealth[length=3mm]}, very thick},
    dim/.style={font=\scriptsize, text=gray!70!black},
    lbl/.style={font=\small\bfseries}
]
    % Encoder
    \node[module, fill=orange!10, draw=orange!70!black, align=center] (enc) at (0,0) {
        {\small\bfseries Encoder}\\[3pt]
        \scriptsize Cross-Attention\\
        \scriptsize $\mathbf{Q}$ da latenti\\
        \scriptsize $\mathbf{K}, \mathbf{V}$ da input
    };

    % Processor
    \node[module, fill=blue!10, draw=blue!70!black, align=center] (proc) at (5.5,0) {
        {\small\bfseries Processor}\\[3pt]
        \scriptsize Self-Attention $\times L$\\
        \scriptsize Solo sui latenti\\
        \scriptsize $\mathbf{Q}, \mathbf{K}, \mathbf{V}$ da latenti
    };

    % Decoder
    \node[module, fill=green!10, draw=green!70!black, align=center] (dec) at (11,0) {
        {\small\bfseries Decoder}\\[3pt]
        \scriptsize Cross-Attention\\
        \scriptsize $\mathbf{Q}$ da query di output\\
        \scriptsize $\mathbf{K}, \mathbf{V}$ da latenti
    };

    % Input
    \node[dim] (in_dim) at (0,2.5) {$\mathbf{x} \in \mathbb{R}^{M \times C}$};
    \node[dim] (lat_init) at (0,-2.5) {$\mathbf{z}^{(0)} \in \mathbb{R}^{N \times D}$};

    % Arrows
    \draw[arr, orange!70!black] (in_dim) -- (enc);
    \draw[arr, orange!70!black] (lat_init) -- (enc);
    \draw[arr] (enc) -- (proc) node[midway, above, dim] {$\mathbb{R}^{N \times D}$};
    \draw[arr] (proc) -- (dec) node[midway, above, dim] {$\mathbb{R}^{N \times D}$};

    % Output query
    \node[dim] (q_dim) at (11,-2.5) {$\mathbf{q} \in \mathbb{R}^{O \times E}$};
    \draw[arr, green!70!black] (q_dim) -- (dec);

    % Output
    \node[dim] (out_dim) at (14.5,0) {$\mathbf{y} \in \mathbb{R}^{O \times F}$};
    \draw[arr, green!70!black] (dec) -- (out_dim);

    % Complexity annotations
    \node[font=\tiny, text=orange!70!black] at (0,-1.7) {$\mathcal{O}(MN)$};
    \node[font=\tiny, text=blue!70!black] at (5.5,-1.7) {$\mathcal{O}(LN^2)$};
    \node[font=\tiny, text=green!70!black] at (11,-1.7) {$\mathcal{O}(ON)$};
\end{tikzpicture}
\end{center}


\subsection{Encoder: Cross-Attention dall'input ai latenti}

L'encoder \`e identico a quello del Perceiver originale. Si tratta di un modulo di \textbf{cross-attention} in cui:
\begin{itemize}
    \item Le \textbf{query} ($\mathbf{Q}$) provengono dall'array latente $\mathbf{z}^{(0)}$ (inizializzato con parametri apprendibili)
    \item Le \textbf{chiavi} ($\mathbf{K}$) e i \textbf{valori} ($\mathbf{V}$) provengono dall'input $\mathbf{x}$
\end{itemize}

\begin{formulabox}[Encoder: Cross-Attention Input $\to$ Latenti]
Date le matrici dei pesi $\mathbf{W}_Q^{\text{enc}} \in \mathbb{R}^{D \times d_k}$, $\mathbf{W}_K^{\text{enc}} \in \mathbb{R}^{C \times d_k}$, $\mathbf{W}_V^{\text{enc}} \in \mathbb{R}^{C \times d_v}$:
\begin{align}
    \mathbf{Q}^{\text{enc}} &= \mathbf{z}^{(0)} \, \mathbf{W}_Q^{\text{enc}} \in \mathbb{R}^{N \times d_k} \quad \text{(query dai latenti)} \\
    \mathbf{K}^{\text{enc}} &= \mathbf{x} \, \mathbf{W}_K^{\text{enc}} \in \mathbb{R}^{M \times d_k} \quad \text{(chiavi dall'input)} \\
    \mathbf{V}^{\text{enc}} &= \mathbf{x} \, \mathbf{W}_V^{\text{enc}} \in \mathbb{R}^{M \times d_v} \quad \text{(valori dall'input)} \\[6pt]
    \text{Attention}^{\text{enc}} &= \operatorname{softmax}\!\left(\frac{\mathbf{Q}^{\text{enc}} (\mathbf{K}^{\text{enc}})^\top}{\sqrt{d_k}}\right) \mathbf{V}^{\text{enc}} \in \mathbb{R}^{N \times d_v}
\end{align}
L'output dell'encoder, dopo residual connection e layer normalization:
\[
\mathbf{z}^{(0)}_{\text{out}} = \text{LayerNorm}\!\left(\mathbf{z}^{(0)} + \text{Attention}^{\text{enc}}\right) \in \mathbb{R}^{N \times D}
\]
\end{formulabox}

L'intuizione \`e che ogni \textbf{latente} (riga di $\mathbf{z}^{(0)}$) ``interroga'' l'intero input per decidere quali informazioni raccogliere e dove memorizzarle. Poich\'e ci sono $N$ latenti e $M$ elementi di input, la matrice di attenzione \`e $N \times M$, con complessit\`a:
\[
\text{Complessit\`a encoder} = \mathcal{O}(M \cdot N \cdot D)
\]
Dato che $N \ll M$ (tipicamente $N = 256$--$1024$ mentre $M$ pu\`o essere $50{,}000$+), questo \`e \textbf{lineare in $M$}, molto pi\`u efficiente della self-attention quadratica $\mathcal{O}(M^2 D)$.

\begin{center}
\begin{tikzpicture}[
    mat/.style={draw, minimum width=1.2cm, minimum height=2cm, align=center, font=\scriptsize},
    op/.style={circle, draw, minimum size=0.7cm, font=\small},
    arr/.style={-{Stealth[length=2.5mm]}, thick}
]
    % Input
    \node[mat, fill=red!10, minimum height=3cm] (x) at (0,0) {$\mathbf{x}$\\$M \times C$};

    % W_K, W_V
    \node[mat, fill=gray!10, minimum height=1.5cm] (wk) at (2.5,1) {$\mathbf{W}_K$\\$C \times d_k$};
    \node[mat, fill=gray!10, minimum height=1.5cm] (wv) at (2.5,-1) {$\mathbf{W}_V$\\$C \times d_v$};

    % K, V
    \node[mat, fill=yellow!15, minimum height=3cm] (K) at (5,1) {$\mathbf{K}$\\$M \times d_k$};
    \node[mat, fill=cyan!10, minimum height=3cm] (V) at (5,-1) {$\mathbf{V}$\\$M \times d_v$};

    % Latent
    \node[mat, fill=blue!15, minimum height=1.5cm] (z) at (0,-4) {$\mathbf{z}^{(0)}$\\$N \times D$};
    \node[mat, fill=gray!10, minimum height=1cm] (wq) at (2.5,-4) {$\mathbf{W}_Q$\\$D \times d_k$};
    \node[mat, fill=yellow!25, minimum height=1.5cm] (Q) at (5,-4) {$\mathbf{Q}$\\$N \times d_k$};

    % Attention matrix
    \node[mat, fill=orange!15, minimum height=1.5cm, minimum width=2.5cm] (att) at (8.5,-1.5) {$\mathbf{A} = \text{softmax}\!\left(\frac{\mathbf{Q}\mathbf{K}^\top}{\sqrt{d_k}}\right)$\\$N \times M$};

    % Output
    \node[mat, fill=blue!20, minimum height=1.5cm] (out) at (12.5,-1.5) {Output\\$N \times d_v$};

    % Arrows
    \draw[arr] (x) -- (wk);
    \draw[arr] (x) -- (wv);
    \draw[arr] (wk) -- (K);
    \draw[arr] (wv) -- (V);
    \draw[arr] (z) -- (wq);
    \draw[arr] (wq) -- (Q);
    \draw[arr] (Q) -- (att);
    \draw[arr] (K) -- (att);
    \draw[arr] (att) -- (out);
    \draw[arr] (V) -- (out);

    % Labels
    \node[font=\scriptsize\bfseries, text=orange!70!black] at (8.5,-3.2) {$N$ latenti $\times$ $M$ input};
\end{tikzpicture}
\end{center}

\begin{notabox}[Nota sulla multi-head attention]
In pratica, sia l'encoder che il decoder usano \textbf{multi-head attention} con $h$ teste. Le dimensioni per-testa sono $d_k = d_v = D/h$. L'output delle teste viene concatenato e proiettato:
\[
\text{MultiHead}(\mathbf{Q}, \mathbf{K}, \mathbf{V}) = \text{Concat}(\text{head}_1, \ldots, \text{head}_h)\mathbf{W}_O
\]
dove $\mathbf{W}_O \in \mathbb{R}^{hd_v \times D}$. Le complessit\`a rimangono le stesse in termini di ordine di grandezza.
\end{notabox}

\begin{notabox}[In parole semplici: l'Encoder]
Immagina di avere una \textbf{biblioteca enorme} (l'input, con $M$ libri) e solo $N$ scaffali (i latenti) per organizzarla. L'encoder \`e il bibliotecario: per ogni scaffale, guarda \emph{tutti} i libri e decide quali mettere l\`i. Lo scaffale ``Animali'' raccoglier\`a informazioni dai libri su gatti, cani, pesci, etc. Lo scaffale ``Colori'' raccoglier\`a informazioni sulle tonalit\`a presenti. Il meccanismo di attenzione \`e esattamente questo: ogni scaffale (latente) ``guarda'' tutti i libri (input) e decide quanto di ciascuno conservare, usando i pesi di attenzione come criterio di rilevanza.

La cosa geniale \`e che non serve leggere tutti i libri per ogni domanda che farai dopo --- basta consultare i $N$ scaffali, che sono molto meno di $M$ libri!
\end{notabox}


\subsection{Processor: Self-Attention nello spazio latente}

Il processor \`e il cuore computazionale del modello. Consiste in $L$ layer di \textbf{self-attention} applicati esclusivamente all'array latente $\mathbf{z} \in \mathbb{R}^{N \times D}$.

\begin{formulabox}[Processor: $L$ layer di Self-Attention sui latenti]
Per ogni layer $\ell = 1, \ldots, L$:
\begin{align}
    \mathbf{Q}^{(\ell)} &= \mathbf{z}^{(\ell-1)} \mathbf{W}_Q^{(\ell)}, \quad
    \mathbf{K}^{(\ell)} = \mathbf{z}^{(\ell-1)} \mathbf{W}_K^{(\ell)}, \quad
    \mathbf{V}^{(\ell)} = \mathbf{z}^{(\ell-1)} \mathbf{W}_V^{(\ell)} \\[4pt]
    \mathbf{a}^{(\ell)} &= \text{LayerNorm}\!\left(\mathbf{z}^{(\ell-1)} + \operatorname{softmax}\!\left(\frac{\mathbf{Q}^{(\ell)} (\mathbf{K}^{(\ell)})^\top}{\sqrt{d_k}}\right) \mathbf{V}^{(\ell)}\right) \\[4pt]
    \mathbf{z}^{(\ell)} &= \text{LayerNorm}\!\left(\mathbf{a}^{(\ell)} + \text{FFN}(\mathbf{a}^{(\ell)})\right)
\end{align}
dove $\text{FFN}(\mathbf{a}) = \text{GELU}(\mathbf{a}\mathbf{W}_1 + \mathbf{b}_1)\mathbf{W}_2 + \mathbf{b}_2$ \`e la rete feed-forward posizionale.
\end{formulabox}

\begin{notabox}[Il processor \`e indipendente da $M$ e da $O$!]
La complessit\`a del processor \`e:
\[
\text{Complessit\`a processor} = \mathcal{O}(L \cdot N^2 \cdot D)
\]
Questa non dipende n\'e dalla dimensione dell'input $M$ n\'e dalla dimensione dell'output $O$. Poich\'e $N$ \`e un iperparametro scelto piccolo (tipicamente $256$--$1024$), il processor \`e computazionalmente molto efficiente. \`E in questo modulo che avviene la maggior parte del ``ragionamento'' del modello: i latenti si scambiano informazione tra loro attraverso self-attention per $L$ passi.
\end{notabox}

\begin{notabox}[Weight sharing nel processor]
Nel Perceiver IO, il processor pu\`o usare \textbf{weight sharing}: i pesi dei layer di self-attention vengono condivisi tra diversi blocchi. Ad esempio, nel setup per ImageNet si usano 8 blocchi di 6 moduli di attention ciascuno, dove i pesi sono condivisi \emph{tra i blocchi} (ma non all'interno dello stesso blocco). Questo riduce significativamente il numero di parametri pur mantenendo la profondit\`a effettiva del modello. Formalmente, se il processor ha $B$ blocchi con $L_b$ layer ciascuno, i parametri unici sono quelli di un singolo blocco ($L_b$ set di pesi), ma l'effettiva profondit\`a \`e $B \times L_b$.
\end{notabox}

\begin{notabox}[In parole semplici: il Processor e il weight sharing]
Il processor \`e il ``cervello pensante'' del Perceiver IO. \`E come un gruppo di $N$ esperti seduti attorno a un tavolo che \textbf{discutono tra loro} per $L$ turni. Ad ogni turno, ogni esperto ascolta tutti gli altri e aggiorna le proprie idee.

Il \textbf{weight sharing} \`e come usare lo stesso ``protocollo di discussione'' per pi\`u turni. Immagina che i primi 6 turni seguano un certo formato di discussione (es.\ brainstorming), poi i turni 7--12 ripetano lo stesso formato (un secondo brainstorming), e cos\`i via per 8 cicli. I ``formati'' sono gli stessi (stessi pesi), ma le discussioni sono diverse perch\'e le idee si sono evolute.

Questo permette di avere $8 \times 6 = 48$ turni di discussione effettivi (molta profondit\`a) usando solo i pesi di 6 turni (pochi parametri).
\end{notabox}


\subsection{Decoder: Cross-Attention dai latenti all'output --- L'INNOVAZIONE CHIAVE}

Il \textbf{decoder} \`e la vera novit\`a del Perceiver IO rispetto al Perceiver originale. Si tratta di un modulo di \textbf{cross-attention} che opera in direzione opposta rispetto all'encoder:

\begin{itemize}
    \item Le \textbf{query} ($\mathbf{Q}$) provengono dall'\textbf{array di query di output} $\mathbf{q} \in \mathbb{R}^{O \times E}$
    \item Le \textbf{chiavi} ($\mathbf{K}$) e i \textbf{valori} ($\mathbf{V}$) provengono dai \textbf{latenti} $\mathbf{z}^{(L)} \in \mathbb{R}^{N \times D}$
\end{itemize}

\begin{formulabox}[Decoder: Cross-Attention Latenti $\to$ Output]
Date le matrici dei pesi $\mathbf{W}_Q^{\text{dec}} \in \mathbb{R}^{E \times d_k}$, $\mathbf{W}_K^{\text{dec}} \in \mathbb{R}^{D \times d_k}$, $\mathbf{W}_V^{\text{dec}} \in \mathbb{R}^{D \times d_v}$:
\begin{align}
    \mathbf{Q}^{\text{dec}} &= \mathbf{q} \, \mathbf{W}_Q^{\text{dec}} \in \mathbb{R}^{O \times d_k} \quad \text{(query dall'output query array)} \\
    \mathbf{K}^{\text{dec}} &= \mathbf{z}^{(L)} \mathbf{W}_K^{\text{dec}} \in \mathbb{R}^{N \times d_k} \quad \text{(chiavi dai latenti)} \\
    \mathbf{V}^{\text{dec}} &= \mathbf{z}^{(L)} \mathbf{W}_V^{\text{dec}} \in \mathbb{R}^{N \times d_v} \quad \text{(valori dai latenti)} \\[6pt]
    \text{Attention}^{\text{dec}} &= \operatorname{softmax}\!\left(\frac{\mathbf{Q}^{\text{dec}} (\mathbf{K}^{\text{dec}})^\top}{\sqrt{d_k}}\right) \mathbf{V}^{\text{dec}} \in \mathbb{R}^{O \times d_v}
\end{align}
L'output finale, eventualmente con un MLP aggiuntivo:
\[
\mathbf{y} = \text{MLP}\!\left(\text{Attention}^{\text{dec}}\right) \in \mathbb{R}^{O \times F}
\]
dove $F$ \`e la dimensione dell'output per punto (es.\ 2 per optical flow, $|\text{vocabolario}|$ per linguaggio, 3 per RGB).
\end{formulabox}

L'aspetto cruciale \`e che \textbf{ogni punto di output interroga indipendentemente} l'array latente. La matrice di attenzione \`e $O \times N$, dove la riga $i$-esima contiene i pesi con cui il punto di output $i$ combina le informazioni dei $N$ latenti.

\begin{center}
\begin{tikzpicture}[
    mat/.style={draw, minimum width=1.3cm, align=center, font=\scriptsize},
    arr/.style={-{Stealth[length=2.5mm]}, thick}
]
    % Output queries
    \node[mat, fill=purple!10, minimum height=3.5cm] (q) at (0,0) {$\mathbf{q}$\\$O \times E$};
    \node[mat, fill=gray!10, minimum height=1.2cm] (wq) at (2.5,0) {$\mathbf{W}_Q^{\text{dec}}$\\$E \times d_k$};
    \node[mat, fill=purple!20, minimum height=3.5cm] (Q) at (5,0) {$\mathbf{Q}^{\text{dec}}$\\$O \times d_k$};

    % Latents
    \node[mat, fill=blue!15, minimum height=1.8cm] (z) at (0,-4.5) {$\mathbf{z}^{(L)}$\\$N \times D$};
    \node[mat, fill=gray!10, minimum height=1cm] (wk) at (2.5,-3.8) {$\mathbf{W}_K^{\text{dec}}$\\$D \times d_k$};
    \node[mat, fill=gray!10, minimum height=1cm] (wv) at (2.5,-5.2) {$\mathbf{W}_V^{\text{dec}}$\\$D \times d_v$};
    \node[mat, fill=yellow!15, minimum height=1.8cm] (K) at (5,-3.8) {$\mathbf{K}^{\text{dec}}$\\$N \times d_k$};
    \node[mat, fill=cyan!10, minimum height=1.8cm] (V) at (5,-5.2) {$\mathbf{V}^{\text{dec}}$\\$N \times d_v$};

    % Attention matrix
    \node[mat, fill=green!15, minimum height=3cm, minimum width=2cm] (att) at (8.5,-2) {$\mathbf{A}^{\text{dec}}$\\$O \times N$\\[4pt]\scriptsize $\operatorname{softmax}\!\left(\frac{\mathbf{Q}\mathbf{K}^\top}{\sqrt{d_k}}\right)$};

    % Output
    \node[mat, fill=teal!15, minimum height=3.5cm, minimum width=1.5cm] (out) at (12,-2) {$\mathbf{y}$\\$O \times d_v$};

    % Arrows
    \draw[arr] (q) -- (wq);
    \draw[arr] (wq) -- (Q);
    \draw[arr] (z) -- (wk);
    \draw[arr] (z) -- (wv);
    \draw[arr] (wk) -- (K);
    \draw[arr] (wv) -- (V);
    \draw[arr] (Q) -- (att);
    \draw[arr] (K) -- (att);
    \draw[arr] (att) -- (out);
    \draw[arr] (V) -- (out);

    % Dimension labels
    \draw[decorate, decoration={brace, amplitude=5pt, mirror}, thick, green!50!black] (7.2,-3.8) -- (7.2,-0.2) node[midway, left=8pt, font=\scriptsize, text=green!50!black] {$O$ output};
    \draw[decorate, decoration={brace, amplitude=5pt}, thick, blue!50!black] (9.8,-3.8) -- (9.8,-0.2) node[midway, right=8pt, font=\scriptsize, text=blue!50!black] {$N$ latenti};
\end{tikzpicture}
\end{center}

La complessit\`a del decoder \`e:
\[
\text{Complessit\`a decoder} = \mathcal{O}(O \cdot N \cdot D)
\]
che \`e \textbf{lineare in $O$}. Questo \`e fondamentale: anche se l'output ha $O = 50{,}000$ punti, il costo per punto \`e solo $\mathcal{O}(N \cdot D)$, perch\'e ogni punto interroga solo $N$ latenti (non tutti gli altri punti di output).

\begin{notabox}[Perch\'e il decoder usa un solo layer di cross-attention?]
Il paper propone un decoder con \textbf{un singolo layer} di cross-attention (senza self-attention tra le query di output). La motivazione \`e duplice:
\begin{enumerate}
    \item \textbf{Efficienza}: se $O$ \`e grande (es.\ $50{,}000$ pixel), fare self-attention tra le query di output avrebbe costo $\mathcal{O}(O^2 D)$, vanificando il vantaggio dell'architettura
    \item \textbf{Sufficienza}: tutta l'informazione necessaria \`e gi\`a stata elaborata nel processor e risiede nei latenti. Il decoder deve solo ``leggere'' dal latente, non fare ulteriore ragionamento
\end{enumerate}
Opzionalmente, il decoder pu\`o includere un piccolo MLP dopo la cross-attention per trasformazioni punto-per-punto.
\end{notabox}

\begin{notabox}[In parole semplici: il Decoder]
Immagina di avere un \textbf{dizionario enciclopedico} (i latenti processati). Per ogni informazione che vuoi ottenere (output query), fai una \textbf{domanda specifica} al dizionario. Se vuoi sapere il colore del pixel alla posizione $(100, 50)$, la tua domanda include ``posizione $(100, 50)$''. Il dizionario cerca le pagine pi\`u rilevanti (i latenti con i pesi di attenzione pi\`u alti) e ti d\`a la risposta.

La cosa fondamentale \`e che \textbf{ogni domanda \`e indipendente}: la risposta alla domanda sul pixel $(100, 50)$ non dipende dalla risposta alla domanda sul pixel $(200, 300)$. Questo \`e ci\`o che rende il decoder cos\`i efficiente e parallelizzabile.

Pensalo come un esame orale: il professore (decoder) fa domande diverse a ogni studente (output query), e ogni studente risponde consultando lo stesso libro di testo (i latenti). Le risposte sono diverse perch\'e le domande sono diverse, non perch\'e i libri cambiano!
\end{notabox}


%% ============================================================
%%  SEZIONE 2.5 --- MECCANISMO DI DECODE DETTAGLIATO
%% ============================================================
\subsection{Meccanismo di Decode dettagliato --- L'innovazione chiave}

Questa sottosezione approfondisce il decoder del Perceiver IO, che rappresenta il contributo principale del paper rispetto al Perceiver originale.

\subsubsection{Output Query Array --- Spiegazione completa}

L'\textbf{output query array} $\mathbf{q} \in \mathbb{R}^{O \times E}$ \`e l'elemento che rende il Perceiver IO cos\`i flessibile. Ogni riga $\mathbf{q}_i \in \mathbb{R}^E$ codifica la domanda: ``Quale informazione devo estrarre dall'array latente per produrre l'output $i$?''

La query per un singolo punto di output pu\`o essere costruita combinando diversi tipi di informazione:

\begin{formulabox}[Struttura generale di una output query]
\[
\mathbf{q}_i = f\!\left(\text{pos}_i, \; \text{task}_i, \; \text{modality}_i, \; \text{input\_feat}_i\right) \in \mathbb{R}^E
\]
dove $f$ pu\`o essere una concatenazione, una somma, o una combinazione appresa. Le componenti sono:
\begin{itemize}
    \item $\text{pos}_i$: \textbf{informazione posizionale} --- dove si trova il punto di output nello spazio (coordinate spaziali, temporali, o nella sequenza)
    \item $\text{task}_i$: \textbf{informazione sul task} --- quale compito il modello deve svolgere per questo punto (es.\ classificazione, ricostruzione, predizione)
    \item $\text{modality}_i$: \textbf{informazione sulla modalit\`a} --- a quale modalit\`a di output appartiene questo punto (es.\ video, audio, label)
    \item $\text{input\_feat}_i$: \textbf{feature dall'input} --- informazione estratta direttamente dall'input correlata a questo punto di output (usata soprattutto per optical flow)
\end{itemize}
\end{formulabox}

Vediamo ora nel dettaglio come le query vengono costruite per ciascun tipo di task.

\subsubsection{Query per la Classificazione}

Per la classificazione (es.\ ImageNet), l'output \`e un \textbf{singolo vettore} di probabilit\`a sulle $K$ classi. Si usa pertanto una \textbf{singola query appresa}:

\begin{formulabox}[Query per classificazione]
\[
\mathbf{q}_{\text{cls}} \in \mathbb{R}^{1 \times E} \quad \text{(parametro apprendibile)}
\]
\begin{itemize}
    \item $O = 1$: una sola query produce un solo output
    \item La query \`e completamente appresa dal training (analogo al token \texttt{[CLS]} di BERT, Devlin et al., 2019)
    \item Output: $\mathbf{y} = \text{MLP}(\text{CrossAttn}(\mathbf{q}_{\text{cls}}, \mathbf{z}^{(L)})) \in \mathbb{R}^{1 \times K}$
    \item La singola query ``riassume'' tutta l'informazione dei latenti in un unico vettore
\end{itemize}
\end{formulabox}

Per il \textbf{multi-task learning}, invece di una singola query, si usano embedding task-specifici. Ogni task ha il proprio embedding appreso $\mathbf{e}_{\text{task}} \in \mathbb{R}^{E_t}$, permettendo a un singolo modello di gestire pi\`u task con query diverse.

\subsubsection{Query per Optical Flow}

Per l'optical flow, l'output deve specificare un vettore di spostamento 2D $(dx, dy)$ per \textbf{ogni pixel} dell'immagine. La costruzione delle query \`e pi\`u elaborata.

\begin{formulabox}[Query per optical flow]
Per ogni pixel alla posizione $(x, y)$:
\[
\mathbf{q}_{(x,y)} = \text{concat}\!\left[\underbrace{\text{FourierPE}\!\left(\frac{x}{W}, \frac{y}{H}\right)}_{\text{encoding posizionale}}, \; \underbrace{\mathbf{f}_{(x,y)}^{\text{input}}}_{\text{feature dall'input}}\right] \in \mathbb{R}^E
\]
\begin{itemize}
    \item $O = H \times W$: una query per ogni pixel (es.\ $O = 11{,}408$ per Sintel alla risoluzione di training)
    \item Il \textbf{Fourier positional encoding} codifica le coordinate $(x, y)$ normalizzate in $[0, 1]$ come combinazione di seni e coseni a diverse frequenze
    \item Le \textbf{feature dell'input} $\mathbf{f}_{(x,y)}^{\text{input}}$ sono estratte dall'input (patch $3 \times 3$ del primo frame) nella posizione $(x, y)$
    \item Output per pixel: $\mathbf{y}_{(x,y)} \in \mathbb{R}^2$ (spostamento orizzontale $dx$ e verticale $dy$)
\end{itemize}
\end{formulabox}

L'intuizione \`e potente: ogni query dice al decoder ``Sono il pixel $(x, y)$ del primo frame, con queste feature locali --- dimmi dove mi trovo nel secondo frame''. Il decoder cerca nei latenti (che contengono informazione globale sulle corrispondenze tra i due frame) la risposta specifica per quel pixel.

\begin{center}
\begin{tikzpicture}[
    arr/.style={-{Stealth[length=2.5mm]}, thick},
    block/.style={draw, rounded corners, minimum width=2.5cm, minimum height=0.8cm, align=center, font=\scriptsize}
]
    % Pixel query
    \node[block, fill=purple!10, minimum width=4cm] (pix) at (0,3) {Pixel $(x, y)$ nel Frame $t$};

    % Fourier PE
    \node[block, fill=yellow!15] (fpe) at (-2.5,1.5) {FourierPE$(x/W, y/H)$};

    % Input features
    \node[block, fill=red!10] (feat) at (2.5,1.5) {Patch $3\times3$ features};

    % Concat
    \node[block, fill=purple!20, minimum width=5cm] (query) at (0,0) {$\mathbf{q}_{(x,y)} = [\text{FourierPE} \| \mathbf{f}_{(x,y)}]$};

    % Decoder cross-attention
    \node[block, fill=green!15, minimum width=5cm] (dec) at (0,-1.5) {Decoder Cross-Attention};

    % Latents
    \node[block, fill=blue!15, minimum width=2.5cm] (lat) at (5,-1.5) {Latenti $\mathbf{z}^{(L)}$\\$256 \times 512$};

    % Output
    \node[block, fill=teal!15] (out) at (0,-3) {$(dx, dy)$ per pixel $(x,y)$};

    \draw[arr] (pix) -- (fpe);
    \draw[arr] (pix) -- (feat);
    \draw[arr] (fpe) -- (query);
    \draw[arr] (feat) -- (query);
    \draw[arr] (query) -- (dec);
    \draw[arr] (lat) -- (dec);
    \draw[arr] (dec) -- (out);

    \node[font=\tiny\itshape, text=gray] at (3, -0.5) {Una query per pixel};
    \node[font=\tiny\itshape, text=gray, align=center] at (3, -2.5) {Ogni pixel ``interroga''\\tutti i 256 latenti};
\end{tikzpicture}
\end{center}

\subsubsection{Query per Masked Language Modeling}

Per il Masked Language Modeling (MLM), il modello deve predire i byte originali nelle posizioni mascherate.

\begin{formulabox}[Query per MLM]
Per ogni posizione mascherata $i$:
\[
\mathbf{q}_i = \text{LearnedPE}(i) \in \mathbb{R}^E
\]
\begin{itemize}
    \item $O$ = numero di posizioni mascherate (tipicamente $\approx 15\%$ di $M = 2048$, quindi $O \approx 307$)
    \item Ogni query \`e un \textbf{embedding posizionale appreso} specifico per quella posizione nella sequenza
    \item Output per posizione: $\mathbf{y}_i \in \mathbb{R}^{260}$ dove $260 = 256$ (valori byte) $+ 4$ (token speciali: PAD, MASK, CLS, SEP)
    \item Loss: \textbf{softmax cross-entropy} per posizione: $\mathcal{L}_i = -\log p(x_i^* \mid \mathbf{q}_i, \mathbf{z}^{(L)})$
\end{itemize}
\end{formulabox}

La differenza rispetto a BERT \`e che BERT produce output per \textbf{tutte} le posizioni (tramite self-attention), mentre il Perceiver IO produce output \textbf{solo per le posizioni richieste dalle query}. Questo \`e pi\`u efficiente perch\'e il decoder ha costo $\mathcal{O}(O \cdot N)$ con $O \ll M$.

\subsubsection{Query per Autoencoding Multimodale}

L'autoencoding multimodale \`e il caso pi\`u complesso e dimostra appieno la flessibilit\`a del meccanismo di query.

\begin{formulabox}[Query per autoencoding multimodale (Kinetics-700)]
Tre tipi di query, una per ogni modalit\`a:

\textbf{Query video} (per ogni patch $(x, y, t)$):
\[
\mathbf{q}_v^{(x,y,t)} = \text{concat}\!\left[\text{FourierPE}(x, y, t), \; \mathbf{e}_{\text{video}}\right] \in \mathbb{R}^E
\]
$O_{\text{video}} = 16 \times 56 \times 56 = 50{,}176$ patch.\\[6pt]

\textbf{Query audio} (per ogni campione temporale $t$):
\[
\mathbf{q}_a^{(t)} = \text{concat}\!\left[\text{FourierPE}(t), \; \mathbf{e}_{\text{audio}}\right] \in \mathbb{R}^E
\]
$O_{\text{audio}} = 30{,}720$ campioni.\\[6pt]

\textbf{Query label} (singola):
\[
\mathbf{q}_l = \mathbf{e}_{\text{label}} \in \mathbb{R}^E
\]
$O_{\text{label}} = 1$.\\[6pt]

\textbf{Totale}: $O = O_{\text{video}} + O_{\text{audio}} + O_{\text{label}} = 50{,}176 + 30{,}720 + 1 = 80{,}897$ punti di output!

Dove $\mathbf{e}_{\text{video}}, \mathbf{e}_{\text{audio}}, \mathbf{e}_{\text{label}} \in \mathbb{R}^{E_m}$ sono \textbf{embedding di modalit\`a appresi} (``flag'' che dicono al decoder quale tipo di output produrre).
\end{formulabox}

\begin{center}
\begin{tikzpicture}[
    qbox/.style={draw, rounded corners, minimum width=5.5cm, minimum height=1.2cm, align=left, font=\scriptsize},
    arr/.style={-{Stealth[length=2mm]}, thick}
]
    % Video queries
    \node[qbox, fill=red!8] (qv) at (0,4) {
        \textbf{Query Video} ($O_v = 50{,}176$)\\
        $\mathbf{q}_v = [\text{FourierPE}(x,y,t) \| \mathbf{e}_{\text{video}}]$\\
        Output: ricostruzione pixel RGB
    };

    % Audio queries
    \node[qbox, fill=yellow!8] (qa) at (0,1.5) {
        \textbf{Query Audio} ($O_a = 30{,}720$)\\
        $\mathbf{q}_a = [\text{FourierPE}(t) \| \mathbf{e}_{\text{audio}}]$\\
        Output: ricostruzione campioni audio
    };

    % Label query
    \node[qbox, fill=green!8] (ql) at (0,-1) {
        \textbf{Query Label} ($O_l = 1$)\\
        $\mathbf{q}_l = \mathbf{e}_{\text{label}}$\\
        Output: logit su 700 classi
    };

    % Decoder
    \node[draw, rounded corners, fill=green!15, minimum width=3.5cm, minimum height=5cm, align=center, font=\small\bfseries] (dec) at (8, 1.5) {Decoder\\Cross-Attention\\(singolo layer)};

    % Latents
    \node[draw, rounded corners, fill=blue!15, minimum width=2.5cm, minimum height=1cm, font=\scriptsize] (lat) at (8, -2) {$\mathbf{z}^{(L)} \in \mathbb{R}^{N \times D}$};

    % Arrows
    \draw[arr] (qv.east) -- (dec.west |- qv.east);
    \draw[arr] (qa.east) -- (dec.west |- qa.east);
    \draw[arr] (ql.east) -- (dec.west |- ql.east);
    \draw[arr] (lat) -- (dec);

    % Total
    \node[draw, dashed, rounded corners, fit=(qv)(qa)(ql), inner sep=8pt, label={[font=\scriptsize\bfseries]above:$O_{\text{totale}} = 80{,}897$}] {};

    % Subsampling note
    \node[draw, rounded corners, fill=orange!10, minimum width=5cm, font=\scriptsize, align=center] at (4,-3) {
        \textbf{Durante il training}: subsample a $512 + 512 + 1 = 1{,}025$ query\\
        \textbf{Durante l'inferenza}: tutte le $80{,}897$ query
    };
\end{tikzpicture}
\end{center}

\begin{notabox}[In parole semplici: l'autoencoding multimodale]
Immagina un \textbf{interprete simultaneo} che ascolta un discorso in una lingua (video + audio) e deve tradurlo in \textbf{tre lingue diverse contemporaneamente} (ricostruzione video, ricostruzione audio, classificazione). Le query di output sono come i \textbf{microfoni di tre canali diversi}: ognuno chiede una ``traduzione'' diversa dallo stesso discorso compreso dall'interprete.

L'interprete (i latenti) ha capito il \emph{significato} del discorso. Il microfono ``video'' chiede: ``Com'era visivamente la scena al tempo $t$, posizione $(x, y)$?''. Il microfono ``audio'' chiede: ``Che suono c'era al campione $t$?''. Il microfono ``label'' chiede: ``Di cosa parlava in generale?''. E l'interprete risponde a ognuno con l'informazione appropriata, pescando dalla sua comprensione globale.
\end{notabox}

\begin{notabox}[Subsampling durante il training per le query multimodali]
Con $O \approx 80{,}000$ punti di output, calcolare il decoder completo ad ogni passo di training sarebbe troppo costoso. Il Perceiver IO adotta una strategia di \textbf{subsampling}:
\begin{itemize}
    \item \textbf{512 pixel video} campionati casualmente
    \item \textbf{512 campioni audio} campionati casualmente
    \item \textbf{1 label} (sempre presente)
    \item Totale per step di training: $O_{\text{train}} = 1{,}025$ invece di $80{,}897$
\end{itemize}
Questo funziona perch\'e le query sono \textbf{indipendenti} tra loro nel decoder (non c'\`e self-attention tra di esse), quindi il gradiente calcolato su un sottoinsieme \`e uno stimatore non distorto del gradiente completo.
\end{notabox}

\subsubsection{Query per StarCraft II}

Per StarCraft II (sostituzione del Transformer delle entit\`a in AlphaStar, Vinyals et al., 2019):

\begin{formulabox}[Query per StarCraft II]
\begin{itemize}
    \item $O = 512$ entit\`a (unit\`a di gioco)
    \item Ogni query \`e costruita dalle \textbf{feature dell'entit\`a}: tipo, posizione, stato, alleanza
    \item $\mathbf{q}_i = \text{EntityFeatures}(i) \in \mathbb{R}^E$ (le feature dell'entit\`a $i$ servono come query)
    \item Output per entit\`a: \textbf{entity embedding} aggiornato
    \item Output aggiuntivo: un \textbf{embedding 1D pooled} (media sulle entit\`a) per le decisioni globali
    \item \textbf{Nessun tuning architetturale} necessario: si sostituisce direttamente il Transformer di AlphaStar
\end{itemize}
\end{formulabox}

\subsubsection{Riepilogo delle query per task}

\begin{center}
\begin{tabular}{p{2.8cm}cccp{3.5cm}}
\toprule
\textbf{Task} & \textbf{$O$} & \textbf{$F$} & \textbf{Query} & \textbf{Note} \\
\midrule
Classificazione & 1 & $K$ classi & Appresa & Come [CLS] di BERT \\
Optical Flow & $H \times W$ & 2 & PE + input feat & Una per pixel \\
MLM (byte) & $\sim$307 & 260 & Learned PE(pos) & Solo pos.\ mascherate \\
Multi-task GLUE & 1 per task & task-dep. & Task emb + pos & 8 task embedding \\
Multimodal (video) & $50{,}176$ & 48 & PE$(x,y,t)$ + flag & Patch $4 \times 4$ \\
Multimodal (audio) & $30{,}720$ & 1 & PE$(t)$ + flag & Campioni raw \\
Multimodal (label) & 1 & 700 & Flag & Classificazione \\
StarCraft II & 512 & entity dim & Entity features & Feature dell'entit\`a \\
\bottomrule
\end{tabular}
\end{center}

\subsubsection{Decoder Cross-Attention --- Esempio numerico passo-passo}

Illustriamo il funzionamento del decoder con un esempio numerico concreto per l'\textbf{optical flow su Sintel}.

\begin{examplebox}[Esempio numerico: decoder per optical flow su Sintel]
Parametri concreti:
\begin{itemize}
    \item $O = 11{,}408$ pixel (risoluzione di training di Sintel: $\sim$107 $\times$ 107 approssimato)
    \item $N = 256$ latenti
    \item $D = 512$ dimensione dei latenti
    \item $d_k = d_v = 64$ (con $h = 8$ teste, dimensione per-testa)
    \item $E = 322$ dimensione delle query (Fourier PE + input features)
\end{itemize}

\textbf{Passo 1: Proiezione delle query}
\[
\mathbf{Q}^{\text{dec}} = \mathbf{q} \, \mathbf{W}_Q^{\text{dec}}: \quad (11{,}408 \times 322) \cdot (322 \times 64) = 11{,}408 \times 64
\]

\textbf{Passo 2: Proiezione delle chiavi e dei valori}
\[
\mathbf{K}^{\text{dec}} = \mathbf{z}^{(L)} \, \mathbf{W}_K^{\text{dec}}: \quad (256 \times 512) \cdot (512 \times 64) = 256 \times 64
\]
\[
\mathbf{V}^{\text{dec}} = \mathbf{z}^{(L)} \, \mathbf{W}_V^{\text{dec}}: \quad (256 \times 512) \cdot (512 \times 64) = 256 \times 64
\]

\textbf{Passo 3: Calcolo degli score di attenzione}
\[
\text{Score} = \frac{\mathbf{Q}^{\text{dec}} (\mathbf{K}^{\text{dec}})^\top}{\sqrt{64}}: \quad \frac{(11{,}408 \times 64) \cdot (64 \times 256)}{\sqrt{64}} = 11{,}408 \times 256
\]
Ogni pixel ha uno score di compatibilit\`a con ciascuno dei 256 latenti.

\textbf{Passo 4: Softmax riga per riga}
\[
\mathbf{A}^{\text{dec}} = \operatorname{softmax}_{\text{riga}}(\text{Score}): \quad 11{,}408 \times 256
\]
Per ogni pixel, i 256 pesi di attenzione sommano a 1.

\textbf{Passo 5: Combinazione pesata dei valori}
\[
\text{Output} = \mathbf{A}^{\text{dec}} \cdot \mathbf{V}^{\text{dec}}: \quad (11{,}408 \times 256) \cdot (256 \times 64) = 11{,}408 \times 64
\]

\textbf{Passo 6: Concatenazione delle 8 teste e proiezione finale}
\[
\text{MultiHead} = \text{Concat}(\text{head}_1, \ldots, \text{head}_8) \, \mathbf{W}_O: \quad (11{,}408 \times 512) \cdot (512 \times 512) = 11{,}408 \times 512
\]

\textbf{Passo 7: Proiezione all'output finale}
\[
\mathbf{y} = \text{MLP}(\text{MultiHead}): \quad 11{,}408 \times 512 \to 11{,}408 \times 2
\]
Ogni pixel ottiene il suo vettore di spostamento $(dx, dy)$.

\textbf{Interpretazione}: L'attenzione permette a ogni pixel di ``interrogare'' i latenti che contengono l'informazione pi\`u rilevante per la sua posizione. Un pixel vicino a un bordo in movimento attender\`a ai latenti che hanno codificato quel bordo durante l'encoding. Pixel vicini produrranno spostamenti simili perch\'e le loro query (con coordinate simili) attiveranno pattern di attenzione simili --- ottenendo cos\`i coerenza spaziale senza alcun inductive bias esplicito.
\end{examplebox}

\subsubsection{Decoder vs Average Pooling del Perceiver: confronto dettagliato}

La differenza tra il decoder del Perceiver IO e l'average pooling del Perceiver originale \`e la \textbf{chiave} per comprendere il contributo del paper.

\begin{center}
\begin{tikzpicture}[
    block/.style={draw, rounded corners, minimum width=3cm, minimum height=1cm, align=center, font=\scriptsize\bfseries},
    arr/.style={-{Stealth[length=2.5mm]}, thick},
    title/.style={font=\small\bfseries}
]
    % --- Perceiver Average Pooling (left) ---
    \begin{scope}[xshift=0cm]
        \node[title] at (2,5.5) {Perceiver: Average Pooling};

        \node[block, fill=blue!15, minimum width=4cm] (lat1) at (2,4) {Latenti $\mathbf{z}^{(L)} \in \mathbb{R}^{N \times D}$};
        \node[block, fill=gray!15] (avg) at (2,2.5) {$\bar{\mathbf{z}} = \frac{1}{N}\sum_{j=1}^{N} \mathbf{z}_j^{(L)}$};
        \node[block, fill=gray!15] (mlp1) at (2,1) {MLP};
        \node[block, fill=teal!10] (out1) at (2,-0.5) {$\mathbf{y} \in \mathbb{R}^{K}$};

        \draw[arr] (lat1) -- (avg);
        \draw[arr] (avg) -- (mlp1);
        \draw[arr] (mlp1) -- (out1);

        \node[font=\tiny, text=red!60!black, align=center] at (2, -1.5) {Output: singolo vettore\\Perde struttura spaziale\\$O = 1$ sempre};
    \end{scope}

    % --- Perceiver IO Decoder (right) ---
    \begin{scope}[xshift=8cm]
        \node[title] at (2,5.5) {Perceiver IO: Cross-Attention Decoder};

        \node[block, fill=blue!15, minimum width=4cm] (lat2) at (2,4) {Latenti $\mathbf{z}^{(L)} \in \mathbb{R}^{N \times D}$};
        \node[block, fill=purple!10, minimum width=2.5cm] (q2) at (5,2.5) {Query $\mathbf{q} \in \mathbb{R}^{O \times E}$};
        \node[block, fill=green!15, minimum width=4cm] (ca) at (2,2.5) {Cross-Attention\\$\mathbf{A} \in \mathbb{R}^{O \times N}$};
        \node[block, fill=gray!15] (mlp2) at (2,1) {MLP (point-wise)};
        \node[block, fill=teal!10, minimum width=4cm] (out2) at (2,-0.5) {$\mathbf{y} \in \mathbb{R}^{O \times F}$};

        \draw[arr] (lat2) -- (ca);
        \draw[arr] (q2) -- (ca);
        \draw[arr] (ca) -- (mlp2);
        \draw[arr] (mlp2) -- (out2);

        \node[font=\tiny, text=green!60!black, align=center] at (2, -1.5) {Output: $O$ vettori\\Preserva struttura\\$O$ arbitrario};
    \end{scope}

    % Separator
    \draw[dashed, gray, thick] (5.5, 6) -- (5.5, -2);
\end{tikzpicture}
\end{center}

\begin{notabox}[Perch\'e il decoder con cross-attention \`e superiore all'average pooling]
\begin{enumerate}
    \item \textbf{Average pooling}: $\bar{\mathbf{z}} = \frac{1}{N}\sum_j \mathbf{z}_j$ tratta tutti i latenti allo stesso modo. L'informazione posizionale e la struttura spaziale vengono completamente \textbf{distrutte}. L'output \`e sempre un singolo vettore, indipendentemente dal task.

    \item \textbf{Cross-attention decoder}: ogni query $\mathbf{q}_i$ produce un pattern di attenzione \textbf{diverso} $\boldsymbol{\alpha}_i \in \mathbb{R}^N$ sui latenti. L'output $\mathbf{y}_i = \sum_j \alpha_{ij} \mathbf{z}_j^{(L)} \mathbf{W}_V$ \`e una combinazione pesata \textbf{specifica} per quel punto. Punti di output diversi possono focalizzarsi su latenti diversi, preservando la struttura spaziale.

    \item \textbf{Generalizzazione}: l'average pooling \`e un caso speciale del decoder cross-attention, ottenuto quando tutte le query sono identiche e i pesi di attenzione sono uniformi ($\alpha_{ij} = 1/N$ per tutti $i, j$).
\end{enumerate}
\end{notabox}


\subsection{Design delle Output Query --- Elemento cruciale}

\begin{figure}[H]
\centering
\includegraphics[width=0.85\textwidth]{paper_figures/perceiver_io_fig3_output_queries.png}
\caption{Costruzione delle output query per i diversi task (Jaegle et al., 2022, Figura 3). Ogni task usa query diverse: position encoding per il linguaggio, coordinate $(x,y)$ + feature per l'optical flow, embedding per task + modality per il multimodal autoencoding, embedding per entity per StarCraft II. Per la classificazione, una singola query appresa \`e sufficiente.}
\label{fig:perceiver_io_queries}
\end{figure}

Le \textbf{output query} $\mathbf{q} \in \mathbb{R}^{O \times E}$ sono l'elemento pi\`u innovativo e flessibile del Perceiver IO. Ogni query codifica la ``domanda'' che un punto di output pone all'array latente: \emph{``Qual \`e l'informazione che mi serve?''}

Il design delle query dipende dal task e pu\`o includere diversi tipi di informazione:

\begin{definition}[Output Query]
Una output query $\mathbf{q}_i \in \mathbb{R}^E$ per il punto di output $i$ \`e un vettore che codifica \textbf{cosa} il modello deve produrre per quel punto. Tipicamente \`e costruita combinando:
\begin{itemize}
    \item \textbf{Position encoding}: codifica della posizione del punto di output (es.\ coordinate $(x,y)$ per pixel)
    \item \textbf{Task embedding}: codifica del tipo di task richiesto (es.\ ``classificazione'' o ``segmentazione'')
    \item \textbf{Modality embedding}: codifica della modalit\`a di output (es.\ ``video'', ``audio'', ``label'')
    \item \textbf{Feature dell'input}: informazioni estratte direttamente dall'input per quel punto
\end{itemize}
\end{definition}

Vediamo nel dettaglio come le query vengono costruite per diversi task:

\begin{center}
\begin{tikzpicture}[
    taskbox/.style={draw, rounded corners, minimum width=4.5cm, minimum height=2.5cm, align=left, font=\scriptsize},
    querybox/.style={draw, rounded corners, minimum width=3cm, minimum height=0.5cm, align=center, font=\tiny\ttfamily},
    arr/.style={-{Stealth[length=2mm]}, thick},
    title/.style={font=\small\bfseries}
]
    % Classification
    \node[taskbox, fill=red!5, draw=red!50!black] (cls) at (0,5.5) {
        \textbf{Classificazione}\\[3pt]
        $O = 1$ (singola query)\\
        Query: vettore appreso\\
        (come [CLS] di BERT)\\
        Output: $\mathbf{y} \in \mathbb{R}^{1 \times K}$
    };
    \node[querybox, fill=red!15] at (0,3.5) {[learned\_embedding]};

    % Optical flow
    \node[taskbox, fill=blue!5, draw=blue!50!black] (flow) at (6,5.5) {
        \textbf{Optical Flow}\\[3pt]
        $O = H \times W$ (per-pixel)\\
        Query: pos\_enc$(x,y)$\\
        \quad + feature dell'input\\
        Output: $\mathbf{y} \in \mathbb{R}^{HW \times 2}$
    };
    \node[querybox, fill=blue!15] at (6,3.5) {[pos\_enc(x,y) $\|$ input\_feat]};

    % Language MLM
    \node[taskbox, fill=green!5, draw=green!50!black] (lang) at (12,5.5) {
        \textbf{Linguaggio (MLM)}\\[3pt]
        $O = $ posizioni mascherate\\
        Query: pos\_enc(posizione)\\
        Output: logit per token\\
        $\mathbf{y} \in \mathbb{R}^{O \times |V|}$
    };
    \node[querybox, fill=green!15] at (12,3.5) {[learned\_pos\_embed(i)]};

    % Multitask
    \node[taskbox, fill=orange!5, draw=orange!50!black] (multi) at (0,0.5) {
        \textbf{Multi-task}\\[3pt]
        $O$ dipende dal task\\
        Query: task\_embed\\
        \quad + pos\_embed\\
        Output: task-specifico
    };
    \node[querybox, fill=orange!15] at (0,-1.5) {[task\_id\_embed $\|$ pos\_embed]};

    % Multimodal autoencoding
    \node[taskbox, fill=purple!5, draw=purple!50!black] (modal) at (6,0.5) {
        \textbf{Autoenc.\ Multimodale}\\[3pt]
        $O = O_v + O_a + 1$\\
        Video: pos$(x,y,t)$ + flag\\
        Audio: pos$(t)$ + flag\\
        Label: flag
    };
    \node[querybox, fill=purple!15, minimum width=4cm] at (6,-1.5) {\begin{tabular}{l}video: [pos(x,y,t) $\|$ is\_video]\\audio: [pos(t) $\|$ is\_audio]\\label: [is\_label]\end{tabular}};

    % StarCraft
    \node[taskbox, fill=teal!5, draw=teal!50!black] (sc) at (12,0.5) {
        \textbf{StarCraft II}\\[3pt]
        $O =$ num entit\`a\\
        Query: entity\_embed\\
        Output: azioni per entit\`a
    };
    \node[querybox, fill=teal!15] at (12,-1.5) {[entity\_type $\|$ entity\_feat]};

\end{tikzpicture}
\end{center}

\begin{examplebox}[Esempio: query per optical flow]
Per il task di optical flow, vogliamo predire lo spostamento $(dx, dy)$ di ogni pixel tra due frame. La query per il pixel alla posizione $(x, y)$ \`e:
\[
\mathbf{q}_{(x,y)} = \text{concat}\!\left[\text{FourierPE}(x, y), \; \mathbf{f}_{(x,y)}\right]
\]
dove $\text{FourierPE}(x, y)$ \`e un encoding posizionale di Fourier (seni e coseni a diverse frequenze) e $\mathbf{f}_{(x,y)}$ sono le feature dell'input nella posizione $(x, y)$ (ottenute dal primo frame). La query dice al decoder: ``Sono il pixel $(x,y)$ del primo frame e ho queste feature locali --- dimmi dove mi trovo nel secondo frame.''
\end{examplebox}

\begin{examplebox}[Esempio: query per classificazione]
Per la classificazione, serve un \emph{singolo} output: il vettore di probabilit\`a sulle classi. Si usa quindi una \textbf{singola query appresa} $\mathbf{q} \in \mathbb{R}^{1 \times E}$, analoga al token \texttt{[CLS]} di BERT:
\[
\mathbf{q}_{\text{cls}} = \text{parametro apprendibile} \in \mathbb{R}^{1 \times E}
\]
L'output \`e $\mathbf{y} = \text{MLP}(\text{CrossAttn}(\mathbf{q}_{\text{cls}}, \mathbf{z}^{(L)})) \in \mathbb{R}^{1 \times K}$, dove $K$ \`e il numero di classi.
\end{examplebox}

In tutti i casi, il pattern \`e lo stesso: la query codifica la \emph{struttura} dell'output desiderato. Per l'optical flow, la struttura \`e spaziale (coordinate $x, y$). Per il masked language modeling, \`e posizionale (quale byte). Per il multimodal autoencoding, \`e modale (video/audio/label). Questa flessibilit\`a \`e l'innovazione chiave del Perceiver IO: lo stesso encoder e processor gestiscono tutti i task, e solo le query cambiano.

\subsection{Complessit\`a computazionale totale}

Il Perceiver IO estende la complessit\`a del Perceiver originale (analizzata nella \hyperref[sec:complessita]{Sezione~\ref*{sec:complessita}}) aggiungendo il termine del decoder. Riassumiamo la complessit\`a dei tre moduli:

\begin{formulabox}[Complessit\`a Perceiver IO]
\begin{align}
    \text{Encoder:} \quad & \mathcal{O}(M \cdot N \cdot D) \\
    \text{Processor:} \quad & \mathcal{O}(L \cdot N^2 \cdot D) \\
    \text{Decoder:} \quad & \mathcal{O}(O \cdot N \cdot D) \\[6pt]
    \textbf{Totale:} \quad & \mathcal{O}\!\left((M + O) \cdot N \cdot D + L \cdot N^2 \cdot D\right)
\end{align}
Semplificando (omettendo $D$):
\[
\boxed{\mathcal{O}\!\left(MN + LN^2 + ON\right)}
\]
\textbf{Lineare} sia in $M$ (dimensione input) che in $O$ (dimensione output)!
\end{formulabox}

Confrontiamo con un Transformer standard che processa la concatenazione di input e output:

\begin{center}
\begin{tabular}{lcc}
\toprule
\textbf{Modello} & \textbf{Complessit\`a} & \textbf{Dipendenza da $M$, $O$} \\
\midrule
Transformer standard & $\mathcal{O}\!\left(L \cdot (M+O)^2 \cdot D\right)$ & Quadratica \\
Perceiver (solo encoder) & $\mathcal{O}\!\left(MN + LN^2\right)$ & Lineare in $M$, no output struct. \\
\textbf{Perceiver IO} & $\mathcal{O}\!\left(MN + LN^2 + ON\right)$ & \textbf{Lineare in $M$ e $O$} \\
\bottomrule
\end{tabular}
\end{center}

\begin{examplebox}[Esempio numerico di complessit\`a]
Consideriamo un task di segmentazione con immagine $224 \times 224$:
\begin{itemize}
    \item $M = 224 \times 224 = 50{,}176$ (pixel di input)
    \item $O = 50{,}176$ (pixel di output, stessa risoluzione)
    \item $N = 512$ (latenti), $D = 1024$, $L = 6$ (layer)
\end{itemize}

\textbf{Perceiver IO:}
\begin{align*}
    \text{Encoder:} &\quad 50{,}176 \times 512 \approx 25.7 \text{M operazioni} \\
    \text{Processor:} &\quad 6 \times 512^2 \approx 1.6 \text{M operazioni} \\
    \text{Decoder:} &\quad 50{,}176 \times 512 \approx 25.7 \text{M operazioni} \\
    \text{Totale:} &\quad \approx 53 \text{M operazioni}
\end{align*}

\textbf{Transformer standard} (self-attention su tutti i $100{,}352$ token):
\begin{align*}
    \text{Per layer:} &\quad (50{,}176 + 50{,}176)^2 \approx 10{,}070 \text{M operazioni} \\
    \text{Totale ($6$ layer):} &\quad \approx 60{,}423 \text{M operazioni}
\end{align*}

Il Perceiver IO \`e circa \textbf{1000 volte pi\`u efficiente}!
\end{examplebox}

\begin{center}
\begin{tikzpicture}
\begin{axis}[
    width=12cm, height=7cm,
    xlabel={Dimensione input/output $M = O$},
    ylabel={Complessit\`a (unit\`a arbitrarie, log scale)},
    ymode=log,
    xmin=100, xmax=60000,
    legend pos=north west,
    legend style={font=\small},
    grid=major,
    thick
]
    % Transformer: L*(M+O)^2 = 6*(2M)^2 = 24*M^2
    \addplot[red, very thick, domain=100:60000, samples=100] {24*x^2};
    \addlegendentry{Transformer: $\mathcal{O}(L(M\!+\!O)^2)$}

    % Perceiver IO: MN + LN^2 + ON = 2*M*512 + 6*512^2 = 1024*M + 1572864
    \addplot[blue, very thick, domain=100:60000, samples=100] {1024*x + 1572864};
    \addlegendentry{Perceiver IO: $\mathcal{O}(MN\!+\!LN^2\!+\!ON)$}

    % Reference line for N=512
    \addplot[dashed, gray, domain=100:60000] {512^2*6};
    \addlegendentry{Solo processor: $\mathcal{O}(LN^2)$}
\end{axis}
\end{tikzpicture}
\captionof{figure}{Confronto della complessit\`a computazionale al crescere della dimensione dell'input/output ($N=512$, $L=6$). Il Perceiver IO scala linearmente, il Transformer quadraticamente.}
\end{center}


%% ============================================================
%%  SEZIONE 3 --- APPLICAZIONI SPERIMENTALI
%% ============================================================
\section{Applicazioni sperimentali}

Una delle caratteristiche pi\`u impressionanti del Perceiver IO \`e la sua \textbf{generalit\`a}: lo \emph{stesso modello architetturale} ottiene risultati competitivi o state-of-the-art su task molto diversi, semplicemente cambiando le output query e la loss.

\subsection{Comprensione del linguaggio: GLUE benchmark}

\subsubsection{Setup dell'input: byte grezzi senza tokenizzazione}

Una scelta radicale del Perceiver IO per i task linguistici \`e operare direttamente sui \textbf{byte UTF-8 grezzi}, senza alcun tipo di tokenizzazione (niente WordPiece, SentencePiece, BPE). Ogni carattere viene rappresentato come il suo valore byte (un intero $0$--$255$), trasformato in un embedding appreso.

\begin{itemize}
    \item \textbf{Input}: sequenza di byte UTF-8, $M = 2048$ byte
    \item \textbf{Vocabolario}: $256$ byte + $4$ token speciali (PAD, MASK, CLS, SEP) = $260$ token totali
    \item \textbf{Embedding}: ogni byte $\to$ vettore appreso $\in \mathbb{R}^C$
    \item $N = 256$ latenti, $D = 1280$
    \item \textbf{Nessuna tokenizzazione}: il modello impara a comporre byte in ``parole'' e strutture sintattiche \emph{interamente} attraverso l'attenzione
\end{itemize}

Questa scelta ha profonde implicazioni:
\begin{enumerate}
    \item \textbf{Universalit\`a linguistica}: funziona per qualsiasi lingua e script, senza bisogno di un tokenizer specifico
    \item \textbf{Robustezza}: gestisce naturalmente errori tipografici, neologismi, codice sorgente
    \item \textbf{Sequenze pi\`u lunghe}: un testo di 512 token WordPiece corrisponde a circa $2048$ byte, quindi la sequenza di byte \`e $4\times$ pi\`u lunga, ma il Perceiver IO la gestisce comunque grazie alla complessit\`a lineare in $M$
\end{enumerate}

\subsubsection{Pre-training con Masked Language Modeling}

Il Perceiver IO viene prima pre-addestrato con \textbf{Masked Language Modeling (MLM)}, analogo a BERT ma a livello di byte:
\begin{enumerate}
    \item Si mascherano casualmente il $15\%$ dei byte di input (mascheramento a livello di \textbf{parola}, non di singolo byte: quando si sceglie di mascherare una parola, tutti i suoi byte vengono mascherati)
    \item Il modello deve predire i byte originali nelle posizioni mascherate
    \item L'\textbf{output query} per ogni posizione mascherata $i$ \`e un \textbf{embedding posizionale appreso} $\mathbf{q}_i = \text{LearnedPE}(i) \in \mathbb{R}^E$
    \item L'output per ogni posizione \`e un vettore di logit su $260$ classi ($256$ byte + $4$ token speciali)
    \item \textbf{Dataset di pre-training}: Wikipedia inglese + C4 (Colossal Clean Crawled Corpus)
\end{enumerate}

\begin{center}
\begin{tikzpicture}[
    byte/.style={draw, minimum width=0.6cm, minimum height=0.6cm, font=\tiny, inner sep=1pt},
    masked/.style={byte, fill=red!20, draw=red!50!black},
    arr/.style={-{Stealth[length=2mm]}, thick}
]
    % Input bytes
    \node[font=\small\bfseries] at (-1.5, 2) {Input:};
    \foreach \i/\c/\s in {0/T/byte, 1/h/byte, 2/e/byte, 3/ /byte, 4/c/masked, 5/a/masked, 6/t/masked, 7/ /byte, 8/s/byte, 9/a/byte, 10/t/byte} {
        \node[\s] (in\i) at (\i*0.8, 2) {\c};
    }

    % Encoder + Processor (simplified)
    \node[draw, rounded corners, fill=blue!10, minimum width=8cm, minimum height=0.8cm] (proc) at (4, 0.5) {Encoder + Processor $\to$ Latenti $\mathbf{z}^{(L)} \in \mathbb{R}^{256 \times 1280}$};

    % Output queries (only for masked positions)
    \node[font=\small\bfseries] at (-1.5, -1) {Query:};
    \node[draw, fill=purple!10, minimum width=0.6cm, minimum height=0.6cm, font=\tiny] (q4) at (4*0.8, -1) {$\mathbf{q}_4$};
    \node[draw, fill=purple!10, minimum width=0.6cm, minimum height=0.6cm, font=\tiny] (q5) at (5*0.8, -1) {$\mathbf{q}_5$};
    \node[draw, fill=purple!10, minimum width=0.6cm, minimum height=0.6cm, font=\tiny] (q6) at (6*0.8, -1) {$\mathbf{q}_6$};

    % Decoder
    \node[draw, rounded corners, fill=green!10, minimum width=3cm, minimum height=0.8cm] (dec) at (4, -2.5) {Decoder (Cross-Attention)};

    % Output
    \node[font=\small\bfseries] at (-1.5, -4) {Output:};
    \node[draw, fill=teal!10, minimum width=0.6cm, minimum height=0.6cm, font=\tiny] (o4) at (4*0.8, -4) {c};
    \node[draw, fill=teal!10, minimum width=0.6cm, minimum height=0.6cm, font=\tiny] (o5) at (5*0.8, -4) {a};
    \node[draw, fill=teal!10, minimum width=0.6cm, minimum height=0.6cm, font=\tiny] (o6) at (6*0.8, -4) {t};

    % Arrows
    \foreach \i in {0,...,10} {
        \draw[arr, gray!50, thin] (in\i) -- (proc);
    }
    \draw[arr] (proc) -- (dec);
    \draw[arr] (q4) -- (dec);
    \draw[arr] (q5) -- (dec);
    \draw[arr] (q6) -- (dec);
    \draw[arr] (dec) -- (o4);
    \draw[arr] (dec) -- (o5);
    \draw[arr] (dec) -- (o6);

    \node[font=\scriptsize, text=red!60!black] at (5*0.8, 2.7) {posizioni mascherate};
\end{tikzpicture}
\end{center}

\subsubsection{Risultati dettagliati su GLUE}

I risultati sono notevoli. Il seguente confronto dettagliato mostra le diverse configurazioni testate:

\begin{center}
\begin{tabular}{lccccc}
\toprule
\textbf{Modello} & \textbf{Tokenizer} & \textbf{$M$} & \textbf{$N$} & \textbf{Param} & \textbf{GLUE Avg} \\
\midrule
BERT Base (SentencePiece) & SentencePiece & 512 & 512 & 110M & 81.0 \\
Perceiver IO Base (SP) & SentencePiece & 512 & 256 & 223M & 81.2 \\
Perceiver IO (UTF-8) & Nessuno (byte) & 2048 & 256 & 201M & 81.0 \\
\textbf{Perceiver IO++ (UTF-8)} & Nessuno (byte) & 2048 & 256 & \textbf{425M} & \textbf{81.8} \\
\midrule
\multicolumn{6}{l}{\scriptsize BERT Base: 12 layer, 109B FLOPs} \\
\multicolumn{6}{l}{\scriptsize Perceiver IO Base (SP): 26 layer, 119B FLOPs} \\
\multicolumn{6}{l}{\scriptsize Perceiver IO (UTF-8): 26 layer, 113B FLOPs} \\
\multicolumn{6}{l}{\scriptsize Perceiver IO++ (UTF-8): 40 layer, 241B FLOPs} \\
\bottomrule
\end{tabular}
\end{center}

Risultati per task individuale della configurazione pi\`u performante:

\begin{center}
\begin{tabular}{lccccccccc}
\toprule
\textbf{Modello} & \textbf{CoLA} & \textbf{SST-2} & \textbf{MRPC} & \textbf{STS-B} & \textbf{QQP} & \textbf{MNLI} & \textbf{QNLI} & \textbf{RTE} & \textbf{Media} \\
\midrule
BERT-Base & 52.1 & 93.5 & 88.9 & 85.8 & 71.2 & 84.6 & 90.5 & 66.4 & 79.1 \\
BERT-Large & 60.5 & 94.9 & 89.3 & 86.5 & 72.1 & 86.7 & 92.7 & 70.1 & 81.0 \\
\textbf{Perceiver IO++} & \textbf{61.3} & 93.7 & \textbf{90.3} & \textbf{87.8} & \textbf{72.3} & 85.7 & 91.2 & \textbf{72.2} & \textbf{81.8} \\
\midrule
\multicolumn{10}{l}{\scriptsize\textit{Perceiver IO++: media 81.8 vs BERT-Large: 81.0 --- senza tokenizzazione!}} \\
\bottomrule
\end{tabular}
\end{center}

\begin{notabox}[Insight chiave sui risultati linguistici]
Il Perceiver IO ottiene 81.8 di media su GLUE contro il 81.0 di BERT-Large, e lo fa operando su \textbf{byte grezzi} senza alcun tokenizer. Questo dimostra che:
\begin{enumerate}
    \item La tokenizzazione \textbf{non \`e necessaria} per ottenere buone performance linguistiche
    \item L'architettura Perceiver IO \`e sufficientemente potente da apprendere la struttura sub-word direttamente dai dati
    \item L'approccio basato su byte \`e potenzialmente pi\`u \textbf{robusto} e \textbf{universale}
    \item La versione con SentencePiece (81.2) e quella con byte grezzi (81.0 con modello base, 81.8 con modello grande) ottengono performance simili, confermando che il byte-level \`e una strategia valida
\end{enumerate}
\end{notabox}

\subsection{Multi-task learning su GLUE}

Un vantaggio unico del design basato su query \`e la facilit\`a con cui si pu\`o fare \textbf{multi-task learning}. Invece di addestrare un modello separato per ciascuno degli 8 task GLUE, si addestra un \textbf{singolo modello} per tutti.

\subsubsection{Costruzione delle query multi-task}

La query per il multi-task combina due componenti:
\begin{itemize}
    \item \textbf{Task embedding} $\mathbf{e}_{\text{task}} \in \mathbb{R}^{E_t}$: un vettore appreso per ogni task (8 vettori totali)
    \item \textbf{Position/output embedding}: un vettore che indica quale output specifico del task si sta richiedendo
\end{itemize}

\[
\mathbf{q}_i^{\text{multi}} = \text{concat}\!\left[\mathbf{e}_{\text{task}}, \; \mathbf{e}_{\text{pos}}(i)\right] \in \mathbb{R}^{E}
\]

\begin{center}
\begin{tikzpicture}[
    comp/.style={draw, rounded corners, minimum height=0.8cm, align=center, font=\scriptsize},
    arr/.style={-{Stealth[length=2mm]}, thick}
]
    % Task embedding
    \node[comp, fill=orange!15, minimum width=2.5cm] (task) at (0,0) {Task Embedding\\$\mathbf{e}_{\text{MNLI}}$};

    % Position embedding
    \node[comp, fill=cyan!15, minimum width=2.5cm] (pos) at (4,0) {Position Embed\\$\mathbf{e}_{\text{pos}}(i)$};

    % Concat
    \node[comp, fill=purple!15, minimum width=5cm] (query) at (2,-1.5) {Output Query: $\mathbf{q}_i = [\mathbf{e}_{\text{MNLI}} \| \mathbf{e}_{\text{pos}}(i)]$};

    % Decoder
    \node[comp, fill=green!10, minimum width=5cm] (dec) at (2,-3) {Decoder Cross-Attention};

    % Output
    \node[comp, fill=teal!10, minimum width=3cm] (out) at (2,-4.5) {Predizione MNLI};

    \draw[arr] (task) -- (query);
    \draw[arr] (pos) -- (query);
    \draw[arr] (query) -- (dec);
    \draw[arr] (dec) -- (out);

    % Other tasks
    \node[font=\scriptsize, text=gray] at (7,-1.5) {\begin{tabular}{l}
    Stessa architettura\\per CoLA, SST-2,\\MRPC, STS-B, \ldots\\Cambia solo $\mathbf{e}_{\text{task}}$
    \end{tabular}};
\end{tikzpicture}
\end{center}

\subsubsection{Risultati multi-task dettagliati}

Il paper esplora diverse strategie di multi-task learning e confronta i risultati:

\begin{center}
\begin{tabular}{lcc}
\toprule
\textbf{Configurazione} & \textbf{GLUE Avg} & \textbf{Note} \\
\midrule
Single-task (8 modelli separati) & 81.0 & Un modello per task \\
Multi-task, decoder condiviso & 81.5 & Stesso decoder per tutti i task \\
Multi-task, decoder task-specifico & 81.8 & Decoder diverso per task \\
Multi-task, query task-specifiche & 81.8 & Stesso decoder, query diverse \\
\bottomrule
\end{tabular}
\end{center}

I risultati sono sorprendenti: il modello multi-task ottiene \textbf{81.8} di media, \emph{identico} al modello single-task pi\`u grande e \textbf{superiore} a BERT-Large (81.0). Questo suggerisce che i diversi task linguistici condividono strutture profonde che il Perceiver IO riesce a catturare in un unico spazio latente.

\begin{notabox}[Perch\'e il multi-task funziona cos\`i bene nel Perceiver IO]
La chiave \`e che encoder e processor sono \textbf{completamente condivisi} tra i task. L'unica differenziazione avviene nelle \textbf{output query}, che contengono un embedding specifico per il task. Questo significa che:
\begin{enumerate}
    \item I latenti catturano una rappresentazione \textbf{universale} del linguaggio
    \item Ogni task ``legge'' da questa rappresentazione ci\`o di cui ha bisogno, attraverso query diverse
    \item I task possono ``aiutarsi a vicenda'' durante il training: un task pu\`o forzare i latenti a catturare feature utili anche per un altro task (transfer implicito)
\end{enumerate}
\end{notabox}


\subsection{Optical Flow (Sintel e KITTI)}

L'\textbf{optical flow} \`e il task di stimare il campo di spostamento denso tra due frame consecutivi di un video. Per ogni pixel $(x, y)$ nel primo frame, il modello deve predire lo spostamento $(dx, dy)$ che lo mappa al frame successivo.

\subsubsection{Setup dettagliato}

\begin{itemize}
    \item \textbf{Input}: due frame consecutivi concatenati lungo la dimensione dei canali
    \item Per ogni pixel, si estraggono \textbf{patch $3 \times 3$} da entrambi i frame, ottenendo un vettore di $3 \times 3 \times 3 \times 2 = 54$ feature per pixel
    \item Si aggiunge un \textbf{position encoding 2D} (Fourier features) a ciascun pixel
    \item $M = H \times W$ pixel totali
    \item \textbf{Output}: $O = H \times W$, con $F = 2$ (spostamento orizzontale e verticale per pixel)
    \item \textbf{Query}: per ogni pixel $(x, y)$, la query combina:
    \begin{itemize}
        \item Fourier position encoding delle coordinate $(x, y)$ normalizzate
        \item Feature estratte dall'input in quella posizione
    \end{itemize}
\end{itemize}

\[
\mathbf{q}_{(x,y)} = \text{concat}\!\left[\text{FourierPE}\!\left(\frac{x}{W}, \frac{y}{H}\right), \; \mathbf{f}_{(x,y)}^{\text{input}}\right]
\]

\subsubsection{Risultati dettagliati}

\begin{center}
\begin{tabular}{lcccc}
\toprule
\textbf{Modello} & \textbf{Sintel Clean} & \textbf{Sintel Final} & \textbf{KITTI-15 Fl-all} & \textbf{Architettura} \\
 & EPE $\downarrow$ & EPE $\downarrow$ & \% $\downarrow$ & \\
\midrule
FlowNet2 (Ilg et al., 2017) & 3.96 & 6.02 & -- & CNN specializzata \\
PWC-Net (Sun et al., 2018) & 3.86 & 5.13 & 9.60 & CNN + cost volume \\
RAFT (Teed \& Deng, 2020) & 1.94 & 3.18 & 5.10 & CNN + iterativa \\
GMA (Jiang et al., 2021) & 1.39 & 2.47 & 5.15 & RAFT + attention \\
\midrule
\textbf{Perceiver IO} & \textbf{1.81} & \textbf{2.42} & \textbf{4.98} & General purpose \\
\bottomrule
\end{tabular}
\end{center}

\begin{notabox}[Risultato notevole per l'optical flow]
Il Perceiver IO raggiunge performance state-of-the-art su Sintel (Final) e KITTI-15, \textbf{senza} utilizzare nessuna delle componenti specializzate tradizionali:
\begin{itemize}
    \item \textbf{Nessun cost volume}: i metodi tradizionali (es.\ PWC-Net, Sun et al., 2018) costruiscono esplicitamente un volume di correlazione tra le feature dei due frame. Il Perceiver IO impara le corrispondenze implicitamente nel processor
    \item \textbf{Nessun warping}: non c'\`e allineamento iterativo dei frame
    \item \textbf{Nessuna struttura gerarchica}: non c'\`e piramide multi-scala come in PWC-Net
    \item \textbf{Nessuna convoluzione}: solo cross-attention e self-attention
    \item \textbf{Nessuna iterazione}: RAFT (Teed \& Deng, 2020) usa un meccanismo iterativo di aggiornamento, il Perceiver IO produce l'output in un singolo forward pass
\end{itemize}
Questo dimostra che un'architettura sufficientemente generale pu\`o sostituire molte componenti specializzate.
\end{notabox}

\begin{center}
\begin{tikzpicture}[
    frame/.style={draw, minimum width=3cm, minimum height=2cm, fill=gray!5},
    arr/.style={-{Stealth[length=2.5mm]}, thick}
]
    % Frame 1
    \node[frame] (f1) at (0,0) {};
    \node[font=\scriptsize\bfseries, above] at (f1.north) {Frame $t$};
    \fill[red] (0.3,0.2) circle (3pt);
    \node[font=\tiny, right] at (0.5,0.2) {$(x,y)$};

    % Frame 2
    \node[frame] (f2) at (4,0) {};
    \node[font=\scriptsize\bfseries, above] at (f2.north) {Frame $t\!+\!1$};
    \fill[blue] (0.9+4,-0.3) circle (3pt);
    \node[font=\tiny, right] at (1.1+4,-0.3) {$(x\!+\!dx, y\!+\!dy)$};

    % Flow arrow
    \draw[arr, red!70!black, dashed] (0.3,0.2) -- (4.9,-0.3) node[midway, above, font=\tiny] {$(dx, dy)$};

    % Perceiver IO pipeline
    \node[draw, rounded corners, fill=orange!10, minimum width=6cm, minimum height=0.8cm] (enc) at (2,-2.5) {Encoder: cross-attention su patch $3\times3$ concatenate};
    \node[draw, rounded corners, fill=blue!10, minimum width=6cm, minimum height=0.8cm] (proc) at (2,-3.7) {Processor: $L$ layer self-attention sui latenti};
    \node[draw, rounded corners, fill=green!10, minimum width=6cm, minimum height=0.8cm] (dec) at (2,-4.9) {Decoder: query = FourierPE$(x,y)$ + input features};
    \node[draw, rounded corners, fill=teal!10, minimum width=3cm, minimum height=0.8cm, font=\scriptsize] (out) at (2,-6.1) {Output: $(dx, dy)$ per pixel};

    \draw[arr] (2,-1) -- (enc);
    \draw[arr] (enc) -- (proc);
    \draw[arr] (proc) -- (dec);
    \draw[arr] (dec) -- (out);
\end{tikzpicture}
\end{center}


\subsection{Autoencoding multimodale (Kinetics-700)}

Il task di autoencoding multimodale \`e forse la dimostrazione pi\`u impressionante della flessibilit\`a del Perceiver IO: un \textbf{singolo modello} riceve in input video, audio e label di classe, e deve ricostruire \textbf{tutte e tre le modalit\`a simultaneamente}.

\subsubsection{Setup dell'input dettagliato}

L'input \`e la concatenazione di tre modalit\`a, serializzate in un unico array 2D con embedding di modalit\`a:
\begin{itemize}
    \item \textbf{Video}: 16 frame a risoluzione $224 \times 224$, suddivisi in patch $4 \times 4$. Ogni frame ha $56 \times 56 = 3{,}136$ patch, per un totale di $16 \times 3{,}136 = 50{,}176$ patch. Ogni patch ha $3 \times 16 = 48$ feature (RGB per i 16 pixel della patch)
    \item \textbf{Audio}: $30{,}720$ campioni audio (circa 2 secondi a 16kHz), ciascuno rappresentato da un singolo scalare
    \item \textbf{Label di classe}: un singolo token one-hot su $700$ classi di Kinetics
\end{itemize}

L'input totale ha $M = 50{,}176 + 30{,}720 + 1 = 80{,}897$ elementi. Ciascun elemento viene arricchito con:
\begin{enumerate}
    \item Un \textbf{embedding di modalit\`a} appreso (``questa \`e una patch video'' / ``questo \`e un campione audio'' / ``questa \`e la label'')
    \item Un \textbf{encoding posizionale} appropriato (Fourier 3D per video, Fourier 1D per audio, nessuno per label)
\end{enumerate}

Gestire un input cos\`i grande con self-attention standard richiederebbe $\mathcal{O}(M^2) \approx 6.5 \times 10^9$ operazioni per layer, impraticabile.

\subsubsection{Design delle output query multimodali}

Per ricostruire le tre modalit\`a, servono query specifiche per ciascuna:

\begin{formulabox}[Query multimodali per autoencoding]
\begin{align}
    \text{Query video:} \quad & \mathbf{q}_v = \text{concat}\!\left[\text{FourierPE}(x, y, t), \; \mathbf{e}_{\text{video}}\right] \in \mathbb{R}^E \\
    \text{Query audio:} \quad & \mathbf{q}_a = \text{concat}\!\left[\text{FourierPE}(t), \; \mathbf{e}_{\text{audio}}\right] \in \mathbb{R}^E \\
    \text{Query label:} \quad & \mathbf{q}_l = \mathbf{e}_{\text{label}} \in \mathbb{R}^E
\end{align}
dove $\mathbf{e}_{\text{video}}, \mathbf{e}_{\text{audio}}, \mathbf{e}_{\text{label}}$ sono embedding appresi specifici per modalit\`a (``flag'' di modalit\`a).
\end{formulabox}

\begin{center}
\begin{tikzpicture}[
    modbox/.style={draw, rounded corners, minimum width=2.2cm, minimum height=1.2cm, align=center, font=\scriptsize},
    querybox/.style={draw, rounded corners, minimum width=3.5cm, minimum height=0.8cm, align=center, font=\tiny},
    arr/.style={-{Stealth[length=2mm]}, thick}
]
    % Input modalities
    \node[modbox, fill=red!10] (vid_in) at (0,3) {Video\\$50{,}176$ patch};
    \node[modbox, fill=yellow!10] (aud_in) at (3,3) {Audio\\$30{,}720$ sample};
    \node[modbox, fill=green!10] (lab_in) at (6,3) {Label\\$1$ token};

    % Encoder
    \node[draw, rounded corners, fill=orange!10, minimum width=8cm, minimum height=0.8cm] (enc) at (3,1.5) {Encoder: concatena e cross-attention};

    % Processor
    \node[draw, rounded corners, fill=blue!10, minimum width=8cm, minimum height=0.8cm] (proc) at (3,0.3) {Processor: self-attention $\times L$ su $N$ latenti};

    % Decoder
    \node[draw, rounded corners, fill=green!10, minimum width=8cm, minimum height=0.8cm] (dec) at (3,-1) {Decoder: cross-attention con query multimodali};

    % Output queries
    \node[querybox, fill=red!5] (qv) at (-2,-1) {$\mathbf{q}_v$: PE$(x,y,t)$ + is\_video};
    \node[querybox, fill=yellow!5] (qa) at (-2,-2) {$\mathbf{q}_a$: PE$(t)$ + is\_audio};
    \node[querybox, fill=green!5] (ql) at (-2,-3) {$\mathbf{q}_l$: is\_label};

    % Output modalities
    \node[modbox, fill=red!15] (vid_out) at (0,-4.5) {Video\\ricostruito};
    \node[modbox, fill=yellow!15] (aud_out) at (3,-4.5) {Audio\\ricostruito};
    \node[modbox, fill=green!15] (lab_out) at (6,-4.5) {Label\\predetta};

    % Arrows
    \draw[arr] (vid_in) -- (enc);
    \draw[arr] (aud_in) -- (enc);
    \draw[arr] (lab_in) -- (enc);
    \draw[arr] (enc) -- (proc);
    \draw[arr] (proc) -- (dec);
    \draw[arr] (qv) -- (dec);
    \draw[arr] (qa) -- (dec);
    \draw[arr] (ql) -- (dec);
    \draw[arr] (dec) -- (vid_out);
    \draw[arr] (dec) -- (aud_out);
    \draw[arr] (dec) -- (lab_out);
\end{tikzpicture}
\end{center}

\subsubsection{Loss multimodale e mascheramento}

La loss per l'autoencoding multimodale \`e una combinazione pesata delle loss per ciascuna modalit\`a:
\[
\mathcal{L}_{\text{multi}} = \lambda_v \mathcal{L}_{\text{video}} + \lambda_a \mathcal{L}_{\text{audio}} + \lambda_l \mathcal{L}_{\text{label}}
\]
dove:
\begin{itemize}
    \item $\mathcal{L}_{\text{video}} = \frac{1}{|\mathcal{S}_v|}\sum_{i \in \mathcal{S}_v} \|\hat{\mathbf{y}}_i^v - \mathbf{y}_i^{v*}\|_2^2$ (\textbf{MSE} sulla ricostruzione dei pixel)
    \item $\mathcal{L}_{\text{audio}} = \frac{1}{|\mathcal{S}_a|}\sum_{i \in \mathcal{S}_a} (\hat{y}_i^a - y_i^{a*})^2$ (\textbf{MSE} sulla ricostruzione dei campioni audio)
    \item $\mathcal{L}_{\text{label}} = -\sum_{k=1}^{700} y_k^* \log \hat{y}_k$ (\textbf{cross-entropy} per la classificazione)
\end{itemize}

\begin{notabox}[Mascheramento della label al 50\%]
Per evitare che il modello ``barini'' (semplicemente copiando la label dall'input all'output senza elaborare video e audio), la \textbf{label di classe viene mascherata il 50\% delle volte} durante il training. Quando la label \`e mascherata:
\begin{itemize}
    \item Il modello deve predire la classe solo dal video e dall'audio
    \item Questo forza i latenti a catturare informazione semantica dalle modalit\`a percettive
    \item Impedisce shortcut dove i latenti memorizzano solo la label ignorando il contenuto
\end{itemize}
\end{notabox}

\subsubsection{Risultati dell'autoencoding multimodale}

I risultati vengono misurati in termini di PSNR (Peak Signal-to-Noise Ratio) per video e audio, e accuracy per la classificazione. I rapporti di compressione testati sono $88\times$, $176\times$ e $352\times$ (controllati dal numero di latenti $N$).

\begin{center}
\begin{tabular}{lcccc}
\toprule
\textbf{Compressione} & \textbf{$N$ latenti} & \textbf{Video PSNR} & \textbf{Audio PSNR} & \textbf{Classif.\ Acc.} \\
\midrule
$88\times$ & $N = 920$ & 33.8 dB & 36.2 dB & 45.6\% \\
$176\times$ & $N = 460$ & 32.5 dB & 34.8 dB & 44.2\% \\
$352\times$ & $N = 230$ & 30.1 dB & 32.1 dB & 41.0\% \\
\bottomrule
\end{tabular}
\end{center}

\begin{notabox}[Propriet\`a notevole dell'autoencoder multimodale]
Il fatto che un \textbf{singolo} array latente $\mathbf{z} \in \mathbb{R}^{N \times D}$ contenga informazioni sufficienti per ricostruire video, audio E label di classe simultaneamente dimostra che:
\begin{enumerate}
    \item I latenti del Perceiver IO formano una \textbf{rappresentazione condivisa} cross-modale
    \item Le diverse modalit\`a possono ``aiutarsi a vicenda'' nella ricostruzione (es.\ il suono di un applauso aiuta a ricostruire il video dell'evento)
    \item Il decoder basato su query \`e sufficientemente flessibile da estrarre informazioni specifiche per ciascuna modalit\`a dallo stesso spazio latente
    \item Il tradeoff compressione-qualit\`a \`e controllabile direttamente dal numero di latenti $N$
\end{enumerate}
\end{notabox}


\subsection{Classificazione ImageNet --- Risultati dettagliati}

Il Perceiver IO viene testato estensivamente sulla classificazione di immagini su ImageNet (Deng et al., 2009). La Tabella~7 del paper originale (Jaegle et al., 2022) confronta numerose configurazioni:

\begin{center}
\begin{tabular}{lcccc}
\toprule
\textbf{Modello} & \textbf{Pre-training} & \textbf{Pos.\ Encoding} & \textbf{Top-1 \%} & \textbf{Note} \\
\midrule
\multicolumn{5}{l}{\textit{Baseline ConvNet}} \\
ResNet-50 (He et al., 2016) & ImageNet & -- & 78.6 & CNN classica \\
NFNet-F6+SAM (Brock et al., 2021) & ImageNet & -- & 86.5 & SOTA CNN \\
\midrule
\multicolumn{5}{l}{\textit{Baseline ViT}} \\
ViT-B/16 (Dosovitskiy et al., 2021) & ImageNet & Learned & 77.9 & Senza JFT \\
ViT-H/14 (Dosovitskiy et al., 2021) & JFT-300M & Learned & 88.6 & Con pre-training \\
\midrule
\multicolumn{5}{l}{\textit{Perceiver (originale)}} \\
Perceiver & ImageNet & 2D Fourier & 78.6 & Come ResNet-50 \\
\midrule
\multicolumn{5}{l}{\textit{Perceiver IO}} \\
Config A (learned pos) & ImageNet & Learned & 67.6 & Performance bassa \\
Config A (2D Fourier) & ImageNet & 2D Fourier & \textbf{79.0} & Sopra Perceiver \\
Config B (learned pos) & JFT-300M & Learned & 72.7 & Migliora con dati \\
Config B (2D Fourier) & JFT-300M & 2D Fourier & \textbf{84.5} & Forte transfer \\
\midrule
\multicolumn{5}{l}{\textit{Con preprocessing convolutivo}} \\
Config + Conv (ImageNet) & ImageNet & Conv embed & 77.4 & Come ViT ibrido \\
Config + Conv (JFT) & JFT-300M & Conv embed & 82.1 & -- \\
Config + Conv (JFT, grande) & JFT-300M & Conv embed & \textbf{86.4} & Vicino a NFNet \\
\bottomrule
\end{tabular}
\end{center}

\begin{notabox}[Osservazioni sui risultati ImageNet]
\begin{enumerate}
    \item Il \textbf{Fourier positional encoding} \`e cruciale: con learned positions la performance crolla (67.6\% vs 79.0\%). Questo perch\'e il Fourier PE codifica l'inductive bias della struttura 2D dell'immagine, che con posizioni apprese deve essere appresa dai dati.
    \item Il \textbf{pre-training su JFT-300M} (dataset di 300M di immagini, 18k classi) porta un grande miglioramento: da 79.0\% a 84.5\%.
    \item L'aggiunta di un \textbf{preprocessing convolutivo} (embedding delle patch con una piccola CNN prima dell'encoder) porta a 86.4\%, quasi raggiungendo NFNet (86.5\%).
    \item Senza alcun inductive bias (n\'e Fourier PE n\'e convoluzioni), il modello ha bisogno di molti pi\`u dati per raggiungere buone performance.
\end{enumerate}
\end{notabox}


\subsection{StarCraft II: sostituzione del Transformer di AlphaStar}

Un'applicazione particolarmente interessante \`e l'uso del Perceiver IO come \textbf{sostituto del Transformer delle entit\`a} in AlphaStar (Vinyals et al., 2019), l'agente di DeepMind per StarCraft II.

\subsubsection{Setup}

\begin{itemize}
    \item In StarCraft II, il numero di entit\`a varia dinamicamente (fino a $\sim$512 unit\`a)
    \item AlphaStar originale usa un Transformer per elaborare le feature delle entit\`a
    \item Il Perceiver IO sostituisce questo Transformer: le feature delle entit\`a vengono usate sia come input sia come query di output
    \item L'output include \textbf{embedding per-entit\`a} (per decisioni a livello di singola unit\`a) e un \textbf{embedding globale} (media sulle entit\`a, per decisioni strategiche)
\end{itemize}

\subsubsection{Risultati dettagliati}

\begin{center}
\begin{tabular}{lccc}
\toprule
\textbf{Modello} & \textbf{Win Rate} & \textbf{FLOPs} & \textbf{Parametri} \\
\midrule
AlphaStar (Transformer) & 87\% & 3.3B & 144M \\
\textbf{AlphaStar (Perceiver IO)} & \textbf{87\%} & \textbf{0.93B} & \textbf{140M} \\
\midrule
\multicolumn{4}{l}{\scriptsize Riduzione FLOPs: $3.5\times$ con performance identica} \\
\bottomrule
\end{tabular}
\end{center}

\begin{notabox}[Punti chiave per StarCraft II]
\begin{enumerate}
    \item \textbf{Win rate identico}: 87\% per entrambe le configurazioni. La qualit\`a delle decisioni non \`e compromessa.
    \item \textbf{3.5x meno FLOPs}: $0.93$B vs $3.3$B. Il Perceiver IO \`e molto pi\`u efficiente perch\'e usa un bottleneck latente di dimensione fissa.
    \item \textbf{Parametri simili}: $140$M vs $144$M. Non c'\`e un risparmio significativo in termini di parametri, ma il risparmio computazionale \`e enorme.
    \item \textbf{Nessun tuning architetturale}: il Perceiver IO \`e stato inserito come ``drop-in replacement'' senza modifiche specifiche per il dominio.
    \item Il numero variabile di entit\`a non \`e un problema: il bottleneck latente ha dimensione fissa $N$, indipendente dal numero di entit\`a.
\end{enumerate}
\end{notabox}


\subsection{AudioSet (classificazione audio)}

Il Perceiver IO \`e stato testato anche su \textbf{AudioSet} (Gemmeke et al., 2017), un dataset di classificazione audio su larga scala con 527 classi:

\begin{itemize}
    \item \textbf{Input}: spettrogramma mel (128 bande di frequenza $\times$ $T$ frame temporali)
    \item \textbf{Query}: singola query appresa (come per ImageNet)
    \item \textbf{Output}: vettore di probabilit\`a su 527 classi (multi-label)
    \item Risultati: \textbf{mAP = 0.421} (competitivo con modelli specializzati per audio)
\end{itemize}


%% ============================================================
%%  SEZIONE 3.5 --- TABELLA RIASSUNTIVA DEI TASK
%% ============================================================
\subsection{Tabella riassuntiva di tutti i task (dall'Appendice del paper)}

La seguente tabella riassume le dimensioni di input e output per tutti i task sperimentali, corrispondente alla Table~5 del paper originale (Jaegle et al., 2022):

\begin{center}
\footnotesize
\begin{tabular}{p{2.5cm}ccccccc}
\toprule
\textbf{Task} & \textbf{$M$} & \textbf{$C$} & \textbf{$N$} & \textbf{$D$} & \textbf{$O$} & \textbf{$F$} & \textbf{$L$} \\
\midrule
GLUE (byte) & 2,048 & 260 & 256 & 1,280 & variabile & task-dep & 26/40 \\
Optical flow & $H \times W$ & 54 & 2,048 & 512 & $H \times W$ & 2 & 24 \\
Multimodal (K700) & $\sim$80k & 48/1/700 & 512 & 1,024 & $\sim$80k & 48/1/700 & 8 \\
ImageNet & 50,176 & 3/768 & 512 & 1,024 & 1 & 1,000 & 48 \\
StarCraft II & $\leq$512 & entity dim & 256 & 512 & $\leq$512 & entity dim & 8 \\
AudioSet & $128 \times T$ & 1 & 512 & 1,024 & 1 & 527 & 8 \\
\bottomrule
\end{tabular}
\end{center}

\begin{notabox}[Struttura degli embedding posizionali e di task (Table 6 del paper)]
Gli encoding posizionali e gli embedding di task usati per i diversi esperimenti:
\begin{itemize}
    \item \textbf{GLUE input}: embedding posizionale appreso (2048 posizioni) + embedding del segmento (per distinguere le due frasi in task di coppia)
    \item \textbf{GLUE output}: embedding posizionale appreso + embedding di task (per multi-task)
    \item \textbf{Optical flow input}: Fourier PE 2D (seni e coseni a frequenze $\{1, 2, 4, \ldots, 64\}$)
    \item \textbf{Optical flow output}: Fourier PE 2D + feature dell'input (patch $3 \times 3$)
    \item \textbf{Multimodal input}: Fourier PE (3D per video, 1D per audio) + embedding di modalit\`a
    \item \textbf{Multimodal output}: Fourier PE (3D/1D) + embedding di modalit\`a
    \item \textbf{ImageNet input}: Fourier PE 2D (o learned, o convolutivo)
    \item \textbf{ImageNet output}: embedding appreso (singola query per classificazione)
    \item \textbf{StarCraft II}: feature dell'entit\`a (tipo, posizione, stato, alleanza)
\end{itemize}
\end{notabox}


%% ============================================================
%%  SEZIONE 4 --- CONFRONTO CON IL PERCEIVER ORIGINALE
%% ============================================================
\section{Confronto con il Perceiver originale}

\begin{center}
\begin{tabular}{p{3.5cm}p{5cm}p{5cm}}
\toprule
\textbf{Caratteristica} & \textbf{Perceiver (2021)} & \textbf{Perceiver IO (2022)} \\
\midrule
\textbf{Flessibilit\`a input} & Input arbitrari (immagini, audio, point cloud, multimodale) & Identica \\
\midrule
\textbf{Flessibilit\`a output} & \textbf{Solo classificazione} (singolo vettore) & \textbf{Output strutturati arbitrari} (per-pixel, sequenze, multimodale) \\
\midrule
\textbf{Meccanismo decoder} & Average pooling $\to$ MLP & Cross-attention con output query \\
\midrule
\textbf{Dimensione output $O$} & Fissa: $O = K$ (classi) & Variabile: $O$ qualsiasi \\
\midrule
\textbf{Weight sharing} & S\`i (encoder ripetuto) & Opzionale \\
\midrule
\textbf{Complessit\`a totale} & $\mathcal{O}(MN + LN^2)$ & $\mathcal{O}(MN + LN^2 + ON)$ \\
\midrule
\textbf{Applicazioni} & Classificazione immagini, audio, point cloud & Classificazione, optical flow, NLU, multi-task, autoencoding multimodale, giochi \\
\bottomrule
\end{tabular}
\end{center}

\begin{center}
\begin{tikzpicture}[
    block/.style={draw, rounded corners, minimum width=2.2cm, minimum height=0.9cm, align=center, font=\tiny\bfseries},
    arr/.style={-{Stealth[length=2mm]}, thick},
    title/.style={font=\normalsize\bfseries}
]
    % Perceiver
    \begin{scope}[xshift=0cm]
        \node[title] at (2,6) {Perceiver};

        \node[block, fill=red!10] (p_in) at (2,5) {Input $\mathbb{R}^{M \times C}$};
        \node[block, fill=orange!15] (p_ca1) at (2,3.8) {Cross-Att};
        \node[block, fill=blue!15] (p_sa1) at (2,2.6) {Self-Att $\times L$};

        % Weight sharing arrows
        \draw[arr, dotted, red!50!black] (p_sa1.west) to[out=180,in=180] (p_ca1.west);
        \node[font=\tiny, text=red!50!black, left=0.5cm of p_sa1] {$\times R$};

        \node[block, fill=gray!15] (p_pool) at (2,1.2) {Avg Pool};
        \node[block, fill=gray!15] (p_mlp) at (2,0) {MLP};
        \node[block, fill=teal!10] (p_out) at (2,-1.2) {Classe $\mathbb{R}^K$};

        \draw[arr] (p_in) -- (p_ca1);
        \draw[arr] (p_ca1) -- (p_sa1);
        \draw[arr] (p_sa1) -- (p_pool);
        \draw[arr] (p_pool) -- (p_mlp);
        \draw[arr] (p_mlp) -- (p_out);
    \end{scope}

    % Arrow between
    \draw[very thick, -{Stealth[length=4mm]}, blue!60!black] (5, 2.5) -- (7, 2.5) node[midway, above, font=\small] {evoluzione};

    % Perceiver IO
    \begin{scope}[xshift=8cm]
        \node[title] at (2,6) {Perceiver IO};

        \node[block, fill=red!10] (io_in) at (2,5) {Input $\mathbb{R}^{M \times C}$};
        \node[block, fill=orange!15] (io_enc) at (2,3.8) {Encoder\\(Cross-Att)};
        \node[block, fill=blue!15] (io_proc) at (2,2.6) {Processor\\(Self-Att $\times L$)};
        \node[block, fill=green!15] (io_dec) at (2,1.2) {Decoder\\(Cross-Att)};
        \node[block, fill=teal!10, minimum width=3cm] (io_out) at (2,-0.2) {Output $\mathbb{R}^{O \times F}$};
        \node[block, fill=purple!10] (io_q) at (5,1.2) {Query\\$\mathbb{R}^{O \times E}$};

        \draw[arr] (io_in) -- (io_enc);
        \draw[arr] (io_enc) -- (io_proc);
        \draw[arr] (io_proc) -- (io_dec);
        \draw[arr] (io_dec) -- (io_out);
        \draw[arr] (io_q) -- (io_dec);
    \end{scope}
\end{tikzpicture}
\end{center}


%% ============================================================
%%  SEZIONE 4.5 --- CONFRONTO COMPLETO CON ALTRE ARCHITETTURE
%% ============================================================
\section{Confronto completo con altre architetture}

\begin{center}
\footnotesize
\begin{tabular}{p{2.8cm}p{2.6cm}p{2.6cm}p{2.6cm}p{2.6cm}}
\toprule
\textbf{Feature} & \textbf{Transformer} (Vaswani et al., 2017) & \textbf{ViT} (Dosovitskiy et al., 2021) & \textbf{Perceiver} (Jaegle et al., 2021) & \textbf{Perceiver IO} (Jaegle et al., 2022) \\
\midrule
\textbf{Complessit\`a} & $\mathcal{O}(n^2 d)$ quadratica & $\mathcal{O}(n^2 d)$ quadratica & $\mathcal{O}(MN + LN^2)$ lineare in $M$ & $\mathcal{O}(MN + LN^2 + ON)$ lineare in $M$ e $O$ \\
\midrule
\textbf{Input} & Sequenze tokenizzate & Patch di immagini & Byte array arbitrari & Byte array arbitrari \\
\midrule
\textbf{Output} & Sequenze (autoregressivo o per-token) & Singola classe (CLS token) & Singola classe (avg pool) & \textbf{Arbitrario} (via query) \\
\midrule
\textbf{Flessibilit\`a output} & Limitata a sequenze & Solo classificazione & Solo classificazione & Per-pixel, sequenze, multi-modale, multi-task \\
\midrule
\textbf{Pos.\ Encoding} & Sinusoidale / appreso & Appreso (per patch) & Fourier / appreso & Fourier / appreso / conv \\
\midrule
\textbf{Bottleneck} & Nessuno & Nessuno & Latenti $\mathbb{R}^{N \times D}$ & Latenti $\mathbb{R}^{N \times D}$ \\
\midrule
\textbf{Scalabilit\`a input} & Limitata ($\leq 512$--$2048$ token) & Limitata (risoluzione fissa) & Illimitata & Illimitata \\
\midrule
\textbf{Multi-task} & Teste separate & Teste separate & Non supportato (no output strutturati) & Naturale (via query) \\
\midrule
\textbf{Multi-modale} & Architetture separate & Solo immagini & Input multimodale, output singolo & Input E output multimodali \\
\midrule
\textbf{Inductive bias} & Nessuno (universale) & Patch 2D & Nessuno & Nessuno (nelle query) \\
\bottomrule
\end{tabular}
\end{center}


%% ============================================================
%%  SEZIONE 5 --- COLLEGAMENTO CON I CONCETTI DEL CORSO
%% ============================================================
\section{Collegamento con i concetti del corso}

Il Perceiver IO non nasce in un vuoto: incorpora e generalizza molte idee che abbiamo studiato durante il corso. Vediamo come ogni concetto si collega.

\subsection{Architettura encoder-decoder (Lezioni 17, 19)}

Il paradigma \textbf{encoder-decoder} \`e un tema ricorrente nel corso:
\begin{itemize}
    \item \textbf{Seq2Seq} (Lezione 17): encoder RNN che comprime l'input in un vettore, decoder RNN che genera l'output. Il Perceiver IO generalizza questo concetto: l'encoder comprime in un \emph{array} latente (non un singolo vettore), e il decoder usa cross-attention (non generazione autoregressiva).
    \item \textbf{Transformer} (Lezione 19): encoder con self-attention + decoder con cross-attention e self-attention causale. Il Perceiver IO semplifica: il decoder non ha self-attention tra le query (troppo costoso) e non \`e autoregressivo.
\end{itemize}

\begin{center}
\begin{tikzpicture}[
    block/.style={draw, rounded corners, minimum width=2.5cm, minimum height=0.7cm, align=center, font=\tiny},
    arr/.style={-{Stealth[length=2mm]}, thick}
]
    % Seq2Seq
    \node[font=\small\bfseries] at (0,3) {Seq2Seq};
    \node[block, fill=orange!10] (s_enc) at (0,2) {RNN Encoder};
    \node[block, fill=gray!10] (s_h) at (0,1) {Vettore $\mathbf{h}$};
    \node[block, fill=green!10] (s_dec) at (0,0) {RNN Decoder\\(autoregressivo)};
    \draw[arr] (s_enc) -- (s_h);
    \draw[arr] (s_h) -- (s_dec);

    % Transformer
    \node[font=\small\bfseries] at (4.5,3) {Transformer};
    \node[block, fill=orange!10] (t_enc) at (4.5,2) {Self-Att Encoder};
    \node[block, fill=gray!10] (t_h) at (4.5,1) {Sequenza $\mathbf{H}$};
    \node[block, fill=green!10] (t_dec) at (4.5,0) {Cross-Att + Self-Att\\Decoder (causale)};
    \draw[arr] (t_enc) -- (t_h);
    \draw[arr] (t_h) -- (t_dec);

    % Perceiver IO
    \node[font=\small\bfseries] at (9,3) {Perceiver IO};
    \node[block, fill=orange!10] (p_enc) at (9,2) {Cross-Att Encoder};
    \node[block, fill=blue!10] (p_lat) at (9,1) {Latenti $\mathbf{z}$\\+ Self-Att Processor};
    \node[block, fill=green!10] (p_dec) at (9,0) {Cross-Att Decoder\\(parallelo, con query)};
    \draw[arr] (p_enc) -- (p_lat);
    \draw[arr] (p_lat) -- (p_dec);

    % Arrows of evolution
    \draw[-{Stealth}, dashed, gray] (1.5,1) -- (3,1);
    \draw[-{Stealth}, dashed, gray] (6.2,1) -- (7.5,1);
\end{tikzpicture}
\end{center}

\subsection{Attention come ``soft dictionary'' (Lezione 18)}

Nella Lezione 18 abbiamo visto l'attenzione come un \textbf{dizionario soft}: le query interrogano un database di coppie chiave-valore. Nel Perceiver IO, questa metafora diventa particolarmente chiara:

\begin{itemize}
    \item \textbf{Encoder}: i latenti (query) interrogano l'input (database) per raccogliere informazioni
    \item \textbf{Decoder}: le output query interrogano i latenti (database processato) per estrarre le risposte
\end{itemize}

Il Perceiver IO rende esplicita la separazione tra ``database'' (latenti) e ``interrogazione'' (query), molto pi\`u di un Transformer standard dove le stesse rappresentazioni servono simultaneamente come query, chiavi e valori.

\subsection{Multi-task learning e transfer (Lezione 20)}

Il \textbf{multi-task learning} trattato nella Lezione 20 trova nel Perceiver IO un'implementazione particolarmente elegante. Nei modelli tradizionali, il multi-task si realizza con:
\begin{itemize}
    \item Trunk condiviso + teste task-specifiche (hard parameter sharing)
    \item O reti separate con regolarizzazione incrociata (soft parameter sharing)
\end{itemize}

Nel Perceiver IO, il multi-task \`e gestito interamente attraverso le \textbf{output query}: encoder e processor sono completamente condivisi, e le query diversificano l'output per task. Non servono teste task-specifiche aggiuntive.

Un esempio concreto: nel benchmark GLUE, il Perceiver IO affronta task diversi --- sentiment analysis (SST-2), similarit\`a semantica (STS-B), inferenza (MNLI) --- con lo \emph{stesso} modello. Per SST-2 (classificazione binaria) si usa una singola query appresa, analoga al token \texttt{[CLS]}. Per STS-B (regressione) si usa una query che produce un singolo scalare. Per MNLI (classificazione a tre classi) la query produce tre logit. L'encoder e il processor restano identici: cambia solo la query di output e la testa finale.

\subsection{Confronto con BERT (Lezione 21)}

La Lezione 21 ha trattato BERT in dettaglio. Il Perceiver IO offre un confronto illuminante:

\begin{center}
\begin{tabular}{lll}
\toprule
\textbf{Aspetto} & \textbf{BERT} & \textbf{Perceiver IO} \\
\midrule
Tokenizzazione & WordPiece (30k vocab) & Nessuna (byte grezzi) \\
Input sequenza & $\leq 512$ token & $\leq 2048$ byte \\
Architettura & Transformer encoder & Encode-Process-Decode \\
Pre-training & MLM + NSP & MLM (su byte) \\
Output struttura & Token [CLS] o per-token & Query-based, arbitrario \\
Multi-task & Fine-tuning separato & Query diverse, modello unico \\
Complessit\`a & $\mathcal{O}(n^2)$ nella sequenza & $\mathcal{O}(nN)$ lineare \\
GLUE score & 81.0 (Large) & 81.8 \\
\bottomrule
\end{tabular}
\end{center}

\subsection{Autoencoding (Lezione 21)}

Il task di \textbf{autoencoding multimodale} del Perceiver IO si collega direttamente agli autoencoder trattati nella Lezione 21:
\begin{itemize}
    \item \textbf{Autoencoder classico}: encoder $\to$ collo di bottiglia $\to$ decoder, con ricostruzione dell'input
    \item \textbf{Perceiver IO autoencoder}: encoder (cross-attention) $\to$ latenti (bottleneck) $\to$ decoder (cross-attention con query), con ricostruzione \emph{multimodale}
\end{itemize}

La differenza chiave \`e che il bottleneck del Perceiver IO \`e \textbf{strutturato} (un array $N \times D$, non un singolo vettore) e il decoder \`e basato su query, permettendo di ricostruire output di dimensione e natura diversa dall'input.

\subsection{Analisi della complessit\`a dell'attention (Lezioni 14, 19)}

Le Lezioni 14 e 19 hanno analizzato la complessit\`a computazionale dell'attenzione. Il Perceiver IO rappresenta una delle soluzioni pi\`u eleganti al problema della complessit\`a quadratica:

\begin{itemize}
    \item \textbf{Approcci linear attention} (es.\ Performer): approssimano la matrice di attenzione
    \item \textbf{Approcci sparse attention} (es.\ Longformer): limitano le connessioni
    \item \textbf{Perceiver IO}: usa un \textbf{bottleneck latente} per disaccoppiare input e output, ottenendo complessit\`a lineare \emph{esatta} (non approssimata)
\end{itemize}


%% ============================================================
%%  SEZIONE 6 --- BACKWARD PASS E TRAINING
%% ============================================================
\section{Backward pass e dettagli di training}

\subsection{Loss per i diversi task}

La funzione di loss dipende dal task specifico:

\begin{formulabox}[Loss per task]
\begin{itemize}
    \item \textbf{Classificazione} (ImageNet, label in multimodal):
    \[
    \mathcal{L}_{\text{cls}} = -\sum_{k=1}^{K} y_k \log \hat{y}_k \quad \text{(cross-entropy)}
    \]

    \item \textbf{Masked Language Modeling}:
    \[
    \mathcal{L}_{\text{MLM}} = -\sum_{i \in \mathcal{M}} \log p(x_i \mid \mathbf{x}_{\setminus\mathcal{M}}) \quad \text{(cross-entropy su byte mascherati)}
    \]

    \item \textbf{Optical flow}:
    \[
    \mathcal{L}_{\text{flow}} = \frac{1}{|\mathcal{P}|}\sum_{(x,y) \in \mathcal{P}} \left\|\hat{\mathbf{f}}_{(x,y)} - \mathbf{f}_{(x,y)}^*\right\|_2 \quad \text{(End-Point Error, EPE)}
    \]
    dove $\hat{\mathbf{f}}$ \`e il flusso predetto e $\mathbf{f}^*$ \`e il ground truth.

    \item \textbf{Autoencoding multimodale}:
    \[
    \mathcal{L}_{\text{multi}} = \lambda_v \mathcal{L}_{\text{video}} + \lambda_a \mathcal{L}_{\text{audio}} + \lambda_l \mathcal{L}_{\text{label}}
    \]
    con $\mathcal{L}_{\text{video}}$ e $\mathcal{L}_{\text{audio}}$ basate su MSE (ricostruzione) e $\mathcal{L}_{\text{label}}$ basata su cross-entropy.
\end{itemize}
\end{formulabox}

\subsection{Backward pass dettagliato per il decoder del Perceiver IO}

Analizziamo nel dettaglio il flusso dei gradienti attraverso il decoder, che \`e il componente nuovo rispetto al Perceiver originale.

\subsubsection{Gradiente dal decoder alla proiezione di output}

Sia $\mathbf{y} = \text{MLP}(\mathbf{o}^{\text{dec}})$ l'output finale, dove $\mathbf{o}^{\text{dec}} \in \mathbb{R}^{O \times d_v}$ \`e l'output della cross-attention del decoder. Il gradiente della loss rispetto a $\mathbf{o}^{\text{dec}}$ si propaga attraverso l'MLP con le regole standard della backpropagation.

\subsubsection{Gradiente attraverso la cross-attention del decoder}

Sia $\mathbf{o}^{\text{dec}} = \mathbf{A}^{\text{dec}} \mathbf{V}^{\text{dec}}$ dove $\mathbf{A}^{\text{dec}} = \operatorname{softmax}(\mathbf{S}/\sqrt{d_k})$ e $\mathbf{S} = \mathbf{Q}^{\text{dec}} (\mathbf{K}^{\text{dec}})^\top$.

\begin{formulabox}[Gradienti del decoder cross-attention]
Dato $\frac{\partial \mathcal{L}}{\partial \mathbf{o}^{\text{dec}}} \in \mathbb{R}^{O \times d_v}$:

\textbf{Gradiente rispetto ai valori} (che fluisce verso i latenti):
\[
\frac{\partial \mathcal{L}}{\partial \mathbf{V}^{\text{dec}}} = (\mathbf{A}^{\text{dec}})^\top \frac{\partial \mathcal{L}}{\partial \mathbf{o}^{\text{dec}}} \in \mathbb{R}^{N \times d_v}
\]

\textbf{Gradiente rispetto alla matrice di attenzione}:
\[
\frac{\partial \mathcal{L}}{\partial \mathbf{A}^{\text{dec}}} = \frac{\partial \mathcal{L}}{\partial \mathbf{o}^{\text{dec}}} (\mathbf{V}^{\text{dec}})^\top \in \mathbb{R}^{O \times N}
\]

\textbf{Gradiente attraverso il softmax} (per la riga $i$):
\[
\frac{\partial \mathcal{L}}{\partial \mathbf{S}_i} = \mathbf{A}_i^{\text{dec}} \odot \left(\frac{\partial \mathcal{L}}{\partial \mathbf{A}_i^{\text{dec}}} - \left\langle \frac{\partial \mathcal{L}}{\partial \mathbf{A}_i^{\text{dec}}}, \mathbf{A}_i^{\text{dec}} \right\rangle \mathbf{1}\right)
\]

\textbf{Gradiente rispetto alle query} (che fluisce verso le output query):
\[
\frac{\partial \mathcal{L}}{\partial \mathbf{Q}^{\text{dec}}} = \frac{1}{\sqrt{d_k}} \frac{\partial \mathcal{L}}{\partial \mathbf{S}} \mathbf{K}^{\text{dec}} \in \mathbb{R}^{O \times d_k}
\]

\textbf{Gradiente rispetto alle chiavi} (che fluisce verso i latenti):
\[
\frac{\partial \mathcal{L}}{\partial \mathbf{K}^{\text{dec}}} = \frac{1}{\sqrt{d_k}} \left(\frac{\partial \mathcal{L}}{\partial \mathbf{S}}\right)^\top \mathbf{Q}^{\text{dec}} \in \mathbb{R}^{N \times d_k}
\]
\end{formulabox}

\begin{notabox}[Due flussi di gradiente dal decoder ai latenti]
Il decoder produce gradienti che fluiscono verso i latenti $\mathbf{z}^{(L)}$ attraverso \textbf{due percorsi}:
\begin{enumerate}
    \item Attraverso i \textbf{valori} $\mathbf{V}^{\text{dec}} = \mathbf{z}^{(L)} \mathbf{W}_V^{\text{dec}}$: il gradiente $\frac{\partial \mathcal{L}}{\partial \mathbf{V}^{\text{dec}}}$ si propaga ai latenti tramite $\frac{\partial \mathcal{L}}{\partial \mathbf{z}^{(L)}} \mathrel{+}= \frac{\partial \mathcal{L}}{\partial \mathbf{V}^{\text{dec}}} (\mathbf{W}_V^{\text{dec}})^\top$
    \item Attraverso le \textbf{chiavi} $\mathbf{K}^{\text{dec}} = \mathbf{z}^{(L)} \mathbf{W}_K^{\text{dec}}$: il gradiente $\frac{\partial \mathcal{L}}{\partial \mathbf{K}^{\text{dec}}}$ si propaga ai latenti tramite $\frac{\partial \mathcal{L}}{\partial \mathbf{z}^{(L)}} \mathrel{+}= \frac{\partial \mathcal{L}}{\partial \mathbf{K}^{\text{dec}}} (\mathbf{W}_K^{\text{dec}})^\top$
\end{enumerate}
Il gradiente combinato $\frac{\partial \mathcal{L}}{\partial \mathbf{z}^{(L)}}$ si propaga poi all'indietro attraverso il processor e l'encoder con le regole standard della backpropagation per self-attention e cross-attention.
\end{notabox}

\subsubsection{Gradiente per le diverse loss}

\begin{formulabox}[Gradiente della loss rispetto all'output del decoder]
\textbf{Cross-Entropy} (classificazione, MLM):
\[
\frac{\partial \mathcal{L}_{\text{CE}}}{\partial \hat{y}_k} = \hat{y}_k - y_k^* \quad \text{(con softmax output)}
\]

\textbf{MSE} (ricostruzione video/audio):
\[
\frac{\partial \mathcal{L}_{\text{MSE}}}{\partial \hat{y}_i} = \frac{2}{|\mathcal{S}|}(\hat{y}_i - y_i^*)
\]

\textbf{EPE} (optical flow):
\[
\frac{\partial \mathcal{L}_{\text{EPE}}}{\partial \hat{\mathbf{f}}_{(x,y)}} = \frac{\hat{\mathbf{f}}_{(x,y)} - \mathbf{f}_{(x,y)}^*}{\left\|\hat{\mathbf{f}}_{(x,y)} - \mathbf{f}_{(x,y)}^*\right\|_2 + \epsilon}
\]
dove $\epsilon$ \`e una piccola costante per la stabilit\`a numerica.
\end{formulabox}

\subsection{Flusso del gradiente completo}

Il backward pass attraversa i tre moduli in ordine inverso: decoder $\to$ processor $\to$ encoder. \`E importante notare che:

\begin{enumerate}
    \item I gradienti dal decoder fluiscono attraverso $\mathbf{z}^{(L)}$ ai parametri del processor e dell'encoder
    \item I gradienti dalle query di output influenzano i parametri del decoder e, indirettamente, l'intero modello
    \item Per l'autoencoding multimodale, i gradienti di ciascuna modalit\`a contribuiscono tutti all'aggiornamento dell'array latente condiviso
    \item Nel caso di weight sharing nel processor, i gradienti di tutti i blocchi che condividono gli stessi pesi vengono \textbf{sommati} prima dell'aggiornamento
\end{enumerate}

\begin{center}
\begin{tikzpicture}[
    block/.style={draw, rounded corners, minimum width=3cm, minimum height=1cm, align=center, font=\small},
    arr/.style={-{Stealth[length=2.5mm]}, thick},
    grad/.style={-{Stealth[length=2.5mm]}, thick, dashed, red!70!black}
]
    \node[block, fill=orange!10] (enc) at (0,0) {Encoder};
    \node[block, fill=blue!10] (proc) at (5,0) {Processor};
    \node[block, fill=green!10] (dec) at (10,0) {Decoder};
    \node[block, fill=red!10] (loss) at (10,-2) {$\mathcal{L}$};

    % Forward
    \draw[arr, blue!60!black] (enc) -- (proc) node[midway, above, font=\scriptsize] {$\mathbf{z}^{(0)}$};
    \draw[arr, blue!60!black] (proc) -- (dec) node[midway, above, font=\scriptsize] {$\mathbf{z}^{(L)}$};
    \draw[arr, blue!60!black] (dec) -- (loss);

    % Backward (gradients)
    \draw[grad] (loss) -- (dec) node[midway, right, font=\scriptsize, text=red!70!black] {$\frac{\partial \mathcal{L}}{\partial \theta_{\text{dec}}}$};
    \draw[grad] (dec) -- (proc) node[midway, below, font=\scriptsize, text=red!70!black] {$\frac{\partial \mathcal{L}}{\partial \mathbf{z}^{(L)}}$};
    \draw[grad] (proc) -- (enc) node[midway, below, font=\scriptsize, text=red!70!black] {$\frac{\partial \mathcal{L}}{\partial \mathbf{z}^{(0)}}$};

    \node[font=\scriptsize, blue!60!black] at (5,0.8) {Forward $\longrightarrow$};
    \node[font=\scriptsize, red!70!black] at (5,-0.8) {$\longleftarrow$ Backward};
\end{tikzpicture}
\end{center}


\subsection{Dettagli dell'ottimizzazione}

Il paper utilizza diverse configurazioni di training a seconda del task:

\begin{center}
\begin{tabular}{lcccc}
\toprule
\textbf{Parametro} & \textbf{Linguaggio} & \textbf{Optical Flow} & \textbf{Multimodale} & \textbf{ImageNet} \\
\midrule
Ottimizzatore & LAMB & Adam & LAMB & LAMB \\
Learning rate & $10^{-3}$ & $2 \times 10^{-4}$ & $10^{-3}$ & $2 \times 10^{-3}$ \\
Batch size & 64 & 8 & 32 & 1024 \\
LR schedule & Cosine decay & Cosine decay & Cosine decay & Flat + Cosine \\
Warmup steps & 10k & 5k & 10k & 10k \\
Dropout & 0.0 & 0.0 & 0.0 & 0.0 \\
Latenti $N$ & 256 & 2048 & 512 & 512 \\
Layer processor $L$ & 26 & 24 & 8 & 48 \\
\bottomrule
\end{tabular}
\end{center}

\subsubsection{Dettagli specifici per ImageNet}

Il training su ImageNet segue una procedura specifica (Jaegle et al., 2022):
\begin{itemize}
    \item \textbf{Ottimizzatore}: LAMB (You et al., 2020), progettato per batch grandi
    \item \textbf{Batch size}: 1024
    \item \textbf{Epoche}: 110 totali
    \item \textbf{Learning rate}: $2 \times 10^{-3}$, costante per le prime 55 epoche, poi cosine decay a zero
    \item \textbf{Weight decay}: 0.1
    \item \textbf{Gradient clipping}: norma massima 10
    \item \textbf{Dropout}: 0 (nessun dropout)
    \item \textbf{Struttura processor}: 8 blocchi di 6 moduli di attention ciascuno, con pesi condivisi tra i blocchi
    \item \textbf{Total depth}: $8 \times 6 = 48$ layer di self-attention effettivi
\end{itemize}

\subsubsection{Pre-training su JFT-300M}

Per le configurazioni con pre-training (Jaegle et al., 2022):
\begin{itemize}
    \item \textbf{Dataset}: JFT-300M (circa 300 milioni di immagini, 18{,}291 classi)
    \item \textbf{Batch size}: 8192
    \item \textbf{Hardware}: 256 TPU v3
    \item \textbf{Durata}: diversi giorni di training
    \item Dopo il pre-training su JFT, il modello viene \textbf{fine-tuned} su ImageNet-1K
    \item La testa di classificazione JFT (18k classi) viene sostituita con una per ImageNet (1000 classi)
\end{itemize}

\subsection{Subsampling delle output query durante il training}

Per task con output molto grandi (es.\ autoencoding multimodale con $O \approx 80{,}000$), calcolare il decoder su tutte le query ad ogni passo di training \`e proibitivo. La soluzione \`e il \textbf{subsampling}:

\begin{enumerate}
    \item Ad ogni passo di training, si campiona un sottoinsieme $\mathcal{S} \subset \{1, \ldots, O\}$ con $|\mathcal{S}| \ll O$
    \item Si calcolano le query di output solo per gli indici in $\mathcal{S}$
    \item La loss viene calcolata solo su questi punti: $\mathcal{L} = \frac{1}{|\mathcal{S}|}\sum_{i \in \mathcal{S}} \ell_i$
    \item Il gradiente \`e un \textbf{stimatore non distorto} del gradiente sulla loss completa
\end{enumerate}

Questa tecnica \`e analoga al \textbf{mini-batch} stochastic gradient descent, ma applicata alla dimensione dell'output anzich\'e alla dimensione dei dati.

\begin{notabox}[Perch\'e il subsampling funziona?]
Il gradiente rispetto ai parametri del decoder e del processor calcolato su un sottoinsieme di query \`e uno stimatore non distorto del gradiente completo, perch\'e ogni query contribuisce in modo indipendente alla loss (non c'\`e self-attention tra le query di output!). Questo \`e possibile proprio grazie al design del decoder con un singolo layer di cross-attention senza interazione tra query.

Formalmente, la loss completa \`e:
\[
\mathcal{L} = \frac{1}{O}\sum_{i=1}^{O} \ell_i(\mathbf{y}_i, \mathbf{y}_i^*)
\]
e lo stimatore mini-batch \`e:
\[
\hat{\mathcal{L}} = \frac{1}{|\mathcal{S}|}\sum_{i \in \mathcal{S}} \ell_i(\mathbf{y}_i, \mathbf{y}_i^*)
\]
con $\mathbb{E}_\mathcal{S}[\hat{\mathcal{L}}] = \mathcal{L}$ se $\mathcal{S}$ \`e campionato uniformemente.
\end{notabox}


%% ============================================================
%%  SEZIONE 7 --- ANALISI FORMALE
%% ============================================================
\section{Analisi formale dell'architettura}

\subsection{Il Perceiver IO come sistema di information retrieval}

L'architettura del Perceiver IO pu\`o essere interpretata come un sistema di \textbf{information retrieval} neurale a due fasi:

\begin{enumerate}
    \item \textbf{Indicizzazione} (Encoder + Processor): l'input viene ``indicizzato'' nell'array latente. I $N$ latenti agiscono come $N$ ``slot di memoria'' che catturano aspetti diversi dell'input. Il processor raffina questa indicizzazione attraverso self-attention.

    \item \textbf{Retrieval} (Decoder): le output query ``interrogano'' l'indice per recuperare le informazioni rilevanti. Ogni query genera un pattern di attenzione diverso sull'array latente, estraendo una combinazione pesata specifica per quel punto di output.
\end{enumerate}

Formalmente, sia $\alpha_{ij}^{\text{dec}}$ il peso di attenzione del decoder tra la query di output $i$ e il latente $j$:
\[
\alpha_{ij}^{\text{dec}} = \frac{\exp\!\left(\frac{\mathbf{q}_i \mathbf{W}_Q (\mathbf{z}_j^{(L)} \mathbf{W}_K)^\top}{\sqrt{d_k}}\right)}{\sum_{j'=1}^{N} \exp\!\left(\frac{\mathbf{q}_i \mathbf{W}_Q (\mathbf{z}_{j'}^{(L)} \mathbf{W}_K)^\top}{\sqrt{d_k}}\right)}
\]

Allora l'output per il punto $i$ \`e:
\[
\mathbf{y}_i = \sum_{j=1}^{N} \alpha_{ij}^{\text{dec}} \, \mathbf{z}_j^{(L)} \mathbf{W}_V^{\text{dec}} + \text{residual + MLP}
\]

Ogni punto di output $\mathbf{y}_i$ \`e una \textbf{combinazione convessa} dei valori dei latenti, con pesi determinati dalla compatibilit\`a tra la query e le chiavi latenti.

\subsection{Espressivit\`a del bottleneck latente}

Un aspetto critico \`e se $N$ latenti sono \emph{sufficienti} per catturare l'informazione necessaria. Possiamo analizzare questo formalmente.

\begin{proposition}[Capacit\`a informativa dei latenti]
L'array latente $\mathbf{z} \in \mathbb{R}^{N \times D}$ ha una capacit\`a informativa di $N \cdot D$ numeri reali. Se l'output richiede $O \cdot F$ numeri reali con $O \cdot F > N \cdot D$, allora il decoder \emph{non pu\`o} ricostruire perfettamente un output arbitrario.
\end{proposition}

Tuttavia, in pratica gli output non sono arbitrari: hanno struttura (correlazioni spaziali, temporali, semantiche) che pu\`o essere sfruttata. Il Perceiver IO funziona bene quando la ``dimensionalit\`a intrinseca'' dell'output \`e molto inferiore a $O \cdot F$.

\begin{examplebox}[Rapporto di compressione nell'autoencoding]
Nell'autoencoding di Kinetics-700:
\begin{itemize}
    \item Input/Output totale: $\sim 80{,}000$ elementi $\times$ feature per elemento
    \item Latenti: $N = 512$, $D = 1024$ $\to$ $512 \times 1024 = 524{,}288$ numeri
    \item Rapporto di compressione: $\approx 176 \times$
\end{itemize}
Il fatto che la ricostruzione sia di buona qualit\`a (PSNR $> 32$ dB) dimostra che i latenti catturano la struttura intrinseca dei dati multimodali.
\end{examplebox}

\subsection{Invarianza e equivarianza}

Il Perceiver IO ha propriet\`a interessanti rispetto a permutazioni dell'input e dell'output:

\begin{itemize}
    \item \textbf{Invarianza dell'encoder rispetto all'ordine dei latenti iniziali}: i latenti sono inizializzati con parametri appresi, quindi l'ordine \`e fisso. Tuttavia, qualsiasi permutazione dei latenti produce un modello funzionalmente equivalente (con parametri permutati).

    \item \textbf{Equivarianza dell'output rispetto alle query}: se permutiamo le output query, l'output viene permutato nello stesso modo. Formalmente:
    \[
    \text{Decoder}(\mathbf{P}\mathbf{q}, \mathbf{z}) = \mathbf{P} \cdot \text{Decoder}(\mathbf{q}, \mathbf{z})
    \]
    dove $\mathbf{P}$ \`e una matrice di permutazione. Questo perch\'e ogni query opera \textbf{indipendentemente}.

    \item \textbf{Sensibilit\`a all'ordine dell'input}: l'encoder \emph{non} \`e invariante all'ordine dell'input per default (l'input pu\`o includere encoding posizionali). Se si aggiungono encoding posizionali, l'ordine \`e codificato esplicitamente; senza di essi, l'encoder tratta l'input come un \emph{insieme} (set).
\end{itemize}

\subsection{Il Perceiver IO come generalizzazione di architetture note}

Il Perceiver IO pu\`o essere visto come una generalizzazione di diverse architetture esistenti:

\begin{proofbox}[Il Perceiver IO generalizza diverse architetture]
\textbf{Caso 1: $N = M$, nessun encoder.} Se il numero di latenti coincide con la dimensione dell'input e si elimina l'encoder (i latenti \emph{sono} l'input), il processor diventa un Transformer encoder standard. Se si aggiunge il decoder con query, si ottiene un Transformer encoder-decoder.

\textbf{Caso 2: $O = 1$, query appresa.} Se c'\`e una sola query di output, il decoder produce un singolo vettore. Se la query \`e costante, questo \`e equivalente a un average pooling appreso (attented pooling). Si ottiene un classificatore come il Perceiver originale.

\textbf{Caso 3: $O = M$, query = input features + PE.} Se le query di output sono costruite dalle stesse feature dell'input con encoding posizionale, si ottiene un modello che per ogni punto di input produce un output (come un U-Net o un modello di segmentazione, ma senza convoluzioni).

\textbf{Caso 4: $N \to \infty$.} Al crescere di $N$, il bottleneck diventa meno restrittivo e il modello converge (in principio) a un Transformer standard encoder-decoder, ma con complessit\`a crescente.
\end{proofbox}


%% ============================================================
%%  SEZIONE 7.5 --- PREPROCESSING CONVOLUTIVO
%% ============================================================
\subsection{Preprocessing convolutivo su ImageNet}

L'Appendice del paper (Jaegle et al., 2022) descrive un'opzione di \textbf{preprocessing convolutivo} per le immagini, ispirata dai ViT ibridi (Dosovitskiy et al., 2021):

\begin{itemize}
    \item Anzich\'e dare pixel grezzi + Fourier PE all'encoder, si preprocessa l'immagine con una piccola CNN
    \item La CNN estrae feature map a risoluzione ridotta (es.\ da $224 \times 224$ a $7 \times 7$)
    \item Queste feature vengono flattened e usate come input per l'encoder del Perceiver IO
    \item Questo introduce un \textbf{inductive bias spaziale} (invarianza a traslazione, localit\`a) che aiuta soprattutto con pochi dati
\end{itemize}

\begin{center}
\begin{tabular}{lccc}
\toprule
\textbf{Configurazione} & \textbf{Pre-training} & \textbf{Top-1 \%} & \textbf{Note} \\
\midrule
Pixel + Fourier PE & ImageNet & 79.0 & Baseline \\
Conv preprocessing & ImageNet & 77.4 & Meno dati, utile \\
Conv preprocessing & JFT-300M & 82.1 & -- \\
Conv preprocessing (grande) & JFT-300M & \textbf{86.4} & Vicino a SOTA \\
\bottomrule
\end{tabular}
\end{center}


%% ============================================================
%%  SEZIONE 8 --- DIAGRAMMA COMPLETO DELL'ARCHITETTURA
%% ============================================================
\section{Diagramma riassuntivo completo dell'architettura}

Flusso completo del Perceiver IO con dimensioni degli array, operazioni e complessit\`a.

\begin{center}
\begin{tikzpicture}[
    block/.style={draw, thick, rounded corners=4pt, minimum width=5cm, minimum height=1.5cm, align=center, font=\small},
    data/.style={draw, thick, minimum width=3.5cm, minimum height=0.8cm, align=center, font=\footnotesize, rounded corners=2pt},
    arr/.style={-{Stealth[length=3mm]}, very thick},
    dim/.style={font=\scriptsize\itshape, text=gray!60!black},
    compl/.style={font=\scriptsize\bfseries, text=red!60!black}
]
    % Input data
    \node[data, fill=red!10, draw=red!50!black] (x) at (0,10) {Input: $\mathbf{x} \in \mathbb{R}^{M \times C}$};
    \node[dim, right=0.5cm of x] {$M$ = num.\ elementi input};

    % Initial latent
    \node[data, fill=blue!10, draw=blue!50!black] (z0) at (0,8) {Latente iniziale: $\mathbf{z}^{(0)} \in \mathbb{R}^{N \times D}$};
    \node[dim, right=0.5cm of z0] {Parametri appresi};

    % Encoder
    \node[block, fill=orange!10, draw=orange!60!black] (enc) at (0,6) {
        \textbf{ENCODER}\\[2pt]
        Cross-Attention\\
        $\mathbf{Q} = \mathbf{z}^{(0)}\mathbf{W}_Q$, $\mathbf{K} = \mathbf{x}\mathbf{W}_K$, $\mathbf{V} = \mathbf{x}\mathbf{W}_V$
    };
    \node[compl, right=2.8cm of enc] {$\mathcal{O}(MND)$};

    % Encoded latent
    \node[data, fill=blue!15, draw=blue!60!black] (z_enc) at (0,3.8) {$\tilde{\mathbf{z}} \in \mathbb{R}^{N \times D}$};

    % Processor
    \node[block, fill=blue!10, draw=blue!70!black, minimum height=2.5cm] (proc) at (0,1.5) {
        \textbf{PROCESSOR}\\[2pt]
        $L$ layer di Self-Attention\\
        $\mathbf{z}^{(\ell)} = \text{SA}(\mathbf{z}^{(\ell-1)}) + \text{FFN}$\\
        Ogni layer: residual + LN
    };
    \node[compl, right=2.8cm of proc] {$\mathcal{O}(LN^2D)$};

    % Processed latent
    \node[data, fill=blue!25, draw=blue!70!black] (zL) at (0,-1) {$\mathbf{z}^{(L)} \in \mathbb{R}^{N \times D}$};

    % Output queries
    \node[data, fill=purple!10, draw=purple!50!black] (q) at (-5.5,-3) {Query: $\mathbf{q} \in \mathbb{R}^{O \times E}$};
    \node[dim, left=0.5cm of q, align=right] {Dipende\\dal task};

    % Decoder
    \node[block, fill=green!10, draw=green!60!black] (dec) at (0,-3) {
        \textbf{DECODER}\\[2pt]
        Cross-Attention\\
        $\mathbf{Q} = \mathbf{q}\mathbf{W}_Q$, $\mathbf{K} = \mathbf{z}^{(L)}\mathbf{W}_K$, $\mathbf{V} = \mathbf{z}^{(L)}\mathbf{W}_V$
    };
    \node[compl, right=2.8cm of dec] {$\mathcal{O}(OND)$};

    % Optional MLP
    \node[data, fill=gray!10, draw=gray!50!black] (mlp) at (0,-5) {MLP point-wise (opzionale)};

    % Output
    \node[data, fill=teal!15, draw=teal!60!black] (y) at (0,-6.5) {Output: $\mathbf{y} \in \mathbb{R}^{O \times F}$};
    \node[dim, right=0.5cm of y] {$O$ = num.\ output, $F$ = feature per output};

    % Arrows
    \draw[arr] (x) -- (enc);
    \draw[arr] (z0) -- (enc);
    \draw[arr] (enc) -- (z_enc);
    \draw[arr] (z_enc) -- (proc);
    \draw[arr] (proc) -- (zL);
    \draw[arr] (zL) -- (dec);
    \draw[arr] (q) -- (dec);
    \draw[arr] (dec) -- (mlp);
    \draw[arr] (mlp) -- (y);

    % Total complexity
    \node[draw, thick, rounded corners, fill=red!5, draw=red!50!black, minimum width=6cm, font=\small, align=center] at (0,-8) {
        Complessit\`a totale: $\mathcal{O}(MN + LN^2 + ON)$\\
        \textbf{Lineare} in $M$ e $O$
    };
\end{tikzpicture}
\end{center}


%% ============================================================
%%  SEZIONE 9 --- DOMANDE D'ESAME
%% ============================================================
\section{Domande d'esame tipiche}

In questa sezione raccogliamo le domande pi\`u probabili sull'architettura Perceiver IO, suddivise per difficolt\`a e tipologia.

\subsection{Domande concettuali}

\begin{enumerate}[label=\textbf{Q\arabic*.}]
    \item \textbf{Qual \`e il limite principale del Perceiver originale che Perceiver IO risolve?}\\
    \textit{Risposta attesa}: Il Perceiver originale pu\`o produrre solo output semplici (classificazione) tramite average pooling. Il Perceiver IO introduce un decoder basato su cross-attention con output query, permettendo output strutturati di dimensione arbitraria.

    \item \textbf{Descrivere il paradigma Read-Process-Write del Perceiver IO.}\\
    \textit{Risposta attesa}: Read = encoder (cross-attention dall'input ai latenti), Process = processor (self-attention sui latenti), Write = decoder (cross-attention dai latenti all'output tramite query).

    \item \textbf{Perch\'e il decoder del Perceiver IO usa un singolo layer di cross-attention senza self-attention tra le query?}\\
    \textit{Risposta attesa}: (1) Efficienza: self-attention tra $O$ query avrebbe costo $\mathcal{O}(O^2)$, vanificando il vantaggio; (2) Sufficienza: l'informazione \`e gi\`a stata elaborata nel processor, il decoder deve solo ``leggere''.

    \item \textbf{Cosa sono le output query e come vengono costruite per diversi task?}\\
    \textit{Risposta attesa}: Le output query codificano ``cosa il modello deve produrre'' per ogni punto di output. Per la classificazione: singola query appresa. Per l'optical flow: position encoding + input features. Per il multi-task: task embedding + position embedding. Per il multimodale: position encoding + modality flag.

    \item \textbf{Come gestisce il Perceiver IO il linguaggio senza tokenizzazione?}\\
    \textit{Risposta attesa}: Opera direttamente su byte UTF-8 ($M = 2048$ byte), con embedding appresi per-byte. La complessit\`a lineare in $M$ rende gestibili le sequenze pi\`u lunghe. Supera BERT su GLUE (81.8 vs 81.0) senza WordPiece.

    \item \textbf{Spiegare come viene costruita la query per l'optical flow e perch\'e include feature dell'input.}\\
    \textit{Risposta attesa}: La query per ogni pixel $(x,y)$ \`e $\mathbf{q}_{(x,y)} = [\text{FourierPE}(x,y) \| \mathbf{f}_{(x,y)}]$, dove $\mathbf{f}_{(x,y)}$ sono le feature della patch $3 \times 3$ del primo frame. Le feature dell'input sono incluse perch\'e forniscono al decoder informazione sul ``contenuto'' del pixel (non solo la posizione), permettendo una ricerca pi\`u mirata nelle corrispondenze tra i frame. Senza queste feature, il decoder conoscerebbe solo la posizione ma non l'aspetto visivo del pixel, rendendo pi\`u difficile trovare la corrispondenza.

    \item \textbf{Perch\'e nell'autoencoding multimodale la label viene mascherata il 50\% delle volte?}\\
    \textit{Risposta attesa}: Per evitare il ``shortcut'' dove il modello copia la label dall'input all'output senza elaborare video e audio. Il mascheramento forza i latenti a catturare informazione semantica dalle modalit\`a percettive, non solo a memorizzare la label.
\end{enumerate}

\subsection{Domande tecniche e di complessit\`a}

\begin{enumerate}[label=\textbf{Q\arabic*.}, resume]
    \item \textbf{Derivare la complessit\`a computazionale totale del Perceiver IO e confrontarla con un Transformer standard.}\\
    \textit{Risposta attesa}: Encoder $\mathcal{O}(MN)$, Processor $\mathcal{O}(LN^2)$, Decoder $\mathcal{O}(ON)$. Totale: $\mathcal{O}(MN + LN^2 + ON)$. Transformer: $\mathcal{O}(L(M+O)^2)$. Il Perceiver IO \`e lineare in $M$ e $O$, il Transformer quadratico.

    \item \textbf{Con $M = 50{,}000$, $O = 50{,}000$, $N = 512$, $L = 6$, calcolare il rapporto tra la complessit\`a del Perceiver IO e quella del Transformer.}\\
    \textit{Risposta attesa}: Perceiver IO: $50{,}000 \times 512 + 6 \times 512^2 + 50{,}000 \times 512 \approx 52.8\text{M}$. Transformer: $6 \times (100{,}000)^2 = 60{,}000\text{M}$. Rapporto: $\approx 1{,}136\times$ pi\`u efficiente.

    \item \textbf{Scrivere le equazioni complete del modulo decoder, specificando le dimensioni di tutte le matrici.}\\
    \textit{Risposta attesa}: $\mathbf{Q} = \mathbf{q}\mathbf{W}_Q^{\text{dec}} \in \mathbb{R}^{O \times d_k}$, $\mathbf{K} = \mathbf{z}^{(L)}\mathbf{W}_K^{\text{dec}} \in \mathbb{R}^{N \times d_k}$, $\mathbf{V} = \mathbf{z}^{(L)}\mathbf{W}_V^{\text{dec}} \in \mathbb{R}^{N \times d_v}$, Attention $= \operatorname{softmax}(\mathbf{QK}^\top/\sqrt{d_k})\mathbf{V} \in \mathbb{R}^{O \times d_v}$. Matrice di attenzione: $O \times N$.

    \item \textbf{Perch\'e il subsampling delle output query durante il training produce un gradiente non distorto?}\\
    \textit{Risposta attesa}: Perch\'e le query operano indipendentemente (non c'\`e self-attention tra di esse nel decoder). Ogni query contribuisce additivamente alla loss. Il campionamento di un sottoinsieme produce uno stimatore non distorto del gradiente totale, analogo al mini-batch SGD.

    \item \textbf{Qual \`e il ruolo del processor e perch\'e la sua complessit\`a \`e indipendente da $M$ e $O$?}\\
    \textit{Risposta attesa}: Il processor applica $L$ layer di self-attention \emph{solo} sui $N$ latenti. Poich\'e $N$ \`e un iperparametro fisso (non dipende dai dati), la complessit\`a $\mathcal{O}(LN^2D)$ \`e costante rispetto a $M$ e $O$. Questo disaccoppia il ``ragionamento'' dalla dimensionalit\`a dell'input/output.

    \item \textbf{Calcolare il numero di parametri del decoder per l'optical flow, con $E = 322$, $D = 512$, $d_k = d_v = 64$, $h = 8$ teste.}\\
    \textit{Risposta attesa}: $\mathbf{W}_Q: 322 \times 64 \times 8 = 164{,}864$. $\mathbf{W}_K: 512 \times 64 \times 8 = 262{,}144$. $\mathbf{W}_V: 512 \times 64 \times 8 = 262{,}144$. $\mathbf{W}_O: 512 \times 512 = 262{,}144$. Totale cross-attention: $\approx 951$K. Pi\`u i parametri dell'MLP di proiezione all'output (da 512 a 2).
\end{enumerate}

\subsection{Domande di confronto e riflessione}

\begin{enumerate}[label=\textbf{Q\arabic*.}, resume]
    \item \textbf{Confrontare il meccanismo di output del Perceiver IO con il token [CLS] di BERT e con il decoder autoregressivo del Transformer.}\\
    \textit{Risposta attesa}: [CLS] di BERT: singolo vettore, solo classificazione. Decoder Transformer: autoregressivo, sequenziale, specifico per sequenze. Perceiver IO: query-based, parallelo, flessibile per output di qualsiasi tipo e dimensione.

    \item \textbf{In che senso il Perceiver IO realizza un autoencoder multimodale? Quali sono le differenze rispetto a un VAE?}\\
    \textit{Risposta attesa}: Il Perceiver IO comprime input multimodali in latenti e ricostruisce tutte le modalit\`a. A differenza del VAE: (1) i latenti non sono vincolati a una distribuzione nota; (2) il bottleneck \`e strutturato ($N \times D$, non un singolo vettore); (3) il decoder \`e basato su query, non su deconvoluzione/RNN.

    \item \textbf{Perch\'e il multi-task learning nel Perceiver IO \`e pi\`u naturale rispetto ai modelli tradizionali?}\\
    \textit{Risposta attesa}: Nei modelli tradizionali servono teste task-specifiche (MLP diversi). Nel Perceiver IO, la specializzazione per task avviene interamente nelle output query: basta usare un diverso task embedding nella query. Encoder e processor sono completamente condivisi, e il decoder \`e generico.

    \item \textbf{Il Perceiver IO pu\`o sostituire un Vision Transformer (ViT)? Quali sono i trade-off?}\\
    \textit{Risposta attesa}: S\`i, per la classificazione di immagini raggiunge performance comparabili (79--84.5\% su ImageNet). Trade-off: il ViT sfrutta l'induttivo bias della struttura a patch (analogo alle CNN), mentre il Perceiver IO \`e pi\`u generico ma potrebbe richiedere pi\`u dati. Il vantaggio del Perceiver IO \`e la flessibilit\`a su task diversi dalla classificazione.

    \item \textbf{Come si relaziona la complessit\`a del Perceiver IO con gli approcci di linear attention (es.\ Performer)?}\\
    \textit{Risposta attesa}: Entrambi ottengono complessit\`a lineare nella lunghezza della sequenza, ma con meccanismi diversi. Il Performer approssima la matrice di attenzione con kernel random features. Il Perceiver IO usa un bottleneck latente \emph{esatto}: non approssima l'attenzione, ma la calcola esattamente su un set ridotto di latenti.

    \item \textbf{Quali sono le limitazioni del Perceiver IO? In quali scenari potrebbe non funzionare bene?}\\
    \textit{Risposta attesa}: (1) Se l'output richiede forte interazione tra i punti di output (es.\ generazione autoregressiva con dipendenze sequenziali), il decoder a singolo layer senza self-attention tra query potrebbe essere limitante. (2) Il numero di latenti $N$ deve essere scelto con cura: troppo piccolo perde informazione, troppo grande vanifica il vantaggio computazionale. (3) Per task molto specifici, architetture specializzate possono ancora superare il Perceiver IO grazie a bias induttivi appropriati.

    \item \textbf{Spiegare come il Perceiver IO gestisce il task di optical flow senza le componenti tradizionali (cost volume, warping, piramide multi-scala).}\\
    \textit{Risposta attesa}: L'input sono le patch $3 \times 3$ di entrambi i frame concatenate. L'encoder proietta tutto nei latenti. Il processor elabora le relazioni globali (corrispondenze tra frame). Il decoder, con query che includono la posizione del pixel e le feature locali, estrae lo spostamento per ogni pixel. Il meccanismo di attenzione nel processor impara implicitamente le corrispondenze che i cost volume calcolano esplicitamente.

    \item \textbf{In che modo l'approccio byte-level del Perceiver IO per il linguaggio si confronta con i modelli sub-word (BPE, WordPiece)?}\\
    \textit{Risposta attesa}: I modelli sub-word comprimono il testo in meno token (es.\ 512 token WordPiece $\approx$ 2048 byte), riducendo la lunghezza della sequenza ma introducendo un vocabolario artificiale. Il Perceiver IO opera su byte ($4\times$ pi\`u lunghi) ma gestisce la lunghezza grazie alla complessit\`a lineare. Vantaggi: universalit\`a linguistica, robustezza a errori, nessun preprocessing. Il risultato su GLUE (81.8 vs 81.0) suggerisce che la tokenizzazione non \`e un prerequisito per buone performance.

    \item \textbf{Confrontare il Perceiver IO con un Transformer standard per StarCraft II. Perch\'e il Perceiver IO \`e pi\`u efficiente mantenendo la stessa qualit\`a?}\\
    \textit{Risposta attesa}: Il Transformer di AlphaStar applica self-attention su tutte le 512 entit\`a, con costo $\mathcal{O}(512^2) = \mathcal{O}(262{,}144)$ per layer. Il Perceiver IO usa un bottleneck latente con $N$ latenti (es.\ $N = 64$), quindi il costo \`e $\mathcal{O}(512 \times 64 + L \times 64^2 + 512 \times 64)$, molto inferiore. La qualit\`a si mantiene perch\'e 64 latenti sono sufficienti a catturare le relazioni strategiche tra le entit\`a.

    \item \textbf{Perch\'e il Fourier positional encoding \`e cos\`i importante per ImageNet e come si confronta con il learned positional encoding?}\\
    \textit{Risposta attesa}: Il Fourier PE codifica esplicitamente la struttura 2D delle immagini tramite seni e coseni a diverse frequenze, fornendo un inductive bias spaziale. Il learned PE deve apprendere questa struttura dai dati. Su ImageNet con solo 1.3M di immagini, il Fourier PE (79.0\%) supera di gran lunga il learned PE (67.6\%). Con pi\`u dati (JFT-300M), il gap si riduce ma non scompare (84.5\% vs 72.7\%).
\end{enumerate}


%% ============================================================
%%  SEZIONE 10 --- ABLATION STUDIES
%% ============================================================
\section{Ablation Studies dal paper originale}

Questa sezione analizza in dettaglio gli \textbf{ablation studies} riportati nel paper di Jaegle et al.\ (2022). Gli ablation studies sono esperimenti in cui si ``rimuove'' o si ``varia'' un componente alla volta per capire il contributo di ciascun elemento. Sono fondamentali per comprendere \emph{perch\'e} il Perceiver IO funziona.

\subsection{Table 7 --- Risultati completi su ImageNet}

La Table 7 del paper \`e la tabella pi\`u completa per la classificazione su ImageNet. Riportiamo qui tutti i risultati con un commento dettagliato per ciascuna riga.

\begin{center}
\footnotesize
\begin{tabular}{p{4.5cm}cccp{4cm}}
\toprule
\textbf{Modello} & \textbf{Pre-train} & \textbf{Top-1 \%} & \textbf{Param (M)} & \textbf{Commento} \\
\midrule
\multicolumn{5}{l}{\textit{Baseline ConvNet}} \\
ResNet-50 & ImageNet & 78.6 & 26 & CNN classica, riferimento storico \\
NFNet-F6+SAM & ImageNet & 86.5 & 438 & SOTA tra le CNN, molto grande \\
\midrule
\multicolumn{5}{l}{\textit{Baseline ViT}} \\
ViT-B/16 & ImageNet & 77.9 & 86 & ViT senza pre-training extra \\
ViT-H/14 & JFT-300M & 88.6 & 632 & SOTA assoluto con enorme pre-training \\
\midrule
\multicolumn{5}{l}{\textit{Perceiver originale (Jaegle et al., 2021)}} \\
Perceiver, 2D Fourier PE & ImageNet & 78.6 & 42.1 & Uguale a ResNet-50! \\
Perceiver, Learned PE & ImageNet & 67.6 & 42.1 & $-11\%$ senza Fourier \\
Perceiver, Conv+MaxPool & ImageNet & 77.4 & -- & Con preprocessing convolutivo \\
\midrule
\multicolumn{5}{l}{\textit{Perceiver IO}} \\
Config A, 2D Fourier PE & ImageNet & \textbf{79.0} & 48.4 & \textbf{+0.4\% vs Perceiver!} \\
Config A, Learned PE & ImageNet & 72.7 & 48.4 & +5.1\% vs Perceiver learned \\
Config A, Conv+MaxPool & ImageNet & 82.1 & -- & Grande miglioramento \\
Config B, 2D Fourier & JFT-300M & 84.5 & 212 & Con pre-training JFT \\
Config B, Conv+MaxPool & JFT-300M & \textbf{86.4} & -- & Quasi NFNet! \\
\bottomrule
\end{tabular}
\end{center}

\begin{notabox}[Analisi dettagliata dei risultati ImageNet --- Cosa ci dicono?]
Questi risultati rivelano diversi insight fondamentali:

\textbf{1. Il decoder migliora SEMPRE le performance.}
\begin{itemize}
    \item Perceiver (avg pool) con 2D Fourier: 78.6\% $\to$ Perceiver IO (cross-att decoder): \textbf{79.0\%}
    \item Perceiver (avg pool) con Learned PE: 67.6\% $\to$ Perceiver IO: \textbf{72.7\%} (+5.1\%)
    \item Perceiver (avg pool) con Conv+MaxPool: 77.4\% $\to$ Perceiver IO: \textbf{82.1\%} (+4.7\%)
\end{itemize}
In \emph{ogni} configurazione, il decoder con cross-attention supera l'average pooling. Questo dimostra che il meccanismo di query \`e superiore anche per la classificazione (dove $O = 1$).

\textbf{2. Il Fourier PE \`e cruciale con pochi dati.}
\begin{itemize}
    \item Con Fourier PE: 79.0\%, senza: 72.7\% ($-6.3\%$)
    \item Il Fourier PE ``regala'' al modello la conoscenza della struttura 2D dell'immagine
    \item Senza di esso, il modello deve impararla dai dati, e 1.3M immagini non bastano
\end{itemize}

\textbf{3. Pi\`u dati compensano la mancanza di inductive bias.}
\begin{itemize}
    \item Con solo ImageNet e Learned PE: 72.7\%
    \item Ma con JFT-300M e Learned PE: miglioramento significativo
    \item Con JFT-300M e 2D Fourier: 84.5\% (eccellente)
\end{itemize}

\textbf{4. Il preprocessing convolutivo \`e il ``boost'' pi\`u grande.}
\begin{itemize}
    \item Config A con Conv+MaxPool: 82.1\% vs 79.0\% senza Conv (+3.1\%)
    \item Config B con Conv+MaxPool e JFT: 86.4\% (a soli 0.1\% da NFNet!)
    \item La convoluzione aggiunge localit\`a e invarianza a traslazione, inductive bias potentissimi per le immagini
\end{itemize}
\end{notabox}

\begin{notabox}[In parole semplici: perch\'e il decoder migliora anche la classificazione?]
Sembra controintuitivo: per la classificazione serve un \textbf{singolo} vettore di output, quindi perch\'e il decoder con cross-attention (che \`e progettato per output multipli) \`e meglio dell'average pooling?

La risposta \`e che l'average pooling tratta tutti i latenti \textbf{allo stesso modo} --- fa una media democratica. Ma non tutti i latenti sono ugualmente importanti per classificare! Alcuni latenti potrebbero contenere informazioni sullo sfondo (irrilevanti), altri sul soggetto principale (cruciali).

Il decoder con cross-attention impara una \textbf{query intelligente} che ``sa'' a quali latenti dare pi\`u peso. \`E come la differenza tra chiedere ``Dimmi tutto quello che sai'' (average pooling) e ``Dimmi solo le cose importanti per capire se questa \`e una foto di un gatto'' (cross-attention con query appresa).
\end{notabox}

\subsection{Table 8 --- Velocit\`a di training}

Un risultato sorprendente \`e che il Perceiver IO non \`e pi\`u lento del Perceiver originale, nonostante abbia un decoder in pi\`u.

\begin{center}
\begin{tabular}{lcc}
\toprule
\textbf{Modello} & \textbf{Steps/sec} & \textbf{Note} \\
\midrule
Perceiver, 2D Fourier PE & 4.73 & Baseline \\
Perceiver IO, 2D Fourier PE & \textbf{4.85} & Pi\`u veloce del Perceiver! \\
Perceiver IO, pretrained (config B) & 6.41 & Architettura pi\`u leggera \\
\bottomrule
\end{tabular}
\end{center}

\begin{notabox}[Perch\'e il Perceiver IO \`e pi\`u veloce del Perceiver nonostante il decoder?]
Questo sembra un paradosso: il Perceiver IO ha un componente in pi\`u (il decoder), eppure \`e leggermente pi\`u veloce. La spiegazione \`e che:
\begin{enumerate}
    \item Il decoder per la classificazione ha $O = 1$ (una sola query), quindi il costo \`e $\mathcal{O}(1 \cdot N \cdot D)$ --- trascurabile
    \item Il Perceiver originale usa weight sharing con encoder ripetuto (cross-attention dall'input ai latenti per $R$ iterazioni), che costa $\mathcal{O}(R \cdot M \cdot N)$
    \item Il Perceiver IO pu\`o evitare le iterazioni dell'encoder, usando un singolo passo di cross-attention + un processor pi\`u profondo
    \item La configurazione pre-trained usa un'architettura pi\`u leggera (meno layer per blocco, meno ripetizioni), risultando ancora pi\`u veloce (6.41 steps/sec)
\end{enumerate}
\end{notabox}

\subsection{Table 9 --- StarCraft II: efficienza computazionale}

\begin{center}
\begin{tabular}{lccc}
\toprule
\textbf{Modello} & \textbf{Win Rate} & \textbf{FLOPs} & \textbf{Parametri} \\
\midrule
Transformer (entity encoder) & 87\% & 3.3B & 144M \\
\textbf{Perceiver IO} & \textbf{87\%} & \textbf{0.93B} & \textbf{140M} \\
\midrule
\multicolumn{4}{l}{Riduzione FLOPs: $\mathbf{3.5\times}$} \\
\bottomrule
\end{tabular}
\end{center}

\begin{examplebox}[Analisi dettagliata: perch\'e $3.5\times$ meno FLOPs?]
Calcoliamo la complessit\`a in dettaglio.

\textbf{Transformer originale di AlphaStar}:
\begin{itemize}
    \item Self-attention su $E = 512$ entit\`a, per $L$ layer
    \item Costo per layer: $\mathcal{O}(E^2 \cdot d) = \mathcal{O}(512^2 \cdot d) = \mathcal{O}(262{,}144 \cdot d)$
    \item Con $L$ layer e FFN: $\sim 3.3$B FLOPs totali
\end{itemize}

\textbf{Perceiver IO}:
\begin{itemize}
    \item Encoder: $\mathcal{O}(E \cdot N \cdot D) = \mathcal{O}(512 \cdot 256 \cdot 512) \approx 67$M
    \item Processor: $\mathcal{O}(L \cdot N^2 \cdot D) = \mathcal{O}(8 \cdot 256^2 \cdot 512) \approx 268$M
    \item Decoder: $\mathcal{O}(E \cdot N \cdot D) = \mathcal{O}(512 \cdot 256 \cdot 512) \approx 67$M
    \item Totale: $\approx 402$M operazioni (solo attenzione) + FFN $\approx 0.93$B
\end{itemize}

Il risparmio viene dal processor: self-attention su $N = 256$ latenti invece di $E = 512$ entit\`a, quindi $256^2 = 65{,}536$ vs $512^2 = 262{,}144$, un fattore $4\times$ solo per la matrice di attenzione. Il bottleneck latente \`e la chiave dell'efficienza.
\end{examplebox}

\subsection{Table 2 --- Ablation multi-task su GLUE}

Questo \`e uno degli ablation pi\`u interessanti del paper. Confronta diverse strategie per gestire il multi-task learning.

\begin{center}
\begin{tabular}{p{5cm}cp{5.5cm}}
\toprule
\textbf{Configurazione} & \textbf{GLUE Avg} & \textbf{Descrizione} \\
\midrule
Single-task query & 81.0 & Modelli separati, una query per task \\
Shared input token & 81.5 & Input condiviso + token speciale per task \\
Task-specific input tokens & 81.8 & Token di input diverso per ogni task \\
Multitask query & 81.8 & Query di output diverse per task \\
\bottomrule
\end{tabular}
\end{center}

\begin{notabox}[Analisi dell'ablation multi-task]
Cosa significano queste righe?

\textbf{Single-task query (81.0)}: 8 modelli separati, ciascuno addestrato su un solo task GLUE. Questo \`e il baseline ``standard''.

\textbf{Shared input token (81.5)}: Un singolo modello per tutti i task. Per indicare quale task eseguire, si aggiunge un token speciale condiviso all'input (simile al token [CLS] di BERT). Risultato: +0.5 rispetto a single-task --- il multi-task gi\`a aiuta!

\textbf{Task-specific input tokens (81.8)}: Come sopra, ma ogni task ha il \textbf{proprio} token di input (un token diverso per CoLA, SST-2, MNLI, etc.). Risultato: +0.8 vs single-task --- differenziare i task all'input aiuta.

\textbf{Multitask query (81.8)}: L'approccio del Perceiver IO: l'input \`e identico per tutti i task, ma le \textbf{output query} contengono un task embedding diverso. Risultato: \textbf{identico} a task-specific input tokens!

\textbf{Insight chiave}: l'approccio basato su query di output \`e efficace quanto avere token di input specifici per task. Ma \`e pi\`u elegante: non richiede di modificare l'input, e la specializzazione avviene interamente nel decoder. Questo suggerisce che le query di output sono un meccanismo sufficiente per differenziare i task, senza bisogno di token [CLS] o simili.
\end{notabox}

\subsection{Ablation sull'Optical Flow}

Il paper presenta anche un ablation sull'optical flow (nell'appendice), confrontando diverse configurazioni di pre/post-processing.

\begin{center}
\begin{tabular}{p{5cm}ccc}
\toprule
\textbf{Configurazione} & \textbf{Sintel Clean} & \textbf{Sintel Final} & \textbf{Note} \\
 & EPE $\downarrow$ & EPE $\downarrow$ & \\
\midrule
Perceiver IO (no pre/post) & 1.81 & 2.42 & Risultato principale \\
+ Conv + MaxPool preprocessing & $\sim$1.85 & $\sim$2.50 & Simile, non migliora \\
+ RAFT-style upsampling & $\sim$1.90 & $\sim$2.55 & Leggermente peggio \\
Senza input features in query & $\sim$2.10 & $\sim$2.80 & Peggiore senza feature \\
\bottomrule
\end{tabular}
\end{center}

\begin{notabox}[Insight chiave sull'ablation dell'optical flow]
Il risultato pi\`u importante \`e che il Perceiver IO raggiunge performance competitive o migliori \textbf{senza alcuna componente specializzata}:
\begin{enumerate}
    \item \textbf{Nessun preprocessing convolutivo} necessario: la cross-attention sull'input grezzo funziona gi\`a benissimo
    \item \textbf{Nessun upsampling specializzato} (come quello di RAFT): il decoder produce direttamente l'output alla risoluzione desiderata
    \item Le \textbf{feature dell'input nelle query} sono importanti: rimuoverle peggiora di $\sim$0.3--0.4 EPE, perch\'e il decoder perde informazione locale sul pixel
    \item Questo dimostra che un'architettura \textbf{general-purpose} pu\`o sostituire ingegneria specifica del dominio
\end{enumerate}
\end{notabox}


%% ============================================================
%%  SEZIONE 11 --- APPROFONDIMENTO LOSS FUNCTIONS
%% ============================================================
\section{Approfondimento: le funzioni di Loss} \label{sec:cross_entropy_loss}

Per ogni dominio sperimentale, il Perceiver IO utilizza una funzione di loss diversa. In questa sezione approfondiamo ciascuna con spiegazioni dettagliate, formule complete, e \textbf{esempi numerici passo-passo}.

\subsection{Cross-Entropy Loss (per classificazione)}

La cross-entropy \`e la loss standard per i problemi di classificazione. Misura la ``distanza'' tra la distribuzione predetta e la distribuzione vera.

\begin{formulabox}[Cross-Entropy Loss --- Formula completa]
Data una predizione $\hat{\mathbf{y}} = [\hat{y}_1, \ldots, \hat{y}_K]$ dopo il softmax e il vettore target one-hot $\mathbf{y}^* = [y_1^*, \ldots, y_K^*]$ (dove $y_c^* = 1$ per la classe corretta $c$ e $y_k^* = 0$ altrimenti):

\[
\mathcal{L}_{\text{CE}} = -\sum_{k=1}^{K} y_k^* \log(\hat{y}_k) = -\log(\hat{y}_c)
\]

dove $c$ \`e l'indice della classe corretta. La seconda uguaglianza vale perch\'e $y_k^* = 0$ per $k \neq c$.

\textbf{Con softmax}: se i logit (output prima del softmax) sono $\mathbf{z} = [z_1, \ldots, z_K]$:
\[
\hat{y}_k = \frac{e^{z_k}}{\sum_{j=1}^K e^{z_j}} \quad \Rightarrow \quad \mathcal{L}_{\text{CE}} = -z_c + \log\left(\sum_{j=1}^K e^{z_j}\right)
\]

\textbf{Usata per}: ImageNet, GLUE (classificazione), StarCraft II, head di classificazione nel multimodale, MLM (per-byte).
\end{formulabox}

\begin{examplebox}[Esempio numerico: Cross-Entropy passo-passo]
Supponiamo un problema di classificazione con $K = 3$ classi (gatto, cane, uccello), e la classe corretta \`e ``gatto'' (classe 0).

\textbf{Passo 1}: Il modello produce i logit:
\[
\mathbf{z} = [2.0, \; 1.0, \; 0.1]
\]

\textbf{Passo 2}: Applico il softmax:
\begin{align*}
\hat{y}_0 &= \frac{e^{2.0}}{e^{2.0} + e^{1.0} + e^{0.1}} = \frac{7.389}{7.389 + 2.718 + 1.105} = \frac{7.389}{11.212} = 0.659 \\
\hat{y}_1 &= \frac{e^{1.0}}{11.212} = \frac{2.718}{11.212} = 0.242 \\
\hat{y}_2 &= \frac{e^{0.1}}{11.212} = \frac{1.105}{11.212} = 0.099
\end{align*}

\textbf{Passo 3}: Calcolo la cross-entropy (la classe corretta \`e 0):
\[
\mathcal{L}_{\text{CE}} = -\log(\hat{y}_0) = -\log(0.659) = 0.417
\]

\textbf{Interpretazione}: Se il modello fosse perfetto ($\hat{y}_0 = 1.0$), la loss sarebbe $-\log(1) = 0$. Se fosse completamente sbagliato ($\hat{y}_0 \to 0$), la loss tenderebbe a $+\infty$. Il valore $0.417$ indica una predizione ``ragionevole ma migliorabile''.

\textbf{Passo 4}: Gradiente rispetto ai logit (fondamentale per il backward pass!):
\[
\frac{\partial \mathcal{L}}{\partial z_k} = \hat{y}_k - y_k^* = \begin{cases} 0.659 - 1 = -0.341 & \text{per } k = 0 \text{ (classe corretta)} \\ 0.242 - 0 = 0.242 & \text{per } k = 1 \\ 0.099 - 0 = 0.099 & \text{per } k = 2 \end{cases}
\]
Il gradiente ``spinge'' il logit della classe corretta verso l'alto ($-0.341$, quindi aggiornamento positivo) e i logit delle classi sbagliate verso il basso.
\end{examplebox}

\subsection{MSE Loss (per ricostruzione)}

La Mean Squared Error \`e la loss standard per problemi di regressione/ricostruzione. Misura la distanza media al quadrato tra predizione e target.

\begin{formulabox}[MSE Loss --- Formula completa]
Dato un vettore predetto $\hat{\mathbf{y}} \in \mathbb{R}^F$ e il target $\mathbf{y}^* \in \mathbb{R}^F$:

\[
\mathcal{L}_{\text{MSE}} = \frac{1}{N_{\text{tot}}} \sum_{i=1}^{N_{\text{punti}}} \sum_{f=1}^{F} \left(\hat{y}_{i,f} - y_{i,f}^*\right)^2
\]

dove $N_{\text{tot}} = N_{\text{punti}} \times F$ \`e il numero totale di valori.

\textbf{Propriet\`a chiave}:
\begin{itemize}
    \item Sempre $\geq 0$, uguale a $0$ solo con ricostruzione perfetta
    \item Penalizza di pi\`u gli errori grandi (a causa del quadrato)
    \item Il gradiente \`e proporzionale all'errore: $\frac{\partial \mathcal{L}}{\partial \hat{y}_{i,f}} = \frac{2}{N_{\text{tot}}}(\hat{y}_{i,f} - y_{i,f}^*)$
\end{itemize}

\textbf{Usata per}: ricostruzione video (pixel RGB) e audio (campioni) nell'autoencoding multimodale.
\end{formulabox}

\begin{examplebox}[Esempio numerico: MSE per ricostruzione video]
Supponiamo di voler ricostruire un singolo pixel RGB.

\textbf{Pixel predetto}: $\hat{\mathbf{y}} = [0.8, \; 0.3, \; 0.5]$ (canali R, G, B normalizzati in $[0, 1]$)

\textbf{Pixel target}: $\mathbf{y}^* = [1.0, \; 0.0, \; 0.5]$

\textbf{Calcolo MSE}:
\begin{align*}
\mathcal{L}_{\text{MSE}} &= \frac{1}{3}\left[(0.8 - 1.0)^2 + (0.3 - 0.0)^2 + (0.5 - 0.5)^2\right] \\
&= \frac{1}{3}\left[(-0.2)^2 + (0.3)^2 + (0.0)^2\right] \\
&= \frac{1}{3}\left[0.04 + 0.09 + 0.00\right] \\
&= \frac{0.13}{3} = 0.0433
\end{align*}

\textbf{Gradiente per ogni canale}:
\begin{align*}
\frac{\partial \mathcal{L}}{\partial \hat{y}_R} &= \frac{2}{3}(0.8 - 1.0) = -0.133 \quad \text{(spinge R verso l'alto)} \\
\frac{\partial \mathcal{L}}{\partial \hat{y}_G} &= \frac{2}{3}(0.3 - 0.0) = +0.200 \quad \text{(spinge G verso il basso)} \\
\frac{\partial \mathcal{L}}{\partial \hat{y}_B} &= \frac{2}{3}(0.5 - 0.5) = 0.000 \quad \text{(B \`e gi\`a perfetto)}
\end{align*}

\textbf{Conversione in PSNR} (Peak Signal-to-Noise Ratio), la metrica usata nel paper:
\[
\text{PSNR} = 10 \cdot \log_{10}\left(\frac{\text{MAX}^2}{\text{MSE}}\right) = 10 \cdot \log_{10}\left(\frac{1^2}{0.0433}\right) = 10 \cdot \log_{10}(23.1) = 10 \cdot 1.363 = 13.6 \text{ dB}
\]
(Un PSNR di 30+ dB \`e considerato buona qualit\`a per il video.)
\end{examplebox}

\subsection{End-Point Error --- EPE (per optical flow)}

L'EPE \`e la metrica standard per valutare l'optical flow. Misura la distanza euclidea media tra il flusso predetto e il flusso ground truth.

\begin{formulabox}[End-Point Error --- Formula completa]
Per ogni pixel $(x, y)$, il flusso predetto \`e $\hat{\mathbf{f}}_{(x,y)} = (\hat{u}, \hat{v})$ e il flusso ground truth \`e $\mathbf{f}^*_{(x,y)} = (u^*, v^*)$, dove $u$ \`e lo spostamento orizzontale e $v$ quello verticale.

\textbf{EPE per singolo pixel}:
\[
\text{EPE}_{(x,y)} = \sqrt{(\hat{u} - u^*)^2 + (\hat{v} - v^*)^2} = \left\|\hat{\mathbf{f}}_{(x,y)} - \mathbf{f}^*_{(x,y)}\right\|_2
\]

\textbf{EPE medio} (usato come loss):
\[
\mathcal{L}_{\text{EPE}} = \frac{1}{|\mathcal{P}|} \sum_{(x,y) \in \mathcal{P}} \sqrt{(\hat{u}_{(x,y)} - u^*_{(x,y)})^2 + (\hat{v}_{(x,y)} - v^*_{(x,y)})^2}
\]
dove $\mathcal{P}$ \`e l'insieme di tutti i pixel validi.

\textbf{Gradiente} (con $\epsilon$ per stabilit\`a numerica):
\[
\frac{\partial \mathcal{L}_{\text{EPE}}}{\partial \hat{u}_{(x,y)}} = \frac{1}{|\mathcal{P}|} \cdot \frac{\hat{u}_{(x,y)} - u^*_{(x,y)}}{\sqrt{(\hat{u}_{(x,y)} - u^*_{(x,y)})^2 + (\hat{v}_{(x,y)} - v^*_{(x,y)})^2} + \epsilon}
\]

\textbf{Nota}: a differenza della MSE, la EPE usa la radice quadrata, quindi penalizza \textbf{meno} gli errori grandi (norma L2 vs norma L2 al quadrato). Questo rende l'ottimizzazione pi\`u robusta a outlier.
\end{formulabox}

\begin{examplebox}[Esempio numerico: EPE passo-passo]
Un pixel nel frame $t$ ha predizione di flusso $\hat{\mathbf{f}} = (2.1, -0.5)$ e ground truth $\mathbf{f}^* = (2.0, -0.3)$.

\textbf{Calcolo EPE}:
\begin{align*}
\text{EPE} &= \sqrt{(2.1 - 2.0)^2 + (-0.5 - (-0.3))^2} \\
&= \sqrt{(0.1)^2 + (-0.2)^2} \\
&= \sqrt{0.01 + 0.04} \\
&= \sqrt{0.05} = 0.224 \text{ pixel}
\end{align*}

\textbf{Interpretazione}: l'errore di $0.224$ pixel significa che la predizione di dove si trova questo pixel nel frame successivo \`e sbagliata di circa un quarto di pixel. Per il benchmark Sintel, il Perceiver IO raggiunge un EPE medio di $1.81$ pixel (clean) e $2.42$ pixel (final) --- errori di pochi pixel su frame di $\sim$400 pixel di larghezza.

\textbf{Confronto}: un EPE di $2.42$ su un'immagine $436 \times 1024$ corrisponde a un errore relativo di circa $\frac{2.42}{436} \approx 0.56\%$, molto accurato!
\end{examplebox}

\subsection{Loss combinata per l'autoencoding multimodale}

Nel task multimodale, il modello deve ottimizzare \textbf{tre loss diverse simultaneamente}.

\begin{formulabox}[Loss combinata multimodale --- Dettaglio completo]
\[
\mathcal{L}_{\text{multi}} = \underbrace{\lambda_v \cdot \frac{1}{|\mathcal{S}_v|}\sum_{i \in \mathcal{S}_v} \|\hat{\mathbf{y}}_i^v - \mathbf{y}_i^{v*}\|_2^2}_{\text{MSE video (ricostruzione pixel)}} + \underbrace{\lambda_a \cdot \frac{1}{|\mathcal{S}_a|}\sum_{i \in \mathcal{S}_a} (\hat{y}_i^a - y_i^{a*})^2}_{\text{MSE audio (ricostruzione campioni)}} + \underbrace{\lambda_l \cdot \left(-\sum_{k=1}^{700} y_k^* \log \hat{y}_k^l\right)}_{\text{CE label (classificazione)}}
\]

dove:
\begin{itemize}
    \item $\mathcal{S}_v$: sottoinsieme di pixel video campionati ($|\mathcal{S}_v| = 512$ durante training)
    \item $\mathcal{S}_a$: sottoinsieme di campioni audio campionati ($|\mathcal{S}_a| = 512$ durante training)
    \item $\lambda_v, \lambda_a, \lambda_l$: pesi per bilanciare le tre loss (iperparametri)
\end{itemize}

\textbf{Tradeoff}: Aumentare $\lambda_v$ e $\lambda_a$ migliora la ricostruzione (PSNR pi\`u alto) ma pu\`o peggiorare la classificazione. Aumentare $\lambda_l$ migliora l'accuracy ma pu\`o peggiorare la ricostruzione. Il paper esplora diversi rapporti di compressione ($88\times$, $176\times$, $352\times$) che influenzano questo tradeoff attraverso il numero di latenti $N$.
\end{formulabox}

\begin{notabox}[In parole semplici: la loss combinata]
Immagina di dover dare un \textbf{voto unico} a uno studente che ha fatto tre prove diverse: un tema (video), un dettato (audio) e un quiz a risposta multipla (classificazione). Non puoi dare tre voti separati --- devi combinarli in uno solo.

I pesi $\lambda_v, \lambda_a, \lambda_l$ sono come i ``crediti'' di ciascuna prova: se il tema vale il doppio del quiz, metti $\lambda_v = 2\lambda_l$. Il voto finale \`e la media pesata.

Il problema \`e che migliorare su una prova pu\`o peggiorare le altre (lo studente ha tempo limitato per studiare = i latenti hanno capacit\`a limitata). Ecco perch\'e il numero di latenti $N$ \`e cruciale: pi\`u latenti $=$ pi\`u ``tempo di studio'' $=$ migliori risultati su tutte le prove.
\end{notabox}


%% ============================================================
%%  SEZIONE 12 --- TOKENIZATION vs BYTE-LEVEL
%% ============================================================
\section{Approfondimento: Tokenizzazione vs.\ processing a livello di byte}

Questa \`e una delle innovazioni pi\`u importanti del Perceiver IO per il linguaggio. La comprensione di questa scelta \`e fondamentale per l'esame.

\subsection{Cos'\`e la tokenizzazione?}

La tokenizzazione \`e il processo di suddivisione del testo in unit\`a elementari (token) che il modello pu\`o elaborare.

\begin{formulabox}[Tipi di tokenizzazione]
\textbf{1. Word-level}: ogni parola \`e un token.
\begin{itemize}
    \item ``Il gatto dorme'' $\to$ [``Il'', ``gatto'', ``dorme''] (3 token)
    \item Problema: vocabolario enorme, parole rare/nuove non gestite (out-of-vocabulary, OOV)
\end{itemize}

\textbf{2. Character-level}: ogni carattere \`e un token.
\begin{itemize}
    \item ``Il gatto dorme'' $\to$ [``I'', ``l'', `` '', ``g'', ``a'', ``t'', ``t'', ``o'', ...] (15 token)
    \item Problema: sequenze molto lunghe, difficile catturare relazioni a lungo raggio
\end{itemize}

\textbf{3. Sub-word (BPE / WordPiece / SentencePiece)}: compromesso tra word e character.
\begin{itemize}
    \item BERT usa \textbf{WordPiece} con vocabolario di $\sim$30{,}000 token
    \item ``unbelievable'' $\to$ [``un'', ``\#\#believ'', ``\#\#able''] (3 token)
    \item ``gatto'' $\to$ [``gatto''] (1 token, parola comune)
    \item ``supercalifragilistichespiralidoso'' $\to$ [``super'', ``\#\#cal'', ``\#\#if'', ...] (molti token)
    \item Vantaggi: vocabolario gestibile, gestisce parole rare
    \item \textbf{Svantaggi}: dipende dalla lingua, richiede training del tokenizer, perde informazione a livello di carattere, overhead ingegneristico
\end{itemize}
\end{formulabox}

\subsection{L'approccio del Perceiver IO: byte grezzi UTF-8}

Per capire la differenza, pensate a come impariamo a leggere. Un bambino legge lettera per lettera (byte-level): \`e lento ma non richiede conoscenza pregressa. Un adulto legge per parole (token-level): \`e veloce ma richiede un vocabolario appreso. Il Perceiver IO torna al livello ``bambino'' (byte) ma compensa con la potenza computazionale del latent bottleneck --- come un bambino con una memoria fotografica che pu\`o processare migliaia di lettere contemporaneamente.

Il Perceiver IO adotta un approccio radicalmente diverso: \textbf{niente tokenizzazione}. L'input \`e la sequenza grezza di byte UTF-8.

\begin{formulabox}[Encoding UTF-8 --- Come funziona]
\textbf{UTF-8} \`e lo standard di codifica dei caratteri. Ogni carattere \`e rappresentato da 1--4 byte:

\begin{center}
\begin{tabular}{lccl}
\toprule
\textbf{Carattere} & \textbf{Byte (decimale)} & \textbf{Byte (hex)} & \textbf{Nota} \\
\midrule
``H'' & 72 & 0x48 & ASCII, 1 byte \\
``e'' & 101 & 0x65 & ASCII, 1 byte \\
``l'' & 108 & 0x6C & ASCII, 1 byte \\
`` '' (spazio) & 32 & 0x20 & ASCII, 1 byte \\
``\`e'' & 195, 168 & 0xC3, 0xA8 & 2 byte (accentata) \\
\bottomrule
\end{tabular}
\end{center}

\textbf{Per il Perceiver IO}:
\begin{itemize}
    \item Vocabolario: $256$ valori byte possibili + $4$ token speciali (PAD, MASK, CLS, SEP) = \textbf{260 token}
    \item Input: $M = 2048$ byte (lunghezza fissa, padding se necessario)
    \item ``Hello'' = $[72, 101, 108, 108, 111]$ (5 byte, 5 ``token'')
    \item ``The cat sat on the mat'' = 22 byte (22 ``token'')
    \item Confronto: BERT userebbe $\sim$6--8 token per la stessa frase
\end{itemize}
\end{formulabox}

\subsection{Perch\'e funziona con il Perceiver IO ma non con i Transformer standard?}

Questa \`e la domanda cruciale. Il motivo \`e la \textbf{complessit\`a computazionale}.

\begin{center}
\begin{tikzpicture}[
    block/.style={draw, rounded corners, minimum width=5.5cm, minimum height=2.5cm, align=center, font=\small},
    arr/.style={-{Stealth[length=2.5mm]}, thick}
]
    % BERT
    \node[block, fill=red!10, draw=red!50!black] (bert) at (0,0) {
        \textbf{BERT / Transformer}\\[4pt]
        $M = 512$ token WordPiece\\
        Complessit\`a per layer:\\
        $\mathcal{O}(M^2) = 512^2 = \mathbf{262{,}144}$\\[3pt]
        Se usasse byte: $M = 2048$\\
        $\mathcal{O}(M^2) = 2048^2 = \mathbf{4{,}194{,}304}$\\
        $16\times$ pi\`u costoso!
    };

    % Perceiver IO
    \node[block, fill=green!10, draw=green!50!black] (pio) at (8,0) {
        \textbf{Perceiver IO}\\[4pt]
        $M = 2048$ byte grezzi\\
        $N = 256$ latenti\\
        Complessit\`a per layer:\\
        $\mathcal{O}(M \cdot N) = 2048 \times 256 = \mathbf{524{,}288}$\\[3pt]
        Solo $2\times$ il costo di BERT\\con token, e $8\times$ meno\\del Transformer su byte!
    };

    % Arrow
    \draw[very thick, -{Stealth[length=4mm]}, blue!60!black] (3.5, 0) -- (4.5, 0);
    \node[font=\small, text=blue!60!black, align=center] at (4, 1) {Il bottleneck\\rende feasible\\il byte-level};
\end{tikzpicture}
\end{center}

\begin{proofbox}[Dimostrazione: perch\'e il byte-level \`e impossibile per i Transformer standard]
Consideriamo un testo di $T = 512$ token WordPiece. In media, ogni token corrisponde a $\sim$4 byte, quindi il testo ha $B = 2048$ byte.

\textbf{Con Transformer standard}:
\begin{itemize}
    \item Su token: complessit\`a $\mathcal{O}(T^2 \cdot d) = \mathcal{O}(512^2 \cdot d) = \mathcal{O}(262{,}144 \cdot d)$
    \item Su byte: complessit\`a $\mathcal{O}(B^2 \cdot d) = \mathcal{O}(2048^2 \cdot d) = \mathcal{O}(4{,}194{,}304 \cdot d)$
    \item Rapporto: $\frac{B^2}{T^2} = \frac{2048^2}{512^2} = 16\times$ pi\`u costoso
\end{itemize}
Anche con hardware moderno, un fattore $16\times$ rende il training proibitivo (BERT gi\`a richiede settimane su cluster di TPU).

\textbf{Con Perceiver IO} ($N = 256$ latenti):
\begin{itemize}
    \item Encoder: $\mathcal{O}(B \cdot N \cdot D) = \mathcal{O}(2048 \cdot 256 \cdot D) = \mathcal{O}(524{,}288 \cdot D)$
    \item Processor: $\mathcal{O}(L \cdot N^2 \cdot D) = \mathcal{O}(L \cdot 65{,}536 \cdot D)$ (indipendente da $B$!)
    \item Il costo dell'encoder \`e solo $2\times$ quello di BERT su token, e il processor non dipende da $B$
\end{itemize}

\textbf{Conclusione}: il bottleneck latente ``assorbe'' il costo aggiuntivo delle sequenze pi\`u lunghe, rendendo il byte-level processing computazionalmente fattibile.
\end{proofbox}

\subsection{Risultati: confronto diretto byte vs.\ token}

\begin{center}
\begin{tabular}{lcccc}
\toprule
\textbf{Modello} & \textbf{Input} & \textbf{Vocabolario} & \textbf{$M$} & \textbf{GLUE Avg} \\
\midrule
BERT-Base & SentencePiece token & $\sim$30{,}000 & 512 & 79.1 \\
BERT-Large & SentencePiece token & $\sim$30{,}000 & 512 & 81.0 \\
\midrule
Perceiver IO (SP) & SentencePiece token & $\sim$30{,}000 & 512 & 81.2 \\
Perceiver IO (byte) & Byte UTF-8 grezzi & 260 & 2048 & 81.0 \\
\textbf{Perceiver IO++ (byte)} & Byte UTF-8 grezzi & \textbf{260} & \textbf{2048} & \textbf{81.8} \\
\bottomrule
\end{tabular}
\end{center}

\begin{notabox}[In parole semplici: byte-level processing]
Invece di leggere un libro \textbf{parola per parola} (tokenizzazione), il Perceiver IO legge \textbf{lettera per lettera} (byte). Sembra pi\`u lento e meno efficiente --- come se un lettore leggesse ``I-l g-a-t-t-o d-o-r-m-e'' invece di ``Il gatto dorme''.

Ma ecco il trucco: il Perceiver IO ha un ``cervello'' (i latenti) di dimensione fissa. Che tu legga per lettere o per parole, il cervello comprime tutto nello stesso spazio. E leggere per lettere ha un vantaggio enorme: \textbf{non devi imparare dove iniziano e finiscono le parole}. Funziona per qualsiasi lingua, qualsiasi alfabeto, qualsiasi testo --- anche emoji, codice sorgente, o lingue mai viste durante l'addestramento.

Il risultato? Il Perceiver IO raggiunge \textbf{81.8 su GLUE} leggendo lettera per lettera, \textbf{superando BERT-Large} (81.0) che legge per parole (token). La tokenizzazione non \`e necessaria!
\end{notabox}


%% ============================================================
%%  SEZIONE 13 --- OUTPUT QUERY SUBSAMPLING
%% ============================================================
\section{Approfondimento: Subsampling delle output query}

Il subsampling delle output query \`e una tecnica cruciale che rende possibile l'addestramento su task con output enormi (come l'autoencoding multimodale).

\subsection{Il problema: output troppo grandi}

\begin{examplebox}[Il problema dell'autoencoding multimodale]
Nell'autoencoding su Kinetics-700:
\begin{itemize}
    \item \textbf{Video}: $16 \times 56 \times 56 = 50{,}176$ patch
    \item \textbf{Audio}: $30{,}720$ campioni
    \item \textbf{Label}: $1$
    \item \textbf{Totale}: $O = 80{,}897$ punti di output
\end{itemize}

Il decoder deve calcolare una cross-attention di dimensione $O \times N = 80{,}897 \times 512$:
\begin{itemize}
    \item Matrice di attenzione: $80{,}897 \times 512 \approx 41$M valori
    \item Per ogni testa: proiezione Q, K, V + softmax + combinazione
    \item Con 8 teste: $\sim$330M operazioni
    \item Pi\`u la memoria per lo storage della matrice di attenzione e dei gradienti
\end{itemize}

\textbf{Questo \`e fattibile} per un singolo forward pass (inferenza), ma per il \textbf{training} serve anche il backward pass, che raddoppia la memoria. E con batch size 32, la memoria totale diventa proibitiva.
\end{examplebox}

\subsection{La soluzione: campionamento casuale}

\begin{formulabox}[Subsampling delle output query --- Procedura]
\textbf{Durante il training}:
\begin{enumerate}
    \item Campiona casualmente $|\mathcal{S}_v| = 512$ indici video da $\{1, \ldots, 50{,}176\}$
    \item Campiona casualmente $|\mathcal{S}_a| = 512$ indici audio da $\{1, \ldots, 30{,}720\}$
    \item Includi sempre la label ($O_l = 1$)
    \item Totale query per step: $O_{\text{train}} = 512 + 512 + 1 = 1{,}025$ invece di $80{,}897$
    \item Calcola il decoder e la loss \textbf{solo} su queste $1{,}025$ query
\end{enumerate}

\textbf{Durante l'inferenza}:
\begin{enumerate}
    \item Decodifica \textbf{tutte} le $80{,}897$ query
    \item Pu\`o essere fatto in batch (es.\ $10{,}000$ query alla volta) se la memoria \`e limitata
\end{enumerate}

\textbf{Risparmio}: $\frac{80{,}897}{1{,}025} \approx 79\times$ meno query per step di training!
\end{formulabox}

\subsection{Perch\'e funziona: prova del gradiente non distorto}

Questa \`e la giustificazione teorica fondamentale del subsampling.

\begin{proofbox}[Dimostrazione: il gradiente del subsampling \`e non distorto]
\textbf{Loss completa}:
\[
\mathcal{L} = \frac{1}{O}\sum_{i=1}^{O} \ell_i(\mathbf{y}_i, \mathbf{y}_i^*)
\]
dove $\ell_i$ \`e la loss per il singolo punto di output $i$ (MSE per video/audio, CE per label).

\textbf{Loss con subsampling}: sia $\mathcal{S} \subset \{1, \ldots, O\}$ un sottoinsieme campionato uniformemente a caso, con $|\mathcal{S}| = S \ll O$:
\[
\hat{\mathcal{L}} = \frac{1}{S}\sum_{i \in \mathcal{S}} \ell_i(\mathbf{y}_i, \mathbf{y}_i^*)
\]

\textbf{Valore atteso della loss campionata}:
\begin{align}
\mathbb{E}_{\mathcal{S}}[\hat{\mathcal{L}}] &= \mathbb{E}_{\mathcal{S}}\left[\frac{1}{S}\sum_{i \in \mathcal{S}} \ell_i\right] \\
&= \frac{1}{S} \sum_{i \in \mathcal{S}} \mathbb{E}[\ell_i] \tag{linearit\`a del valore atteso} \\
&= \frac{1}{S} \cdot S \cdot \frac{1}{O}\sum_{i=1}^{O} \ell_i \tag{ogni $\ell_i$ ha probabilit\`a $S/O$ di essere selezionato} \\
&= \frac{1}{O}\sum_{i=1}^{O} \ell_i = \mathcal{L} \tag{definizione di $\mathcal{L}$}
\end{align}

Quindi $\mathbb{E}[\hat{\mathcal{L}}] = \mathcal{L}$: la loss campionata \`e uno \textbf{stimatore non distorto} della loss completa.

\textbf{Per il gradiente}: poich\'e la derivata \`e un operatore lineare e il valore atteso commuta con la derivata (sotto condizioni di regolarit\`a):
\[
\mathbb{E}_{\mathcal{S}}\left[\nabla_\theta \hat{\mathcal{L}}\right] = \nabla_\theta \mathbb{E}_{\mathcal{S}}[\hat{\mathcal{L}}] = \nabla_\theta \mathcal{L}
\]

Quindi anche il \textbf{gradiente campionato} \`e uno stimatore non distorto del gradiente completo!

\textbf{Condizione cruciale}: questo funziona perch\'e le query di output sono \textbf{indipendenti} nel decoder (non c'\`e self-attention tra di esse). Se ci fosse interazione tra le query, la loss $\ell_i$ dipenderebbe da \emph{quali altre query} sono nel batch, e il campionamento introdurrebbe un bias.
\end{proofbox}

\begin{notabox}[In parole semplici: il subsampling]
Immagina di essere un professore che deve correggere 80{,}000 esercizi di 80{,}000 studenti. Correggerli tutti ad ogni iterazione di feedback \`e impossibile. Allora fai cos\`i: ad ogni giro, \textbf{scegli a caso 1{,}000 esercizi} e dai feedback solo su quelli.

Funziona perch\'e ogni esercizio \`e \textbf{indipendente} dagli altri: il voto dell'esercizio 42 non dipende dal voto dell'esercizio 7{,}891. Quindi il feedback medio sui 1{,}000 esercizi campionati \`e rappresentativo del feedback che daresti sull'intero set di 80{,}000.

Dopo molte iterazioni, ogni esercizio sar\`a stato corretto pi\`u volte, e lo studente (il modello) avr\`a migliorato su tutti e 80{,}000 gli esercizi, anche se ad ogni giro ne ha visti solo 1{,}000.

\textbf{Attenzione}: questo NON funzionerebbe se il voto di un esercizio dipendesse dagli altri (come in un esame di gruppo). Ecco perch\'e \`e fondamentale che il decoder del Perceiver IO non abbia self-attention tra le query!
\end{notabox}


%% ============================================================
%%  SEZIONE 14 --- WALKTHROUGH DETTAGLIATI
%% ============================================================
\section{Walkthrough dettagliati per ogni dominio sperimentale}

In questa sezione seguiamo il flusso dei dati attraverso l'intero Perceiver IO per tre domini specifici, passo dopo passo, con tutte le dimensioni esplicite.

\subsection{Walkthrough: Masked Language Modeling}

Seguiamo il flusso completo per la frase ``The cat sat on the [MASK]''.

\begin{examplebox}[Walkthrough completo: MLM su ``The cat sat on the {[}MASK{]}'']
\textbf{Passo 0: Encoding in byte UTF-8}

La frase ``The cat sat on the [MASK]'' (senza [MASK], che verr\`a sostituito):

Testo originale: ``The cat sat on the mat''

Byte UTF-8: $[84, 104, 101, 32, 99, 97, 116, 32, 115, 97, 116, 32, 111, 110, 32, 116, 104, 101, 32, 109, 97, 116]$ (22 byte)

Con padding a $M = 2048$: i restanti $2048 - 22 = 2026$ posizioni vengono riempite con il token PAD (indice 256).

\textbf{Passo 1: Mascheramento (15\% delle parole)}

Si maschera la parola ``mat'' $\to$ i byte alle posizioni 19, 20, 21 vengono sostituiti con il token MASK (indice 257).

Input mascherato: $[84, 104, 101, 32, 99, 97, 116, 32, 115, 97, 116, 32, 111, 110, 32, 116, 104, 101, 32, \textbf{257, 257, 257}, 256, 256, \ldots]$

\textbf{Passo 2: Embedding}

Ogni byte $\to$ vettore appreso $\in \mathbb{R}^{C}$ ($C = 1280$ nel paper).

$\mathbf{x} \in \mathbb{R}^{2048 \times 1280}$

Si aggiunge un embedding posizionale appreso per ogni posizione $i \in \{0, \ldots, 2047\}$:
$\mathbf{x}_i = \text{ByteEmbed}(\text{byte}_i) + \text{PosEmbed}(i)$

\textbf{Passo 3: Encoder (cross-attention)}

$N = 256$ latenti iniziali appresi $\mathbf{z}^{(0)} \in \mathbb{R}^{256 \times 1280}$.

Cross-attention: Q dai latenti, K e V dall'input.
\begin{itemize}
    \item $\mathbf{Q} = \mathbf{z}^{(0)} \mathbf{W}_Q \in \mathbb{R}^{256 \times 64}$ (per testa, 8 teste)
    \item $\mathbf{K} = \mathbf{x} \mathbf{W}_K \in \mathbb{R}^{2048 \times 64}$
    \item $\mathbf{V} = \mathbf{x} \mathbf{W}_V \in \mathbb{R}^{2048 \times 64}$
    \item Attenzione: $\mathbb{R}^{256 \times 2048}$ (ogni latente guarda tutti i 2048 byte)
    \item Output: $\mathbb{R}^{256 \times 1280}$
\end{itemize}

I 256 latenti ora contengono un riassunto dell'intera sequenza di 2048 byte.

\textbf{Passo 4: Processor (26 layer di self-attention)}

Per $\ell = 1, \ldots, 26$:
\begin{itemize}
    \item Self-attention tra i 256 latenti: $\mathbb{R}^{256 \times 256}$ matrice di attenzione
    \item FFN: $1280 \to 5120 \to 1280$ con GELU
    \item Residual + LayerNorm
\end{itemize}

Dopo 26 layer: $\mathbf{z}^{(26)} \in \mathbb{R}^{256 \times 1280}$ (latenti ``raffinati'').

\textbf{Passo 5: Decoder}

Output query: solo per le 3 posizioni mascherate (19, 20, 21).

$\mathbf{q} = [\text{LearnedPE}(19), \text{LearnedPE}(20), \text{LearnedPE}(21)] \in \mathbb{R}^{3 \times 1280}$

Cross-attention: Q dalle query, K e V dai latenti.
\begin{itemize}
    \item $\mathbf{Q}^{\text{dec}} \in \mathbb{R}^{3 \times 64}$ (per testa)
    \item $\mathbf{K}^{\text{dec}}, \mathbf{V}^{\text{dec}} \in \mathbb{R}^{256 \times 64}$
    \item Attenzione: $\mathbb{R}^{3 \times 256}$ (3 query $\times$ 256 latenti)
    \item Output: $\mathbb{R}^{3 \times 1280}$
\end{itemize}

\textbf{Passo 6: Proiezione all'output}

MLP: $1280 \to 260$ (260 = 256 byte + 4 token speciali)

$\mathbf{y} \in \mathbb{R}^{3 \times 260}$

Softmax per riga $\to$ distribuzione di probabilit\`a su 260 token per ciascuna delle 3 posizioni.

\textbf{Passo 7: Loss}

Target: byte originali di ``mat'' = $[109, 97, 116]$

$\mathcal{L} = -\frac{1}{3}\left[\log p(\text{byte } 109 \mid \text{pos } 19) + \log p(\text{byte } 97 \mid \text{pos } 20) + \log p(\text{byte } 116 \mid \text{pos } 21)\right]$
\end{examplebox}

\subsection{Walkthrough: Optical Flow su Sintel}

\begin{examplebox}[Walkthrough completo: optical flow su Sintel]
\textbf{Passo 0: Preparazione dell'input}

Due frame consecutivi di Sintel, risoluzione $384 \times 160$ (largura $\times$ altezza, risoluzione di training).

\textbf{Passo 1: Estrazione patch}

Per ogni pixel $(x, y)$, si estrae una patch $3 \times 3$ da \textbf{entrambi} i frame:
\begin{itemize}
    \item Frame 1: patch $3 \times 3 \times 3$ (RGB) = $27$ valori
    \item Frame 2: patch $3 \times 3 \times 3$ = $27$ valori
    \item Concatenazione: $27 + 27 = 54$ feature per pixel
\end{itemize}

$M = 384 \times 160 = 61{,}440$ pixel

\textbf{Passo 2: Position encoding}

Fourier PE 2D per ogni pixel $(x, y)$ normalizzato in $[-1, 1]$:
\begin{itemize}
    \item Frequenze: $\{1, 2, 4, 8, 16, 32, 64\}$ ($B = 7$ bande)
    \item Per ogni dimensione e frequenza: $\sin(\pi f \cdot x)$ e $\cos(\pi f \cdot x)$
    \item Dimensione PE: $2 \times 7 \times 2 = 28$ per pixel
\end{itemize}

Input totale per pixel: $54 + 28 = 82$ feature $\to$ $C = 82$

$\mathbf{x} \in \mathbb{R}^{61{,}440 \times 82}$

\textbf{Passo 3: Encoder}

$N = 2{,}048$ latenti (pi\`u del solito, per catturare le corrispondenze fine tra pixel).

$\mathbf{z}^{(0)} \in \mathbb{R}^{2{,}048 \times 512}$ (latenti iniziali appresi)

Cross-attention: matrice $2{,}048 \times 61{,}440$ (ogni latente guarda tutti i pixel).

\textbf{Passo 4: Processor (24 layer)}

Self-attention tra i 2{,}048 latenti: matrice $2{,}048 \times 2{,}048$.

I latenti scambiano informazione sulle corrispondenze globali tra i due frame. Dopo 24 layer, le relazioni tra pixel lontani sono state elaborate.

$\mathbf{z}^{(24)} \in \mathbb{R}^{2{,}048 \times 512}$

\textbf{Passo 5: Decoder}

Per ogni pixel $(x, y)$, la query \`e:
\[
\mathbf{q}_{(x,y)} = \text{concat}[\text{FourierPE}(x, y), \mathbf{f}_{(x,y)}^{\text{input}}] \in \mathbb{R}^{E}
\]
dove $\mathbf{f}_{(x,y)}^{\text{input}}$ sono le feature dal primo frame (non dal secondo!).

$O = 61{,}440$ query (una per pixel).

Cross-attention: matrice $61{,}440 \times 2{,}048$.

Output: $\mathbb{R}^{61{,}440 \times 512}$ $\to$ MLP $\to$ $\mathbb{R}^{61{,}440 \times 2}$

\textbf{Passo 6: Output}

Per ogni pixel: $(dx, dy) \in \mathbb{R}^2$ = vettore di spostamento.

\textbf{Passo 7: Loss}

$\mathcal{L} = \frac{1}{61{,}440}\sum_{(x,y)} \sqrt{(dx_{\text{pred}} - dx_{\text{true}})^2 + (dy_{\text{pred}} - dy_{\text{true}})^2}$
\end{examplebox}

\subsection{Walkthrough: Autoencoding multimodale}

\begin{examplebox}[Walkthrough completo: autoencoding multimodale su Kinetics-700]
\textbf{Passo 0: Preparazione dell'input multimodale}

\textbf{Video}: 16 frame a $224 \times 224$, suddivisi in patch $4 \times 4$.
\begin{itemize}
    \item Ogni frame: $\frac{224}{4} \times \frac{224}{4} = 56 \times 56 = 3{,}136$ patch
    \item 16 frame: $16 \times 3{,}136 = 50{,}176$ patch
    \item Ogni patch: $4 \times 4 \times 3 = 48$ feature (RGB per ogni pixel della patch)
    \item + Fourier PE 3D (per $x$, $y$, $t$) + embedding di modalit\`a (``is\_video'')
\end{itemize}

\textbf{Audio}: $30{,}720$ campioni (circa 2 secondi a 16kHz).
\begin{itemize}
    \item Ogni campione: 1 valore scalare
    \item + Fourier PE 1D (per $t$) + embedding di modalit\`a (``is\_audio'')
\end{itemize}

\textbf{Label}: 1 token = vettore one-hot su 700 classi.
\begin{itemize}
    \item + embedding di modalit\`a (``is\_label'')
    \item \textbf{Mascherata al 50\%}: met\`a delle volte il token label viene sostituito con zeri
\end{itemize}

\textbf{Passo 1: Padding e serializzazione}

Tutti gli elementi vengono portati alla stessa dimensione di canale ($C = 704$) tramite padding con zeri, poi concatenati:

$\mathbf{x} = [\underbrace{\mathbf{x}_{\text{video}}}_{50{,}176 \times 704}; \; \underbrace{\mathbf{x}_{\text{audio}}}_{30{,}720 \times 704}; \; \underbrace{\mathbf{x}_{\text{label}}}_{1 \times 704}] \in \mathbb{R}^{80{,}897 \times 704}$

\textbf{Passo 2: Encoder}

$N = 512$ latenti, $D = 1{,}024$.

Cross-attention: $512 \times 80{,}897$ (ogni latente guarda TUTTO: video, audio e label insieme).

$\mathbf{z}^{(0)}_{\text{out}} \in \mathbb{R}^{512 \times 1{,}024}$

\textbf{Passo 3: Processor (8 layer)}

Self-attention tra 512 latenti per 8 layer. I latenti integrano informazione cross-modale.

$\mathbf{z}^{(8)} \in \mathbb{R}^{512 \times 1{,}024}$

\textbf{Passo 4: Decoder (con subsampling)}

\textbf{Durante il training}: campiona 512 query video + 512 query audio + 1 query label = 1{,}025 totali.

Ciascuna query:
\begin{itemize}
    \item Video: $\mathbf{q}_v^{(x,y,t)} = [\text{FourierPE}(x, y, t) \| \mathbf{e}_{\text{video}}]$
    \item Audio: $\mathbf{q}_a^{(t)} = [\text{FourierPE}(t) \| \mathbf{e}_{\text{audio}}]$
    \item Label: $\mathbf{q}_l = \mathbf{e}_{\text{label}}$
\end{itemize}

Cross-attention: $1{,}025 \times 512$ (1{,}025 query $\times$ 512 latenti).

\textbf{Passo 5: Output}
\begin{itemize}
    \item 512 pixel ricostruiti: $\hat{\mathbf{y}}_v \in \mathbb{R}^{512 \times 48}$ (RGB per patch $4 \times 4$)
    \item 512 campioni audio ricostruiti: $\hat{\mathbf{y}}_a \in \mathbb{R}^{512 \times 1}$
    \item 1 vettore di logit: $\hat{\mathbf{y}}_l \in \mathbb{R}^{1 \times 700}$
\end{itemize}

\textbf{Passo 6: Loss combinata}

$\mathcal{L} = \lambda_v \cdot \text{MSE}_{\text{video}} + \lambda_a \cdot \text{MSE}_{\text{audio}} + \lambda_l \cdot \text{CE}_{\text{label}}$
\end{examplebox}

\begin{center}
\begin{tikzpicture}[
    stepbox/.style={draw, rounded corners, minimum width=10cm, minimum height=1cm, align=center, font=\scriptsize},
    arr/.style={-{Stealth[length=2mm]}, thick},
    dim/.style={font=\tiny\itshape, text=gray!60!black}
]
    \node[stepbox, fill=red!10] (s0) at (0,8) {\textbf{Input}: Video ($50{,}176$ patch) + Audio ($30{,}720$ campioni) + Label ($1$ token)};
    \node[dim, right=0.2cm of s0] {$80{,}897 \times 704$};

    \node[stepbox, fill=orange!10] (s1) at (0,6.5) {\textbf{Encoder}: Cross-attention, $N=512$ latenti guardano tutto l'input};
    \node[dim, right=0.2cm of s1] {$512 \times 1{,}024$};

    \node[stepbox, fill=blue!10] (s2) at (0,5) {\textbf{Processor}: 8 layer self-attention sui 512 latenti};
    \node[dim, right=0.2cm of s2] {$512 \times 1{,}024$};

    \node[stepbox, fill=green!10] (s3) at (0,3.5) {\textbf{Decoder}: 1{,}025 query (512 video + 512 audio + 1 label)};
    \node[dim, right=0.2cm of s3] {$1{,}025 \times 1{,}024$};

    \node[stepbox, fill=teal!10] (s4) at (0,2) {\textbf{Output}: 512 pixel + 512 audio + 1 label $\to$ Loss combinata};

    \draw[arr] (s0) -- (s1);
    \draw[arr] (s1) -- (s2);
    \draw[arr] (s2) -- (s3);
    \draw[arr] (s3) -- (s4);

    \node[draw, dashed, rounded corners, fill=orange!5, minimum width=4cm, font=\tiny, align=center] at (6, 3.5) {Subsampling:\\$80{,}897 \to 1{,}025$\\($79\times$ risparmio)};
\end{tikzpicture}
\end{center}


%% ============================================================
%%  SEZIONE 15 --- TABELLE MANCANTI E RISULTATI ADDIZIONALI
%% ============================================================
\section{Tabelle e risultati addizionali dal paper}

In questa sezione raccogliamo risultati e tabelle addizionali dal paper che completano il quadro sperimentale.

\subsection{Table 3 --- Risultati Optical Flow dettagliati per dataset}

Il paper riporta risultati sia sui \textbf{training set} (``train'') sia sui \textbf{benchmark pubblici} (``test'') per Sintel e KITTI. I risultati pi\`u importanti sono quelli sui set di test, che vengono valutati su un server esterno.

\begin{center}
\begin{tabular}{lcccccc}
\toprule
\textbf{Modello} & \multicolumn{2}{c}{\textbf{Sintel Clean}} & \multicolumn{2}{c}{\textbf{Sintel Final}} & \multicolumn{2}{c}{\textbf{KITTI-15}} \\
 & train & test & train & test & train & test \\
 & EPE & EPE & EPE & EPE & EPE & Fl-all \\
\midrule
FlowNet2 & 2.02 & 3.96 & 3.14 & 6.02 & -- & -- \\
PWC-Net+ & 1.71 & 3.45 & 2.34 & 4.60 & -- & 7.72 \\
RAFT & 0.76 & 1.94 & 1.22 & 3.18 & 0.63 & 5.10 \\
GMA & 0.62 & 1.39 & 1.06 & 2.47 & -- & 5.15 \\
\midrule
\textbf{Perceiver IO} & 0.94 & \textbf{1.81} & 1.66 & \textbf{2.42} & -- & \textbf{4.98} \\
\bottomrule
\end{tabular}
\end{center}

\begin{notabox}[Analisi train vs.\ test per l'optical flow]
Un'osservazione interessante: sul \textbf{training set}, GMA e RAFT hanno EPE inferiore al Perceiver IO (es.\ $0.62$ vs $0.94$ su Sintel Clean train). Ma sul \textbf{test set}, il Perceiver IO \`e molto competitivo ($1.81$ vs $1.39$ per GMA).

Questo suggerisce che il Perceiver IO potrebbe avere una \textbf{migliore generalizzazione}: la differenza train-test \`e minore ($0.94 \to 1.81$, fattore $1.93\times$) rispetto a GMA ($0.62 \to 1.39$, fattore $2.24\times$) e RAFT ($0.76 \to 1.94$, fattore $2.55\times$). L'architettura general-purpose potrebbe essere meno soggetta all'overfitting.
\end{notabox}

\subsection{Table 4 --- Risultati multimodali per diversi rapporti di compressione}

Questa tabella mostra come la qualit\`a della ricostruzione e della classificazione varia con il numero di latenti (e quindi con il rapporto di compressione).

\begin{center}
\begin{tabular}{lccccc}
\toprule
\textbf{Compressione} & \textbf{$N$} & \textbf{Video PSNR} & \textbf{Audio PSNR} & \textbf{Top-1 Acc.} & \textbf{Top-5 Acc.} \\
\midrule
$88\times$ & 920 & 33.8 dB & 36.2 dB & 45.6\% & 73.4\% \\
$176\times$ & 460 & 32.5 dB & 34.8 dB & 44.2\% & 72.1\% \\
$352\times$ & 230 & 30.1 dB & 32.1 dB & 41.0\% & 68.2\% \\
\midrule
Solo classificazione & 512 & -- & -- & 68.0\% & -- \\
\bottomrule
\end{tabular}
\end{center}

\begin{notabox}[Tradeoff compressione-qualit\`a]
I risultati mostrano un tradeoff chiaro:
\begin{enumerate}
    \item \textbf{Pi\`u latenti = migliore qualit\`a}: passando da $N = 230$ ($352\times$ compressione) a $N = 920$ ($88\times$), il PSNR video migliora di $3.7$ dB e l'accuracy di $4.6\%$
    \item \textbf{L'accuracy multimodale \`e inferiore alla sola classificazione}: $45.6\%$ vs $68.0\%$. Questo perch\'e i latenti devono contenere informazione per \emph{tre} task simultaneamente, non solo per la classificazione
    \item \textbf{Il PSNR audio \`e pi\`u alto del PSNR video}: ricostruire l'audio \`e ``pi\`u facile'' perch\'e ha meno informazione (1D vs 3D) e meno variabilit\`a
\end{enumerate}
\end{notabox}

\subsection{Risultati AudioSet completi}

\begin{center}
\begin{tabular}{lcc}
\toprule
\textbf{Modello} & \textbf{mAP} & \textbf{Architettura} \\
\midrule
ResNet-50 (audio) & 0.314 & CNN 2D su spettrogramma \\
AST (Audio Spectrogram Transformer) & 0.459 & ViT su spettrogramma \\
\textbf{Perceiver IO} & 0.421 & General purpose \\
\bottomrule
\end{tabular}
\end{center}

\begin{notabox}[AudioSet: discussione]
Il Perceiver IO raggiunge un mAP di $0.421$, superiore a ResNet-50 ma inferiore ad AST ($0.459$). Questo \`e notevole considerando che:
\begin{itemize}
    \item Il Perceiver IO \textbf{non ha componenti specifiche per l'audio} (no mel-filterbank specifiche, no data augmentation audio-specifica come SpecAugment)
    \item AST usa un ViT pre-addestrato su ImageNet e poi fine-tuned su AudioSet, sfruttando il transfer learning da immagini ad audio
    \item Il risultato dimostra che un'architettura general-purpose \`e competitiva anche su task audio-specifici
\end{itemize}
\end{notabox}


%% ============================================================
%%  SEZIONE 16 --- RIASSUNTO FINALE (era sezione 10)
%% ============================================================
\section{Riassunto finale}

\begin{formulabox}[Punti chiave del Perceiver IO]
\begin{enumerate}
    \item \textbf{Architettura Encode-Process-Decode}: encoder (cross-attention input$\to$latenti), processor (self-attention sui latenti), decoder (cross-attention latenti$\to$output con query)

    \item \textbf{Output query}: il meccanismo fondamentale che permette output strutturati arbitrari. Ogni query codifica ``cosa voglio sapere'' e interroga indipendentemente i latenti

    \item \textbf{Complessit\`a}: $\mathcal{O}(MN + LN^2 + ON)$, lineare in $M$ (input) e $O$ (output), dove $N \ll M, O$ \`e il numero di latenti

    \item \textbf{Generalit\`a}: un'unica architettura per classificazione, optical flow, NLU (su byte grezzi), multi-task, autoencoding multimodale, e giochi (StarCraft II)

    \item \textbf{Risultati}: supera BERT su GLUE senza tokenizzazione (81.8 vs 81.0); SOTA su Sintel per optical flow (2.42 EPE); autoencoding video+audio+label simultaneo; 84.5\% su ImageNet con pre-training; 87\% win rate su StarCraft II con 3.5x meno FLOPs

    \item \textbf{Design delle query}: la scelta della struttura delle query \`e il principale ``design decision'' task-specifico. Tutto il resto dell'architettura rimane invariato

    \item \textbf{Decoder singolo layer}: il decoder usa un singolo layer di cross-attention senza self-attention tra le query, per efficienza ($\mathcal{O}(ON)$ vs $\mathcal{O}(O^2)$) e perch\'e tutta l'informazione \`e gi\`a nei latenti

    \item \textbf{Subsampling}: durante il training, le output query possono essere campionate casualmente, producendo un gradiente non distorto grazie all'indipendenza delle query
\end{enumerate}
\end{formulabox}

\begin{center}
\begin{tikzpicture}[
    concept/.style={draw, rounded corners, minimum width=2.5cm, minimum height=1cm, align=center, font=\scriptsize\bfseries},
    arr/.style={-{Stealth[length=2mm]}, thick, gray}
]
    % Central node
    \node[concept, fill=blue!20, draw=blue!60!black, minimum width=3.5cm, minimum height=1.5cm] (pio) at (0,0) {Perceiver IO\\Encode-Process-Decode};

    % Surrounding concepts
    \node[concept, fill=red!10] (cls) at (-5, 2.5) {Classificazione\\(singola query)};
    \node[concept, fill=orange!10] (flow) at (-5, 0) {Optical Flow\\(query per-pixel)};
    \node[concept, fill=yellow!10] (lang) at (-5, -2.5) {Linguaggio\\(query per-byte)};
    \node[concept, fill=green!10] (multi) at (5, 2.5) {Multi-task\\(task embedding)};
    \node[concept, fill=purple!10] (modal) at (5, 0) {Multimodale\\(modality flag)};
    \node[concept, fill=teal!10] (game) at (5, -2.5) {StarCraft II\\(entity query)};

    % Arrows
    \draw[arr] (pio) -- (cls);
    \draw[arr] (pio) -- (flow);
    \draw[arr] (pio) -- (lang);
    \draw[arr] (pio) -- (multi);
    \draw[arr] (pio) -- (modal);
    \draw[arr] (pio) -- (game);

    % Central annotation
    \node[font=\tiny\itshape, text=blue!60!black] at (0, -1.2) {Stessa architettura, diverse query};
\end{tikzpicture}
\end{center}

\vspace{1cm}

\begin{notabox}[Messaggio chiave per l'esame]
Il Perceiver IO dimostra che un'architettura basata su \textbf{attenzione latente} e \textbf{output query} pu\`o essere sufficientemente generale da competere con modelli specializzati su una vasta gamma di task. L'idea fondamentale \`e semplice: \textbf{disaccoppiare} la dimensione dell'input e dell'output dal costo computazionale attraverso un bottleneck latente, e usare \textbf{query specifiche per task} per estrarre l'informazione necessaria. Questa architettura rappresenta un passo importante verso modelli di deep learning veramente \textbf{general-purpose}.

I tre punti da ricordare assolutamente:
\begin{enumerate}
    \item \textbf{Il decoder \`e l'innovazione}: la cross-attention con output query \`e ci\`o che distingue il Perceiver IO dal Perceiver originale e permette output strutturati
    \item \textbf{La complessit\`a \`e lineare}: $\mathcal{O}(MN + LN^2 + ON)$ con $N \ll M, O$, ottenuta attraverso il bottleneck latente
    \item \textbf{Le query codificano il task}: posizione, task, modalit\`a e feature dell'input --- tutto \`e nelle query
\end{enumerate}
\end{notabox}


%% ============================================================
%%  APPENDICE --- CONTENUTI AGGIUNTIVI DAL PAPER
%% ============================================================
% Sezione di appendice P2 (non usa \appendix, numerazione araba)
\section{Appendice: dettagli aggiuntivi dal paper}

\subsection{Configurazione completa per ogni esperimento}

\begin{center}
\footnotesize
\begin{tabular}{p{2.5cm}p{1.5cm}p{1.5cm}p{1.5cm}p{1.5cm}p{1.5cm}p{1.5cm}}
\toprule
\textbf{Parametro} & \textbf{GLUE} & \textbf{GLUE++} & \textbf{Flow} & \textbf{Multimod.} & \textbf{ImageNet} & \textbf{SC2} \\
\midrule
Latenti $N$ & 256 & 256 & 2,048 & 512 & 512 & 256 \\
Dim.\ latente $D$ & 1,280 & 1,536 & 512 & 1,024 & 1,024 & 512 \\
Num.\ teste $h$ & 8 & 8 & 16 & 8 & 8 & 8 \\
Layer totali & 26 & 40 & 24 & 8 & 48 & 8 \\
Blocchi condivisi & 1 & 1 & 1 & 1 & 8 & 1 \\
Layer per blocco & 26 & 40 & 24 & 8 & 6 & 8 \\
FFN hidden dim & 5,120 & 6,144 & 2,048 & 4,096 & 4,096 & 2,048 \\
Parametri totali & 201M & 425M & $\sim$170M & $\sim$100M & $\sim$300M & 140M \\
\midrule
Ottimizzatore & LAMB & LAMB & Adam & LAMB & LAMB & Adam \\
Learning rate & $10^{-3}$ & $10^{-3}$ & $2\!\times\!10^{-4}$ & $10^{-3}$ & $2\!\times\!10^{-3}$ & $10^{-4}$ \\
Batch size & 64 & 64 & 8 & 32 & 1024 & 128 \\
Weight decay & 0.1 & 0.1 & 0.0 & 0.1 & 0.1 & 0.0 \\
Grad clip norm & 10 & 10 & 1 & 10 & 10 & -- \\
Dropout & 0.0 & 0.0 & 0.0 & 0.0 & 0.0 & 0.0 \\
\bottomrule
\end{tabular}
\end{center}

\subsection{Dettagli sull'encoding posizionale di Fourier}

Il Perceiver IO usa \textbf{Fourier position encodings} (Jaegle et al., 2021) per codificare le posizioni dell'input e delle query. Per una posizione $\mathbf{p} = (p_1, \ldots, p_d)$ normalizzata in $[-1, 1]^d$:

\begin{formulabox}[Fourier Position Encoding]
\[
\text{FourierPE}(\mathbf{p}) = \left[\sin(\pi f_1 p_1), \cos(\pi f_1 p_1), \ldots, \sin(\pi f_B p_d), \cos(\pi f_B p_d)\right]
\]
dove $f_1, \ldots, f_B$ sono le frequenze (tipicamente $f_b = 2^{b-1}$ per $b = 1, \ldots, B$, con $B$ il numero di bande).

Dimensione del Fourier PE: $2 \times B \times d$ (2 per seno/coseno, $B$ bande, $d$ dimensioni spaziali).

Esempio per un pixel $(x, y)$ in un'immagine con $B = 6$ bande:
\begin{itemize}
    \item $d = 2$ (coordinate $x$ e $y$)
    \item Dimensione PE = $2 \times 6 \times 2 = 24$
    \item Frequenze: $f \in \{1, 2, 4, 8, 16, 32\}$
    \item PE = $[\sin(\pi \cdot 1 \cdot x), \cos(\pi \cdot 1 \cdot x), \sin(\pi \cdot 2 \cdot x), \ldots, \cos(\pi \cdot 32 \cdot y)]$
\end{itemize}
\end{formulabox}

\subsection{Dettagli su StarCraft II}

La sostituzione del Transformer di entit\`a in AlphaStar (Vinyals et al., 2019) procede come segue:

\begin{enumerate}
    \item Le \textbf{feature delle entit\`a} (tipo di unit\`a, posizione $(x, y)$, punti vita, stato, alleanza, etc.) vengono proiettate in uno spazio embedding
    \item Queste feature servono sia come \textbf{input} per l'encoder del Perceiver IO, sia come \textbf{query di output} per il decoder
    \item L'encoder comprime le feature delle entit\`a nei latenti
    \item Il processor elabora le relazioni tra i latenti (corrispondenti a relazioni strategiche tra gruppi di entit\`a)
    \item Il decoder restituisce un embedding aggiornato per ciascuna entit\`a
    \item Un embedding globale (media sulle entit\`a) viene usato per le decisioni strategiche di alto livello
    \item Il numero di entit\`a pu\`o variare (da poche decine a 512), ma il bottleneck latente ha dimensione fissa --- questo \`e particolarmente vantaggioso in un gioco dove il numero di unit\`a cambia continuamente
\end{enumerate}

\subsection{Perch\'e LAMB invece di Adam?}

Il paper usa l'ottimizzatore \textbf{LAMB} (You et al., 2020) per la maggior parte degli esperimenti. LAMB \`e una variante di Adam progettata per il training con \textbf{batch grandi}:

\begin{itemize}
    \item \textbf{Adam} normalizza i gradienti per parametro (adattivo), ma non tiene conto della scala dei pesi
    \item \textbf{LAMB} aggiunge una normalizzazione basata sul rapporto tra la norma dei pesi e la norma dell'aggiornamento (``layer-wise adaptive rate'')
    \item Questo permette di scalare il batch size senza degradare la convergenza
    \item Per il Perceiver IO con batch 1024 su ImageNet o batch 64 per GLUE, LAMB \`e pi\`u stabile di Adam
    \item L'eccezione \`e l'optical flow (batch 8) dove Adam funziona meglio
\end{itemize}

\subsection{AudioSet: dettagli aggiuntivi}

Per la classificazione audio su AudioSet (Gemmeke et al., 2017):
\begin{itemize}
    \item L'input \`e uno spettrogramma mel con 128 bande di frequenza
    \item I segmenti audio hanno durata di 10 secondi
    \item Il modello riceve in input $128 \times T$ valori, dove $T$ \`e il numero di frame temporali
    \item Si usa Fourier PE 2D (una dimensione per frequenza, una per tempo)
    \item La classificazione \`e multi-label (un clip pu\`o avere pi\`u classi)
    \item La metrica \`e mAP (mean Average Precision)
    \item Il Perceiver IO raggiunge mAP = 0.421, competitivo con modelli audio-specifici
\end{itemize}




\newpage
\part*{Appendici --- Fondamenti di Deep Learning}
\addcontentsline{toc}{part}{Appendici --- Fondamenti di Deep Learning}
\appendix

\section{Il Perceptrone} \label{app:perceptron}

Il \textbf{perceptrone} \`e il mattone fondamentale di tutte le reti neurali. Proposto da Frank Rosenblatt nel 1958, \`e il modello pi\`u semplice di neurone artificiale: riceve degli input numerici, li combina linearmente e produce un output binario tramite una funzione di attivazione a gradino. Comprendere il perceptrone \`e essenziale per capire come funzionano le reti pi\`u complesse che costituiscono il Perceiver (\hyperref[sec:architettura]{Sezione~3}).

\subsection{Struttura del perceptrone}

Un perceptrone riceve $n$ input $x_1, x_2, \ldots, x_n$, ciascuno moltiplicato per un \textbf{peso} $w_1, w_2, \ldots, w_n$. I prodotti vengono sommati insieme a un termine di \textbf{bias} $b$, e il risultato passa attraverso una \textbf{funzione di attivazione} $f$ (tipicamente la funzione a gradino di Heaviside):

\begin{formulabox}[Formula del perceptrone]
\[
    y = f\!\left(\sum_{i=1}^{n} w_i \, x_i + b\right)
\]
dove la funzione di attivazione a gradino \`e:
\[
    f(z) = \begin{cases} 1 & \text{se } z \geq 0 \\ 0 & \text{se } z < 0 \end{cases}
\]
\end{formulabox}

In parole semplici: il perceptrone calcola una \textbf{somma ponderata} degli input (ogni input contribuisce in proporzione al suo peso), aggiunge il bias, e poi ``decide'' se attivarsi (output 1) o no (output 0). I pesi determinano l'\textbf{importanza} di ciascun input: un peso grande e positivo significa che quell'input spinge fortemente verso l'attivazione; un peso negativo significa che quell'input ``inibisce'' l'attivazione.

Il \textbf{bias} $b$ \`e come una \textbf{soglia}: sposta il punto di decisione. Un bias positivo rende il neurone pi\`u facile da attivare; un bias negativo rende pi\`u difficile l'attivazione.

\bigskip

% TikZ: Perceptron diagram
\begin{figure}[H]
\centering
\begin{tikzpicture}[
    inputnode/.style={circle, draw, fill=blue!10, minimum size=0.9cm, font=\small},
    neuronnode/.style={circle, draw, fill=orange!20, minimum size=1.4cm, font=\small},
    outputnode/.style={circle, draw, fill=green!15, minimum size=0.9cm, font=\small},
    arrowstyle/.style={-{Stealth[length=2.5mm]}, thick},
    >=Stealth
]

% Input nodes
\node[inputnode] (x1) at (0, 2.5) {$x_1$};
\node[inputnode] (x2) at (0, 1) {$x_2$};
\node[inputnode] (x3) at (0, -0.5) {$x_3$};
\node[inputnode] (b) at (0, -2) {$1$};
\node[font=\footnotesize, left=0.1cm of b] {bias};

% Neuron
\node[neuronnode] (neuron) at (4, 0.75) {$\displaystyle\sum$};

% Activation
\node[neuronnode, fill=red!15] (act) at (7, 0.75) {$f$};

% Output
\node[outputnode] (out) at (9.5, 0.75) {$y$};

% Arrows with weights
\draw[arrowstyle] (x1) -- (neuron) node[midway, above, font=\footnotesize] {$w_1$};
\draw[arrowstyle] (x2) -- (neuron) node[midway, above, font=\footnotesize] {$w_2$};
\draw[arrowstyle] (x3) -- (neuron) node[midway, below, font=\footnotesize] {$w_3$};
\draw[arrowstyle] (b) -- (neuron) node[midway, below, font=\footnotesize] {$b$};

\draw[arrowstyle] (neuron) -- (act) node[midway, above, font=\footnotesize] {$z = \sum w_i x_i + b$};
\draw[arrowstyle] (act) -- (out) node[midway, above, font=\footnotesize] {$f(z)$};

% Labels
\node[font=\small\bfseries, blue!60!black] at (0, 3.5) {Input};
\node[font=\small\bfseries, orange!70!black] at (4, 2.3) {Somma ponderata};
\node[font=\small\bfseries, red!60!black] at (7, 2.3) {Attivazione};
\node[font=\small\bfseries, green!50!black] at (9.5, 1.7) {Output};

\end{tikzpicture}
\caption{Struttura del perceptrone. Gli input $x_i$ vengono moltiplicati per i pesi $w_i$, sommati con il bias $b$, e passati attraverso la funzione di attivazione $f$.}
\label{fig:perceptron}
\end{figure}

\subsection{Esempio pratico: decidere se indossare una giacca}

\begin{examplebox}[Esempio: giacca s\`i o no?]
Vogliamo decidere se indossare una giacca basandoci su due fattori:
\begin{itemize}
    \item $x_1 = T$ = temperatura (normalizzata: 0 = freddo, 1 = caldo). Oggi: $T = 0.5$ (mite).
    \item $x_2 = P$ = piove? (0 = no, 1 = s\`i). Oggi: $P = 0$ (non piove).
\end{itemize}

I pesi e il bias sono:
\begin{itemize}
    \item $w_1 = -0.7$ (temperatura alta $\to$ meno bisogno di giacca, peso negativo)
    \item $w_2 = 1.0$ (pioggia $\to$ pi\`u bisogno di giacca, peso positivo)
    \item $b = 0.1$ (leggera tendenza a indossare la giacca)
\end{itemize}

\textbf{Calcolo}:
\begin{align*}
    z &= w_1 \cdot x_1 + w_2 \cdot x_2 + b \\
      &= (-0.7) \cdot 0.5 + 1.0 \cdot 0 + 0.1 \\
      &= -0.35 + 0 + 0.1 = -0.25
\end{align*}

Poich\'e $z = -0.25 < 0$, la funzione a gradino restituisce $y = 0$: \textbf{non indossare la giacca}.

Se invece piovesse ($P = 1$): $z = -0.35 + 1.0 + 0.1 = 0.75 > 0 \implies y = 1$: \textbf{indossa la giacca}.
\end{examplebox}


\subsection{Training del perceptrone}

L'addestramento del perceptrone avviene con la \textbf{regola di aggiornamento del perceptrone}: dopo ogni esempio, se il modello sbaglia, i pesi vengono corretti nella direzione giusta.

\begin{formulabox}[Regola di aggiornamento del perceptrone]
\[
    w_i^{\text{new}} = w_i^{\text{old}} + \lambda \cdot (\text{target} - \text{output}) \cdot x_i
\]
\[
    b^{\text{new}} = b^{\text{old}} + \lambda \cdot (\text{target} - \text{output})
\]
dove $\lambda$ \`e il \textbf{learning rate} (tipicamente un piccolo valore come 0.1 o 0.01).
\end{formulabox}

Il \textbf{learning rate} $\lambda$ controlla la velocit\`a di apprendimento: troppo grande porta a oscillazioni, troppo piccolo rende l'apprendimento lento. Il \textbf{teorema di convergenza del perceptrone} garantisce che, se i dati sono linearmente separabili, l'algoritmo converge in un numero finito di passi.

\begin{examplebox}[Esempio numerico di training]
Supponiamo $x_1 = 1$, $x_2 = 0$, target = 1, output = 0, $\lambda = 0.1$:
\begin{align*}
    \text{errore} &= \text{target} - \text{output} = 1 - 0 = 1 \\
    w_1^{\text{new}} &= w_1^{\text{old}} + 0.1 \cdot 1 \cdot 1 = w_1^{\text{old}} + 0.1 \\
    w_2^{\text{new}} &= w_2^{\text{old}} + 0.1 \cdot 1 \cdot 0 = w_2^{\text{old}} \quad \text{(invariato, perch\'e } x_2 = 0\text{)}
\end{align*}
Il peso $w_1$ aumenta perch\'e l'input $x_1$ avrebbe dovuto contribuire all'attivazione ma non lo ha fatto sufficientemente.
\end{examplebox}


\subsection{Limiti del perceptrone: il problema XOR}

Il perceptrone pu\`o risolvere solo problemi \textbf{linearmente separabili}, cio\`e problemi dove le due classi possono essere separate da un iperpiano (una retta in 2D). L'esempio classico di un problema NON linearmente separabile \`e lo \textbf{XOR}:

\begin{center}
\begin{tabular}{cc|c}
\toprule
$x_1$ & $x_2$ & XOR \\
\midrule
0 & 0 & 0 \\
0 & 1 & 1 \\
1 & 0 & 1 \\
1 & 1 & 0 \\
\bottomrule
\end{tabular}
\end{center}

Non esiste nessuna retta nel piano $(x_1, x_2)$ che separi i punti con output 0 dai punti con output 1. Questo limite fu evidenziato da Minsky e Papert (1969) e fu superato solo con l'introduzione dei \textbf{perceptroni multi-strato} (reti feedforward), descritti nell'Appendice~\ref{app:feedforward}.

% TikZ: XOR problem visualization
\begin{figure}[H]
\centering
\begin{tikzpicture}[scale=2]
    % Axes
    \draw[-{Stealth}, thick] (-0.3, 0) -- (1.5, 0) node[right] {$x_1$};
    \draw[-{Stealth}, thick] (0, -0.3) -- (0, 1.5) node[above] {$x_2$};

    % Points XOR = 0 (red circles)
    \node[circle, fill=red!70, draw=red!90!black, minimum size=0.35cm, label={[font=\footnotesize]below left:{0}}] at (0, 0) {};
    \node[circle, fill=red!70, draw=red!90!black, minimum size=0.35cm, label={[font=\footnotesize]above right:{0}}] at (1, 1) {};

    % Points XOR = 1 (blue squares)
    \node[draw=blue!90!black, fill=blue!50, minimum size=0.3cm, label={[font=\footnotesize]above left:{1}}] at (0, 1) {};
    \node[draw=blue!90!black, fill=blue!50, minimum size=0.3cm, label={[font=\footnotesize]below right:{1}}] at (1, 0) {};

    % No separating line annotation
    \node[font=\footnotesize\itshape, text=gray] at (0.5, -0.5) {Nessuna retta pu\`o separare};
    \node[font=\footnotesize\itshape, text=gray] at (0.5, -0.7) {le classi 0 e 1};
\end{tikzpicture}
\caption{Il problema XOR non \`e linearmente separabile: non esiste una retta che separi i punti rossi (XOR=0) dai punti blu (XOR=1).}
\label{fig:xor}
\end{figure}



\newpage
%% ============================================================
%%  APPENDICE B --- RETI FEEDFORWARD
%% ============================================================
\section{Reti Neurali Feedforward (MLP)} \label{app:feedforward}

Le \textbf{reti neurali feedforward} (o Multi-Layer Perceptron, MLP) superano i limiti del singolo perceptrone (Appendice~\ref{app:perceptron}) aggiungendo uno o pi\`u \textbf{strati nascosti} (hidden layers) tra input e output. Con almeno un hidden layer e funzioni di attivazione non lineari, un MLP \`e un \textbf{approssimatore universale}: pu\`o approssimare qualsiasi funzione continua (teorema di Cybenko, 1989). Il Feed-Forward Network del Perceiver (Sezione~\ref{sec:latent_transformer}) \`e esattamente un MLP a due strati.

\subsection{Architettura}

Una rete feedforward \`e organizzata in \textbf{strati} (layer):
\begin{enumerate}
    \item \textbf{Input Layer}: riceve i dati ($n$ neuroni per $n$ feature).
    \item \textbf{Hidden Layers}: uno o pi\`u strati intermedi. Ogni neurone calcola una somma ponderata degli input dallo strato precedente, seguita da una funzione di attivazione non lineare.
    \item \textbf{Output Layer}: produce il risultato finale (un neurone per regressione, $K$ neuroni per classificazione a $K$ classi).
\end{enumerate}

Il termine ``feedforward'' indica che l'informazione fluisce in una sola direzione: dall'input all'output, senza cicli o connessioni all'indietro.

% TikZ: Feedforward network
\begin{figure}[H]
\centering
\begin{tikzpicture}[
    neuron/.style={circle, draw, minimum size=0.7cm, font=\footnotesize},
    input/.style={neuron, fill=blue!15},
    hidden/.style={neuron, fill=orange!20},
    output/.style={neuron, fill=green!15},
    arrowstyle/.style={-{Stealth[length=2mm]}, thick, gray!70},
    >=Stealth
]

% Input layer
\foreach \i in {1,2,3} {
    \node[input] (I\i) at (0, {3 - \i * 1.2}) {$x_{\i}$};
}

% Hidden layer 1
\foreach \j in {1,2,3,4} {
    \node[hidden] (H1\j) at (3, {3.8 - \j * 1.1}) {$h^1_{\j}$};
}

% Hidden layer 2
\foreach \k in {1,2,3} {
    \node[hidden] (H2\k) at (6, {3 - \k * 1.2}) {$h^2_{\k}$};
}

% Output layer
\foreach \l in {1,2} {
    \node[output] (O\l) at (9, {2.2 - \l * 1.2}) {$y_{\l}$};
}

% Connections input -> hidden 1
\foreach \i in {1,2,3} {
    \foreach \j in {1,2,3,4} {
        \draw[arrowstyle] (I\i) -- (H1\j);
    }
}

% Connections hidden 1 -> hidden 2
\foreach \j in {1,2,3,4} {
    \foreach \k in {1,2,3} {
        \draw[arrowstyle] (H1\j) -- (H2\k);
    }
}

% Connections hidden 2 -> output
\foreach \k in {1,2,3} {
    \foreach \l in {1,2} {
        \draw[arrowstyle] (H2\k) -- (O\l);
    }
}

% Labels
\node[font=\small\bfseries, blue!60!black] at (0, 3.5) {Input};
\node[font=\small\bfseries, orange!70!black] at (3, 3.5) {Hidden 1};
\node[font=\small\bfseries, orange!70!black] at (6, 3.5) {Hidden 2};
\node[font=\small\bfseries, green!50!black] at (9, 3.5) {Output};

\end{tikzpicture}
\caption{Rete neurale feedforward con 3 input, 2 hidden layer (4 e 3 neuroni) e 2 output. Ogni freccia rappresenta un peso apprendibile.}
\label{fig:feedforward_net}
\end{figure}


\subsection{Forward propagation}

Per ogni neurone $j$ nello strato $\ell$, il calcolo procede in due fasi:

\begin{formulabox}[Forward propagation]
\textbf{1. Input netto} (somma ponderata):
\[
    z_j^{(\ell)} = \sum_{i} w_{ji}^{(\ell)} \, a_i^{(\ell-1)} + b_j^{(\ell)}
\]
\textbf{2. Attivazione}:
\[
    a_j^{(\ell)} = \sigma\!\left(z_j^{(\ell)}\right)
\]
dove $\sigma$ \`e una funzione di attivazione non lineare (ReLU, sigmoid, tanh, GELU\ldots), $a_i^{(\ell-1)}$ sono le attivazioni dello strato precedente, $w_{ji}^{(\ell)}$ sono i pesi, e $b_j^{(\ell)}$ \`e il bias.

In notazione matriciale per l'intero strato:
\[
    \mathbf{a}^{(\ell)} = \sigma\!\left(\mathbf{W}^{(\ell)} \mathbf{a}^{(\ell-1)} + \mathbf{b}^{(\ell)}\right)
\]
\end{formulabox}


\subsection{Loss function: Mean Squared Error}

Per la regressione, la funzione di costo pi\`u comune \`e il \textbf{Mean Squared Error} (MSE):

\begin{formulabox}[Mean Squared Error]
\[
    \mathcal{L}_{\text{MSE}} = \frac{1}{n} \sum_{i=1}^{n} \left(y_i^{\text{pred}} - y_i^{\text{true}}\right)^2
\]
dove $n$ \`e il numero di campioni. Per la classificazione multi-classe si usa la cross-entropy (come nel Perceiver, \hyperref[sec:forward]{Sezione~5}).
\end{formulabox}


\subsection{Backpropagation}

La \textbf{backpropagation} (Rumelhart, Hinton \& Williams, 1986) \`e l'algoritmo per calcolare i gradienti della loss rispetto a tutti i pesi della rete, applicando la \textbf{chain rule} del calcolo differenziale strato per strato dall'output verso l'input.

\begin{formulabox}[Backpropagation: passi chiave]

\textbf{Passo 1 --- Gradiente dell'errore rispetto all'output}:
\[
    \delta_j^{(L)} = \frac{\partial \mathcal{L}}{\partial a_j^{(L)}}
\]
Per MSE: $\delta_j^{(L)} = \frac{2}{n}(a_j^{(L)} - y_j^{\text{true}})$.

\textbf{Passo 2 --- Gradiente rispetto all'input netto} (chain rule con la derivata dell'attivazione):
\[
    \delta_j^{(\ell)} = \frac{\partial \mathcal{L}}{\partial z_j^{(\ell)}} = \frac{\partial \mathcal{L}}{\partial a_j^{(\ell)}} \cdot \sigma'\!\left(z_j^{(\ell)}\right)
\]

\textbf{Passo 3 --- Gradiente rispetto ai pesi}:
\[
    \frac{\partial \mathcal{L}}{\partial w_{ji}^{(\ell)}} = \delta_j^{(\ell)} \cdot a_i^{(\ell-1)}
\]

\textbf{Passo 4 --- Gradiente rispetto al bias}:
\[
    \frac{\partial \mathcal{L}}{\partial b_j^{(\ell)}} = \delta_j^{(\ell)}
\]

\textbf{Passo 5 --- Propagazione all'indietro} (per lo strato precedente):
\[
    \frac{\partial \mathcal{L}}{\partial a_i^{(\ell-1)}} = \sum_j w_{ji}^{(\ell)} \cdot \delta_j^{(\ell)}
\]

\textbf{Passo 6 --- Aggiornamento pesi con gradient descent}:
\[
    w_{ji}^{(\ell)} \leftarrow w_{ji}^{(\ell)} - \eta \cdot \frac{\partial \mathcal{L}}{\partial w_{ji}^{(\ell)}}
\]
dove $\eta$ \`e il learning rate.
\end{formulabox}

\begin{notabox}[In parole semplici: la backpropagation come una catena di responsabilit\`a]
Immagina una catena di montaggio dove ogni operaio (neurone) modifica il pezzo. Se il prodotto finale \`e difettoso (loss alta), il controllo qualit\`a (backpropagation) risale la catena chiedendo a ogni operaio: ``Quanto sei responsabile di questo difetto?''. L'operaio che ha contribuito di pi\`u al difetto riceve la correzione pi\`u grande (gradiente pi\`u grande). La chain rule \`e il meccanismo che permette di ``scomporre'' la responsabilit\`a strato per strato.
\end{notabox}

\begin{examplebox}[Esempio numerico: backpropagation su rete semplice]
Rete con 1 input $x = 2$, 1 hidden neuron con ReLU, 1 output con identit\`a. Pesi: $w_1 = 0.5$, $w_2 = -0.3$, bias = 0. Target: $y_{\text{true}} = 1$.

\textbf{Forward}: $z_1 = 0.5 \times 2 = 1.0$, $a_1 = \text{ReLU}(1.0) = 1.0$, output: $\hat{y} = -0.3 \times 1.0 = -0.3$.

\textbf{Loss MSE}: $\mathcal{L} = (-0.3 - 1)^2 = 1.69$.

\textbf{Backward}:
\begin{itemize}
    \item $\frac{\partial \mathcal{L}}{\partial \hat{y}} = 2(-0.3 - 1) = -2.6$
    \item $\frac{\partial \mathcal{L}}{\partial w_2} = -2.6 \times a_1 = -2.6 \times 1.0 = -2.6$
    \item $\frac{\partial \mathcal{L}}{\partial a_1} = -2.6 \times w_2 = -2.6 \times (-0.3) = 0.78$
    \item $\frac{\partial \mathcal{L}}{\partial z_1} = 0.78 \times \text{ReLU}'(1.0) = 0.78 \times 1 = 0.78$
    \item $\frac{\partial \mathcal{L}}{\partial w_1} = 0.78 \times x = 0.78 \times 2 = 1.56$
\end{itemize}

\textbf{Aggiornamento} ($\eta = 0.01$): $w_2 \leftarrow -0.3 - 0.01 \times (-2.6) = -0.274$, $w_1 \leftarrow 0.5 - 0.01 \times 1.56 = 0.4844$.
\end{examplebox}



\newpage
%% ============================================================
%%  APPENDICE C --- RNN
%% ============================================================
\section{Reti Neurali Ricorrenti (RNN)} \label{app:rnn}

Le \textbf{reti neurali ricorrenti} (RNN) sono architetture progettate per elaborare \textbf{dati sequenziali}: testo (sequenza di parole), audio (sequenza di campioni), serie temporali (sequenza di misurazioni). A differenza delle reti feedforward (Appendice~\ref{app:feedforward}), che processano un input alla volta senza memoria, le RNN mantengono uno \textbf{stato nascosto} (hidden state) che accumula informazione dai passi precedenti.

Le RNN sono un precursore fondamentale del Transformer e del Perceiver: il Perceiver eredita l'idea dello stato interno (l'array latente $\mathbf{z}$, \hyperref[sec:panoramica]{Sezione~3.1}) e dell'elaborazione iterativa con weight sharing (Sezione~\ref{sec:iterative}), ma elimina la sequenzialit\`a grazie all'attenzione.

\begin{figure}[H]
\centering
\includegraphics[width=0.85\textwidth]{paper_figures/rnn_architecture_original.png}
\caption{Architettura di una RNN (da Olah, 2015). A sinistra: la cella con il loop ricorrente. A destra: la stessa rete ``srotolata'' nel tempo --- ogni cella condivide gli stessi pesi e passa lo stato nascosto $\mathbf{h}_t$ al passo successivo.}
\label{fig:rnn_original_arch}
\end{figure}

Come si vede dalla Figura~\ref{fig:rnn_original_arch}, la struttura della RNN \`e sorprendentemente semplice: una singola cella $A$ riceve ad ogni passo l'input $\mathbf{x}_t$ e lo stato nascosto precedente $\mathbf{h}_{t-1}$, e produce un nuovo stato $\mathbf{h}_t$ e un output $\mathbf{y}_t$. Lo ``srotolamento'' a destra mostra che la stessa cella (con gli stessi pesi) viene riutilizzata ad ogni passo temporale.

\subsection{Il loop nello stato nascosto}

L'idea fondamentale delle RNN \`e un \textbf{loop}: lo stato nascosto $\mathbf{h}_t$ al tempo $t$ dipende sia dall'input corrente $\mathbf{x}_t$ sia dallo stato nascosto precedente $\mathbf{h}_{t-1}$. Questa dipendenza crea una ``memoria'' delle informazioni passate.

Le RNN utilizzano tre matrici di pesi:
\begin{itemize}
    \item $\mathbf{W} \in \mathbb{R}^{d_h \times d_h}$: pesi per lo stato nascosto (da $\mathbf{h}_{t-1}$ a $\mathbf{h}_t$)
    \item $\mathbf{U} \in \mathbb{R}^{d_h \times d_x}$: pesi per l'input (da $\mathbf{x}_t$ a $\mathbf{h}_t$)
    \item $\mathbf{V} \in \mathbb{R}^{d_y \times d_h}$: pesi per l'output (da $\mathbf{h}_t$ a $\mathbf{y}_t$)
\end{itemize}

\begin{formulabox}[Equazioni fondamentali della RNN]
\textbf{Stato nascosto}:
\[
    \mathbf{h}_t = \tanh\!\left(\mathbf{W} \cdot \mathbf{h}_{t-1} + \mathbf{U} \cdot \mathbf{x}_t + \mathbf{b}_h\right)
\]
\textbf{Output}:
\[
    \mathbf{y}_t = \mathbf{V} \cdot \mathbf{h}_t + \mathbf{b}_y
\]
I pesi $\mathbf{W}$, $\mathbf{U}$, $\mathbf{V}$ sono \textbf{condivisi} in tutti i passi temporali (weight sharing, stessa idea del Perceiver nella Sezione~\ref{sec:iterative}).
\end{formulabox}

\subsection{Srotolamento nel tempo (Unrolling)}

Per addestrare una RNN, si ``srotola'' il loop nel tempo, ottenendo una catena di copie della stessa rete, una per ogni passo temporale. Questo \`e concettualmente identico a come il Perceiver ``srotola'' le sue iterazioni di cross-attention.

% TikZ: RNN unrolled
\begin{figure}[H]
\centering
\begin{tikzpicture}[
    cell/.style={draw, rounded corners=3pt, fill=orange!15, minimum width=1.6cm, minimum height=1.0cm, font=\small\bfseries},
    io/.style={draw, rounded corners=2pt, minimum width=1.0cm, minimum height=0.6cm, font=\small},
    arrowstyle/.style={-{Stealth[length=2.5mm]}, thick},
    >=Stealth
]

% RNN cells
% Cell t-1
\node[cell] (h0) at (0, 0) {RNN};
\node[io, fill=blue!10] (x0) at (0, -1.8) {$\mathbf{x}_{t-1}$};
\node[io, fill=green!10] (y0) at (0, 1.8) {$\mathbf{y}_{t-1}$};
\draw[arrowstyle] (x0) -- (h0);
\draw[arrowstyle] (h0) -- (y0);
% Cell t
\node[cell] (h1) at (3.5, 0) {RNN};
\node[io, fill=blue!10] (x1) at (3.5, -1.8) {$\mathbf{x}_{t}$};
\node[io, fill=green!10] (y1) at (3.5, 1.8) {$\mathbf{y}_{t}$};
\draw[arrowstyle] (x1) -- (h1);
\draw[arrowstyle] (h1) -- (y1);
% Cell t+1
\node[cell] (h2) at (7, 0) {RNN};
\node[io, fill=blue!10] (x2) at (7, -1.8) {$\mathbf{x}_{t+1}$};
\node[io, fill=green!10] (y2) at (7, 1.8) {$\mathbf{y}_{t+1}$};
\draw[arrowstyle] (x2) -- (h2);
\draw[arrowstyle] (h2) -- (y2);

% Horizontal arrows (hidden state flow)
\draw[arrowstyle, red!60!black] (-1.8, 0) -- (h0.west) node[midway, above, font=\footnotesize] {$\mathbf{h}_{t-2}$};
\draw[arrowstyle, red!60!black] (h0.east) -- (h1.west) node[midway, above, font=\footnotesize] {$\mathbf{h}_{t-1}$};
\draw[arrowstyle, red!60!black] (h1.east) -- (h2.west) node[midway, above, font=\footnotesize] {$\mathbf{h}_{t}$};
\draw[arrowstyle, red!60!black] (h2.east) -- (8.8, 0) node[midway, above, font=\footnotesize] {$\mathbf{h}_{t+1}$};

% Weight sharing annotation
\draw[decorate, decoration={brace, amplitude=5pt}, thick, blue!60!black]
    (7.8, -2.3) -- (-0.8, -2.3) node[midway, below=6pt, font=\small, blue!60!black] {Stessi pesi $\mathbf{W}$, $\mathbf{U}$, $\mathbf{V}$ in tutti i passi};

\end{tikzpicture}
\caption{RNN srotolata nel tempo: ogni cella condivide gli stessi pesi. Lo stato nascosto $\mathbf{h}_t$ (frecce rosse) fluisce da un passo al successivo, accumulando informazione.}
\label{fig:rnn_unrolled}
\end{figure}


\subsection{Esempio numerico: forward pass di una RNN}

\begin{examplebox}[Forward pass RNN su una sequenza di 3 elementi]
Prendiamo una RNN minuscola con $d_h = 2$ (stato nascosto) e $d_x = 1$ (un input scalare).

\textbf{Pesi} (condivisi in tutti i passi):
\[
\mathbf{W} = \begin{pmatrix} 0.5 & 0.1 \\ -0.2 & 0.4 \end{pmatrix}, \quad
\mathbf{U} = \begin{pmatrix} 0.6 \\ 0.3 \end{pmatrix}, \quad
\mathbf{b}_h = \begin{pmatrix} 0 \\ 0 \end{pmatrix}
\]

\textbf{Sequenza di input}: $x_1 = 1.0$, $x_2 = 0.5$, $x_3 = -0.3$. Stato iniziale: $\mathbf{h}_0 = \begin{pmatrix} 0 \\ 0 \end{pmatrix}$.

\textbf{Passo $t = 1$} ($x_1 = 1.0$):
\begin{align*}
    \mathbf{W} \cdot \mathbf{h}_0 + \mathbf{U} \cdot x_1 &= \begin{pmatrix} 0 \\ 0 \end{pmatrix} + \begin{pmatrix} 0.6 \\ 0.3 \end{pmatrix} = \begin{pmatrix} 0.6 \\ 0.3 \end{pmatrix} \\
    \mathbf{h}_1 &= \tanh\begin{pmatrix} 0.6 \\ 0.3 \end{pmatrix} = \begin{pmatrix} 0.537 \\ 0.291 \end{pmatrix}
\end{align*}

\textbf{Passo $t = 2$} ($x_2 = 0.5$): ora $\mathbf{h}_1$ entra nel calcolo --- la RNN ``ricorda'' il passo precedente:
\begin{align*}
    \mathbf{W} \cdot \mathbf{h}_1 + \mathbf{U} \cdot x_2 &= \begin{pmatrix} 0.5 \cdot 0.537 + 0.1 \cdot 0.291 \\ -0.2 \cdot 0.537 + 0.4 \cdot 0.291 \end{pmatrix} + \begin{pmatrix} 0.3 \\ 0.15 \end{pmatrix} = \begin{pmatrix} 0.598 \\ 0.159 \end{pmatrix} \\
    \mathbf{h}_2 &= \tanh\begin{pmatrix} 0.598 \\ 0.159 \end{pmatrix} = \begin{pmatrix} 0.535 \\ 0.158 \end{pmatrix}
\end{align*}

\textbf{Passo $t = 3$} ($x_3 = -0.3$):
\begin{align*}
    \mathbf{W} \cdot \mathbf{h}_2 + \mathbf{U} \cdot x_3 &= \begin{pmatrix} 0.5 \cdot 0.535 + 0.1 \cdot 0.158 \\ -0.2 \cdot 0.535 + 0.4 \cdot 0.158 \end{pmatrix} + \begin{pmatrix} -0.18 \\ -0.09 \end{pmatrix} = \begin{pmatrix} 0.104 \\ -0.134 \end{pmatrix} \\
    \mathbf{h}_3 &= \tanh\begin{pmatrix} 0.104 \\ -0.134 \end{pmatrix} = \begin{pmatrix} 0.104 \\ -0.133 \end{pmatrix}
\end{align*}

\textbf{Osservazione chiave}: lo stato $\mathbf{h}_3$ dipende da \emph{tutti e tre} gli input $x_1, x_2, x_3$ --- la RNN ha accumulato informazione passo dopo passo. Ma nota come l'influenza di $x_1$ su $\mathbf{h}_3$ \`e ormai molto attenuata: questo \`e il vanishing gradient in azione, anche dopo soli 3 passi.
\end{examplebox}


\subsection{Vanishing gradient problem}

Il problema pi\`u grave delle RNN \`e il \textbf{vanishing gradient}: quando la sequenza \`e lunga, i gradienti che fluiscono all'indietro attraverso molti passi temporali tendono a \textbf{zero}, rendendo impossibile apprendere dipendenze a lungo raggio.

\begin{formulabox}[Vanishing gradient nelle RNN]
Il gradiente della loss rispetto allo stato nascosto al tempo $t$ richiede di moltiplicare ripetutamente la matrice $\mathbf{W}$:
\[
    \frac{\partial \mathbf{h}_T}{\partial \mathbf{h}_t} = \prod_{k=t+1}^{T} \frac{\partial \mathbf{h}_k}{\partial \mathbf{h}_{k-1}} = \prod_{k=t+1}^{T} \mathbf{W}^\top \cdot \text{diag}\!\left(\tanh'(\mathbf{z}_k)\right)
\]
Se gli autovalori di $\mathbf{W}$ sono $< 1$: il prodotto tende a $\mathbf{0}$ \textbf{(vanishing gradient)}.

Se gli autovalori di $\mathbf{W}$ sono $> 1$: il prodotto esplode \textbf{(exploding gradient)}.
\end{formulabox}

L'\textbf{exploding gradient} si pu\`o mitigare con il \textbf{gradient clipping}: se la norma del gradiente supera una soglia $\tau$, si riscala:
\[
    \mathbf{g} \leftarrow \frac{\tau}{\|\mathbf{g}\|} \cdot \mathbf{g} \quad \text{se } \|\mathbf{g}\| > \tau
\]

Il vanishing gradient \`e invece pi\`u insidioso e ha motivato l'introduzione di LSTM (Appendice~\ref{app:lstm}) e GRU (Appendice~\ref{app:gru}), e successivamente del Transformer (Appendice~\ref{app:transformer}), che elimina completamente la ricorrenza.



\newpage
%% ============================================================
%%  APPENDICE D --- LSTM
%% ============================================================
\section{LSTM (Long Short-Term Memory)} \label{app:lstm}

Le \textbf{LSTM} (Hochreiter \& Schmidhuber, 1997) sono state progettate per risolvere il problema del vanishing gradient (Appendice~\ref{app:rnn}) introducendo un meccanismo di \textbf{porte} (gates) che regola il flusso dell'informazione. L'idea chiave \`e una \textbf{cella di memoria} $\mathbf{C}_t$ che pu\`o mantenere informazione per lunghi periodi, come un ``nastro trasportatore'' che attraversa tutta la sequenza.

\begin{figure}[H]
\centering
\includegraphics[width=0.85\textwidth]{paper_figures/lstm_architecture_original.png}
\caption{Architettura della cella LSTM srotolata nel tempo (da Olah, 2015). Si vede la linea orizzontale superiore (stato della cella $\mathbf{C}_t$) che scorre come un ``nastro trasportatore'', modulata dalle tre porte: forget ($\sigma$), input ($\sigma$ + $\tanh$) e output ($\sigma$). A differenza della RNN semplice, l'informazione pu\`o fluire quasi inalterata attraverso molti passi.}
\label{fig:lstm_original_arch}
\end{figure}

Come si vede dalla Figura~\ref{fig:lstm_original_arch}, la LSTM \`e significativamente pi\`u complessa di una RNN semplice: ha 4 componenti interagenti (3 porte sigmoide e un candidato $\tanh$) che controllano il flusso dell'informazione. La linea orizzontale in alto (lo stato della cella $\mathbf{C}_t$) \`e la chiave: l'informazione pu\`o scorrere lungo questa linea con modifiche minime, risolvendo il problema del vanishing gradient.

\subsection{Architettura della cella LSTM}

La cella LSTM ha 4 componenti fondamentali, ognuna con il suo ruolo:

\begin{formulabox}[Le 4 componenti della LSTM]

\textbf{1. Forget Gate} (``Cosa dimenticare?''):
\[
    \mathbf{f}_t = \sigma\!\left(\mathbf{W}_f \cdot [\mathbf{h}_{t-1}, \mathbf{x}_t] + \mathbf{b}_f\right)
\]
Output: valori tra 0 e 1. Un valore vicino a 0 significa ``dimentica questa informazione''; un valore vicino a 1 significa ``ricorda''.

\textbf{2. Input Gate} (``Cosa aggiornare?''):
\[
    \mathbf{i}_t = \sigma\!\left(\mathbf{W}_i \cdot [\mathbf{h}_{t-1}, \mathbf{x}_t] + \mathbf{b}_i\right)
\]
Decide quali valori della cella verranno aggiornati con nuova informazione.

\textbf{3. Aggiornamento dello stato della cella}:
\[
    \tilde{\mathbf{C}}_t = \tanh\!\left(\mathbf{W}_C \cdot [\mathbf{h}_{t-1}, \mathbf{x}_t] + \mathbf{b}_C\right) \quad \text{(candidato)}
\]
\[
    \mathbf{C}_t = \mathbf{f}_t \odot \mathbf{C}_{t-1} + \mathbf{i}_t \odot \tilde{\mathbf{C}}_t
\]
Lo stato si aggiorna in due passi: prima si ``dimentica'' (moltiplicando per $\mathbf{f}_t$), poi si ``aggiunge'' nuova informazione (moltiplicando il candidato per $\mathbf{i}_t$).

\textbf{4. Output Gate} (``Cosa produrre in output?''):
\[
    \mathbf{o}_t = \sigma\!\left(\mathbf{W}_o \cdot [\mathbf{h}_{t-1}, \mathbf{x}_t] + \mathbf{b}_o\right)
\]
\[
    \mathbf{h}_t = \mathbf{o}_t \odot \tanh(\mathbf{C}_t)
\]
L'output gate filtra lo stato della cella per produrre lo stato nascosto.
\end{formulabox}

La funzione \textbf{sigmoide} $\sigma(x) = \frac{1}{1 + e^{-x}}$ \`e utilizzata come ``porta'' perch\'e produce valori tra 0 e 1: 0 = porta chiusa (nessuna informazione passa), 1 = porta aperta (tutta l'informazione passa). Valori intermedi modulano parzialmente il flusso.

% TikZ: LSTM cell
\begin{figure}[H]
\centering
\begin{tikzpicture}[
    gate/.style={draw, rounded corners=2pt, fill=#1, minimum width=1.2cm, minimum height=0.7cm, font=\footnotesize\bfseries},
    gate/.default={yellow!20},
    op/.style={circle, draw, fill=gray!10, minimum size=0.5cm, font=\footnotesize},
    arrowstyle/.style={-{Stealth[length=2mm]}, thick},
    line/.style={thick},
    >=Stealth
]

% Cell state line (the highway)
\draw[line, red!60!black] (-1, 3) -- (10, 3);
\node[font=\footnotesize, red!60!black] at (-1.5, 3) {$\mathbf{C}_{t-1}$};
\node[font=\footnotesize, red!60!black] at (10.5, 3) {$\mathbf{C}_t$};

% Hidden state line
\draw[line, blue!60!black] (-1, 0) -- (1, 0);
\node[font=\footnotesize, blue!60!black] at (-1.5, 0) {$\mathbf{h}_{t-1}$};

% Forget gate
\node[gate=red!15] (fg) at (2, 1.5) {$\sigma$};
\node[font=\footnotesize, below=0.1cm of fg] {Forget};
\node[op] (fgmul) at (2, 3) {$\times$};
\draw[arrowstyle] (fg) -- (fgmul);
\draw[line] (1, 0) |- (fg.west);

% Input gate
\node[gate=green!15] (ig) at (4.5, 1.5) {$\sigma$};
\node[font=\footnotesize, below=0.1cm of ig] {Input};
\node[gate=cyan!15] (cand) at (4.5, 0) {$\tanh$};
\node[font=\footnotesize, below=0.1cm of cand] {Candidato};
\node[op] (igmul) at (4.5, 2.3) {$\times$};
\node[op] (add) at (4.5, 3) {$+$};

\draw[arrowstyle] (ig) -- (igmul);
\draw[arrowstyle] (cand) -- (igmul);
\draw[arrowstyle] (igmul) -- (add);

% Output gate
\node[gate=purple!15] (og) at (7.5, 1.5) {$\sigma$};
\node[font=\footnotesize, below=0.1cm of og] {Output};
\node[gate=cyan!15] (otanh) at (7.5, 3) {$\tanh$};
\node[op] (ogmul) at (8.5, 1.5) {$\times$};

\draw[arrowstyle] (og) -- (ogmul);
\draw[arrowstyle] (otanh.east) -| (ogmul.north);

% Output hidden state
\draw[arrowstyle, blue!60!black] (ogmul.east) -- (10, 1.5);
\node[font=\footnotesize, blue!60!black] at (10.5, 1.5) {$\mathbf{h}_t$};

% Input x_t
\node[draw, fill=blue!10, rounded corners, minimum width=1cm, font=\footnotesize] (xt) at (4.5, -1.5) {$\mathbf{x}_t$};
\draw[arrowstyle] (xt) -- (cand);

% Connect input to gates
\draw[line, gray] (1, 0) -- (1, -0.8) -| (fg.south);
\draw[line, gray] (1, -0.8) -| (ig.south);
\draw[line, gray] (1, -0.8) -| (og.south);
\draw[line, gray] (xt.west) -| (1, -0.8);

\end{tikzpicture}
\caption{Cella LSTM. Lo stato della cella $\mathbf{C}_t$ (linea rossa) scorre come un ``nastro trasportatore'', modulato dalle tre porte: forget (rosso), input (verde) e output (viola).}
\label{fig:lstm_cell}
\end{figure}

\begin{examplebox}[Esempio intuitivo: LSTM che legge un testo]
Immagina una LSTM che legge la frase: ``Il gatto, che era nero e aveva gli occhi verdi, \textbf{corse} via''. La cella di memoria deve ricordare ``il gatto'' (soggetto singolare) attraverso tutta la frase subordinata. Il \textbf{forget gate} decide di non dimenticare il soggetto (chiude la porta a 1). L'\textbf{input gate} aggiunge informazioni sugli aggettivi ma non sovrascrive il soggetto. Quando arriva il verbo, l'\textbf{output gate} produce lo stato nascosto che codifica ``soggetto singolare'', permettendo di concordare correttamente il verbo.
\end{examplebox}

\subsection{Esempio numerico: forward pass di una LSTM}

\begin{examplebox}[Forward pass LSTM --- gate per gate]
LSTM con $d_h = 2$, $d_x = 1$. Stato iniziale: $\mathbf{h}_{t-1} = \begin{pmatrix} 0.2 \\ -0.1 \end{pmatrix}$, $\mathbf{C}_{t-1} = \begin{pmatrix} 0.5 \\ 0.3 \end{pmatrix}$, input $x_t = 0.7$.

Vettore concatenato: $[\mathbf{h}_{t-1}, x_t] = \begin{pmatrix} 0.2 & -0.1 & 0.7 \end{pmatrix}$.

\textbf{Passo 1 --- Forget Gate} (``Cosa dimenticare della cella?'').
Con $\mathbf{W}_f = \begin{pmatrix} 0.3 & 0.1 & 0.2 \\ 0.1 & 0.4 & -0.1 \end{pmatrix}$, $\mathbf{b}_f = \begin{pmatrix} 0.5 \\ 0.5 \end{pmatrix}$:
\begin{align*}
    \mathbf{f}_t &= \sigma\!\left(\begin{pmatrix} 0.06 + (-0.01) + 0.14 + 0.5 \\ 0.02 + (-0.04) + (-0.07) + 0.5 \end{pmatrix}\right) = \sigma\!\begin{pmatrix} 0.69 \\ 0.41 \end{pmatrix} = \begin{pmatrix} 0.666 \\ 0.601 \end{pmatrix}
\end{align*}
Valori $\approx 0.6$--$0.7$: la LSTM mantiene buona parte della cella passata.

\textbf{Passo 2 --- Input Gate} (``Cosa aggiungere di nuovo?'').
Con $\mathbf{W}_i = \begin{pmatrix} -0.2 & 0.3 & 0.5 \\ 0.4 & -0.1 & 0.2 \end{pmatrix}$, $\mathbf{b}_i = \begin{pmatrix} 0 \\ 0 \end{pmatrix}$:
\begin{align*}
    \mathbf{i}_t &= \sigma\!\begin{pmatrix} -0.04 + (-0.03) + 0.35 \\ 0.08 + 0.01 + 0.14 \end{pmatrix} = \sigma\!\begin{pmatrix} 0.28 \\ 0.23 \end{pmatrix} = \begin{pmatrix} 0.570 \\ 0.557 \end{pmatrix}
\end{align*}

\textbf{Passo 3 --- Candidato} $\tilde{\mathbf{C}}_t$.
Con $\mathbf{W}_C = \begin{pmatrix} 0.1 & -0.3 & 0.6 \\ 0.2 & 0.5 & -0.2 \end{pmatrix}$, $\mathbf{b}_C = \begin{pmatrix} 0 \\ 0 \end{pmatrix}$:
\begin{align*}
    \tilde{\mathbf{C}}_t &= \tanh\!\begin{pmatrix} 0.02 + 0.03 + 0.42 \\ 0.04 + (-0.05) + (-0.14) \end{pmatrix} = \tanh\!\begin{pmatrix} 0.47 \\ -0.15 \end{pmatrix} = \begin{pmatrix} 0.438 \\ -0.149 \end{pmatrix}
\end{align*}

\textbf{Passo 4 --- Aggiornamento della cella}:
\begin{align*}
    \mathbf{C}_t &= \mathbf{f}_t \odot \mathbf{C}_{t-1} + \mathbf{i}_t \odot \tilde{\mathbf{C}}_t \\
    &= \begin{pmatrix} 0.666 \\ 0.601 \end{pmatrix} \odot \begin{pmatrix} 0.5 \\ 0.3 \end{pmatrix} + \begin{pmatrix} 0.570 \\ 0.557 \end{pmatrix} \odot \begin{pmatrix} 0.438 \\ -0.149 \end{pmatrix} \\
    &= \begin{pmatrix} 0.333 \\ 0.180 \end{pmatrix} + \begin{pmatrix} 0.250 \\ -0.083 \end{pmatrix} = \begin{pmatrix} 0.583 \\ 0.097 \end{pmatrix}
\end{align*}

\textbf{Passo 5 --- Output Gate e stato nascosto}:
Con $\mathbf{W}_o = \begin{pmatrix} 0.2 & -0.1 & 0.3 \\ -0.3 & 0.2 & 0.1 \end{pmatrix}$, $\mathbf{b}_o = \begin{pmatrix} 0 \\ 0 \end{pmatrix}$:
\begin{align*}
    \mathbf{o}_t &= \sigma\!\begin{pmatrix} 0.04 + 0.01 + 0.21 \\ -0.06 + (-0.02) + 0.07 \end{pmatrix} = \sigma\!\begin{pmatrix} 0.26 \\ -0.01 \end{pmatrix} = \begin{pmatrix} 0.565 \\ 0.498 \end{pmatrix} \\
    \mathbf{h}_t &= \mathbf{o}_t \odot \tanh(\mathbf{C}_t) = \begin{pmatrix} 0.565 \\ 0.498 \end{pmatrix} \odot \begin{pmatrix} 0.523 \\ 0.097 \end{pmatrix} = \begin{pmatrix} 0.296 \\ 0.048 \end{pmatrix}
\end{align*}

\textbf{Interpretazione}: la cella $\mathbf{C}_t = (0.583, 0.097)$ ha mantenuto buona parte della prima componente ($0.5 \to 0.583$) ma la seconda \`e scesa ($0.3 \to 0.097$). Lo stato nascosto $\mathbf{h}_t = (0.296, 0.048)$ \`e una versione ``filtrata'' della cella: l'output gate decide cosa rendere visibile.
\end{examplebox}

\begin{notabox}[Perch\'e la LSTM risolve il vanishing gradient]
Nella LSTM, il gradiente pu\`o fluire direttamente attraverso lo stato della cella $\mathbf{C}_t$ grazie alla connessione additiva: $\mathbf{C}_t = \mathbf{f}_t \odot \mathbf{C}_{t-1} + \ldots$. Se il forget gate \`e vicino a 1, il gradiente passa quasi inalterato: $\frac{\partial \mathbf{C}_t}{\partial \mathbf{C}_{t-1}} = \mathbf{f}_t \approx \mathbf{1}$. Questa \`e la stessa intuizione delle residual connections (Sezione~\ref{sec:residual}) usate nel Perceiver.
\end{notabox}



\newpage
%% ============================================================
%%  APPENDICE E --- GRU
%% ============================================================
\section{GRU (Gated Recurrent Unit)} \label{app:gru}

La \textbf{GRU} (Cho et al., 2014) \`e una semplificazione della LSTM (Appendice~\ref{app:lstm}) che combina forget gate e input gate in un'unica \textbf{update gate} e unifica lo stato della cella e lo stato nascosto. Risultato: solo \textbf{2 porte} invece di 3, meno parametri, addestramento pi\`u veloce, prestazioni spesso comparabili.

\begin{figure}[H]
\centering
\includegraphics[width=0.75\textwidth]{paper_figures/gru_architecture_original.png}
\caption{Architettura della cella GRU con le equazioni (da Olah, 2015). Si vedono le due porte: reset gate $\mathbf{r}_t$ e update gate $\mathbf{z}_t$. Lo stato nascosto viene aggiornato tramite interpolazione lineare: $(1 - \mathbf{z}_t) \cdot \mathbf{h}_{t-1} + \mathbf{z}_t \cdot \tilde{\mathbf{h}}_t$.}
\label{fig:gru_original_arch}
\end{figure}

Come si vede dalla Figura~\ref{fig:gru_original_arch}, la GRU \`e notevolmente pi\`u semplice della LSTM: solo 2 porte (reset e update) contro 3, e un unico stato $\mathbf{h}_t$ invece di due ($\mathbf{h}_t$ e $\mathbf{C}_t$). L'update gate $\mathbf{z}_t$ controlla il bilanciamento tra ``ricordare il passato'' e ``usare la nuova informazione'' tramite un'elegante interpolazione lineare.

\subsection{Le 2 porte della GRU}

\begin{formulabox}[Equazioni della GRU]
\textbf{1. Reset Gate} (``Quanto dello stato passato \`e rilevante?''):
\[
    \mathbf{r}_t = \sigma\!\left(\mathbf{W}_r \cdot [\mathbf{h}_{t-1}, \mathbf{x}_t] + \mathbf{b}_r\right)
\]
Se $\mathbf{r}_t \approx 0$, lo stato passato viene ``resettato'' (ignorato): utile quando l'input corrente \`e cos\`i diverso che il passato non \`e pi\`u rilevante.

\textbf{2. Update Gate} (``Quanto aggiornare lo stato?''):
\[
    \mathbf{z}_t = \sigma\!\left(\mathbf{W}_z \cdot [\mathbf{h}_{t-1}, \mathbf{x}_t] + \mathbf{b}_z\right)
\]
Controlla il bilanciamento tra vecchio stato e nuovo candidato.

\textbf{3. Stato candidato}:
\[
    \tilde{\mathbf{h}}_t = \tanh\!\left(\mathbf{W} \cdot [\mathbf{r}_t \odot \mathbf{h}_{t-1}, \mathbf{x}_t] + \mathbf{b}\right)
\]
Il reset gate $\mathbf{r}_t$ modula quanto dello stato passato entra nel calcolo del candidato.

\textbf{4. Nuovo stato nascosto}:
\[
    \mathbf{h}_t = (1 - \mathbf{z}_t) \odot \mathbf{h}_{t-1} + \mathbf{z}_t \odot \tilde{\mathbf{h}}_t
\]
Se $\mathbf{z}_t \approx 0$: $\mathbf{h}_t \approx \mathbf{h}_{t-1}$ (mantieni il vecchio stato).
Se $\mathbf{z}_t \approx 1$: $\mathbf{h}_t \approx \tilde{\mathbf{h}}_t$ (usa il nuovo candidato).
\end{formulabox}

% TikZ: GRU cell
\begin{figure}[H]
\centering
\begin{tikzpicture}[
    gate/.style={draw, rounded corners=2pt, fill=#1, minimum width=1.3cm, minimum height=0.7cm, font=\footnotesize\bfseries},
    gate/.default={yellow!20},
    op/.style={circle, draw, fill=gray!10, minimum size=0.5cm, font=\footnotesize},
    arrowstyle/.style={-{Stealth[length=2mm]}, thick},
    line/.style={thick},
    >=Stealth
]

% Hidden state flow
\draw[line, blue!60!black] (-1, 2) -- (1, 2);
\node[font=\footnotesize, blue!60!black] at (-1.8, 2) {$\mathbf{h}_{t-1}$};

% Reset gate
\node[gate=red!15] (rg) at (2, 0.5) {$\sigma$};
\node[font=\footnotesize] at (2, -0.1) {Reset $\mathbf{r}_t$};

% Update gate
\node[gate=green!15] (zg) at (5, 0.5) {$\sigma$};
\node[font=\footnotesize] at (5, -0.1) {Update $\mathbf{z}_t$};

% Candidate
\node[gate=cyan!15] (cand) at (3.5, 2) {$\tanh$};
\node[font=\footnotesize] at (3.5, 1.3) {Candidato $\tilde{\mathbf{h}}_t$};

% Multiply reset gate with h
\node[op] (rmul) at (2, 2) {$\times$};
\draw[arrowstyle] (rg) -- (rmul);
\draw[arrowstyle] (rmul) -- (cand);

% Linear interpolation
\node[op] (zmul1) at (6.5, 2) {$\times$};
\node[op] (zmul2) at (6.5, 3.5) {$\times$};
\node[op] (add) at (8, 2.75) {$+$};

% 1 - z_t branch
\node[font=\footnotesize] at (5.5, 3.8) {$1 - \mathbf{z}_t$};

\draw[arrowstyle] (cand) -- (zmul1);
\draw[line] (zg.north) |- (zmul1.south);
\draw[arrowstyle] (zmul1) -- (add);

% h_{t-1} branch going up
\draw[line, blue!60!black] (1, 2) -- (1, 3.5) -- (zmul2.west);
\draw[line] (zg.north) |- (zmul2.south);
\draw[arrowstyle] (zmul2) -- (add);

% Output
\draw[arrowstyle, blue!60!black] (add) -- (9.5, 2.75);
\node[font=\footnotesize, blue!60!black] at (10.2, 2.75) {$\mathbf{h}_t$};

% Input
\node[draw, fill=blue!10, rounded corners, minimum width=1cm, font=\footnotesize] (xt) at (3.5, -1) {$\mathbf{x}_t$};
\draw[line, gray] (xt) -- (3.5, 0.5);
\draw[line, gray] (3.5, 0.5) -| (rg.south);
\draw[line, gray] (3.5, 0.5) -| (zg.south);

\end{tikzpicture}
\caption{Cella GRU. Solo 2 porte (reset e update) contro le 3 della LSTM. L'update gate controlla il bilanciamento tra vecchio stato e nuovo candidato tramite interpolazione lineare.}
\label{fig:gru_cell}
\end{figure}

\subsection{Esempio numerico: forward pass di una GRU}

Per rendere concrete le formule, calcoliamo un passo della GRU con $d_h = 2$ (dimensione stato nascosto) e $d_x = 1$ (un solo input).

\begin{examplebox}[Forward pass GRU --- passo per passo]
\textbf{Dati iniziali}: stato nascosto $\mathbf{h}_{t-1} = \begin{pmatrix} 0.5 \\ -0.3 \end{pmatrix}$, input $x_t = 0.8$.

Il vettore concatenato $[\mathbf{h}_{t-1}, x_t]$ ha dimensione 3:
$\begin{pmatrix} 0.5 & -0.3 & 0.8 \end{pmatrix}$

\textbf{Passo 1 --- Reset Gate} $\mathbf{r}_t$. Supponiamo:
$\mathbf{W}_r = \begin{pmatrix} 0.4 & 0.1 & 0.3 \\ 0.2 & -0.5 & 0.1 \end{pmatrix}$, $\mathbf{b}_r = \begin{pmatrix} 0 \\ 0 \end{pmatrix}$
\begin{align*}
    \mathbf{W}_r \cdot [\mathbf{h}_{t-1}, x_t] &= \begin{pmatrix} 0.4 \cdot 0.5 + 0.1 \cdot (-0.3) + 0.3 \cdot 0.8 \\ 0.2 \cdot 0.5 + (-0.5) \cdot (-0.3) + 0.1 \cdot 0.8 \end{pmatrix} = \begin{pmatrix} 0.41 \\ 0.33 \end{pmatrix} \\
    \mathbf{r}_t &= \sigma\!\begin{pmatrix} 0.41 \\ 0.33 \end{pmatrix} = \begin{pmatrix} 0.601 \\ 0.582 \end{pmatrix}
\end{align*}

\textbf{Passo 2 --- Update Gate} $\mathbf{z}_t$. Supponiamo:
$\mathbf{W}_z = \begin{pmatrix} -0.2 & 0.3 & 0.5 \\ 0.1 & 0.4 & -0.1 \end{pmatrix}$, $\mathbf{b}_z = \begin{pmatrix} 0 \\ 0 \end{pmatrix}$
\begin{align*}
    \mathbf{W}_z \cdot [\mathbf{h}_{t-1}, x_t] &= \begin{pmatrix} -0.2 \cdot 0.5 + 0.3 \cdot (-0.3) + 0.5 \cdot 0.8 \\ 0.1 \cdot 0.5 + 0.4 \cdot (-0.3) + (-0.1) \cdot 0.8 \end{pmatrix} = \begin{pmatrix} 0.21 \\ -0.15 \end{pmatrix} \\
    \mathbf{z}_t &= \sigma\!\begin{pmatrix} 0.21 \\ -0.15 \end{pmatrix} = \begin{pmatrix} 0.552 \\ 0.463 \end{pmatrix}
\end{align*}

\textbf{Passo 3 --- Stato candidato} $\tilde{\mathbf{h}}_t$. Prima si modula lo stato passato con il reset gate:
\[
    \mathbf{r}_t \odot \mathbf{h}_{t-1} = \begin{pmatrix} 0.601 \\ 0.582 \end{pmatrix} \odot \begin{pmatrix} 0.5 \\ -0.3 \end{pmatrix} = \begin{pmatrix} 0.301 \\ -0.175 \end{pmatrix}
\]
Poi si concatena con $x_t$ e si applica $\tanh$. Supponiamo:
$\mathbf{W} = \begin{pmatrix} 0.3 & -0.2 & 0.6 \\ 0.1 & 0.5 & -0.3 \end{pmatrix}$, $\mathbf{b} = \begin{pmatrix} 0 \\ 0 \end{pmatrix}$
\begin{align*}
    \mathbf{W} \cdot [\mathbf{r}_t \odot \mathbf{h}_{t-1}, x_t] &= \begin{pmatrix} 0.3 \cdot 0.301 + (-0.2) \cdot (-0.175) + 0.6 \cdot 0.8 \\ 0.1 \cdot 0.301 + 0.5 \cdot (-0.175) + (-0.3) \cdot 0.8 \end{pmatrix} = \begin{pmatrix} 0.605 \\ -0.297 \end{pmatrix} \\
    \tilde{\mathbf{h}}_t &= \tanh\!\begin{pmatrix} 0.605 \\ -0.297 \end{pmatrix} = \begin{pmatrix} 0.541 \\ -0.289 \end{pmatrix}
\end{align*}

\textbf{Passo 4 --- Nuovo stato nascosto} $\mathbf{h}_t$ (interpolazione lineare):
\begin{align*}
    \mathbf{h}_t &= (1 - \mathbf{z}_t) \odot \mathbf{h}_{t-1} + \mathbf{z}_t \odot \tilde{\mathbf{h}}_t \\
    &= \begin{pmatrix} 0.448 \\ 0.537 \end{pmatrix} \odot \begin{pmatrix} 0.5 \\ -0.3 \end{pmatrix} + \begin{pmatrix} 0.552 \\ 0.463 \end{pmatrix} \odot \begin{pmatrix} 0.541 \\ -0.289 \end{pmatrix} \\
    &= \begin{pmatrix} 0.224 \\ -0.161 \end{pmatrix} + \begin{pmatrix} 0.299 \\ -0.134 \end{pmatrix} = \begin{pmatrix} 0.523 \\ -0.295 \end{pmatrix}
\end{align*}

\textbf{Interpretazione}: $\mathbf{z}_t \approx (0.55, 0.46)$ --- la GRU sta ``mescolando'' circa met\`a vecchio stato e met\`a nuovo candidato. Lo stato finale $\mathbf{h}_t = (0.523, -0.295)$ \`e vicino sia al vecchio $\mathbf{h}_{t-1} = (0.5, -0.3)$ sia al candidato $\tilde{\mathbf{h}}_t = (0.541, -0.289)$: un aggiornamento graduale.
\end{examplebox}


\subsection{Confronto GRU vs LSTM}

\begin{center}
\begin{tabular}{@{}lcc@{}}
\toprule
\textbf{Propriet\`a} & \textbf{LSTM} & \textbf{GRU} \\
\midrule
Numero di porte & 3 (forget, input, output) & 2 (reset, update) \\
Stato della cella & Separato ($\mathbf{C}_t$ e $\mathbf{h}_t$) & Unificato ($\mathbf{h}_t$) \\
Parametri per unit\`a & $4 \times d_h \times (d_h + d_x)$ & $3 \times d_h \times (d_h + d_x)$ \\
Velocit\`a di training & Pi\`u lento & Pi\`u veloce ($\sim$25\% meno parametri) \\
Prestazioni & Migliori su sequenze lunghe & Comparabili, a volte migliori \\
\bottomrule
\end{tabular}
\end{center}

\begin{notabox}[Quando usare GRU vs LSTM]
In generale, la GRU \`e preferita quando: (i) i dati sono pochi (meno parametri = meno overfitting); (ii) la velocit\`a di training \`e importante; (iii) le sequenze non sono estremamente lunghe. La LSTM \`e preferita per sequenze molto lunghe dove il forget gate separato offre un controllo pi\`u fine. Entrambe sono state in gran parte sostituite dai Transformer (Appendice~\ref{app:transformer}) che offrono parallelizzazione e migliori prestazioni.
\end{notabox}



\newpage
%% ============================================================
%%  APPENDICE F --- CNN
%% ============================================================
\section{Reti Neurali Convoluzionali (CNN)} \label{app:cnn}

Un'immagine digitale \`e una griglia di pixel: migliaia o milioni di numeri disposti in righe e colonne. Un MLP classico (Appendice~\ref{app:feedforward}) tratterebbe ogni pixel come un input indipendente, collegandolo a ogni neurone dello strato successivo. Per un'immagine $224 \times 224$ a colori, il solo primo strato richiederebbe oltre 7 miliardi di connessioni --- un approccio computazionalmente insostenibile che, per di pi\`u, ignora completamente il fatto che i pixel vicini sono correlati tra loro. Le \textbf{Convolutional Neural Networks} (CNN) nascono proprio per risolvere questo problema: sfruttano la \textbf{struttura spaziale} dell'immagine tramite piccoli filtri che scorrono localmente, riducendo drasticamente i parametri e catturando i pattern visivi in modo naturale.

Come discusso nella \hyperref[sec:specializzate]{Sezione~1.1}, le CNN sono l'architettura dominante per la visione artificiale. Tuttavia, questa specializzazione ha un prezzo: le CNN funzionano bene solo su dati con struttura a griglia. Il Perceiver (\hyperref[sec:architettura]{Sezione~3}) adotta una strategia opposta, trattando ogni pixel come un elemento di una sequenza generica, senza assumere nulla sulla geometria dei dati. Il risultato \`e un'architettura pi\`u flessibile che, pur rinunciando ai vantaggi strutturali delle CNN, raggiunge prestazioni competitive (78.0\% su ImageNet).

Per capire concretamente come funziona una CNN, seguiamo passo per passo l'architettura \textbf{LeNet-5} (LeCun et al., 1998) --- la prima CNN di successo, usata per riconoscere cifre scritte a mano. Nonostante sia un modello storico, contiene \emph{tutti} gli ingredienti fondamentali delle CNN moderne. La Figura~\ref{fig:lenet5_full_arch} mostra l'architettura completa che andremo a spiegare pezzo per pezzo.

\begin{figure}[H]
\centering
\includegraphics[width=0.9\textwidth]{paper_figures/lenet5_paper_original.png}
\caption{Architettura di LeNet-5 dal paper originale (LeCun et al., 1998). L'input $32 \times 32$ attraversa strati di convoluzione (C1, C3, C5) e subsampling/pooling (S2, S4), con le dimensioni delle feature map indicate (es.\ C1: 6 feature map $28 \times 28$). I layer finali sono fully connected (F6: 84 neuroni, Output: 10 classi).}
\label{fig:lenet5_full_arch}
\end{figure}

\begin{figure}[H]
\centering
\includegraphics[width=0.85\textwidth]{paper_figures/cnn_architecture_clean.jpg}
\caption{Schema generico di una CNN. La pipeline si divide in due fasi: \textbf{Feature Learning} (blocchi ripetuti di convoluzione + ReLU + pooling che estraggono features di complessit\`a crescente) e \textbf{Classification} (flatten + strati fully connected + softmax che producono le probabilit\`a di classe).}
\label{fig:cnn_clean_arch}
\end{figure}

L'idea di fondo \`e semplice: l'immagine entra da sinistra (Figura~\ref{fig:lenet5_full_arch}) e viene progressivamente ``compressa'' in rappresentazioni sempre pi\`u piccole e astratte. Un primo blocco di \emph{convoluzione + subsampling} estrae 6 feature map $14 \times 14$; un secondo blocco le raffina in 16 feature map $5 \times 5$; infine, questi valori vengono appiattiti e passati a strati fully connected che producono la classificazione. Lo schema generico di Figura~\ref{fig:cnn_clean_arch} evidenzia le due fasi fondamentali che ritroviamo in \emph{ogni} CNN: \textbf{Feature Learning} (i blocchi convoluzionali estraggono automaticamente features di complessit\`a crescente) e \textbf{Classification} (strati densi che prendono la decisione finale). Vediamo ora ogni pezzo nel dettaglio.


\subsection{Input: l'immagine come matrice di numeri}

Il punto di partenza \`e un'immagine in scala di grigi di $32 \times 32$ pixel. Si pu\`o pensare a questa immagine come a un \textbf{foglio di calcolo} con 32 righe e 32 colonne: ogni cella contiene un numero che rappresenta la luminosit\`a di quel punto (da 0 = nero a 255 = bianco). Formalmente, la rappresentazione \`e una matrice $32 \times 32 \times 1$ (un solo canale). Un'immagine a colori avrebbe 3 ``fogli'' sovrapposti --- uno per il rosso, uno per il verde, uno per il blu (RGB) --- ma per le cifre scritte a mano basta la luminanza.

\begin{figure}[H]
\centering
\includegraphics[width=0.35\textwidth]{paper_figures/lenet_01_input_32x32.png}
\caption{Input di LeNet-5: un'immagine $32 \times 32 \times 1$ (singolo canale di luminanza).}
\label{fig:cnn_input}
\end{figure}

Per il computer, dunque, non esiste alcuna ``immagine'' --- solo una griglia di 1\,024 numeri. Tornando all'analogia del foglio di calcolo: se ingrandissimo l'immagine fino a vedere i singoli pixel, ogni quadratino corrisponderebbe esattamente a una cella con il suo valore numerico (Figura~\ref{fig:cnn_pixel_matrix}). Il compito della CNN \`e imparare a estrarre significato da questa griglia di numeri.

\begin{figure}[H]
\centering
\includegraphics[width=0.5\textwidth]{paper_figures/cnn_frames/06_rgb_three_channels.png}
\caption{Un'immagine e il suo canale numerico: a sinistra l'immagine pixelata, a destra la matrice numerica dove ogni cella contiene il valore di luminanza (0--255).}
\label{fig:cnn_pixel_matrix}
\end{figure}


\subsection{Lo strato convoluzionale: il cuore della CNN}

Il componente chiave di una CNN \`e lo \textbf{strato convoluzionale}. L'idea \`e semplice: invece di connettere ogni pixel a ogni neurone (come farebbe un MLP --- il che richiederebbe un numero enorme di parametri), si usa un piccolo \textbf{filtro} (o \textbf{kernel}) che ``scorre'' sull'immagine e guarda solo una regione locale alla volta.

L'analogia pi\`u intuitiva \`e quella di una \textbf{lente d'ingrandimento} che scorre sistematicamente sull'immagine: ad ogni posizione, la lente esamina solo un piccolo riquadro di pixel, decide se il pattern che vede corrisponde a qualcosa di interessante (un bordo, un angolo, una curva), e annota il risultato prima di spostarsi alla posizione successiva. Il ``tipo'' di pattern cercato \`e determinato dai valori numerici del filtro.

Pi\`u formalmente, un kernel \`e una piccola matrice di pesi (in LeNet-5, $5 \times 5 = 25$ pesi). Ad ogni posizione, il kernel calcola il \textbf{prodotto scalare} tra i propri pesi e i pixel sottostanti, producendo un singolo numero. Scorrendo su tutta l'immagine, il kernel produce una \textbf{feature map}: una nuova ``immagine'' in cui ogni valore indica quanto fortemente quel pattern \`e presente nella regione corrispondente.

\begin{figure}[H]
\centering
\includegraphics[width=0.38\textwidth]{paper_figures/cnn_frames/10_kernel_mini_matrix.png}
\caption{Un kernel (feature detector): una piccola matrice $5 \times 5$ che scorre sull'immagine. Ad ogni posizione calcola il prodotto scalare con i pixel sottostanti, producendo un valore nella feature map di output.}
\label{fig:cnn_kernel}
\end{figure}

\begin{examplebox}[Convoluzione passo per passo]
Consideriamo un'immagine $5 \times 5$ (un solo canale) e un filtro $3 \times 3$:

\[
\text{Immagine} = \begin{pmatrix}
1 & 0 & 1 & 2 & 1 \\
0 & 1 & 2 & 0 & 0 \\
1 & 2 & 1 & 0 & 1 \\
0 & 0 & 1 & 1 & 2 \\
2 & 1 & 0 & 1 & 0
\end{pmatrix}
\qquad
\text{Filtro} = \begin{pmatrix}
1 & 0 & -1 \\
1 & 0 & -1 \\
1 & 0 & -1
\end{pmatrix}
\]

\textbf{Posizione (0,0)}: il filtro si sovrappone ai pixel $[0{:}3, 0{:}3]$:
\begin{align*}
    \text{output}_{0,0} &= 1{\cdot}1 + 0{\cdot}0 + 1{\cdot}(-1) + 0{\cdot}1 + 1{\cdot}0 + 2{\cdot}(-1) + 1{\cdot}1 + 2{\cdot}0 + 1{\cdot}(-1) \\
    &= 1 + 0 - 1 + 0 + 0 - 2 + 1 + 0 - 1 = \mathbf{-2}
\end{align*}

\textbf{Posizione (0,1)}: il filtro si sposta di 1 pixel a destra (stride = 1):
\begin{align*}
    \text{output}_{0,1} &= 0{\cdot}1 + 1{\cdot}0 + 2{\cdot}(-1) + 1{\cdot}1 + 2{\cdot}0 + 0{\cdot}(-1) + 2{\cdot}1 + 1{\cdot}0 + 0{\cdot}(-1) \\
    &= 0 + 0 - 2 + 1 + 0 + 0 + 2 + 0 + 0 = \mathbf{1}
\end{align*}

La feature map risultante ha dimensione $(5-3+1) \times (5-3+1) = 3 \times 3$.

Questo filtro \`e un \textbf{rilevatore di bordi verticali}: ha valori positivi a sinistra e negativi a destra, producendo valori alti dove c'\`e una transizione netta da sinistra a destra.
\end{examplebox}

Nell'esempio precedente, la feature map di output \`e pi\`u piccola dell'input: da $5 \times 5$ si \`e passati a $3 \times 3$. Questa riduzione non \`e casuale, ma segue una formula precisa che dipende dalla dimensione del filtro, dal passo con cui lo si sposta (\emph{stride}) e dall'eventuale aggiunta di pixel di bordo (\emph{padding}).

\begin{formulabox}[Dimensione dell'output di una convoluzione]
Data un'immagine $H \times W$ e un filtro $k \times k$ con stride $s$ e padding $p$:
\[
    H_{\text{out}} = \frac{H - k + 2p}{s} + 1 \qquad W_{\text{out}} = \frac{W - k + 2p}{s} + 1
\]
In LeNet-5, C1 ha $k = 5$, $s = 1$, $p = 0$: $\frac{32 - 5 + 0}{1} + 1 = 28$. Quindi l'output \`e $28 \times 28$.
\end{formulabox}

Se diversi kernel rilevano pattern diversi, \`e naturale usarne pi\`u di uno. In LeNet-5, il primo strato (C1) usa \textbf{6 kernel} $5 \times 5$, ciascuno che produce una feature map $28 \times 28$. Il risultato \`e un volume $28 \times 28 \times 6$ (Figura~\ref{fig:cnn_conv1}).

\begin{figure}[H]
\centering
\includegraphics[width=0.5\textwidth]{paper_figures/lenet_02_conv1_6filters.png}
\caption{Strato C1 di LeNet-5: 6 filtri $5 \times 5$ con stride 1 trasformano l'input $32 \times 32 \times 1$ in un volume $28 \times 28 \times 6$. Ogni filtro produce una feature map diversa.}
\label{fig:cnn_conv1}
\end{figure}

Ma che tipo di pattern rilevano concretamente questi 6 kernel? Ogni kernel, in virt\`u dei propri valori numerici, \`e sensibile a un tipo diverso di struttura locale. Ad esempio, un kernel con valori positivi a sinistra e negativi a destra (come quello dell'esempio precedente) risponde ai bordi verticali; uno con valori positivi in alto e negativi in basso rileva i bordi orizzontali; e cos\`i via per le diagonali (Figura~\ref{fig:cnn_edges}).

\begin{figure}[H]
\centering
\includegraphics[width=0.7\textwidth]{paper_figures/cnn_frames/15_edge_detection_result.png}
\caption{Diversi kernel applicati alla stessa immagine: ciascuno (con i propri valori numerici mostrati) evidenzia un tipo specifico di bordo. L'immagine originale \`e al centro.}
\label{fig:cnn_edges}
\end{figure}

Per rendere l'idea ancora pi\`u concreta, osserviamo il risultato effettivo del primo strato su una cifra reale: le 6 feature map prodotte per la cifra ``6'' (Figura~\ref{fig:cnn_6fmaps}) mostrano chiaramente come ciascun kernel abbia estratto un aspetto diverso della forma --- bordi orizzontali, verticali e angoli emergono in modo netto dalle diverse mappe.

\begin{figure}[H]
\centering
\includegraphics[width=0.55\textwidth]{paper_figures/cnn_frames/19_horizontal_vertical_corners.png}
\caption{Le 6 feature map del primo strato convoluzionale per la cifra ``6'' (input mostrato in basso). Ogni kernel rileva un pattern diverso: bordi orizzontali, verticali, angoli.}
\label{fig:cnn_6fmaps}
\end{figure}

Dopo la convoluzione si applica una \textbf{funzione di attivazione} non-lineare. In LeNet-5 si usa $\tanh$; nelle CNN moderne si preferisce $\text{ReLU}(x) = \max(0, x)$ (Sezione~\ref{sec:relu}). La non-linearit\`a \`e fondamentale: senza di essa, impilare pi\`u strati convoluzionali sarebbe matematicamente equivalente a un singolo strato lineare.


\subsection{Lo strato di pooling: ridurre le dimensioni}

Dopo la convoluzione, uno strato di \textbf{pooling} riduce la dimensione spaziale delle feature map, mantenendo le informazioni pi\`u importanti. In LeNet-5 si usa un pooling $2 \times 2$ con stride 2: per ogni finestra di $2 \times 2$ pixel, si mantiene un solo valore rappresentativo (tipicamente il massimo, detto \textbf{max-pooling}).

\begin{examplebox}[Max-Pooling $2 \times 2$]
Data una feature map $4 \times 4$, il max-pooling $2 \times 2$ con stride 2 seleziona il massimo in ogni finestra:

\[
\begin{pmatrix}
\mathbf{6} & 2 & 3 & \mathbf{5} \\
1 & 4 & 0 & 2 \\
\mathbf{8} & 3 & \mathbf{7} & 1 \\
0 & 5 & 2 & 4
\end{pmatrix}
\xrightarrow{\text{MaxPool } 2 \times 2}
\begin{pmatrix}
6 & 5 \\
8 & 7
\end{pmatrix}
\]

La dimensione si dimezza: da $4 \times 4$ a $2 \times 2$. Questo: (1) \textbf{riduce} il costo computazionale degli strati successivi, (2) introduce \textbf{invarianza traslazionale locale} (piccoli spostamenti dell'input non cambiano l'output), e (3) \textbf{seleziona} l'attivazione pi\`u forte in ogni regione.
\end{examplebox}

In LeNet-5, lo strato S2 applica pooling $2 \times 2$ con stride 2 alle 6 feature map di C1, riducendo le dimensioni da $28 \times 28 \times 6$ a $14 \times 14 \times 6$ (Figura~\ref{fig:cnn_pool1}).

\begin{figure}[H]
\centering
\includegraphics[width=0.55\textwidth]{paper_figures/lenet_03_conv1_pool1.png}
\caption{Strati C1 + S2 di LeNet-5: la convoluzione produce $28 \times 28 \times 6$, il pooling $2 \times 2$ riduce a $14 \times 14 \times 6$.}
\label{fig:cnn_pool1}
\end{figure}

L'effetto visivo \`e notevole (Figura~\ref{fig:cnn_pool_visual}): la feature map risultante \`e pi\`u piccola, eppure la struttura della cifra rimane perfettamente riconoscibile. Il pooling ha eliminato i dettagli ridondanti conservando l'informazione essenziale.

\begin{figure}[H]
\centering
\includegraphics[width=0.5\textwidth]{paper_figures/cnn_frames/24_conv_vs_pooling_compare.png}
\caption{Effetto del pooling sulla cifra ``5'': a sinistra la feature map dopo convoluzione (dimensione maggiore), a destra dopo pooling (dimensione ridotta ma struttura preservata).}
\label{fig:cnn_pool_visual}
\end{figure}


\subsection{Secondo blocco e gerarchia delle features}

Il passo successivo \`e ripetere la coppia convoluzione + pooling. In LeNet-5, lo strato C3 applica \textbf{16 filtri} $5 \times 5$ al volume $14 \times 14 \times 6$, producendo $10 \times 10 \times 16$. Lo strato S4 riduce con pooling a $5 \times 5 \times 16$ (Figura~\ref{fig:cnn_conv2}).

\begin{figure}[H]
\centering
\includegraphics[width=0.55\textwidth]{paper_figures/lenet_04_conv2_16filters.png}
\caption{Strati C3 + S4: 16 filtri $5 \times 5$ producono $10 \times 10 \times 16$, ridotto dal pooling a $5 \times 5 \times 16$. Le feature map sono ora molto piccole ma semanticamente ricche.}
\label{fig:cnn_conv2}
\end{figure}

Perch\'e servono pi\`u blocchi? L'analogia con la lettura \`e illuminante. Immaginiamo di leggere un testo: il primo livello di elaborazione riconosce i \textbf{tratti elementari} --- linee verticali, curve, angoli --- proprio come i nostri occhi distinguono i singoli tratti di inchiostro. Il secondo livello combina questi tratti in \textbf{strutture intermedie}: un tratto verticale accanto a una curva forma una ``d'', due linee che si incrociano formano una ``x''. Un eventuale terzo livello riconoscerebbe combinazioni ancora pi\`u complesse. In una CNN avviene lo stesso: i kernel dei primi strati rilevano features \textbf{semplici} (bordi, linee), ma i kernel del secondo strato operano sulle feature map del primo --- e quindi rilevano \textbf{combinazioni} di features semplici: angoli (combinazione di bordi orizzontali e verticali), curve, parti di cifre. Questo fenomeno si chiama \textbf{gerarchia delle rappresentazioni} ed \`e il principio fondamentale del deep learning (Figura~\ref{fig:cnn_hierarchy}).

\begin{figure}[H]
\centering
\includegraphics[width=0.7\textwidth]{paper_figures/cnn_frames/28_low_mid_high_features.png}
\caption{Gerarchia delle features apprese da una CNN: i primi strati (Low) rilevano bordi e linee semplici, gli strati intermedi (Mid) combinano pattern geometrici, e gli strati profondi (High) riconoscono strutture complesse.}
\label{fig:cnn_hierarchy}
\end{figure}


\subsection{Flatten e classificazione: dagli strati FC all'output}

A questo punto le 16 feature map $5 \times 5$ contengono una rappresentazione compatta e astratta dell'immagine originale: da 1\,024 pixel grezzi siamo passati a 400 valori ($5 \times 5 \times 16$) che codificano features di alto livello. Il passo finale \`e \textbf{classificare} questi 400 valori.

Le feature map vengono \emph{appiattite} (\textbf{flatten}) in un vettore di 400 elementi (Figura~\ref{fig:cnn_flatten}), che viene poi passato a strati \textbf{fully connected} (FC) --- identici a quelli di un MLP (Appendice~\ref{app:feedforward}).

\begin{figure}[H]
\centering
\includegraphics[width=0.45\textwidth]{paper_figures/lenet_05_flatten_400.png}
\caption{Flatten: le feature map $5 \times 5 \times 16$ vengono appiattite in un vettore di 400 elementi.}
\label{fig:cnn_flatten}
\end{figure}

In LeNet-5, gli strati FC sono: $400 \to 120 \to 84 \to 10$ (Figura~\ref{fig:cnn_fc}). L'ultimo strato ha 10 neuroni, uno per ogni cifra (0--9). Si applica la softmax (Sezione~\ref{sec:softmax}) per ottenere le probabilit\`a.

\begin{figure}[H]
\centering
\includegraphics[width=0.5\textwidth]{paper_figures/lenet_06_fc_layers.png}
\caption{Strati FC di LeNet-5: $400 \to 120 \to 84 \to 10$ (uno per cifra 0--9).}
\label{fig:cnn_fc}
\end{figure}

Mettendo insieme tutti i pezzi, la pipeline completa per una cifra concreta (Figura~\ref{fig:cnn_pipeline_full}) rende visibile l'intera trasformazione: le feature map diventano progressivamente pi\`u piccole e astratte, fino a condensarsi in un vettore di valori pronti per la classificazione.

\begin{figure}[H]
\centering
\includegraphics[width=0.75\textwidth]{paper_figures/cnn_10_classifier_fc.png}
\caption{Pipeline completa di feature extraction per la cifra ``0'': input $\to$ CONV+POOL (6 feature map) $\to$ CONV+POOL (16 feature map) $\to$ FC $\to$ classificazione.}
\label{fig:cnn_pipeline_full}
\end{figure}


\subsection{Conteggio dei parametri}

L'intero modello LeNet-5 ha solo $\sim$60\,806 parametri (Figura~\ref{fig:cnn_params}) --- un numero minuscolo per gli standard odierni, ma sufficiente per raggiungere $>$99\% di accuratezza sulle cifre.

\begin{figure}[H]
\centering
\includegraphics[width=0.75\textwidth]{paper_figures/lenet_08_total_params.png}
\caption{LeNet-5 con il conteggio totale dei parametri: $\sim$60\,806.}
\label{fig:cnn_params}
\end{figure}

\begin{examplebox}[Calcolo dei parametri di C1]
Lo strato C1 ha 6 filtri $5 \times 5 \times 1$, ciascuno con un bias:
\[
    \text{Parametri C1} = 6 \times (5 \times 5 \times 1 + 1) = 6 \times 26 = 156
\]
Confronto: un MLP fully connected che collegasse tutti i $32 \times 32 = 1\,024$ pixel ai $28 \times 28 \times 6 = 4\,704$ neuroni di output richiederebbe $1\,024 \times 4\,704 = 4\,816\,896$ parametri --- \textbf{30\,000 volte di pi\`u}. La convoluzione \`e efficiente perch\'e riusa lo stesso kernel in ogni posizione (\textbf{weight sharing}).
\end{examplebox}


\subsection{Perch\'e le CNN funzionano: i due inductive bias}

Le CNN sfruttano due propriet\`a fondamentali delle immagini:

\begin{enumerate}
    \item \textbf{Localit\`a}: i pixel vicini sono correlati. Un occhio \`e un pattern locale; ha senso analizzare l'immagine con filtri piccoli piuttosto che collegare ogni pixel a ogni neurone. Questo riduce drasticamente il numero di parametri.
    \item \textbf{Invarianza alla traslazione}: un gatto resta un gatto sia in alto a sinistra che in basso a destra. La \emph{condivisione dei pesi} (lo stesso kernel si applica ovunque) garantisce questa propriet\`a.
\end{enumerate}

Questi due \textbf{inductive bias} sono un'arma a doppio taglio. Da un lato, rendono le CNN straordinariamente efficienti per le immagini: con meno parametri e meno dati di addestramento, una CNN impara a ``vedere'' molto rapidamente, proprio perch\'e le sue assunzioni sulla struttura dei dati sono corrette. Dall'altro, queste stesse assunzioni la \emph{imprigionano}: una CNN non sa gestire dati senza struttura a griglia (audio, nuvole di punti 3D, input multimodale), perch\'e i concetti di ``vicinanza spaziale'' e ``traslazione'' non hanno significato in quei domini.

Il Perceiver adotta la strategia opposta: rinuncia completamente a entrambi gli inductive bias, trattando l'input come una sequenza generica di elementi senza relazioni geometriche predefinite. Il prezzo \`e che deve ``riscoprire da solo'', durante l'addestramento, quelle regolarit\`a che una CNN conosce per costruzione --- e per farlo ha bisogno di pi\`u dati e pi\`u calcolo. Il vantaggio \`e che la stessa architettura funziona su immagini, audio, video e nuvole di punti senza alcuna modifica strutturale.


\subsection{Training di una CNN}

L'addestramento segue lo schema standard:
\begin{enumerate}
    \item \textbf{Forward pass}: l'immagine attraversa tutti gli strati, producendo le probabilit\`a di classificazione.
    \item \textbf{Loss}: si calcola la cross-entropy tra le probabilit\`a predette e la classe vera.
    \item \textbf{Backpropagation}: si calcolano i gradienti di tutti i parametri (pesi dei filtri, bias, pesi FC).
    \item \textbf{Aggiornamento}: si aggiornano i parametri con gradient descent (o Adam, LAMB, ecc.).
\end{enumerate}

I parametri apprendibili sono i \textbf{valori dei filtri}: la rete impara automaticamente quali filtri sono utili durante l'addestramento. Nessuno programma manualmente un kernel ``rileva bordi verticali'' --- la backpropagation trova da sola i valori ottimali.

\begin{figure}[H]
\centering
\includegraphics[width=0.8\textwidth]{paper_figures/docx_cnn_pipeline.jpg}
\caption{Pipeline completa di una CNN generica: convoluzione (CONV), attivazione (ReLU), pooling (POOL) e strati fully connected (FC).}
\label{fig:cnn_pipeline_general}
\end{figure}


\newpage
%% ============================================================
%%  APPENDICE G --- RESNET
%% ============================================================
\section{ResNet e Residual Connections} \label{app:resnet}

Le \textbf{Residual Networks} (ResNet, He et al., 2016) hanno rivoluzionato il deep learning dimostrando che reti \emph{molto} profonde (50, 101, fino a 152 strati) possono essere addestrate efficacemente grazie alle \textbf{skip connections} (connessioni residue). Le residual connections sono un componente fondamentale sia del Transformer (Sezione~\ref{sec:prerequisiti}) sia del Perceiver (Sezione~\ref{sec:residual}).

\begin{figure}[H]
\centering
\includegraphics[width=0.8\textwidth]{paper_figures/docx_resnet50_arch.png}
\caption{Architettura di ResNet-50. Si osservano i diversi stage della rete, ciascuno composto da blocchi residui (bottleneck blocks) con convoluzioni $1 \times 1$, $3 \times 3$ e $1 \times 1$. Le dimensioni delle feature map si riducono progressivamente mentre il numero di canali aumenta.}
\label{fig:docx_resnet50_arch}
\end{figure}

Come si vede dalla Figura~\ref{fig:docx_resnet50_arch}, ResNet-50 \`e organizzata in stage successivi, ognuno con un numero crescente di canali e dimensioni spaziali decrescenti. Le skip connections, visibili come frecce che ``scavalcano'' i blocchi convoluzionali, permettono al gradiente di fluire direttamente attraverso la rete, risolvendo il problema della degradazione descritto di seguito.

\subsection{Il problema della degradazione}

Prima di ResNet, si osservava un fenomeno controintuitivo: aggiungere pi\`u strati a una rete peggiorava le prestazioni, anche sul training set. Questo \textbf{non} \`e solo vanishing gradient (Appendice~\ref{app:rnn}): \`e un \textbf{problema di ottimizzazione}. Una rete a 56 strati funziona peggio di una a 20 strati, nonostante la rete pi\`u profonda contenga la soluzione della rete pi\`u bassa (gli strati aggiuntivi potrebbero semplicemente implementare la funzione identit\`a). Ma in pratica l'ottimizzatore non riesce a trovare questa soluzione.

\subsection{La soluzione: skip connections}

\begin{formulabox}[Blocco residuo]
\[
    \mathbf{y} = F(\mathbf{x}) + \mathbf{x}
\]
dove $F$ \`e una funzione (tipicamente 2-3 strati convoluzionali) e $\mathbf{x}$ \`e l'input. Se la trasformazione ottimale \`e vicina all'identit\`a, il blocco deve apprendere solo $F(\mathbf{x}) \approx \mathbf{0}$, che \`e molto pi\`u facile che apprendere una trasformazione completa da zero.
\end{formulabox}

% TikZ: Residual block
\begin{figure}[H]
\centering
\begin{tikzpicture}[
    block/.style={draw, rounded corners=3pt, fill=#1, minimum width=2.5cm, minimum height=0.7cm, font=\small},
    block/.default={orange!15},
    arrowstyle/.style={-{Stealth[length=2.5mm]}, thick},
    >=Stealth
]

% Input
\node[font=\small\bfseries] (x) at (0, 0) {$\mathbf{x}$};

% Branch point
\coordinate (branch) at (1.5, 0);

% Main path
\node[block=orange!15] (conv1) at (4, 0) {Conv $3 \times 3$ + BN + ReLU};
\node[block=orange!15] (conv2) at (4, -1.5) {Conv $3 \times 3$ + BN};

% Addition
\node[circle, draw, fill=green!15, minimum size=0.6cm, font=\small] (add) at (7, -1.5) {$+$};

% ReLU after addition
\node[block=yellow!15, minimum width=1.5cm] (relu) at (9.5, -1.5) {ReLU};

% Output
\node[font=\small\bfseries] (y) at (11.5, -1.5) {$\mathbf{y}$};

% Arrows - main path
\draw[arrowstyle] (x) -- (branch);
\draw[arrowstyle] (branch) -- (conv1);
\draw[arrowstyle] (conv1) -- (conv2);
\draw[arrowstyle] (conv2) -- (add);
\draw[arrowstyle] (add) -- (relu);
\draw[arrowstyle] (relu) -- (y);

% Skip connection
\draw[arrowstyle, blue!60!black, thick] (branch) -- (1.5, -1.5) -- (add.west);
\node[font=\footnotesize, blue!60!black] at (0.7, -0.75) {skip};
\node[font=\footnotesize, blue!60!black] at (0.3, -1.0) {connection};

% Labels
\node[font=\footnotesize\itshape, above=0.1cm of conv1] {$F(\mathbf{x})$};

\end{tikzpicture}
\caption{Blocco residuo di ResNet. L'input $\mathbf{x}$ viene sommato all'output della trasformazione $F(\mathbf{x})$, formando $\mathbf{y} = F(\mathbf{x}) + \mathbf{x}$.}
\label{fig:residual_block}
\end{figure}


\subsection{Bottleneck block}

Per reti molto profonde (ResNet-50 e oltre), si usa un \textbf{bottleneck block} con tre convoluzioni:

\begin{formulabox}[Bottleneck Block]
\begin{enumerate}
    \item \textbf{Conv $1 \times 1$} (riduzione): da 256 canali a 64 canali. Riduce la dimensionalit\`a.
    \item \textbf{Conv $3 \times 3$}: operazione convoluzionale nello spazio ridotto a 64 canali.
    \item \textbf{Conv $1 \times 1$} (espansione): da 64 canali a 256 canali. Ripristina la dimensione.
\end{enumerate}
Questo schema riduce i FLOP di circa $3\times$ rispetto a due convoluzioni $3 \times 3$ dirette a 256 canali.
\end{formulabox}

L'idea del bottleneck \`e concettualmente simile al \textbf{latent bottleneck del Perceiver} (\hyperref[sec:latent_bottleneck]{Sezione~1.3}): comprimere l'informazione in uno spazio pi\`u piccolo, elaborarla, e poi ri-espanderla.

\subsection{Esempio numerico: perch\'e la skip connection salva il gradiente}

\begin{examplebox}[Gradient flow: senza vs con residual connection]
Consideriamo un blocco con input $x = 1.5$ e una trasformazione $F(x) = W_2 \cdot \text{ReLU}(W_1 \cdot x)$ con $W_1 = 0.3$ e $W_2 = 0.4$.

\textbf{Senza residual connection} (rete classica): $y = F(x)$.
\begin{align*}
    z_1 &= W_1 \cdot x = 0.3 \cdot 1.5 = 0.45 \\
    a_1 &= \text{ReLU}(0.45) = 0.45 \\
    y &= W_2 \cdot a_1 = 0.4 \cdot 0.45 = 0.18
\end{align*}
Il gradiente $\frac{\partial y}{\partial x} = W_2 \cdot W_1 = 0.4 \cdot 0.3 = \mathbf{0.12}$.

Con 10 blocchi impilati: il gradiente si moltiplica $\to 0.12^{10} = 6.2 \times 10^{-10}$ --- \textbf{vanishing!}

\textbf{Con residual connection}: $y = F(x) + x$.
\begin{align*}
    y &= F(x) + x = 0.18 + 1.5 = 1.68
\end{align*}
Il gradiente: $\frac{\partial y}{\partial x} = \frac{\partial F(x)}{\partial x} + \frac{\partial x}{\partial x} = 0.12 + \mathbf{1} = \mathbf{1.12}$.

Con 10 blocchi: il termine $+1$ in ogni blocco impedisce che il gradiente svanisca. Il prodotto di $(1 + \epsilon_k)$ per 10 blocchi rimane vicino a 1, invece di collassare a zero.

\textbf{Conclusione}: la skip connection aggiunge un ``percorso diretto'' per il gradiente (il termine $+1$), che garantisce che l'informazione possa fluire dalla loss fino ai primi strati senza attenuarsi.
\end{examplebox}


\subsection{Confronto CNN tradizionale vs ResNet}

\begin{center}
\small
\begin{tabular}{@{}p{3.5cm}p{4cm}p{4.5cm}@{}}
\toprule
\textbf{Propriet\`a} & \textbf{CNN tradizionale} & \textbf{ResNet} \\
\midrule
Profondit\`a max & $\sim$20 strati & 152+ strati \\
Flusso gradiente & Si attenua con la profondit\`a & Skip preserva il gradiente \\
Ottimizzazione & Difficile per reti profonde & $F(\mathbf{x}) \approx 0$ facile \\
Top-1 ImageNet & 72.0\% (VGG-19) & 76.0\% (ResNet-50) \\
\bottomrule
\end{tabular}
\end{center}



\newpage
%% ============================================================
%%  APPENDICE H --- TRANSFORMER DETTAGLIATO
%% ============================================================
\section{Il Transformer (Vaswani et al., 2017) --- Dettagliato} \label{app:transformer}

Il \textbf{Transformer} \`e l'architettura che ha rivoluzionato il deep learning a partire dal 2017. Il Perceiver (\hyperref[sec:architettura]{Sezione~3}) \`e una generalizzazione del Transformer che disaccoppia la complessit\`a dall'input tramite il latent bottleneck. Questa appendice descrive in modo dettagliato e discorsivo tutti i componenti del Transformer, con esempi numerici passo per passo.

\begin{figure}[H]
\centering
\includegraphics[width=0.85\textwidth]{paper_figures/docx_transformer_full_zoom.jpg}
\caption{Architettura completa del Transformer (Vaswani et al., 2017) con zoom sui componenti chiave. A sinistra: architettura encoder-decoder completa. Al centro: dettaglio del Multi-Head Attention. A destra: dettaglio della Scaled Dot-Product Attention.}
\label{fig:docx_transformer_full_zoom}
\end{figure}

Come si vede dalla Figura~\ref{fig:docx_transformer_full_zoom}, il Transformer \`e composto da un encoder (a sinistra) e un decoder (a destra), ciascuno formato da una pila di blocchi identici. I due componenti chiave ingranditi --- Multi-Head Attention e Scaled Dot-Product Attention --- sono il cuore del meccanismo di attenzione che permette a ogni elemento della sequenza di interagire con tutti gli altri.

\subsection{Motivazione: superare i limiti delle RNN}

Le RNN (Appendice~\ref{app:rnn}) hanno due limiti fondamentali:
\begin{enumerate}
    \item \textbf{Sequenzialit\`a}: ogni passo dipende dal precedente, impedendo la parallelizzazione. Per una sequenza di 1000 token, servono 1000 passi sequenziali.
    \item \textbf{Vanishing gradient}: nonostante LSTM e GRU, le dipendenze a lunghissimo raggio rimangono problematiche.
\end{enumerate}

Il Transformer risolve entrambi i problemi: l'attenzione permette a qualsiasi token di interagire direttamente con qualsiasi altro token (distanza 1 in termini di flusso del gradiente), e tutti i token sono elaborati in parallelo.


\subsection{Architettura completa encoder-decoder}

Il Transformer originale (per la traduzione automatica) ha due componenti principali:

\begin{itemize}
    \item \textbf{Encoder}: una pila di $N = 6$ blocchi identici. Ogni blocco contiene:
    \begin{enumerate}
        \item Multi-Head Self-Attention
        \item Feed-Forward Network (MLP a 2 strati)
    \end{enumerate}
    Con residual connection e layer normalization dopo ogni sotto-layer.

    \item \textbf{Decoder}: una pila di $N = 6$ blocchi. Ogni blocco contiene:
    \begin{enumerate}
        \item Masked Multi-Head Self-Attention (mascherata)
        \item Multi-Head Cross-Attention (query dal decoder, K/V dall'encoder)
        \item Feed-Forward Network
    \end{enumerate}
\end{itemize}

% TikZ: Transformer architecture
\begin{figure}[H]
\centering
\begin{tikzpicture}[
    block/.style={draw, rounded corners=3pt, fill=#1, minimum width=3.5cm, minimum height=0.6cm, font=\footnotesize},
    block/.default={orange!15},
    smallblock/.style={draw, rounded corners=2pt, fill=#1, minimum width=2.8cm, minimum height=0.5cm, font=\scriptsize},
    smallblock/.default={gray!10},
    arrowstyle/.style={-{Stealth[length=2mm]}, thick},
    braceann/.style={decorate, decoration={brace, amplitude=5pt}, thick},
    >=Stealth,
    scale=0.9, transform shape
]

% Encoder
\node[font=\large\bfseries] at (0, 7.8) {Encoder};

\node[block=blue!10] (epe) at (0, 0) {Positional Encoding};
\node[block=orange!15] (esa) at (0, 1.5) {Multi-Head Self-Attention};
\node[smallblock=gray!10] (ean) at (0, 2.3) {Add \& Norm};
\node[block=green!15] (effn) at (0, 3.3) {Feed-Forward Network};
\node[smallblock=gray!10] (ean2) at (0, 4.1) {Add \& Norm};

% Encoder brace
\draw[braceann, blue!50!black] (-2.3, 1.2) -- (-2.3, 4.4) node[midway, left=6pt, font=\footnotesize, blue!50!black] {$\times N$};

% Input embedding
\node[block=red!10, minimum width=3cm] (emb) at (0, -1) {Input Embedding};
\node[font=\footnotesize] at (0, -1.8) {Input tokens};

\draw[arrowstyle] (emb) -- (epe);
\draw[arrowstyle] (epe) -- (esa);
\draw[arrowstyle] (esa) -- (ean);
\draw[arrowstyle] (ean) -- (effn);
\draw[arrowstyle] (effn) -- (ean2);

% Encoder output arrow
\draw[arrowstyle, red!60!black, thick] (ean2.north) -- (0, 5.5) -- (5, 5.5);
\node[font=\scriptsize, red!60!black] at (2.5, 5.8) {Encoder output (K, V per cross-attention)};

% Decoder
\node[font=\large\bfseries] at (8, 7.8) {Decoder};

\node[block=blue!10] (dpe) at (8, 0) {Positional Encoding};
\node[block=orange!15] (dmsa) at (8, 1.5) {Masked Self-Attention};
\node[smallblock=gray!10] (dan1) at (8, 2.3) {Add \& Norm};
\node[block=purple!15] (dca) at (8, 3.3) {Cross-Attention};
\node[smallblock=gray!10] (dan2) at (8, 4.1) {Add \& Norm};
\node[block=green!15] (dffn) at (8, 5.1) {Feed-Forward Network};
\node[smallblock=gray!10] (dan3) at (8, 5.9) {Add \& Norm};

% Decoder brace
\draw[braceann, purple!50!black] (10.3, 1.2) -- (10.3, 6.2) node[midway, right=6pt, font=\footnotesize, purple!50!black] {$\times N$};

% Decoder input
\node[block=red!10, minimum width=3cm] (demb) at (8, -1) {Output Embedding};
\node[font=\footnotesize] at (8, -1.8) {Output tokens (shifted)};

\draw[arrowstyle] (demb) -- (dpe);
\draw[arrowstyle] (dpe) -- (dmsa);
\draw[arrowstyle] (dmsa) -- (dan1);
\draw[arrowstyle] (dan1) -- (dca);

% Cross-attention K, V from encoder
\draw[arrowstyle, red!60!black] (5, 5.5) -- (dca.west);

\draw[arrowstyle] (dca) -- (dan2);
\draw[arrowstyle] (dan2) -- (dffn);
\draw[arrowstyle] (dffn) -- (dan3);

% Output
\node[block=cyan!15, minimum width=2.5cm] (linear) at (8, 6.8) {Linear + Softmax};
\draw[arrowstyle] (dan3) -- (linear);

\end{tikzpicture}
\caption{Architettura completa del Transformer (Vaswani et al., 2017). L'encoder elabora la sequenza sorgente; il decoder genera la sequenza target token per token, con cross-attention sull'encoder.}
\label{fig:transformer_full}
\end{figure}

\begin{figure}[H]
\centering
\includegraphics[width=0.5\textwidth]{paper_figures/docx_transformer_encoder_highlight.png}
\caption{Architettura del Transformer con l'\textbf{encoder} evidenziato in rosso. L'encoder \`e composto da $N$ blocchi identici, ciascuno con Multi-Head Self-Attention e Feed-Forward Network, con residual connections e layer normalization.}
\label{fig:docx_transformer_encoder}
\end{figure}

Osservando il diagramma in Figura~\ref{fig:docx_transformer_encoder}, notiamo che l'encoder (evidenziato in rosso) riceve la sequenza sorgente e la elabora attraverso blocchi ripetuti di self-attention e feed-forward network. L'output dell'encoder fornisce le rappresentazioni contestualizzate che il decoder utilizzer\`a come chiavi e valori nella cross-attention.

\begin{figure}[H]
\centering
\includegraphics[width=0.5\textwidth]{paper_figures/docx_transformer_decoder_highlight.png}
\caption{Architettura del Transformer con il \textbf{decoder} evidenziato in rosso. Il decoder aggiunge la Masked Self-Attention e la Cross-Attention rispetto all'encoder, permettendo la generazione autoregressiva condizionata sulla sequenza sorgente.}
\label{fig:docx_transformer_decoder}
\end{figure}

Osservando la Figura~\ref{fig:docx_transformer_decoder}, il decoder (evidenziato in rosso) presenta una struttura pi\`u complessa dell'encoder: oltre alla self-attention (mascherata per impedire di ``guardare il futuro''), include un blocco di cross-attention che permette a ogni token generato di accedere alle rappresentazioni dell'encoder. Questo meccanismo di cross-attention \`e lo stesso principio utilizzato dal Perceiver.

\subsection{Self-Attention con esempio numerico passo passo}

Vediamo un esempio completo e dettagliato di self-attention con una sequenza di 3 token in dimensione $d = 4$.

\begin{examplebox}[Esempio numerico: self-attention passo per passo]
Sequenza di input (3 token, $d = 4$):
\[
\mathbf{X} = \begin{pmatrix} 1 & 0 & 1 & 0 \\ 0 & 2 & 0 & 2 \\ 1 & 1 & 1 & 1 \end{pmatrix}
\]

Matrici di proiezione semplificate ($d_k = d_v = 2$):
\[
\mathbf{W}_Q = \begin{pmatrix} 1 & 0 \\ 0 & 1 \\ 1 & 0 \\ 0 & 1 \end{pmatrix}, \quad
\mathbf{W}_K = \begin{pmatrix} 0 & 1 \\ 1 & 0 \\ 0 & 1 \\ 1 & 0 \end{pmatrix}, \quad
\mathbf{W}_V = \begin{pmatrix} 1 & 1 \\ 0 & 0 \\ 0 & 1 \\ 1 & 0 \end{pmatrix}
\]

\textbf{Passo 1: Calcolo Q, K, V}:
\begin{align*}
\mathbf{Q} &= \mathbf{X} \cdot \mathbf{W}_Q = \begin{pmatrix} 2 & 0 \\ 0 & 4 \\ 2 & 2 \end{pmatrix}, \quad
\mathbf{K} = \mathbf{X} \cdot \mathbf{W}_K = \begin{pmatrix} 0 & 2 \\ 4 & 0 \\ 2 & 2 \end{pmatrix} \\
\mathbf{V} &= \mathbf{X} \cdot \mathbf{W}_V = \begin{pmatrix} 1 & 2 \\ 2 & 0 \\ 2 & 2 \end{pmatrix}
\end{align*}

\textbf{Passo 2: Score} ($\mathbf{Q} \cdot \mathbf{K}^\top$):
\[
\mathbf{S} = \mathbf{Q} \cdot \mathbf{K}^\top = \begin{pmatrix} 0 & 8 & 4 \\ 0 & 0 & 8 \\ 8 & 8 & 8 \end{pmatrix}
\]

\textbf{Passo 3: Scaling} (dividi per $\sqrt{d_k} = \sqrt{2} \approx 1.414$):
\[
\frac{\mathbf{S}}{\sqrt{2}} \approx \begin{pmatrix} 0 & 5.66 & 2.83 \\ 0 & 0 & 5.66 \\ 5.66 & 5.66 & 5.66 \end{pmatrix}
\]

\textbf{Passo 4: Softmax} (per riga):
\[
\mathbf{A} \approx \begin{pmatrix} 0.003 & 0.946 & 0.051 \\ 0.003 & 0.003 & 0.994 \\ 0.333 & 0.333 & 0.334 \end{pmatrix}
\]
Il token 1 attende quasi esclusivamente al token 2; il token 2 attende al token 3; il token 3 attende uniformemente a tutti.

\textbf{Passo 5: Output} ($\mathbf{A} \cdot \mathbf{V}$):
\[
\text{Output} = \mathbf{A} \cdot \mathbf{V} \approx \begin{pmatrix} 2.0 & 0.1 \\ 2.0 & 2.0 \\ 1.7 & 1.3 \end{pmatrix}
\]
\end{examplebox}


\subsection{Multi-Head Attention con esempio numerico}

\begin{figure}[H]
\centering
\includegraphics[width=0.7\textwidth]{paper_figures/docx_sdpa_multihead.png}
\caption{Diagrammi dettagliati dei meccanismi di attenzione. A sinistra: Scaled Dot-Product Attention, che calcola il prodotto scalare tra query e chiavi, scala per $\sqrt{d_k}$, applica la softmax e moltiplica per i valori. A destra: Multi-Head Attention, che esegue $h$ attenzioni in parallelo su sottospazi diversi e concatena i risultati.}
\label{fig:docx_sdpa_multihead}
\end{figure}

Il diagramma in Figura~\ref{fig:docx_sdpa_multihead} mostra chiaramente come la Scaled Dot-Product Attention (a sinistra) costituisca il blocco fondamentale, mentre la Multi-Head Attention (a destra) esegua multiple istanze di questa operazione in parallelo su proiezioni diverse dell'input, concatenando poi i risultati per ottenere una rappresentazione pi\`u ricca.

\begin{examplebox}[Multi-Head Attention: idea e dimensioni]
Con $d = 512$ e $H = 8$ teste:
\begin{itemize}
    \item Ogni testa opera su un sottospazio di $d_k = d_v = 512/8 = 64$ dimensioni.
    \item La testa $h$ ha le proprie matrici $\mathbf{W}_Q^h, \mathbf{W}_K^h, \mathbf{W}_V^h \in \mathbb{R}^{512 \times 64}$.
    \item Ogni testa produce un output $\in \mathbb{R}^{n \times 64}$.
    \item I risultati vengono concatenati: $[\text{head}_1; \ldots; \text{head}_8] \in \mathbb{R}^{n \times 512}$.
    \item La matrice di proiezione $\mathbf{W}_O \in \mathbb{R}^{512 \times 512}$ produce l'output finale.
\end{itemize}

Ogni testa pu\`o specializzarsi: una testa potrebbe catturare relazioni soggetto-verbo, un'altra relazioni di prossimit\`a, un'altra pattern di negazione. La concatenazione combina tutte queste ``prospettive'' diverse.
\end{examplebox}


\subsection{Masked attention nel decoder}

\begin{figure}[H]
\centering
\includegraphics[width=0.5\textwidth]{paper_figures/docx_sdpa_mask.png}
\caption{Scaled Dot-Product Attention con evidenziato il passaggio di \textbf{Mask} (freccia rossa). La maschera viene applicata dopo lo scaling e prima della softmax: i valori corrispondenti a posizioni future vengono impostati a $-\infty$, in modo che la softmax li trasformi in zero.}
\label{fig:docx_sdpa_mask}
\end{figure}

Come evidenziato dalla freccia rossa in Figura~\ref{fig:docx_sdpa_mask}, il passaggio di mascheramento (Mask) \`e un'operazione opzionale che si inserisce tra lo scaling e la softmax. Nel decoder, questa maschera triangolare inferiore impedisce al modello di ``barare'' guardando i token futuri durante la generazione autoregressiva.

Nel decoder, durante la generazione, il token alla posizione $t$ pu\`o vedere solo i token alle posizioni $1, 2, \ldots, t$ (non quelli futuri). Questo si implementa con una \textbf{maschera triangolare inferiore}:

\begin{formulabox}[Masked Self-Attention]
\[
    \text{MaskedAttention}(\mathbf{Q}, \mathbf{K}, \mathbf{V}) = \text{softmax}\!\left(\frac{\mathbf{Q}\mathbf{K}^\top}{\sqrt{d_k}} + \mathbf{M}\right) \mathbf{V}
\]
dove $M_{ij} = 0$ se $i \geq j$ (posizioni visibili) e $M_{ij} = -\infty$ se $i < j$ (posizioni future mascherate). La softmax trasforma $-\infty$ in 0, impedendo l'attenzione sul futuro.
\end{formulabox}


\subsection{Cross-attention encoder-decoder}

Nel decoder, la \textbf{cross-attention} permette a ogni token generato di ``guardare'' l'intera sequenza sorgente elaborata dall'encoder:
\begin{itemize}
    \item \textbf{Query}: dallo stato del decoder (ci\`o che stiamo generando)
    \item \textbf{Key, Value}: dall'output dell'encoder (la sequenza sorgente)
\end{itemize}

Questo \`e lo \emph{stesso meccanismo} usato nel Perceiver (Sezione~\ref{sec:cross_attention}), dove i latenti fanno le query e l'input fornisce chiavi e valori.


\subsection{Layer Normalization pre-norm}

Il Transformer originale usa \textbf{post-norm}: $\text{LN}(\mathbf{x} + f(\mathbf{x}))$. Le architetture pi\`u recenti (GPT-2, Perceiver) usano \textbf{pre-norm}: $\mathbf{x} + f(\text{LN}(\mathbf{x}))$. La pre-norm \`e pi\`u stabile per reti molto profonde perch\'e ogni sotto-blocco riceve input normalizzati, evitando che le attivazioni crescano in modo incontrollato.


\subsection{Positional Encoding sinusoidale}

L'attenzione \`e invariante per permutazione: senza informazione posizionale, il Transformer non sa distinguere ``il gatto mangia il topo'' da ``il topo mangia il gatto''. Il positional encoding risolve questo problema.

\begin{formulabox}[Positional Encoding sinusoidale]
\[
    \text{PE}(pos, 2i) = \sin\!\left(\frac{pos}{10000^{2i/d}}\right), \quad
    \text{PE}(pos, 2i+1) = \cos\!\left(\frac{pos}{10000^{2i/d}}\right)
\]
Le frequenze decrescono geometricamente: le prime dimensioni oscillano rapidamente (distinguono posizioni vicine), le ultime oscillano lentamente (catturano la struttura globale).

\textbf{Propriet\`a chiave}: $\text{PE}(pos + k)$ pu\`o essere espresso come trasformazione lineare di $\text{PE}(pos)$, permettendo al modello di apprendere relazioni posizionali relative.
\end{formulabox}

% TikZ: Positional encoding visualization
\begin{figure}[H]
\centering
\begin{tikzpicture}
\begin{axis}[
    width=12cm, height=5cm,
    xlabel={Posizione ($pos$)},
    ylabel={Valore PE},
    legend pos=outer north east,
    grid=major,
    domain=0:50,
    samples=200,
    ymin=-1.2, ymax=1.2
]
\addplot[blue, thick] {sin(deg(x / 1))};
\addlegendentry{$i=0$ (alta freq.)}
\addplot[red, thick] {sin(deg(x / 10))};
\addlegendentry{$i=2$ (media freq.)}
\addplot[green!60!black, thick] {sin(deg(x / 100))};
\addlegendentry{$i=4$ (bassa freq.)}
\end{axis}
\end{tikzpicture}
\caption{Visualizzazione del positional encoding sinusoidale per diverse dimensioni. Le dimensioni con indice basso ($i=0$) oscillano rapidamente, distinguendo posizioni vicine; le dimensioni con indice alto ($i=4$) oscillano lentamente.}
\label{fig:pe_sinusoidal}
\end{figure}



\newpage
%% ============================================================
%%  APPENDICE I --- VISION TRANSFORMER (ViT)
%% ============================================================
\section{Vision Transformer (ViT)} \label{app:vit}

Il \textbf{Vision Transformer} (ViT, Dosovitskiy et al., 2021) ha dimostrato che il Transformer (Appendice~\ref{app:transformer}), originariamente progettato per il testo, pu\`o essere applicato con successo alle immagini. Come discusso nella Sezione~\ref{sec:limite_transformer}, il ViT suddivide l'immagine in \textbf{patch} e le tratta come token, applicando un Transformer encoder standard. Il Perceiver (\hyperref[sec:architettura]{Sezione~3}) va oltre, eliminando anche la necessit\`a delle patch e operando direttamente sui pixel.

\begin{figure}[H]
\centering
\includegraphics[width=0.8\textwidth]{paper_figures/docx_vit_architecture.png}
\caption{Pipeline completa del Vision Transformer (ViT). L'immagine viene suddivisa in patch di dimensione fissa, ciascuna patch viene appiattita e proiettata linearmente in un embedding. Si aggiunge un token CLS e il positional encoding, poi la sequenza viene elaborata da un Transformer encoder standard. L'output del token CLS viene usato per la classificazione tramite un MLP head.}
\label{fig:docx_vit_architecture}
\end{figure}

Come si vede dalla Figura~\ref{fig:docx_vit_architecture}, il ViT trasforma un problema di visione artificiale in un problema di elaborazione di sequenze: l'immagine viene ``tokenizzata'' in patch, e da quel punto in poi il processing \`e identico a quello di un Transformer per il testo. Questa intuizione \`e fondamentale per comprendere come il Perceiver generalizza ulteriormente l'approccio, eliminando anche la necessit\`a di suddividere l'immagine in patch.

\subsection{Motivazione}

Le CNN (Appendice~\ref{app:cnn}) sono ottime per le immagini grazie al loro inductive bias (localit\`a, invarianza traslazionale), ma non catturano facilmente le dipendenze globali: un pixel in alto a sinistra non ``vede'' direttamente un pixel in basso a destra senza molti strati intermedi. Il ViT usa l'attenzione per creare \textbf{dipendenze globali fin dal primo strato}: ogni patch pu\`o interagire con qualsiasi altra patch.

\subsection{Pipeline completa}

\begin{enumerate}
    \item \textbf{Patch embedding}: l'immagine $224 \times 224 \times 3$ viene suddivisa in patch di $16 \times 16$ pixel:
    \[
        \text{Numero di patch} = \frac{224}{16} \times \frac{224}{16} = 14 \times 14 = 196
    \]

    \item \textbf{Proiezione lineare}: ogni patch di $16 \times 16 \times 3 = 768$ pixel viene proiettata in un vettore di $D = 768$ dimensioni tramite una matrice $\mathbf{E} \in \mathbb{R}^{768 \times 768}$:
    \[
        \mathbf{z}_i = \text{flatten}(\text{patch}_i) \cdot \mathbf{E} \in \mathbb{R}^{768}
    \]

    \item \textbf{Token CLS}: si aggiunge un token speciale $\mathbf{x}_{\text{class}} \in \mathbb{R}^{768}$ (parametro apprendibile) all'inizio della sequenza. Dopo il Transformer, questo token conterr\`a una rappresentazione globale dell'intera immagine, usata per la classificazione.

    \item \textbf{Positional encoding} (appreso): si sommano embedding posizionali apprendibili $\mathbf{E}_{\text{pos}} \in \mathbb{R}^{197 \times 768}$ (197 = 196 patch + 1 CLS token).

    \item \textbf{Transformer Encoder}: $L = 12$ blocchi di [Multi-Head Self-Attention + FFN], identici al Transformer standard (Appendice~\ref{app:transformer}).

    \item \textbf{Classification head}: un MLP (2 strati) applicato all'output del token CLS:
    \[
        \hat{y} = \text{MLP}(\text{LN}(\mathbf{z}_{\text{CLS}}^{(L)}))
    \]
\end{enumerate}

\begin{formulabox}[Formula compatta del ViT]
\[
    \mathbf{z}_0 = [\mathbf{x}_{\text{class}};\; \mathbf{x}_1 \mathbf{E};\; \mathbf{x}_2 \mathbf{E};\; \ldots;\; \mathbf{x}_P \mathbf{E}] + \mathbf{E}_{\text{pos}} \in \mathbb{R}^{(P+1) \times D}
\]
dove $P = 196$ \`e il numero di patch, $\mathbf{x}_i$ \`e la patch $i$ appiattita, $\mathbf{E} \in \mathbb{R}^{(16^2 \cdot 3) \times D}$ \`e la matrice di proiezione, e $\mathbf{E}_{\text{pos}} \in \mathbb{R}^{(P+1) \times D}$ sono i positional embedding appresi.
\end{formulabox}


\begin{examplebox}[Patch embedding: esempio numerico con un'immagine $4 \times 4$]
Per semplicit\`a, consideriamo un'immagine in scala di grigi $4 \times 4 \times 1$ con patch di $2 \times 2$:

\[
\text{Immagine} = \begin{pmatrix}
10 & 20 & 30 & 40 \\
50 & 60 & 70 & 80 \\
90 & 100 & 110 & 120 \\
130 & 140 & 150 & 160
\end{pmatrix}
\]

\textbf{Passo 1 --- Suddivisione in patch} $2 \times 2$: otteniamo $\frac{4}{2} \times \frac{4}{2} = 4$ patch:
\[
P_1 = \begin{pmatrix} 10 & 20 \\ 50 & 60 \end{pmatrix}, \quad
P_2 = \begin{pmatrix} 30 & 40 \\ 70 & 80 \end{pmatrix}, \quad
P_3 = \begin{pmatrix} 90 & 100 \\ 130 & 140 \end{pmatrix}, \quad
P_4 = \begin{pmatrix} 110 & 120 \\ 150 & 160 \end{pmatrix}
\]

\textbf{Passo 2 --- Flatten}: ogni patch $2 \times 2$ viene appiattita in un vettore di $2 \times 2 = 4$ elementi:
\[
\mathbf{x}_1 = (10, 20, 50, 60), \quad \mathbf{x}_2 = (30, 40, 70, 80), \quad \ldots
\]

\textbf{Passo 3 --- Proiezione lineare}: moltiplichiamo per $\mathbf{E} \in \mathbb{R}^{4 \times 3}$ (proiettiamo da 4 a $D = 3$). Con:
\[
\mathbf{E} = \begin{pmatrix} 0.1 & 0.0 & -0.1 \\ 0.0 & 0.1 & 0.1 \\ -0.1 & 0.0 & 0.1 \\ 0.0 & -0.1 & 0.0 \end{pmatrix}
\]
\begin{align*}
    \mathbf{z}_1 &= \mathbf{x}_1 \cdot \mathbf{E} = (10 \cdot 0.1 + 20 \cdot 0 + 50 \cdot (-0.1) + 60 \cdot 0, \;\ldots) = (-4.0, -4.0, 4.0)
\end{align*}

\textbf{Passo 4 --- Token CLS + positional encoding}: si prepende un token apprendibile $\mathbf{x}_{\text{CLS}}$ e si sommano i positional embedding. La sequenza finale ha 5 token (1 CLS + 4 patch), ciascuno un vettore di dimensione $D = 3$.

Nel ViT reale: 196 patch da 768 dimensioni ciascuna $\to$ sequenza $197 \times 768$ per il Transformer encoder.
\end{examplebox}


% TikZ: ViT pipeline
\begin{figure}[H]
\centering
\begin{tikzpicture}[
    block/.style={draw, rounded corners=3pt, fill=#1, minimum width=2.0cm, minimum height=0.7cm, font=\footnotesize},
    block/.default={orange!15},
    patchbox/.style={draw, fill=#1, minimum width=0.5cm, minimum height=0.5cm},
    arrowstyle/.style={-{Stealth[length=2.5mm]}, thick},
    >=Stealth,
    scale=0.85, transform shape
]

% Image
\node[draw, fill=red!10, minimum width=2.5cm, minimum height=2.5cm] (img) at (0, 0) {};
\node[font=\footnotesize] at (0, -1.8) {$224 \times 224$};

% Grid on image
\foreach \i in {0,1,2,3} {
    \foreach \j in {0,1,2,3} {
        \draw[gray, thin] (-1.2 + \i*0.6, -1.2 + \j*0.6) rectangle (-0.6 + \i*0.6, -0.6 + \j*0.6);
    }
}

% Arrow
\draw[arrowstyle] (1.5, 0) -- (3, 0) node[midway, above, font=\scriptsize] {Patch $16{\times}16$};

% Patch tokens
\foreach \i/\c in {0/red!20, 1/blue!20, 2/green!20, 3/yellow!20, 4/purple!15} {
    \node[patchbox=\c] (p\i) at (3.8 + \i*0.7, 0) {};
}
\node[font=\scriptsize] at (5.2, -0.5) {$\cdots$};
\node[patchbox=orange!20] (pN) at (6.5, 0) {};

% CLS token
\node[patchbox=cyan!30, minimum width=0.5cm] (cls) at (3.1, 0) {};
\node[font=\tiny] at (3.1, 0.4) {CLS};

% Linear projection
\draw[arrowstyle] (7.2, 0) -- (8, 0) node[midway, above, font=\scriptsize] {$\mathbf{E}$};

% Positional embedding
\node[block=blue!10, minimum width=2cm] (pe) at (9.5, -1.2) {$+\;\mathbf{E}_{\text{pos}}$};
\draw[arrowstyle] (9.5, -0.5) -- (pe);

% Transformer encoder
\node[block=orange!20, minimum width=3cm, minimum height=1.5cm] (transf) at (9.5, 1.5) {\begin{tabular}{c} Transformer \\ Encoder $\times L$ \end{tabular}};
\draw[arrowstyle] (pe.east) -| (11.5, 0) |- (transf.east);

% MLP head
\node[block=green!15, minimum width=2cm] (mlp) at (13.5, 1.5) {MLP Head};
\draw[arrowstyle] (transf.east) -- (mlp.west);

% Output
\node[block=purple!15, minimum width=1.5cm] (out) at (16, 1.5) {Classe};
\draw[arrowstyle] (mlp) -- (out);

% Annotation
\node[font=\scriptsize, text=gray] at (13.5, 0.7) {Solo CLS token};

\end{tikzpicture}
\caption{Pipeline del Vision Transformer (ViT). L'immagine viene suddivisa in patch $16 \times 16$, proiettate linearmente, arricchite con positional encoding, e elaborate da un Transformer encoder.}
\label{fig:vit_pipeline}
\end{figure}


\subsection{ViT vs CNN: confronto dettagliato}

\begin{center}
\begin{tabular}{@{}p{3cm}p{5cm}p{5cm}@{}}
\toprule
\textbf{Aspetto} & \textbf{CNN (es.\ ResNet)} & \textbf{ViT} \\
\midrule
Campo recettivo & Locale (kernel $3{\times}3$), cresce con la profondit\`a & Globale fin dal primo strato \\
Inductive bias & Localit\`a, invarianza traslazionale & Quasi nessuno (solo patch 2D) \\
Dati necessari & Funziona bene con pochi dati & Richiede grandi dataset (JFT-300M) \\
Scalabilit\`a & Saturazione con pi\`u dati & Migliora continuamente \\
Complessit\`a & $\mathcal{O}(M \cdot k^2 \cdot C)$ per strato & $\mathcal{O}(P^2 \cdot D)$ per strato \\
Prestazioni ImageNet & ResNet-50: 77.6\% & ViT-B/16: 77.9\% (con pretraining) \\
\bottomrule
\end{tabular}
\end{center}

\begin{notabox}[title={ViT, Perceiver e inductive bias}]
Il ViT riduce l'inductive bias rispetto alle CNN, ma ne mantiene ancora uno: la suddivisione in patch impone una griglia regolare 2D. Il \textbf{Perceiver} elimina anche questo: opera su singoli pixel (o qualsiasi altra unit\`a elementare), senza assumere alcuna struttura spaziale. Questo spiega perch\'e il Perceiver \`e completamente invariante alla permutazione (\hyperref[sec:imagenet_perm]{Sezione~10.2}), mentre il ViT perde 15 punti su ImageNet permutato.
\end{notabox}

\begin{examplebox}[Confronto dimensioni: ViT vs Perceiver]
Su un'immagine $224 \times 224$:
\begin{itemize}
    \item \textbf{ViT}: 196 patch $\to$ matrice attenzione $196 \times 196 = 38{,}416$ elementi.
    \item \textbf{Perceiver}: 50{,}176 pixel, matrice cross-attention $512 \times 50{,}176 \approx 25.7$M elementi (pi\`u grande, ma comunque lineare in $M$ per $N$ fisso).
    \item \textbf{Transformer diretto su pixel}: matrice $50{,}176 \times 50{,}176 \approx 2.5 \times 10^9$ elementi (impossibile).
\end{itemize}
Il ViT aggira il problema riducendo la risoluzione; il Perceiver lo risolve strutturalmente con il latent bottleneck.
\end{examplebox}




\end{document}
